\documentclass[10pt,openany]{book}
% One editorial system for the five existing Studium collections.
% Load after \documentclass, before book-specific configuration.
\usepackage[utf8]{inputenc}
\usepackage[T1]{fontenc}
\usepackage{lmodern}
\DeclareUnicodeCharacter{2212}{\ensuremath{-}}
\DeclareUnicodeCharacter{2192}{\ensuremath{\to}}
\DeclareUnicodeCharacter{03B1}{\ensuremath{\alpha}}
\usepackage[spanish,es-nodecimaldot]{babel}
\spanishplainpercent
\AtBeginDocument{\spanishdeactivate{"~<>}}
\usepackage[a4paper,margin=2.5cm,headheight=14pt]{geometry}
\usepackage{amsmath,amssymb,mathtools,booktabs,array,longtable,graphicx,xcolor,tikz,enumitem,microtype,titlesec,fancyhdr,listings,etoolbox}
\usetikzlibrary{arrows.meta,calc,positioning,patterns}
\usepackage[breakable,skins]{tcolorbox}
\usepackage[labelfont=bf]{caption}
\usepackage[hyphens]{url}
\usepackage[hidelinks,unicode,hypertexnames=false]{hyperref}
\hypersetup{pdfauthor={Daniel Miguel Tejedor}}
\addto\captionsspanish{\renewcommand{\contentsname}{Índice general}\renewcommand{\tablename}{Cuadro}}
\AtBeginDocument{\normalfont\DeclareFontShape{T1}{lmr}{bx}{sc}{<->ssub*lmr/m/sc}{}}
\definecolor{studiumInk}{HTML}{142C40}
\definecolor{studiumRule}{HTML}{142C40}
\definecolor{studiumBlue}{RGB}{24,74,130}
\definecolor{studiumCyan}{RGB}{0,184,205}
\definecolor{studiumLight}{RGB}{232,240,250}
\definecolor{studiumTipBack}{RGB}{232,244,236}
\definecolor{studiumTipFrame}{RGB}{27,107,58}
\definecolor{studiumDefBack}{RGB}{232,240,250}
\definecolor{studiumDefFrame}{RGB}{24,74,130}
\definecolor{studiumWorkedBack}{RGB}{255,244,230}
\definecolor{studiumWorkedFrame}{RGB}{138,78,12}
\definecolor{studiumCheckBack}{RGB}{243,236,248}
\definecolor{studiumCheckFrame}{RGB}{96,48,130}
\definecolor{studiumCodeBack}{RGB}{245,245,242}
\definecolor{studiumCodeRule}{RGB}{48,48,52}
\titleformat{\chapter}[display]{\normalfont\filright}{\large\mdseries\scshape\chaptertitlename\ \thechapter}{1ex}{\normalfont\huge\bfseries}
\titlespacing*{\chapter}{0pt}{2.5ex plus 1ex minus .2ex}{2.3ex}
\titleformat{name=\chapter,numberless}[display]{\normalfont\filright}{}{0pt}{\normalfont\huge\bfseries}
\titleformat{\section}{\normalfont\Large\bfseries}{\thesection}{.8em}{}
\titleformat{\subsection}{\normalfont\large\bfseries}{\thesubsection}{.7em}{}
\titleformat{\subsubsection}{\normalfont\normalsize\bfseries}{\thesubsubsection}{.7em}{}
\makeatletter
\renewcommand*\l@section{\@dottedtocline{1}{1.5em}{3.5em}}
\@openrightfalse
\let\cleardoublepage\clearpage
\makeatother
\setcounter{secnumdepth}{2}\setcounter{tocdepth}{1}
\setlength{\parindent}{1.5em}\setlength{\parskip}{3pt plus .5pt}
\setlength{\emergencystretch}{2em}\raggedbottom
\setlist{itemsep=2pt,topsep=4pt,leftmargin=*}
\setlist[description]{style=nextline,leftmargin=0pt,labelwidth=0pt,labelindent=0pt}
\allowdisplaybreaks[2]
\newcounter{studiumfigure}
\makeatletter\@addtoreset{studiumfigure}{chapter}\makeatother
\renewcommand{\thestudiumfigure}{\thechapter.\arabic{studiumfigure}}
\newcommand{\studiumcaption}[2]{\refstepcounter{studiumfigure}\label{#1}\addcontentsline{lof}{figure}{\protect\numberline{\thestudiumfigure}#2}{\small\textbf{Figura~\thestudiumfigure.} #2\par}}
\newcommand{\Raire}{287.05287}
\newcommand{\gzero}{9.80665}
\tcbset{studiumcommon/.style={breakable,boxrule=.65pt,arc=.7pt,left=1.7mm,right=1.7mm,top=1.3mm,bottom=1.3mm,before skip=10pt,after skip=10pt,fonttitle=\bfseries,coltitle=white,coltext=black},
 studiumtip/.style={studiumcommon,colback=studiumTipBack,colframe=studiumTipFrame,colbacktitle=studiumTipFrame},
 studiumdef/.style={studiumcommon,colback=studiumDefBack,colframe=studiumDefFrame,colbacktitle=studiumDefFrame},
 studiumworked/.style={studiumcommon,colback=studiumWorkedBack,colframe=studiumWorkedFrame,colbacktitle=studiumWorkedFrame},
 studiumcheck/.style={studiumcommon,colback=studiumCheckBack,colframe=studiumCheckFrame,colbacktitle=studiumCheckFrame}}
\newtcolorbox{studiumworked}[1]{studiumworked,title={#1}}
\newcommand{\studiumchapterguide}[1]{\par\noindent\textcolor{studiumInk}{\textbf{Orientación del capítulo.}}\enspace\textit{#1}\par\medskip}
\newcommand{\objetivos}[1]{\studiumchapterguide{#1}}
\newcommand{\caja}[2]{\begin{tcolorbox}[studiumdef,title={#1}]#2\end{tcolorbox}}
\newcommand{\resumen}[1]{\begin{tcolorbox}[studiumcheck,title={Síntesis del capítulo}]#1\end{tcolorbox}}
\newcommand{\studiumchaptercheck}{\begin{tcolorbox}[studiumcheck,title={Comprobación final}]
\begin{enumerate}
\item Explica el modelo del capítulo y las hipótesis que permiten utilizarlo.
\item Reproduce un desarrollo o un ejemplo sin consultar su solución y comprueba unidades, signos y condiciones de validez.
\item Resuelve las preguntas o ejercicios disponibles y distingue el resultado calculado de su interpretación física.
\end{enumerate}\end{tcolorbox}}
\newcommand{\figura}[2]{\begin{figure}[tbp]\centering\includegraphics[width=.83\linewidth]{#1.pdf}\caption{#2}\label{fig:#1}\end{figure}\noindent La figura~\ref{fig:#1} ofrece una representación del modelo discutido.}
\newsavebox{\studiumMathBox}
\newcommand{\fitmath}[1]{\sbox{\studiumMathBox}{$\displaystyle #1$}\ifdim\wd\studiumMathBox>\linewidth\resizebox{\linewidth}{!}{\usebox{\studiumMathBox}}\else\usebox{\studiumMathBox}\fi}
\pagestyle{fancy}\fancyhf{}
\fancyhead[LE]{\footnotesize\thepage\hspace{1.5em}\mdseries\scshape\nouppercase{\leftmark}}
\fancyhead[RO]{\footnotesize\mdseries\scshape\nouppercase{\leftmark}\hspace{1.5em}\thepage}
\renewcommand{\headrulewidth}{.4pt}\renewcommand{\footrulewidth}{0pt}
\fancypagestyle{plain}{\fancyhf{}\fancyfoot[C]{\small\thepage}\renewcommand{\headrulewidth}{0pt}}
\lstdefinestyle{studium}{basicstyle=\ttfamily\small,columns=fullflexible,keepspaces=true,showstringspaces=false,breaklines=true,breakatwhitespace=false,frame=leftline,framerule=1.15pt,rulecolor=\color{studiumCodeRule},backgroundcolor=\color{studiumCodeBack},xleftmargin=1.15em,framexleftmargin=.75em,aboveskip=.9em,belowskip=.7em,tabsize=4,keywordstyle=\ttfamily\bfseries,commentstyle=\ttfamily\itshape,stringstyle=\ttfamily}
\lstset{style=studium}
% Shared cover: every letter remains editable vector text; photographs keep original pixels.
\providecommand{\studiumcoverphoto}{}
\providecommand{\studiumcoveredition}{Primera edición}
\providecommand{\studiumcovercategory}{MANUAL UNIVERSITARIO}
\providecommand{\studiumphotocredit}{Crédito y procedencia en el proyecto del libro.}
\providecommand{\studiumphototop}{0}
\newsavebox{\studiumPhotoBox}
\newcommand{\studiumcover}[3]{%
\hypersetup{pageanchor=false}\begin{titlepage}\thispagestyle{empty}
\begin{tikzpicture}[remember picture,overlay,x=1pt,y=-1pt]
\begin{scope}[shift={(current page.north west)}]
\fill[studiumInk] (0,0) rectangle (595.276,841.89);
\begin{scope}
\clip (0,89.5) rectangle (595.276,585.5);
\sbox{\studiumPhotoBox}{\includegraphics[width=\paperwidth]{\studiumcoverphoto}}%
\ifdim\ht\studiumPhotoBox<496pt
 \node[anchor=north,inner sep=0pt] at (297.638,89.5){\includegraphics[height=496pt]{\studiumcoverphoto}};
\else
 \node[anchor=north,inner sep=0pt] at (297.638,\studiumphototop){\usebox{\studiumPhotoBox}};
\fi
\end{scope}
\fill[white] (0,0) rectangle (595.276,89.5);
\node[anchor=west,text=studiumInk,font=\sffamily\bfseries\fontsize{18}{20}\selectfont,inner sep=0pt] at (51.223,38.5){STUDIUM};
\fill[studiumCyan] (51.223,61.5) rectangle (80.8,63.7);
\node[anchor=west,text=studiumInk,font=\sffamily\fontsize{9}{11}\selectfont,inner sep=0pt] at (89.49,65.5){BIBLIOTECA DE INGENIERÍA};
\node[anchor=west,text=studiumCyan,font=\sffamily\bfseries\fontsize{10}{12}\selectfont,inner sep=0pt] at (51.223,608){#3};
\node[anchor=north west,text=white,text width=493pt,font=\rmfamily\bfseries\fontsize{32}{35}\selectfont,inner sep=0pt] at (51.223,632){#1};
\node[anchor=north west,text=white!82!studiumInk,text width=493pt,font=\sffamily\fontsize{11}{14}\selectfont,inner sep=0pt] at (51.223,722){#2};
\node[anchor=west,text=white,font=\sffamily\bfseries\fontsize{10}{12}\selectfont,inner sep=0pt] at (51.223,763){GRADO EN INGENIERÍA AEROESPACIAL};
\node[anchor=west,text=white!82!studiumInk,font=\sffamily\fontsize{10}{12}\selectfont,inner sep=0pt] at (51.223,778){Universidad de León};
\draw[white!75!studiumInk,line width=.4pt] (51.223,792) -- (544.05,792);
\node[anchor=west,text=white!82!studiumInk,font=\sffamily\fontsize{10}{12}\selectfont,inner sep=0pt] at (51.223,815){2026};
\node[anchor=east,text=white,font=\sffamily\bfseries\fontsize{10}{12}\selectfont,inner sep=0pt] at (544.05,815){\studiumcoveredition};
\end{scope}\end{tikzpicture}\null\end{titlepage}\hypersetup{pageanchor=true}}
\newcommand{\portada}[2]{\studiumcover{#1}{#2}{\studiumcovercategory}}

\providecommand{\studiumvolumetype}{manual}
\newcommand{\studiumreadingguide}{%
\chapter*{Cómo utilizar este volumen}\addcontentsline{toc}{chapter}{Cómo utilizar este volumen}
\ifdefstring{\studiumvolumetype}{manual}{\noindent\textbf{Ruta de estudio.} Empieza por la orientación del capítulo y sus conceptos. Escribe las hipótesis antes de utilizar las ecuaciones y sigue los desarrollos con sus unidades. Después intenta los problemas y las preguntas de repaso sin consultar las soluciones. Cierra cada capítulo con la comprobación final.}{%
\ifdefstring{\studiumvolumetype}{problemas}{\noindent\textbf{Ruta de trabajo.} Lee el enunciado y copia sus datos con unidades. Define las incógnitas, el modelo y las hipótesis. Intenta resolver el problema antes de consultar el desarrollo. Compara después el resultado, la interpretación física y las comprobaciones.}{\noindent\textbf{Ruta de entrenamiento.} Lee las instrucciones particulares de cada simulacro y reserva el tiempo propuesto. Trabaja primero con los enunciados. Después utiliza las soluciones indicadas en el volumen para revisar el planteamiento, los cálculos y las comprobaciones.}}

\noindent\textbf{Recuadros.} El azul identifica definiciones y resultados; el verde, consejos; el naranja, problemas resueltos; y el violeta, síntesis y autoevaluación. Los apartados técnicos conservan los nombres propios de cada asignatura.

\noindent\textbf{Problemas y simulacros.} Separa enunciado, datos, hipótesis, modelo, desarrollo, resultado y comprobación. En los simulacros trabaja primero con los enunciados y consulta después las soluciones. Sus tiempos y puntuaciones son propuestas de entrenamiento cuando así se indican, sin sustituir las instrucciones oficiales del curso.

\noindent\textbf{Consulta.} El índice general permite localizar los temas. La notación, los formularios, las soluciones y la bibliografía mantienen sus apartados al final del manual cuando forman parte de él.
\section*{Créditos y alcance editorial}\addcontentsline{toc}{section}{Créditos y alcance editorial}
\noindent Daniel Miguel Tejedor. Colección Studium. Edición de estudio independiente para el Grado en Ingeniería Aeroespacial, Universidad de León.

\noindent La compilación y la unificación editorial no añaden una aprobación institucional ni una nueva verificación científica. Cada libro conserva las hipótesis, límites y comprobaciones de sus modelos; los informes registran su alcance y sus pendientes.

\noindent\textbf{Imagen de portada.} \studiumphotocredit
}

% Book-specific metadata; all visual rules live in shared/latex/.
\renewcommand{\studiumcoverphoto}{../figures/cover.jpeg}
\renewcommand{\studiumcoveredition}{Primera edición}
\renewcommand{\studiumphotocredit}{Fotografía conservada del PDF de clase aportado por el usuario. La identificación y el crédito de origen siguen pendientes.}
\renewcommand{\studiumphototop}{0}
\graphicspath{{../figures/}}

\hypersetup{pdftitle={Turbomáquinas}}
\renewcommand{\studiumvolumetype}{manual}
\begin{document}
\studiumcover{TURBOMÁQUINAS}{Fundamentos, análisis y diseño de turbomáquinas aeronáuticas}{MANUAL UNIVERSITARIO}
\frontmatter
\studiumreadingguide
\chapter*{Presentación y alcance}
\addcontentsline{toc}{chapter}{Presentación y alcance}
Este manual de estudio desarrolla el programa de Turbomáquinas de cuarto curso del Grado en Ingeniería Aeroespacial de la Universidad de León. Se ha contrastado con la guía docente 2026/2027 de la asignatura 0710340, aportada por el usuario y leída en sus cinco páginas, y con los siete documentos docentes del archivo de Ágora. No constituye una publicación oficial de la Universidad ni implica respaldo de su profesorado.

La guía, p. 2, confirma los cinco bloques que estructuran las partes I a V. La parte VI reúne la práctica y los simulacros originales. La revisión termodinámica inicial es instrumental. El modelo dinámico del capítulo 10 y la aproximación modal del capítulo 13 son ampliaciones expresamente identificadas.

\begin{longtable}{p{.16\linewidth}p{.50\linewidth}p{.20\linewidth}}
\textbf{Bloque} & \textbf{Contenido oficial, guía 2026/2027, p. 2} & \textbf{Capítulos}\\\hline
I & Generalidades y teorema fundamental de las turbomáquinas & 1--3\\
II & Turbinas de acción y de reacción, estudio bidimensional & 4--6\\
III & Triángulos de velocidad y problema del arranque en compresores axiales & 7--10\\
IV & Equilibrio radial y esfuerzos en turbomáquinas axiales & 11--14\\
V & Compresores radiales & 15--17\\
Práctica & Desarrollo del trabajo de representación a distintos radios según el enunciado aportado & 18\\
Entrenamiento & Simulacros originales, sin atribución al profesorado & 19\\
\end{longtable}
La guía establece 6 ECTS, primer semestre y 150 horas de dedicación total (pp. 1 y 3). Las pruebas mixtas tienen un peso del 70 por ciento y los trabajos del 30 por ciento (p. 3), con las condiciones adicionales de p. 4. Los tiempos, rúbricas y datos de los simulacros de este manual son propuestas de entrenamiento, no reglas oficiales de examen.

\section*{Cómo estudiar y comprobar}
Lea primero el convenio de signos y el cambio de referencia del capítulo 2. En cada problema, declare estaciones y distinga magnitudes estáticas, totales absolutas y totales relativas. Los ejercicios resueltos muestran modelo, hipótesis, datos con unidades, desarrollo y una interpretación. Los propuestos son casos de transferencia con otros datos; comparten familias físicas de los resueltos y exigen predecir tendencias antes de consultar la solución. Hay 72 casos resueltos, 72 ejercicios de transferencia y seis problemas integrados de examen. Son 78 familias originales de cálculo con 150 conjuntos de datos, no 150 familias independientes.

Se muestran cuatro niveles: básico, intermedio, avanzado y de examen. Las soluciones de los ejercicios propuestos se reúnen en el apéndice; el cuaderno de problemas separado incluye las 150 soluciones. Los identificadores locales de los problemas quedan definidos por sus tablas de datos y no sustituyen la notación general del texto.

\section*{Estado científico y derechos}
Los pasos numéricos han sido reproducidos por Studium. Se han añadido evaluaciones con unidades, comprobaciones de dominio, identidades simbólicas y residuos seleccionados de conservación. Esto no equivale a validación experimental ni a revisión académica experta. La auditoría del framework conserva advertencias de fundamentación formal y exige revisión.

Los apuntes y los dos documentos de Wuolah se han consultado de forma privada. No se incluyen sus páginas, figuras, enunciados ni soluciones. Los problemas y figuras del manual son originales. Los documentos de Wuolah han servido únicamente para identificar competencias generales de estudio. Las notas abiertas del MIT se citan como contraste, con atribución a Manuel Martinez-Sanchez y a la base docente de Jack L. Kerrebrock; no se reproducen sus imágenes de terceros. Las notas abiertas del MIT consultadas se ofrecen bajo CC BY-NC-SA 4.0; el expediente conserva su atribución y las condiciones de uso. No se han redistribuido los PDF aportados ni se altera la titularidad de materiales de terceros.

\section*{Convención común}
El eje axial apunta aguas abajo y la dirección tangencial positiva coincide con el giro declarado. El trabajo comunicado al fluido es positivo en compresión. El trabajo extraído se expresa positivo en turbina. Alfa se mide con velocidad absoluta; beta con velocidad relativa. Ambos son ángulos algebraicos desde el eje. La letra de reacción se distingue de la constante del gas y de la solidez. Una ecuación reducida solo se aplica después de comprobar sus hipótesis.
\clearpage
\tableofcontents
\listoffigures
\mainmatter

\part{Principios básicos y análisis de turbomáquinas}

\chapter{Fundamentos, clasificación y estados del gas}
\label{chap:c01}

\studiumchapterguide{Objetivo: desarrollar fundamentos, clasificación y estados del gas, declarar hipótesis y comprobar condiciones físicas.}
% studium:end-record eyJpZCI6ICJQQVItMDAwMSIsICJzZWN0aW9uIjogImMwMSJ9

\section{Introducción y objetivos}

Una turbomáquina intercambia energía con una corriente continua mediante un rotor. Su interés aeronáutico no procede únicamente de producir un incremento o una caída de presión: permite hacerlo con caudales elevados y una masa compatible con el vuelo. El compresor prepara el aire para la combustión; la turbina transforma parte de la energía del gas en potencia para los ejes. El ventilador, aunque eleve poco la presión, pertenece a la misma familia física.
El objetivo es distinguir máquinas motoras y generadoras, formular balances con signos coherentes y convertir correctamente entre estados estáticos y totales. Se presuponen termodinámica de sistemas abiertos, cálculo diferencial y mecánica de fluidos. Al terminar, el lector debe poder explicar por qué un estator cambia la presión sin aportar trabajo y por qué un incremento de temperatura estática no demuestra, por sí solo, que haya trabajo de eje.
El contenido pertenece al bloque I del programa aportado. La revisión termodinámica es instrumental: no constituye un tema adicional atribuido a la guía docente. La clasificación de figura~\ref{fig:clasificacion} organiza tres decisiones distintas: sentido de intercambio energético, trayectoria meridiana y grado de compresibilidad.
% studium:end-record eyJpZCI6ICJQQVItMDAwMiIsICJzZWN0aW9uIjogImMwMSJ9

\par\vspace{\baselineskip}
\noindent\begin{minipage}{\linewidth}
\centering
\begingroup
\setbox0=\hbox{\begin{tikzpicture}
\node[draw] (a) at (3.5,3) {Turbomáquinas};\node[draw,align=center,text width=2.8cm,minimum height=1.35cm] (b) at (0,1) {Energía\\motora /\\generadora};\node[draw,align=center,text width=2.8cm,minimum height=1.35cm] (c) at (3.5,1) {Trayectoria\\axial / radial /\\mixta};\node[draw,align=center,text width=2.8cm,minimum height=1.35cm] (d) at (7,1) {Modelo\\compresible /\\incompresible};\draw[->] (a)--(b);\draw[->] (a)--(c);\draw[->] (a)--(d);
\end{tikzpicture}}%
\count0=\ht0
\advance\count0 by \dp0
\count2=\wd0
\count4=\dimexpr\textheight/100*38\relax
\count6=\linewidth
\ifnum\count0<1
\count0=1
\fi
\ifnum\count2<1
\count2=1
\fi
\divide\count0 by 4096
\divide\count2 by 4096
\divide\count4 by 4096
\divide\count6 by 4096
\ifnum\count0<1
\count0=1
\fi
\ifnum\count2<1
\count2=1
\fi
\ifnum\count4<1
\count4=1
\fi
\ifnum\count6<1
\count6=1
\fi
\count8=\count2
\ifnum\count0>\count4
\count8=\numexpr\count2*\count4/\count0\relax
\fi
\ifnum\count8>\count6
\resizebox{\linewidth}{!}{\begin{tikzpicture}
\node[draw] (a) at (3.5,3) {Turbomáquinas};\node[draw,align=center,text width=2.8cm,minimum height=1.35cm] (b) at (0,1) {Energía\\motora /\\generadora};\node[draw,align=center,text width=2.8cm,minimum height=1.35cm] (c) at (3.5,1) {Trayectoria\\axial / radial /\\mixta};\node[draw,align=center,text width=2.8cm,minimum height=1.35cm] (d) at (7,1) {Modelo\\compresible /\\incompresible};\draw[->] (a)--(b);\draw[->] (a)--(c);\draw[->] (a)--(d);
\end{tikzpicture}}%
\else
\ifnum\count0>\count4
\resizebox{!}{\dimexpr\textheight/100*38\relax}{\begin{tikzpicture}
\node[draw] (a) at (3.5,3) {Turbomáquinas};\node[draw,align=center,text width=2.8cm,minimum height=1.35cm] (b) at (0,1) {Energía\\motora /\\generadora};\node[draw,align=center,text width=2.8cm,minimum height=1.35cm] (c) at (3.5,1) {Trayectoria\\axial / radial /\\mixta};\node[draw,align=center,text width=2.8cm,minimum height=1.35cm] (d) at (7,1) {Modelo\\compresible /\\incompresible};\draw[->] (a)--(b);\draw[->] (a)--(c);\draw[->] (a)--(d);
\end{tikzpicture}}%
\else
\begin{tikzpicture}
\node[draw] (a) at (3.5,3) {Turbomáquinas};\node[draw,align=center,text width=2.8cm,minimum height=1.35cm] (b) at (0,1) {Energía\\motora /\\generadora};\node[draw,align=center,text width=2.8cm,minimum height=1.35cm] (c) at (3.5,1) {Trayectoria\\axial / radial /\\mixta};\node[draw,align=center,text width=2.8cm,minimum height=1.35cm] (d) at (7,1) {Modelo\\compresible /\\incompresible};\draw[->] (a)--(b);\draw[->] (a)--(c);\draw[->] (a)--(d);
\end{tikzpicture}%
\fi
\fi
\endgroup
\par\nopagebreak
\studiumcaption{fig:clasificacion}{Criterios independientes de clasificación. Esquema original.}
\end{minipage}
\par\vspace{\baselineskip}
% studium:end-record eyJpZCI6ICJGSUctMDAyMiIsICJzZWN0aW9uIjogImMwMSJ9

\section{Clasificar sin confundir los criterios}

Se denomina generadora a la máquina que entrega energía mecánica al fluido; compresores, bombas y ventiladores son ejemplos. Una máquina motora entrega energía del fluido al eje: es el caso de una turbina. En este manual el trabajo específico comunicado al fluido se escribe \(w_f\) y es positivo en un compresor. La potencia extraída útil de una turbina se define positiva como \(P_t=-\dot m w_f\), antes de descontar las pérdidas mecánicas indicadas en cada problema.
La dirección axial o radial describe la trayectoria media a través del rotor. Una componente tangencial importante no convierte una máquina axial en radial. En una máquina mixta cambia de forma significativa tanto el radio como la coordenada axial. La clasificación hidráulica o térmica se refiere a la importancia de la variación de densidad en el modelo. Un ventilador puede analizarse como incompresible a pesar de trabajar con aire; una bomba sometida a cavitación exige un modelo que incluya cambio de fase.
Una máquina de desplazamiento positivo intercambia energía mediante volúmenes que se confinan y desplazan. La turbomáquina utiliza fuerzas aerodinámicas sobre álabes y un flujo continuo. No conviene identificar continuidad del flujo con ausencia de fluctuaciones: las estelas, el paso de álabes y la interacción rotor-estator introducen periodicidad incluso en el punto nominal.
% studium:end-record eyJpZCI6ICJQQVItMDAwMyIsICJzZWN0aW9uIjogImMwMSJ9

\section{Magnitudes estáticas y totales}

La temperatura estática \(T\) caracteriza la energía interna local del gas; la velocidad absoluta \(C\) caracteriza su energía cinética por unidad de masa. La entalpía total absoluta es

\[
\fitmath{h_0=h+\frac{C^2}{2}.}
\]
Si el gas es caloríficamente perfecto, con \(c_p\) constante en el intervalo considerado, se obtiene

\[
\fitmath{T_0=T+\frac{C^2}{2c_p},\qquad c_p=\frac{\gamma R_g}{\gamma-1}.}
\]
La temperatura total representa la temperatura alcanzada al frenar adiabáticamente una corriente sin trabajo. La presión total se define mediante un frenado adicionalmente reversible hasta reposo. Su relación local con la presión estática usa una transformación isentrópica de referencia, aunque el flujo real entre dos estaciones tenga pérdidas:

\[
\fitmath{\frac{T_0}{T}=1+\frac{\gamma-1}{2}M^2,\qquad
\frac{p_0}{p}=\left(1+\frac{\gamma-1}{2}M^2\right)^{\gamma/(\gamma-1)},\qquad M=\frac{C}{\sqrt{\gamma R_g T}}.}
\]
El exponente contiene una propiedad del gas, no un rendimiento de la máquina. No se introduce una eficiencia en la relación de remanso local. La eficiencia relaciona estados de entrada y salida y sus contrapartidas ideales; el cambio estático-total en una misma estación es otra operación.
Para un gas con calor específico variable debe usarse \(h(T)\). Se resuelve \(h(T_0)-h(T)=C^2/2\) y la presión total se calcula imponiendo igualdad de entropía entre el estado local y el estado de remanso. Mantener un valor constante de \(\gamma\) a temperaturas de turbina es una simplificación educativa que se declara, no una propiedad universal del aire o de los productos de combustión.
% studium:end-record eyJpZCI6ICJQQVItMDAwNCIsICJzZWN0aW9uIjogImMwMSJ9

\section{Balances y segunda ley}

Un volumen de control estacionario con una entrada y una salida, sin cambios de cota apreciables, cumple

\[
\fitmath{q+w_f=h_{0,\mathrm{sal}}-h_{0,\mathrm{ent}}.}
\]
El signo de \(q\) es positivo hacia el fluido. Si la carcasa es adiabática, \(q=0\), pero todavía puede crecer la entropía por viscosidad, mezcla y ondas de choque. Adiabático e isentrópico solo coinciden cuando, además, no hay generación de entropía. Para gas perfecto:

\[
\fitmath{s_2-s_1=c_p\ln\frac{T_2}{T_1}-R_g\ln\frac{p_2}{p_1}.}
\]
En un estator adiabático sin trabajo, \(h_0\) permanece constante y la presión total cae si hay irreversibilidad. Puede aumentar la presión estática si la velocidad disminuye. En un rotor compresor aumenta \(h_0\); en un rotor turbina disminuye. Esta separación entre energía y degradación de energía debe comprobarse antes de atribuir cualquier cambio de presión a una pérdida.
La continuidad se escribe \(\dot m=\int_A\rho C_n\,dA\). Solo si densidad y componente normal pueden sustituirse por valores uniformes compatibles se reduce a \(\dot m=\rho C_n A\). En un anillo axial, la velocidad que atraviesa el área frontal es \(C_x\), no el módulo absoluto cuando existe remolino.
% studium:end-record eyJpZCI6ICJQQVItMDAwNSIsICJzZWN0aW9uIjogImMwMSJ9

\section{Método de resolución y admisibilidad}

Se comienza dibujando las estaciones y declarando qué magnitudes son estáticas, totales absolutas o totales relativas. Después se comprueba que \(c_p\), \(R_g\) y \(\gamma\) sean compatibles. La temperatura estática obtenida de una total debe ser positiva; una presión negativa suele revelar unidades mezcladas o una relación aplicada al revés. Todas las temperaturas en cocientes termodinámicos se expresan en kelvin.
Un ejemplo conceptual: dos conductos adiabáticos pueden tener el mismo \(T_0\) y distinta \(p_0\). El de menor presión total tiene más entropía, pero no ha perdido necesariamente energía total hacia el exterior. La energía se conserva; la posibilidad de convertirla en trabajo se degrada. Esta distinción reaparecerá en cada modelo de pérdida.
% studium:end-record eyJpZCI6ICJQQVItMDAwNiIsICJzZWN0aW9uIjogImMwMSJ9

\section{Hipótesis, errores y formulario}

Los modelos iniciales son estacionarios, unidimensionales y de gas perfecto; no describen combustión, disociación, fugas ni refrigeración salvo indicación expresa. Se emplean propiedades locales de remanso. Una medición con sonda requiere, en la práctica, correcciones de recuperación y calibración que quedan fuera de estos ejercicios.
Errores habituales: usar grados Celsius en una razón de temperaturas, introducir \(C\) en continuidad axial cuando existe \(C_\theta\), llamar pérdida al descenso de temperatura estática de una tobera y suponer que todo proceso adiabático es reversible. Autocomprobación: explique qué puede cambiar y qué debe conservarse al atravesar un difusor adiabático real.
Formulario: \(h_0=h+C^2/2\), \(a=\sqrt{\gamma R_gT}\), \(\dot m=\rho C_xA\), \(q+w_f=\Delta h_0\). Cada expresión requiere las hipótesis del apartado en que se ha deducido.
Referencias y trazabilidad: ULE, Generalidades, pp. \(2-16\), para clasificación y planteamiento del balance; MIT 16.50, lección 24, pp. \(1-4\), como contraste del intercambio dinámico. La revisión termodinámica y los ejemplos son elaboración original. Véase el registro de fuentes para distinguir material consultado y bibliografía recomendada.
% studium:end-record eyJpZCI6ICJQQVItMDAwNyIsICJzZWN0aW9uIjogImMwMSJ9

\begin{tcolorbox}[title={Desarrollos del modelo}, studiumdef, breakable]

\noindent\textbf{Modelo desarrollado: Estado estático desde remanso}

\begin{equation}
\rho=\frac{p_0}{R_g T_0}\left(1+\frac{\gamma-1}{2}\mathrm{Ma}^2\right)^{-\frac{1}{\gamma-1}}
\label{eq:EC-C01}
\end{equation}

\noindent\textit{Significado.} Densidad estática a partir de la presión y temperatura totales locales, con frenado isentrópico local. Aquí $p_0$ y $T_0$ son totales; no intercambiarlas con magnitudes estáticas.

\noindent\textbf{Hipótesis}
\begin{itemize}
\item Gas perfecto y propiedades constantes.
\item Remanso definido localmente; no se presume isentropía entre estaciones.
\end{itemize}

\noindent\textbf{Principios}
\begin{itemize}
\item Balance y definiciones físicas desarrollados en el capítulo; propiedades y dominio declarados en el problema.
\end{itemize}

\noindent\textbf{Aplicabilidad}
\begin{itemize}
\item Modelo educativo del caso C01-01-R
\end{itemize}

\noindent\textbf{Limitaciones}
\begin{itemize}
\item Modelo educativo; la aplicabilidad a una máquina real requiere revisión experta.
\end{itemize}

\noindent\textbf{Pasos}
\begin{enumerate}
\item
El remanso local fija la temperatura estática.

\[
\fitmath{T=\frac{T_0}{\frac{\mathrm{Ma}^{2} \left(\gamma - 1\right)}{2} + 1}}
\]
\item
Se aplica el frenado isentrópico local para recuperar la presión.

\[
\fitmath{p=p_0 \left(\frac{\mathrm{Ma}^{2} \left(\gamma - 1\right)}{2} + 1\right)^{- \frac{\gamma}{\gamma - 1}}}
\]
\item
La ecuación de estado usa las dos magnitudes estáticas.

\[
\fitmath{\rho=\frac{p_0 \left(\frac{\mathrm{Ma}^{2} \left(\gamma - 1\right)}{2} + 1\right)^{- \frac{1}{\gamma - 1}}}{R_g T_0}}
\]
\end{enumerate}

{\noindent\footnotesize verificada dimensionalmente.\par}
\end{tcolorbox}

\begin{tcolorbox}[title={Problema resuelto 1}, studiumworked, breakable]

\noindent\textbf{Enunciado}
\(C01-01-R \cdot\) Estado estático desde remanso. Una estación de ensayo tiene los datos totales y Mach indicados. Obtenga temperatura, presión y densidad estáticas; la incógnita final es la densidad. Datos: \(T_0=318\, \mathrm{K}\); \(p_0=154000\, \mathrm{Pa}\); \(\mathrm{Ma}=0.62\); \(\gamma=1.4\); \(R_g=287\, \mathrm{J}\,\mathrm{kg}^{-1}\,\mathrm{K}^{-1}\).

\noindent\textbf{Datos}
\begin{itemize}
\item $T_0 = 318\,\mathrm{K}$ — temperatura total
\item $p_0 = 154000\,\mathrm{Pa}$ — presión total
\item $\mathrm{Ma} = 0.62$ — Mach absoluto
\item $\gamma = 1.4$ — razón de calores
\item $R_g = 287\,\mathrm{J}\,\mathrm{kg}^{-1}\,\mathrm{K}^{-1}$ — constante del gas
\end{itemize}

\noindent\textbf{Incógnita}
$\rho$, en $\mathrm{kg\,m^{-3}}$; también las magnitudes intermedias desarrolladas.

\noindent\textbf{Modelo físico}

\noindent\textit{Estado estático desde remanso}

\[
\fitmath{\rho=\frac{p_0 \left(\frac{\mathrm{Ma}^{2} \left(\gamma - 1\right)}{2} + 1\right)^{- \frac{1}{\gamma - 1}}}{R_g T_0}}
\]

\noindent\textbf{Hipótesis}
\begin{itemize}
\item Gas perfecto y propiedades constantes.
\item Remanso definido localmente; no se presume isentropía entre estaciones.
\end{itemize}

\noindent\textbf{Desarrollo paso a paso}

\begin{enumerate}

\item El remanso local fija la temperatura estática.

\[\fitmath{T=\frac{T_0}{1+\frac{\left(\gamma-1\right)\cdot {\mathrm{Ma}}^{2}}{2}}}\]

\[\fitmath{T=\frac{(318)}{1+\frac{\left((1.4)-1\right)\cdot {(0.62)}^{2}}{2}}\approx 295.2975\,\mathrm{K}}\]

\item Se aplica el frenado isentrópico local para recuperar la presión.

\[\fitmath{p=\frac{p_0}{{\left(1+\frac{\left(\gamma-1\right)\cdot {\mathrm{Ma}}^{2}}{2}\right)}^{\frac{\gamma}{\gamma-1}}}}\]

\[\fitmath{p=\frac{(154000)}{{\left(1+\frac{\left((1.4)-1\right)\cdot {(0.62)}^{2}}{2}\right)}^{\frac{(1.4)}{(1.4)-1}}}\approx 118832.5\,\mathrm{Pa}}\]

\item La ecuación de estado usa las dos magnitudes estáticas.

\[\fitmath{\rho=\frac{p_0}{R_g\cdot T_0\cdot {\left(1+\frac{\left(\gamma-1\right)\cdot {\mathrm{Ma}}^{2}}{2}\right)}^{\frac{1}{\gamma-1}}}}\]

\[\fitmath{\rho=\frac{(154000)}{(287)\cdot (318)\cdot {\left(1+\frac{\left((1.4)-1\right)\cdot {(0.62)}^{2}}{2}\right)}^{\frac{1}{(1.4)-1}}}\approx 1.402147\,\mathrm{kg/m^{3}}}\]

\end{enumerate}


\noindent\textbf{Sustitución numérica}

\[
\fitmath{\frac{\left(154000\,\mathrm{Pa}\right) \left(\frac{\left(0.62\right)^{2} \left(\left(1.4\right) - 1\right)}{2} + 1\right)^{- \frac{1}{\left(1.4\right) - 1}}}{\left(287\,\mathrm{J}\,\mathrm{kg}^{-1}\,\mathrm{K}^{-1}\right) \left(318\,\mathrm{K}\right)}}
\]

\noindent\textbf{Resultado}

\[
\fitmath{\rho=1.402146609\times10^{0}\,\mathrm{kg}\,\mathrm{m}^{-3}}
\]

\noindent\textbf{Interpretación}
La densidad total no sustituye a la estática en el área de paso. Verifique que T y p son menores que sus valores totales.

\noindent\textbf{Comprobación de los pasos}
\begin{itemize}
\item Paso 2: $295.298$ — reproducido
\item Paso 4: $118832$ — reproducido
\item Paso 6: $1.40215$ — reproducido
\end{itemize}

\noindent\textbf{Comprobación de la respuesta} la aritmética del último paso reproduce la respuesta enunciada a la precisión con que está escrita. Es una reproducción aritmética; no valida el modelo físico ni sus hipótesis.

\noindent\textbf{Limitaciones}
\begin{itemize}
\item Modelo educativo; la aplicabilidad a una máquina real requiere revisión experta.
\end{itemize}

\noindent\textbf{Errores frecuentes}
\begin{itemize}
\item Introducir Celsius en un cociente.
\item Usar presión total y temperatura estática en la misma ecuación de estado.
\end{itemize}
% studium:end-record eyJpZCI6ICJQUkItMDAwMSIsICJzZWN0aW9uIjogImMwMSJ9
\end{tcolorbox}

\begin{tcolorbox}[title={Problema resuelto 2}, studiumworked, breakable]

\noindent\textbf{Enunciado}
\(C01-02-R \cdot\) Difusor adiabático. Un difusor reduce el módulo de velocidad. Halle el aumento de temperatura estática y la temperatura de salida. Datos: \(T_e=291\, \mathrm{K}\); \(V_e=218\, \mathrm{m}\,\mathrm{s}^{-1}\); \(V_s=94\, \mathrm{m}\,\mathrm{s}^{-1}\); \(c_p=1004.5\, \mathrm{J}\,\mathrm{kg}^{-1}\,\mathrm{K}^{-1}\).

\noindent\textbf{Datos}
\begin{itemize}
\item $T_e = 291\,\mathrm{K}$ — temperatura estática de entrada
\item $V_e = 218\,\mathrm{m}\,\mathrm{s}^{-1}$ — velocidad de entrada
\item $V_s = 94\,\mathrm{m}\,\mathrm{s}^{-1}$ — velocidad de salida
\item $c_p = 1004.5\,\mathrm{J}\,\mathrm{kg}^{-1}\,\mathrm{K}^{-1}$ — calor específico
\end{itemize}

\noindent\textbf{Incógnita}
$T_s$, en $\mathrm{K}$; también las magnitudes intermedias desarrolladas.

\noindent\textbf{Modelo físico}

\noindent\textit{Difusor adiabático}

\[
\fitmath{T_s=T_e + \frac{V_e^{2} - V_s^{2}}{2 c_p}}
\]

\noindent\textbf{Hipótesis}
\begin{itemize}
\item Estacionario, adiabático, sin trabajo ni cambio de cota.
\end{itemize}

\noindent\textbf{Desarrollo paso a paso}

\begin{enumerate}

\item Sin trabajo ni calor, la entalpía total permanece constante.

\[\fitmath{\mathrm{dT}=\frac{{\mathrm{Ve}}^{2}-{\mathrm{Vs}}^{2}}{2\cdot \mathrm{cp}}}\]

\[\fitmath{\mathrm{dT}=\frac{{(218)}^{2}-{(94)}^{2}}{2\cdot (1004.5)}\approx 19.25734\,\mathrm{K}}\]

\item Se suma el aumento estático al estado de entrada.

\[\fitmath{\mathrm{Ts}=\mathrm{Te}+\frac{{\mathrm{Ve}}^{2}-{\mathrm{Vs}}^{2}}{2\cdot \mathrm{cp}}}\]

\[\fitmath{\mathrm{Ts}=(291)+\frac{{(218)}^{2}-{(94)}^{2}}{2\cdot (1004.5)}\approx 310.2573\,\mathrm{K}}\]

\end{enumerate}


\noindent\textbf{Sustitución numérica}

\[
\fitmath{\left(291\,\mathrm{K}\right) + \frac{\left(\left(218\,\mathrm{m}\,\mathrm{s}^{-1}\right)\right)^{2} - \left(\left(94\,\mathrm{m}\,\mathrm{s}^{-1}\right)\right)^{2}}{2 \left(1004.5\,\mathrm{J}\,\mathrm{kg}^{-1}\,\mathrm{K}^{-1}\right)}}
\]

\noindent\textbf{Resultado}

\[
\fitmath{T_s=3.102573420\times10^{2}\,\mathrm{K}}
\]

\noindent\textbf{Interpretación}
La temperatura estática aumenta aunque no se suministra trabajo al gas. La presión final no puede calcularse sin una condición de pérdida.

\noindent\textbf{Comprobación de los pasos}
\begin{itemize}
\item Paso 2: $19.2573$ — reproducido
\item Paso 4: $310.257$ — reproducido
\end{itemize}

\noindent\textbf{Comprobación de la respuesta} la aritmética del último paso reproduce la respuesta enunciada a la precisión con que está escrita. Es una reproducción aritmética; no valida el modelo físico ni sus hipótesis.

\noindent\textbf{Limitaciones}
\begin{itemize}
\item Modelo educativo; la aplicabilidad a una máquina real requiere revisión experta.
\end{itemize}

\noindent\textbf{Errores frecuentes}
\begin{itemize}
\item Interpretar el aumento estático como trabajo de eje.
\end{itemize}
% studium:end-record eyJpZCI6ICJQUkItMDAwMyIsICJzZWN0aW9uIjogImMwMSJ9
\end{tcolorbox}

\begin{tcolorbox}[title={Problema resuelto 3}, studiumworked, breakable]

\noindent\textbf{Enunciado}
\(C01-03-R \cdot\) Entropía en fila fija. Una fila fija conserva temperatura total y reduce la presión total con la razón indicada. Halle la generación específica de entropía. Datos: \(\frac{p_{0s}}{p_{0e}}=0.963\); \(R_g=287\, \mathrm{J}\,\mathrm{kg}^{-1}\,\mathrm{K}^{-1}\).

\noindent\textbf{Datos}
\begin{itemize}
\item $\frac{p_{0s}}{p_{0e}} = 0.963$ — presión total salida/entrada
\item $R_g = 287\,\mathrm{J}\,\mathrm{kg}^{-1}\,\mathrm{K}^{-1}$ — constante del gas
\end{itemize}

\noindent\textbf{Incógnita}
$\Delta s$, en $\mathrm{J\,kg^{-1}\,K^{-1}}$; también las magnitudes intermedias desarrolladas.

\noindent\textbf{Modelo físico}

\noindent\textit{Entropía en fila fija}

\[
\fitmath{\Delta s=- R_g \log{\left(\frac{p_{0s}}{p_{0e}} \right)}}
\]

\noindent\textbf{Hipótesis}
\begin{itemize}
\item Gas perfecto, adiabático, sin trabajo ni mezcla de especies.
\end{itemize}

\noindent\textbf{Desarrollo paso a paso}

\begin{enumerate}

\item La igualdad de temperaturas totales elimina el término térmico de la entropía.

\[\fitmath{L=-\ln\left(\mathrm{ratio}\right)}\]

\[\fitmath{L=-\ln\left((0.963)\right)\approx 0.03770187\,}\]

\item El descenso de presión produce generación positiva.

\[\fitmath{\mathrm{ds}=-R_g\cdot \ln\left(\mathrm{ratio}\right)}\]

\[\fitmath{\mathrm{ds}=-(287)\cdot \ln\left((0.963)\right)\approx 10.82044\,\mathrm{J/(kg\,K)}}\]

\end{enumerate}


\noindent\textbf{Sustitución numérica}

\[
\fitmath{- \left(287\,\mathrm{J}\,\mathrm{kg}^{-1}\,\mathrm{K}^{-1}\right) \log{\left(\left(0.963\right) \right)}}
\]

\noindent\textbf{Resultado}

\[
\fitmath{\Delta s=1.082043588\times10^{1}\,\mathrm{J}\,\mathrm{kg}^{-1}\,\mathrm{K}^{-1}}
\]

\noindent\textbf{Interpretación}
El signo debe ser positivo para la fila adiabática pasiva. No se ha perdido energía total hacia el ambiente.

\noindent\textbf{Comprobación de los pasos}
\begin{itemize}
\item Paso 2: $0.0377019$ — reproducido
\item Paso 4: $10.8204$ — reproducido
\end{itemize}

\noindent\textbf{Comprobación de la respuesta} la aritmética del último paso reproduce la respuesta enunciada a la precisión con que está escrita. Es una reproducción aritmética; no valida el modelo físico ni sus hipótesis.

\noindent\textbf{Limitaciones}
\begin{itemize}
\item Modelo educativo; la aplicabilidad a una máquina real requiere revisión experta.
\end{itemize}

\noindent\textbf{Errores frecuentes}
\begin{itemize}
\item Invertir la razón y obtener entropía negativa.
\end{itemize}
% studium:end-record eyJpZCI6ICJQUkItMDAwNSIsICJzZWN0aW9uIjogImMwMSJ9
\end{tcolorbox}

\begin{tcolorbox}[title={Problema resuelto 4}, studiumworked, breakable]

\noindent\textbf{Enunciado}
\(C01-04-R \cdot\) Anillo con remolino. Dimensione el área frontal anular a partir de un módulo absoluto y un ángulo con el eje. Halle también la altura a radio medio fijado; esa altura es la incógnita final. Datos: \(\dot m=19.6\, \mathrm{kg}\,\mathrm{s}^{-1}\); \(\rho=1.38\, \mathrm{kg}\,\mathrm{m}^{-3}\); \(C_e=186\, \mathrm{m}\,\mathrm{s}^{-1}\); \(\alpha_e=0.5585054\, \mathrm{rad}\); \(r_m=0.29\, \mathrm{m}\).

\noindent\textbf{Datos}
\begin{itemize}
\item $\dot m = 19.6\,\mathrm{kg}\,\mathrm{s}^{-1}$ — caudal
\item $\rho = 1.38\,\mathrm{kg}\,\mathrm{m}^{-3}$ — densidad estática
\item $C_e = 186\,\mathrm{m}\,\mathrm{s}^{-1}$ — módulo absoluto
\item $\alpha_e = 0.55850536063819\,\mathrm{rad}$ — ángulo absoluto con el eje
\item $r_m = 0.29\,\mathrm{m}$ — radio medio aritmético
\end{itemize}

\noindent\textbf{Incógnita}
$b$, en $\mathrm{m}$; también las magnitudes intermedias desarrolladas.

\noindent\textbf{Modelo físico}

\noindent\textit{Anillo con remolino}

\[
\fitmath{b=\frac{\dot m}{2 \pi C_e \rho r_m \cos{\left(\alpha_e \right)}}}
\]

\noindent\textbf{Hipótesis}
\begin{itemize}
\item Flujo axial medio uniforme, bloqueo despreciable.
\end{itemize}

\noindent\textbf{Desarrollo paso a paso}

\begin{enumerate}

\item La proyección axial atraviesa el área frontal.

\[\fitmath{\mathrm{Cx}=\mathrm{Vel}\cdot \cos\left(\mathrm{ang}\right)}\]

\[\fitmath{\mathrm{Cx}=(186)\cdot \cos\left((0.5585053606)\right)\approx 157.7369\,\mathrm{m/s}}\]

\item La continuidad fija el área anular.

\[\fitmath{\mathrm{Aa}=\frac{\mathrm{md}}{\rho\cdot \mathrm{Vel}\cdot \cos\left(\mathrm{ang}\right)}}\]

\[\fitmath{\mathrm{Aa}=\frac{(19.6)}{(1.38)\cdot (186)\cdot \cos\left((0.5585053606)\right)}\approx 0.09004167\,\mathrm{m^{2}}}\]

\item La identidad geométrica A=2 pi rm b fija la altura.

\[\fitmath{b=\frac{\mathrm{md}}{\rho\cdot \mathrm{Vel}\cdot \cos\left(\mathrm{ang}\right)\cdot 2\cdot \pi\cdot \mathrm{rm}}}\]

\[\fitmath{b=\frac{(19.6)}{(1.38)\cdot (186)\cdot \cos\left((0.5585053606)\right)\cdot 2\cdot \pi\cdot (0.29)}\approx 0.04941578\,\mathrm{m}}\]

\end{enumerate}


\noindent\textbf{Sustitución numérica}

\[
\fitmath{\frac{\left(19.6\,\mathrm{kg}\,\mathrm{s}^{-1}\right)}{2 \pi \left(186\,\mathrm{m}\,\mathrm{s}^{-1}\right) \left(1.38\,\mathrm{kg}\,\mathrm{m}^{-3}\right) \left(0.29\,\mathrm{m}\right) \cos{\left(\left(0.55850536063819\,\mathrm{rad}\right) \right)}}}
\]

\noindent\textbf{Resultado}

\[
\fitmath{b=4.941578411\times10^{-2}\,\mathrm{m}}
\]

\noindent\textbf{Interpretación}
Usar el módulo de velocidad reduciría indebidamente el área. Reconstruya el caudal con el área obtenida.

\noindent\textbf{Comprobación de los pasos}
\begin{itemize}
\item Paso 2: $157.737$ — reproducido
\item Paso 4: $0.0900417$ — reproducido
\item Paso 6: $0.0494158$ — reproducido
\end{itemize}

\noindent\textbf{Comprobación de la respuesta} la aritmética del último paso reproduce la respuesta enunciada a la precisión con que está escrita. Es una reproducción aritmética; no valida el modelo físico ni sus hipótesis.

\noindent\textbf{Limitaciones}
\begin{itemize}
\item Modelo educativo; la aplicabilidad a una máquina real requiere revisión experta.
\end{itemize}

\noindent\textbf{Errores frecuentes}
\begin{itemize}
\item Usar el seno del ángulo axial.
\item Confundir altura y diámetro.
\end{itemize}
% studium:end-record eyJpZCI6ICJQUkItMDAwNyIsICJzZWN0aW9uIjogImMwMSJ9
\end{tcolorbox}

\begin{tcolorbox}[title={Ejercicios propuestos}, studiumworked, breakable]

\noindent\textbf{1 (Básico, Desarrollo)}
\(C01-01-P \cdot\) Estado estático desde remanso. Ejercicio de transferencia. Una estación de ensayo tiene los datos totales y Mach indicados. Obtenga temperatura, presión y densidad estáticas; la incógnita final es la densidad. Antes de consultar la solución, prediga el efecto de cambiar el primer dato respecto del caso resuelto de esta familia. Justifique la tendencia y explique qué restricción podría convertirse en limitante. Datos: \(T_0=338.67\, \mathrm{K}\); \(p_0=154000\, \mathrm{Pa}\); \(\mathrm{Ma}=0.62\); \(\gamma=1.4\); \(R_g=287\, \mathrm{J}\,\mathrm{kg}^{-1}\,\mathrm{K}^{-1}\).

\noindent\textbf{2 (Intermedio, Desarrollo)}
\(C01-02-P \cdot\) Difusor adiabático. Ejercicio de transferencia. Un difusor reduce el módulo de velocidad. Halle el aumento de temperatura estática y la temperatura de salida. Antes de consultar la solución, prediga el efecto de cambiar el primer dato respecto del caso resuelto de esta familia. Justifique la tendencia y explique qué restricción podría convertirse en limitante. Datos: \(T_e=309.915\, \mathrm{K}\); \(V_e=218\, \mathrm{m}\,\mathrm{s}^{-1}\); \(V_s=94\, \mathrm{m}\,\mathrm{s}^{-1}\); \(c_p=1004.5\, \mathrm{J}\,\mathrm{kg}^{-1}\,\mathrm{K}^{-1}\).

\noindent\textbf{3 (Avanzado, Desarrollo)}
\(C01-03-P \cdot\) Entropía en fila fija. Ejercicio de transferencia. Una fila fija conserva temperatura total y reduce la presión total con la razón indicada. Halle la generación específica de entropía. Antes de consultar la solución, prediga el efecto de cambiar el primer dato respecto del caso resuelto de esta familia. Justifique la tendencia y explique qué restricción podría convertirse en limitante. Datos: \(\frac{p_{0s}}{p_{0e}}=0.947\); \(R_g=287\, \mathrm{J}\,\mathrm{kg}^{-1}\,\mathrm{K}^{-1}\).

\noindent\textbf{4 (Nivel de examen, Desarrollo)}
\(C01-04-P \cdot\) Anillo con remolino. Ejercicio de transferencia. Dimensione el área frontal anular a partir de un módulo absoluto y un ángulo con el eje. Halle también la altura a radio medio fijado; esa altura es la incógnita final. Antes de consultar la solución, prediga el efecto de cambiar el primer dato respecto del caso resuelto de esta familia. Justifique la tendencia y explique qué restricción podría convertirse en limitante. Datos: \(\dot m=20.874\, \mathrm{kg}\,\mathrm{s}^{-1}\); \(\rho=1.38\, \mathrm{kg}\,\mathrm{m}^{-3}\); \(C_e=186\, \mathrm{m}\,\mathrm{s}^{-1}\); \(\alpha_e=0.5585054\, \mathrm{rad}\); \(r_m=0.29\, \mathrm{m}\).
\end{tcolorbox}


\studiumchaptercheck

\chapter{Momento angular, Euler y triángulos de velocidades}
\label{chap:c02}

\studiumchapterguide{Objetivo: desarrollar momento angular, euler y triángulos de velocidades, declarar hipótesis y comprobar condiciones físicas.}
% studium:end-record eyJpZCI6ICJQQVItMDAwOCIsICJzZWN0aW9uIjogImMwMiJ9

\section{Introducción y objetivos}

La ecuación de Euler une mecánica y termodinámica: el par ejercido sobre una corriente cambia su momento angular y, al actuar sobre una fila móvil, cambia su energía total. El objetivo es deducir esa relación, identificar su signo en ambas familias de máquinas y cerrar un triángulo sin recurrir a una figura memorizada. Pertenece al bloque I y proporciona el lenguaje de los bloques II, III y V.
Un dibujo de álabes sugiere cómo se desvía el flujo, pero no garantiza que el gas siga exactamente la tangente del metal. Incidencia y desviación se analizarán más adelante. Aquí los triángulos representan velocidades del gas medidas o calculadas en estaciones definidas, no necesariamente ángulos geométricos del perfil.
% studium:end-record eyJpZCI6ICJQQVItMDAwOSIsICJzZWN0aW9uIjogImMwMiJ9

\section{Volumen de control y par}

Se toma un volumen anular fijo que corta la corriente antes y después del rotor. El eje \(x\) coincide con el eje de giro; \(\theta\) es positivo en el sentido de la velocidad angular \(\Omega\). En régimen estacionario, el balance integral del momento angular respecto de ese eje es

\[
\fitmath{\tau_{x,\mathrm{ext}}=\int_{A_s}\rho C_n rC_\theta\,dA-\int_{A_e}\rho C_n rC_\theta\,dA.}
\]
Las presiones sobre caras axisimétricas no producen momento alrededor del eje. Si el único par relevante es el de los álabes y cada corte puede representarse por una corriente uniforme a un radio eficaz, resulta

\[
\fitmath{\tau_{r\to f}=\dot m\left(r_s C_{\theta s}-r_e C_{\theta e}\right).}
\]
La reacción sobre el rotor tiene el signo opuesto. Multiplicando por \(\Omega\) se obtiene la potencia transferida al fluido. Al dividir por el caudal y definir \(U=\Omega r\):

\[
\fitmath{w_f=U_s C_{\theta s}-U_e C_{\theta e}.}
\]
No es necesario que el gas sea perfecto ni que el proceso sea isentrópico para obtener el trabajo de Euler. Sí hacen falta un régimen medio compatible con el balance y una representación adecuada de los flujos de momento angular. Si existen fugas, sangrados o esfuerzos tangenciales adicionales, deben incluirse.
% studium:end-record eyJpZCI6ICJQQVItMDAxMCIsICJzZWN0aW9uIjogImMwMiJ9

\section{Del par a la entalpía total}


Para un rotor adiabático sin variación importante de energía potencial,



\[
\fitmath{h_{0s}-h_{0e}=w_f=\Omega(r_sC_{\theta s}-r_eC_{\theta e}).}
\]
En un compresor axial de radio aproximadamente constante, \(w_f=U(C_{\theta 2}-C_{\theta 1})>0\). La turbina se describirá con 1 antes del estator, 2 entre estator y rotor y 3 después del rotor: \(w_t=U(C_{\theta 2}-C_{\theta 3})>0\), mientras \(w_f=-w_t\). Nunca se cambia el sentido positivo de \(\theta\) para hacer positivo un resultado.
El estator puede ejercer un par sobre el fluido; su soporte transmite la reacción a la carcasa. Sin embargo, su velocidad angular es cero y no transmite potencia de eje. Por ello es posible cambiar el remolino sin modificar la entalpía total absoluta. Confundir par con potencia conduce a atribuir trabajo al estator.
% studium:end-record eyJpZCI6ICJQQVItMDAxMSIsICJzZWN0aW9uIjogImMwMiJ9

\section{Construcción vectorial y ángulos}

La velocidad periférica es \(\mathbf U=\Omega r\,\mathbf e_\theta\). La velocidad relativa al rotor cumple

\[
\fitmath{\mathbf C=\mathbf W+\mathbf U,\qquad W_x=C_x,\qquad W_\theta=C_\theta-U.}
\]
Se miden ángulos algebraicos desde el eje axial hacia \(+\theta\):

\[
\fitmath{\tan\alpha=\frac{C_\theta}{C_x},\qquad \tan\beta=\frac{W_\theta}{C_x}.}
\]
Con \(C_x>0\), los ángulos quedan entre menos noventa y noventa grados; el signo de \(\beta\) puede ser negativo. Para cálculos generales se utiliza la función de dos argumentos para conservar el cuadrante. En figura~\ref{fig:triangulo} el extremo de \(\mathbf C\) se alcanza sumando primero \(\mathbf U\) y después \(\mathbf W\). El triángulo es una suma vectorial, no una colección de longitudes que se suman escalarmente.
Algunas referencias expresan \(\beta\) como un módulo positivo hacia el lado contrario al giro o miden ángulos desde la tangente. Antes de trasladar una fórmula se reconstruyen sus componentes. Si el ángulo está referido a la tangente, seno y coseno intercambian sus papeles; si cambia el sentido positivo, cambia el signo de la componente. No existe una conversión basada exclusivamente en el nombre Hill\(-\)Peterson o Mattingly: manda el dibujo y la definición de la edición consultada.
% studium:end-record eyJpZCI6ICJQQVItMDAxMiIsICJzZWN0aW9uIjogImMwMiJ9

\par\vspace{\baselineskip}
\noindent\begin{minipage}{\linewidth}
\centering
\begingroup
\setbox0=\hbox{\begin{tikzpicture}
\draw[->,gray] (0,0)--(5,0) node[right] {$+x$};\draw[->,gray] (0,-.5)--(0,3.8) node[above] {$+\theta$};\draw[->,blue,very thick] (0,0)--(4,1.2) node[right] {$\mathbf C$};\draw[->,orange,very thick] (0,0)--(0,3.1) node[left] {$\mathbf U$};\draw[->,teal,very thick] (0,3.1)--(4,1.2) node[midway,above] {$\mathbf W$};\draw[->] (1,0) arc (0:16.7:1) node[right] {$\alpha$};\draw[dashed] (0,3.1)--(2,3.1);\draw[->] (1.2,3.1) arc (0:-25.4:1.2) node[right] {$\beta<0$};\node at (3,-.7) {$\mathbf C=\mathbf U+\mathbf W$};
\end{tikzpicture}}%
\count0=\ht0
\advance\count0 by \dp0
\count2=\wd0
\count4=\dimexpr\textheight/100*38\relax
\count6=\linewidth
\ifnum\count0<1
\count0=1
\fi
\ifnum\count2<1
\count2=1
\fi
\divide\count0 by 4096
\divide\count2 by 4096
\divide\count4 by 4096
\divide\count6 by 4096
\ifnum\count0<1
\count0=1
\fi
\ifnum\count2<1
\count2=1
\fi
\ifnum\count4<1
\count4=1
\fi
\ifnum\count6<1
\count6=1
\fi
\count8=\count2
\ifnum\count0>\count4
\count8=\numexpr\count2*\count4/\count0\relax
\fi
\ifnum\count8>\count6
\resizebox{\linewidth}{!}{\begin{tikzpicture}
\draw[->,gray] (0,0)--(5,0) node[right] {$+x$};\draw[->,gray] (0,-.5)--(0,3.8) node[above] {$+\theta$};\draw[->,blue,very thick] (0,0)--(4,1.2) node[right] {$\mathbf C$};\draw[->,orange,very thick] (0,0)--(0,3.1) node[left] {$\mathbf U$};\draw[->,teal,very thick] (0,3.1)--(4,1.2) node[midway,above] {$\mathbf W$};\draw[->] (1,0) arc (0:16.7:1) node[right] {$\alpha$};\draw[dashed] (0,3.1)--(2,3.1);\draw[->] (1.2,3.1) arc (0:-25.4:1.2) node[right] {$\beta<0$};\node at (3,-.7) {$\mathbf C=\mathbf U+\mathbf W$};
\end{tikzpicture}}%
\else
\ifnum\count0>\count4
\resizebox{!}{\dimexpr\textheight/100*38\relax}{\begin{tikzpicture}
\draw[->,gray] (0,0)--(5,0) node[right] {$+x$};\draw[->,gray] (0,-.5)--(0,3.8) node[above] {$+\theta$};\draw[->,blue,very thick] (0,0)--(4,1.2) node[right] {$\mathbf C$};\draw[->,orange,very thick] (0,0)--(0,3.1) node[left] {$\mathbf U$};\draw[->,teal,very thick] (0,3.1)--(4,1.2) node[midway,above] {$\mathbf W$};\draw[->] (1,0) arc (0:16.7:1) node[right] {$\alpha$};\draw[dashed] (0,3.1)--(2,3.1);\draw[->] (1.2,3.1) arc (0:-25.4:1.2) node[right] {$\beta<0$};\node at (3,-.7) {$\mathbf C=\mathbf U+\mathbf W$};
\end{tikzpicture}}%
\else
\begin{tikzpicture}
\draw[->,gray] (0,0)--(5,0) node[right] {$+x$};\draw[->,gray] (0,-.5)--(0,3.8) node[above] {$+\theta$};\draw[->,blue,very thick] (0,0)--(4,1.2) node[right] {$\mathbf C$};\draw[->,orange,very thick] (0,0)--(0,3.1) node[left] {$\mathbf U$};\draw[->,teal,very thick] (0,3.1)--(4,1.2) node[midway,above] {$\mathbf W$};\draw[->] (1,0) arc (0:16.7:1) node[right] {$\alpha$};\draw[dashed] (0,3.1)--(2,3.1);\draw[->] (1.2,3.1) arc (0:-25.4:1.2) node[right] {$\beta<0$};\node at (3,-.7) {$\mathbf C=\mathbf U+\mathbf W$};
\end{tikzpicture}%
\fi
\fi
\endgroup
\par\nopagebreak
\studiumcaption{fig:triangulo}{Suma vectorial local, eje x axial y theta positivo con el giro. Las flechas conservan sus orígenes y extremos.}
\end{minipage}
\par\vspace{\baselineskip}
% studium:end-record eyJpZCI6ICJGSUctMDAwMiIsICJzZWN0aW9uIjogImMwMiJ9

\section{Rotalpía y cambio de referencia}


Sustituyendo \(W^2=C^2+U^2-2UC_\theta\) se obtiene la identidad



\[
\fitmath{I=h+\frac{W^2}{2}-\frac{U^2}{2}=h_0-UC_\theta.}
\]
La conservación de \(I\) entre entrada y salida del mismo tubo de corriente de un rotor adiabático estacionario en su referencia se relaciona con Euler. Cuando el radio es constante, \(h+W^2/2\) permanece constante. Cuando el radio cambia, ese término por sí solo no se conserva: falta la contribución centrífuga \(-U^2/2\).
Las cantidades totales relativas se definen con \(W\), no con \(C\). En una misma estación, \(T_{0r}=T+W^2/(2c_p)\) y \(p_{0r}=p(T_{0r}/T)^{\gamma/(\gamma-1)}\). Una comparación de presiones totales pertenecientes a referencias distintas no es un coeficiente de pérdida. Primero deben transformarse a una misma referencia física.
% studium:end-record eyJpZCI6ICJQQVItMDAxMyIsICJzZWN0aW9uIjogImMwMiJ9

\section{Cierre, límites y errores}

Un procedimiento robusto calcula componentes, verifica \(C_x=W_x\), reconstruye módulos y solo entonces obtiene ángulos. La comprobación independiente del trabajo utiliza \(\Delta h_0\), si se conocen temperaturas totales; otra usa par y velocidad angular. Dos operaciones numéricas equivalentes no son dos experimentos independientes, pero sí detectan errores de implementación.
La ecuación media no determina por sí sola la distribución radial de velocidades, las pérdidas ni la geometría de álabes. Existen muchas máquinas con el mismo trabajo de Euler y eficiencias distintas. Errores habituales: usar velocidad relativa en el flujo de momento angular del volumen fijo, omitir el remolino de entrada en una máquina radial o suponer que \(U\) es constante a lo largo de una pala.
Formulario: \(\tau=\dot m\Delta(rC_\theta)\), \(w_f=\Delta(UC_\theta)\), \(\mathbf C=\mathbf W+\mathbf U\), \(I=h_0-UC_\theta\). Autocomprobación: una fila fija cambia \(C_\theta\); explique cómo puede transmitir par sin transmitir potencia.
Referencias: ULE, Generalidades, pp. \(15-18\); MIT 16.50, lección 24, pp. \(2-4\). Las deducciones se desarrollan de nuevo con el convenio algebraico de este manual. No se reutilizan las ilustraciones de las fuentes.
% studium:end-record eyJpZCI6ICJQQVItMDAxNCIsICJzZWN0aW9uIjogImMwMiJ9

\begin{tcolorbox}[title={Desarrollos del modelo}, studiumdef, breakable]

\noindent\textbf{Modelo desarrollado: Par de un compresor}

\begin{equation}
\dot W=\Omega\,\dot m\,r\,(C_{\theta s}-C_{\theta e})
\label{eq:EC-C02}
\end{equation}

\noindent\textit{Significado.} Potencia de eje comunicada al fluido a radio constante. $C_{\theta s}$ y $C_{\theta e}$ son componentes tangenciales absolutas; el signo indica compresor o turbina.

\noindent\textbf{Hipótesis}
\begin{itemize}
\item Estacionario y sin pares adicionales.
\item Radio constante y flujo uniforme en los cortes.
\end{itemize}

\noindent\textbf{Principios}
\begin{itemize}
\item Balance y definiciones físicas desarrollados en el capítulo; propiedades y dominio declarados en el problema.
\end{itemize}

\noindent\textbf{Aplicabilidad}
\begin{itemize}
\item Modelo educativo del caso C02-01-R
\end{itemize}

\noindent\textbf{Limitaciones}
\begin{itemize}
\item Modelo educativo; la aplicabilidad a una máquina real requiere revisión experta.
\end{itemize}

\noindent\textbf{Pasos}
\begin{enumerate}
\item
Se calcula el flujo neto de momento angular.

\[
\fitmath{M=\dot m r \left(- C_{\theta e} + C_{\theta s}\right)}
\]
\item
Se multiplica el par por la velocidad angular.

\[
\fitmath{\dot{W}=\Omega \dot m r \left(- C_{\theta e} + C_{\theta s}\right)}
\]
\end{enumerate}

{\noindent\footnotesize verificada dimensionalmente.\par}
\end{tcolorbox}

\begin{tcolorbox}[title={Problema resuelto 1}, studiumworked, breakable]

\noindent\textbf{Enunciado}
\(C02-01-R \cdot\) Par de un compresor. Calcule el par comunicado al fluido y la potencia ideal de un rotor a radio constante. La incógnita final es la potencia. Datos: \(\dot m=22\, \mathrm{kg}\,\mathrm{s}^{-1}\); \(r=0.31\, \mathrm{m}\); \(\Omega=970\, \mathrm{s}^{-1}\); \(C_{\theta e}=24\, \mathrm{m}\,\mathrm{s}^{-1}\); \(C_{\theta s}=109\, \mathrm{m}\,\mathrm{s}^{-1}\).

\noindent\textbf{Datos}
\begin{itemize}
\item $\dot m = 22\,\mathrm{kg}\,\mathrm{s}^{-1}$ — caudal
\item $r = 0.31\,\mathrm{m}$ — radio
\item $\Omega = 970\,\mathrm{s}^{-1}$ — velocidad angular en rad/s
\item $C_{\theta e} = 24\,\mathrm{m}\,\mathrm{s}^{-1}$ — remolino de entrada
\item $C_{\theta s} = 109\,\mathrm{m}\,\mathrm{s}^{-1}$ — remolino de salida
\end{itemize}

\noindent\textbf{Incógnita}
$\dot{W}$, en $\mathrm{W}$; también las magnitudes intermedias desarrolladas.

\noindent\textbf{Modelo físico}

\noindent\textit{Par de un compresor}

\[
\fitmath{\dot{W}=\Omega \dot m r \left(- C_{\theta e} + C_{\theta s}\right)}
\]

\noindent\textbf{Hipótesis}
\begin{itemize}
\item Estacionario y sin pares adicionales.
\item Radio constante y flujo uniforme en los cortes.
\end{itemize}

\noindent\textbf{Desarrollo paso a paso}

\begin{enumerate}

\item Se calcula el flujo neto de momento angular.

\[\fitmath{\tau=\mathrm{md}\cdot \mathrm{rad}\cdot \left(\mathrm{cs}-\mathrm{ce}\right)}\]

\[\fitmath{\tau=(22)\cdot (0.31)\cdot \left((109)-(24)\right)\approx 579.7\,\mathrm{N\,m}}\]

\item Se multiplica el par por la velocidad angular.

\[\fitmath{\mathrm{Pow}=\mathrm{md}\cdot \mathrm{rad}\cdot \left(\mathrm{cs}-\mathrm{ce}\right)\cdot \Omega}\]

\[\fitmath{\mathrm{Pow}=(22)\cdot (0.31)\cdot \left((109)-(24)\right)\cdot (970)\approx 562309\,\mathrm{W}}\]

\end{enumerate}


\noindent\textbf{Sustitución numérica}

\[
\fitmath{\left(970\,\mathrm{s}^{-1}\right) \left(22\,\mathrm{kg}\,\mathrm{s}^{-1}\right) \left(0.31\,\mathrm{m}\right) \left(- \left(24\,\mathrm{m}\,\mathrm{s}^{-1}\right) + \left(109\,\mathrm{m}\,\mathrm{s}^{-1}\right)\right)}
\]

\noindent\textbf{Resultado}

\[
\fitmath{\dot{W}=5.623090000\times10^{5}\,\mathrm{W}}
\]

\noindent\textbf{Interpretación}
La potencia sobre el fluido es positiva; la reacción sobre el rotor es negativa.

\noindent\textbf{Comprobación de los pasos}
\begin{itemize}
\item Paso 2: $579.7$ — reproducido
\item Paso 4: $562309$ — reproducido
\end{itemize}

\noindent\textbf{Comprobación de la respuesta} la aritmética del último paso reproduce la respuesta enunciada a la precisión con que está escrita. Es una reproducción aritmética; no valida el modelo físico ni sus hipótesis.

\noindent\textbf{Limitaciones}
\begin{itemize}
\item Modelo educativo; la aplicabilidad a una máquina real requiere revisión experta.
\end{itemize}

\noindent\textbf{Errores frecuentes}
\begin{itemize}
\item Usar velocidades relativas en el volumen fijo.
\end{itemize}
% studium:end-record eyJpZCI6ICJQUkItMDAwOSIsICJzZWN0aW9uIjogImMwMiJ9
\end{tcolorbox}

\begin{tcolorbox}[title={Problema resuelto 2}, studiumworked, breakable]

\noindent\textbf{Enunciado}
\(C02-02-R \cdot\) Triángulo y referencia relativa. Una estación tiene las componentes absolutas y velocidad periférica indicadas. Obtenga la componente relativa, el módulo relativo y la temperatura total relativa. Datos: \(C_x=158\, \mathrm{m}\,\mathrm{s}^{-1}\); \(C_\theta=76\, \mathrm{m}\,\mathrm{s}^{-1}\); \(U=284\, \mathrm{m}\,\mathrm{s}^{-1}\); \(T_e=306\, \mathrm{K}\); \(c_p=1004.5\, \mathrm{J}\,\mathrm{kg}^{-1}\,\mathrm{K}^{-1}\).

\noindent\textbf{Datos}
\begin{itemize}
\item $C_x = 158\,\mathrm{m}\,\mathrm{s}^{-1}$ — componente axial
\item $C_\theta = 76\,\mathrm{m}\,\mathrm{s}^{-1}$ — componente tangencial absoluta
\item $U = 284\,\mathrm{m}\,\mathrm{s}^{-1}$ — velocidad periférica
\item $T_e = 306\,\mathrm{K}$ — temperatura estática
\item $c_p = 1004.5\,\mathrm{J}\,\mathrm{kg}^{-1}\,\mathrm{K}^{-1}$ — calor específico
\end{itemize}

\noindent\textbf{Incógnita}
$T_{0r}$, en $\mathrm{K}$; también las magnitudes intermedias desarrolladas.

\noindent\textbf{Modelo físico}

\noindent\textit{Triángulo y referencia relativa}

\[
\fitmath{T_{0r}=T_e + \frac{C_x^{2} + \left(C_\theta - U\right)^{2}}{2 c_p}}
\]

\noindent\textbf{Hipótesis}
\begin{itemize}
\item Mismo punto y referencia de rotor con U tangencial.
\end{itemize}

\noindent\textbf{Desarrollo paso a paso}

\begin{enumerate}

\item La transformación vectorial resta U a la componente tangencial.

\[\fitmath{\mathrm{wt}=\mathrm{ct}-\mathrm{uu}}\]

\[\fitmath{\mathrm{wt}=(76)-(284)\approx -208\,\mathrm{m/s}}\]

\item El módulo usa componentes ortogonales.

\[\fitmath{\mathrm{Wr}=\sqrt{{\mathrm{cx}}^{2}+{\left(\mathrm{ct}-\mathrm{uu}\right)}^{2}}}\]

\[\fitmath{\mathrm{Wr}=\sqrt{{(158)}^{2}+{\left((76)-(284)\right)}^{2}}\approx 261.2049\,\mathrm{m/s}}\]

\item El remanso relativo se construye con W.

\[\fitmath{\mathrm{Tzr}=\mathrm{Te}+\frac{{\mathrm{cx}}^{2}+{\left(\mathrm{ct}-\mathrm{uu}\right)}^{2}}{2\cdot \mathrm{cp}}}\]

\[\fitmath{\mathrm{Tzr}=(306)+\frac{{(158)}^{2}+{\left((76)-(284)\right)}^{2}}{2\cdot (1004.5)}\approx 339.9612\,\mathrm{K}}\]

\end{enumerate}


\noindent\textbf{Sustitución numérica}

\[
\fitmath{\left(306\,\mathrm{K}\right) + \frac{\left(\left(158\,\mathrm{m}\,\mathrm{s}^{-1}\right)\right)^{2} + \left(\left(76\,\mathrm{m}\,\mathrm{s}^{-1}\right) - \left(284\,\mathrm{m}\,\mathrm{s}^{-1}\right)\right)^{2}}{2 \left(1004.5\,\mathrm{J}\,\mathrm{kg}^{-1}\,\mathrm{K}^{-1}\right)}}
\]

\noindent\textbf{Resultado}

\[
\fitmath{T_{0r}=3.399611747\times10^{2}\,\mathrm{K}}
\]

\noindent\textbf{Interpretación}
El beta relativo es negativo. La temperatura total relativa no tiene por qué coincidir con la absoluta.

\noindent\textbf{Comprobación de los pasos}
\begin{itemize}
\item Paso 2: $-208$ — reproducido
\item Paso 4: $261.205$ — reproducido
\item Paso 6: $339.961$ — reproducido
\end{itemize}

\noindent\textbf{Comprobación de la respuesta} la aritmética del último paso reproduce la respuesta enunciada a la precisión con que está escrita. Es una reproducción aritmética; no valida el modelo físico ni sus hipótesis.

\noindent\textbf{Limitaciones}
\begin{itemize}
\item Modelo educativo; la aplicabilidad a una máquina real requiere revisión experta.
\end{itemize}

\noindent\textbf{Errores frecuentes}
\begin{itemize}
\item Restar módulos de velocidades.
\end{itemize}
% studium:end-record eyJpZCI6ICJQUkItMDAxMSIsICJzZWN0aW9uIjogImMwMiJ9
\end{tcolorbox}

\begin{tcolorbox}[title={Problema resuelto 3}, studiumworked, breakable]

\noindent\textbf{Enunciado}
\(C02-03-R \cdot\) Euler con radios distintos. Un rotor radial trabaja con prerrotación no nula. Calcule el trabajo comunicado al fluido y el incremento de temperatura total. Datos: \(\Omega=1650\, \mathrm{s}^{-1}\); \(r_e=0.075\, \mathrm{m}\); \(r_s=0.19\, \mathrm{m}\); \(C_{\theta e}=31\, \mathrm{m}\,\mathrm{s}^{-1}\); \(C_{\theta s}=247\, \mathrm{m}\,\mathrm{s}^{-1}\); \(c_p=1004.5\, \mathrm{J}\,\mathrm{kg}^{-1}\,\mathrm{K}^{-1}\).

\noindent\textbf{Datos}
\begin{itemize}
\item $\Omega = 1650\,\mathrm{s}^{-1}$ — velocidad angular
\item $r_e = 0.075\,\mathrm{m}$ — radio entrada
\item $r_s = 0.19\,\mathrm{m}$ — radio salida
\item $C_{\theta e} = 31\,\mathrm{m}\,\mathrm{s}^{-1}$ — remolino entrada
\item $C_{\theta s} = 247\,\mathrm{m}\,\mathrm{s}^{-1}$ — remolino salida
\item $c_p = 1004.5\,\mathrm{J}\,\mathrm{kg}^{-1}\,\mathrm{K}^{-1}$ — calor específico
\end{itemize}

\noindent\textbf{Incógnita}
$\Delta T_0$, en $\mathrm{K}$; también las magnitudes intermedias desarrolladas.

\noindent\textbf{Modelo físico}

\noindent\textit{Euler con radios distintos}

\[
\fitmath{\Delta T_0=\frac{\Omega \left(- C_{\theta e} r_e + C_{\theta s} r_s\right)}{c_p}}
\]

\noindent\textbf{Hipótesis}
\begin{itemize}
\item Adiabático, propiedades constantes.
\item Valores medios compatibles con el flujo de momento angular.
\end{itemize}

\noindent\textbf{Desarrollo paso a paso}

\begin{enumerate}

\item La velocidad periférica de salida depende de su radio.

\[\fitmath{\mathrm{Us}=\Omega\cdot \mathrm{rs}}\]

\[\fitmath{\mathrm{Us}=(1650)\cdot (0.19)\approx 313.5\,\mathrm{m/s}}\]

\item Se conserva el término de prerrotación en Euler.

\[\fitmath{\mathrm{wc}=\Omega\cdot \left(\mathrm{rs}\cdot \mathrm{cs}-\mathrm{re}\cdot \mathrm{ce}\right)}\]

\[\fitmath{\mathrm{wc}=(1650)\cdot \left((0.19)\cdot (247)-(0.075)\cdot (31)\right)\approx 73598.25\,\mathrm{J/kg}}\]

\item Se transforma trabajo en aumento de remanso.

\[\fitmath{\mathrm{dTz}=\frac{\Omega\cdot \left(\mathrm{rs}\cdot \mathrm{cs}-\mathrm{re}\cdot \mathrm{ce}\right)}{\mathrm{cp}}}\]

\[\fitmath{\mathrm{dTz}=\frac{(1650)\cdot \left((0.19)\cdot (247)-(0.075)\cdot (31)\right)}{(1004.5)}\approx 73.26854\,\mathrm{K}}\]

\end{enumerate}


\noindent\textbf{Sustitución numérica}

\[
\fitmath{\frac{\left(1650\,\mathrm{s}^{-1}\right) \left(- \left(31\,\mathrm{m}\,\mathrm{s}^{-1}\right) \left(0.075\,\mathrm{m}\right) + \left(247\,\mathrm{m}\,\mathrm{s}^{-1}\right) \left(0.19\,\mathrm{m}\right)\right)}{\left(1004.5\,\mathrm{J}\,\mathrm{kg}^{-1}\,\mathrm{K}^{-1}\right)}}
\]

\noindent\textbf{Resultado}

\[
\fitmath{\Delta T_0=7.326854156\times10^{1}\,\mathrm{K}}
\]

\noindent\textbf{Interpretación}
Omitir el remolino de entrada sobreestima el trabajo en este caso. Compare ambas evaluaciones.

\noindent\textbf{Comprobación de los pasos}
\begin{itemize}
\item Paso 2: $313.5$ — reproducido
\item Paso 4: $73598.2$ — reproducido
\item Paso 6: $73.2685$ — reproducido
\end{itemize}

\noindent\textbf{Comprobación de la respuesta} la aritmética del último paso reproduce la respuesta enunciada a la precisión con que está escrita. Es una reproducción aritmética; no valida el modelo físico ni sus hipótesis.

\noindent\textbf{Limitaciones}
\begin{itemize}
\item Modelo educativo; la aplicabilidad a una máquina real requiere revisión experta.
\end{itemize}

\noindent\textbf{Errores frecuentes}
\begin{itemize}
\item Usar el radio de salida para ambos extremos.
\end{itemize}
% studium:end-record eyJpZCI6ICJQUkItMDAxMyIsICJzZWN0aW9uIjogImMwMiJ9
\end{tcolorbox}

\begin{tcolorbox}[title={Problema resuelto 4}, studiumworked, breakable]

\noindent\textbf{Enunciado}
\(C02-04-R \cdot\) Comprobación de rotalpía. Para un rotor se conocen estados cinemáticos de entrada y salida. Calcule el cambio de entalpía estática exigido por rotalpía. Datos: \(U_e=192\, \mathrm{m}\,\mathrm{s}^{-1}\); \(U_s=331\, \mathrm{m}\,\mathrm{s}^{-1}\); \(W_e=246\, \mathrm{m}\,\mathrm{s}^{-1}\); \(W_s=214\, \mathrm{m}\,\mathrm{s}^{-1}\).

\noindent\textbf{Datos}
\begin{itemize}
\item $U_e = 192\,\mathrm{m}\,\mathrm{s}^{-1}$ — velocidad periférica entrada
\item $U_s = 331\,\mathrm{m}\,\mathrm{s}^{-1}$ — velocidad periférica salida
\item $W_e = 246\,\mathrm{m}\,\mathrm{s}^{-1}$ — velocidad relativa entrada
\item $W_s = 214\,\mathrm{m}\,\mathrm{s}^{-1}$ — velocidad relativa salida
\end{itemize}

\noindent\textbf{Incógnita}
$\Delta h$, en $\mathrm{J\,kg^{-1}}$; también las magnitudes intermedias desarrolladas.

\noindent\textbf{Modelo físico}

\noindent\textit{Comprobación de rotalpía}

\[
\fitmath{\Delta h=- \frac{U_e^{2}}{2} + \frac{U_s^{2}}{2} + \frac{W_e^{2}}{2} - \frac{W_s^{2}}{2}}
\]

\noindent\textbf{Hipótesis}
\begin{itemize}
\item Rotor adiabático, estacionario en referencia móvil, sin cambio de cota.
\end{itemize}

\noindent\textbf{Desarrollo paso a paso}

\begin{enumerate}

\item Se calcula la contribución de cambio de radio.

\[\fitmath{\mathrm{du}=\frac{{\mathrm{us}}^{2}-{\mathrm{ue}}^{2}}{2}}\]

\[\fitmath{\mathrm{du}=\frac{{(331)}^{2}-{(192)}^{2}}{2}\approx 36348.5\,\mathrm{J/kg}}\]

\item Se evalúa por separado la variación de energía cinética relativa.

\[\fitmath{\mathrm{ref}=\frac{{\mathrm{ws}}^{2}-{\mathrm{we}}^{2}}{2}}\]

\[\fitmath{\mathrm{ref}=\frac{{(214)}^{2}-{(246)}^{2}}{2}\approx -7360\,\mathrm{J/kg}}\]

\item De I constante se despeja la variación estática.

\[\fitmath{\mathrm{dh}=\frac{{\mathrm{us}}^{2}-{\mathrm{ue}}^{2}-{\mathrm{ws}}^{2}+{\mathrm{we}}^{2}}{2}}\]

\[\fitmath{\mathrm{dh}=\frac{{(331)}^{2}-{(192)}^{2}-{(214)}^{2}+{(246)}^{2}}{2}\approx 43708.5\,\mathrm{J/kg}}\]

\end{enumerate}


\noindent\textbf{Sustitución numérica}

\[
\fitmath{- \frac{\left(\left(192\,\mathrm{m}\,\mathrm{s}^{-1}\right)\right)^{2}}{2} + \frac{\left(\left(331\,\mathrm{m}\,\mathrm{s}^{-1}\right)\right)^{2}}{2} + \frac{\left(\left(246\,\mathrm{m}\,\mathrm{s}^{-1}\right)\right)^{2}}{2} - \frac{\left(\left(214\,\mathrm{m}\,\mathrm{s}^{-1}\right)\right)^{2}}{2}}
\]

\noindent\textbf{Resultado}

\[
\fitmath{\Delta h=4.370850000\times10^{4}\,\mathrm{J}\,\mathrm{kg}^{-1}}
\]

\noindent\textbf{Interpretación}
El término centrífugo no puede omitirse cuando U cambia. Reintroduzca dh en la identidad de rotalpía y compruebe residuo cero.

\noindent\textbf{Comprobación de los pasos}
\begin{itemize}
\item Paso 2: $36348.5$ — reproducido
\item Paso 4: $-7360$ — reproducido
\item Paso 6: $43708.5$ — reproducido
\end{itemize}

\noindent\textbf{Comprobación de la respuesta} la aritmética del último paso reproduce la respuesta enunciada a la precisión con que está escrita. Es una reproducción aritmética; no valida el modelo físico ni sus hipótesis.

\noindent\textbf{Limitaciones}
\begin{itemize}
\item Modelo educativo; la aplicabilidad a una máquina real requiere revisión experta.
\end{itemize}

\noindent\textbf{Errores frecuentes}
\begin{itemize}
\item Conservar h+W al cuadrado/2 a radio variable.
\end{itemize}
% studium:end-record eyJpZCI6ICJQUkItMDAxNSIsICJzZWN0aW9uIjogImMwMiJ9
\end{tcolorbox}

\begin{tcolorbox}[title={Ejercicios propuestos}, studiumworked, breakable]

\noindent\textbf{1 (Básico, Desarrollo)}
\(C02-01-P \cdot\) Par de un compresor. Ejercicio de transferencia. Calcule el par comunicado al fluido y la potencia ideal de un rotor a radio constante. La incógnita final es la potencia. Antes de consultar la solución, prediga el efecto de cambiar el primer dato respecto del caso resuelto de esta familia. Justifique la tendencia y explique qué restricción podría convertirse en limitante. Datos: \(\dot m=23.43\, \mathrm{kg}\,\mathrm{s}^{-1}\); \(r=0.31\, \mathrm{m}\); \(\Omega=970\, \mathrm{s}^{-1}\); \(C_{\theta e}=24\, \mathrm{m}\,\mathrm{s}^{-1}\); \(C_{\theta s}=109\, \mathrm{m}\,\mathrm{s}^{-1}\).

\noindent\textbf{2 (Intermedio, Desarrollo)}
\(C02-02-P \cdot\) Triángulo y referencia relativa. Ejercicio de transferencia. Una estación tiene las componentes absolutas y velocidad periférica indicadas. Obtenga la componente relativa, el módulo relativo y la temperatura total relativa. Antes de consultar la solución, prediga el efecto de cambiar el primer dato respecto del caso resuelto de esta familia. Justifique la tendencia y explique qué restricción podría convertirse en limitante. Datos: \(C_x=168.27\, \mathrm{m}\,\mathrm{s}^{-1}\); \(C_\theta=76\, \mathrm{m}\,\mathrm{s}^{-1}\); \(U=284\, \mathrm{m}\,\mathrm{s}^{-1}\); \(T_e=306\, \mathrm{K}\); \(c_p=1004.5\, \mathrm{J}\,\mathrm{kg}^{-1}\,\mathrm{K}^{-1}\).

\noindent\textbf{3 (Avanzado, Desarrollo)}
\(C02-03-P \cdot\) Euler con radios distintos. Ejercicio de transferencia. Un rotor radial trabaja con prerrotación no nula. Calcule el trabajo comunicado al fluido y el incremento de temperatura total. Antes de consultar la solución, prediga el efecto de cambiar el primer dato respecto del caso resuelto de esta familia. Justifique la tendencia y explique qué restricción podría convertirse en limitante. Datos: \(\Omega=1757.25\, \mathrm{s}^{-1}\); \(r_e=0.075\, \mathrm{m}\); \(r_s=0.19\, \mathrm{m}\); \(C_{\theta e}=31\, \mathrm{m}\,\mathrm{s}^{-1}\); \(C_{\theta s}=247\, \mathrm{m}\,\mathrm{s}^{-1}\); \(c_p=1004.5\, \mathrm{J}\,\mathrm{kg}^{-1}\,\mathrm{K}^{-1}\).

\noindent\textbf{4 (Nivel de examen, Desarrollo)}
\(C02-04-P \cdot\) Comprobación de rotalpía. Ejercicio de transferencia. Para un rotor se conocen estados cinemáticos de entrada y salida. Calcule el cambio de entalpía estática exigido por rotalpía. Antes de consultar la solución, prediga el efecto de cambiar el primer dato respecto del caso resuelto de esta familia. Justifique la tendencia y explique qué restricción podría convertirse en limitante. Datos: \(U_e=204.48\, \mathrm{m}\,\mathrm{s}^{-1}\); \(U_s=331\, \mathrm{m}\,\mathrm{s}^{-1}\); \(W_e=246\, \mathrm{m}\,\mathrm{s}^{-1}\); \(W_s=214\, \mathrm{m}\,\mathrm{s}^{-1}\).
\end{tcolorbox}


\studiumchaptercheck

\chapter{Rendimientos, pérdidas y semejanza}
\label{chap:c03}

\studiumchapterguide{Objetivo: desarrollar rendimientos, pérdidas y semejanza, declarar hipótesis y comprobar condiciones físicas.}
% studium:end-record eyJpZCI6ICJQQVItMDAxNSIsICJzZWN0aW9uIjogImMwMyJ9

\section{Introducción y objetivos}

Conocer el trabajo de Euler no permite obtener directamente la relación de presiones de una máquina real. Hace falta describir la generación de entropía. Este capítulo fija las definiciones de rendimiento y enseña a relacionarlas con los balances, antes de introducir coeficientes de diseño. Es contenido del bloque I aplicado después a los bloques II, III y V.
El lector debe distinguir entre un rendimiento global, un rendimiento de escalón y una eficiencia politrópica. También debe reconocer qué parámetros deben conservarse para comparar máquinas semejantes y por qué una curva adimensional no constituye, por sí sola, un mapa experimental.
% studium:end-record eyJpZCI6ICJQQVItMDAxNiIsICJzZWN0aW9uIjogImMwMyJ9

\section{Compresión real e ideal}

Se comparan entrada y salida totales absolutas. El estado ideal \(02s\) tiene la misma presión total de salida real y la entropía total de entrada. Se define

\[
\fitmath{\eta_{c,\mathrm{tt}}=\frac{h_{02s}-h_{01}}{h_{02}-h_{01}},\qquad
\pi_c=\frac{p_{02}}{p_{01}}>1.}
\]
Para gas caloríficamente perfecto:

\[
\fitmath{T_{02}=T_{01}\left[1+\frac{\pi_c^{(\gamma-1)/\gamma}-1}{\eta_{c,\mathrm{tt}}}\right].}
\]
El trabajo ideal es menor que el real: el denominador representa el coste energético efectivo. Si se conoce el incremento de temperatura y se busca la presión, se despeja

\[
\fitmath{\pi_c=\left[1+\eta_{c,\mathrm{tt}}\frac{T_{02}-T_{01}}{T_{01}}\right]^{\gamma/(\gamma-1)}.}
\]
No se multiplica todo el incremento de presión por el rendimiento. La no linealidad de la relación isentrópica impide ese atajo.
% studium:end-record eyJpZCI6ICJQQVItMDAxNyIsICJzZWN0aW9uIjogImMwMyJ9

\section{Expansión y energía de salida}

En una turbina el trabajo positivo extraído es el descenso real de entalpía total. El rendimiento total-total se define como

\[
\fitmath{\eta_{t,\mathrm{tt}}=\frac{h_{01}-h_{03}}{h_{01}-h_{03s}},\qquad
T_{03}=T_{01}\left[1-\eta_{t,\mathrm{tt}}\left(1-\left(\frac{p_{03}}{p_{01}}\right)^{(\gamma-1)/\gamma}\right)\right].}
\]
El rendimiento total-estático compara el trabajo real con una expansión ideal hasta la presión estática de salida, sin acreditar como trabajo útil la energía cinética que abandona el escalón. Por tanto su denominador es diferente. Cuando un escalón descarga en otro, parte de esa velocidad puede aprovecharse; cuando descarga en un escape, la definición de utilidad depende del sistema.
Debe precisarse si la eficiencia incluye refrigeración, fugas y pérdidas de cojinete. Se adoptará \(P_{eje,t}=\eta_m\dot m w_t\) para una turbina y \(P_{eje,c}=\dot m w_f/\eta_m\) para un compresor. No se aplica dos veces el rendimiento mecánico si ya estaba incorporado en la potencia suministrada.
% studium:end-record eyJpZCI6ICJQQVItMDAxOCIsICJzZWN0aW9uIjogImMwMyJ9

\section{Eficiencia politrópica}

Considérese una sucesión de compresiones diferenciales. Para una transformación isentrópica infinitesimal \(dh_s=v\,\mathrm{d}p\). Con gas perfecto, su definición conduce a

\[
\fitmath{\begin{aligned}
\eta_{p,c} &= \frac{(\gamma-1)\,d\ln p_0}{\gamma\,d\ln T_0},\qquad
\frac{T_{02}}{T_{01}} \\
&= \pi_c^{(\gamma-1)/(\gamma\eta_{p,c})}.
\end{aligned}}
\]
La integración exige que \(\eta_{p,c}\) y las propiedades del gas sean constantes. Para expansión, con la convención habitual,

\[
\fitmath{\frac{T_{03}}{T_{01}}=\left(\frac{p_{03}}{p_{01}}\right)^{\eta_{p,t}(\gamma-1)/\gamma}.}
\]
En compresión, una eficiencia politrópica fija produce una eficiencia isentrópica global menor al aumentar la relación de presiones. En expansión la relación entre ambas cambia de sentido. La figura~\ref{fig:eficiencia} muestra curvas calculadas exclusivamente con estas definiciones; no representa un catálogo de máquinas.
Las relaciones de presión de los escalones se multiplican, mientras los incrementos de entalpía total se suman. Los rendimientos de escalón no se multiplican para obtener el rendimiento total. Cada escalón recibe una temperatura diferente y su trabajo ideal debe referirse al estado que realmente le corresponde.
% studium:end-record eyJpZCI6ICJQQVItMDAxOSIsICJzZWN0aW9uIjogImMwMyJ9

\par\vspace{\baselineskip}
\noindent\begin{minipage}{\linewidth}
\centering
\begingroup
\setbox0=\hbox{\begin{tikzpicture}
\draw[->,gray] (0,0)--(7,0);\draw[->,gray] (0,0)--(0,3.8);\draw (0.0,0)--(0.0,-.08) node[below,font=\scriptsize] {2};\draw (0,0.0)--(-.08,0.0) node[left,font=\scriptsize] {0.75};\draw (1.5,0)--(1.5,-.08) node[below,font=\scriptsize] {6.5};\draw (0,0.75)--(-.08,0.75) node[left,font=\scriptsize] {0.792};\draw (3.0,0)--(3.0,-.08) node[below,font=\scriptsize] {11};\draw (0,1.5)--(-.08,1.5) node[left,font=\scriptsize] {0.833};\draw (4.5,0)--(4.5,-.08) node[below,font=\scriptsize] {15.5};\draw (0,2.25)--(-.08,2.25) node[left,font=\scriptsize] {0.875};\draw (6.0,0)--(6.0,-.08) node[below,font=\scriptsize] {20};\draw (0,3.0)--(-.08,3.0) node[left,font=\scriptsize] {0.917};\draw[blue,very thick] plot coordinates {(0.00000,2.51660) (0.10000,2.47840) (0.20000,2.44455) (0.30000,2.41414) (0.40000,2.38653) (0.50000,2.36122) (0.60000,2.33785) (0.70000,2.31614) (0.80000,2.29586) (0.90000,2.27684) (1.00000,2.25891) (1.10000,2.24196) (1.20000,2.22589) (1.30000,2.21060) (1.40000,2.19602) (1.50000,2.18209) (1.60000,2.16875) (1.70000,2.15595) (1.80000,2.14365) (1.90000,2.13180) (2.00000,2.12038) (2.10000,2.10936) (2.20000,2.09871) (2.30000,2.08839) (2.40000,2.07840) (2.50000,2.06872) (2.60000,2.05931) (2.70000,2.05017) (2.80000,2.04129) (2.90000,2.03264) (3.00000,2.02422) (3.10000,2.01601) (3.20000,2.00801) (3.30000,2.00020) (3.40000,1.99257) (3.50000,1.98512) (3.60000,1.97783) (3.70000,1.97070) (3.80000,1.96373) (3.90000,1.95690) (4.00000,1.95021) (4.10000,1.94366) (4.20000,1.93723) (4.30000,1.93093) (4.40000,1.92475) (4.50000,1.91868) (4.60000,1.91272) (4.70000,1.90687) (4.80000,1.90112) (4.90000,1.89546) (5.00000,1.88991) (5.10000,1.88444) (5.20000,1.87907) (5.30000,1.87378) (5.40000,1.86857) (5.50000,1.86345) (5.60000,1.85840) (5.70000,1.85343) (5.80000,1.84854) (5.90000,1.84371) (6.00000,1.83895)};\node[below,align=center] at (3.5,-.65) {Relación de presión};\node[rotate=90] at (-1,2) {$\eta_{c,\mathrm{tt}}$};
\end{tikzpicture}}%
\count0=\ht0
\advance\count0 by \dp0
\count2=\wd0
\count4=\dimexpr\textheight/100*38\relax
\count6=\linewidth
\ifnum\count0<1
\count0=1
\fi
\ifnum\count2<1
\count2=1
\fi
\divide\count0 by 4096
\divide\count2 by 4096
\divide\count4 by 4096
\divide\count6 by 4096
\ifnum\count0<1
\count0=1
\fi
\ifnum\count2<1
\count2=1
\fi
\ifnum\count4<1
\count4=1
\fi
\ifnum\count6<1
\count6=1
\fi
\count8=\count2
\ifnum\count0>\count4
\count8=\numexpr\count2*\count4/\count0\relax
\fi
\ifnum\count8>\count6
\resizebox{\linewidth}{!}{\begin{tikzpicture}
\draw[->,gray] (0,0)--(7,0);\draw[->,gray] (0,0)--(0,3.8);\draw (0.0,0)--(0.0,-.08) node[below,font=\scriptsize] {2};\draw (0,0.0)--(-.08,0.0) node[left,font=\scriptsize] {0.75};\draw (1.5,0)--(1.5,-.08) node[below,font=\scriptsize] {6.5};\draw (0,0.75)--(-.08,0.75) node[left,font=\scriptsize] {0.792};\draw (3.0,0)--(3.0,-.08) node[below,font=\scriptsize] {11};\draw (0,1.5)--(-.08,1.5) node[left,font=\scriptsize] {0.833};\draw (4.5,0)--(4.5,-.08) node[below,font=\scriptsize] {15.5};\draw (0,2.25)--(-.08,2.25) node[left,font=\scriptsize] {0.875};\draw (6.0,0)--(6.0,-.08) node[below,font=\scriptsize] {20};\draw (0,3.0)--(-.08,3.0) node[left,font=\scriptsize] {0.917};\draw[blue,very thick] plot coordinates {(0.00000,2.51660) (0.10000,2.47840) (0.20000,2.44455) (0.30000,2.41414) (0.40000,2.38653) (0.50000,2.36122) (0.60000,2.33785) (0.70000,2.31614) (0.80000,2.29586) (0.90000,2.27684) (1.00000,2.25891) (1.10000,2.24196) (1.20000,2.22589) (1.30000,2.21060) (1.40000,2.19602) (1.50000,2.18209) (1.60000,2.16875) (1.70000,2.15595) (1.80000,2.14365) (1.90000,2.13180) (2.00000,2.12038) (2.10000,2.10936) (2.20000,2.09871) (2.30000,2.08839) (2.40000,2.07840) (2.50000,2.06872) (2.60000,2.05931) (2.70000,2.05017) (2.80000,2.04129) (2.90000,2.03264) (3.00000,2.02422) (3.10000,2.01601) (3.20000,2.00801) (3.30000,2.00020) (3.40000,1.99257) (3.50000,1.98512) (3.60000,1.97783) (3.70000,1.97070) (3.80000,1.96373) (3.90000,1.95690) (4.00000,1.95021) (4.10000,1.94366) (4.20000,1.93723) (4.30000,1.93093) (4.40000,1.92475) (4.50000,1.91868) (4.60000,1.91272) (4.70000,1.90687) (4.80000,1.90112) (4.90000,1.89546) (5.00000,1.88991) (5.10000,1.88444) (5.20000,1.87907) (5.30000,1.87378) (5.40000,1.86857) (5.50000,1.86345) (5.60000,1.85840) (5.70000,1.85343) (5.80000,1.84854) (5.90000,1.84371) (6.00000,1.83895)};\node[below,align=center] at (3.5,-.65) {Relación de presión};\node[rotate=90] at (-1,2) {$\eta_{c,\mathrm{tt}}$};
\end{tikzpicture}}%
\else
\ifnum\count0>\count4
\resizebox{!}{\dimexpr\textheight/100*38\relax}{\begin{tikzpicture}
\draw[->,gray] (0,0)--(7,0);\draw[->,gray] (0,0)--(0,3.8);\draw (0.0,0)--(0.0,-.08) node[below,font=\scriptsize] {2};\draw (0,0.0)--(-.08,0.0) node[left,font=\scriptsize] {0.75};\draw (1.5,0)--(1.5,-.08) node[below,font=\scriptsize] {6.5};\draw (0,0.75)--(-.08,0.75) node[left,font=\scriptsize] {0.792};\draw (3.0,0)--(3.0,-.08) node[below,font=\scriptsize] {11};\draw (0,1.5)--(-.08,1.5) node[left,font=\scriptsize] {0.833};\draw (4.5,0)--(4.5,-.08) node[below,font=\scriptsize] {15.5};\draw (0,2.25)--(-.08,2.25) node[left,font=\scriptsize] {0.875};\draw (6.0,0)--(6.0,-.08) node[below,font=\scriptsize] {20};\draw (0,3.0)--(-.08,3.0) node[left,font=\scriptsize] {0.917};\draw[blue,very thick] plot coordinates {(0.00000,2.51660) (0.10000,2.47840) (0.20000,2.44455) (0.30000,2.41414) (0.40000,2.38653) (0.50000,2.36122) (0.60000,2.33785) (0.70000,2.31614) (0.80000,2.29586) (0.90000,2.27684) (1.00000,2.25891) (1.10000,2.24196) (1.20000,2.22589) (1.30000,2.21060) (1.40000,2.19602) (1.50000,2.18209) (1.60000,2.16875) (1.70000,2.15595) (1.80000,2.14365) (1.90000,2.13180) (2.00000,2.12038) (2.10000,2.10936) (2.20000,2.09871) (2.30000,2.08839) (2.40000,2.07840) (2.50000,2.06872) (2.60000,2.05931) (2.70000,2.05017) (2.80000,2.04129) (2.90000,2.03264) (3.00000,2.02422) (3.10000,2.01601) (3.20000,2.00801) (3.30000,2.00020) (3.40000,1.99257) (3.50000,1.98512) (3.60000,1.97783) (3.70000,1.97070) (3.80000,1.96373) (3.90000,1.95690) (4.00000,1.95021) (4.10000,1.94366) (4.20000,1.93723) (4.30000,1.93093) (4.40000,1.92475) (4.50000,1.91868) (4.60000,1.91272) (4.70000,1.90687) (4.80000,1.90112) (4.90000,1.89546) (5.00000,1.88991) (5.10000,1.88444) (5.20000,1.87907) (5.30000,1.87378) (5.40000,1.86857) (5.50000,1.86345) (5.60000,1.85840) (5.70000,1.85343) (5.80000,1.84854) (5.90000,1.84371) (6.00000,1.83895)};\node[below,align=center] at (3.5,-.65) {Relación de presión};\node[rotate=90] at (-1,2) {$\eta_{c,\mathrm{tt}}$};
\end{tikzpicture}}%
\else
\begin{tikzpicture}
\draw[->,gray] (0,0)--(7,0);\draw[->,gray] (0,0)--(0,3.8);\draw (0.0,0)--(0.0,-.08) node[below,font=\scriptsize] {2};\draw (0,0.0)--(-.08,0.0) node[left,font=\scriptsize] {0.75};\draw (1.5,0)--(1.5,-.08) node[below,font=\scriptsize] {6.5};\draw (0,0.75)--(-.08,0.75) node[left,font=\scriptsize] {0.792};\draw (3.0,0)--(3.0,-.08) node[below,font=\scriptsize] {11};\draw (0,1.5)--(-.08,1.5) node[left,font=\scriptsize] {0.833};\draw (4.5,0)--(4.5,-.08) node[below,font=\scriptsize] {15.5};\draw (0,2.25)--(-.08,2.25) node[left,font=\scriptsize] {0.875};\draw (6.0,0)--(6.0,-.08) node[below,font=\scriptsize] {20};\draw (0,3.0)--(-.08,3.0) node[left,font=\scriptsize] {0.917};\draw[blue,very thick] plot coordinates {(0.00000,2.51660) (0.10000,2.47840) (0.20000,2.44455) (0.30000,2.41414) (0.40000,2.38653) (0.50000,2.36122) (0.60000,2.33785) (0.70000,2.31614) (0.80000,2.29586) (0.90000,2.27684) (1.00000,2.25891) (1.10000,2.24196) (1.20000,2.22589) (1.30000,2.21060) (1.40000,2.19602) (1.50000,2.18209) (1.60000,2.16875) (1.70000,2.15595) (1.80000,2.14365) (1.90000,2.13180) (2.00000,2.12038) (2.10000,2.10936) (2.20000,2.09871) (2.30000,2.08839) (2.40000,2.07840) (2.50000,2.06872) (2.60000,2.05931) (2.70000,2.05017) (2.80000,2.04129) (2.90000,2.03264) (3.00000,2.02422) (3.10000,2.01601) (3.20000,2.00801) (3.30000,2.00020) (3.40000,1.99257) (3.50000,1.98512) (3.60000,1.97783) (3.70000,1.97070) (3.80000,1.96373) (3.90000,1.95690) (4.00000,1.95021) (4.10000,1.94366) (4.20000,1.93723) (4.30000,1.93093) (4.40000,1.92475) (4.50000,1.91868) (4.60000,1.91272) (4.70000,1.90687) (4.80000,1.90112) (4.90000,1.89546) (5.00000,1.88991) (5.10000,1.88444) (5.20000,1.87907) (5.30000,1.87378) (5.40000,1.86857) (5.50000,1.86345) (5.60000,1.85840) (5.70000,1.85343) (5.80000,1.84854) (5.90000,1.84371) (6.00000,1.83895)};\node[below,align=center] at (3.5,-.65) {Relación de presión};\node[rotate=90] at (-1,2) {$\eta_{c,\mathrm{tt}}$};
\end{tikzpicture}%
\fi
\fi
\endgroup
\par\nopagebreak
\studiumcaption{fig:eficiencia}{Conversión calculada para eficiencia politrópica de compresor 0,90 y gamma 1,4. Sin datos de máquinas.}
\end{minipage}
\par\vspace{\baselineskip}
% studium:end-record eyJpZCI6ICJGSUctMDAwMyIsICJzZWN0aW9uIjogImMwMyJ9

\section{Pérdidas y variables adimensionales}

En una fila fija adiabática, con la misma entalpía total, el incremento de entropía puede escribirse \(\Delta s=-R_g\ln(p_{02}/p_{01})\). Una pérdida de presión total no equivale a una pérdida de presión estática: un difusor puede recuperar esta última mientras pierde la primera. En el rotor se debe definir un coeficiente en referencia relativa, con un denominador y unas estaciones especificados.
Para un escalón axial al radio de referencia se adoptan

\[
\fitmath{\phi=\frac{C_x}{U},\qquad \psi_c=\frac{w_f}{U^2},\qquad \psi_t=\frac{w_t}{U^2}.}
\]
Ambas cargas se expresan positivas para su familia; no se utiliza un único \(\psi\) con signo oculto. También se requieren Mach, Reynolds y relaciones geométricas. En máquinas semejantes, con densidad casi constante y efectos viscosos comparables, \(Q\) escala como \(ND^3\), el trabajo como \(N^2D^2\) y la potencia como \(\rho N^3D^5\). Esas leyes no sustituyen la semejanza compresible.
En gas perfecto de composición fija, las variables corregidas usuales son \(N_c=N/\sqrt{T_{01}/T_{\mathrm{ref}}}\) y \(\dot m_c=\dot m\sqrt{T_{01}/T_{\mathrm{ref}}}/(p_{01}/p_{\mathrm{ref}})\). Las referencias deben declararse. Cambiar composición o Reynolds puede impedir que dos puntos con igual caudal y velocidad corregidos sean exactamente semejantes.
% studium:end-record eyJpZCI6ICJQQVItMDAyMCIsICJzZWN0aW9uIjogImMwMyJ9

\section{Criterio físico y formulario}

Una eficiencia inferior a uno es necesaria en estos modelos adiabáticos, pero no basta para demostrar que un diseño sea realizable. Deben comprobarse Mach relativo, difusión, tensiones y estabilidad. Un resultado superior a uno puede señalar estados de referencia mal elegidos, magnitudes estáticas mezcladas con totales o calor no incluido en el balance.
Errores habituales: aplicar la definición de compresor a una turbina, confundir una pérdida de presión con una pérdida de energía, multiplicar eficiencias de escalón y usar leyes de afinidad incompresibles a través de regímenes transónicos. Autocomprobación: conserve \(\eta_{p,c}\) y aumente \(\pi_c\); justifique la tendencia de \(\eta_{c,\mathrm{tt}}\).
Formulario: \(w_c=c_p\Delta T_0\), \(w_t=-c_p\Delta T_0\), \(\eta_{c,\mathrm{tt}}=w_{cs}/w_c\), \(\eta_{t,\mathrm{tt}}=w_t/w_{ts}\), \(\phi=C_x/U\), \(\psi=|w_f|/U^2\). Referencias: ULE Generalidades, pp. \(9-12\) y 19; Diseño de compresores, pp. \(3-8\); contraste de la definición politrópica en MIT 16.50, lección 26, p. 4.
% studium:end-record eyJpZCI6ICJQQVItMDAyMSIsICJzZWN0aW9uIjogImMwMyJ9

\begin{tcolorbox}[title={Desarrollos del modelo}, studiumdef, breakable]

\noindent\textbf{Modelo desarrollado: Compresión total-total}

\begin{equation}
w_c=\frac{T_{0e}\,c_p}{\eta_{tt}}\left(\pi_c^{\frac{\gamma-1}{\gamma}}-1\right)
\label{eq:EC-C03}
\end{equation}

\noindent\textit{Significado.} Trabajo específico real de un compresor total-total. $\eta_{tt}$ es la eficiencia total-total del escalón, distinta de una eficiencia politrópica.

\noindent\textbf{Hipótesis}
\begin{itemize}
\item Gas caloríficamente perfecto y proceso adiabático.
\end{itemize}

\noindent\textbf{Principios}
\begin{itemize}
\item Balance y definiciones físicas desarrollados en el capítulo; propiedades y dominio declarados en el problema.
\end{itemize}

\noindent\textbf{Aplicabilidad}
\begin{itemize}
\item Modelo educativo del caso C03-01-R
\end{itemize}

\noindent\textbf{Limitaciones}
\begin{itemize}
\item Modelo educativo; la aplicabilidad a una máquina real requiere revisión experta.
\end{itemize}

\noindent\textbf{Pasos}
\begin{enumerate}
\item
La referencia isentrópica utiliza la misma presión de salida.

\[
\fitmath{\Delta T_{0s}=T_{0e} \left(\pi_c^{\frac{\gamma - 1}{\gamma}} - 1\right)}
\]
\item
El trabajo real es mayor que el ideal.

\[
\fitmath{w_c=\frac{T_{0e} c_p \left(\pi_c^{\frac{\gamma - 1}{\gamma}} - 1\right)}{\eta}}
\]
\end{enumerate}

{\noindent\footnotesize verificada dimensionalmente.\par}
\end{tcolorbox}

\begin{tcolorbox}[title={Problema resuelto 1}, studiumworked, breakable]

\noindent\textbf{Enunciado}
\(C03-01-R \cdot\) Compresión total-total. Calcule temperatura total de salida y trabajo de una compresión con relación de presión y rendimiento dados. Datos: \(T_{0e}=294\, \mathrm{K}\); \(\pi_c=4.6\); \(\eta=0.84\); \(\gamma=1.4\); \(c_p=1004.5\, \mathrm{J}\,\mathrm{kg}^{-1}\,\mathrm{K}^{-1}\).

\noindent\textbf{Datos}
\begin{itemize}
\item $T_{0e} = 294\,\mathrm{K}$ — temperatura total entrada
\item $\pi_c = 4.6$ — relación presión total
\item $\eta = 0.84$ — rendimiento total-total
\item $\gamma = 1.4$ — razón de calores
\item $c_p = 1004.5\,\mathrm{J}\,\mathrm{kg}^{-1}\,\mathrm{K}^{-1}$ — calor específico
\end{itemize}

\noindent\textbf{Incógnita}
$w_c$, en $\mathrm{J\,kg^{-1}}$; también las magnitudes intermedias desarrolladas.

\noindent\textbf{Modelo físico}

\noindent\textit{Compresión total-total}

\[
\fitmath{w_c=\frac{T_{0e} c_p \left(\pi_c^{\frac{\gamma - 1}{\gamma}} - 1\right)}{\eta}}
\]

\noindent\textbf{Hipótesis}
\begin{itemize}
\item Gas caloríficamente perfecto y proceso adiabático.
\end{itemize}

\noindent\textbf{Desarrollo paso a paso}

\begin{enumerate}

\item Temperatura total de salida

\[\fitmath{\mathrm{Tout}=\mathrm{Tin}\cdot \left(1+\frac{{\mathrm{pr}}^{\frac{\gamma-1}{\gamma}}-1}{\eta}\right)}\]

\[\fitmath{\mathrm{Tout}=(294)\cdot \left(1+\frac{{(4.6)}^{\frac{(1.4)-1}{(1.4)}}-1}{(0.84)}\right)\approx 485.2868\,\mathrm{K}}\]

\item La referencia isentrópica utiliza la misma presión de salida.

\[\fitmath{\mathrm{dTs}=\mathrm{Tin}\cdot \left({\mathrm{pr}}^{\frac{\gamma-1}{\gamma}}-1\right)}\]

\[\fitmath{\mathrm{dTs}=(294)\cdot \left({(4.6)}^{\frac{(1.4)-1}{(1.4)}}-1\right)\approx 160.6809\,\mathrm{K}}\]

\item El trabajo real es mayor que el ideal.

\[\fitmath{\mathrm{wc}=\frac{\mathrm{cp}\cdot \mathrm{Tin}\cdot \left({\mathrm{pr}}^{\frac{\gamma-1}{\gamma}}-1\right)}{\eta}}\]

\[\fitmath{\mathrm{wc}=\frac{(1004.5)\cdot (294)\cdot \left({(4.6)}^{\frac{(1.4)-1}{(1.4)}}-1\right)}{(0.84)}\approx 192147.6\,\mathrm{J/kg}}\]

\end{enumerate}


\noindent\textbf{Sustitución numérica}

\[
\fitmath{\frac{\left(294\,\mathrm{K}\right) \left(1004.5\,\mathrm{J}\,\mathrm{kg}^{-1}\,\mathrm{K}^{-1}\right) \left(\left(4.6\right)^{\frac{\left(1.4\right) - 1}{\left(1.4\right)}} - 1\right)}{\left(0.84\right)}}
\]

\noindent\textbf{Resultado}

\[
\fitmath{w_c=1.921475762\times10^{5}\,\mathrm{J}\,\mathrm{kg}^{-1}}
\]

\noindent\textbf{Temperatura total de salida}
\[
\fitmath{T_{0s}=T_{0e}+w_c/c_p=485.2867856\,\mathrm{K}}
\]

\noindent\textbf{Comprobación energética.} $c_p(T_{0s}-T_{0e})=w_c$; no confundir el incremento isentrópico con la temperatura real de salida.

\noindent\textbf{Interpretación}
La eficiencia divide el trabajo ideal en compresión. El resultado debe superar el coste isentrópico.

\noindent\textbf{Comprobación de los pasos}
\begin{itemize}
\item Paso 2: $160.681$ — reproducido
\item Paso 4: $192148$ — reproducido
\end{itemize}

\noindent\textbf{Comprobación de la respuesta} la aritmética del último paso reproduce la respuesta enunciada a la precisión con que está escrita. Es una reproducción aritmética; no valida el modelo físico ni sus hipótesis.

\noindent\textbf{Limitaciones}
\begin{itemize}
\item Modelo educativo; la aplicabilidad a una máquina real requiere revisión experta.
\end{itemize}

\noindent\textbf{Errores frecuentes}
\begin{itemize}
\item Multiplicar por eta para calcular el trabajo real.
\end{itemize}
% studium:end-record eyJpZCI6ICJQUkItMDAxNyIsICJzZWN0aW9uIjogImMwMyJ9
\end{tcolorbox}

\begin{tcolorbox}[title={Problema resuelto 2}, studiumworked, breakable]

\noindent\textbf{Enunciado}
\(C03-02-R \cdot\) Expansión total-total. Determine el trabajo positivo extraído por una turbina con razón de presión salida/entrada dada. Datos: \(T_{0e}=1435\, \mathrm{K}\); \(\pi_t=0.43\); \(\eta=0.89\); \(\gamma=1.33\); \(c_p=1156.697\, \mathrm{J}\,\mathrm{kg}^{-1}\,\mathrm{K}^{-1}\).

\noindent\textbf{Datos}
\begin{itemize}
\item $T_{0e} = 1435\,\mathrm{K}$ — temperatura total entrada
\item $\pi_t = 0.43$ — presión total salida/entrada
\item $\eta = 0.89$ — rendimiento total-total
\item $\gamma = 1.33$ — razón de calores
\item $c_p = 1156.696969697\,\mathrm{J}\,\mathrm{kg}^{-1}\,\mathrm{K}^{-1}$ — calor específico compatible con R=287
\end{itemize}

\noindent\textbf{Incógnita}
$w_t$, en $\mathrm{J\,kg^{-1}}$; también las magnitudes intermedias desarrolladas.

\noindent\textbf{Modelo físico}

\noindent\textit{Expansión total-total}

\[
\fitmath{w_t=T_{0e} c_p \eta \left(1 - \pi_t^{\frac{\gamma - 1}{\gamma}}\right)}
\]

\noindent\textbf{Hipótesis}
\begin{itemize}
\item Sin sangrados ni refrigeración, gas perfecto constante.
\end{itemize}

\noindent\textbf{Desarrollo paso a paso}

\begin{enumerate}

\item La caída ideal se refiere a igual presión final.

\[\fitmath{\mathrm{dTs}=\mathrm{Tin}\cdot \left(1-{\mathrm{pr}}^{\frac{\gamma-1}{\gamma}}\right)}\]

\[\fitmath{\mathrm{dTs}=(1435)\cdot \left(1-{(0.43)}^{\frac{(1.33)-1}{(1.33)}}\right)\approx 271.1206\,\mathrm{K}}\]

\item La caída real es una fracción de la ideal.

\[\fitmath{\mathrm{wt}=\eta\cdot \mathrm{cp}\cdot \mathrm{Tin}\cdot \left(1-{\mathrm{pr}}^{\frac{\gamma-1}{\gamma}}\right)}\]

\[\fitmath{\mathrm{wt}=(0.89)\cdot (1156.69697)\cdot (1435)\cdot \left(1-{(0.43)}^{\frac{(1.33)-1}{(1.33)}}\right)\approx 279107.9\,\mathrm{J/kg}}\]

\end{enumerate}


\noindent\textbf{Sustitución numérica}

\[
\fitmath{\left(1435\,\mathrm{K}\right) \left(1156.696969697\,\mathrm{J}\,\mathrm{kg}^{-1}\,\mathrm{K}^{-1}\right) \left(0.89\right) \left(1 - \left(0.43\right)^{\frac{\left(1.33\right) - 1}{\left(1.33\right)}}\right)}
\]

\noindent\textbf{Resultado}

\[
\fitmath{w_t=2.791078880\times10^{5}\,\mathrm{J}\,\mathrm{kg}^{-1}}
\]

\noindent\textbf{Interpretación}
El trabajo útil idealizado es positivo y menor que la caída isentrópica disponible.

\noindent\textbf{Comprobación de los pasos}
\begin{itemize}
\item Paso 2: $271.121$ — reproducido
\item Paso 4: $279108$ — reproducido
\end{itemize}

\noindent\textbf{Comprobación de la respuesta} la aritmética del último paso reproduce la respuesta enunciada a la precisión con que está escrita. Es una reproducción aritmética; no valida el modelo físico ni sus hipótesis.

\noindent\textbf{Limitaciones}
\begin{itemize}
\item Modelo educativo; la aplicabilidad a una máquina real requiere revisión experta.
\end{itemize}

\noindent\textbf{Errores frecuentes}
\begin{itemize}
\item Usar presión entrada/salida sin transformar la fórmula.
\end{itemize}
% studium:end-record eyJpZCI6ICJQUkItMDAxOSIsICJzZWN0aW9uIjogImMwMyJ9
\end{tcolorbox}

\begin{tcolorbox}[title={Problema resuelto 3}, studiumworked, breakable]

\noindent\textbf{Enunciado}
\(C03-03-R \cdot\) Conversión politrópica. Un compresor tiene eficiencia politrópica uniforme. Calcule la razón de temperaturas y su eficiencia isentrópica global. Datos: \(\pi_c=12.4\); \(\eta_p=0.905\); \(\gamma=1.4\).

\noindent\textbf{Datos}
\begin{itemize}
\item $\pi_c = 12.4$ — relación presión total
\item $\eta_p = 0.905$ — eficiencia politrópica
\item $\gamma = 1.4$ — razón de calores
\end{itemize}

\noindent\textbf{Incógnita}
$\eta$ (adimensional); también las magnitudes intermedias desarrolladas.

\noindent\textbf{Modelo físico}

\noindent\textit{Conversión politrópica}

\[
\fitmath{\eta=\frac{\pi_c^{\frac{\gamma - 1}{\gamma}} - 1}{\pi_c^{\frac{\gamma - 1}{\eta_p \gamma}} - 1}}
\]

\noindent\textbf{Hipótesis}
\begin{itemize}
\item Eficiencia politrópica y propiedades constantes.
\end{itemize}

\noindent\textbf{Desarrollo paso a paso}

\begin{enumerate}

\item Se integra la relación diferencial politrópica.

\[\fitmath{\mathrm{Tr}={\mathrm{pr}}^{\frac{\gamma-1}{\gamma\cdot \mathrm{ep}}}}\]

\[\fitmath{\mathrm{Tr}={(12.4)}^{\frac{(1.4)-1}{(1.4)\cdot (0.905)}}\approx 2.214115\,}\]

\item Se compara el trabajo isentrópico global con el real.

\[\fitmath{\mathrm{eta}=\frac{{\mathrm{pr}}^{\frac{\gamma-1}{\gamma}}-1}{{\mathrm{pr}}^{\frac{\gamma-1}{\gamma\cdot \mathrm{ep}}}-1}}\]

\[\fitmath{\mathrm{eta}=\frac{{(12.4)}^{\frac{(1.4)-1}{(1.4)}}-1}{{(12.4)}^{\frac{(1.4)-1}{(1.4)\cdot (0.905)}}-1}\approx 0.8673654\,}\]

\end{enumerate}


\noindent\textbf{Sustitución numérica}

\[
\fitmath{\frac{\left(12.4\right)^{\frac{\left(1.4\right) - 1}{\left(1.4\right)}} - 1}{\left(12.4\right)^{\frac{\left(1.4\right) - 1}{\left(0.905\right) \left(1.4\right)}} - 1}}
\]

\noindent\textbf{Resultado}

\[
\fitmath{\eta=8.673653999\times10^{-1}}
\]

\noindent\textbf{Interpretación}
La eficiencia global es menor que la politrópica en este modelo de compresión.

\noindent\textbf{Comprobación de los pasos}
\begin{itemize}
\item Paso 2: $2.21412$ — reproducido
\item Paso 4: $0.867365$ — reproducido
\end{itemize}

\noindent\textbf{Comprobación de la respuesta} la aritmética del último paso reproduce la respuesta enunciada a la precisión con que está escrita. Es una reproducción aritmética; no valida el modelo físico ni sus hipótesis.

\noindent\textbf{Limitaciones}
\begin{itemize}
\item Modelo educativo; la aplicabilidad a una máquina real requiere revisión experta.
\end{itemize}

\noindent\textbf{Errores frecuentes}
\begin{itemize}
\item Elevar la eficiencia de escalón al número de escalones.
\end{itemize}
% studium:end-record eyJpZCI6ICJQUkItMDAyMSIsICJzZWN0aW9uIjogImMwMyJ9
\end{tcolorbox}

\begin{tcolorbox}[title={Problema resuelto 4}, studiumworked, breakable]

\noindent\textbf{Enunciado}
\(C03-04-R \cdot\) Corrección de caudal. Convierta el caudal medido a las referencias de presión y temperatura especificadas. Datos: \(\dot m=31.7\, \mathrm{kg}\,\mathrm{s}^{-1}\); \(T_{0e}=271\, \mathrm{K}\); \(p_{0e}=82500\, \mathrm{Pa}\); \(T_{\mathrm{ref}}=288.15\, \mathrm{K}\); \(p_{\mathrm{ref}}=101325\, \mathrm{Pa}\).

\noindent\textbf{Datos}
\begin{itemize}
\item $\dot m = 31.7\,\mathrm{kg}\,\mathrm{s}^{-1}$ — caudal real
\item $T_{0e} = 271\,\mathrm{K}$ — temperatura total entrada
\item $p_{0e} = 82500\,\mathrm{Pa}$ — presión total entrada
\item $T_{\mathrm{ref}} = 288.15\,\mathrm{K}$ — temperatura referencia
\item $p_{\mathrm{ref}} = 101325\,\mathrm{Pa}$ — presión referencia
\end{itemize}

\noindent\textbf{Incógnita}
$\dot{m}_c$, en $\mathrm{kg\,s^{-1}}$; también las magnitudes intermedias desarrolladas.

\noindent\textbf{Modelo físico}

\noindent\textit{Corrección de caudal}

\[
\fitmath{\dot{m}_c=\frac{\dot m p_{\mathrm{ref}} \sqrt{\frac{T_{0e}}{T_{\mathrm{ref}}}}}{p_{0e}}}
\]

\noindent\textbf{Hipótesis}
\begin{itemize}
\item Mismo gas y referencias declaradas.
\end{itemize}

\noindent\textbf{Desarrollo paso a paso}

\begin{enumerate}

\item La corrección térmica es la raíz de la razón de temperaturas.

\[\fitmath{\theta=\sqrt{\frac{\mathrm{Tin}}{\mathrm{Tr}}}}\]

\[\fitmath{\theta=\sqrt{\frac{(271)}{(288.15)}}\approx 0.9697847\,}\]

\item La presión corrige la capacidad de paso del gas.

\[\fitmath{\mathrm{mc}=\frac{\mathrm{md}\cdot \sqrt{\frac{\mathrm{Tin}}{\mathrm{Tr}}}}{\frac{\mathrm{pin}}{\mathrm{pr}}}}\]

\[\fitmath{\mathrm{mc}=\frac{(31.7)\cdot \sqrt{\frac{(271)}{(288.15)}}}{\frac{(82500)}{(101325)}}\approx 37.75698\,\mathrm{kg/s}}\]

\end{enumerate}


\noindent\textbf{Sustitución numérica}

\[
\fitmath{\frac{\left(31.7\,\mathrm{kg}\,\mathrm{s}^{-1}\right) \left(101325\,\mathrm{Pa}\right) \sqrt{\frac{\left(271\,\mathrm{K}\right)}{\left(288.15\,\mathrm{K}\right)}}}{\left(82500\,\mathrm{Pa}\right)}}
\]

\noindent\textbf{Resultado}

\[
\fitmath{\dot{m}_c=3.775698084\times10^{1}\,\mathrm{kg}\,\mathrm{s}^{-1}}
\]

\noindent\textbf{Interpretación}
El caudal corregido es una variable de comparación. No es una nueva medición de masa real.

\noindent\textbf{Comprobación de los pasos}
\begin{itemize}
\item Paso 2: $0.969785$ — reproducido
\item Paso 4: $37.757$ — reproducido
\end{itemize}

\noindent\textbf{Comprobación de la respuesta} la aritmética del último paso reproduce la respuesta enunciada a la precisión con que está escrita. Es una reproducción aritmética; no valida el modelo físico ni sus hipótesis.

\noindent\textbf{Limitaciones}
\begin{itemize}
\item Modelo educativo; la aplicabilidad a una máquina real requiere revisión experta.
\end{itemize}

\noindent\textbf{Errores frecuentes}
\begin{itemize}
\item Invertir la corrección de presión.
\end{itemize}
% studium:end-record eyJpZCI6ICJQUkItMDAyMyIsICJzZWN0aW9uIjogImMwMyJ9
\end{tcolorbox}

\begin{tcolorbox}[title={Ejercicios propuestos}, studiumworked, breakable]

\noindent\textbf{1 (Básico, Desarrollo)}
\(C03-01-P \cdot\) Compresión total-total. Ejercicio de transferencia. Calcule temperatura total de salida y trabajo de una compresión con relación de presión y rendimiento dados. Antes de consultar la solución, prediga el efecto de cambiar el primer dato respecto del caso resuelto de esta familia. Justifique la tendencia y explique qué restricción podría convertirse en limitante. Datos: \(T_{0e}=313.11\, \mathrm{K}\); \(\pi_c=4.6\); \(\eta=0.84\); \(\gamma=1.4\); \(c_p=1004.5\, \mathrm{J}\,\mathrm{kg}^{-1}\,\mathrm{K}^{-1}\).

\noindent\textbf{2 (Intermedio, Desarrollo)}
\(C03-02-P \cdot\) Expansión total-total. Ejercicio de transferencia. Determine el trabajo positivo extraído por una turbina con razón de presión salida/entrada dada. Antes de consultar la solución, prediga el efecto de cambiar el primer dato respecto del caso resuelto de esta familia. Justifique la tendencia y explique qué restricción podría convertirse en limitante. Datos: \(T_{0e}=1528.275\, \mathrm{K}\); \(\pi_t=0.43\); \(\eta=0.89\); \(\gamma=1.33\); \(c_p=1156.697\, \mathrm{J}\,\mathrm{kg}^{-1}\,\mathrm{K}^{-1}\).

\noindent\textbf{3 (Avanzado, Desarrollo)}
\(C03-03-P \cdot\) Conversión politrópica. Ejercicio de transferencia. Un compresor tiene eficiencia politrópica uniforme. Calcule la razón de temperaturas y su eficiencia isentrópica global. Antes de consultar la solución, prediga el efecto de cambiar el primer dato respecto del caso resuelto de esta familia. Justifique la tendencia y explique qué restricción podría convertirse en limitante. Datos: \(\pi_c=13.206\); \(\eta_p=0.905\); \(\gamma=1.4\).

\noindent\textbf{4 (Nivel de examen, Desarrollo)}
\(C03-04-P \cdot\) Corrección de caudal. Ejercicio de transferencia. Convierta el caudal medido a las referencias de presión y temperatura especificadas. Antes de consultar la solución, prediga el efecto de cambiar el primer dato respecto del caso resuelto de esta familia. Justifique la tendencia y explique qué restricción podría convertirse en limitante. Datos: \(\dot m=33.7605\, \mathrm{kg}\,\mathrm{s}^{-1}\); \(T_{0e}=271\, \mathrm{K}\); \(p_{0e}=82500\, \mathrm{Pa}\); \(T_{\mathrm{ref}}=288.15\, \mathrm{K}\); \(p_{\mathrm{ref}}=101325\, \mathrm{Pa}\).
\end{tcolorbox}


\studiumchaptercheck

\part{Estudio bidimensional de turbinas axiales}

\chapter{Turbinas axiales de acción}
\label{chap:c04}

\studiumchapterguide{Objetivo: desarrollar turbinas axiales de acción, declarar hipótesis y comprobar condiciones físicas.}
% studium:end-record eyJpZCI6ICJQQVItMDAyMiIsICJzZWN0aW9uIjogImMwNCJ9

\section{Introducción y objetivos}

En una turbina de acción ideal la caída de entalpía estática tiene lugar en el estator. El rotor desvía la corriente y extrae trabajo sin caída estática ideal. Esta descripción no significa que no haya fuerzas de presión sobre las caras del álabe, ni que un rotor real carezca de pérdidas. El objetivo es relacionar esa arquitectura con sus triángulos, su trabajo y la energía residual de salida.
Se utiliza 1 a la entrada del estator, 2 a la entrada del rotor y 3 a su salida. Los ángulos se mantienen algebraicos y \(\Omega>0\). El capítulo pertenece al bloque II; la optimización ideal que se introduce sirve para interpretar límites, no para seleccionar una turbina comercial.
% studium:end-record eyJpZCI6ICJQQVItMDAyMyIsICJzZWN0aW9uIjogImMwNCJ9

\section{Qué hace cada fila}

El estator adiabático mantiene \(h_{02}=h_{01}\) y acelera el gas al transformar entalpía estática en velocidad. Introduce normalmente una componente tangencial positiva. El rotor reduce el momento angular de la corriente, recibe el par opuesto y produce trabajo positivo \(w_t\). Si el radio permanece constante:

\[
\fitmath{w_t=U(C_{\theta2}-C_{\theta3}),\qquad
h_{03}=h_{02}-w_t.}
\]
En el rotor ideal a radio constante, \(h_2+W_2^2/2=h_3+W_3^2/2\). Si \(h_2=h_3\), entonces \(W_2=W_3\). Con componente axial constante, las componentes tangenciales relativas son opuestas. Se deduce

\[
\fitmath{W_{\theta3}=-W_{\theta2},\qquad C_{\theta3}=2U-C_{\theta2}.}
\]
La figura~\ref{fig:accion} muestra esta construcción para un caso calculado. Al cambiar \(U\), cambia el remolino de salida, aunque se mantenga la velocidad de entrada al rotor. La forma del triángulo contiene una decisión de diseño y no una identidad universal.
% studium:end-record eyJpZCI6ICJQQVItMDAyNCIsICJzZWN0aW9uIjogImMwNCJ9

\par\vspace{\baselineskip}
\noindent\begin{minipage}{\linewidth}
\centering
\begingroup
\setbox0=\hbox{\begin{tikzpicture}
\draw[->,gray] (0,0)--(2.5,0);\draw[->,blue,thick] (0,0)--(1.8,4.6) node[right] {$C_2$};\draw[->,orange,thick] (0,0)--(0,2.2) node[left] {$U$};\draw[->,teal,thick] (0,2.2)--(1.8,4.6) node[midway,left] {$W_2$};\begin{scope}[xshift=5cm]\draw[->,gray] (0,0)--(2.5,0);\draw[->,blue,thick] (0,0)--(1.8,-.2) node[right] {$C_3$};\draw[->,orange,thick] (0,0)--(0,2.2) node[left] {$U$};\draw[->,teal,thick] (0,2.2)--(1.8,-.2) node[midway,right] {$W_3$};\end{scope}\node at (1,-.8) {Entrada 2};\node at (6,-.8) {Salida 3};
\end{tikzpicture}}%
\count0=\ht0
\advance\count0 by \dp0
\count2=\wd0
\count4=\dimexpr\textheight/100*38\relax
\count6=\linewidth
\ifnum\count0<1
\count0=1
\fi
\ifnum\count2<1
\count2=1
\fi
\divide\count0 by 4096
\divide\count2 by 4096
\divide\count4 by 4096
\divide\count6 by 4096
\ifnum\count0<1
\count0=1
\fi
\ifnum\count2<1
\count2=1
\fi
\ifnum\count4<1
\count4=1
\fi
\ifnum\count6<1
\count6=1
\fi
\count8=\count2
\ifnum\count0>\count4
\count8=\numexpr\count2*\count4/\count0\relax
\fi
\ifnum\count8>\count6
\resizebox{\linewidth}{!}{\begin{tikzpicture}
\draw[->,gray] (0,0)--(2.5,0);\draw[->,blue,thick] (0,0)--(1.8,4.6) node[right] {$C_2$};\draw[->,orange,thick] (0,0)--(0,2.2) node[left] {$U$};\draw[->,teal,thick] (0,2.2)--(1.8,4.6) node[midway,left] {$W_2$};\begin{scope}[xshift=5cm]\draw[->,gray] (0,0)--(2.5,0);\draw[->,blue,thick] (0,0)--(1.8,-.2) node[right] {$C_3$};\draw[->,orange,thick] (0,0)--(0,2.2) node[left] {$U$};\draw[->,teal,thick] (0,2.2)--(1.8,-.2) node[midway,right] {$W_3$};\end{scope}\node at (1,-.8) {Entrada 2};\node at (6,-.8) {Salida 3};
\end{tikzpicture}}%
\else
\ifnum\count0>\count4
\resizebox{!}{\dimexpr\textheight/100*38\relax}{\begin{tikzpicture}
\draw[->,gray] (0,0)--(2.5,0);\draw[->,blue,thick] (0,0)--(1.8,4.6) node[right] {$C_2$};\draw[->,orange,thick] (0,0)--(0,2.2) node[left] {$U$};\draw[->,teal,thick] (0,2.2)--(1.8,4.6) node[midway,left] {$W_2$};\begin{scope}[xshift=5cm]\draw[->,gray] (0,0)--(2.5,0);\draw[->,blue,thick] (0,0)--(1.8,-.2) node[right] {$C_3$};\draw[->,orange,thick] (0,0)--(0,2.2) node[left] {$U$};\draw[->,teal,thick] (0,2.2)--(1.8,-.2) node[midway,right] {$W_3$};\end{scope}\node at (1,-.8) {Entrada 2};\node at (6,-.8) {Salida 3};
\end{tikzpicture}}%
\else
\begin{tikzpicture}
\draw[->,gray] (0,0)--(2.5,0);\draw[->,blue,thick] (0,0)--(1.8,4.6) node[right] {$C_2$};\draw[->,orange,thick] (0,0)--(0,2.2) node[left] {$U$};\draw[->,teal,thick] (0,2.2)--(1.8,4.6) node[midway,left] {$W_2$};\begin{scope}[xshift=5cm]\draw[->,gray] (0,0)--(2.5,0);\draw[->,blue,thick] (0,0)--(1.8,-.2) node[right] {$C_3$};\draw[->,orange,thick] (0,0)--(0,2.2) node[left] {$U$};\draw[->,teal,thick] (0,2.2)--(1.8,-.2) node[midway,right] {$W_3$};\end{scope}\node at (1,-.8) {Entrada 2};\node at (6,-.8) {Salida 3};
\end{tikzpicture}%
\fi
\fi
\endgroup
\par\nopagebreak
\studiumcaption{fig:accion}{Acción ideal calculada: $U=220\,\mathrm{m/s}$, $C_x=180\,\mathrm{m/s}$, $C_{\theta 2}=460\,\mathrm{m/s}$. Escala común en ambos triángulos.}
\end{minipage}
\par\vspace{\baselineskip}
% studium:end-record eyJpZCI6ICJGSUctMDAwNCIsICJzZWN0aW9uIjogImMwNCJ9

\section{Trabajo y óptimo ideal}


Sustituyendo el cierre de impulso ideal en Euler:



\[
\fitmath{w_t=2U(C_{\theta2}-U).}
\]
Para una velocidad \(C_2\) y un ángulo \(\alpha_2\) fijados, \(C_{\theta2}=C_2\sin\alpha_2\). Derivando respecto de \(U\), el máximo trabajo se obtiene en

\[
\fitmath{U_{opt}=\frac{C_{\theta2}}{2},\qquad
w_{t,max}=\frac{C_{\theta2}^2}{2}.}
\]
En ese punto \(C_{\theta3}=0\): no queda remolino, pero permanece \(C_x\). Por ello la energía cinética de salida no se anula. No debe identificarse la razón \(w_t/(C_2^2/2)\) con un rendimiento termodinámico total-total: solo mide la utilización de la energía cinética incidente bajo las hipótesis establecidas.
Con \(U\) demasiado pequeño, el gas conserva mucho momento angular; con \(U\) demasiado grande, el rotor puede comunicar momento angular al gas y deja de producir trabajo. El intervalo ideal de trabajo positivo es \(0<U<C_{\theta2}\). La condición mecánica, el caudal y el acoplamiento al compresor limitan además los valores posibles de \(U\).
% studium:end-record eyJpZCI6ICJQQVItMDAyNSIsICJzZWN0aW9uIjogImMwNCJ9

\section{Acción real y grado de reacción}


La reacción se define para la turbina por



\[
\fitmath{\mathcal R_t=\frac{h_2-h_3}{h_1-h_3}.}
\]
Si el escalón es repetitivo en velocidad absoluta, de modo que \(C_1=C_3\), la caída estática total del escalón coincide con \(w_t\). Con radio y componente axial constantes se obtiene

\[
\fitmath{\mathcal R_t=1-\frac{C_{\theta2}+C_{\theta3}}{2U}.}
\]
Esta fórmula compacta necesita todas esas hipótesis. Si \(C_1\ne C_3\), la definición estática sigue siendo válida, pero su denominador incluye el cambio de energía cinética entre extremos. Una turbina con \(\mathcal R_t=0\) ideal no implica igualdad de presiones totales relativas en el rotor real.
Un modelo simplificado puede especificar un cociente \(k=W_3/W_2\) y una dirección de salida. En tal caso debe recalcularse \(C_{\theta3}\), el trabajo y las temperaturas. No se impone además \(h_2=h_3\) si el conjunto de condiciones ya resulta incompatible. Un único coeficiente de velocidad no representa todas las pérdidas de perfil, secundarias, de punta y de mezcla.
% studium:end-record eyJpZCI6ICJQQVItMDAyNiIsICJzZWN0aW9uIjogImMwNCJ9

\section{Carga de la cascada y límites}

La fuerza tangencial sobre una fila se relaciona con el cambio de cantidad de movimiento. Para velocidad axial y densidad aproximadamente constantes, un parámetro de carga tipo Zweifel puede expresarse mediante una fuerza tangencial por paso, dividida por presión dinámica de salida y cuerda axial. Aquí se emplea

\[
\fitmath{Z=2\frac{s}{c_x}\cos^2 a_s\,|\tan a_e-\tan a_s|.}
\]
Los ángulos \(a_e,a_s\) pertenecen a la referencia de la fila: absolutos para estator y relativos para rotor. Esta expresión de bajo Mach es una estimación de carga; su valor recomendado depende de familia de perfil y condiciones. No se usa como prueba de una eficiencia ni como garantía de ausencia de separación.
En turbinas de acción, el gran giro relativo puede causar separación a pesar de que la caída de presión global sea favorable. La idea de que una turbina nunca pierde sustentación es incorrecta. El bloqueo de garganta y los esfuerzos térmicos son restricciones diferentes que se examinarán en los capítulos siguientes.
% studium:end-record eyJpZCI6ICJQQVItMDAyNyIsICJzZWN0aW9uIjogImMwNCJ9

\section{Errores y síntesis}

Errores habituales: suponer que toda la energía cinética incidente se transforma en trabajo, invertir la dirección de \(W_3\), fijar simultáneamente demasiadas condiciones ideales y reales y aplicar la expresión abreviada de reacción a un escalón no repetitivo. Autocomprobación: construya dos triángulos con igual \(C_{\theta2}\) y distinto \(U\); identifique cuándo cambia el signo del trabajo.
Formulario: \(w_t=U(C_{\theta2}-C_{\theta3})\); para acción ideal con \(C_x\) constante, \(C_{\theta3}=2U-C_{\theta2}\), \(w_t=2U(C_{\theta2}-U)\). La reacción general se calcula con entalpías estáticas, no con una regla visual.
Referencias: ULE Turbinas axiales, pp. \(2-6\) y \(17-22\). MIT 16.50, lección 27, como contraste de la definición de reacción; la lección 28 aporta el contexto del parámetro de carga. Las ecuaciones y las figuras se desarrollan aquí con notación propia.
% studium:end-record eyJpZCI6ICJQQVItMDAyOCIsICJzZWN0aW9uIjogImMwNCJ9

\begin{tcolorbox}[title={Desarrollos del modelo}, studiumdef, breakable]

\noindent\textbf{Modelo desarrollado: Acción ideal}

\begin{equation}
\Delta T_0=\frac{2\,U\,(C_\theta-U)}{c_p}
\label{eq:EC-C04}
\end{equation}

\noindent\textit{Significado.} Descenso de temperatura total en una acción ideal con radio y componente axial constantes. $\Delta T_0$ es una caída positiva; no equivale a $T_{\mathrm{salida}}-T_{\mathrm{entrada}}$.

\noindent\textbf{Hipótesis}
\begin{itemize}
\item Acción ideal, radio y componente axial constantes.
\end{itemize}

\noindent\textbf{Principios}
\begin{itemize}
\item Balance y definiciones físicas desarrollados en el capítulo; propiedades y dominio declarados en el problema.
\end{itemize}

\noindent\textbf{Aplicabilidad}
\begin{itemize}
\item Modelo educativo del caso C04-01-R
\end{itemize}

\noindent\textbf{Limitaciones}
\begin{itemize}
\item Modelo educativo; la aplicabilidad a una máquina real requiere revisión experta.
\end{itemize}

\noindent\textbf{Pasos}
\begin{enumerate}
\item
La simetría relativa da la componente tangencial de salida.

\[
\fitmath{C_{\theta s}=- C_\theta + 2 U}
\]
\item
Euler proporciona trabajo positivo de turbina.

\[
\fitmath{w_t=2 U \left(C_\theta - U\right)}
\]
\item
Se divide por el calor específico.

\[
\fitmath{\Delta T_0=\frac{2 U \left(C_\theta - U\right)}{c_p}}
\]
\end{enumerate}

{\noindent\footnotesize verificada dimensionalmente.\par}
\end{tcolorbox}

\begin{tcolorbox}[title={Problema resuelto 1}, studiumworked, breakable]

\noindent\textbf{Enunciado}
\(C04-01-R \cdot\) Acción ideal. Una turbina de acción ideal recibe las componentes y velocidad periférica dadas. Halle remolino de salida, trabajo y descenso de temperatura total. Datos: \(C_\theta=462\, \mathrm{m}\,\mathrm{s}^{-1}\); \(U=211\, \mathrm{m}\,\mathrm{s}^{-1}\); \(c_p=1156.697\, \mathrm{J}\,\mathrm{kg}^{-1}\,\mathrm{K}^{-1}\).

\noindent\textbf{Datos}
\begin{itemize}
\item $C_\theta = 462\,\mathrm{m}\,\mathrm{s}^{-1}$ — remolino entrada rotor
\item $U = 211\,\mathrm{m}\,\mathrm{s}^{-1}$ — velocidad periférica
\item $c_p = 1156.696969697\,\mathrm{J}\,\mathrm{kg}^{-1}\,\mathrm{K}^{-1}$ — calor específico
\end{itemize}

\noindent\textbf{Incógnita}
$\Delta T_0$, en $\mathrm{K}$; también las magnitudes intermedias desarrolladas.

\noindent\textbf{Modelo físico}

\noindent\textit{Acción ideal}

\[
\fitmath{\Delta T_0=\frac{2 U \left(C_\theta - U\right)}{c_p}}
\]

\noindent\textbf{Hipótesis}
\begin{itemize}
\item Acción ideal, radio y componente axial constantes.
\end{itemize}

\noindent\textbf{Desarrollo paso a paso}

\begin{enumerate}

\item La simetría relativa da la componente tangencial de salida.

\[\fitmath{\mathrm{cts}=2\cdot \mathrm{uu}-\mathrm{ct}}\]

\[\fitmath{\mathrm{cts}=2\cdot (211)-(462)\approx -40\,\mathrm{m/s}}\]

\item Euler proporciona trabajo positivo de turbina.

\[\fitmath{\mathrm{wt}=2\cdot \mathrm{uu}\cdot \left(\mathrm{ct}-\mathrm{uu}\right)}\]

\[\fitmath{\mathrm{wt}=2\cdot (211)\cdot \left((462)-(211)\right)\approx 105922\,\mathrm{J/kg}}\]

\item Se divide por el calor específico.

\[\fitmath{\mathrm{dTz}=\frac{2\cdot \mathrm{uu}\cdot \left(\mathrm{ct}-\mathrm{uu}\right)}{\mathrm{cp}}}\]

\[\fitmath{\mathrm{dTz}=\frac{2\cdot (211)\cdot \left((462)-(211)\right)}{(1156.69697)}\approx 91.57282\,\mathrm{K}}\]

\end{enumerate}


\noindent\textbf{Sustitución numérica}

\[
\fitmath{\frac{2 \left(211\,\mathrm{m}\,\mathrm{s}^{-1}\right) \left(\left(462\,\mathrm{m}\,\mathrm{s}^{-1}\right) - \left(211\,\mathrm{m}\,\mathrm{s}^{-1}\right)\right)}{\left(1156.696969697\,\mathrm{J}\,\mathrm{kg}^{-1}\,\mathrm{K}^{-1}\right)}}
\]

\noindent\textbf{Resultado}

\[
\fitmath{\Delta T_0=9.157281706\times10^{1}\,\mathrm{K}}
\]

\noindent\textbf{Interpretación}
El remolino de salida es negativo con estos datos. Mantener su signo es necesario para cerrar Euler.

\noindent\textbf{Comprobación de los pasos}
\begin{itemize}
\item Paso 2: $-40$ — reproducido
\item Paso 4: $105922$ — reproducido
\item Paso 6: $91.5728$ — reproducido
\end{itemize}

\noindent\textbf{Comprobación de la respuesta} la aritmética del último paso reproduce la respuesta enunciada a la precisión con que está escrita. Es una reproducción aritmética; no valida el modelo físico ni sus hipótesis.

\noindent\textbf{Limitaciones}
\begin{itemize}
\item Modelo educativo; la aplicabilidad a una máquina real requiere revisión experta.
\end{itemize}

\noindent\textbf{Errores frecuentes}
\begin{itemize}
\item Tomar valor absoluto del remolino antes de aplicar Euler.
\end{itemize}
% studium:end-record eyJpZCI6ICJQUkItMDAyNSIsICJzZWN0aW9uIjogImMwNCJ9
\end{tcolorbox}

\begin{tcolorbox}[title={Problema resuelto 2}, studiumworked, breakable]

\noindent\textbf{Enunciado}
\(C04-02-R \cdot\) Óptimo de utilización. Para una velocidad de entrada al rotor y ángulo fijados, obtenga la velocidad periférica de máximo trabajo ideal y ese trabajo. Datos: \(C_e=590\, \mathrm{m}\,\mathrm{s}^{-1}\); \(\alpha_e=1.117011\, \mathrm{rad}\).

\noindent\textbf{Datos}
\begin{itemize}
\item $C_e = 590\,\mathrm{m}\,\mathrm{s}^{-1}$ — módulo entrada rotor
\item $\alpha_e = 1.1170107212764\,\mathrm{rad}$ — ángulo absoluto axial
\end{itemize}

\noindent\textbf{Incógnita}
$w_{\max}$, en $\mathrm{J\,kg^{-1}}$; también las magnitudes intermedias desarrolladas.

\noindent\textbf{Modelo físico}

\noindent\textit{Óptimo de utilización}

\[
\fitmath{w_{\max}=\frac{C_e^{2} \sin^{2}{\left(\alpha_e \right)}}{2}}
\]

\noindent\textbf{Hipótesis}
\begin{itemize}
\item Velocidad y orientación incidentes fijadas; rotor de acción ideal.
\end{itemize}

\noindent\textbf{Desarrollo paso a paso}

\begin{enumerate}

\item Se proyecta la velocidad tangencial.

\[\fitmath{\mathrm{ct}=\mathrm{Vel}\cdot \sin\left(\mathrm{ang}\right)}\]

\[\fitmath{\mathrm{ct}=(590)\cdot \sin\left((1.117010721)\right)\approx 530.2885\,\mathrm{m/s}}\]

\item El óptimo de la parábola está en U=Ctheta/2.

\[\fitmath{\mathrm{uopt}=\frac{\mathrm{Vel}\cdot \sin\left(\mathrm{ang}\right)}{2}}\]

\[\fitmath{\mathrm{uopt}=\frac{(590)\cdot \sin\left((1.117010721)\right)}{2}\approx 265.1442\,\mathrm{m/s}}\]

\item Se evalúa el máximo ideal.

\[\fitmath{\mathrm{wmax}=\frac{{\mathrm{Vel}}^{2}\cdot {\sin\left(\mathrm{ang}\right)}^{2}}{2}}\]

\[\fitmath{\mathrm{wmax}=\frac{{(590)}^{2}\cdot {\sin\left((1.117010721)\right)}^{2}}{2}\approx 140602.9\,\mathrm{J/kg}}\]

\end{enumerate}


\noindent\textbf{Sustitución numérica}

\[
\fitmath{\frac{\left(\left(590\,\mathrm{m}\,\mathrm{s}^{-1}\right)\right)^{2} \sin^{2}{\left(\left(1.1170107212764\,\mathrm{rad}\right) \right)}}{2}}
\]

\noindent\textbf{Resultado}

\[
\fitmath{w_{\max}=1.406029399\times10^{5}\,\mathrm{J}\,\mathrm{kg}^{-1}}
\]

\noindent\textbf{Interpretación}
Queda velocidad axial a la salida incluso en el óptimo. La utilización cinética no es una eficiencia total-total.

\noindent\textbf{Comprobación de los pasos}
\begin{itemize}
\item Paso 2: $530.288$ — reproducido
\item Paso 4: $265.144$ — reproducido
\item Paso 6: $140603$ — reproducido
\end{itemize}

\noindent\textbf{Comprobación de la respuesta} la aritmética del último paso reproduce la respuesta enunciada a la precisión con que está escrita. Es una reproducción aritmética; no valida el modelo físico ni sus hipótesis.

\noindent\textbf{Limitaciones}
\begin{itemize}
\item Modelo educativo; la aplicabilidad a una máquina real requiere revisión experta.
\end{itemize}

\noindent\textbf{Errores frecuentes}
\begin{itemize}
\item Suponer que el óptimo elimina toda velocidad de salida.
\end{itemize}
% studium:end-record eyJpZCI6ICJQUkItMDAyNyIsICJzZWN0aW9uIjogImMwNCJ9
\end{tcolorbox}

\begin{tcolorbox}[title={Problema resuelto 3}, studiumworked, breakable]

\noindent\textbf{Enunciado}
\(C04-03-R \cdot\) Rotor con reducción de velocidad relativa. Un rotor recibe W tangencial de entrada positiva y descarga con dirección relativa simétrica y módulo reducido por k. Halle la potencia extraída. Datos: \(C_x=176\, \mathrm{m}\,\mathrm{s}^{-1}\); \(C_\theta=493\, \mathrm{m}\,\mathrm{s}^{-1}\); \(U=238\, \mathrm{m}\,\mathrm{s}^{-1}\); \(k_W=0.92\); \(\dot m=17.8\, \mathrm{kg}\,\mathrm{s}^{-1}\).

\noindent\textbf{Datos}
\begin{itemize}
\item $C_x = 176\,\mathrm{m}\,\mathrm{s}^{-1}$ — componente axial entrada
\item $C_\theta = 493\,\mathrm{m}\,\mathrm{s}^{-1}$ — remolino entrada
\item $U = 238\,\mathrm{m}\,\mathrm{s}^{-1}$ — velocidad periférica
\item $k_W = 0.92$ — cociente Wsal/Went
\item $\dot m = 17.8\,\mathrm{kg}\,\mathrm{s}^{-1}$ — caudal
\end{itemize}

\noindent\textbf{Incógnita}
$\dot{W}$, en $\mathrm{W}$; también las magnitudes intermedias desarrolladas.

\noindent\textbf{Modelo físico}

\noindent\textit{Rotor con reducción de velocidad relativa}

\[
\fitmath{\dot{W}=\dot m U \left(C_\theta + k_W \left(C_\theta - U\right) - U\right)}
\]

\noindent\textbf{Hipótesis}
\begin{itemize}
\item Radio constante; k y dirección relativa son datos del modelo.
\end{itemize}

\noindent\textbf{Desarrollo paso a paso}

\begin{enumerate}

\item La salida tangencial absoluta combina U y la componente relativa reducida.

\[\fitmath{\mathrm{cts}=\mathrm{uu}-\mathrm{kk}\cdot \left(\mathrm{ct}-\mathrm{uu}\right)}\]

\[\fitmath{\mathrm{cts}=(238)-(0.92)\cdot \left((493)-(238)\right)\approx 3.4\,\mathrm{m/s}}\]

\item Se calcula el trabajo de Euler.

\[\fitmath{\mathrm{wt}=\mathrm{uu}\cdot \left(\mathrm{ct}-\left(\mathrm{uu}-\mathrm{kk}\cdot \left(\mathrm{ct}-\mathrm{uu}\right)\right)\right)}\]

\[\fitmath{\mathrm{wt}=(238)\cdot \left((493)-\left((238)-(0.92)\cdot \left((493)-(238)\right)\right)\right)\approx 116524.8\,\mathrm{J/kg}}\]

\item El caudal convierte trabajo en potencia.

\[\fitmath{\mathrm{Pow}=\mathrm{md}\cdot \mathrm{uu}\cdot \left(\mathrm{ct}-\left(\mathrm{uu}-\mathrm{kk}\cdot \left(\mathrm{ct}-\mathrm{uu}\right)\right)\right)}\]

\[\fitmath{\mathrm{Pow}=(17.8)\cdot (238)\cdot \left((493)-\left((238)-(0.92)\cdot \left((493)-(238)\right)\right)\right)\approx 2074141\,\mathrm{W}}\]

\end{enumerate}


\noindent\textbf{Sustitución numérica}

\[
\fitmath{\left(17.8\,\mathrm{kg}\,\mathrm{s}^{-1}\right) \left(238\,\mathrm{m}\,\mathrm{s}^{-1}\right) \left(\left(493\,\mathrm{m}\,\mathrm{s}^{-1}\right) + \left(0.92\right) \left(\left(493\,\mathrm{m}\,\mathrm{s}^{-1}\right) - \left(238\,\mathrm{m}\,\mathrm{s}^{-1}\right)\right) - \left(238\,\mathrm{m}\,\mathrm{s}^{-1}\right)\right)}
\]

\noindent\textbf{Resultado}

\[
\fitmath{\dot{W}=2.074141440\times10^{6}\,\mathrm{W}}
\]

\noindent\textbf{Interpretación}
La componente axial de salida es k veces la de entrada al mantener el ángulo relativo simétrico. No se impone además acción estática perfecta.

\noindent\textbf{Comprobación de los pasos}
\begin{itemize}
\item Paso 2: $3.4$ — reproducido
\item Paso 4: $116525$ — reproducido
\item Paso 6: $2.07414\times 10^{6}$ — reproducido
\end{itemize}

\noindent\textbf{Comprobación de la respuesta} la aritmética del último paso reproduce la respuesta enunciada a la precisión con que está escrita. Es una reproducción aritmética; no valida el modelo físico ni sus hipótesis.

\noindent\textbf{Limitaciones}
\begin{itemize}
\item Modelo educativo; la aplicabilidad a una máquina real requiere revisión experta.
\end{itemize}

\noindent\textbf{Errores frecuentes}
\begin{itemize}
\item Mantener simultáneamente Cx constante y escala k de todo el vector.
\end{itemize}
% studium:end-record eyJpZCI6ICJQUkItMDAyOSIsICJzZWN0aW9uIjogImMwNCJ9
\end{tcolorbox}

\begin{tcolorbox}[title={Problema resuelto 4}, studiumworked, breakable]

\noindent\textbf{Enunciado}
\(C04-04-R \cdot\) Carga de cascada de turbina. Estime el parámetro de carga Z de una fila con ángulos algebraicos dados, en el modelo de baja compresibilidad definido en el capítulo. Datos: \(\alpha_e=0.8203047\, \mathrm{rad}\); \(\alpha_s=-1.012291\, \mathrm{rad}\); \(s=0.024\, \mathrm{m}\); \(c_x=0.049\, \mathrm{m}\).

\noindent\textbf{Datos}
\begin{itemize}
\item $\alpha_e = 0.82030474843733\,\mathrm{rad}$ — ángulo entrada en referencia de fila
\item $\alpha_s = -1.0122909661567\,\mathrm{rad}$ — ángulo salida
\item $s = 0.024\,\mathrm{m}$ — paso
\item $c_x = 0.049\,\mathrm{m}$ — cuerda axial
\end{itemize}

\noindent\textbf{Incógnita}
$Z$ (adimensional); también las magnitudes intermedias desarrolladas.

\noindent\textbf{Modelo físico}

\noindent\textit{Carga de cascada de turbina}

\[
\fitmath{Z=\frac{2 s \left(\tan{\left(\alpha_e \right)} - \tan{\left(\alpha_s \right)}\right) \cos^{2}{\left(\alpha_s \right)}}{c_x}}
\]

\noindent\textbf{Hipótesis}
\begin{itemize}
\item Densidad y velocidad axial aproximadamente constantes; ángulos algebraicos.
\end{itemize}

\noindent\textbf{Desarrollo paso a paso}

\begin{enumerate}

\item Se calcula la razón de paso sobre cuerda axial.

\[\fitmath{\mathrm{sc}=\frac{\mathrm{ss}}{\mathrm{cc}}}\]

\[\fitmath{\mathrm{sc}=\frac{(0.024)}{(0.049)}\approx 0.4897959\,}\]

\item Se evalúa la carga a partir del cambio tangencial normalizado.

\[\fitmath{Z=2\cdot \frac{\mathrm{ss}}{\mathrm{cc}}\cdot {\cos\left(\mathrm{ps}\right)}^{2}\cdot \left(\tan\left(\mathrm{pe}\right)-\tan\left(\mathrm{ps}\right)\right)}\]

\[\fitmath{Z=2\cdot \frac{(0.024)}{(0.049)}\cdot {\cos\left((-1.012290966)\right)}^{2}\cdot \left(\tan\left((0.8203047484)\right)-\tan\left((-1.012290966)\right)\right)\approx 0.7352166\,}\]

\end{enumerate}


\noindent\textbf{Sustitución numérica}

\[
\fitmath{\frac{2 \left(0.024\,\mathrm{m}\right) \left(\tan{\left(\left(0.82030474843733\,\mathrm{rad}\right) \right)} - \tan{\left(\left(-1.0122909661567\,\mathrm{rad}\right) \right)}\right) \cos^{2}{\left(\left(-1.0122909661567\,\mathrm{rad}\right) \right)}}{\left(0.049\,\mathrm{m}\right)}}
\]

\noindent\textbf{Resultado}

\[
\fitmath{Z=7.352166148\times10^{-1}}
\]

\noindent\textbf{Interpretación}
El resultado es un indicador de carga del modelo. No permite asignar una pérdida ni una eficiencia sin correlación contrastada.

\noindent\textbf{Comprobación de los pasos}
\begin{itemize}
\item Paso 2: $0.489796$ — reproducido
\item Paso 4: $0.735217$ — reproducido
\end{itemize}

\noindent\textbf{Comprobación de la respuesta} la aritmética del último paso reproduce la respuesta enunciada a la precisión con que está escrita. Es una reproducción aritmética; no valida el modelo físico ni sus hipótesis.

\noindent\textbf{Limitaciones}
\begin{itemize}
\item Modelo educativo; la aplicabilidad a una máquina real requiere revisión experta.
\end{itemize}

\noindent\textbf{Errores frecuentes}
\begin{itemize}
\item Usar cuerda real donde el parámetro pide cuerda axial.
\end{itemize}
% studium:end-record eyJpZCI6ICJQUkItMDAzMSIsICJzZWN0aW9uIjogImMwNCJ9
\end{tcolorbox}

\begin{tcolorbox}[title={Ejercicios propuestos}, studiumworked, breakable]

\noindent\textbf{1 (Básico, Desarrollo)}
\(C04-01-P \cdot\) Acción ideal. Ejercicio de transferencia. Una turbina de acción ideal recibe las componentes y velocidad periférica dadas. Halle remolino de salida, trabajo y descenso de temperatura total. Antes de consultar la solución, prediga el efecto de cambiar el primer dato respecto del caso resuelto de esta familia. Justifique la tendencia y explique qué restricción podría convertirse en limitante. Datos: \(C_\theta=492.03\, \mathrm{m}\,\mathrm{s}^{-1}\); \(U=211\, \mathrm{m}\,\mathrm{s}^{-1}\); \(c_p=1156.697\, \mathrm{J}\,\mathrm{kg}^{-1}\,\mathrm{K}^{-1}\).

\noindent\textbf{2 (Intermedio, Desarrollo)}
\(C04-02-P \cdot\) Óptimo de utilización. Ejercicio de transferencia. Para una velocidad de entrada al rotor y ángulo fijados, obtenga la velocidad periférica de máximo trabajo ideal y ese trabajo. Antes de consultar la solución, prediga el efecto de cambiar el primer dato respecto del caso resuelto de esta familia. Justifique la tendencia y explique qué restricción podría convertirse en limitante. Datos: \(C_e=628.35\, \mathrm{m}\,\mathrm{s}^{-1}\); \(\alpha_e=1.117011\, \mathrm{rad}\).

\noindent\textbf{3 (Avanzado, Desarrollo)}
\(C04-03-P \cdot\) Rotor con reducción de velocidad relativa. Ejercicio de transferencia. Un rotor recibe W tangencial de entrada positiva y descarga con dirección relativa simétrica y módulo reducido por k. Halle la potencia extraída. Antes de consultar la solución, prediga el efecto de cambiar el primer dato respecto del caso resuelto de esta familia. Justifique la tendencia y explique qué restricción podría convertirse en limitante. Datos: \(C_x=187.44\, \mathrm{m}\,\mathrm{s}^{-1}\); \(C_\theta=493\, \mathrm{m}\,\mathrm{s}^{-1}\); \(U=238\, \mathrm{m}\,\mathrm{s}^{-1}\); \(k_W=0.92\); \(\dot m=17.8\, \mathrm{kg}\,\mathrm{s}^{-1}\).

\noindent\textbf{4 (Nivel de examen, Desarrollo)}
\(C04-04-P \cdot\) Carga de cascada de turbina. Ejercicio de transferencia. Estime el parámetro de carga Z de una fila con ángulos algebraicos dados, en el modelo de baja compresibilidad definido en el capítulo. Antes de consultar la solución, prediga el efecto de cambiar el primer dato respecto del caso resuelto de esta familia. Justifique la tendencia y explique qué restricción podría convertirse en limitante. Datos: \(\alpha_e=0.8736246\, \mathrm{rad}\); \(\alpha_s=-1.012291\, \mathrm{rad}\); \(s=0.024\, \mathrm{m}\); \(c_x=0.049\, \mathrm{m}\).
\end{tcolorbox}


\studiumchaptercheck

\chapter{Turbinas de reacción y rendimiento del escalón}
\label{chap:c05}

\studiumchapterguide{Objetivo: desarrollar turbinas de reacción y rendimiento del escalón, declarar hipótesis y comprobar condiciones físicas.}
% studium:end-record eyJpZCI6ICJQQVItMDAyOSIsICJzZWN0aW9uIjogImMwNSJ9

\section{Introducción y objetivos}

Una turbina de reacción reparte la caída estática entre estator y rotor. El rotor acelera la corriente en su referencia y, simultáneamente, reduce su entalpía total absoluta al extraer trabajo. Comprender ambas afirmaciones requiere separar las referencias. El objetivo es diseñar los triángulos a partir de carga, flujo y reacción, y cerrar sus estados termodinámicos.
El capítulo pertenece al bloque II. No se presume que una reacción del cincuenta por ciento sea óptima para toda máquina. Su interés reside en la simetría que aparece bajo condiciones concretas y en el reparto de aceleraciones que permite estudiar.
% studium:end-record eyJpZCI6ICJQQVItMDAzMCIsICJzZWN0aW9uIjogImMwNSJ9

\section{Deducción de la reacción reducida}


Con radio constante, el rotor adiabático conserva \(h+W^2/2\). Por tanto



\[
\fitmath{h_2-h_3=\frac{W_3^2-W_2^2}{2}.}
\]
Al expandir las componentes y mantener \(C_x\) constante:

\[
\fitmath{h_2-h_3=(C_{\theta2}-C_{\theta3})\left[U-\frac{C_{\theta2}+C_{\theta3}}{2}\right].}
\]
Si además \(C_1^2=C_3^2\), el descenso estático del escalón coincide con el trabajo de Euler. Dividiendo se recupera la reacción reducida. Definiendo \(\psi_t=w_t/U^2\), se dispone de dos ecuaciones lineales:

\[
\fitmath{\frac{C_{\theta2}-C_{\theta3}}{U}=\psi_t,\qquad
\frac{C_{\theta2}+C_{\theta3}}{U}=2(1-\mathcal R_t).}
\]
Así,

\[
\fitmath{\frac{C_{\theta2}}U=1-\mathcal R_t+\frac{\psi_t}{2},\qquad
\frac{C_{\theta3}}U=1-\mathcal R_t-\frac{\psi_t}{2}.}
\]
No hay que resolver una ecuación trigonométrica con ramas ambiguas: se calculan primero las componentes. Después se obtiene \(\alpha_i=\arctan(C_{\theta i}/C_x)\) y \(\beta_i=\arctan[(C_{\theta i}-U)/C_x]\).
% studium:end-record eyJpZCI6ICJQQVItMDAzMSIsICJzZWN0aW9uIjogImMwNSJ9

\section{Simetría y condiciones de salida}

Para \(\mathcal R_t=1/2\), la suma tangencial es \(U\). Entonces \(W_{\theta2}=-C_{\theta3}\) y \(W_{\theta3}=-C_{\theta2}\). Con igual \(C_x\), los módulos se intercambian: \(W_2=C_3\) y \(W_3=C_2\). Esta es la simetría de un escalón ideal de reacción media; los álabes de rotor y estator no pasan a ser necesariamente intercambiables por geometría o solicitación.
Si se exige salida sin remolino, \(C_{\theta3}=0\), la carga queda ligada a la reacción por \(\psi_t=2(1-\mathcal R_t)\). En ese caso no son tres parámetros independientes. La figura~\ref{fig:reaccion} compara el reparto de caída estática para distintas reacciones. Si se mantiene la condición de descarga axial, el cambio de carga es una consecuencia del cierre de los triángulos.
Una reacción negativa indica subida estática en el rotor respecto de una caída estática positiva de escalón. No debe eliminarse cambiando un signo sin explicación. Puede ser consecuencia de la ley radial o de un diseño con difusión relativa desfavorable y requiere una evaluación aerodinámica específica.
% studium:end-record eyJpZCI6ICJQQVItMDAzMiIsICJzZWN0aW9uIjogImMwNSJ9

\par\vspace{\baselineskip}
\noindent\begin{minipage}{\linewidth}
\centering
\begingroup
\setbox0=\hbox{\begin{tikzpicture}
\foreach \x/\rr in {0/0,3/.5,6/.7}{\fill[orange!35] (\x,0) rectangle (\x+1.5,3*\rr);\fill[blue!25] (\x,3*\rr) rectangle (\x+1.5,3);\node[below] at (\x+.75,-.3) {$\mathcal R=\rr$};}\node[orange!80!black] at (3.7,1) {Rotor};\node[blue] at (3.7,2.4) {Estator};
\end{tikzpicture}}%
\count0=\ht0
\advance\count0 by \dp0
\count2=\wd0
\count4=\dimexpr\textheight/100*38\relax
\count6=\linewidth
\ifnum\count0<1
\count0=1
\fi
\ifnum\count2<1
\count2=1
\fi
\divide\count0 by 4096
\divide\count2 by 4096
\divide\count4 by 4096
\divide\count6 by 4096
\ifnum\count0<1
\count0=1
\fi
\ifnum\count2<1
\count2=1
\fi
\ifnum\count4<1
\count4=1
\fi
\ifnum\count6<1
\count6=1
\fi
\count8=\count2
\ifnum\count0>\count4
\count8=\numexpr\count2*\count4/\count0\relax
\fi
\ifnum\count8>\count6
\resizebox{\linewidth}{!}{\begin{tikzpicture}
\foreach \x/\rr in {0/0,3/.5,6/.7}{\fill[orange!35] (\x,0) rectangle (\x+1.5,3*\rr);\fill[blue!25] (\x,3*\rr) rectangle (\x+1.5,3);\node[below] at (\x+.75,-.3) {$\mathcal R=\rr$};}\node[orange!80!black] at (3.7,1) {Rotor};\node[blue] at (3.7,2.4) {Estator};
\end{tikzpicture}}%
\else
\ifnum\count0>\count4
\resizebox{!}{\dimexpr\textheight/100*38\relax}{\begin{tikzpicture}
\foreach \x/\rr in {0/0,3/.5,6/.7}{\fill[orange!35] (\x,0) rectangle (\x+1.5,3*\rr);\fill[blue!25] (\x,3*\rr) rectangle (\x+1.5,3);\node[below] at (\x+.75,-.3) {$\mathcal R=\rr$};}\node[orange!80!black] at (3.7,1) {Rotor};\node[blue] at (3.7,2.4) {Estator};
\end{tikzpicture}}%
\else
\begin{tikzpicture}
\foreach \x/\rr in {0/0,3/.5,6/.7}{\fill[orange!35] (\x,0) rectangle (\x+1.5,3*\rr);\fill[blue!25] (\x,3*\rr) rectangle (\x+1.5,3);\node[below] at (\x+.75,-.3) {$\mathcal R=\rr$};}\node[orange!80!black] at (3.7,1) {Rotor};\node[blue] at (3.7,2.4) {Estator};
\end{tikzpicture}%
\fi
\fi
\endgroup
\par\nopagebreak
\studiumcaption{fig:reaccion}{Caída estática ideal de igual magnitud total, repartida entre filas. No representa caída de presión.}
\end{minipage}
\par\vspace{\baselineskip}
% studium:end-record eyJpZCI6ICJGSUctMDAwNSIsICJzZWN0aW9uIjogImMwNSJ9

\section{Cierre termodinámico por estaciones}

Se parte de \(T_{01},p_{01}\), se determina \(T_{02}=T_{01}\) y se aplica la pérdida de presión total del estator con una definición explícita. Con \(C_2\),

\[
\fitmath{T_2=T_{02}-\frac{C_2^2}{2c_p},\qquad p_2=p_{02}\left(\frac{T_2}{T_{02}}\right)^{\gamma/(\gamma-1)}.}
\]
Euler proporciona \(T_{03}=T_{02}-w_t/c_p\). Para obtener \(p_{03}\) hace falta una eficiencia del escalón o un modelo de pérdida del rotor; no basta con conservar rotalpía. Usando eficiencia total-total:

\[
\fitmath{\frac{p_{03}}{p_{01}}=\left[1-\frac{w_t}{\eta_{t,\mathrm{tt}}c_pT_{01}}\right]^{\gamma/(\gamma-1)}.}
\]
La base debe ser positiva. A continuación se resta la energía cinética de salida para hallar \(T_3\) y se transforma localmente a \(p_3\). La continuidad produce el área en cada corte usando la densidad estática y \(C_x\). Una expansión tiende a aumentar el área necesaria si no crece suficientemente la componente axial.
No se debe imponer, además, una caída independiente de \(p_{0r}\) si ya se ha fijado la presión final mediante eficiencia global, salvo que se compruebe la compatibilidad. Un problema bien planteado tiene incógnitas y ecuaciones independientes; añadir pérdidas sin contar grados de libertad puede sobredefinirlo.
Caso completo original por estaciones: turbina axial. Se usa un radio medio fijo y velocidad axial constante. Datos: U=310 m/s, Cx=190 m/s, caudal 22 kg/s, radio medio 0.31 m, gamma=1.33 y cp=1156.69697 J/(kg K). Los remolinos absolutos en 1, 2 y 3 son -40, 450, -40 m/s. La presión total relativa del rotor de salida es 0,98 veces la de entrada, y la presión total absoluta del estator de salida es 0,985 veces la de entrada. Estas son pérdidas prescritas del caso educativo; no se añade además una eficiencia global independiente.
Procedimiento: se obtiene el trabajo con Euler y se conserva la entalpía total absoluta en el estator. Se resta la energía cinética absoluta para hallar cada temperatura estática. En el rotor se transforma al remanso relativo, se aplica su razón de pérdida y se transforma de vuelta al estado absoluto. Se aplica la pérdida absoluta del estator en la fila que corresponde. Las presiones estáticas se calculan desde el estado de remanso local, no mediante una isentropía ficticia de todo el escalón.

\[
\fitmath{\begin{array}{c|rrrr}\text{Est.}&T\,[\mathrm K]&T_0\,[\mathrm K]&p\,[\mathrm{kPa}]&p_0\,[\mathrm{kPa}]\\\hline1&1433.7&1450&764.375&800\\2&1346.86&1450&585.295&788\\3&1302.38&1318.68&500.98&526.727\\\end{array}}
\]

\[
\fitmath{\begin{array}{c|rrrrr}\text{Est.}&W_\theta\,[\mathrm{m/s}]&\alpha\,[{}^\circ]&\beta\,[{}^\circ]&T_{0r}\,[\mathrm K]&p_{0r}\,[\mathrm{kPa}]\\\hline1&-350&-11.8887&-61.5044&1502.26&922.71\\2&140&67.1094&36.3844&1370.94&628.62\\3&-350&-11.8887&-61.5044&1370.94&616.048\\\end{array}}
\]
En todos los cortes la componente axial es la prescrita. Las áreas se obtienen con la densidad estática y esa componente, y valen 0.0623309, 0.0764714, 0.0863911 metros cuadrados. Las alturas correspondientes son 32.0009, 39.2606, 44.3535 milímetros.
Resultados integrados: trabajo 151.9 kJ/kg, eficiencia total-total resultante 0.91945 y reacción estática 0.33871. La eficiencia se ha deducido después de las pérdidas de fila y no se ha impuesto como un dato adicional. La entropía aumenta y los estados permanecen positivos.
Comprobación: el salto de entalpía total coincide con Euler, el remanso relativo se conserva a través del rotor a radio fijo y cada área reconstruye el mismo caudal. Los residuos se guardan en el expediente; son inferiores a una millonésima en sus unidades coherentes. El cierre de balances no valida las razones de pérdida prescritas ni la aptitud aerodinámica de una geometría real.
% studium:end-record eyJpZCI6ICJQQVItMDAzMyIsICJzZWN0aW9uIjogImMwNSJ9

\section{Pérdidas, mezcla y evaluación del diseño}

La generación de entropía surge en capas límite, estelas, flujos secundarios, fugas y mezcla de refrigerante. Las curvas de eficiencia de una familia experimental dependen de Mach, Reynolds, holguras y método de promediado. Por ello no se digitalizan las curvas de los apuntes para presentarlas como resultados de una máquina original.
La evaluación preliminar compara trabajo, caudal, velocidad de punta, Mach de salida de cada fila y altura de pala. Una mayor carga reduce el número de escalones, pero puede elevar velocidad, giro y pérdidas. Una mayor reacción modifica las velocidades relativas y la temperatura de recuperación que ve el rotor. El mejor compromiso exige datos experimentales o simulación contrastada.
% studium:end-record eyJpZCI6ICJQQVItMDAzNCIsICJzZWN0aW9uIjogImMwNSJ9

\section{Errores y formulario}

Errores habituales: atribuir a una reacción de medio la mitad exacta de la caída de presión, confundir el descenso estático con el total y utilizar una eficiencia total-estática dentro de una fórmula total-total. Autocomprobación: dos escalones tienen la misma reacción pero distinto \(\psi_t\); construya sus triángulos y explique por qué su descarga no tiene por qué ser axial.
Formulario: \(\mathcal R_t=(h_2-h_3)/(h_1-h_3)\); las expresiones lineales de componentes solo corresponden al escalón de extremos con igual energía cinética. \(T_{03}=T_{01}-w_t/c_p\). Referencias: ULE Turbinas axiales, pp. \(4-8\) y \(13-16\); MIT 16.50, lección 27, pp. \(2-4\), para contraste de reacción. Las soluciones numéricas se recalculan independientemente.
% studium:end-record eyJpZCI6ICJQQVItMDAzNSIsICJzZWN0aW9uIjogImMwNSJ9

\begin{tcolorbox}[title={Desarrollos del modelo}, studiumdef, breakable]

\noindent\textbf{Modelo desarrollado: Triángulos desde carga y reacción}

\begin{equation}
W_{\theta 3}=U\left(-\frac{\psi}{2}-\mathcal R\right)
\label{eq:EC-C05}
\end{equation}

\noindent\textit{Significado.} Componente relativa tangencial a la salida del rotor. Es $W_{\theta 3}$, no $C_{\theta 3}$; el signo depende de $\psi$ y $\mathcal R$.

\noindent\textbf{Hipótesis}
\begin{itemize}
\item Radio y componente axial constantes; igual energía cinética en extremos del escalón.
\end{itemize}

\noindent\textbf{Principios}
\begin{itemize}
\item Balance y definiciones físicas desarrollados en el capítulo; propiedades y dominio declarados en el problema.
\end{itemize}

\noindent\textbf{Aplicabilidad}
\begin{itemize}
\item Modelo educativo del caso C05-01-R
\end{itemize}

\noindent\textbf{Limitaciones}
\begin{itemize}
\item Modelo educativo; la aplicabilidad a una máquina real requiere revisión experta.
\end{itemize}

\noindent\textbf{Pasos}
\begin{enumerate}
\item
La suma y diferencia de remolinos dan la entrada.

\[
\fitmath{C_{\theta2}=U \left(\frac{\psi}{2} - \mathcal R + 1\right)}
\]
\item
Se obtiene la descarga absoluta.

\[
\fitmath{C_{\theta3}=U \left(- \frac{\psi}{2} - \mathcal R + 1\right)}
\]
\item
La referencia relativa resta U a la descarga.

\[
\fitmath{W_{\theta 3}=U \left(- \frac{\psi}{2} - \mathcal R\right)}
\]
\end{enumerate}

{\noindent\footnotesize verificada dimensionalmente.\par}
\end{tcolorbox}

\begin{tcolorbox}[title={Problema resuelto 1}, studiumworked, breakable]

\noindent\textbf{Enunciado}
\(C05-01-R \cdot\) Triángulos desde carga y reacción. Diseñe las componentes del escalón repetitivo con carga, reacción y velocidad periférica dadas. Obtenga la componente relativa tangencial a la salida del rotor. Datos: \(U=305\, \mathrm{m}\,\mathrm{s}^{-1}\); \(\psi=1.42\); \(\mathcal R=0.46\).

\noindent\textbf{Datos}
\begin{itemize}
\item $U = 305\,\mathrm{m}\,\mathrm{s}^{-1}$ — velocidad periférica
\item $\psi = 1.42$ — carga turbina
\item $\mathcal R = 0.46$ — reacción
\end{itemize}

\noindent\textbf{Incógnita}
$W_{\theta 3}$, en $\mathrm{m\,s^{-1}}$; también las magnitudes intermedias desarrolladas.

\noindent\textbf{Modelo físico}

\noindent\textit{Triángulos desde carga y reacción}

\[
\fitmath{W_{\theta 3}=U \left(- \frac{\psi}{2} - \mathcal R\right)}
\]

\noindent\textbf{Hipótesis}
\begin{itemize}
\item Radio y componente axial constantes; igual energía cinética en extremos del escalón.
\end{itemize}

\noindent\textbf{Desarrollo paso a paso}

\begin{enumerate}

\item La suma y diferencia de remolinos dan la entrada.

\[\fitmath{ct2=\mathrm{uu}\cdot \left(1-\mathrm{rea}+\frac{\mathrm{psi}}{2}\right)}\]

\[\fitmath{ct2=(305)\cdot \left(1-(0.46)+\frac{(1.42)}{2}\right)\approx 381.25\,\mathrm{m/s}}\]

\item Se obtiene la descarga absoluta.

\[\fitmath{ct3=\mathrm{uu}\cdot \left(1-\mathrm{rea}-\frac{\mathrm{psi}}{2}\right)}\]

\[\fitmath{ct3=(305)\cdot \left(1-(0.46)-\frac{(1.42)}{2}\right)\approx -51.85\,\mathrm{m/s}}\]

\item La referencia relativa resta U a la descarga.

\[\fitmath{wt3=\mathrm{uu}\cdot \left(-\mathrm{rea}-\frac{\mathrm{psi}}{2}\right)}\]

\[\fitmath{wt3=(305)\cdot \left(-(0.46)-\frac{(1.42)}{2}\right)\approx -356.85\,\mathrm{m/s}}\]

\end{enumerate}


\noindent\textbf{Sustitución numérica}

\[
\fitmath{\left(305\,\mathrm{m}\,\mathrm{s}^{-1}\right) \left(- \frac{\left(1.42\right)}{2} - \left(0.46\right)\right)}
\]

\noindent\textbf{Resultado}

\[
\fitmath{W_{\theta 3}=-3.568500000\times10^{2}\,\mathrm{m}\,\mathrm{s}^{-1}}
\]

\noindent\textbf{Interpretación}
La componente relativa de salida es negativa; su módulo no debe sustituir su signo en la suma vectorial.

\noindent\textbf{Comprobación de los pasos}
\begin{itemize}
\item Paso 2: $381.25$ — reproducido
\item Paso 4: $-51.85$ — reproducido
\item Paso 6: $-356.85$ — reproducido
\end{itemize}

\noindent\textbf{Comprobación de la respuesta} la aritmética del último paso reproduce la respuesta enunciada a la precisión con que está escrita. Es una reproducción aritmética; no valida el modelo físico ni sus hipótesis.

\noindent\textbf{Limitaciones}
\begin{itemize}
\item Modelo educativo; la aplicabilidad a una máquina real requiere revisión experta.
\end{itemize}

\noindent\textbf{Errores frecuentes}
\begin{itemize}
\item Elegir carga, reacción y descarga axial como tres datos independientes.
\end{itemize}
% studium:end-record eyJpZCI6ICJQUkItMDAzMyIsICJzZWN0aW9uIjogImMwNSJ9
\end{tcolorbox}

\begin{tcolorbox}[title={Problema resuelto 2}, studiumworked, breakable]

\noindent\textbf{Enunciado}
\(C05-02-R \cdot\) Reacción general no repetitiva. Calcule la reacción estática de una turbina cuando las velocidades absolutas de entrada y salida del escalón difieren. Datos: \(U=287\, \mathrm{m}\,\mathrm{s}^{-1}\); \(C_{\theta2}=421\, \mathrm{m}\,\mathrm{s}^{-1}\); \(C_{\theta3}=-38\, \mathrm{m}\,\mathrm{s}^{-1}\); \(C_1=127\, \mathrm{m}\,\mathrm{s}^{-1}\); \(C_2=467\, \mathrm{m}\,\mathrm{s}^{-1}\); \(C_3=203\, \mathrm{m}\,\mathrm{s}^{-1}\).

\noindent\textbf{Datos}
\begin{itemize}
\item $U = 287\,\mathrm{m}\,\mathrm{s}^{-1}$ — velocidad periférica
\item $C_{\theta 2} = 421\,\mathrm{m}\,\mathrm{s}^{-1}$ — remolino entrada rotor
\item $C_{\theta 3} = -38\,\mathrm{m}\,\mathrm{s}^{-1}$ — remolino salida rotor
\item $c_{1} = 127\,\mathrm{m}\,\mathrm{s}^{-1}$ — módulo entrada escalón
\item $c_{2} = 467\,\mathrm{m}\,\mathrm{s}^{-1}$ — módulo entrada rotor
\item $c_{3} = 203\,\mathrm{m}\,\mathrm{s}^{-1}$ — módulo salida escalón
\end{itemize}

\noindent\textbf{Incógnita}
$\mathcal{R}$ (adimensional); también las magnitudes intermedias desarrolladas.

\noindent\textbf{Modelo físico}

\noindent\textit{Reacción general no repetitiva}

\[
\fitmath{\mathcal R=\frac{- \frac{c_{2}^{2}}{2} + \frac{c_{3}^{2}}{2} + U \left(C_{\theta 2} - C_{\theta 3}\right)}{- \frac{c_{1}^{2}}{2} + \frac{c_{3}^{2}}{2} + U \left(C_{\theta 2} - C_{\theta 3}\right)}}
\]

\noindent\textbf{Hipótesis}
\begin{itemize}
\item Rotor y estator adiabáticos; radios constantes para Euler.
\end{itemize}

\noindent\textbf{Desarrollo paso a paso}

\begin{enumerate}

\item Euler da el trabajo extraído.

\[\fitmath{\mathrm{wt}=\mathrm{uu}\cdot \left(\mathrm{ct}_{2}-\mathrm{ct}_{3}\right)}\]

\[\fitmath{\mathrm{wt}=(287)\cdot \left((421)-(-38)\right)\approx 131733\,\mathrm{J/kg}}\]

\item La caída estática del rotor incluye el cambio cinético.

\[\fitmath{\mathrm{dhr}=\mathrm{uu}\cdot \left(\mathrm{ct}_{2}-\mathrm{ct}_{3}\right)+\frac{{c_{3}}^{2}-{c_{2}}^{2}}{2}}\]

\[\fitmath{\mathrm{dhr}=(287)\cdot \left((421)-(-38)\right)+\frac{{(203)}^{2}-{(467)}^{2}}{2}\approx 43293\,\mathrm{J/kg}}\]

\item La caída estática del escalón tiene extremos 1 y 3.

\[\fitmath{\mathrm{rea}=\frac{\mathrm{uu}\cdot \left(\mathrm{ct}_{2}-\mathrm{ct}_{3}\right)+\frac{{c_{3}}^{2}-{c_{2}}^{2}}{2}}{\mathrm{uu}\cdot \left(\mathrm{ct}_{2}-\mathrm{ct}_{3}\right)+\frac{{c_{3}}^{2}-{c_{1}}^{2}}{2}}}\]

\[\fitmath{\mathrm{rea}=\frac{(287)\cdot \left((421)-(-38)\right)+\frac{{(203)}^{2}-{(467)}^{2}}{2}}{(287)\cdot \left((421)-(-38)\right)+\frac{{(203)}^{2}-{(127)}^{2}}{2}}\approx 0.3000769\,}\]

\end{enumerate}


\noindent\textbf{Sustitución numérica}

\[
\fitmath{\frac{- \frac{\left(\left(467\,\mathrm{m}\,\mathrm{s}^{-1}\right)\right)^{2}}{2} + \frac{\left(\left(203\,\mathrm{m}\,\mathrm{s}^{-1}\right)\right)^{2}}{2} + \left(287\,\mathrm{m}\,\mathrm{s}^{-1}\right) \left(\left(421\,\mathrm{m}\,\mathrm{s}^{-1}\right) - \left(-38\,\mathrm{m}\,\mathrm{s}^{-1}\right)\right)}{- \frac{\left(\left(127\,\mathrm{m}\,\mathrm{s}^{-1}\right)\right)^{2}}{2} + \frac{\left(\left(203\,\mathrm{m}\,\mathrm{s}^{-1}\right)\right)^{2}}{2} + \left(287\,\mathrm{m}\,\mathrm{s}^{-1}\right) \left(\left(421\,\mathrm{m}\,\mathrm{s}^{-1}\right) - \left(-38\,\mathrm{m}\,\mathrm{s}^{-1}\right)\right)}}
\]

\noindent\textbf{Resultado}

\[
\fitmath{\mathcal R=3.000769375\times10^{-1}}
\]

\noindent\textbf{Interpretación}
La reacción general difiere de la fórmula reducida si cambia la energía cinética entre extremos.

\noindent\textbf{Comprobación de los pasos}
\begin{itemize}
\item Paso 2: $131733$ — reproducido
\item Paso 4: $43293$ — reproducido
\item Paso 6: $0.300077$ — reproducido
\end{itemize}

\noindent\textbf{Comprobación de la respuesta} la aritmética del último paso reproduce la respuesta enunciada a la precisión con que está escrita. Es una reproducción aritmética; no valida el modelo físico ni sus hipótesis.

\noindent\textbf{Limitaciones}
\begin{itemize}
\item Modelo educativo; la aplicabilidad a una máquina real requiere revisión experta.
\end{itemize}

\noindent\textbf{Errores frecuentes}
\begin{itemize}
\item Usar caída total en el denominador de una reacción definida como estática.
\end{itemize}
% studium:end-record eyJpZCI6ICJQUkItMDAzNSIsICJzZWN0aW9uIjogImMwNSJ9
\end{tcolorbox}

\begin{tcolorbox}[title={Problema resuelto 3}, studiumworked, breakable]

\noindent\textbf{Enunciado}
\(C05-03-R \cdot\) Presión de descarga de turbina. Dados trabajo y rendimiento total-total de un escalón, halle la presión total de salida. Datos: \(T_{0e}=1510\, \mathrm{K}\); \(p_{0e}=1240000\, \mathrm{Pa}\); \(w_t=156000\, \mathrm{J}\,\mathrm{kg}^{-1}\); \(\eta=0.885\); \(c_p=1156.697\, \mathrm{J}\,\mathrm{kg}^{-1}\,\mathrm{K}^{-1}\); \(\gamma=1.33\).

\noindent\textbf{Datos}
\begin{itemize}
\item $T_{0e} = 1510\,\mathrm{K}$ — temperatura total entrada
\item $p_{0e} = 1240000\,\mathrm{Pa}$ — presión total entrada
\item $w_t = 156000\,\mathrm{J}\,\mathrm{kg}^{-1}$ — trabajo extraído
\item $\eta = 0.885$ — rendimiento total-total
\item $c_p = 1156.696969697\,\mathrm{J}\,\mathrm{kg}^{-1}\,\mathrm{K}^{-1}$ — calor específico
\item $\gamma = 1.33$ — razón de calores
\end{itemize}

\noindent\textbf{Incógnita}
$p_{0s}$, en $\mathrm{Pa}$; también las magnitudes intermedias desarrolladas.

\noindent\textbf{Modelo físico}

\noindent\textit{Presión de descarga de turbina}

\[
\fitmath{p_{0s}=p_{0e} \left(1 - \frac{w_t}{T_{0e} c_p \eta}\right)^{\frac{\gamma}{\gamma - 1}}}
\]

\noindent\textbf{Hipótesis}
\begin{itemize}
\item Sin refrigeración ni sangrado; gas perfecto constante.
\end{itemize}

\noindent\textbf{Desarrollo paso a paso}

\begin{enumerate}

\item Se determina la caída de temperatura ideal equivalente.

\[\fitmath{\mathrm{dTs}=\frac{\mathrm{wt}}{\eta\cdot \mathrm{cp}}}\]

\[\fitmath{\mathrm{dTs}=\frac{(156000)}{(0.885)\cdot (1156.69697)}\approx 152.3918\,\mathrm{K}}\]

\item Se calcula la razón de temperaturas isentrópicas.

\[\fitmath{\mathrm{rr}=1-\frac{\mathrm{wt}}{\eta\cdot \mathrm{cp}\cdot \mathrm{Tin}}}\]

\[\fitmath{\mathrm{rr}=1-\frac{(156000)}{(0.885)\cdot (1156.69697)\cdot (1510)}\approx 0.8990782\,}\]

\item La razón de presiones corresponde al estado ideal equivalente.

\[\fitmath{\mathrm{pout}=\mathrm{pin}\cdot {\left(1-\frac{\mathrm{wt}}{\eta\cdot \mathrm{cp}\cdot \mathrm{Tin}}\right)}^{\frac{\gamma}{\gamma-1}}}\]

\[\fitmath{\mathrm{pout}=(1240000)\cdot {\left(1-\frac{(156000)}{(0.885)\cdot (1156.69697)\cdot (1510)}\right)}^{\frac{(1.33)}{(1.33)-1}}\approx 807628.4\,\mathrm{Pa}}\]

\end{enumerate}


\noindent\textbf{Sustitución numérica}

\[
\fitmath{\left(1240000\,\mathrm{Pa}\right) \left(1 - \frac{\left(156000\,\mathrm{J}\,\mathrm{kg}^{-1}\right)}{\left(1510\,\mathrm{K}\right) \left(1156.696969697\,\mathrm{J}\,\mathrm{kg}^{-1}\,\mathrm{K}^{-1}\right) \left(0.885\right)}\right)^{\frac{\left(1.33\right)}{\left(1.33\right) - 1}}}
\]

\noindent\textbf{Resultado}

\[
\fitmath{p_{0s}=8.076283789\times10^{5}\,\mathrm{Pa}}
\]

\noindent\textbf{Interpretación}
La base debe ser positiva. La presión resultante es total absoluta, no estática.

\noindent\textbf{Comprobación de los pasos}
\begin{itemize}
\item Paso 2: $152.392$ — reproducido
\item Paso 4: $0.899078$ — reproducido
\item Paso 6: $807628$ — reproducido
\end{itemize}

\noindent\textbf{Comprobación de la respuesta} la aritmética del último paso reproduce la respuesta enunciada a la precisión con que está escrita. Es una reproducción aritmética; no valida el modelo físico ni sus hipótesis.

\noindent\textbf{Limitaciones}
\begin{itemize}
\item Modelo educativo; la aplicabilidad a una máquina real requiere revisión experta.
\end{itemize}

\noindent\textbf{Errores frecuentes}
\begin{itemize}
\item Introducir eficiencia total-estática en una expresión total-total.
\end{itemize}
% studium:end-record eyJpZCI6ICJQUkItMDAzNyIsICJzZWN0aW9uIjogImMwNSJ9
\end{tcolorbox}

\begin{tcolorbox}[title={Problema resuelto 4}, studiumworked, breakable]

\noindent\textbf{Enunciado}
\(C05-04-R \cdot\) Descarga axial y potencia. Una turbina repetitiva de reacción dada descarga sin remolino. Calcule su carga y potencia útil. Datos: \(\mathcal R=0.52\); \(U=322\, \mathrm{m}\,\mathrm{s}^{-1}\); \(\dot m=26\, \mathrm{kg}\,\mathrm{s}^{-1}\); \(\eta_m=0.975\).

\noindent\textbf{Datos}
\begin{itemize}
\item $\mathcal R = 0.52$ — reacción
\item $U = 322\,\mathrm{m}\,\mathrm{s}^{-1}$ — velocidad periférica
\item $\dot m = 26\,\mathrm{kg}\,\mathrm{s}^{-1}$ — caudal
\item $\eta_m = 0.975$ — rendimiento mecánico
\end{itemize}

\noindent\textbf{Incógnita}
$\dot{W}$, en $\mathrm{W}$; también las magnitudes intermedias desarrolladas.

\noindent\textbf{Modelo físico}

\noindent\textit{Descarga axial y potencia}

\[
\fitmath{\dot{W}=2 \eta_m \dot m U^{2} \left(1 - \mathcal R\right)}
\]

\noindent\textbf{Hipótesis}
\begin{itemize}
\item Escalón repetitivo, radio y Cx constantes, eficiencia mecánica separada.
\end{itemize}

\noindent\textbf{Desarrollo paso a paso}

\begin{enumerate}

\item La descarga axial liga la carga a la reacción.

\[\fitmath{\psi=2\cdot \left(1-\mathrm{rea}\right)}\]

\[\fitmath{\psi=2\cdot \left(1-(0.52)\right)\approx 0.96\,}\]

\item El trabajo se obtiene de la carga.

\[\fitmath{\mathrm{wt}=2\cdot \left(1-\mathrm{rea}\right)\cdot {\mathrm{uu}}^{2}}\]

\[\fitmath{\mathrm{wt}=2\cdot \left(1-(0.52)\right)\cdot {(322)}^{2}\approx 99536.64\,\mathrm{J/kg}}\]

\item Se aplica una sola vez el rendimiento mecánico.

\[\fitmath{\mathrm{Pow}=\mathrm{em}\cdot \mathrm{md}\cdot 2\cdot \left(1-\mathrm{rea}\right)\cdot {\mathrm{uu}}^{2}}\]

\[\fitmath{\mathrm{Pow}=(0.975)\cdot (26)\cdot 2\cdot \left(1-(0.52)\right)\cdot {(322)}^{2}\approx 2523254\,\mathrm{W}}\]

\end{enumerate}


\noindent\textbf{Sustitución numérica}

\[
\fitmath{2 \left(0.975\right) \left(26\,\mathrm{kg}\,\mathrm{s}^{-1}\right) \left(\left(322\,\mathrm{m}\,\mathrm{s}^{-1}\right)\right)^{2} \left(1 - \left(0.52\right)\right)}
\]

\noindent\textbf{Resultado}

\[
\fitmath{\dot{W}=2.523253824\times10^{6}\,\mathrm{W}}
\]

\noindent\textbf{Interpretación}
La condición de descarga impide variar carga y reacción independientemente.

\noindent\textbf{Comprobación de los pasos}
\begin{itemize}
\item Paso 2: $0.96$ — reproducido
\item Paso 4: $99536.6$ — reproducido
\item Paso 6: $2.52325\times 10^{6}$ — reproducido
\end{itemize}

\noindent\textbf{Comprobación de la respuesta} la aritmética del último paso reproduce la respuesta enunciada a la precisión con que está escrita. Es una reproducción aritmética; no valida el modelo físico ni sus hipótesis.

\noindent\textbf{Limitaciones}
\begin{itemize}
\item Modelo educativo; la aplicabilidad a una máquina real requiere revisión experta.
\end{itemize}

\noindent\textbf{Errores frecuentes}
\begin{itemize}
\item Multiplicar por eficiencia mecánica dos veces.
\end{itemize}
% studium:end-record eyJpZCI6ICJQUkItMDAzOSIsICJzZWN0aW9uIjogImMwNSJ9
\end{tcolorbox}

\begin{tcolorbox}[title={Ejercicios propuestos}, studiumworked, breakable]

\noindent\textbf{1 (Básico, Desarrollo)}
\(C05-01-P \cdot\) Triángulos desde carga y reacción. Ejercicio de transferencia. Diseñe las componentes del escalón repetitivo con carga, reacción y velocidad periférica dadas. Obtenga la componente relativa tangencial a la salida del rotor. Antes de consultar la solución, prediga el efecto de cambiar el primer dato respecto del caso resuelto de esta familia. Justifique la tendencia y explique qué restricción podría convertirse en limitante. Datos: \(U=324.825\, \mathrm{m}\,\mathrm{s}^{-1}\); \(\psi=1.42\); \(\mathcal R=0.46\).

\noindent\textbf{2 (Intermedio, Desarrollo)}
\(C05-02-P \cdot\) Reacción general no repetitiva. Ejercicio de transferencia. Calcule la reacción estática de una turbina cuando las velocidades absolutas de entrada y salida del escalón difieren. Antes de consultar la solución, prediga el efecto de cambiar el primer dato respecto del caso resuelto de esta familia. Justifique la tendencia y explique qué restricción podría convertirse en limitante. Datos: \(U=305.655\, \mathrm{m}\,\mathrm{s}^{-1}\); \(C_{\theta2}=421\, \mathrm{m}\,\mathrm{s}^{-1}\); \(C_{\theta3}=-38\, \mathrm{m}\,\mathrm{s}^{-1}\); \(C_1=127\, \mathrm{m}\,\mathrm{s}^{-1}\); \(C_2=467\, \mathrm{m}\,\mathrm{s}^{-1}\); \(C_3=203\, \mathrm{m}\,\mathrm{s}^{-1}\).

\noindent\textbf{3 (Avanzado, Desarrollo)}
\(C05-03-P \cdot\) Presión de descarga de turbina. Ejercicio de transferencia. Dados trabajo y rendimiento total-total de un escalón, halle la presión total de salida. Antes de consultar la solución, prediga el efecto de cambiar el primer dato respecto del caso resuelto de esta familia. Justifique la tendencia y explique qué restricción podría convertirse en limitante. Datos: \(T_{0e}=1608.15\, \mathrm{K}\); \(p_{0e}=1240000\, \mathrm{Pa}\); \(w_t=156000\, \mathrm{J}\,\mathrm{kg}^{-1}\); \(\eta=0.885\); \(c_p=1156.697\, \mathrm{J}\,\mathrm{kg}^{-1}\,\mathrm{K}^{-1}\); \(\gamma=1.33\).

\noindent\textbf{4 (Nivel de examen, Desarrollo)}
\(C05-04-P \cdot\) Descarga axial y potencia. Ejercicio de transferencia. Una turbina repetitiva de reacción dada descarga sin remolino. Calcule su carga y potencia útil. Antes de consultar la solución, prediga el efecto de cambiar el primer dato respecto del caso resuelto de esta familia. Justifique la tendencia y explique qué restricción podría convertirse en limitante. Datos: \(\mathcal R=0.565\); \(U=322\, \mathrm{m}\,\mathrm{s}^{-1}\); \(\dot m=26\, \mathrm{kg}\,\mathrm{s}^{-1}\); \(\eta_m=0.975\).
\end{tcolorbox}


\studiumchaptercheck

\chapter{Caudal, bloqueo sónico y diseño de turbinas}
\label{chap:c06}

\studiumchapterguide{Objetivo: desarrollar caudal, bloqueo sónico y diseño de turbinas, declarar hipótesis y comprobar condiciones físicas.}
% studium:end-record eyJpZCI6ICJQQVItMDAzNiIsICJzZWN0aW9uIjogImMwNiJ9

\section{Introducción y objetivos}

Una turbina no puede aumentar indefinidamente el caudal reduciendo la presión de descarga. Cuando la garganta alcanza condiciones sónicas, su capacidad de paso queda limitada para unas condiciones totales de entrada dadas. El objetivo es deducir ese máximo, relacionarlo con área y orientación y combinarlo con trabajo y velocidad de giro en un diseño preliminar.
Se trata del bloque II. El modelo de garganta es unidimensional e isentrópico; su uso en una cascada real necesita área efectiva, coeficiente de descarga y estado total en la referencia adecuada. Se distinguirá entre bloqueo de una fila y limitación de un conjunto de escalones.
% studium:end-record eyJpZCI6ICJQQVItMDAzNyIsICJzZWN0aW9uIjogImMwNiJ9

\section{Función de gasto}


La continuidad normal a una sección y las relaciones de remanso dan



\[
\fitmath{\frac{\dot m}{A}=\frac{p_0}{\sqrt{T_0}}\sqrt{\frac{\gamma}{R_g}}\,M\left(1+\frac{\gamma-1}{2}M^2\right)^{-\frac{\gamma+1}{2(\gamma-1)}}.}
\]
La función que multiplica \(p_0/\sqrt{T_0}\) crece en la rama subsónica y decrece en la supersónica. Derivando su logaritmo:

\[
\fitmath{\frac{d\ln F}{dM}=\frac1M-\frac{(\gamma+1)M/2}{1+(\gamma-1)M^2/2}.}
\]
El cero positivo es \(M=1\). Para estados totales y área fijados, el gasto máximo ideal es

\[
\fitmath{\frac{\dot m_*}{A_*}=\frac{p_0}{\sqrt{T_0}}\sqrt{\frac{\gamma}{R_g}}\left(\frac{2}{\gamma+1}\right)^{\frac{\gamma+1}{2(\gamma-1)}}.}
\]
Las unidades son masa por área y tiempo. Un máximo matemático no prueba que la geometría alcance ese estado sin pérdidas o separación: describe la capacidad del modelo.
% studium:end-record eyJpZCI6ICJQQVItMDAzOCIsICJzZWN0aW9uIjogImMwNiJ9

\section{Presión crítica y orientación}

La razón de presión crítica local ideal es \(p_*/p_0=[2/(\gamma+1)]^{\gamma/(\gamma-1)}\). En una tobera convergente con condiciones de depósito y contrapresión prescritas, la reducción suficiente de contrapresión permite alcanzar el bloqueo. Una tobera convergente-divergente puede desarrollar además regiones supersónicas y ondas de choque; el área de salida y la contrapresión deciden qué solución es físicamente compatible.
El área de garganta se mide normal a la velocidad. El área anular frontal se mide normal al eje. Si la velocidad en el corte forma un ángulo \(\alpha\), \(C_n=C\cos\alpha\) respecto del área axial. Una corriente con gran remolino puede alcanzar Mach alto y mantener un Mach axial moderado. No se calcula el gasto de garganta multiplicando el área frontal por el módulo de velocidad sin la proyección correspondiente.
En el rotor la garganta está fija en la referencia móvil, por lo que se consideran velocidad y condiciones totales relativas. La rotalpía aporta la relación energética, pero el radio variable introduce el término centrífugo. El modelo de área crítica de una fila debe especificar qué radio y qué estado relativo se usan.
% studium:end-record eyJpZCI6ICJQQVItMDAzOSIsICJzZWN0aW9uIjogImMwNiJ9

\section{Diseño preliminar y número de escalones}

Dado el trabajo total requerido \(w_{t,tot}\), una velocidad periférica factible y una carga objetivo, una estimación inicial es \(n\approx w_{t,tot}/(\psi_tU^2)\). El entero elegido obliga a recalcular el trabajo de cada escalón, su temperatura y su relación de presión. No es correcto redondear el número y mantener inalterada la carga si el trabajo total debe cerrarse exactamente.
El procedimiento comienza por balance de potencia con el compresor, sigue por velocidad de giro y radios y termina por reparto de carga, reacción y área. Después se comprueba Mach, capacidad de garganta, tensiones y refrigeración. La figura~\ref{fig:bloqueo} muestra la función de gasto calculada; la rama supersónica se incluye como solución del modelo, sin sugerir que toda fila opere en ella.
La condición de un escalón afecta a los siguientes. Si una fila está bloqueada, el caudal corregido puede permanecer casi fijo mientras cambia la distribución de presiones aguas abajo. El grupo debe satisfacer continuidad, balance energético y compatibilidad de velocidades de eje. Un cálculo fila por fila sin esa compatibilidad puede producir una colección de estados individualmente posibles pero una máquina imposible.
% studium:end-record eyJpZCI6ICJQQVItMDA0MCIsICJzZWN0aW9uIjogImMwNiJ9

\par\vspace{\baselineskip}
\noindent\begin{minipage}{\linewidth}
\centering
\begingroup
\setbox0=\hbox{\begin{tikzpicture}
\draw[->,gray] (0,0)--(7,0);\draw[->,gray] (0,0)--(0,3.8);\draw (0.0,0)--(0.0,-.08) node[below,font=\scriptsize] {0};\draw (0,0.0)--(-.08,0.0) node[left,font=\scriptsize] {0};\draw (1.5,0)--(1.5,-.08) node[below,font=\scriptsize] {0.625};\draw (0,0.75)--(-.08,0.75) node[left,font=\scriptsize] {0.25};\draw (3.0,0)--(3.0,-.08) node[below,font=\scriptsize] {1.25};\draw (0,1.5)--(-.08,1.5) node[left,font=\scriptsize] {0.5};\draw (4.5,0)--(4.5,-.08) node[below,font=\scriptsize] {1.88};\draw (0,2.25)--(-.08,2.25) node[left,font=\scriptsize] {0.75};\draw (6.0,0)--(6.0,-.08) node[below,font=\scriptsize] {2.5};\draw (0,3.0)--(-.08,3.0) node[left,font=\scriptsize] {1};\draw[blue,very thick] plot coordinates {(0.00000,0.00000) (0.10000,0.21410) (0.20000,0.42691) (0.30000,0.63714) (0.40000,0.84355) (0.50000,1.04494) (0.60000,1.24017) (0.70000,1.42820) (0.80000,1.60804) (0.90000,1.77883) (1.00000,1.93981) (1.10000,2.09031) (1.20000,2.22981) (1.30000,2.35787) (1.40000,2.47421) (1.50000,2.57862) (1.60000,2.67103) (1.70000,2.75146) (1.80000,2.82004) (1.90000,2.87699) (2.00000,2.92260) (2.10000,2.95724) (2.20000,2.98136) (2.30000,2.99543) (2.40000,3.00000) (2.50000,2.99563) (2.60000,2.98291) (2.70000,2.96246) (2.80000,2.93490) (2.90000,2.90085) (3.00000,2.86093) (3.10000,2.81576) (3.20000,2.76591) (3.30000,2.71198) (3.40000,2.65452) (3.50000,2.59404) (3.60000,2.53104) (3.70000,2.46601) (3.80000,2.39936) (3.90000,2.33152) (4.00000,2.26285) (4.10000,2.19370) (4.20000,2.12439) (4.30000,2.05519) (4.40000,1.98636) (4.50000,1.91813) (4.60000,1.85070) (4.70000,1.78424) (4.80000,1.71891) (4.90000,1.65484) (5.00000,1.59215) (5.10000,1.53091) (5.20000,1.47122) (5.30000,1.41313) (5.40000,1.35668) (5.50000,1.30191) (5.60000,1.24885) (5.70000,1.19749) (5.80000,1.14785) (5.90000,1.09993) (6.00000,1.05369)};\node[below,align=center] at (3.5,-.65) {Mach};\node[rotate=90] at (-1,2) {$F(M)/F(1)$};
\end{tikzpicture}}%
\count0=\ht0
\advance\count0 by \dp0
\count2=\wd0
\count4=\dimexpr\textheight/100*38\relax
\count6=\linewidth
\ifnum\count0<1
\count0=1
\fi
\ifnum\count2<1
\count2=1
\fi
\divide\count0 by 4096
\divide\count2 by 4096
\divide\count4 by 4096
\divide\count6 by 4096
\ifnum\count0<1
\count0=1
\fi
\ifnum\count2<1
\count2=1
\fi
\ifnum\count4<1
\count4=1
\fi
\ifnum\count6<1
\count6=1
\fi
\count8=\count2
\ifnum\count0>\count4
\count8=\numexpr\count2*\count4/\count0\relax
\fi
\ifnum\count8>\count6
\resizebox{\linewidth}{!}{\begin{tikzpicture}
\draw[->,gray] (0,0)--(7,0);\draw[->,gray] (0,0)--(0,3.8);\draw (0.0,0)--(0.0,-.08) node[below,font=\scriptsize] {0};\draw (0,0.0)--(-.08,0.0) node[left,font=\scriptsize] {0};\draw (1.5,0)--(1.5,-.08) node[below,font=\scriptsize] {0.625};\draw (0,0.75)--(-.08,0.75) node[left,font=\scriptsize] {0.25};\draw (3.0,0)--(3.0,-.08) node[below,font=\scriptsize] {1.25};\draw (0,1.5)--(-.08,1.5) node[left,font=\scriptsize] {0.5};\draw (4.5,0)--(4.5,-.08) node[below,font=\scriptsize] {1.88};\draw (0,2.25)--(-.08,2.25) node[left,font=\scriptsize] {0.75};\draw (6.0,0)--(6.0,-.08) node[below,font=\scriptsize] {2.5};\draw (0,3.0)--(-.08,3.0) node[left,font=\scriptsize] {1};\draw[blue,very thick] plot coordinates {(0.00000,0.00000) (0.10000,0.21410) (0.20000,0.42691) (0.30000,0.63714) (0.40000,0.84355) (0.50000,1.04494) (0.60000,1.24017) (0.70000,1.42820) (0.80000,1.60804) (0.90000,1.77883) (1.00000,1.93981) (1.10000,2.09031) (1.20000,2.22981) (1.30000,2.35787) (1.40000,2.47421) (1.50000,2.57862) (1.60000,2.67103) (1.70000,2.75146) (1.80000,2.82004) (1.90000,2.87699) (2.00000,2.92260) (2.10000,2.95724) (2.20000,2.98136) (2.30000,2.99543) (2.40000,3.00000) (2.50000,2.99563) (2.60000,2.98291) (2.70000,2.96246) (2.80000,2.93490) (2.90000,2.90085) (3.00000,2.86093) (3.10000,2.81576) (3.20000,2.76591) (3.30000,2.71198) (3.40000,2.65452) (3.50000,2.59404) (3.60000,2.53104) (3.70000,2.46601) (3.80000,2.39936) (3.90000,2.33152) (4.00000,2.26285) (4.10000,2.19370) (4.20000,2.12439) (4.30000,2.05519) (4.40000,1.98636) (4.50000,1.91813) (4.60000,1.85070) (4.70000,1.78424) (4.80000,1.71891) (4.90000,1.65484) (5.00000,1.59215) (5.10000,1.53091) (5.20000,1.47122) (5.30000,1.41313) (5.40000,1.35668) (5.50000,1.30191) (5.60000,1.24885) (5.70000,1.19749) (5.80000,1.14785) (5.90000,1.09993) (6.00000,1.05369)};\node[below,align=center] at (3.5,-.65) {Mach};\node[rotate=90] at (-1,2) {$F(M)/F(1)$};
\end{tikzpicture}}%
\else
\ifnum\count0>\count4
\resizebox{!}{\dimexpr\textheight/100*38\relax}{\begin{tikzpicture}
\draw[->,gray] (0,0)--(7,0);\draw[->,gray] (0,0)--(0,3.8);\draw (0.0,0)--(0.0,-.08) node[below,font=\scriptsize] {0};\draw (0,0.0)--(-.08,0.0) node[left,font=\scriptsize] {0};\draw (1.5,0)--(1.5,-.08) node[below,font=\scriptsize] {0.625};\draw (0,0.75)--(-.08,0.75) node[left,font=\scriptsize] {0.25};\draw (3.0,0)--(3.0,-.08) node[below,font=\scriptsize] {1.25};\draw (0,1.5)--(-.08,1.5) node[left,font=\scriptsize] {0.5};\draw (4.5,0)--(4.5,-.08) node[below,font=\scriptsize] {1.88};\draw (0,2.25)--(-.08,2.25) node[left,font=\scriptsize] {0.75};\draw (6.0,0)--(6.0,-.08) node[below,font=\scriptsize] {2.5};\draw (0,3.0)--(-.08,3.0) node[left,font=\scriptsize] {1};\draw[blue,very thick] plot coordinates {(0.00000,0.00000) (0.10000,0.21410) (0.20000,0.42691) (0.30000,0.63714) (0.40000,0.84355) (0.50000,1.04494) (0.60000,1.24017) (0.70000,1.42820) (0.80000,1.60804) (0.90000,1.77883) (1.00000,1.93981) (1.10000,2.09031) (1.20000,2.22981) (1.30000,2.35787) (1.40000,2.47421) (1.50000,2.57862) (1.60000,2.67103) (1.70000,2.75146) (1.80000,2.82004) (1.90000,2.87699) (2.00000,2.92260) (2.10000,2.95724) (2.20000,2.98136) (2.30000,2.99543) (2.40000,3.00000) (2.50000,2.99563) (2.60000,2.98291) (2.70000,2.96246) (2.80000,2.93490) (2.90000,2.90085) (3.00000,2.86093) (3.10000,2.81576) (3.20000,2.76591) (3.30000,2.71198) (3.40000,2.65452) (3.50000,2.59404) (3.60000,2.53104) (3.70000,2.46601) (3.80000,2.39936) (3.90000,2.33152) (4.00000,2.26285) (4.10000,2.19370) (4.20000,2.12439) (4.30000,2.05519) (4.40000,1.98636) (4.50000,1.91813) (4.60000,1.85070) (4.70000,1.78424) (4.80000,1.71891) (4.90000,1.65484) (5.00000,1.59215) (5.10000,1.53091) (5.20000,1.47122) (5.30000,1.41313) (5.40000,1.35668) (5.50000,1.30191) (5.60000,1.24885) (5.70000,1.19749) (5.80000,1.14785) (5.90000,1.09993) (6.00000,1.05369)};\node[below,align=center] at (3.5,-.65) {Mach};\node[rotate=90] at (-1,2) {$F(M)/F(1)$};
\end{tikzpicture}}%
\else
\begin{tikzpicture}
\draw[->,gray] (0,0)--(7,0);\draw[->,gray] (0,0)--(0,3.8);\draw (0.0,0)--(0.0,-.08) node[below,font=\scriptsize] {0};\draw (0,0.0)--(-.08,0.0) node[left,font=\scriptsize] {0};\draw (1.5,0)--(1.5,-.08) node[below,font=\scriptsize] {0.625};\draw (0,0.75)--(-.08,0.75) node[left,font=\scriptsize] {0.25};\draw (3.0,0)--(3.0,-.08) node[below,font=\scriptsize] {1.25};\draw (0,1.5)--(-.08,1.5) node[left,font=\scriptsize] {0.5};\draw (4.5,0)--(4.5,-.08) node[below,font=\scriptsize] {1.88};\draw (0,2.25)--(-.08,2.25) node[left,font=\scriptsize] {0.75};\draw (6.0,0)--(6.0,-.08) node[below,font=\scriptsize] {2.5};\draw (0,3.0)--(-.08,3.0) node[left,font=\scriptsize] {1};\draw[blue,very thick] plot coordinates {(0.00000,0.00000) (0.10000,0.21410) (0.20000,0.42691) (0.30000,0.63714) (0.40000,0.84355) (0.50000,1.04494) (0.60000,1.24017) (0.70000,1.42820) (0.80000,1.60804) (0.90000,1.77883) (1.00000,1.93981) (1.10000,2.09031) (1.20000,2.22981) (1.30000,2.35787) (1.40000,2.47421) (1.50000,2.57862) (1.60000,2.67103) (1.70000,2.75146) (1.80000,2.82004) (1.90000,2.87699) (2.00000,2.92260) (2.10000,2.95724) (2.20000,2.98136) (2.30000,2.99543) (2.40000,3.00000) (2.50000,2.99563) (2.60000,2.98291) (2.70000,2.96246) (2.80000,2.93490) (2.90000,2.90085) (3.00000,2.86093) (3.10000,2.81576) (3.20000,2.76591) (3.30000,2.71198) (3.40000,2.65452) (3.50000,2.59404) (3.60000,2.53104) (3.70000,2.46601) (3.80000,2.39936) (3.90000,2.33152) (4.00000,2.26285) (4.10000,2.19370) (4.20000,2.12439) (4.30000,2.05519) (4.40000,1.98636) (4.50000,1.91813) (4.60000,1.85070) (4.70000,1.78424) (4.80000,1.71891) (4.90000,1.65484) (5.00000,1.59215) (5.10000,1.53091) (5.20000,1.47122) (5.30000,1.41313) (5.40000,1.35668) (5.50000,1.30191) (5.60000,1.24885) (5.70000,1.19749) (5.80000,1.14785) (5.90000,1.09993) (6.00000,1.05369)};\node[below,align=center] at (3.5,-.65) {Mach};\node[rotate=90] at (-1,2) {$F(M)/F(1)$};
\end{tikzpicture}%
\fi
\fi
\endgroup
\par\nopagebreak
\studiumcaption{fig:bloqueo}{Función de gasto ideal normalizada, gamma 1,33. Máximo en Mach uno; la rama supersónica exige geometría compatible.}
\end{minipage}
\par\vspace{\baselineskip}
% studium:end-record eyJpZCI6ICJGSUctMDAwNiIsICJzZWN0aW9uIjogImMwNiJ9

\section{Compromiso con altura y solicitación}

Para caudal fijado, un gran ángulo de descarga reduce la proyección axial y puede exigir mayor área anular. A igual radio medio, aumenta la altura de pala. Esa altura incrementa la tracción centrífuga y puede reducir el margen mecánico. El criterio de trabajo por escalón debe evaluarse junto con la capacidad de paso, no de forma aislada.
La magnitud \(AN^2\) combina área anular y velocidad de giro; aparece al relacionar la tensión centrífuga de una pala uniforme con los radios. Es un indicador preliminar, no un límite material universal. Su valor admisible depende de densidad del material, temperatura, geometría de raíz, vida y margen requerido.
% studium:end-record eyJpZCI6ICJQQVItMDA0MSIsICJzZWN0aW9uIjogImMwNiJ9

\section{Errores y formulario}

Errores habituales: declarar bloqueo porque el Mach absoluto es uno aunque la garganta sea móvil, usar la misma presión total a través de una fila real con pérdidas y extrapolar la función ideal a un flujo separado. Autocomprobación: a igual \(p_0\) y área crítica, explique por qué aumentar \(T_0\) reduce el caudal máximo de masa y no necesariamente el caudal volumétrico.
Formulario: función de gasto anterior, \(M_*=1\), \(A_{an}=\pi(r_t^2-r_h^2)\) y \(w_{t,tot}=\sum_jw_{tj}\). Referencias: ULE Turbinas axiales, pp. 2 y \(9-18\); MIT 16.50, lección 28, contexto del vínculo entre caudal, geometría y esfuerzos. La deducción de la función de gasto y las curvas son originales para este manual.
% studium:end-record eyJpZCI6ICJQQVItMDA0MiIsICJzZWN0aW9uIjogImMwNiJ9

\begin{tcolorbox}[title={Desarrollos del modelo}, studiumdef, breakable]

\noindent\textbf{Modelo desarrollado: Área crítica}

\begin{equation}
A^*=\frac{\dot m\,\sqrt{T_{0e}}}{p_{0e}}\left(\frac{2}{\gamma+1}\right)^{-\frac{\gamma+1}{2(\gamma-1)}}\sqrt{\frac{R_g}{\gamma}}
\label{eq:EC-C06}
\end{equation}

\noindent\textit{Significado.} Área de garganta sónica ideal para gas perfecto con las condiciones de entrada declaradas.

\noindent\textbf{Hipótesis}
\begin{itemize}
\item Garganta ideal unidimensional isentrópica en la referencia indicada.
\end{itemize}

\noindent\textbf{Principios}
\begin{itemize}
\item Balance y definiciones físicas desarrollados en el capítulo; propiedades y dominio declarados en el problema.
\end{itemize}

\noindent\textbf{Aplicabilidad}
\begin{itemize}
\item Modelo educativo del caso C06-01-R
\end{itemize}

\noindent\textbf{Limitaciones}
\begin{itemize}
\item Modelo educativo; la aplicabilidad a una máquina real requiere revisión experta.
\end{itemize}

\noindent\textbf{Pasos}
\begin{enumerate}
\item
Se evalúa el cuadrado del factor crítico a M=1.

\[
\fitmath{F=\frac{\gamma \left(\frac{2}{\gamma + 1}\right)^{\frac{\gamma + 1}{\gamma - 1}}}{R_g}}
\]
\item
Se calcula el flujo másico máximo por área.

\[
\fitmath{G^{*}=\frac{p_{0e} \left(\frac{2}{\gamma + 1}\right)^{\frac{\gamma + 1}{2 \gamma - 2}} \sqrt{\frac{\gamma}{R_g}}}{\sqrt{T_{0e}}}}
\]
\item
Se dimensiona el área normal efectiva.

\[
\fitmath{A^{*}=\frac{\sqrt{T_{0e}} \dot m \left(\frac{2}{\gamma + 1}\right)^{- \frac{\gamma + 1}{2 \gamma - 2}}}{p_{0e} \sqrt{\frac{\gamma}{R_g}}}}
\]
\end{enumerate}

{\noindent\footnotesize verificada dimensionalmente.\par}
\end{tcolorbox}

\begin{tcolorbox}[title={Problema resuelto 1}, studiumworked, breakable]

\noindent\textbf{Enunciado}
\(C06-01-R \cdot\) Área crítica. Determine el gasto máximo por área y el área crítica necesaria para el caudal dado. Datos: \(p_{0e}=680000\, \mathrm{Pa}\); \(T_{0e}=1320\, \mathrm{K}\); \(\gamma=1.33\); \(R_g=287\, \mathrm{J}\,\mathrm{kg}^{-1}\,\mathrm{K}^{-1}\); \(\dot m=24\, \mathrm{kg}\,\mathrm{s}^{-1}\).

\noindent\textbf{Datos}
\begin{itemize}
\item $p_{0e} = 680000\,\mathrm{Pa}$ — presión total garganta de referencia
\item $T_{0e} = 1320\,\mathrm{K}$ — temperatura total
\item $\gamma = 1.33$ — razón de calores
\item $R_g = 287\,\mathrm{J}\,\mathrm{kg}^{-1}\,\mathrm{K}^{-1}$ — constante del gas
\item $\dot m = 24\,\mathrm{kg}\,\mathrm{s}^{-1}$ — caudal
\end{itemize}

\noindent\textbf{Incógnita}
$A^{*}$, en $\mathrm{m^{2}}$; también las magnitudes intermedias desarrolladas.

\noindent\textbf{Modelo físico}

\noindent\textit{Área crítica}

\[
\fitmath{A^{*}=\frac{\sqrt{T_{0e}} \dot m \left(\frac{2}{\gamma + 1}\right)^{- \frac{\gamma + 1}{2 \gamma - 2}}}{p_{0e} \sqrt{\frac{\gamma}{R_g}}}}
\]

\noindent\textbf{Hipótesis}
\begin{itemize}
\item Garganta ideal unidimensional isentrópica en la referencia indicada.
\end{itemize}

\noindent\textbf{Desarrollo paso a paso}

\begin{enumerate}

\item Se evalúa el cuadrado del factor crítico a M=1.

\[\fitmath{F=\frac{\gamma}{R_g}\cdot {\left(\frac{2}{\gamma+1}\right)}^{\frac{\gamma+1}{\gamma-1}}}\]

\[\fitmath{F=\frac{(1.33)}{(287)}\cdot {\left(\frac{2}{(1.33)+1}\right)}^{\frac{(1.33)+1}{(1.33)-1}}\approx 0.001576408\,\mathrm{s^{2}\,K/m^{2}}}\]

\item Se calcula el flujo másico máximo por área.

\[\fitmath{\mathrm{flux}=\frac{\mathrm{pin}}{\sqrt{\mathrm{Tin}}}\cdot \sqrt{\frac{\gamma}{R_g}}\cdot {\left(\frac{2}{\gamma+1}\right)}^{\frac{\gamma+1}{2\cdot \left(\gamma-1\right)}}}\]

\[\fitmath{\mathrm{flux}=\frac{(680000)}{\sqrt{(1320)}}\cdot \sqrt{\frac{(1.33)}{(287)}}\cdot {\left(\frac{2}{(1.33)+1}\right)}^{\frac{(1.33)+1}{2\cdot \left((1.33)-1\right)}}\approx 743.1153\,\mathrm{kg/(m^{2}\,s)}}\]

\item Se dimensiona el área normal efectiva.

\[\fitmath{\mathrm{Ast}=\frac{\mathrm{md}}{\frac{\mathrm{pin}}{\sqrt{\mathrm{Tin}}}\cdot \sqrt{\frac{\gamma}{R_g}}\cdot {\left(\frac{2}{\gamma+1}\right)}^{\frac{\gamma+1}{2\cdot \left(\gamma-1\right)}}}}\]

\[\fitmath{\mathrm{Ast}=\frac{(24)}{\frac{(680000)}{\sqrt{(1320)}}\cdot \sqrt{\frac{(1.33)}{(287)}}\cdot {\left(\frac{2}{(1.33)+1}\right)}^{\frac{(1.33)+1}{2\cdot \left((1.33)-1\right)}}}\approx 0.03229647\,\mathrm{m^{2}}}\]

\end{enumerate}


\noindent\textbf{Sustitución numérica}

\[
\fitmath{\frac{\sqrt{\left(1320\,\mathrm{K}\right)} \left(24\,\mathrm{kg}\,\mathrm{s}^{-1}\right) \left(\frac{2}{\left(1.33\right) + 1}\right)^{- \frac{\left(1.33\right) + 1}{2 \left(1.33\right) - 2}}}{\left(680000\,\mathrm{Pa}\right) \sqrt{\frac{\left(1.33\right)}{\left(287\,\mathrm{J}\,\mathrm{kg}^{-1}\,\mathrm{K}^{-1}\right)}}}}
\]

\noindent\textbf{Resultado}

\[
\fitmath{A^{*}=3.229646860\times10^{-2}\,\mathrm{m}^{2}}
\]

\noindent\textbf{Interpretación}
El área es normal al flujo y efectiva. Un coeficiente de descarga inferior a uno exigiría mayor área geométrica.

\noindent\textbf{Comprobación de los pasos}
\begin{itemize}
\item Paso 2: $0.00157641$ — reproducido
\item Paso 4: $743.115$ — reproducido
\item Paso 6: $0.0322965$ — reproducido
\end{itemize}

\noindent\textbf{Comprobación de la respuesta} la aritmética del último paso reproduce la respuesta enunciada a la precisión con que está escrita. Es una reproducción aritmética; no valida el modelo físico ni sus hipótesis.

\noindent\textbf{Limitaciones}
\begin{itemize}
\item Modelo educativo; la aplicabilidad a una máquina real requiere revisión experta.
\end{itemize}

\noindent\textbf{Errores frecuentes}
\begin{itemize}
\item Usar área frontal anular como garganta sin proyección.
\end{itemize}
% studium:end-record eyJpZCI6ICJQUkItMDA0MSIsICJzZWN0aW9uIjogImMwNiJ9
\end{tcolorbox}

\begin{tcolorbox}[title={Problema resuelto 2}, studiumworked, breakable]

\noindent\textbf{Enunciado}
\(C06-02-R \cdot\) Cociente crítico. Calcule la presión crítica estática de una tobera convergente ideal. Datos: \(p_{0e}=915000\, \mathrm{Pa}\); \(\gamma=1.32\).

\noindent\textbf{Datos}
\begin{itemize}
\item $p_{0e} = 915000\,\mathrm{Pa}$ — presión total depósito
\item $\gamma = 1.32$ — razón de calores
\end{itemize}

\noindent\textbf{Incógnita}
$p^{*}$, en $\mathrm{Pa}$; también las magnitudes intermedias desarrolladas.

\noindent\textbf{Modelo físico}

\noindent\textit{Cociente crítico}

\[
\fitmath{p^{*}=p_{0e} \left(\frac{2}{\gamma + 1}\right)^{\frac{\gamma}{\gamma - 1}}}
\]

\noindent\textbf{Hipótesis}
\begin{itemize}
\item Depósito y tobera convergente ideal; gas perfecto.
\end{itemize}

\noindent\textbf{Desarrollo paso a paso}

\begin{enumerate}

\item Se aplica la razón isentrópica a Mach uno.

\[\fitmath{\mathrm{ratio}={\left(\frac{2}{\gamma+1}\right)}^{\frac{\gamma}{\gamma-1}}}\]

\[\fitmath{\mathrm{ratio}={\left(\frac{2}{(1.32)+1}\right)}^{\frac{(1.32)}{(1.32)-1}}\approx 0.5421392\,}\]

\item Se obtiene la presión de garganta crítica.

\[\fitmath{\mathrm{pstar}=\mathrm{pin}\cdot {\left(\frac{2}{\gamma+1}\right)}^{\frac{\gamma}{\gamma-1}}}\]

\[\fitmath{\mathrm{pstar}=(915000)\cdot {\left(\frac{2}{(1.32)+1}\right)}^{\frac{(1.32)}{(1.32)-1}}\approx 496057.3\,\mathrm{Pa}}\]

\end{enumerate}


\noindent\textbf{Sustitución numérica}

\[
\fitmath{\left(915000\,\mathrm{Pa}\right) \left(\frac{2}{\left(1.32\right) + 1}\right)^{\frac{\left(1.32\right)}{\left(1.32\right) - 1}}}
\]

\noindent\textbf{Resultado}

\[
\fitmath{p^{*}=4.960573496\times10^{5}\,\mathrm{Pa}}
\]

\noindent\textbf{Interpretación}
Una contrapresión inferior permite bloqueo en la tobera ideal indicada; no define toda la expansión de una cascada real.

\noindent\textbf{Comprobación de los pasos}
\begin{itemize}
\item Paso 2: $0.542139$ — reproducido
\item Paso 4: $496057$ — reproducido
\end{itemize}

\noindent\textbf{Comprobación de la respuesta} la aritmética del último paso reproduce la respuesta enunciada a la precisión con que está escrita. Es una reproducción aritmética; no valida el modelo físico ni sus hipótesis.

\noindent\textbf{Limitaciones}
\begin{itemize}
\item Modelo educativo; la aplicabilidad a una máquina real requiere revisión experta.
\end{itemize}

\noindent\textbf{Errores frecuentes}
\begin{itemize}
\item Interpretar la razón crítica como pérdida de presión total.
\end{itemize}
% studium:end-record eyJpZCI6ICJQUkItMDA0MyIsICJzZWN0aW9uIjogImMwNiJ9
\end{tcolorbox}

\begin{tcolorbox}[title={Problema resuelto 3}, studiumworked, breakable]

\noindent\textbf{Enunciado}
\(C06-03-R \cdot\) Reparto entero de escalones. Se ha elegido un número entero de escalones de igual trabajo. Calcule la carga real que cierra exactamente el trabajo global. Datos: \(w_t=384000\, \mathrm{J}\,\mathrm{kg}^{-1}\); \(N_{\mathrm{esc}}=3\); \(U=338\, \mathrm{m}\,\mathrm{s}^{-1}\).

\noindent\textbf{Datos}
\begin{itemize}
\item $w_t = 384000\,\mathrm{J}\,\mathrm{kg}^{-1}$ — trabajo total
\item $N_{\mathrm{esc}} = 3$ — número de escalones
\item $U = 338\,\mathrm{m}\,\mathrm{s}^{-1}$ — velocidad periférica común de referencia
\end{itemize}

\noindent\textbf{Incógnita}
$\psi$ (adimensional); también las magnitudes intermedias desarrolladas.

\noindent\textbf{Modelo físico}

\noindent\textit{Reparto entero de escalones}

\[
\fitmath{\psi=\frac{w_t}{N_{\mathrm{esc}} U^{2}}}
\]

\noindent\textbf{Hipótesis}
\begin{itemize}
\item Trabajo igual por escalón y misma U de referencia.
\end{itemize}

\noindent\textbf{Desarrollo paso a paso}

\begin{enumerate}

\item El reparto conserva la suma del trabajo.

\[\fitmath{\mathrm{we}=\frac{\mathrm{wt}}{\mathrm{nn}}}\]

\[\fitmath{\mathrm{we}=\frac{(384000)}{(3)}\approx 128000\,\mathrm{J/kg}}\]

\item Se calcula la carga efectiva tras elegir el entero.

\[\fitmath{\mathrm{psi}=\frac{\mathrm{wt}}{\mathrm{nn}\cdot {\mathrm{uu}}^{2}}}\]

\[\fitmath{\mathrm{psi}=\frac{(384000)}{(3)\cdot {(338)}^{2}}\approx 1.120409\,}\]

\end{enumerate}


\noindent\textbf{Sustitución numérica}

\[
\fitmath{\frac{\left(384000\,\mathrm{J}\,\mathrm{kg}^{-1}\right)}{\left(3\right) \left(\left(338\,\mathrm{m}\,\mathrm{s}^{-1}\right)\right)^{2}}}
\]

\noindent\textbf{Resultado}

\[
\fitmath{\psi=1.120408949\times10^{0}}
\]

\noindent\textbf{Interpretación}
El entero elegido obliga a recalcular psi. La factibilidad aerodinámica se revisa después.

\noindent\textbf{Comprobación de los pasos}
\begin{itemize}
\item Paso 2: $128000$ — reproducido
\item Paso 4: $1.12041$ — reproducido
\end{itemize}

\noindent\textbf{Comprobación de la respuesta} la aritmética del último paso reproduce la respuesta enunciada a la precisión con que está escrita. Es una reproducción aritmética; no valida el modelo físico ni sus hipótesis.

\noindent\textbf{Limitaciones}
\begin{itemize}
\item Modelo educativo; la aplicabilidad a una máquina real requiere revisión experta.
\end{itemize}

\noindent\textbf{Errores frecuentes}
\begin{itemize}
\item Redondear n y conservar una carga previa incompatible.
\end{itemize}
% studium:end-record eyJpZCI6ICJQUkItMDA0NSIsICJzZWN0aW9uIjogImMwNiJ9
\end{tcolorbox}

\begin{tcolorbox}[title={Problema resuelto 4}, studiumworked, breakable]

\noindent\textbf{Enunciado}
\(C06-04-R \cdot\) Caudal de garganta oblicua. Una sección axial opera con módulo de velocidad sónico y ángulo dado. Halle el caudal usando su proyección axial. Datos: \(\rho=1.92\, \mathrm{kg}\,\mathrm{m}^{-3}\); \(C=626\, \mathrm{m}\,\mathrm{s}^{-1}\); \(\alpha_e=1.064651\, \mathrm{rad}\); \(A=0.045\, \mathrm{m}^{2}\).

\noindent\textbf{Datos}
\begin{itemize}
\item $\rho = 1.92\,\mathrm{kg}\,\mathrm{m}^{-3}$ — densidad estática
\item $C = 626\,\mathrm{m}\,\mathrm{s}^{-1}$ — módulo sónico local
\item $\alpha_e = 1.0646508437165\,\mathrm{rad}$ — ángulo con el eje
\item $A = 0.045\,\mathrm{m}^{2}$ — área anular frontal
\end{itemize}

\noindent\textbf{Incógnita}
$\dot{m}$, en $\mathrm{kg\,s^{-1}}$; también las magnitudes intermedias desarrolladas.

\noindent\textbf{Modelo físico}

\noindent\textit{Caudal de garganta oblicua}

\[
\fitmath{\dot m=A \rho C \cos{\left(\alpha_e \right)}}
\]

\noindent\textbf{Hipótesis}
\begin{itemize}
\item Corte axial con densidad y velocidad uniformes.
\end{itemize}

\noindent\textbf{Desarrollo paso a paso}

\begin{enumerate}

\item Se proyecta el módulo sobre la normal axial.

\[\fitmath{\mathrm{cx}=\mathrm{vel}\cdot \cos\left(\mathrm{ang}\right)}\]

\[\fitmath{\mathrm{cx}=(626)\cdot \cos\left((1.064650844)\right)\approx 303.4908\,\mathrm{m/s}}\]

\item La continuidad usa esa componente.

\[\fitmath{\mathrm{md}=\rho\cdot \mathrm{vel}\cdot \cos\left(\mathrm{ang}\right)\cdot \mathrm{area}}\]

\[\fitmath{\mathrm{md}=(1.92)\cdot (626)\cdot \cos\left((1.064650844)\right)\cdot (0.045)\approx 26.22161\,\mathrm{kg/s}}\]

\end{enumerate}


\noindent\textbf{Sustitución numérica}

\[
\fitmath{\left(0.045\,\mathrm{m}^{2}\right) \left(1.92\,\mathrm{kg}\,\mathrm{m}^{-3}\right) \left(626\,\mathrm{m}\,\mathrm{s}^{-1}\right) \cos{\left(\left(1.0646508437165\,\mathrm{rad}\right) \right)}}
\]

\noindent\textbf{Resultado}

\[
\fitmath{\dot m=2.622160704\times10^{1}\,\mathrm{kg}\,\mathrm{s}^{-1}}
\]

\noindent\textbf{Interpretación}
Mach axial y Mach del módulo son distintos. La garganta real puede tener otra orientación y área.

\noindent\textbf{Comprobación de los pasos}
\begin{itemize}
\item Paso 2: $303.491$ — reproducido
\item Paso 4: $26.2216$ — reproducido
\end{itemize}

\noindent\textbf{Comprobación de la respuesta} la aritmética del último paso reproduce la respuesta enunciada a la precisión con que está escrita. Es una reproducción aritmética; no valida el modelo físico ni sus hipótesis.

\noindent\textbf{Limitaciones}
\begin{itemize}
\item Modelo educativo; la aplicabilidad a una máquina real requiere revisión experta.
\end{itemize}

\noindent\textbf{Errores frecuentes}
\begin{itemize}
\item Omitir el coseno por estar el módulo a Mach uno.
\end{itemize}
% studium:end-record eyJpZCI6ICJQUkItMDA0NyIsICJzZWN0aW9uIjogImMwNiJ9
\end{tcolorbox}

\begin{tcolorbox}[title={Ejercicios propuestos}, studiumworked, breakable]

\noindent\textbf{1 (Básico, Desarrollo)}
\(C06-01-P \cdot\) Área crítica. Ejercicio de transferencia. Determine el gasto máximo por área y el área crítica necesaria para el caudal dado. Antes de consultar la solución, prediga el efecto de cambiar el primer dato respecto del caso resuelto de esta familia. Justifique la tendencia y explique qué restricción podría convertirse en limitante. Datos: \(p_{0e}=724200\, \mathrm{Pa}\); \(T_{0e}=1320\, \mathrm{K}\); \(\gamma=1.33\); \(R_g=287\, \mathrm{J}\,\mathrm{kg}^{-1}\,\mathrm{K}^{-1}\); \(\dot m=24\, \mathrm{kg}\,\mathrm{s}^{-1}\).

\noindent\textbf{2 (Intermedio, Desarrollo)}
\(C06-02-P \cdot\) Cociente crítico. Ejercicio de transferencia. Calcule la presión crítica estática de una tobera convergente ideal. Antes de consultar la solución, prediga el efecto de cambiar el primer dato respecto del caso resuelto de esta familia. Justifique la tendencia y explique qué restricción podría convertirse en limitante. Datos: \(p_{0e}=974475\, \mathrm{Pa}\); \(\gamma=1.32\).

\noindent\textbf{3 (Avanzado, Desarrollo)}
\(C06-03-P \cdot\) Reparto entero de escalones. Ejercicio de transferencia. Se ha elegido un número entero de escalones de igual trabajo. Calcule la carga real que cierra exactamente el trabajo global. Antes de consultar la solución, prediga el efecto de cambiar el primer dato respecto del caso resuelto de esta familia. Justifique la tendencia y explique qué restricción podría convertirse en limitante. Datos: \(w_t=408960\, \mathrm{J}\,\mathrm{kg}^{-1}\); \(N_{\mathrm{esc}}=3\); \(U=338\, \mathrm{m}\,\mathrm{s}^{-1}\).

\noindent\textbf{4 (Nivel de examen, Desarrollo)}
\(C06-04-P \cdot\) Caudal de garganta oblicua. Ejercicio de transferencia. Una sección axial opera con módulo de velocidad sónico y ángulo dado. Halle el caudal usando su proyección axial. Antes de consultar la solución, prediga el efecto de cambiar el primer dato respecto del caso resuelto de esta familia. Justifique la tendencia y explique qué restricción podría convertirse en limitante. Datos: \(\rho=2.0448\, \mathrm{kg}\,\mathrm{m}^{-3}\); \(C=626\, \mathrm{m}\,\mathrm{s}^{-1}\); \(\alpha_e=1.064651\, \mathrm{rad}\); \(A=0.045\, \mathrm{m}^{2}\).
\end{tcolorbox}


\studiumchaptercheck

\part{Estudio bidimensional de compresores axiales}

\chapter{Escalón de compresor axial al radio medio}
\label{chap:c07}

\studiumchapterguide{Objetivo: desarrollar escalón de compresor axial al radio medio, declarar hipótesis y comprobar condiciones físicas.}
% studium:end-record eyJpZCI6ICJQQVItMDA0MyIsICJzZWN0aW9uIjogImMwNyJ9

\section{Introducción y objetivos}

Un escalón compresor utiliza el rotor para incrementar la entalpía total y el estator para recuperar presión estática y orientar la corriente. La dificultad aerodinámica es que ambas filas suelen difundir la corriente en su referencia. El capítulo enseña a construir un escalón desde trabajo, flujo y reacción y a comprobar que los triángulos coincidan con el balance termodinámico. Pertenece al bloque III.
El diseño al radio medio es un primer nivel de descripción. No representa por sí solo el comportamiento de punta y raíz. Los datos elegidos deben permitir después una distribución radial, una geometría de anillo y una velocidad de giro compatibles.
% studium:end-record eyJpZCI6ICJQQVItMDA0NCIsICJzZWN0aW9uIjogImMwNyJ9

\section{Estaciones y componentes}

Se usa 1 antes del rotor, 2 después del rotor y 3 después del estator. Se mantiene \(\mathbf C=\mathbf W+\mathbf U\). Al radio medio constante:

\[
\fitmath{w_c=U(C_{\theta2}-C_{\theta1}),\qquad
\psi_c=\frac{C_{\theta2}-C_{\theta1}}U,\qquad \phi=\frac{C_x}U.}
\]
El rotor eleva la energía total aun cuando reduzca la velocidad relativa. El estator no entrega trabajo: \(h_{03}=h_{02}\). Si su salida es más axial y disminuye el módulo absoluto, aumenta la entalpía estática. La figura~\ref{fig:compresor} distingue las estaciones y las referencias; las flechas corresponden a componentes algebraicas.
Con \(C_x\) constante, el giro absoluto es \(\alpha_2-\alpha_1\), mientras el giro relativo es \(\beta_2-\beta_1\). No se suman ángulos de referencias distintas para calcular la deflexión de una fila. Si una referencia utiliza módulos positivos de beta, las fórmulas de giro deben transformarse antes de compararlas.
% studium:end-record eyJpZCI6ICJQQVItMDA0NSIsICJzZWN0aW9uIjogImMwNyJ9

\par\vspace{\baselineskip}
\noindent\begin{minipage}{\linewidth}
\centering
\begingroup
\setbox0=\hbox{\begin{tikzpicture}
\draw[thick] (-.8,0)--(8.8,0);\draw[thick] (-.8,3)--(8.8,3);\fill[orange!25] (.5,.4) rectangle (2.5,2.6);\fill[blue!20] (4.5,.4) rectangle (6.5,2.6);\node at (1.5,1.5) {Rotor};\node at (5.5,1.5) {Estator};\foreach \x/\n in {-.3/1,3.5/2,7.7/3}{\draw[dashed] (\x,-.2)--(\x,3.2);\node[above] at (\x,3.2) {\n};}\draw[->,thick] (7,1.5)--(8.5,1.5) node[right] {$+x$};\draw[->,orange,thick] (1.5,3.5)--(1.5,4.3) node[above] {$U>0$};\node at (3.5,-.8) {$\Delta h_{0,R}>0,\quad\Delta h_{0,S}=0$};
\end{tikzpicture}}%
\count0=\ht0
\advance\count0 by \dp0
\count2=\wd0
\count4=\dimexpr\textheight/100*38\relax
\count6=\linewidth
\ifnum\count0<1
\count0=1
\fi
\ifnum\count2<1
\count2=1
\fi
\divide\count0 by 4096
\divide\count2 by 4096
\divide\count4 by 4096
\divide\count6 by 4096
\ifnum\count0<1
\count0=1
\fi
\ifnum\count2<1
\count2=1
\fi
\ifnum\count4<1
\count4=1
\fi
\ifnum\count6<1
\count6=1
\fi
\count8=\count2
\ifnum\count0>\count4
\count8=\numexpr\count2*\count4/\count0\relax
\fi
\ifnum\count8>\count6
\resizebox{\linewidth}{!}{\begin{tikzpicture}
\draw[thick] (-.8,0)--(8.8,0);\draw[thick] (-.8,3)--(8.8,3);\fill[orange!25] (.5,.4) rectangle (2.5,2.6);\fill[blue!20] (4.5,.4) rectangle (6.5,2.6);\node at (1.5,1.5) {Rotor};\node at (5.5,1.5) {Estator};\foreach \x/\n in {-.3/1,3.5/2,7.7/3}{\draw[dashed] (\x,-.2)--(\x,3.2);\node[above] at (\x,3.2) {\n};}\draw[->,thick] (7,1.5)--(8.5,1.5) node[right] {$+x$};\draw[->,orange,thick] (1.5,3.5)--(1.5,4.3) node[above] {$U>0$};\node at (3.5,-.8) {$\Delta h_{0,R}>0,\quad\Delta h_{0,S}=0$};
\end{tikzpicture}}%
\else
\ifnum\count0>\count4
\resizebox{!}{\dimexpr\textheight/100*38\relax}{\begin{tikzpicture}
\draw[thick] (-.8,0)--(8.8,0);\draw[thick] (-.8,3)--(8.8,3);\fill[orange!25] (.5,.4) rectangle (2.5,2.6);\fill[blue!20] (4.5,.4) rectangle (6.5,2.6);\node at (1.5,1.5) {Rotor};\node at (5.5,1.5) {Estator};\foreach \x/\n in {-.3/1,3.5/2,7.7/3}{\draw[dashed] (\x,-.2)--(\x,3.2);\node[above] at (\x,3.2) {\n};}\draw[->,thick] (7,1.5)--(8.5,1.5) node[right] {$+x$};\draw[->,orange,thick] (1.5,3.5)--(1.5,4.3) node[above] {$U>0$};\node at (3.5,-.8) {$\Delta h_{0,R}>0,\quad\Delta h_{0,S}=0$};
\end{tikzpicture}}%
\else
\begin{tikzpicture}
\draw[thick] (-.8,0)--(8.8,0);\draw[thick] (-.8,3)--(8.8,3);\fill[orange!25] (.5,.4) rectangle (2.5,2.6);\fill[blue!20] (4.5,.4) rectangle (6.5,2.6);\node at (1.5,1.5) {Rotor};\node at (5.5,1.5) {Estator};\foreach \x/\n in {-.3/1,3.5/2,7.7/3}{\draw[dashed] (\x,-.2)--(\x,3.2);\node[above] at (\x,3.2) {\n};}\draw[->,thick] (7,1.5)--(8.5,1.5) node[right] {$+x$};\draw[->,orange,thick] (1.5,3.5)--(1.5,4.3) node[above] {$U>0$};\node at (3.5,-.8) {$\Delta h_{0,R}>0,\quad\Delta h_{0,S}=0$};
\end{tikzpicture}%
\fi
\fi
\endgroup
\par\nopagebreak
\studiumcaption{fig:compresor}{Esquema de rotor y estator de compresor. Estaciones absolutas 1, 2 y 3. Sin escala geométrica.}
\end{minipage}
\par\vspace{\baselineskip}
% studium:end-record eyJpZCI6ICJGSUctMDAwNyIsICJzZWN0aW9uIjogImMwNyJ9

\section{Reacción del escalón repetitivo}


La reacción de compresor se define por



\[
\fitmath{\mathcal R_c=\frac{h_2-h_1}{h_3-h_1}.}
\]
Si \(C_3^2=C_1^2\), el incremento estático del escalón coincide con \(w_c\). Combinando Euler con la definición de entalpía total:

\[
\fitmath{h_2-h_1=U(C_{\theta2}-C_{\theta1})-\frac{C_2^2-C_1^2}{2}.}
\]
Con igual componente axial a ambos lados del rotor,

\[
\fitmath{\mathcal R_c=1-\frac{C_{\theta1}+C_{\theta2}}{2U}.}
\]
Así,

\[
\fitmath{\frac{C_{\theta1}}U=1-\mathcal R_c-\frac{\psi_c}{2},\qquad
\frac{C_{\theta2}}U=1-\mathcal R_c+\frac{\psi_c}{2}.}
\]
Una reacción de medio proporciona, bajo esas hipótesis, triángulos simétricos. No basta con mantener constante el radio medio para obtenerla: también se necesita el reparto de remolino indicado. La reacción debe comprobarse con entalpías si cambian las velocidades entre extremos.
% studium:end-record eyJpZCI6ICJQQVItMDA0NiIsICJzZWN0aW9uIjogImMwNyJ9

\section{De la carga a la presión}

Si se fijan \(U\) y \(\psi_c\), el incremento de temperatura total es \(\Delta T_0=\psi_cU^2/c_p\). Una eficiencia total-total del escalón permite hallar

\[
\fitmath{\pi_{c,esc}=\left(1+\eta_{c,\mathrm{tt}}\frac{\psi_cU^2}{c_pT_{01}}\right)^{\gamma/(\gamma-1)}.}
\]
A igual carga y velocidad periférica, la relación de presiones disminuye cuando aumenta la temperatura de entrada. Por eso una sucesión de escalones de igual trabajo no produce relaciones de presión iguales. Si se pretende repetición geométrica y aerodinámica, deben examinarse la variación de densidad y Mach y la disponibilidad de área.
Para completar las estaciones se calcula \(C_i^2=C_x^2+C_{\theta i}^2\), luego \(T_i=T_{0i}-C_i^2/(2c_p)\). Las presiones estáticas requieren además \(p_{0i}\). Una eficiencia global determina el extremo final, pero no separa automáticamente la pérdida del rotor de la del estator. Es preciso un reparto adicional físicamente consistente si se quiere una tabla completa de presiones intermedias.
Caso completo original por estaciones: compresor axial. Se usa un radio medio fijo y velocidad axial constante. Datos: U=280 m/s, Cx=160 m/s, caudal 20 kg/s, radio medio 0.28 m, gamma=1.4 y cp=1004.5 J/(kg K). Los remolinos absolutos en 1, 2 y 3 son 40, 150, 40 m/s. La presión total relativa del rotor de salida es 0,98 veces la de entrada, y la presión total absoluta del estator de salida es 0,985 veces la de entrada. Estas son pérdidas prescritas del caso educativo; no se añade además una eficiencia global independiente.
Procedimiento: se obtiene el trabajo con Euler y se conserva la entalpía total absoluta en el estator. Se resta la energía cinética absoluta para hallar cada temperatura estática. En el rotor se transforma al remanso relativo, se aplica su razón de pérdida y se transforma de vuelta al estado absoluto. Se aplica la pérdida absoluta del estator en la fila que corresponde. Las presiones estáticas se calculan desde el estado de remanso local, no mediante una isentropía ficticia de todo el escalón.

\[
\fitmath{\begin{array}{c|rrrr}\text{Est.}&T\,[\mathrm K]&T_0\,[\mathrm K]&p\,[\mathrm{kPa}]&p_0\,[\mathrm{kPa}]\\\hline1&286.461&300&93.583&110\\2&306.72&330.662&116.491&151.544\\3&317.123&330.662&128.952&149.271\\\end{array}}
\]

\[
\fitmath{\begin{array}{c|rrrrr}\text{Est.}&W_\theta\,[\mathrm{m/s}]&\alpha\,[{}^\circ]&\beta\,[{}^\circ]&T_{0r}\,[\mathrm K]&p_{0r}\,[\mathrm{kPa}]\\\hline1&-240&14.0362&-56.3099&327.875&150.122\\2&-130&43.1524&-39.0939&327.875&147.12\\3&-240&14.0362&-56.3099&358.537&198.152\\\end{array}}
\]
En todos los cortes la componente axial es la prescrita. Las áreas se obtienen con la densidad estática y esa componente, y valen 0.109815, 0.094459, 0.0882251 metros cuadrados. Las alturas correspondientes son 62.4198, 53.6915, 50.1481 milímetros.
Resultados integrados: trabajo 30.8 kJ/kg, eficiencia total-total resultante 0.891732 y reacción estática 0.660714. La eficiencia se ha deducido después de las pérdidas de fila y no se ha impuesto como un dato adicional. La entropía aumenta y los estados permanecen positivos.
Comprobación: el salto de entalpía total coincide con Euler, el remanso relativo se conserva a través del rotor a radio fijo y cada área reconstruye el mismo caudal. Los residuos se guardan en el expediente; son inferiores a una millonésima en sus unidades coherentes. El cierre de balances no valida las razones de pérdida prescritas ni la aptitud aerodinámica de una geometría real.
% studium:end-record eyJpZCI6ICJQQVItMDA0NyIsICJzZWN0aW9uIjogImMwNyJ9

\section{Comprobación aerodinámica preliminar}


El Mach relativo de entrada al rotor es



\[
\fitmath{M_{r1}=\frac{\sqrt{C_x^2+(C_{\theta1}-U)^2}}{\sqrt{\gamma R_g T_1}}.}
\]
Un Mach absoluto subsónico no garantiza entrada relativa subsónica. La velocidad periférica puede dominar en punta. Asimismo, una reacción de medio no garantiza difusión moderada: debe evaluarse \(W_2/W_1\) y un parámetro de carga de cascada.
Si se eleva \(\psi_c\), aumenta el cambio tangencial exigido. Para igual solidez puede crecer la difusión y la sensibilidad a incidencia. Si se reduce \(\phi\), los ángulos se alejan del eje y la operación puede aproximarse a separación. Ninguna de estas tendencias produce por sí sola un límite universal; la decisión necesita datos de perfil y del conjunto.
% studium:end-record eyJpZCI6ICJQQVItMDA0OCIsICJzZWN0aW9uIjogImMwNyJ9

\section{Procedimiento y errores}

El orden recomendado es: seleccionar un convenio y estaciones; fijar parámetros independientes; resolver componentes; calcular módulos y ángulos; aplicar energía y eficiencia; determinar áreas; comprobar difusión, Mach y reacción. Solo después se estima la geometría metálica. Autocomprobación: mantenga \(\psi_c\) y \(U\), cambie \(\mathcal R_c\) y observe qué permanece constante y qué cambia en los triángulos.
Errores habituales: asignar trabajo al estator, utilizar \(W\) en Euler, suponer que radio constante implica reacción de medio y aplicar eficiencia global para inventar presiones intermedias. Formulario: \(w_c=U\Delta C_\theta\), \(\Delta T_0=w_c/c_p\), \(\pi_c=[1+\eta_cw_c/(c_pT_{01})]^{\gamma/(\gamma-1)}\).
Referencias: ULE Diseño de compresores, pp. \(2-5\) y \(18-24\); MIT 16.50, lección 24, pp. \(2-4\). Desarrollo y casos numéricos originales. La hipótesis de repetición se verifica expresamente en cada caso.
% studium:end-record eyJpZCI6ICJQQVItMDA0OSIsICJzZWN0aW9uIjogImMwNyJ9

\begin{tcolorbox}[title={Desarrollos del modelo}, studiumdef, breakable]

\noindent\textbf{Modelo desarrollado: Trabajo y presión del escalón}

\begin{equation}
\pi_c=\left(1+\frac{\eta_{tt}\,\psi\,U^2}{T_{0e}\,c_p}\right)^{\frac{\gamma}{\gamma-1}}
\label{eq:EC-C07}
\end{equation}

\noindent\textit{Significado.} Relación de presiones totales de un escalón de compresor. $\psi$ es la carga adimensional y $\eta_{tt}$ la eficiencia total-total.

\noindent\textbf{Hipótesis}
\begin{itemize}
\item Escalón adiabático; eficiencia global y gas perfecto constantes.
\end{itemize}

\noindent\textbf{Principios}
\begin{itemize}
\item Balance y definiciones físicas desarrollados en el capítulo; propiedades y dominio declarados en el problema.
\end{itemize}

\noindent\textbf{Aplicabilidad}
\begin{itemize}
\item Modelo educativo del caso C07-01-R
\end{itemize}

\noindent\textbf{Limitaciones}
\begin{itemize}
\item Modelo educativo; la aplicabilidad a una máquina real requiere revisión experta.
\end{itemize}

\noindent\textbf{Pasos}
\begin{enumerate}
\item
La carga define el trabajo comunicado.

\[
\fitmath{w_c=\psi U^{2}}
\]
\item
El balance energético fija el incremento térmico.

\[
\fitmath{\Delta T_0=\frac{\psi U^{2}}{c_p}}
\]
\item
El trabajo ideal equivalente fija la relación de presión.

\[
\fitmath{\pi_c=\left(1 + \frac{\eta \psi U^{2}}{T_{0e} c_p}\right)^{\frac{\gamma}{\gamma - 1}}}
\]
\end{enumerate}

{\noindent\footnotesize verificada dimensionalmente.\par}
\end{tcolorbox}

\begin{tcolorbox}[title={Problema resuelto 1}, studiumworked, breakable]

\noindent\textbf{Enunciado}
\(C07-01-R \cdot\) Trabajo y presión del escalón. Dados los parámetros del escalón compresor, calcule trabajo, incremento total de temperatura y relación de presión. Datos: \(U=273\, \mathrm{m}\,\mathrm{s}^{-1}\); \(\psi=0.39\); \(T_{0e}=301\, \mathrm{K}\); \(\eta=0.88\); \(c_p=1004.5\, \mathrm{J}\,\mathrm{kg}^{-1}\,\mathrm{K}^{-1}\); \(\gamma=1.4\).

\noindent\textbf{Datos}
\begin{itemize}
\item $U = 273\,\mathrm{m}\,\mathrm{s}^{-1}$ — velocidad periférica
\item $\psi = 0.39$ — carga
\item $T_{0e} = 301\,\mathrm{K}$ — temperatura total entrada
\item $\eta = 0.88$ — eficiencia total-total
\item $c_p = 1004.5\,\mathrm{J}\,\mathrm{kg}^{-1}\,\mathrm{K}^{-1}$ — calor específico
\item $\gamma = 1.4$ — razón de calores
\end{itemize}

\noindent\textbf{Incógnita}
$\pi_c$ (adimensional); también las magnitudes intermedias desarrolladas.

\noindent\textbf{Modelo físico}

\noindent\textit{Trabajo y presión del escalón}

\[
\fitmath{\pi_c=\left(1 + \frac{\eta \psi U^{2}}{T_{0e} c_p}\right)^{\frac{\gamma}{\gamma - 1}}}
\]

\noindent\textbf{Hipótesis}
\begin{itemize}
\item Escalón adiabático; eficiencia global y gas perfecto constantes.
\end{itemize}

\noindent\textbf{Desarrollo paso a paso}

\begin{enumerate}

\item La carga define el trabajo comunicado.

\[\fitmath{\mathrm{wc}=\psi\cdot {\mathrm{uu}}^{2}}\]

\[\fitmath{\mathrm{wc}=(0.39)\cdot {(273)}^{2}\approx 29066.31\,\mathrm{J/kg}}\]

\item El balance energético fija el incremento térmico.

\[\fitmath{\mathrm{dT}=\frac{\psi\cdot {\mathrm{uu}}^{2}}{\mathrm{cp}}}\]

\[\fitmath{\mathrm{dT}=\frac{(0.39)\cdot {(273)}^{2}}{(1004.5)}\approx 28.9361\,\mathrm{K}}\]

\item El trabajo ideal equivalente fija la relación de presión.

\[\fitmath{\mathrm{pr}={\left(1+\frac{\eta\cdot \psi\cdot {\mathrm{uu}}^{2}}{\mathrm{cp}\cdot \mathrm{Tin}}\right)}^{\frac{\gamma}{\gamma-1}}}\]

\[\fitmath{\mathrm{pr}={\left(1+\frac{(0.88)\cdot (0.39)\cdot {(273)}^{2}}{(1004.5)\cdot (301)}\right)}^{\frac{(1.4)}{(1.4)-1}}\approx 1.328739\,}\]

\end{enumerate}


\noindent\textbf{Sustitución numérica}

\[
\fitmath{\left(1 + \frac{\left(0.88\right) \left(0.39\right) \left(\left(273\,\mathrm{m}\,\mathrm{s}^{-1}\right)\right)^{2}}{\left(301\,\mathrm{K}\right) \left(1004.5\,\mathrm{J}\,\mathrm{kg}^{-1}\,\mathrm{K}^{-1}\right)}\right)^{\frac{\left(1.4\right)}{\left(1.4\right) - 1}}}
\]

\noindent\textbf{Resultado}

\[
\fitmath{\pi_c=1.328739105\times10^{0}}
\]

\noindent\textbf{Interpretación}
La relación debe ser mayor que uno. Compare el trabajo reconstruido desde la temperatura con Euler.

\noindent\textbf{Comprobación de los pasos}
\begin{itemize}
\item Paso 2: $29066.3$ — reproducido
\item Paso 4: $28.9361$ — reproducido
\item Paso 6: $1.32874$ — reproducido
\end{itemize}

\noindent\textbf{Comprobación de la respuesta} la aritmética del último paso reproduce la respuesta enunciada a la precisión con que está escrita. Es una reproducción aritmética; no valida el modelo físico ni sus hipótesis.

\noindent\textbf{Limitaciones}
\begin{itemize}
\item Modelo educativo; la aplicabilidad a una máquina real requiere revisión experta.
\end{itemize}

\noindent\textbf{Errores frecuentes}
\begin{itemize}
\item Aplicar la razón de temperaturas reales como si fuera isentrópica.
\end{itemize}
% studium:end-record eyJpZCI6ICJQUkItMDA0OSIsICJzZWN0aW9uIjogImMwNyJ9
\end{tcolorbox}

\begin{tcolorbox}[title={Problema resuelto 2}, studiumworked, breakable]

\noindent\textbf{Enunciado}
\(C07-02-R \cdot\) Reacción y remolinos. Construya un rotor compresor repetitivo a partir de reacción y carga. La incógnita final es el ángulo relativo de salida, algebraico en radianes. Datos: \(U=296\, \mathrm{m}\,\mathrm{s}^{-1}\); \(\psi=0.43\); \(\mathcal R=0.55\); \(C_x=164\, \mathrm{m}\,\mathrm{s}^{-1}\).

\noindent\textbf{Datos}
\begin{itemize}
\item $U = 296\,\mathrm{m}\,\mathrm{s}^{-1}$ — velocidad periférica
\item $\psi = 0.43$ — carga
\item $\mathcal R = 0.55$ — reacción
\item $C_x = 164\,\mathrm{m}\,\mathrm{s}^{-1}$ — velocidad axial
\end{itemize}

\noindent\textbf{Incógnita}
$\beta_2$, en $\mathrm{rad}$; también las magnitudes intermedias desarrolladas.

\noindent\textbf{Modelo físico}

\noindent\textit{Reacción y remolinos}

\[
\fitmath{\beta_2=\operatorname{atan}{\left(\frac{U \left(\frac{\psi}{2} - \mathcal R\right)}{C_x} \right)}}
\]

\noindent\textbf{Hipótesis}
\begin{itemize}
\item Radio y velocidad axial constantes; C3 al cuadrado igual a C1 al cuadrado.
\end{itemize}

\noindent\textbf{Desarrollo paso a paso}

\begin{enumerate}

\item La suma y diferencia fijan el remolino de entrada.

\[\fitmath{ct1=\mathrm{uu}\cdot \left(1-\mathrm{rea}-\frac{\mathrm{psi}}{2}\right)}\]

\[\fitmath{ct1=(296)\cdot \left(1-(0.55)-\frac{(0.43)}{2}\right)\approx 69.56\,\mathrm{m/s}}\]

\item Se calcula la componente absoluta de salida.

\[\fitmath{ct2=\mathrm{uu}\cdot \left(1-\mathrm{rea}+\frac{\mathrm{psi}}{2}\right)}\]

\[\fitmath{ct2=(296)\cdot \left(1-(0.55)+\frac{(0.43)}{2}\right)\approx 196.84\,\mathrm{m/s}}\]

\item Se resta U y se obtiene el ángulo desde el eje.

\[\fitmath{beta2=\arctan\left(\frac{\mathrm{uu}\cdot \left(-\mathrm{rea}+\frac{\mathrm{psi}}{2}\right)}{\mathrm{cx}}\right)}\]

\[\fitmath{beta2=\arctan\left(\frac{(296)\cdot \left(-(0.55)+\frac{(0.43)}{2}\right)}{(164)}\right)\approx -0.54382\,\mathrm{rad}}\]

\end{enumerate}


\noindent\textbf{Sustitución numérica}

\[
\fitmath{\operatorname{atan}{\left(\frac{\left(296\,\mathrm{m}\,\mathrm{s}^{-1}\right) \left(\frac{\left(0.43\right)}{2} - \left(0.55\right)\right)}{\left(164\,\mathrm{m}\,\mathrm{s}^{-1}\right)} \right)}}
\]

\noindent\textbf{Resultado}

\[
\fitmath{\beta_2=-5.438199954\times10^{-1}\,\mathrm{rad}}
\]

\noindent\textbf{Interpretación}
La beta negativa se conserva en la construcción. La reacción dada solo coincide con la definición bajo la hipótesis repetitiva.

\noindent\textbf{Comprobación de los pasos}
\begin{itemize}
\item Paso 2: $69.56$ — reproducido
\item Paso 4: $196.84$ — reproducido
\item Paso 6: $-0.54382$ — reproducido
\end{itemize}

\noindent\textbf{Comprobación de la respuesta} la aritmética del último paso reproduce la respuesta enunciada a la precisión con que está escrita. Es una reproducción aritmética; no valida el modelo físico ni sus hipótesis.

\noindent\textbf{Limitaciones}
\begin{itemize}
\item Modelo educativo; la aplicabilidad a una máquina real requiere revisión experta.
\end{itemize}

\noindent\textbf{Errores frecuentes}
\begin{itemize}
\item Usar el módulo positivo de beta de otro convenio.
\end{itemize}
% studium:end-record eyJpZCI6ICJQUkItMDA1MSIsICJzZWN0aW9uIjogImMwNyJ9
\end{tcolorbox}

\begin{tcolorbox}[title={Problema resuelto 3}, studiumworked, breakable]

\noindent\textbf{Enunciado}
\(C07-03-R \cdot\) Mach relativo. Calcule temperatura estática de entrada y Mach relativo del rotor a partir de remanso y componentes absolutas. Datos: \(T_0=309\, \mathrm{K}\); \(C_x=172\, \mathrm{m}\,\mathrm{s}^{-1}\); \(C_\theta=43\, \mathrm{m}\,\mathrm{s}^{-1}\); \(U=319\, \mathrm{m}\,\mathrm{s}^{-1}\); \(c_p=1004.5\, \mathrm{J}\,\mathrm{kg}^{-1}\,\mathrm{K}^{-1}\); \(\gamma=1.4\); \(R_g=287\, \mathrm{J}\,\mathrm{kg}^{-1}\,\mathrm{K}^{-1}\).

\noindent\textbf{Datos}
\begin{itemize}
\item $T_0 = 309\,\mathrm{K}$ — temperatura total absoluta
\item $C_x = 172\,\mathrm{m}\,\mathrm{s}^{-1}$ — componente axial
\item $C_\theta = 43\,\mathrm{m}\,\mathrm{s}^{-1}$ — remolino absoluto
\item $U = 319\,\mathrm{m}\,\mathrm{s}^{-1}$ — velocidad periférica
\item $c_p = 1004.5\,\mathrm{J}\,\mathrm{kg}^{-1}\,\mathrm{K}^{-1}$ — calor específico
\item $\gamma = 1.4$ — razón de calores
\item $R_g = 287\,\mathrm{J}\,\mathrm{kg}^{-1}\,\mathrm{K}^{-1}$ — constante del gas
\end{itemize}

\noindent\textbf{Incógnita}
$\mathrm{Ma}_r$ (adimensional); también las magnitudes intermedias desarrolladas.

\noindent\textbf{Modelo físico}

\noindent\textit{Mach relativo}

\[
\fitmath{\mathrm{Ma}_r=\frac{\sqrt{C_x^{2} + \left(C_\theta - U\right)^{2}}}{\sqrt{R_g \gamma \left(T_0 - \frac{C_\theta^{2} + C_x^{2}}{2 c_p}\right)}}}
\]

\noindent\textbf{Hipótesis}
\begin{itemize}
\item Gas perfecto y referencia local de rotor.
\end{itemize}

\noindent\textbf{Desarrollo paso a paso}

\begin{enumerate}

\item La energía cinética absoluta determina la temperatura estática.

\[\fitmath{T=T_0-\frac{{\mathrm{cx}}^{2}+{\mathrm{ct}}^{2}}{2\cdot \mathrm{cp}}}\]

\[\fitmath{T=(309)-\frac{{(172)}^{2}+{(43)}^{2}}{2\cdot (1004.5)}\approx 293.3539\,\mathrm{K}}\]

\item El módulo relativo se construye vectorialmente.

\[\fitmath{W=\sqrt{{\mathrm{cx}}^{2}+{\left(\mathrm{ct}-\mathrm{uu}\right)}^{2}}}\]

\[\fitmath{W=\sqrt{{(172)}^{2}+{\left((43)-(319)\right)}^{2}}\approx 325.2076\,\mathrm{m/s}}\]

\item El sonido depende de temperatura estática.

\[\fitmath{\mathrm{Mr}=\frac{\sqrt{{\mathrm{cx}}^{2}+{\left(\mathrm{ct}-\mathrm{uu}\right)}^{2}}}{\sqrt{\gamma\cdot R_g\cdot \left(T_0-\frac{{\mathrm{cx}}^{2}+{\mathrm{ct}}^{2}}{2\cdot \mathrm{cp}}\right)}}}\]

\[\fitmath{\mathrm{Mr}=\frac{\sqrt{{(172)}^{2}+{\left((43)-(319)\right)}^{2}}}{\sqrt{(1.4)\cdot (287)\cdot \left((309)-\frac{{(172)}^{2}+{(43)}^{2}}{2\cdot (1004.5)}\right)}}\approx 0.9472395\,}\]

\end{enumerate}


\noindent\textbf{Sustitución numérica}

\[
\fitmath{\frac{\sqrt{\left(\left(172\,\mathrm{m}\,\mathrm{s}^{-1}\right)\right)^{2} + \left(\left(43\,\mathrm{m}\,\mathrm{s}^{-1}\right) - \left(319\,\mathrm{m}\,\mathrm{s}^{-1}\right)\right)^{2}}}{\sqrt{\left(287\,\mathrm{J}\,\mathrm{kg}^{-1}\,\mathrm{K}^{-1}\right) \left(1.4\right) \left(\left(309\,\mathrm{K}\right) - \frac{\left(\left(43\,\mathrm{m}\,\mathrm{s}^{-1}\right)\right)^{2} + \left(\left(172\,\mathrm{m}\,\mathrm{s}^{-1}\right)\right)^{2}}{2 \left(1004.5\,\mathrm{J}\,\mathrm{kg}^{-1}\,\mathrm{K}^{-1}\right)}\right)}}}
\]

\noindent\textbf{Resultado}

\[
\fitmath{\mathrm{Ma}_r=9.472395356\times10^{-1}}
\]

\noindent\textbf{Interpretación}
El Mach relativo puede superar al absoluto. La propiedad acústica utiliza la misma temperatura estática en ambas referencias.

\noindent\textbf{Comprobación de los pasos}
\begin{itemize}
\item Paso 2: $293.354$ — reproducido
\item Paso 4: $325.208$ — reproducido
\item Paso 6: $0.94724$ — reproducido
\end{itemize}

\noindent\textbf{Comprobación de la respuesta} la aritmética del último paso reproduce la respuesta enunciada a la precisión con que está escrita. Es una reproducción aritmética; no valida el modelo físico ni sus hipótesis.

\noindent\textbf{Limitaciones}
\begin{itemize}
\item Modelo educativo; la aplicabilidad a una máquina real requiere revisión experta.
\end{itemize}

\noindent\textbf{Errores frecuentes}
\begin{itemize}
\item Usar temperatura total en la velocidad del sonido.
\end{itemize}
% studium:end-record eyJpZCI6ICJQUkItMDA1MyIsICJzZWN0aW9uIjogImMwNyJ9
\end{tcolorbox}

\begin{tcolorbox}[title={Problema resuelto 4}, studiumworked, breakable]

\noindent\textbf{Enunciado}
\(C07-04-R \cdot\) Cierre de reacción por entalpías. Un rotor tiene los remolinos indicados y velocidad axial constante; el escalón es repetitivo. Calcule su reacción desde incrementos estáticos. Datos: \(U=284\, \mathrm{m}\,\mathrm{s}^{-1}\); \(C_{\theta e}=61\, \mathrm{m}\,\mathrm{s}^{-1}\); \(C_{\theta s}=172\, \mathrm{m}\,\mathrm{s}^{-1}\).

\noindent\textbf{Datos}
\begin{itemize}
\item $U = 284\,\mathrm{m}\,\mathrm{s}^{-1}$ — velocidad periférica
\item $C_{\theta e} = 61\,\mathrm{m}\,\mathrm{s}^{-1}$ — remolino entrada rotor
\item $C_{\theta s} = 172\,\mathrm{m}\,\mathrm{s}^{-1}$ — remolino salida rotor
\end{itemize}

\noindent\textbf{Incógnita}
$\mathcal{R}$ (adimensional); también las magnitudes intermedias desarrolladas.

\noindent\textbf{Modelo físico}

\noindent\textit{Cierre de reacción por entalpías}

\[
\fitmath{\mathcal R=\frac{\frac{C_{\theta e}^{2}}{2} - \frac{C_{\theta s}^{2}}{2} + U \left(- C_{\theta e} + C_{\theta s}\right)}{U \left(- C_{\theta e} + C_{\theta s}\right)}}
\]

\noindent\textbf{Hipótesis}
\begin{itemize}
\item Radio y Cx constantes, igual velocidad absoluta entre extremos del escalón.
\end{itemize}

\noindent\textbf{Desarrollo paso a paso}

\begin{enumerate}

\item Euler da el incremento total de entalpía.

\[\fitmath{\mathrm{dhz}=\mathrm{uu}\cdot \left(\mathrm{cs}-\mathrm{ce}\right)}\]

\[\fitmath{\mathrm{dhz}=(284)\cdot \left((172)-(61)\right)\approx 31524\,\mathrm{J/kg}}\]

\item Se resta el incremento de energía cinética absoluta.

\[\fitmath{\mathrm{dh}=\mathrm{uu}\cdot \left(\mathrm{cs}-\mathrm{ce}\right)-\frac{{\mathrm{cs}}^{2}-{\mathrm{ce}}^{2}}{2}}\]

\[\fitmath{\mathrm{dh}=(284)\cdot \left((172)-(61)\right)-\frac{{(172)}^{2}-{(61)}^{2}}{2}\approx 18592.5\,\mathrm{J/kg}}\]

\item Se forma la razón de incremento estático del rotor y del escalón.

\[\fitmath{\mathrm{rea}=\frac{\mathrm{uu}\cdot \left(\mathrm{cs}-\mathrm{ce}\right)-\frac{{\mathrm{cs}}^{2}-{\mathrm{ce}}^{2}}{2}}{\mathrm{uu}\cdot \left(\mathrm{cs}-\mathrm{ce}\right)}}\]

\[\fitmath{\mathrm{rea}=\frac{(284)\cdot \left((172)-(61)\right)-\frac{{(172)}^{2}-{(61)}^{2}}{2}}{(284)\cdot \left((172)-(61)\right)}\approx 0.5897887\,}\]

\end{enumerate}


\noindent\textbf{Sustitución numérica}

\[
\fitmath{\frac{\frac{\left(\left(61\,\mathrm{m}\,\mathrm{s}^{-1}\right)\right)^{2}}{2} - \frac{\left(\left(172\,\mathrm{m}\,\mathrm{s}^{-1}\right)\right)^{2}}{2} + \left(284\,\mathrm{m}\,\mathrm{s}^{-1}\right) \left(- \left(61\,\mathrm{m}\,\mathrm{s}^{-1}\right) + \left(172\,\mathrm{m}\,\mathrm{s}^{-1}\right)\right)}{\left(284\,\mathrm{m}\,\mathrm{s}^{-1}\right) \left(- \left(61\,\mathrm{m}\,\mathrm{s}^{-1}\right) + \left(172\,\mathrm{m}\,\mathrm{s}^{-1}\right)\right)}}
\]

\noindent\textbf{Resultado}

\[
\fitmath{\mathcal R=5.897887324\times10^{-1}}
\]

\noindent\textbf{Interpretación}
El resultado debe coincidir con 1-(ce+cs)/(2U). Esa identidad es una comprobación algebraica, no experimental.

\noindent\textbf{Comprobación de los pasos}
\begin{itemize}
\item Paso 2: $31524$ — reproducido
\item Paso 4: $18592.5$ — reproducido
\item Paso 6: $0.589789$ — reproducido
\end{itemize}

\noindent\textbf{Comprobación de la respuesta} la aritmética del último paso reproduce la respuesta enunciada a la precisión con que está escrita. Es una reproducción aritmética; no valida el modelo físico ni sus hipótesis.

\noindent\textbf{Limitaciones}
\begin{itemize}
\item Modelo educativo; la aplicabilidad a una máquina real requiere revisión experta.
\end{itemize}

\noindent\textbf{Errores frecuentes}
\begin{itemize}
\item Confundir el incremento estático de rotor y el total.
\end{itemize}
% studium:end-record eyJpZCI6ICJQUkItMDA1NSIsICJzZWN0aW9uIjogImMwNyJ9
\end{tcolorbox}

\begin{tcolorbox}[title={Ejercicios propuestos}, studiumworked, breakable]

\noindent\textbf{1 (Básico, Desarrollo)}
\(C07-01-P \cdot\) Trabajo y presión del escalón. Ejercicio de transferencia. Dados los parámetros del escalón compresor, calcule trabajo, incremento total de temperatura y relación de presión. Antes de consultar la solución, prediga el efecto de cambiar el primer dato respecto del caso resuelto de esta familia. Justifique la tendencia y explique qué restricción podría convertirse en limitante. Datos: \(U=290.745\, \mathrm{m}\,\mathrm{s}^{-1}\); \(\psi=0.39\); \(T_{0e}=301\, \mathrm{K}\); \(\eta=0.88\); \(c_p=1004.5\, \mathrm{J}\,\mathrm{kg}^{-1}\,\mathrm{K}^{-1}\); \(\gamma=1.4\).

\noindent\textbf{2 (Intermedio, Desarrollo)}
\(C07-02-P \cdot\) Reacción y remolinos. Ejercicio de transferencia. Construya un rotor compresor repetitivo a partir de reacción y carga. La incógnita final es el ángulo relativo de salida, algebraico en radianes. Antes de consultar la solución, prediga el efecto de cambiar el primer dato respecto del caso resuelto de esta familia. Justifique la tendencia y explique qué restricción podría convertirse en limitante. Datos: \(U=315.24\, \mathrm{m}\,\mathrm{s}^{-1}\); \(\psi=0.43\); \(\mathcal R=0.55\); \(C_x=164\, \mathrm{m}\,\mathrm{s}^{-1}\).

\noindent\textbf{3 (Avanzado, Desarrollo)}
\(C07-03-P \cdot\) Mach relativo. Ejercicio de transferencia. Calcule temperatura estática de entrada y Mach relativo del rotor a partir de remanso y componentes absolutas. Antes de consultar la solución, prediga el efecto de cambiar el primer dato respecto del caso resuelto de esta familia. Justifique la tendencia y explique qué restricción podría convertirse en limitante. Datos: \(T_0=329.085\, \mathrm{K}\); \(C_x=172\, \mathrm{m}\,\mathrm{s}^{-1}\); \(C_\theta=43\, \mathrm{m}\,\mathrm{s}^{-1}\); \(U=319\, \mathrm{m}\,\mathrm{s}^{-1}\); \(c_p=1004.5\, \mathrm{J}\,\mathrm{kg}^{-1}\,\mathrm{K}^{-1}\); \(\gamma=1.4\); \(R_g=287\, \mathrm{J}\,\mathrm{kg}^{-1}\,\mathrm{K}^{-1}\).

\noindent\textbf{4 (Nivel de examen, Desarrollo)}
\(C07-04-P \cdot\) Cierre de reacción por entalpías. Ejercicio de transferencia. Un rotor tiene los remolinos indicados y velocidad axial constante; el escalón es repetitivo. Calcule su reacción desde incrementos estáticos. Antes de consultar la solución, prediga el efecto de cambiar el primer dato respecto del caso resuelto de esta familia. Justifique la tendencia y explique qué restricción podría convertirse en limitante. Datos: \(U=302.46\, \mathrm{m}\,\mathrm{s}^{-1}\); \(C_{\theta e}=61\, \mathrm{m}\,\mathrm{s}^{-1}\); \(C_{\theta s}=172\, \mathrm{m}\,\mathrm{s}^{-1}\).
\end{tcolorbox}


\studiumchaptercheck

\chapter{Cascadas: difusión, incidencia, desviación y pérdidas}
\label{chap:c08}

\studiumchapterguide{Objetivo: desarrollar cascadas: difusión, incidencia, desviación y pérdidas, declarar hipótesis y comprobar condiciones físicas.}
% studium:end-record eyJpZCI6ICJQQVItMDA1MCIsICJzZWN0aW9uIjogImMwOCJ9

\section{Introducción y objetivos}

Los triángulos fijan las direcciones medias que necesita el escalón, pero el perfil debe conseguir esas direcciones con pérdidas aceptables. El paso del modelo cinemático a una cascada introduce solidez, incidencia, desviación y difusión. El objetivo es separar variables de flujo y de metal, usar indicadores de carga con cautela y reconocer cuándo una correlación requiere datos que el problema no suministra.
Este capítulo pertenece al bloque III y prepara la geometría del bloque IV. Las correlaciones docentes se presentan como herramientas de diseño preliminar; no sustituyen ensayos ni un cálculo viscoso contrastado.
% studium:end-record eyJpZCI6ICJQQVItMDA1MSIsICJzZWN0aW9uIjogImMwOCJ9

\section{Cuerda, paso y solidez}

La cuerda \(c\) es la distancia recta entre borde de ataque y borde de salida. La cuerda axial es \(c_x=c\cos\chi\), donde \(\chi\) es el calado de la cuerda respecto del eje axial. El paso circunferencial de una fila de \(Z_b\) álabes al radio \(r\) es

\[
\fitmath{s=\frac{2\pi r}{Z_b},\qquad \sigma=\frac{c}{s}.}
\]
Una misma fila tiene paso creciente con el radio; la solidez solo permanece constante si la cuerda crece de forma correspondiente. La figura~\ref{fig:cascada} utiliza una línea media analítica original y distingue el paso de la cuerda y de su proyección axial.
Aumentar solidez reparte la carga entre más superficie, pero incrementa área mojada y puede añadir fricción, masa y bloqueo. Reducirla demasiado exige una carga elevada por álabe. El equilibrio óptimo depende del perfil, Mach, Reynolds, incidencia y espesor, de modo que no se prescribe un valor universal.
% studium:end-record eyJpZCI6ICJQQVItMDA1MiIsICJzZWN0aW9uIjogImMwOCJ9

\par\vspace{\baselineskip}
\noindent\begin{minipage}{\linewidth}
\centering
\begingroup
\setbox0=\hbox{\begin{tikzpicture}
\foreach \y in {0,1.7,3.4}{\draw[blue,thick] (0,\y) .. controls (1.2,\y+1.2) and (2.7,\y+.9) .. (4,\y+.5);\draw[blue,thick] (0,\y) .. controls (1.5,\y+.4) and (2.7,\y+.4) .. (4,\y+.5);}\draw[<->] (-.6,0)--(-.6,1.7) node[midway,left] {$s$};\draw[<->] (0,-.6)--(4,-.6) node[midway,below] {$c_x$};\draw[<->] (0,4.3)--(4,4.8) node[midway,above] {$c$};\node at (6,2) {$\sigma=c/s$};
\end{tikzpicture}}%
\count0=\ht0
\advance\count0 by \dp0
\count2=\wd0
\count4=\dimexpr\textheight/100*38\relax
\count6=\linewidth
\ifnum\count0<1
\count0=1
\fi
\ifnum\count2<1
\count2=1
\fi
\divide\count0 by 4096
\divide\count2 by 4096
\divide\count4 by 4096
\divide\count6 by 4096
\ifnum\count0<1
\count0=1
\fi
\ifnum\count2<1
\count2=1
\fi
\ifnum\count4<1
\count4=1
\fi
\ifnum\count6<1
\count6=1
\fi
\count8=\count2
\ifnum\count0>\count4
\count8=\numexpr\count2*\count4/\count0\relax
\fi
\ifnum\count8>\count6
\resizebox{\linewidth}{!}{\begin{tikzpicture}
\foreach \y in {0,1.7,3.4}{\draw[blue,thick] (0,\y) .. controls (1.2,\y+1.2) and (2.7,\y+.9) .. (4,\y+.5);\draw[blue,thick] (0,\y) .. controls (1.5,\y+.4) and (2.7,\y+.4) .. (4,\y+.5);}\draw[<->] (-.6,0)--(-.6,1.7) node[midway,left] {$s$};\draw[<->] (0,-.6)--(4,-.6) node[midway,below] {$c_x$};\draw[<->] (0,4.3)--(4,4.8) node[midway,above] {$c$};\node at (6,2) {$\sigma=c/s$};
\end{tikzpicture}}%
\else
\ifnum\count0>\count4
\resizebox{!}{\dimexpr\textheight/100*38\relax}{\begin{tikzpicture}
\foreach \y in {0,1.7,3.4}{\draw[blue,thick] (0,\y) .. controls (1.2,\y+1.2) and (2.7,\y+.9) .. (4,\y+.5);\draw[blue,thick] (0,\y) .. controls (1.5,\y+.4) and (2.7,\y+.4) .. (4,\y+.5);}\draw[<->] (-.6,0)--(-.6,1.7) node[midway,left] {$s$};\draw[<->] (0,-.6)--(4,-.6) node[midway,below] {$c_x$};\draw[<->] (0,4.3)--(4,4.8) node[midway,above] {$c$};\node at (6,2) {$\sigma=c/s$};
\end{tikzpicture}}%
\else
\begin{tikzpicture}
\foreach \y in {0,1.7,3.4}{\draw[blue,thick] (0,\y) .. controls (1.2,\y+1.2) and (2.7,\y+.9) .. (4,\y+.5);\draw[blue,thick] (0,\y) .. controls (1.5,\y+.4) and (2.7,\y+.4) .. (4,\y+.5);}\draw[<->] (-.6,0)--(-.6,1.7) node[midway,left] {$s$};\draw[<->] (0,-.6)--(4,-.6) node[midway,below] {$c_x$};\draw[<->] (0,4.3)--(4,4.8) node[midway,above] {$c$};\node at (6,2) {$\sigma=c/s$};
\end{tikzpicture}%
\fi
\fi
\endgroup
\par\nopagebreak
\studiumcaption{fig:cascada}{Cascada original: paso, cuerda y cuerda axial distintos. Línea media educativa.}
\end{minipage}
\par\vspace{\baselineskip}
% studium:end-record eyJpZCI6ICJGSUctMDAwOCIsICJzZWN0aW9uIjogImMwOCJ9

\section{Ángulos de gas y de metal}

Se denota por \(a\) el ángulo del gas en la referencia de la fila y por \(\kappa\) el ángulo metálico de la línea media. Se define incidencia algebraica \(i=a_e-\kappa_e\) y desviación algebraica \(\delta=a_s-\kappa_s\). En rotor se usa \(a=\beta\), en estator \(a=\alpha\). Con esta definición el signo de desviación es inequívoco, aunque difiera del módulo positivo usado en otra referencia.
Los ángulos geométricos se deducen de las tangentes de la línea media en sus extremos. El calado de la cuerda no coincide en general con su promedio aritmético. Ese promedio aparece para ciertos perfiles de curvatura simétrica y ciertos convenios, pero no es una identidad geométrica de todas las palas.
La incidencia cambia si cambia \(C_x/U\), el remolino de entrada o la velocidad de giro. La desviación resulta de la respuesta del flujo: no se impone como una orientación exacta sin un modelo. La condición de incidencia cero puede usarse en una geometría educativa inicial; no demuestra que sea la incidencia de mínima pérdida de un perfil real.
% studium:end-record eyJpZCI6ICJQQVItMDA1MyIsICJzZWN0aW9uIjogImMwOCJ9

\section{Factor de difusión}

Un indicador clásico de cascada de compresor se escribe, con magnitudes de velocidad positivas y componentes tangenciales algebraicas,

\[
\fitmath{D=1-\frac{V_s}{V_e}+\frac{|V_{\theta s}-V_{\theta e}|}{2\sigma V_e}.}
\]
En rotor \(V=W\), en estator \(V=C\). El primer término mide la deceleración global; el segundo representa la carga asociada al giro. El valor de \(D\) crece al aumentar la difusión o reducir la solidez para un giro dado. Es una aproximación que resume el efecto sobre la capa límite, no una medida directa de la presión mínima sobre el perfil.
Los límites docentes utilizados en diseño preliminar deben ligarse a la familia de perfiles y al rango de ensayo. No se declara que \(D\) por debajo de un umbral garantice flujo adherido. En los ejercicios, un valor límite suministrado es una condición del modelo del enunciado, no una recomendación de certificación aeronáutica.
Si se despeja la solidez para un \(D_{\mathrm{lim}}\),

\[
\fitmath{\sigma=\frac{|\Delta V_\theta|}{2V_e(D_{\mathrm{lim}}-1+V_s/V_e)}.}
\]
El denominador debe ser positivo. Si no lo es, aumentar solidez no puede satisfacer ese límite con el cociente de velocidades propuesto; hay que redistribuir el trabajo o modificar los triángulos. Esta comprobación evita aceptar una solidez negativa como una solución algebraica válida.
% studium:end-record eyJpZCI6ICJQQVItMDA1NCIsICJzZWN0aW9uIjogImMwOCJ9

\section{Pérdida y uso prudente de correlaciones}

Para una fila fija puede definirse \(Y=(p_{0e}-p_{0s})/(p_{0e}-p_e)\). Otras referencias usan presión dinámica de salida o distinta referencia. Un valor de \(Y\) carece de significado si se omite su denominador. Para un rotor se necesitan presiones totales relativas y estaciones consistentes.
Los apuntes incluyen reglas de desviación tipo Carter. Sus coeficientes y sus ángulos dependen de la convención y de la geometría. En este manual no se transforma una regla particular de diapositiva en una ley universal. El apéndice práctico utiliza inicialmente incidencia y desviación prescritas para separar la construcción geométrica del modelo aerodinámico. Una aplicación de correlación debe registrar edición, ecuación, rango de validez y convención de signos; queda pendiente la calibración de perfiles reales.
La pérdida aumenta fuera del rango de incidencia favorable, pero no tiene por qué ser simétrica respecto de cero. En flujo transónico aparecen pérdidas de choque e interacción con capa límite; no se incorporan aumentando arbitrariamente un coeficiente de una cascada subsónica.
Ejemplo original de la regla docente de desviación. En la diapositiva de compresor se utiliza una forma particular de Carter. Para no mezclarla con nuestros ángulos algebraicos, se definen aquí magnitudes positivas de orientaciones relativas del mismo lado del eje: kappa de entrada y salida y el ángulo de gas de salida a. La regla de este ejemplo es

\[
\fitmath{\delta_m=\frac{\kappa_{e,m}-\kappa_{s,m}}{4\sqrt{\sigma}},\qquad a_{s,m}=\kappa_{s,m}+\delta_m.}
\]
Al sustituir y despejar la orientación metálica de salida:

\[
\fitmath{\kappa_{s,m}=\frac{a_{s,m}-\kappa_{e,m}/(4\sqrt{\sigma})}{1-1/(4\sqrt{\sigma})}.}
\]
Para una entrada metálica de 58 grados, salida del gas de 37 grados y solidez 1,15, se obtiene una salida metálica de 30.61609 grados y desviación de 6.38391 grados. Si las betas son negativas en nuestro convenio, beta de salida es menos 37 grados y kappa de salida es menos 30.61609 grados; por ello la desviación algebraica es menos 6.38391 grados. No se cambia el convenio del capítulo para aplicar una correlación expresada mediante módulos.
Si el dato es cuerda axial sobre paso, la solidez depende además del calado y se necesita iterar el cierre geométrico. La regla anterior es la forma docente concreta utilizada en este ejemplo, no una ley universal. El coeficiente y su validez para un perfil real requieren contrastar familia, Mach, Reynolds y condiciones del ensayo. Para una sección Hermite no se identifica automáticamente el calado con el promedio de ángulos metálicos.
% studium:end-record eyJpZCI6ICJQQVItMDA1NSIsICJzZWN0aW9uIjogImMwOCJ9

\section{Errores, método y síntesis}

Un cálculo de fila debe documentar referencia, ángulos del gas, ángulos metálicos, cuerda, paso, solidez, criterio de difusión y definición de pérdida. Después se verifica que el número entero de álabes elegido conserve aproximadamente la solidez requerida y se recalcula el valor real. El redondeo es una decisión geométrica, no una operación invisible.
Errores habituales: usar \(C\) para la difusión de un rotor, igualar ángulo de salida del gas y del metal sin hipótesis, confundir \(c\) con \(c_x\) y mezclar definiciones de pérdida. Autocomprobación: explique por qué un aumento de solidez puede reducir el indicador de carga y, a la vez, aumentar otra contribución de pérdida.
Formulario: \(s=2\pi r/Z_b\), \(\sigma=c/s\), \(i=a_e-\kappa_e\), \(\delta=a_s-\kappa_s\), factor \(D\) anterior. Referencias: ULE Diseño de compresores, pp. \(6-18\) y 28; Esfuerzos y refrigeración, pp. \(3-6\), para nomenclatura; MIT 16.50, lección 26, pp. \(2-3\), para contraste del uso del factor de difusión.
% studium:end-record eyJpZCI6ICJQQVItMDA1NiIsICJzZWN0aW9uIjogImMwOCJ9

\begin{tcolorbox}[title={Desarrollos del modelo}, studiumdef, breakable]

\noindent\textbf{Modelo desarrollado: Indicador de difusión}

\begin{equation}
D=\frac{\Delta V_\theta}{2\,\sigma\,V_e}+1-\frac{V_s}{V_e}
\label{eq:EC-C08}
\end{equation}

\noindent\textit{Significado.} Factor de difusión de Lieblein en la referencia de la fila. $\Delta V_\theta$ es el cambio tangencial con el signo usado en el caso; la definición clásica usa su valor absoluto.

\noindent\textbf{Hipótesis}
\begin{itemize}
\item Factor de difusión definido en el capítulo, referencia relativa.
\end{itemize}

\noindent\textbf{Principios}
\begin{itemize}
\item Balance y definiciones físicas desarrollados en el capítulo; propiedades y dominio declarados en el problema.
\end{itemize}

\noindent\textbf{Aplicabilidad}
\begin{itemize}
\item Modelo educativo del caso C08-01-R
\end{itemize}

\noindent\textbf{Limitaciones}
\begin{itemize}
\item Modelo educativo; la aplicabilidad a una máquina real requiere revisión experta.
\end{itemize}

\noindent\textbf{Pasos}
\begin{enumerate}
\item
Se calcula la contribución de deceleración.

\[
\fitmath{1 - \frac{W_s}{W_e}=1 - \frac{W_s}{W_e}}
\]
\item
La carga tangencial añade la segunda contribución.

\[
\fitmath{D=\frac{|\Delta W_\theta|}{2 \sigma W_e} + 1 - \frac{W_s}{W_e}}
\]
\end{enumerate}

{\noindent\footnotesize verificada dimensionalmente.\par}
\end{tcolorbox}

\begin{tcolorbox}[title={Problema resuelto 1}, studiumworked, breakable]

\noindent\textbf{Enunciado}
\(C08-01-R \cdot\) Indicador de difusión. Estime D para una fila móvil con módulos relativos y cambio tangencial dados. Datos: \(W_e=282\, \mathrm{m}\,\mathrm{s}^{-1}\); \(W_s=216\, \mathrm{m}\,\mathrm{s}^{-1}\); \(|\Delta W_\theta|=97\, \mathrm{m}\,\mathrm{s}^{-1}\); \(\sigma=1.24\).

\noindent\textbf{Datos}
\begin{itemize}
\item $W_e = 282\,\mathrm{m}\,\mathrm{s}^{-1}$ — módulo relativo entrada
\item $W_s = 216\,\mathrm{m}\,\mathrm{s}^{-1}$ — módulo relativo salida
\item $|\Delta W_\theta| = 97\,\mathrm{m}\,\mathrm{s}^{-1}$ — magnitud de cambio tangencial relativo
\item $\sigma = 1.24$ — solidez
\end{itemize}

\noindent\textbf{Incógnita}
$D$ (adimensional); también las magnitudes intermedias desarrolladas.

\noindent\textbf{Modelo físico}

\noindent\textit{Indicador de difusión}

\[
\fitmath{D=\frac{|\Delta W_\theta|}{2 \sigma W_e} + 1 - \frac{W_s}{W_e}}
\]

\noindent\textbf{Hipótesis}
\begin{itemize}
\item Factor de difusión definido en el capítulo, referencia relativa.
\end{itemize}

\noindent\textbf{Desarrollo paso a paso}

\begin{enumerate}

\item Se calcula la contribución de deceleración.

\[\fitmath{\mathrm{dec}=1-\frac{\mathrm{vs}}{\mathrm{ve}}}\]

\[\fitmath{\mathrm{dec}=1-\frac{(216)}{(282)}\approx 0.2340426\,}\]

\item La carga tangencial añade la segunda contribución.

\[\fitmath{D=1-\frac{\mathrm{vs}}{\mathrm{ve}}+\frac{\mathrm{dv}}{2\cdot \mathrm{sol}\cdot \mathrm{ve}}}\]

\[\fitmath{D=1-\frac{(216)}{(282)}+\frac{(97)}{2\cdot (1.24)\cdot (282)}\approx 0.3727408\,}\]

\end{enumerate}


\noindent\textbf{Sustitución numérica}

\[
\fitmath{\frac{\left(97\,\mathrm{m}\,\mathrm{s}^{-1}\right)}{2 \left(1.24\right) \left(282\,\mathrm{m}\,\mathrm{s}^{-1}\right)} + 1 - \frac{\left(216\,\mathrm{m}\,\mathrm{s}^{-1}\right)}{\left(282\,\mathrm{m}\,\mathrm{s}^{-1}\right)}}
\]

\noindent\textbf{Resultado}

\[
\fitmath{D=3.727407916\times10^{-1}}
\]

\noindent\textbf{Interpretación}
D es un indicador preliminar. La aceptabilidad exige criterio de perfil y régimen declarado.

\noindent\textbf{Comprobación de los pasos}
\begin{itemize}
\item Paso 2: $0.234043$ — reproducido
\item Paso 4: $0.372741$ — reproducido
\end{itemize}

\noindent\textbf{Comprobación de la respuesta} la aritmética del último paso reproduce la respuesta enunciada a la precisión con que está escrita. Es una reproducción aritmética; no valida el modelo físico ni sus hipótesis.

\noindent\textbf{Limitaciones}
\begin{itemize}
\item Modelo educativo; la aplicabilidad a una máquina real requiere revisión experta.
\end{itemize}

\noindent\textbf{Errores frecuentes}
\begin{itemize}
\item Usar velocidades absolutas para un rotor.
\end{itemize}
% studium:end-record eyJpZCI6ICJQUkItMDA1NyIsICJzZWN0aW9uIjogImMwOCJ9
\end{tcolorbox}

\begin{tcolorbox}[title={Problema resuelto 2}, studiumworked, breakable]

\noindent\textbf{Enunciado}
\(C08-02-R \cdot\) Solidez mínima del modelo. Un diseño impone D no superior al valor indicado. Calcule la solidez que satisface exactamente el límite y compruebe la positividad del denominador. Datos: \(W_e=301\, \mathrm{m}\,\mathrm{s}^{-1}\); \(W_s=232\, \mathrm{m}\,\mathrm{s}^{-1}\); \(|\Delta W_\theta|=112\, \mathrm{m}\,\mathrm{s}^{-1}\); \(D_{\mathrm{lim}}=0.47\).

\noindent\textbf{Datos}
\begin{itemize}
\item $W_e = 301\,\mathrm{m}\,\mathrm{s}^{-1}$ — módulo relativo entrada
\item $W_s = 232\,\mathrm{m}\,\mathrm{s}^{-1}$ — módulo relativo salida
\item $|\Delta W_\theta| = 112\,\mathrm{m}\,\mathrm{s}^{-1}$ — cambio tangencial de magnitud positiva
\item $D_{\mathrm{lim}} = 0.47$ — límite educativo de difusión
\end{itemize}

\noindent\textbf{Incógnita}
$\sigma$ (adimensional); también las magnitudes intermedias desarrolladas.

\noindent\textbf{Modelo físico}

\noindent\textit{Solidez mínima del modelo}

\[
\fitmath{\sigma=\frac{|\Delta W_\theta|}{2 W_e \left(D_{\mathrm{lim}} - 1 + \frac{W_s}{W_e}\right)}}
\]

\noindent\textbf{Hipótesis}
\begin{itemize}
\item Criterio de difusión prescrito como dato; sin garantía de pérdida real.
\end{itemize}

\noindent\textbf{Desarrollo paso a paso}

\begin{enumerate}

\item Se determina el margen restante después de la deceleración.

\[\fitmath{\mathrm{den}=\mathrm{lim}-1+\frac{\mathrm{vs}}{\mathrm{ve}}}\]

\[\fitmath{\mathrm{den}=(0.47)-1+\frac{(232)}{(301)}\approx 0.2407641\,}\]

\item Se despeja la solidez.

\[\fitmath{\mathrm{sol}=\frac{\mathrm{dv}}{2\cdot \mathrm{ve}\cdot \left(\mathrm{lim}-1+\frac{\mathrm{vs}}{\mathrm{ve}}\right)}}\]

\[\fitmath{\mathrm{sol}=\frac{(112)}{2\cdot (301)\cdot \left((0.47)-1+\frac{(232)}{(301)}\right)}\approx 0.7727335\,}\]

\end{enumerate}


\noindent\textbf{Sustitución numérica}

\[
\fitmath{\frac{\left(112\,\mathrm{m}\,\mathrm{s}^{-1}\right)}{2 \left(301\,\mathrm{m}\,\mathrm{s}^{-1}\right) \left(\left(0.47\right) - 1 + \frac{\left(232\,\mathrm{m}\,\mathrm{s}^{-1}\right)}{\left(301\,\mathrm{m}\,\mathrm{s}^{-1}\right)}\right)}}
\]

\noindent\textbf{Resultado}

\[
\fitmath{\sigma=7.727335449\times10^{-1}}
\]

\noindent\textbf{Interpretación}
Si el denominador fuera no positivo, este giro y cociente de módulos no podrían cumplir el límite por aumentar solidez.

\noindent\textbf{Comprobación de los pasos}
\begin{itemize}
\item Paso 2: $0.240764$ — reproducido
\item Paso 4: $0.772734$ — reproducido
\end{itemize}

\noindent\textbf{Comprobación de la respuesta} la aritmética del último paso reproduce la respuesta enunciada a la precisión con que está escrita. Es una reproducción aritmética; no valida el modelo físico ni sus hipótesis.

\noindent\textbf{Limitaciones}
\begin{itemize}
\item Modelo educativo; la aplicabilidad a una máquina real requiere revisión experta.
\end{itemize}

\noindent\textbf{Errores frecuentes}
\begin{itemize}
\item Aceptar una solidez negativa.
\end{itemize}
% studium:end-record eyJpZCI6ICJQUkItMDA1OSIsICJzZWN0aW9uIjogImMwOCJ9
\end{tcolorbox}

\begin{tcolorbox}[title={Problema resuelto 3}, studiumworked, breakable]

\noindent\textbf{Enunciado}
\(C08-03-R \cdot\) Número entero y solidez real. Se eligen Z álabes enteros para un radio y una cuerda dados. Calcule paso y solidez real tras la elección. Datos: \(r=0.263\, \mathrm{m}\); \(c=0.052\, \mathrm{m}\); \(Z=41\).

\noindent\textbf{Datos}
\begin{itemize}
\item $r = 0.263\,\mathrm{m}$ — radio
\item $c = 0.052\,\mathrm{m}$ — cuerda
\item $Z = 41$ — número entero de álabes
\end{itemize}

\noindent\textbf{Incógnita}
$\sigma$ (adimensional); también las magnitudes intermedias desarrolladas.

\noindent\textbf{Modelo físico}

\noindent\textit{Número entero y solidez real}

\[
\fitmath{\sigma=\frac{c Z}{2 \pi r}}
\]

\noindent\textbf{Hipótesis}
\begin{itemize}
\item Fila uniforme al radio de referencia.
\end{itemize}

\noindent\textbf{Desarrollo paso a paso}

\begin{enumerate}

\item La circunferencia se reparte en pasos iguales.

\[\fitmath{\mathrm{ss}=\frac{2\cdot \pi\cdot \mathrm{rad}}{\mathrm{zz}}}\]

\[\fitmath{\mathrm{ss}=\frac{2\cdot \pi\cdot (0.263)}{(41)}\approx 0.04030434\,\mathrm{m}}\]

\item La solidez usa el paso realmente obtenido.

\[\fitmath{\mathrm{sol}=\frac{\mathrm{cc}\cdot \mathrm{zz}}{2\cdot \pi\cdot \mathrm{rad}}}\]

\[\fitmath{\mathrm{sol}=\frac{(0.052)\cdot (41)}{2\cdot \pi\cdot (0.263)}\approx 1.290184\,}\]

\end{enumerate}


\noindent\textbf{Sustitución numérica}

\[
\fitmath{\frac{\left(0.052\,\mathrm{m}\right) \left(41\right)}{2 \pi \left(0.263\,\mathrm{m}\right)}}
\]

\noindent\textbf{Resultado}

\[
\fitmath{\sigma=1.290183797\times10^{0}}
\]

\noindent\textbf{Interpretación}
El redondeo del número de álabes modifica la solidez; debe evaluarse después el parámetro de carga.

\noindent\textbf{Comprobación de los pasos}
\begin{itemize}
\item Paso 2: $0.0403043$ — reproducido
\item Paso 4: $1.29018$ — reproducido
\end{itemize}

\noindent\textbf{Comprobación de la respuesta} la aritmética del último paso reproduce la respuesta enunciada a la precisión con que está escrita. Es una reproducción aritmética; no valida el modelo físico ni sus hipótesis.

\noindent\textbf{Limitaciones}
\begin{itemize}
\item Modelo educativo; la aplicabilidad a una máquina real requiere revisión experta.
\end{itemize}

\noindent\textbf{Errores frecuentes}
\begin{itemize}
\item Conservar exactamente la solidez objetivo después del redondeo.
\end{itemize}
% studium:end-record eyJpZCI6ICJQUkItMDA2MSIsICJzZWN0aW9uIjogImMwOCJ9
\end{tcolorbox}

\begin{tcolorbox}[title={Problema resuelto 4}, studiumworked, breakable]

\noindent\textbf{Enunciado}
\(C08-04-R \cdot\) Incidencia y desviación algebraicas. En una fila fija se miden ángulos de gas y se conocen los metálicos. Determine incidencia y desviación; la incógnita final es delta. Datos: \(\alpha_e=0.8552113\, \mathrm{rad}\); \(\alpha_s=0.3839724\, \mathrm{rad}\); \(\kappa_e=0.8028515\, \mathrm{rad}\); \(\kappa_s=0.3141593\, \mathrm{rad}\).

\noindent\textbf{Datos}
\begin{itemize}
\item $\alpha_e = 0.85521133347722\,\mathrm{rad}$ — ángulo gas entrada
\item $\alpha_s = 0.38397243543875\,\mathrm{rad}$ — ángulo gas salida
\item $\kappa_e = 0.80285145591739\,\mathrm{rad}$ — ángulo metal entrada
\item $\kappa_s = 0.31415926535898\,\mathrm{rad}$ — ángulo metal salida
\end{itemize}

\noindent\textbf{Incógnita}
$\delta$, en $\mathrm{rad}$; también las magnitudes intermedias desarrolladas.

\noindent\textbf{Modelo físico}

\noindent\textit{Incidencia y desviación algebraicas}

\[
\fitmath{\delta=\alpha_s - \kappa_s}
\]

\noindent\textbf{Hipótesis}
\begin{itemize}
\item Ángulos desde el eje y mismo sentido positivo.
\end{itemize}

\noindent\textbf{Desarrollo paso a paso}

\begin{enumerate}

\item La incidencia compara gas y metal en entrada.

\[\fitmath{\mathrm{inc}=\mathrm{ae}-\mathrm{ke}}\]

\[\fitmath{\mathrm{inc}=(0.8552113335)-(0.8028514559)\approx 0.05235988\,\mathrm{rad}}\]

\item La desviación compara gas y metal en salida según nuestro convenio.

\[\fitmath{\mathrm{delta}=\mathrm{asv}-\mathrm{ks}}\]

\[\fitmath{\mathrm{delta}=(0.3839724354)-(0.3141592654)\approx 0.06981317\,\mathrm{rad}}\]

\end{enumerate}


\noindent\textbf{Sustitución numérica}

\[
\fitmath{\left(0.38397243543875\,\mathrm{rad}\right) - \left(0.31415926535898\,\mathrm{rad}\right)}
\]

\noindent\textbf{Resultado}

\[
\fitmath{\delta=6.981317008\times10^{-2}\,\mathrm{rad}}
\]

\noindent\textbf{Interpretación}
Los signos dependen de una definición explícita. Para expresar en grados multiplique los radianes por 180/pi.

\noindent\textbf{Comprobación de los pasos}
\begin{itemize}
\item Paso 2: $0.0523599$ — reproducido
\item Paso 4: $0.0698132$ — reproducido
\end{itemize}

\noindent\textbf{Comprobación de la respuesta} la aritmética del último paso reproduce la respuesta enunciada a la precisión con que está escrita. Es una reproducción aritmética; no valida el modelo físico ni sus hipótesis.

\noindent\textbf{Limitaciones}
\begin{itemize}
\item Modelo educativo; la aplicabilidad a una máquina real requiere revisión experta.
\end{itemize}

\noindent\textbf{Errores frecuentes}
\begin{itemize}
\item Comparar ángulo absoluto con beta relativa.
\end{itemize}
% studium:end-record eyJpZCI6ICJQUkItMDA2MyIsICJzZWN0aW9uIjogImMwOCJ9
\end{tcolorbox}

\begin{tcolorbox}[title={Ejercicios propuestos}, studiumworked, breakable]

\noindent\textbf{1 (Básico, Desarrollo)}
\(C08-01-P \cdot\) Indicador de difusión. Ejercicio de transferencia. Estime D para una fila móvil con módulos relativos y cambio tangencial dados. Antes de consultar la solución, prediga el efecto de cambiar el primer dato respecto del caso resuelto de esta familia. Justifique la tendencia y explique qué restricción podría convertirse en limitante. Datos: \(W_e=300.33\, \mathrm{m}\,\mathrm{s}^{-1}\); \(W_s=216\, \mathrm{m}\,\mathrm{s}^{-1}\); \(|\Delta W_\theta|=97\, \mathrm{m}\,\mathrm{s}^{-1}\); \(\sigma=1.24\).

\noindent\textbf{2 (Intermedio, Desarrollo)}
\(C08-02-P \cdot\) Solidez mínima del modelo. Ejercicio de transferencia. Un diseño impone D no superior al valor indicado. Calcule la solidez que satisface exactamente el límite y compruebe la positividad del denominador. Antes de consultar la solución, prediga el efecto de cambiar el primer dato respecto del caso resuelto de esta familia. Justifique la tendencia y explique qué restricción podría convertirse en limitante. Datos: \(W_e=320.565\, \mathrm{m}\,\mathrm{s}^{-1}\); \(W_s=232\, \mathrm{m}\,\mathrm{s}^{-1}\); \(|\Delta W_\theta|=112\, \mathrm{m}\,\mathrm{s}^{-1}\); \(D_{\mathrm{lim}}=0.47\).

\noindent\textbf{3 (Avanzado, Desarrollo)}
\(C08-03-P \cdot\) Número entero y solidez real. Ejercicio de transferencia. Se eligen Z álabes enteros para un radio y una cuerda dados. Calcule paso y solidez real tras la elección. Antes de consultar la solución, prediga el efecto de cambiar el primer dato respecto del caso resuelto de esta familia. Justifique la tendencia y explique qué restricción podría convertirse en limitante. Datos: \(r=0.280095\, \mathrm{m}\); \(c=0.052\, \mathrm{m}\); \(Z=41\).

\noindent\textbf{4 (Nivel de examen, Desarrollo)}
\(C08-04-P \cdot\) Incidencia y desviación algebraicas. Ejercicio de transferencia. En una fila fija se miden ángulos de gas y se conocen los metálicos. Determine incidencia y desviación; la incógnita final es delta. Antes de consultar la solución, prediga el efecto de cambiar el primer dato respecto del caso resuelto de esta familia. Justifique la tendencia y explique qué restricción podría convertirse en limitante. Datos: \(\alpha_e=0.9108001\, \mathrm{rad}\); \(\alpha_s=0.3839724\, \mathrm{rad}\); \(\kappa_e=0.8028515\, \mathrm{rad}\); \(\kappa_s=0.3141593\, \mathrm{rad}\).
\end{tcolorbox}


\studiumchaptercheck

\chapter{Compresión multiescalón, anillo y acoplamiento}
\label{chap:c09}

\studiumchapterguide{Objetivo: desarrollar compresión multiescalón, anillo y acoplamiento, declarar hipótesis y comprobar condiciones físicas.}
% studium:end-record eyJpZCI6ICJQQVItMDA1NyIsICJzZWN0aW9uIjogImMwOSJ9

\section{Introducción y objetivos}

Una sucesión de escalones no es una repetición independiente del primero. La temperatura y la densidad cambian, las áreas se adaptan y todos los rotores del mismo eje comparten velocidad angular. El objetivo es calcular la evolución media, dimensionar el anillo y comprobar el balance de potencia de un grupo. Pertenece al bloque III.
El cálculo preliminar debe identificar qué se conserva entre escalones: velocidad de giro, radio medio, carga, incremento de temperatura o relación de presiones. Esas condiciones no pueden imponerse todas de forma arbitraria.
% studium:end-record eyJpZCI6ICJQQVItMDA1OCIsICJzZWN0aW9uIjogImMwOSJ9

\section{Acumulación de temperatura y presión}


Para \(n\) escalones adiabáticos,



\[
\fitmath{w_{c,tot}=\sum_{j=1}^n w_{cj}=c_p(T_{0,n+1}-T_{01}),\qquad
\pi_{c,tot}=\prod_{j=1}^n\pi_{cj}.}
\]
Si cada escalón entrega el mismo trabajo, la temperatura total crece linealmente en el índice del escalón. Si cada uno tiene igual eficiencia y trabajo, su relación de presión cambia porque el cociente \(w_{cj}/(c_pT_{0j})\) disminuye. Una compresión de relación constante por escalón, en cambio, exige trabajos crecientes si la temperatura aumenta.
Con eficiencia politrópica constante,

\[
\fitmath{T_{0,n+1}=T_{01}\pi_{c,tot}^{(\gamma-1)/(\gamma\eta_{p,c})}.}
\]
Esta expresión permite estimar la temperatura final y el trabajo global, pero no determina la distribución de carga óptima ni las dimensiones de cada escalón. Es necesario un reparto y una comprobación por estaciones.
% studium:end-record eyJpZCI6ICJQQVItMDA1OSIsICJzZWN0aW9uIjogImMwOSJ9

\section{Área anular y altura}


Para flujo axial medio uniforme,



\[
\fitmath{A_j=\frac{\dot m}{\rho_j C_{xj}},\qquad \rho_j=\frac{p_j}{R_gT_j}.}
\]
Si se usa \(r_m=(r_t+r_h)/2\) y \(b=r_t-r_h\), la identidad exacta es

\[
\fitmath{A=\pi(r_t^2-r_h^2)=2\pi r_m b.}
\]
Por tanto \(b=A/(2\pi r_m)\), \(r_t=r_m+b/2\) y \(r_h=r_m-b/2\). Debe verificarse \(r_h>0\). No se confunde este radio medio aritmético con un radio de igual área o con el radio medio cuadrático usado en otros métodos de promediado.
La figura~\ref{fig:anillo} muestra alternativas a radio medio, punta o raíz constantes. Al comprimir, la densidad creciente tiende a reducir el área si se conserva \(C_x\). Una altura final pequeña aumenta la importancia relativa de la holgura de punta y del bloqueo por paredes.
% studium:end-record eyJpZCI6ICJQQVItMDA2MCIsICJzZWN0aW9uIjogImMwOSJ9

\par\vspace{\baselineskip}
\noindent\begin{minipage}{\linewidth}
\centering
\begingroup
\setbox0=\hbox{\begin{tikzpicture}
\draw[dashed,->] (-.4,0)--(8,0) node[right] {$x$};\draw[thick] (0,1.1)--(7,1.6);\draw[thick] (0,3.2)--(7,2.7);\draw[dashed] (0,2.15)--(7,2.15);\draw[<->] (1,1.17)--(1,3.13) node[midway,right] {$b$};\draw[<->] (5,0)--(5,2.15) node[midway,left] {$r_m$};\node at (4,3.8) {$A=\pi(r_t^2-r_h^2)=2\pi r_mb$};
\end{tikzpicture}}%
\count0=\ht0
\advance\count0 by \dp0
\count2=\wd0
\count4=\dimexpr\textheight/100*38\relax
\count6=\linewidth
\ifnum\count0<1
\count0=1
\fi
\ifnum\count2<1
\count2=1
\fi
\divide\count0 by 4096
\divide\count2 by 4096
\divide\count4 by 4096
\divide\count6 by 4096
\ifnum\count0<1
\count0=1
\fi
\ifnum\count2<1
\count2=1
\fi
\ifnum\count4<1
\count4=1
\fi
\ifnum\count6<1
\count6=1
\fi
\count8=\count2
\ifnum\count0>\count4
\count8=\numexpr\count2*\count4/\count0\relax
\fi
\ifnum\count8>\count6
\resizebox{\linewidth}{!}{\begin{tikzpicture}
\draw[dashed,->] (-.4,0)--(8,0) node[right] {$x$};\draw[thick] (0,1.1)--(7,1.6);\draw[thick] (0,3.2)--(7,2.7);\draw[dashed] (0,2.15)--(7,2.15);\draw[<->] (1,1.17)--(1,3.13) node[midway,right] {$b$};\draw[<->] (5,0)--(5,2.15) node[midway,left] {$r_m$};\node at (4,3.8) {$A=\pi(r_t^2-r_h^2)=2\pi r_mb$};
\end{tikzpicture}}%
\else
\ifnum\count0>\count4
\resizebox{!}{\dimexpr\textheight/100*38\relax}{\begin{tikzpicture}
\draw[dashed,->] (-.4,0)--(8,0) node[right] {$x$};\draw[thick] (0,1.1)--(7,1.6);\draw[thick] (0,3.2)--(7,2.7);\draw[dashed] (0,2.15)--(7,2.15);\draw[<->] (1,1.17)--(1,3.13) node[midway,right] {$b$};\draw[<->] (5,0)--(5,2.15) node[midway,left] {$r_m$};\node at (4,3.8) {$A=\pi(r_t^2-r_h^2)=2\pi r_mb$};
\end{tikzpicture}}%
\else
\begin{tikzpicture}
\draw[dashed,->] (-.4,0)--(8,0) node[right] {$x$};\draw[thick] (0,1.1)--(7,1.6);\draw[thick] (0,3.2)--(7,2.7);\draw[dashed] (0,2.15)--(7,2.15);\draw[<->] (1,1.17)--(1,3.13) node[midway,right] {$b$};\draw[<->] (5,0)--(5,2.15) node[midway,left] {$r_m$};\node at (4,3.8) {$A=\pi(r_t^2-r_h^2)=2\pi r_mb$};
\end{tikzpicture}%
\fi
\fi
\endgroup
\par\nopagebreak
\studiumcaption{fig:anillo}{Vista meridiana de un anillo que reduce área. Radio medio aritmético y altura de pala.}
\end{minipage}
\par\vspace{\baselineskip}
% studium:end-record eyJpZCI6ICJGSUctMDAwOSIsICJzZWN0aW9uIjogImMwOSJ9

\section{Bloqueo geométrico y aerodinámico}

El área efectiva puede modelarse como \(A_{ef}=B A_{geom}\), donde \(0<B\leq1\) representa una fracción de paso disponible. Con esa definición, \(\dot m=\rho C_xB A_{geom}\). Algunos textos llaman bloqueo a \(1-B\); debe registrarse cuál de las dos magnitudes se proporciona.
El espesor de álabes ocupa parte de la sección; las capas límite y las estelas alteran la distribución de flujo y reducen la capacidad efectiva. Un coeficiente único es una aproximación. No se usa además otra corrección que contabilice dos veces el mismo efecto sin comprobarlo.
La condición de \(C_x\) constante exige una geometría adaptada a la densidad y al bloqueo. No puede mantenerse simultáneamente área constante, densidad creciente y caudal constante sin cambiar \(C_x\). La continuidad es el criterio que decide si una hipótesis de diseño repetitivo es compatible con el anillo.
% studium:end-record eyJpZCI6ICJQQVItMDA2MSIsICJzZWN0aW9uIjogImMwOSJ9

\section{Eje común y balance de potencia}


Para un grupo simple sin accesorios ni sangrado:



\[
\fitmath{\eta_m\dot m_t w_t=\dot m_c w_c.}
\]
Si hay combustión, refrigeración y sangrados, el caudal de la turbina no coincide necesariamente con el del compresor. El balance requiere los caudales reales de cada componente y el trabajo de accesorios. La velocidad angular común no significa velocidades periféricas iguales: \(U=\Omega r\) depende del radio.
La potencia de diseño fija una primera condición, pero el equilibrio operativo requiere también compatibilidad de caudal y presión entre componentes. Un punto que cierra potencia puede no cerrar la capacidad de paso. Esa diferencia motiva los mapas y el cálculo fuera de diseño.
% studium:end-record eyJpZCI6ICJQQVItMDA2MiIsICJzZWN0aW9uIjogImMwOSJ9

\section{Sensibilidad y selección del número de escalones}

Una estimación por \(n=w_{c,tot}/(\psi_cU^2)\) ayuda a explorar tamaños. Tras elegir el entero se distribuye el trabajo y se recalcula cada escalón. No debe suponerse que todos conservan igual Mach o igual reacción solo porque comparten carga adimensional.
La incertidumbre de eficiencia afecta temperatura, densidad, área y potencia a la vez. Un estudio de sensibilidad debe recalcular la cadena completa. Variar únicamente la presión de salida y dejar fijos los demás estados puede violar continuidad o energía. Los problemas avanzados del capítulo obligan a seguir esa dependencia.
% studium:end-record eyJpZCI6ICJQQVItMDA2MyIsICJzZWN0aW9uIjogImMwOSJ9

\section{Errores y formulario}

Errores habituales: multiplicar trabajos o sumar relaciones de presión, usar densidad total para dimensionar área estática, identificar todos los radios medios y omitir el caudal refrigerante en el balance de turbina. Autocomprobación: con caudal y \(C_x\) fijos, duplique la densidad y deduzca cómo cambia la altura a radio medio constante.
Formulario: \(\pi_{tot}=\prod\pi_j\), \(w_{tot}=\sum w_j\), \(A=\dot m/(\rho C_xB)\), \(A=2\pi r_mb\). Referencias: ULE Diseño de compresores, pp. \(20-30\) y 36; MIT 16.50, lección 26, pp. \(4-5\). Las geometrías y balances numéricos del manual son originales y están limitados al modelo medio.
% studium:end-record eyJpZCI6ICJQQVItMDA2NCIsICJzZWN0aW9uIjogImMwOSJ9

\begin{tcolorbox}[title={Desarrollos del modelo}, studiumdef, breakable]

\noindent\textbf{Modelo desarrollado: Compresión de varios escalones}

\begin{equation}
\dot W=\frac{T_{0e}\,c_p\,\dot m}{\eta_m}\left(\pi_c^{\frac{\gamma-1}{\eta_p\,\gamma}}-1\right)
\label{eq:EC-C09}
\end{equation}

\noindent\textit{Significado.} Potencia de eje de una compresión politrópica sin sangrado. $\eta_p$ es la eficiencia politrópica, no el rendimiento propulsivo de una hélice.

\noindent\textbf{Hipótesis}
\begin{itemize}
\item Propiedades y eficiencia politrópica constantes, sin sangrado.
\end{itemize}

\noindent\textbf{Principios}
\begin{itemize}
\item Balance y definiciones físicas desarrollados en el capítulo; propiedades y dominio declarados en el problema.
\end{itemize}

\noindent\textbf{Aplicabilidad}
\begin{itemize}
\item Modelo educativo del caso C09-01-R
\end{itemize}

\noindent\textbf{Limitaciones}
\begin{itemize}
\item Modelo educativo; la aplicabilidad a una máquina real requiere revisión experta.
\end{itemize}

\noindent\textbf{Pasos}
\begin{enumerate}
\item
Se integra el modelo politrópico global.

\[
\fitmath{T_{0s}=T_{0e} \pi_c^{\frac{\gamma - 1}{\eta_p \gamma}}}
\]
\item
La potencia de eje supera la potencia comunicada al fluido.

\[
\fitmath{\dot{W}=\frac{T_{0e} c_p \dot m \left(\pi_c^{\frac{\gamma - 1}{\eta_p \gamma}} - 1\right)}{\eta_m}}
\]
\end{enumerate}

{\noindent\footnotesize verificada dimensionalmente.\par}
\end{tcolorbox}

\begin{tcolorbox}[title={Problema resuelto 1}, studiumworked, breakable]

\noindent\textbf{Enunciado}
\(C09-01-R \cdot\) Compresión de varios escalones. Una compresión politrópica global alimenta un eje. Halle temperatura final y potencia requerida en el eje. Datos: \(T_{0e}=289\, \mathrm{K}\); \(\pi_c=9.3\); \(\eta_p=0.9\); \(\gamma=1.4\); \(c_p=1004.5\, \mathrm{J}\,\mathrm{kg}^{-1}\,\mathrm{K}^{-1}\); \(\dot m=23\, \mathrm{kg}\,\mathrm{s}^{-1}\); \(\eta_m=0.98\).

\noindent\textbf{Datos}
\begin{itemize}
\item $T_{0e} = 289\,\mathrm{K}$ — temperatura total entrada
\item $\pi_c = 9.3$ — relación global de presión
\item $\eta_p = 0.9$ — eficiencia politrópica
\item $\gamma = 1.4$ — razón de calores
\item $c_p = 1004.5\,\mathrm{J}\,\mathrm{kg}^{-1}\,\mathrm{K}^{-1}$ — calor específico
\item $\dot m = 23\,\mathrm{kg}\,\mathrm{s}^{-1}$ — caudal
\item $\eta_m = 0.98$ — rendimiento mecánico compresor
\end{itemize}

\noindent\textbf{Incógnita}
$\dot{W}$, en $\mathrm{W}$; también las magnitudes intermedias desarrolladas.

\noindent\textbf{Modelo físico}

\noindent\textit{Compresión de varios escalones}

\[
\fitmath{\dot{W}=\frac{T_{0e} c_p \dot m \left(\pi_c^{\frac{\gamma - 1}{\eta_p \gamma}} - 1\right)}{\eta_m}}
\]

\noindent\textbf{Hipótesis}
\begin{itemize}
\item Propiedades y eficiencia politrópica constantes, sin sangrado.
\end{itemize}

\noindent\textbf{Desarrollo paso a paso}

\begin{enumerate}

\item Se integra el modelo politrópico global.

\[\fitmath{\mathrm{Tout}=\mathrm{Tin}\cdot {\mathrm{pr}}^{\frac{\gamma-1}{\gamma\cdot \mathrm{ep}}}}\]

\[\fitmath{\mathrm{Tout}=(289)\cdot {(9.3)}^{\frac{(1.4)-1}{(1.4)\cdot (0.9)}}\approx 586.6144\,\mathrm{K}}\]

\item La potencia de eje supera la potencia comunicada al fluido.

\[\fitmath{\mathrm{Pow}=\frac{\mathrm{md}\cdot \mathrm{cp}\cdot \mathrm{Tin}\cdot \left({\mathrm{pr}}^{\frac{\gamma-1}{\gamma\cdot \mathrm{ep}}}-1\right)}{\mathrm{em}}}\]

\[\fitmath{\mathrm{Pow}=\frac{(23)\cdot (1004.5)\cdot (289)\cdot \left({(9.3)}^{\frac{(1.4)-1}{(1.4)\cdot (0.9)}}-1\right)}{(0.98)}\approx 7016260\,\mathrm{W}}\]

\end{enumerate}


\noindent\textbf{Sustitución numérica}

\[
\fitmath{\frac{\left(289\,\mathrm{K}\right) \left(1004.5\,\mathrm{J}\,\mathrm{kg}^{-1}\,\mathrm{K}^{-1}\right) \left(23\,\mathrm{kg}\,\mathrm{s}^{-1}\right) \left(\left(9.3\right)^{\frac{\left(1.4\right) - 1}{\left(0.9\right) \left(1.4\right)}} - 1\right)}{\left(0.98\right)}}
\]

\noindent\textbf{Resultado}

\[
\fitmath{\dot{W}=7.016260352\times10^{6}\,\mathrm{W}}
\]

\noindent\textbf{Interpretación}
No se multiplica la eficiencia de cada escalón. El caudal y el incremento total de entalpía fijan la potencia.

\noindent\textbf{Comprobación de los pasos}
\begin{itemize}
\item Paso 2: $586.614$ — reproducido
\item Paso 4: $7.01626\times 10^{6}$ — reproducido
\end{itemize}

\noindent\textbf{Comprobación de la respuesta} la aritmética del último paso reproduce la respuesta enunciada a la precisión con que está escrita. Es una reproducción aritmética; no valida el modelo físico ni sus hipótesis.

\noindent\textbf{Limitaciones}
\begin{itemize}
\item Modelo educativo; la aplicabilidad a una máquina real requiere revisión experta.
\end{itemize}

\noindent\textbf{Errores frecuentes}
\begin{itemize}
\item Usar el producto de eficiencias como eficiencia global.
\end{itemize}
% studium:end-record eyJpZCI6ICJQUkItMDA2NSIsICJzZWN0aW9uIjogImMwOSJ9
\end{tcolorbox}

\begin{tcolorbox}[title={Problema resuelto 2}, studiumworked, breakable]

\noindent\textbf{Enunciado}
\(C09-02-R \cdot\) Altura de anillo con bloqueo. Dimensione el área geométrica y la altura del anillo con un factor B de área disponible. Datos: \(\dot m=27\, \mathrm{kg}\,\mathrm{s}^{-1}\); \(\rho=3.6\, \mathrm{kg}\,\mathrm{m}^{-3}\); \(C_x=151\, \mathrm{m}\,\mathrm{s}^{-1}\); \(B=0.955\); \(r_m=0.285\, \mathrm{m}\).

\noindent\textbf{Datos}
\begin{itemize}
\item $\dot m = 27\,\mathrm{kg}\,\mathrm{s}^{-1}$ — caudal
\item $\rho = 3.6\,\mathrm{kg}\,\mathrm{m}^{-3}$ — densidad estática
\item $C_x = 151\,\mathrm{m}\,\mathrm{s}^{-1}$ — velocidad axial efectiva
\item $B = 0.955$ — fracción de área disponible
\item $r_m = 0.285\,\mathrm{m}$ — radio medio aritmético
\end{itemize}

\noindent\textbf{Incógnita}
$b$, en $\mathrm{m}$; también las magnitudes intermedias desarrolladas.

\noindent\textbf{Modelo físico}

\noindent\textit{Altura de anillo con bloqueo}

\[
\fitmath{b=\frac{\dot m}{2 \pi B C_x \rho r_m}}
\]

\noindent\textbf{Hipótesis}
\begin{itemize}
\item Flujo uniforme en el área disponible.
\end{itemize}

\noindent\textbf{Desarrollo paso a paso}

\begin{enumerate}

\item La continuidad se aplica al área efectiva B A.

\[\fitmath{\mathrm{area}=\frac{\mathrm{md}}{\rho\cdot \mathrm{cx}\cdot \mathrm{bb}}}\]

\[\fitmath{\mathrm{area}=\frac{(27)}{(3.6)\cdot (151)\cdot (0.955)}\approx 0.05200929\,\mathrm{m^{2}}}\]

\item Se convierte el área geométrica en altura.

\[\fitmath{b=\frac{\mathrm{md}}{\rho\cdot \mathrm{cx}\cdot \mathrm{bb}\cdot 2\cdot \pi\cdot \mathrm{rm}}}\]

\[\fitmath{b=\frac{(27)}{(3.6)\cdot (151)\cdot (0.955)\cdot 2\cdot \pi\cdot (0.285)}\approx 0.02904399\,\mathrm{m}}\]

\end{enumerate}


\noindent\textbf{Sustitución numérica}

\[
\fitmath{\frac{\left(27\,\mathrm{kg}\,\mathrm{s}^{-1}\right)}{2 \pi \left(0.955\right) \left(151\,\mathrm{m}\,\mathrm{s}^{-1}\right) \left(3.6\,\mathrm{kg}\,\mathrm{m}^{-3}\right) \left(0.285\,\mathrm{m}\right)}}
\]

\noindent\textbf{Resultado}

\[
\fitmath{b=2.904398583\times10^{-2}\,\mathrm{m}}
\]

\noindent\textbf{Interpretación}
B inferior a uno exige mayor área geométrica. La raíz debe permanecer positiva al restar b/2 al radio medio.

\noindent\textbf{Comprobación de los pasos}
\begin{itemize}
\item Paso 2: $0.0520093$ — reproducido
\item Paso 4: $0.029044$ — reproducido
\end{itemize}

\noindent\textbf{Comprobación de la respuesta} la aritmética del último paso reproduce la respuesta enunciada a la precisión con que está escrita. Es una reproducción aritmética; no valida el modelo físico ni sus hipótesis.

\noindent\textbf{Limitaciones}
\begin{itemize}
\item Modelo educativo; la aplicabilidad a una máquina real requiere revisión experta.
\end{itemize}

\noindent\textbf{Errores frecuentes}
\begin{itemize}
\item Confundir B con la fracción bloqueada 1-B.
\end{itemize}
% studium:end-record eyJpZCI6ICJQUkItMDA2NyIsICJzZWN0aW9uIjogImMwOSJ9
\end{tcolorbox}

\begin{tcolorbox}[title={Problema resuelto 3}, studiumworked, breakable]

\noindent\textbf{Enunciado}
\(C09-03-R \cdot\) Trabajo por escalón y razón de presión. Un escalón tardío recibe gas más caliente. Calcule su relación de presión para un trabajo y eficiencia fijos. Datos: \(T_{0e}=487\, \mathrm{K}\); \(w_c=29800\, \mathrm{J}\,\mathrm{kg}^{-1}\); \(\eta=0.87\); \(\gamma=1.4\); \(c_p=1004.5\, \mathrm{J}\,\mathrm{kg}^{-1}\,\mathrm{K}^{-1}\).

\noindent\textbf{Datos}
\begin{itemize}
\item $T_{0e} = 487\,\mathrm{K}$ — temperatura total entrada
\item $w_c = 29800\,\mathrm{J}\,\mathrm{kg}^{-1}$ — trabajo específico
\item $\eta = 0.87$ — eficiencia total-total
\item $\gamma = 1.4$ — razón de calores
\item $c_p = 1004.5\,\mathrm{J}\,\mathrm{kg}^{-1}\,\mathrm{K}^{-1}$ — calor específico
\end{itemize}

\noindent\textbf{Incógnita}
$\pi_c$ (adimensional); también las magnitudes intermedias desarrolladas.

\noindent\textbf{Modelo físico}

\noindent\textit{Trabajo por escalón y razón de presión}

\[
\fitmath{\pi_c=\left(1 + \frac{\eta w_c}{T_{0e} c_p}\right)^{\frac{\gamma}{\gamma - 1}}}
\]

\noindent\textbf{Hipótesis}
\begin{itemize}
\item Gas perfecto constante y escalón adiabático.
\end{itemize}

\noindent\textbf{Desarrollo paso a paso}

\begin{enumerate}

\item Se calcula el aumento relativo de temperatura ideal.

\[\fitmath{\mathrm{rr}=\frac{\mathrm{eta}\cdot \mathrm{wc}}{\mathrm{cp}\cdot \mathrm{Tin}}}\]

\[\fitmath{\mathrm{rr}=\frac{(0.87)\cdot (29800)}{(1004.5)\cdot (487)}\approx 0.05299765\,}\]

\item Se recupera la relación de presión.

\[\fitmath{\mathrm{pr}={\left(1+\frac{\mathrm{eta}\cdot \mathrm{wc}}{\mathrm{cp}\cdot \mathrm{Tin}}\right)}^{\frac{\gamma}{\gamma-1}}}\]

\[\fitmath{\mathrm{pr}={\left(1+\frac{(0.87)\cdot (29800)}{(1004.5)\cdot (487)}\right)}^{\frac{(1.4)}{(1.4)-1}}\approx 1.198108\,}\]

\end{enumerate}


\noindent\textbf{Sustitución numérica}

\[
\fitmath{\left(1 + \frac{\left(0.87\right) \left(29800\,\mathrm{J}\,\mathrm{kg}^{-1}\right)}{\left(487\,\mathrm{K}\right) \left(1004.5\,\mathrm{J}\,\mathrm{kg}^{-1}\,\mathrm{K}^{-1}\right)}\right)^{\frac{\left(1.4\right)}{\left(1.4\right) - 1}}}
\]

\noindent\textbf{Resultado}

\[
\fitmath{\pi_c=1.198107832\times10^{0}}
\]

\noindent\textbf{Interpretación}
A mayor temperatura de entrada y mismo trabajo, la razón de presión es menor.

\noindent\textbf{Comprobación de los pasos}
\begin{itemize}
\item Paso 2: $0.0529977$ — reproducido
\item Paso 4: $1.19811$ — reproducido
\end{itemize}

\noindent\textbf{Comprobación de la respuesta} la aritmética del último paso reproduce la respuesta enunciada a la precisión con que está escrita. Es una reproducción aritmética; no valida el modelo físico ni sus hipótesis.

\noindent\textbf{Limitaciones}
\begin{itemize}
\item Modelo educativo; la aplicabilidad a una máquina real requiere revisión experta.
\end{itemize}

\noindent\textbf{Errores frecuentes}
\begin{itemize}
\item Asignar igual razón de presión a todos los escalones de igual trabajo.
\end{itemize}
% studium:end-record eyJpZCI6ICJQUkItMDA2OSIsICJzZWN0aW9uIjogImMwOSJ9
\end{tcolorbox}

\begin{tcolorbox}[title={Problema resuelto 4}, studiumworked, breakable]

\noindent\textbf{Enunciado}
\(C09-04-R \cdot\) Acoplamiento de eje con caudales distintos. El compresor y la turbina comparten eje con rendimiento mecánico dado. Calcule el trabajo que debe extraer cada kilogramo de gas en la turbina. Datos: \(\dot{m}_c=25\, \mathrm{kg}\,\mathrm{s}^{-1}\); \(\dot{m}_t=25.6\, \mathrm{kg}\,\mathrm{s}^{-1}\); \(w_c=214000\, \mathrm{J}\,\mathrm{kg}^{-1}\); \(\dot{W}_{\mathrm{acc}}=85000\, \mathrm{W}\); \(\eta_m=0.975\).

\noindent\textbf{Datos}
\begin{itemize}
\item $\dot{m}_c = 25\,\mathrm{kg}\,\mathrm{s}^{-1}$ — caudal compresor
\item $\dot{m}_t = 25.6\,\mathrm{kg}\,\mathrm{s}^{-1}$ — caudal turbina
\item $w_c = 214000\,\mathrm{J}\,\mathrm{kg}^{-1}$ — trabajo comunicado al aire
\item $\dot{W}_{\mathrm{acc}} = 85000\,\mathrm{W}$ — potencia de accesorios
\item $\eta_m = 0.975$ — eficiencia de transmisión desde turbina
\end{itemize}

\noindent\textbf{Incógnita}
$w_t$, en $\mathrm{J\,kg^{-1}}$; también las magnitudes intermedias desarrolladas.

\noindent\textbf{Modelo físico}

\noindent\textit{Acoplamiento de eje con caudales distintos}

\[
\fitmath{w_t=\frac{\dot{W}_{\mathrm{acc}} + \dot{m}_c w_c}{\eta_m \dot{m}_t}}
\]

\noindent\textbf{Hipótesis}
\begin{itemize}
\item wc ya excluye pérdidas de transmisión; em se aplica una vez al lado turbina.
\end{itemize}

\noindent\textbf{Desarrollo paso a paso}

\begin{enumerate}

\item Se suma la demanda de potencia de aire y accesorios.

\[\fitmath{\mathrm{need}=\mathrm{mc}\cdot \mathrm{wc}+\mathrm{acc}}\]

\[\fitmath{\mathrm{need}=(25)\cdot (214000)+(85000)\approx 5435000\,\mathrm{W}}\]

\item El caudal de turbina y la transmisión fijan su trabajo.

\[\fitmath{\mathrm{wt}=\frac{\mathrm{mc}\cdot \mathrm{wc}+\mathrm{acc}}{\mathrm{em}\cdot \mathrm{mt}}}\]

\[\fitmath{\mathrm{wt}=\frac{(25)\cdot (214000)+(85000)}{(0.975)\cdot (25.6)}\approx 217748.4\,\mathrm{J/kg}}\]

\end{enumerate}


\noindent\textbf{Sustitución numérica}

\[
\fitmath{\frac{\left(85000\,\mathrm{W}\right) + \left(25\,\mathrm{kg}\,\mathrm{s}^{-1}\right) \left(214000\,\mathrm{J}\,\mathrm{kg}^{-1}\right)}{\left(0.975\right) \left(25.6\,\mathrm{kg}\,\mathrm{s}^{-1}\right)}}
\]

\noindent\textbf{Resultado}

\[
\fitmath{w_t=2.177483974\times10^{5}\,\mathrm{J}\,\mathrm{kg}^{-1}}
\]

\noindent\textbf{Interpretación}
Usar el caudal del compresor en la turbina omite la diferencia de masas. El reparto de accesorios está explicitado en este modelo.

\noindent\textbf{Comprobación de los pasos}
\begin{itemize}
\item Paso 2: $5435000$ — reproducido
\item Paso 4: $217748$ — reproducido
\end{itemize}

\noindent\textbf{Comprobación de la respuesta} la aritmética del último paso reproduce la respuesta enunciada a la precisión con que está escrita. Es una reproducción aritmética; no valida el modelo físico ni sus hipótesis.

\noindent\textbf{Limitaciones}
\begin{itemize}
\item Modelo educativo; la aplicabilidad a una máquina real requiere revisión experta.
\end{itemize}

\noindent\textbf{Errores frecuentes}
\begin{itemize}
\item Suponer igualdad de caudales sin revisar combustible y sangrados.
\end{itemize}
% studium:end-record eyJpZCI6ICJQUkItMDA3MSIsICJzZWN0aW9uIjogImMwOSJ9
\end{tcolorbox}

\begin{tcolorbox}[title={Ejercicios propuestos}, studiumworked, breakable]

\noindent\textbf{1 (Básico, Desarrollo)}
\(C09-01-P \cdot\) Compresión de varios escalones. Ejercicio de transferencia. Una compresión politrópica global alimenta un eje. Halle temperatura final y potencia requerida en el eje. Antes de consultar la solución, prediga el efecto de cambiar el primer dato respecto del caso resuelto de esta familia. Justifique la tendencia y explique qué restricción podría convertirse en limitante. Datos: \(T_{0e}=307.785\, \mathrm{K}\); \(\pi_c=9.3\); \(\eta_p=0.9\); \(\gamma=1.4\); \(c_p=1004.5\, \mathrm{J}\,\mathrm{kg}^{-1}\,\mathrm{K}^{-1}\); \(\dot m=23\, \mathrm{kg}\,\mathrm{s}^{-1}\); \(\eta_m=0.98\).

\noindent\textbf{2 (Intermedio, Desarrollo)}
\(C09-02-P \cdot\) Altura de anillo con bloqueo. Ejercicio de transferencia. Dimensione el área geométrica y la altura del anillo con un factor B de área disponible. Antes de consultar la solución, prediga el efecto de cambiar el primer dato respecto del caso resuelto de esta familia. Justifique la tendencia y explique qué restricción podría convertirse en limitante. Datos: \(\dot m=28.755\, \mathrm{kg}\,\mathrm{s}^{-1}\); \(\rho=3.6\, \mathrm{kg}\,\mathrm{m}^{-3}\); \(C_x=151\, \mathrm{m}\,\mathrm{s}^{-1}\); \(B=0.955\); \(r_m=0.285\, \mathrm{m}\).

\noindent\textbf{3 (Avanzado, Desarrollo)}
\(C09-03-P \cdot\) Trabajo por escalón y razón de presión. Ejercicio de transferencia. Un escalón tardío recibe gas más caliente. Calcule su relación de presión para un trabajo y eficiencia fijos. Antes de consultar la solución, prediga el efecto de cambiar el primer dato respecto del caso resuelto de esta familia. Justifique la tendencia y explique qué restricción podría convertirse en limitante. Datos: \(T_{0e}=518.655\, \mathrm{K}\); \(w_c=29800\, \mathrm{J}\,\mathrm{kg}^{-1}\); \(\eta=0.87\); \(\gamma=1.4\); \(c_p=1004.5\, \mathrm{J}\,\mathrm{kg}^{-1}\,\mathrm{K}^{-1}\).

\noindent\textbf{4 (Nivel de examen, Desarrollo)}
\(C09-04-P \cdot\) Acoplamiento de eje con caudales distintos. Ejercicio de transferencia. El compresor y la turbina comparten eje con rendimiento mecánico dado. Calcule el trabajo que debe extraer cada kilogramo de gas en la turbina. Antes de consultar la solución, prediga el efecto de cambiar el primer dato respecto del caso resuelto de esta familia. Justifique la tendencia y explique qué restricción podría convertirse en limitante. Datos: \(\dot{m}_c=26.625\, \mathrm{kg}\,\mathrm{s}^{-1}\); \(\dot{m}_t=25.6\, \mathrm{kg}\,\mathrm{s}^{-1}\); \(w_c=214000\, \mathrm{J}\,\mathrm{kg}^{-1}\); \(\dot{W}_{\mathrm{acc}}=85000\, \mathrm{W}\); \(\eta_m=0.975\).
\end{tcolorbox}


\studiumchaptercheck

\chapter{Arranque, pérdida rotativa, bombeo y mapas}
\label{chap:c10}

\studiumchapterguide{Objetivo: desarrollar arranque, pérdida rotativa, bombeo y mapas, declarar hipótesis y comprobar condiciones físicas.}
% studium:end-record eyJpZCI6ICJQQVItMDA2NSIsICJzZWN0aW9uIjogImMxMCJ9

\section{Introducción y objetivos}

El punto nominal no describe todo el funcionamiento de un compresor. Durante arranque, aceleración o cambios de carga, la relación entre velocidad axial y velocidad periférica cambia y las filas pueden trabajar a incidencias alejadas del diseño. El objetivo es interpretar esa variación y distinguir pérdida de perfil, pérdida rotativa y bombeo. Es contenido explícito del bloque III.
Se introduce además un modelo dinámico mínimo del sistema de compresión como ampliación identificada. Sirve para entender por qué el volumen aguas abajo y la respuesta de la instalación importan; no predice una línea de bombeo real sin calibración.
% studium:end-record eyJpZCI6ICJQQVItMDA2NiIsICJzZWN0aW9uIjogImMxMCJ9

\section{Incidencia fuera de diseño}


En el rotor,



\[
\fitmath{\beta_1=\arctan\frac{C_{\theta1}-\Omega r}{C_x},\qquad i=\beta_1-\kappa_1.}
\]
Si el remolino de entrada es nulo, \(\beta_1=\arctan(-U/C_x)\). Reducir \(C_x/U\) lleva el ángulo relativo hacia valores más negativos con nuestro convenio. El efecto sobre incidencia debe evaluarse con el ángulo metálico y su signo; no se declara incidencia positiva o negativa copiando una figura que mida beta en sentido contrario.
La distribución de densidad a baja velocidad de giro difiere de la nominal. Un anillo diseñado para gran aumento de densidad puede producir una relación de velocidades axiales poco favorable en las filas delanteras y traseras durante el arranque. La geometría fija no se adapta por sí sola a esa nueva distribución.
% studium:end-record eyJpZCI6ICJQQVItMDA2NyIsICJzZWN0aW9uIjogImMxMCJ9

\section{Recursos para ampliar el margen}

Los estatores variables modifican el remolino que recibe el rotor siguiente y ajustan incidencia. Los sangrados intermedios alteran el caudal que deben admitir los escalones posteriores. Los ejes independientes permiten velocidades de giro diferentes en grupos distintos. Estas medidas tienen costes de complejidad, rendimiento y control y se seleccionan para el conjunto de la misión.
No se proporciona aquí una secuencia de arranque de un motor real: depende de fabricante, límites térmicos y procedimientos operativos. El modelo docente analiza compatibilidad aerodinámica. Una predicción de temperatura o margen en un ejercicio no es una instrucción de operación.
% studium:end-record eyJpZCI6ICJQQVItMDA2OCIsICJzZWN0aW9uIjogImMxMCJ9

\section{Fenómenos distintos}

La separación sobre un perfil es un fenómeno local. La pérdida rotativa organiza regiones de flujo degradado que se desplazan alrededor del anillo; puede existir caudal medio globalmente positivo. El bombeo es una oscilación del sistema de compresión, con cambios importantes de caudal y presión y posible inversión de flujo. No son sinónimos ni se deducen de un único triángulo medio.
Un rotor de turbina puede tener separación o bloqueo aunque no presente el mismo mecanismo de inestabilidad de un sistema compresor-plenum. La frase de apuntes que contrapone turbina y surge debe interpretarse en ese contexto, sin concluir que la turbina esté libre de toda inestabilidad.
% studium:end-record eyJpZCI6ICJQQVItMDA2OSIsICJzZWN0aW9uIjogImMxMCJ9

\section{Lectura de mapas y definición de margen}

Un mapa representa una variable de presión frente a caudal, con líneas de velocidad corregida y contornos de eficiencia. Las magnitudes de entrada y las referencias de corrección deben estar declaradas. La figura~\ref{fig:mapa} es un esquema cualitativo original; sus curvas no corresponden a ensayos ni permiten leer un margen certificado.
En este manual, cuando se pide un margen de caudal a igual velocidad corregida, se define explícitamente

\[
\fitmath{SM_m=100\frac{\dot m_{c,\mathrm{op}}-\dot m_{c,\mathrm{lim}}}{\dot m_{c,\mathrm{op}}}.}
\]
Existen otras definiciones que combinan presión y caudal. Un número de margen no es comparable entre fuentes si cambia la definición. La línea de bombeo se determina para un sistema y un criterio de detección, no mediante una parábola elegida por conveniencia gráfica.
El punto operativo es la intersección compatible entre componente y sistema. Al cerrar una válvula, la resistencia del sistema cambia; el punto se desplaza. La evolución temporal depende de inercias y capacidades de almacenamiento, por lo que una curva estacionaria por sí sola no describe el ciclo de bombeo.
% studium:end-record eyJpZCI6ICJQQVItMDA3MCIsICJzZWN0aW9uIjogImMxMCJ9

\par\vspace{\baselineskip}
\noindent\begin{minipage}{\linewidth}
\centering
\begingroup
\setbox0=\hbox{\begin{tikzpicture}
\draw[->] (0,0)--(7.4,0) node[right] {$\dot m_c$};\draw[->] (0,0)--(0,4) node[above] {$\pi_c$};\foreach \x/\y in {1/1.1,1.5/2,2.1/2.9}{\draw[blue,thick] (\x,\y) .. controls (\x+1,\y+.5) and (\x+3,\y+.3) .. (\x+3.8,\y-1);}\draw[red,dashed] (.8,.6) .. controls (1.2,1.8) and (1.7,2.7) .. (2.3,3.5);\node[red] at (2,4.5) {Límite ilustrativo};\draw[->,teal,thick] (2,.6)--(4.8,2.8) node[right] {Operación};\node at (6,3.8) {$N_c$ creciente};
\end{tikzpicture}}%
\count0=\ht0
\advance\count0 by \dp0
\count2=\wd0
\count4=\dimexpr\textheight/100*38\relax
\count6=\linewidth
\ifnum\count0<1
\count0=1
\fi
\ifnum\count2<1
\count2=1
\fi
\divide\count0 by 4096
\divide\count2 by 4096
\divide\count4 by 4096
\divide\count6 by 4096
\ifnum\count0<1
\count0=1
\fi
\ifnum\count2<1
\count2=1
\fi
\ifnum\count4<1
\count4=1
\fi
\ifnum\count6<1
\count6=1
\fi
\count8=\count2
\ifnum\count0>\count4
\count8=\numexpr\count2*\count4/\count0\relax
\fi
\ifnum\count8>\count6
\resizebox{\linewidth}{!}{\begin{tikzpicture}
\draw[->] (0,0)--(7.4,0) node[right] {$\dot m_c$};\draw[->] (0,0)--(0,4) node[above] {$\pi_c$};\foreach \x/\y in {1/1.1,1.5/2,2.1/2.9}{\draw[blue,thick] (\x,\y) .. controls (\x+1,\y+.5) and (\x+3,\y+.3) .. (\x+3.8,\y-1);}\draw[red,dashed] (.8,.6) .. controls (1.2,1.8) and (1.7,2.7) .. (2.3,3.5);\node[red] at (2,4.5) {Límite ilustrativo};\draw[->,teal,thick] (2,.6)--(4.8,2.8) node[right] {Operación};\node at (6,3.8) {$N_c$ creciente};
\end{tikzpicture}}%
\else
\ifnum\count0>\count4
\resizebox{!}{\dimexpr\textheight/100*38\relax}{\begin{tikzpicture}
\draw[->] (0,0)--(7.4,0) node[right] {$\dot m_c$};\draw[->] (0,0)--(0,4) node[above] {$\pi_c$};\foreach \x/\y in {1/1.1,1.5/2,2.1/2.9}{\draw[blue,thick] (\x,\y) .. controls (\x+1,\y+.5) and (\x+3,\y+.3) .. (\x+3.8,\y-1);}\draw[red,dashed] (.8,.6) .. controls (1.2,1.8) and (1.7,2.7) .. (2.3,3.5);\node[red] at (2,4.5) {Límite ilustrativo};\draw[->,teal,thick] (2,.6)--(4.8,2.8) node[right] {Operación};\node at (6,3.8) {$N_c$ creciente};
\end{tikzpicture}}%
\else
\begin{tikzpicture}
\draw[->] (0,0)--(7.4,0) node[right] {$\dot m_c$};\draw[->] (0,0)--(0,4) node[above] {$\pi_c$};\foreach \x/\y in {1/1.1,1.5/2,2.1/2.9}{\draw[blue,thick] (\x,\y) .. controls (\x+1,\y+.5) and (\x+3,\y+.3) .. (\x+3.8,\y-1);}\draw[red,dashed] (.8,.6) .. controls (1.2,1.8) and (1.7,2.7) .. (2.3,3.5);\node[red] at (2,4.5) {Límite ilustrativo};\draw[->,teal,thick] (2,.6)--(4.8,2.8) node[right] {Operación};\node at (6,3.8) {$N_c$ creciente};
\end{tikzpicture}%
\fi
\fi
\endgroup
\par\nopagebreak
\studiumcaption{fig:mapa}{Mapa cualitativo sin escala numérica ni datos experimentales. La línea límite es ilustrativa y no predice bombeo.}
\end{minipage}
\par\vspace{\baselineskip}
% studium:end-record eyJpZCI6ICJGSUctMDAxMCIsICJzZWN0aW9uIjogImMxMCJ9

\section{Ampliación: un modelo dinámico mínimo}

Como ejercicio de razonamiento, considérese una inercia de conducto \(L\), una capacidad de plenum \(K\), una característica de compresor \(\Delta p_c(\dot m)\) y una descarga \(\dot m_v(p)\):

\[
\fitmath{L\frac{d\dot m}{\mathrm{d}t}=\Delta p_c(\dot m)-p,\qquad
K\frac{\mathrm{d}p}{\mathrm{d}t}=\dot m-\dot m_v(p).}
\]
Las unidades de \(L\) y \(K\) se deducen de estas ecuaciones; no son longitudes y volúmenes sin conversión. Linealizando alrededor de un equilibrio y llamando \(a=d\Delta p_c/d\dot m\), \(b=d\dot m_v/\mathrm{d}p\), la matriz es

\[
\fitmath{\begin{pmatrix}a/L&-1/L\\1/K&-b/K\end{pmatrix}.}
\]
La estabilidad lineal de este sistema de dos estados requiere traza negativa y determinante positivo. No basta con mirar el signo de la pendiente del compresor sin considerar los demás parámetros. La pérdida rotativa y las no linealidades están excluidas; por ello no se presenta este modelo como simulador de motor.
% studium:end-record eyJpZCI6ICJQQVItMDA3MSIsICJzZWN0aW9uIjogImMxMCJ9

\section{Errores y formulario}

Errores habituales: utilizar el mapa esquemático como base de diseño, mezclar caudal real y corregido, tratar bombeo y pérdida rotativa como el mismo fenómeno y asignar un signo de incidencia sin reconstruir componentes. Autocomprobación: dos instalaciones usan el mismo compresor, pero distinto plenum; explique por qué pueden responder de manera distinta a una perturbación.
Formulario: \(N_c=N/\sqrt{T_{01}/T_{\mathrm{ref}}}\), \(\dot m_c=\dot m\sqrt{T_{01}/T_{\mathrm{ref}}}/(p_{01}/p_{\mathrm{ref}})\) y el margen de caudal definido arriba. Referencias: ULE Generalidades, pp. \(20-22\); MIT 16.50, lección 26, p. 5, para contexto del arranque. El modelo dinámico es una ampliación original explícita, sometida a comprobación algebraica y dimensional, pendiente de validación para aplicaciones reales.
% studium:end-record eyJpZCI6ICJQQVItMDA3MiIsICJzZWN0aW9uIjogImMxMCJ9

\begin{tcolorbox}[title={Desarrollos del modelo}, studiumdef, breakable]

\noindent\textbf{Modelo desarrollado: Incidencia durante variación de caudal}

\begin{equation}
i=-\kappa-\arctan\frac{U}{C_x}
\label{eq:EC-C10}
\end{equation}

\noindent\textit{Significado.} Incidencia con prerrotación nula. $\kappa$ es el ángulo metálico de entrada; los ángulos se miden desde el eje hacia el giro positivo.

\noindent\textbf{Hipótesis}
\begin{itemize}
\item Velocidad axial positiva, sin prerrotación, metal fijo.
\end{itemize}

\noindent\textbf{Principios}
\begin{itemize}
\item Balance y definiciones físicas desarrollados en el capítulo; propiedades y dominio declarados en el problema.
\end{itemize}

\noindent\textbf{Aplicabilidad}
\begin{itemize}
\item Modelo educativo del caso C10-01-R
\end{itemize}

\noindent\textbf{Limitaciones}
\begin{itemize}
\item Modelo educativo; la aplicabilidad a una máquina real requiere revisión experta.
\end{itemize}

\noindent\textbf{Pasos}
\begin{enumerate}
\item
El triángulo relativo utiliza tangencial -U.

\[
\fitmath{\beta=- \operatorname{atan}{\left(\frac{U}{C_x} \right)}}
\]
\item
Se compara con la dirección del metal.

\[
\fitmath{i=- \kappa_e - \operatorname{atan}{\left(\frac{U}{C_x} \right)}}
\]
\end{enumerate}

{\noindent\footnotesize verificada dimensionalmente.\par}
\end{tcolorbox}

\begin{tcolorbox}[title={Problema resuelto 1}, studiumworked, breakable]

\noindent\textbf{Enunciado}
\(C10-01-R \cdot\) Incidencia durante variación de caudal. Un rotor sin prerrotación conserva velocidad periférica y ángulo metálico. Calcule su beta de entrada y la incidencia a la nueva velocidad axial. Datos: \(U=247\, \mathrm{m}\,\mathrm{s}^{-1}\); \(C_x=119\, \mathrm{m}\,\mathrm{s}^{-1}\); \(\kappa_e=-1.064651\, \mathrm{rad}\).

\noindent\textbf{Datos}
\begin{itemize}
\item $U = 247\,\mathrm{m}\,\mathrm{s}^{-1}$ — velocidad periférica
\item $C_x = 119\,\mathrm{m}\,\mathrm{s}^{-1}$ — velocidad axial nueva
\item $\kappa_e = -1.0646508437165\,\mathrm{rad}$ — ángulo metálico de entrada
\end{itemize}

\noindent\textbf{Incógnita}
$i$, en $\mathrm{rad}$; también las magnitudes intermedias desarrolladas.

\noindent\textbf{Modelo físico}

\noindent\textit{Incidencia durante variación de caudal}

\[
\fitmath{i=- \kappa_e - \operatorname{atan}{\left(\frac{U}{C_x} \right)}}
\]

\noindent\textbf{Hipótesis}
\begin{itemize}
\item Velocidad axial positiva, sin prerrotación, metal fijo.
\end{itemize}

\noindent\textbf{Desarrollo paso a paso}

\begin{enumerate}

\item El triángulo relativo utiliza tangencial -U.

\[\fitmath{\beta=\arctan\left(\frac{-\mathrm{uu}}{\mathrm{cx}}\right)}\]

\[\fitmath{\beta=\arctan\left(\frac{-(247)}{(119)}\right)\approx -1.12183\,\mathrm{rad}}\]

\item Se compara con la dirección del metal.

\[\fitmath{\mathrm{inc}=\arctan\left(\frac{-\mathrm{uu}}{\mathrm{cx}}\right)-\mathrm{kap}}\]

\[\fitmath{\mathrm{inc}=\arctan\left(\frac{-(247)}{(119)}\right)-(-1.064650844)\approx -0.05717871\,\mathrm{rad}}\]

\end{enumerate}


\noindent\textbf{Sustitución numérica}

\[
\fitmath{- \left(-1.0646508437165\,\mathrm{rad}\right) - \operatorname{atan}{\left(\frac{\left(247\,\mathrm{m}\,\mathrm{s}^{-1}\right)}{\left(119\,\mathrm{m}\,\mathrm{s}^{-1}\right)} \right)}}
\]

\noindent\textbf{Resultado}

\[
\fitmath{i=-5.717871161\times10^{-2}\,\mathrm{rad}}
\]

\noindent\textbf{Interpretación}
El signo corresponde exclusivamente al convenio algebraico de este manual. Debe compararse con una incidencia nominal calculada igual.

\noindent\textbf{Comprobación de los pasos}
\begin{itemize}
\item Paso 2: $-1.12183$ — reproducido
\item Paso 4: $-0.0571787$ — reproducido
\end{itemize}

\noindent\textbf{Comprobación de la respuesta} la aritmética del último paso reproduce la respuesta enunciada a la precisión con que está escrita. Es una reproducción aritmética; no valida el modelo físico ni sus hipótesis.

\noindent\textbf{Limitaciones}
\begin{itemize}
\item Modelo educativo; la aplicabilidad a una máquina real requiere revisión experta.
\end{itemize}

\noindent\textbf{Errores frecuentes}
\begin{itemize}
\item Conservar beta al cambiar Cx/U.
\end{itemize}
% studium:end-record eyJpZCI6ICJQUkItMDA3MyIsICJzZWN0aW9uIjogImMxMCJ9
\end{tcolorbox}

\begin{tcolorbox}[title={Problema resuelto 2}, studiumworked, breakable]

\noindent\textbf{Enunciado}
\(C10-02-R \cdot\) Margen de caudal definido. A igual velocidad corregida, un punto operativo y un límite prescrito tienen los caudales indicados. Calcule el margen según la definición del capítulo. Datos: \(\dot{m}_{c,\mathrm{op}}=28.4\, \mathrm{kg}\,\mathrm{s}^{-1}\); \(\dot{m}_{c,\mathrm{lim}}=23.9\, \mathrm{kg}\,\mathrm{s}^{-1}\).

\noindent\textbf{Datos}
\begin{itemize}
\item $\dot{m}_{c,\mathrm{op}} = 28.4\,\mathrm{kg}\,\mathrm{s}^{-1}$ — caudal corregido operativo
\item $\dot{m}_{c,\mathrm{lim}} = 23.9\,\mathrm{kg}\,\mathrm{s}^{-1}$ — caudal corregido límite
\end{itemize}

\noindent\textbf{Incógnita}
$\mathrm{SM}_m$ (adimensional); también las magnitudes intermedias desarrolladas.

\noindent\textbf{Modelo físico}

\noindent\textit{Margen de caudal definido}

\[
\fitmath{\mathrm{SM}_m=\frac{- 100 \dot{m}_{c,\mathrm{lim}} + 100 \dot{m}_{c,\mathrm{op}}}{\dot{m}_{c,\mathrm{op}}}}
\]

\noindent\textbf{Hipótesis}
\begin{itemize}
\item Caudales corregidos a las mismas referencias.
\end{itemize}

\noindent\textbf{Desarrollo paso a paso}

\begin{enumerate}

\item Se calcula la separación de caudal.

\[\fitmath{\mathrm{dm}=\mathrm{op}-\mathrm{lim}}\]

\[\fitmath{\mathrm{dm}=(28.4)-(23.9)\approx 4.5\,\mathrm{kg/s}}\]

\item Se normaliza por el caudal operativo y se expresa en porcentaje.

\[\fitmath{\mathrm{SM}=\frac{100\cdot \left(\mathrm{op}-\mathrm{lim}\right)}{\mathrm{op}}}\]

\[\fitmath{\mathrm{SM}=\frac{100\cdot \left((28.4)-(23.9)\right)}{(28.4)}\approx 15.84507\,}\]

\end{enumerate}


\noindent\textbf{Sustitución numérica}

\[
\fitmath{\frac{- 100 \left(23.9\,\mathrm{kg}\,\mathrm{s}^{-1}\right) + 100 \left(28.4\,\mathrm{kg}\,\mathrm{s}^{-1}\right)}{\left(28.4\,\mathrm{kg}\,\mathrm{s}^{-1}\right)}}
\]

\noindent\textbf{Resultado}

\[
\fitmath{\mathrm{SM}_m=1.584507042\times10^{1}}
\]

\noindent\textbf{Interpretación}
El porcentaje solo es comparable con la misma definición y velocidad corregida. El límite es un dato, no una predicción del modelo.

\noindent\textbf{Comprobación de los pasos}
\begin{itemize}
\item Paso 2: $4.5$ — reproducido
\item Paso 4: $15.8451$ — reproducido
\end{itemize}

\noindent\textbf{Comprobación de la respuesta} la aritmética del último paso reproduce la respuesta enunciada a la precisión con que está escrita. Es una reproducción aritmética; no valida el modelo físico ni sus hipótesis.

\noindent\textbf{Limitaciones}
\begin{itemize}
\item Modelo educativo; la aplicabilidad a una máquina real requiere revisión experta.
\end{itemize}

\noindent\textbf{Errores frecuentes}
\begin{itemize}
\item Normalizar con el límite y llamar igual al margen resultante.
\end{itemize}
% studium:end-record eyJpZCI6ICJQUkItMDA3NSIsICJzZWN0aW9uIjogImMxMCJ9
\end{tcolorbox}

\begin{tcolorbox}[title={Problema resuelto 3}, studiumworked, breakable]

\noindent\textbf{Enunciado}
\(C10-03-R \cdot\) Corrección de velocidad. Halle la velocidad corregida y la velocidad periférica física al radio dado. La incógnita final es U real. Datos: \(N=196.6667\, \mathrm{s}^{-1}\); \(T_{0e}=259\, \mathrm{K}\); \(T_{\mathrm{ref}}=288.15\, \mathrm{K}\); \(r=0.276\, \mathrm{m}\).

\noindent\textbf{Datos}
\begin{itemize}
\item $N = 196.66666666667\,\mathrm{s}^{-1}$ — giro físico en revoluciones por segundo; equivale a 11800 rpm
\item $T_{0e} = 259\,\mathrm{K}$ — temperatura total entrada
\item $T_{\mathrm{ref}} = 288.15\,\mathrm{K}$ — temperatura referencia
\item $r = 0.276\,\mathrm{m}$ — radio
\end{itemize}

\noindent\textbf{Incógnita}
$U$, en $\mathrm{m\,s^{-1}}$; también las magnitudes intermedias desarrolladas.

\noindent\textbf{Modelo físico}

\noindent\textit{Corrección de velocidad}

\[
\fitmath{U=2 \pi N r}
\]

\noindent\textbf{Hipótesis}
\begin{itemize}
\item Giro en revoluciones por segundo y radio fijo.
\end{itemize}

\noindent\textbf{Desarrollo paso a paso}

\begin{enumerate}

\item La corrección compara velocidades a la referencia térmica.

\[\fitmath{\mathrm{Nc}=\frac{\mathrm{nn}}{\sqrt{\frac{\mathrm{Tin}}{\mathrm{Tr}}}}}\]

\[\fitmath{\mathrm{Nc}=\frac{(196.6666667)}{\sqrt{\frac{(259)}{(288.15)}}}\approx 207.4389\,\mathrm{s^{-1}}}\]

\item El giro en revoluciones por segundo se convierte a velocidad angular para U.

\[\fitmath{U=2\cdot \pi\cdot \mathrm{nn}\cdot \mathrm{rad}}\]

\[\fitmath{U=2\cdot \pi\cdot (196.6666667)\cdot (0.276)\approx 341.0513\,\mathrm{m/s}}\]

\end{enumerate}


\noindent\textbf{Sustitución numérica}

\[
\fitmath{2 \pi \left(196.66666666667\,\mathrm{s}^{-1}\right) \left(0.276\,\mathrm{m}\right)}
\]

\noindent\textbf{Resultado}

\[
\fitmath{U=3.410512985\times10^{2}\,\mathrm{m}\,\mathrm{s}^{-1}}
\]

\noindent\textbf{Interpretación}
El dato se expresa en revoluciones por segundo; U se calcula con N físico, no con N corregido.

\noindent\textbf{Comprobación de los pasos}
\begin{itemize}
\item Paso 2: $207.439$ — reproducido
\item Paso 4: $341.051$ — reproducido
\end{itemize}

\noindent\textbf{Comprobación de la respuesta} la aritmética del último paso reproduce la respuesta enunciada a la precisión con que está escrita. Es una reproducción aritmética; no valida el modelo físico ni sus hipótesis.

\noindent\textbf{Limitaciones}
\begin{itemize}
\item Modelo educativo; la aplicabilidad a una máquina real requiere revisión experta.
\end{itemize}

\noindent\textbf{Errores frecuentes}
\begin{itemize}
\item Usar N corregido en una velocidad mecánica real.
\end{itemize}
% studium:end-record eyJpZCI6ICJQUkItMDA3NyIsICJzZWN0aW9uIjogImMxMCJ9
\end{tcolorbox}

\begin{tcolorbox}[title={Problema resuelto 4}, studiumworked, breakable]

\noindent\textbf{Enunciado}
\(C10-04-R \cdot\) Estabilidad lineal del modelo ampliado. Para el modelo de dos estados del capítulo, calcule traza y determinante de la matriz lineal. El resultado final es el determinante; interprete estabilidad. Datos: \(a_c=-1800\, \mathrm{Pa}\,\mathrm{s}\,\mathrm{kg}^{-1}\); \(b_v=0.00016\, \mathrm{kg}\,\mathrm{s}^{-1}\,\mathrm{Pa}^{-1}\); \(L_c=260\, \mathrm{Pa}\,\mathrm{s}^{2}\,\mathrm{kg}^{-1}\); \(K_p=0.0018\, \mathrm{kg}\,\mathrm{Pa}^{-1}\).

\noindent\textbf{Datos}
\begin{itemize}
\item $a_c = -1800\,\mathrm{Pa}\,\mathrm{s}\,\mathrm{kg}^{-1}$ — pendiente de compresor
\item $b_v = 0.00016\,\mathrm{kg}\,\mathrm{s}^{-1}\,\mathrm{Pa}^{-1}$ — pendiente de descarga
\item $L_c = 260\,\mathrm{Pa}\,\mathrm{s}^{2}\,\mathrm{kg}^{-1}$ — inercia del modelo
\item $K_p = 0.0018\,\mathrm{kg}\,\mathrm{Pa}^{-1}$ — capacidad del modelo
\end{itemize}

\noindent\textbf{Incógnita}
$\det(J)$, en $\mathrm{s^{-2}}$; también las magnitudes intermedias desarrolladas.

\noindent\textbf{Modelo físico}

\noindent\textit{Estabilidad lineal del modelo ampliado}

\[
\fitmath{\det(J)=\frac{- a_c b_v + 1}{K_p L_c}}
\]

\noindent\textbf{Hipótesis}
\begin{itemize}
\item Linealización local; parámetros positivos L,K; modelo pedagógico sin retardo.
\end{itemize}

\noindent\textbf{Desarrollo paso a paso}

\begin{enumerate}

\item La traza suma los términos diagonales.

\[\fitmath{\mathrm{tr}=\frac{\mathrm{aa}}{\mathrm{ll}}-\frac{\mathrm{bb}}{\mathrm{kk}}}\]

\[\fitmath{\mathrm{tr}=\frac{(-1800)}{(260)}-\frac{(0.00016)}{(0.0018)}\approx -7.011966\,\mathrm{s^{-1}}}\]

\item El determinante incluye el acoplamiento cruzado.

\[\fitmath{\mathrm{det}=\frac{1-\mathrm{aa}\cdot \mathrm{bb}}{\mathrm{ll}\cdot \mathrm{kk}}}\]

\[\fitmath{\mathrm{det}=\frac{1-(-1800)\cdot (0.00016)}{(260)\cdot (0.0018)}\approx 2.752137\,\mathrm{s^{-2}}}\]

\end{enumerate}


\noindent\textbf{Sustitución numérica}

\[
\fitmath{\frac{- \left(-1800\,\mathrm{Pa}\,\mathrm{s}\,\mathrm{kg}^{-1}\right) \left(0.00016\,\mathrm{kg}\,\mathrm{s}^{-1}\,\mathrm{Pa}^{-1}\right) + 1}{\left(0.0018\,\mathrm{kg}\,\mathrm{Pa}^{-1}\right) \left(260\,\mathrm{Pa}\,\mathrm{s}^{2}\,\mathrm{kg}^{-1}\right)}}
\]

\noindent\textbf{Resultado}

\[
\fitmath{\det(J)=2.752136752\times10^{0}\,\mathrm{s}^{-2}}
\]

\noindent\textbf{Interpretación}
Traza negativa y determinante positivo dan estabilidad lineal del modelo de dos estados. Esto no predice pérdida rotativa ni estabilidad de un motor real.

\noindent\textbf{Comprobación de los pasos}
\begin{itemize}
\item Paso 2: $-7.01197$ — reproducido
\item Paso 4: $2.75214$ — reproducido
\end{itemize}

\noindent\textbf{Comprobación de la respuesta} la aritmética del último paso reproduce la respuesta enunciada a la precisión con que está escrita. Es una reproducción aritmética; no valida el modelo físico ni sus hipótesis.

\noindent\textbf{Limitaciones}
\begin{itemize}
\item Modelo educativo; la aplicabilidad a una máquina real requiere revisión experta.
\end{itemize}

\noindent\textbf{Errores frecuentes}
\begin{itemize}
\item Tomar L como longitud sin unidades equivalentes.
\end{itemize}
% studium:end-record eyJpZCI6ICJQUkItMDA3OSIsICJzZWN0aW9uIjogImMxMCJ9
\end{tcolorbox}

\begin{tcolorbox}[title={Ejercicios propuestos}, studiumworked, breakable]

\noindent\textbf{1 (Básico, Desarrollo)}
\(C10-01-P \cdot\) Incidencia durante variación de caudal. Ejercicio de transferencia. Un rotor sin prerrotación conserva velocidad periférica y ángulo metálico. Calcule su beta de entrada y la incidencia a la nueva velocidad axial. Antes de consultar la solución, prediga el efecto de cambiar el primer dato respecto del caso resuelto de esta familia. Justifique la tendencia y explique qué restricción podría convertirse en limitante. Datos: \(U=263.055\, \mathrm{m}\,\mathrm{s}^{-1}\); \(C_x=119\, \mathrm{m}\,\mathrm{s}^{-1}\); \(\kappa_e=-1.064651\, \mathrm{rad}\).

\noindent\textbf{2 (Intermedio, Desarrollo)}
\(C10-02-P \cdot\) Margen de caudal definido. Ejercicio de transferencia. A igual velocidad corregida, un punto operativo y un límite prescrito tienen los caudales indicados. Calcule el margen según la definición del capítulo. Antes de consultar la solución, prediga el efecto de cambiar el primer dato respecto del caso resuelto de esta familia. Justifique la tendencia y explique qué restricción podría convertirse en limitante. Datos: \(\dot{m}_{c,\mathrm{op}}=30.246\, \mathrm{kg}\,\mathrm{s}^{-1}\); \(\dot{m}_{c,\mathrm{lim}}=23.9\, \mathrm{kg}\,\mathrm{s}^{-1}\).

\noindent\textbf{3 (Nivel de examen, Desarrollo)}
\(C10-04-P \cdot\) Estabilidad lineal del modelo ampliado. Ejercicio de transferencia. Para el modelo de dos estados del capítulo, calcule traza y determinante de la matriz lineal. El resultado final es el determinante; interprete estabilidad. Antes de consultar la solución, prediga el efecto de cambiar el primer dato respecto del caso resuelto de esta familia. Justifique la tendencia y explique qué restricción podría convertirse en limitante. Datos: \(a_c=-1917\, \mathrm{Pa}\,\mathrm{s}\,\mathrm{kg}^{-1}\); \(b_v=0.00016\, \mathrm{kg}\,\mathrm{s}^{-1}\,\mathrm{Pa}^{-1}\); \(L_c=260\, \mathrm{Pa}\,\mathrm{s}^{2}\,\mathrm{kg}^{-1}\); \(K_p=0.0018\, \mathrm{kg}\,\mathrm{Pa}^{-1}\).

\noindent\textbf{4 (Avanzado, Desarrollo)}
\(C10-03-P \cdot\) Corrección de velocidad. Ejercicio de transferencia. Halle la velocidad corregida y la velocidad periférica física al radio dado. La incógnita final es U real. Antes de consultar la solución, prediga el efecto de cambiar el primer dato respecto del caso resuelto de esta familia. Justifique la tendencia y explique qué restricción podría convertirse en limitante. Datos: \(N=209.45\, \mathrm{s}^{-1}\); \(T_{0e}=259\, \mathrm{K}\); \(T_{\mathrm{ref}}=288.15\, \mathrm{K}\); \(r=0.276\, \mathrm{m}\).
\end{tcolorbox}


\studiumchaptercheck

\part{Estudio tridimensional de turbomáquinas axiales}

\chapter{Equilibrio radial y leyes de remolino}
\label{chap:c11}

\studiumchapterguide{Objetivo: desarrollar equilibrio radial y leyes de remolino, declarar hipótesis y comprobar condiciones físicas.}
% studium:end-record eyJpZCI6ICJQQVItMDA3MyIsICJzZWN0aW9uIjogImMxMSJ9

\section{Introducción y objetivos}

El radio medio permite estimar un escalón, pero cada sección de una pala tiene distinta velocidad periférica. El equilibrio radial relaciona el gradiente de presión con el remolino y limita las distribuciones que pueden elegirse. El objetivo es deducir esa condición, combinarla con energía y construir leyes radiales coherentes. Pertenece al bloque IV.
El modelo de equilibrio radial simple desprecia aceleración meridiana radial y curvatura importante de las líneas de corriente. No reproduce los flujos secundarios, las capas límite de cubo y carcasa ni la migración radial de pérdidas. Su utilidad es comprobar consistencia antes de un cálculo tridimensional más completo.
% studium:end-record eyJpZCI6ICJQQVItMDA3NCIsICJzZWN0aW9uIjogImMxMSJ9

\section{Balance radial}

En coordenadas cilíndricas, un flujo axisimétrico con velocidad radial despreciable necesita aceleración centrípeta \(-C_\theta^2/r\). El balance de cantidad de movimiento radial conduce a

\[
\fitmath{\frac{\mathrm{d}p}{dr}=\rho\frac{C_\theta^2}{r}.}
\]
La presión crece hacia fuera incluso si el remolino tiene signo negativo, pues aparece su cuadrado. El gradiente no cambia de signo al invertir el sentido de giro de la corriente. Una presión uniforme en todo el radio es incompatible con remolino no nulo dentro de este modelo simple.
La relación de Gibbs es \(\mathrm{d}h=T\,\mathrm{d}s+\mathrm{d}p/\rho\). Si \(h_0=h+(C_x^2+C_\theta^2)/2\), su derivada radial, combinada con el equilibrio, produce

\[
\fitmath{\frac{dh_0}{dr}=T\frac{\mathrm{d}s}{dr}+C_x\frac{dC_x}{dr}+\frac{C_\theta}{r}\frac{d(rC_\theta)}{dr}.}
\]
Esta forma muestra qué hipótesis adicionales se necesitan para afirmar que \(C_x\) es constante con el radio. No es una consecuencia del equilibrio radial por sí solo.
% studium:end-record eyJpZCI6ICJQQVItMDA3NSIsICJzZWN0aW9uIjogImMxMSJ9

\section{Vórtice libre}

Si \(rC_\theta=K_v\), el último término se anula. Si además \(h_0\) y \(s\) son uniformes radialmente, resulta \(C_xdC_x/dr=0\), es decir, \(C_x\) constante en una región de flujo axial no nulo. La ley de vórtice libre debe aplicarse por estación; los valores de \(K_v\) pueden cambiar al atravesar un rotor o un estator.
Para densidad constante como aproximación de integración,

\[
\fitmath{p(r)-p(r_m)=\frac{\rho K_v^2}{2}\left(\frac1{r_m^2}-\frac1{r^2}\right).}
\]
En gas compresible la densidad debe calcularse con el estado termodinámico y no extraerse de la integral sin justificarlo. La solución anterior es útil como comprobación de signo y dimensiones, no como descripción completa de un escalón caliente.
Si la entrada y salida del rotor obedecen vórtice libre, Euler da \(w_f=\Omega(K_{v,s}-K_{v,e})\), independiente del radio. Eso proporciona trabajo radialmente uniforme, pero no reacción uniforme. La figura~\ref{fig:radial} presenta velocidades calculadas para una ley libre y muestra cómo cambia la orientación relativa.
% studium:end-record eyJpZCI6ICJQQVItMDA3NiIsICJzZWN0aW9uIjogImMxMSJ9

\par\vspace{\baselineskip}
\noindent\begin{minipage}{\linewidth}
\centering
\begingroup
\setbox0=\hbox{\begin{tikzpicture}
\draw[->,gray] (0,0)--(7,0);\draw[->,gray] (0,0)--(0,3.8);\draw (0.0,0)--(0.0,-.08) node[below,font=\scriptsize] {0.75};\draw (0,0.0)--(-.08,0.0) node[left,font=\scriptsize] {0};\draw (1.5,0)--(1.5,-.08) node[below,font=\scriptsize] {0.875};\draw (0,0.75)--(-.08,0.75) node[left,font=\scriptsize] {0.375};\draw (3.0,0)--(3.0,-.08) node[below,font=\scriptsize] {1};\draw (0,1.5)--(-.08,1.5) node[left,font=\scriptsize] {0.75};\draw (4.5,0)--(4.5,-.08) node[below,font=\scriptsize] {1.12};\draw (0,2.25)--(-.08,2.25) node[left,font=\scriptsize] {1.12};\draw (6.0,0)--(6.0,-.08) node[below,font=\scriptsize] {1.25};\draw (0,3.0)--(-.08,3.0) node[left,font=\scriptsize] {1.5};\draw[blue,very thick] plot coordinates {(0.00000,1.33333) (0.10000,1.31868) (0.20000,1.30435) (0.30000,1.29032) (0.40000,1.27660) (0.50000,1.26316) (0.60000,1.25000) (0.70000,1.23711) (0.80000,1.22449) (0.90000,1.21212) (1.00000,1.20000) (1.10000,1.18812) (1.20000,1.17647) (1.30000,1.16505) (1.40000,1.15385) (1.50000,1.14286) (1.60000,1.13208) (1.70000,1.12150) (1.80000,1.11111) (1.90000,1.10092) (2.00000,1.09091) (2.10000,1.08108) (2.20000,1.07143) (2.30000,1.06195) (2.40000,1.05263) (2.50000,1.04348) (2.60000,1.03448) (2.70000,1.02564) (2.80000,1.01695) (2.90000,1.00840) (3.00000,1.00000) (3.10000,0.99174) (3.20000,0.98361) (3.30000,0.97561) (3.40000,0.96774) (3.50000,0.96000) (3.60000,0.95238) (3.70000,0.94488) (3.80000,0.93750) (3.90000,0.93023) (4.00000,0.92308) (4.10000,0.91603) (4.20000,0.90909) (4.30000,0.90226) (4.40000,0.89552) (4.50000,0.88889) (4.60000,0.88235) (4.70000,0.87591) (4.80000,0.86957) (4.90000,0.86331) (5.00000,0.85714) (5.10000,0.85106) (5.20000,0.84507) (5.30000,0.83916) (5.40000,0.83333) (5.50000,0.82759) (5.60000,0.82192) (5.70000,0.81633) (5.80000,0.81081) (5.90000,0.80537) (6.00000,0.80000)};\node[below,align=center] at (3.5,-.65) {Radio $r/r_m$};\node[rotate=90] at (-1,2) {Velocidad / $(\Omega r_m)$};\draw[orange,thick] plot coordinates {(0.00000,1.50000) (0.10000,1.51667) (0.20000,1.53333) (0.30000,1.55000) (0.40000,1.56667) (0.50000,1.58333) (0.60000,1.60000) (0.70000,1.61667) (0.80000,1.63333) (0.90000,1.65000) (1.00000,1.66667) (1.10000,1.68333) (1.20000,1.70000) (1.30000,1.71667) (1.40000,1.73333) (1.50000,1.75000) (1.60000,1.76667) (1.70000,1.78333) (1.80000,1.80000) (1.90000,1.81667) (2.00000,1.83333) (2.10000,1.85000) (2.20000,1.86667) (2.30000,1.88333) (2.40000,1.90000) (2.50000,1.91667) (2.60000,1.93333) (2.70000,1.95000) (2.80000,1.96667) (2.90000,1.98333) (3.00000,2.00000) (3.10000,2.01667) (3.20000,2.03333) (3.30000,2.05000) (3.40000,2.06667) (3.50000,2.08333) (3.60000,2.10000) (3.70000,2.11667) (3.80000,2.13333) (3.90000,2.15000) (4.00000,2.16667) (4.10000,2.18333) (4.20000,2.20000) (4.30000,2.21667) (4.40000,2.23333) (4.50000,2.25000) (4.60000,2.26667) (4.70000,2.28333) (4.80000,2.30000) (4.90000,2.31667) (5.00000,2.33333) (5.10000,2.35000) (5.20000,2.36667) (5.30000,2.38333) (5.40000,2.40000) (5.50000,2.41667) (5.60000,2.43333) (5.70000,2.45000) (5.80000,2.46667) (5.90000,2.48333) (6.00000,2.50000)};\node[orange] at (5,2.9) {$U/(\Omega r_m)$};\node[blue] at (4.5,1.3) {$C_\theta/(\Omega r_m)$};
\end{tikzpicture}}%
\count0=\ht0
\advance\count0 by \dp0
\count2=\wd0
\count4=\dimexpr\textheight/100*38\relax
\count6=\linewidth
\ifnum\count0<1
\count0=1
\fi
\ifnum\count2<1
\count2=1
\fi
\divide\count0 by 4096
\divide\count2 by 4096
\divide\count4 by 4096
\divide\count6 by 4096
\ifnum\count0<1
\count0=1
\fi
\ifnum\count2<1
\count2=1
\fi
\ifnum\count4<1
\count4=1
\fi
\ifnum\count6<1
\count6=1
\fi
\count8=\count2
\ifnum\count0>\count4
\count8=\numexpr\count2*\count4/\count0\relax
\fi
\ifnum\count8>\count6
\resizebox{\linewidth}{!}{\begin{tikzpicture}
\draw[->,gray] (0,0)--(7,0);\draw[->,gray] (0,0)--(0,3.8);\draw (0.0,0)--(0.0,-.08) node[below,font=\scriptsize] {0.75};\draw (0,0.0)--(-.08,0.0) node[left,font=\scriptsize] {0};\draw (1.5,0)--(1.5,-.08) node[below,font=\scriptsize] {0.875};\draw (0,0.75)--(-.08,0.75) node[left,font=\scriptsize] {0.375};\draw (3.0,0)--(3.0,-.08) node[below,font=\scriptsize] {1};\draw (0,1.5)--(-.08,1.5) node[left,font=\scriptsize] {0.75};\draw (4.5,0)--(4.5,-.08) node[below,font=\scriptsize] {1.12};\draw (0,2.25)--(-.08,2.25) node[left,font=\scriptsize] {1.12};\draw (6.0,0)--(6.0,-.08) node[below,font=\scriptsize] {1.25};\draw (0,3.0)--(-.08,3.0) node[left,font=\scriptsize] {1.5};\draw[blue,very thick] plot coordinates {(0.00000,1.33333) (0.10000,1.31868) (0.20000,1.30435) (0.30000,1.29032) (0.40000,1.27660) (0.50000,1.26316) (0.60000,1.25000) (0.70000,1.23711) (0.80000,1.22449) (0.90000,1.21212) (1.00000,1.20000) (1.10000,1.18812) (1.20000,1.17647) (1.30000,1.16505) (1.40000,1.15385) (1.50000,1.14286) (1.60000,1.13208) (1.70000,1.12150) (1.80000,1.11111) (1.90000,1.10092) (2.00000,1.09091) (2.10000,1.08108) (2.20000,1.07143) (2.30000,1.06195) (2.40000,1.05263) (2.50000,1.04348) (2.60000,1.03448) (2.70000,1.02564) (2.80000,1.01695) (2.90000,1.00840) (3.00000,1.00000) (3.10000,0.99174) (3.20000,0.98361) (3.30000,0.97561) (3.40000,0.96774) (3.50000,0.96000) (3.60000,0.95238) (3.70000,0.94488) (3.80000,0.93750) (3.90000,0.93023) (4.00000,0.92308) (4.10000,0.91603) (4.20000,0.90909) (4.30000,0.90226) (4.40000,0.89552) (4.50000,0.88889) (4.60000,0.88235) (4.70000,0.87591) (4.80000,0.86957) (4.90000,0.86331) (5.00000,0.85714) (5.10000,0.85106) (5.20000,0.84507) (5.30000,0.83916) (5.40000,0.83333) (5.50000,0.82759) (5.60000,0.82192) (5.70000,0.81633) (5.80000,0.81081) (5.90000,0.80537) (6.00000,0.80000)};\node[below,align=center] at (3.5,-.65) {Radio $r/r_m$};\node[rotate=90] at (-1,2) {Velocidad / $(\Omega r_m)$};\draw[orange,thick] plot coordinates {(0.00000,1.50000) (0.10000,1.51667) (0.20000,1.53333) (0.30000,1.55000) (0.40000,1.56667) (0.50000,1.58333) (0.60000,1.60000) (0.70000,1.61667) (0.80000,1.63333) (0.90000,1.65000) (1.00000,1.66667) (1.10000,1.68333) (1.20000,1.70000) (1.30000,1.71667) (1.40000,1.73333) (1.50000,1.75000) (1.60000,1.76667) (1.70000,1.78333) (1.80000,1.80000) (1.90000,1.81667) (2.00000,1.83333) (2.10000,1.85000) (2.20000,1.86667) (2.30000,1.88333) (2.40000,1.90000) (2.50000,1.91667) (2.60000,1.93333) (2.70000,1.95000) (2.80000,1.96667) (2.90000,1.98333) (3.00000,2.00000) (3.10000,2.01667) (3.20000,2.03333) (3.30000,2.05000) (3.40000,2.06667) (3.50000,2.08333) (3.60000,2.10000) (3.70000,2.11667) (3.80000,2.13333) (3.90000,2.15000) (4.00000,2.16667) (4.10000,2.18333) (4.20000,2.20000) (4.30000,2.21667) (4.40000,2.23333) (4.50000,2.25000) (4.60000,2.26667) (4.70000,2.28333) (4.80000,2.30000) (4.90000,2.31667) (5.00000,2.33333) (5.10000,2.35000) (5.20000,2.36667) (5.30000,2.38333) (5.40000,2.40000) (5.50000,2.41667) (5.60000,2.43333) (5.70000,2.45000) (5.80000,2.46667) (5.90000,2.48333) (6.00000,2.50000)};\node[orange] at (5,2.9) {$U/(\Omega r_m)$};\node[blue] at (4.5,1.3) {$C_\theta/(\Omega r_m)$};
\end{tikzpicture}}%
\else
\ifnum\count0>\count4
\resizebox{!}{\dimexpr\textheight/100*38\relax}{\begin{tikzpicture}
\draw[->,gray] (0,0)--(7,0);\draw[->,gray] (0,0)--(0,3.8);\draw (0.0,0)--(0.0,-.08) node[below,font=\scriptsize] {0.75};\draw (0,0.0)--(-.08,0.0) node[left,font=\scriptsize] {0};\draw (1.5,0)--(1.5,-.08) node[below,font=\scriptsize] {0.875};\draw (0,0.75)--(-.08,0.75) node[left,font=\scriptsize] {0.375};\draw (3.0,0)--(3.0,-.08) node[below,font=\scriptsize] {1};\draw (0,1.5)--(-.08,1.5) node[left,font=\scriptsize] {0.75};\draw (4.5,0)--(4.5,-.08) node[below,font=\scriptsize] {1.12};\draw (0,2.25)--(-.08,2.25) node[left,font=\scriptsize] {1.12};\draw (6.0,0)--(6.0,-.08) node[below,font=\scriptsize] {1.25};\draw (0,3.0)--(-.08,3.0) node[left,font=\scriptsize] {1.5};\draw[blue,very thick] plot coordinates {(0.00000,1.33333) (0.10000,1.31868) (0.20000,1.30435) (0.30000,1.29032) (0.40000,1.27660) (0.50000,1.26316) (0.60000,1.25000) (0.70000,1.23711) (0.80000,1.22449) (0.90000,1.21212) (1.00000,1.20000) (1.10000,1.18812) (1.20000,1.17647) (1.30000,1.16505) (1.40000,1.15385) (1.50000,1.14286) (1.60000,1.13208) (1.70000,1.12150) (1.80000,1.11111) (1.90000,1.10092) (2.00000,1.09091) (2.10000,1.08108) (2.20000,1.07143) (2.30000,1.06195) (2.40000,1.05263) (2.50000,1.04348) (2.60000,1.03448) (2.70000,1.02564) (2.80000,1.01695) (2.90000,1.00840) (3.00000,1.00000) (3.10000,0.99174) (3.20000,0.98361) (3.30000,0.97561) (3.40000,0.96774) (3.50000,0.96000) (3.60000,0.95238) (3.70000,0.94488) (3.80000,0.93750) (3.90000,0.93023) (4.00000,0.92308) (4.10000,0.91603) (4.20000,0.90909) (4.30000,0.90226) (4.40000,0.89552) (4.50000,0.88889) (4.60000,0.88235) (4.70000,0.87591) (4.80000,0.86957) (4.90000,0.86331) (5.00000,0.85714) (5.10000,0.85106) (5.20000,0.84507) (5.30000,0.83916) (5.40000,0.83333) (5.50000,0.82759) (5.60000,0.82192) (5.70000,0.81633) (5.80000,0.81081) (5.90000,0.80537) (6.00000,0.80000)};\node[below,align=center] at (3.5,-.65) {Radio $r/r_m$};\node[rotate=90] at (-1,2) {Velocidad / $(\Omega r_m)$};\draw[orange,thick] plot coordinates {(0.00000,1.50000) (0.10000,1.51667) (0.20000,1.53333) (0.30000,1.55000) (0.40000,1.56667) (0.50000,1.58333) (0.60000,1.60000) (0.70000,1.61667) (0.80000,1.63333) (0.90000,1.65000) (1.00000,1.66667) (1.10000,1.68333) (1.20000,1.70000) (1.30000,1.71667) (1.40000,1.73333) (1.50000,1.75000) (1.60000,1.76667) (1.70000,1.78333) (1.80000,1.80000) (1.90000,1.81667) (2.00000,1.83333) (2.10000,1.85000) (2.20000,1.86667) (2.30000,1.88333) (2.40000,1.90000) (2.50000,1.91667) (2.60000,1.93333) (2.70000,1.95000) (2.80000,1.96667) (2.90000,1.98333) (3.00000,2.00000) (3.10000,2.01667) (3.20000,2.03333) (3.30000,2.05000) (3.40000,2.06667) (3.50000,2.08333) (3.60000,2.10000) (3.70000,2.11667) (3.80000,2.13333) (3.90000,2.15000) (4.00000,2.16667) (4.10000,2.18333) (4.20000,2.20000) (4.30000,2.21667) (4.40000,2.23333) (4.50000,2.25000) (4.60000,2.26667) (4.70000,2.28333) (4.80000,2.30000) (4.90000,2.31667) (5.00000,2.33333) (5.10000,2.35000) (5.20000,2.36667) (5.30000,2.38333) (5.40000,2.40000) (5.50000,2.41667) (5.60000,2.43333) (5.70000,2.45000) (5.80000,2.46667) (5.90000,2.48333) (6.00000,2.50000)};\node[orange] at (5,2.9) {$U/(\Omega r_m)$};\node[blue] at (4.5,1.3) {$C_\theta/(\Omega r_m)$};
\end{tikzpicture}}%
\else
\begin{tikzpicture}
\draw[->,gray] (0,0)--(7,0);\draw[->,gray] (0,0)--(0,3.8);\draw (0.0,0)--(0.0,-.08) node[below,font=\scriptsize] {0.75};\draw (0,0.0)--(-.08,0.0) node[left,font=\scriptsize] {0};\draw (1.5,0)--(1.5,-.08) node[below,font=\scriptsize] {0.875};\draw (0,0.75)--(-.08,0.75) node[left,font=\scriptsize] {0.375};\draw (3.0,0)--(3.0,-.08) node[below,font=\scriptsize] {1};\draw (0,1.5)--(-.08,1.5) node[left,font=\scriptsize] {0.75};\draw (4.5,0)--(4.5,-.08) node[below,font=\scriptsize] {1.12};\draw (0,2.25)--(-.08,2.25) node[left,font=\scriptsize] {1.12};\draw (6.0,0)--(6.0,-.08) node[below,font=\scriptsize] {1.25};\draw (0,3.0)--(-.08,3.0) node[left,font=\scriptsize] {1.5};\draw[blue,very thick] plot coordinates {(0.00000,1.33333) (0.10000,1.31868) (0.20000,1.30435) (0.30000,1.29032) (0.40000,1.27660) (0.50000,1.26316) (0.60000,1.25000) (0.70000,1.23711) (0.80000,1.22449) (0.90000,1.21212) (1.00000,1.20000) (1.10000,1.18812) (1.20000,1.17647) (1.30000,1.16505) (1.40000,1.15385) (1.50000,1.14286) (1.60000,1.13208) (1.70000,1.12150) (1.80000,1.11111) (1.90000,1.10092) (2.00000,1.09091) (2.10000,1.08108) (2.20000,1.07143) (2.30000,1.06195) (2.40000,1.05263) (2.50000,1.04348) (2.60000,1.03448) (2.70000,1.02564) (2.80000,1.01695) (2.90000,1.00840) (3.00000,1.00000) (3.10000,0.99174) (3.20000,0.98361) (3.30000,0.97561) (3.40000,0.96774) (3.50000,0.96000) (3.60000,0.95238) (3.70000,0.94488) (3.80000,0.93750) (3.90000,0.93023) (4.00000,0.92308) (4.10000,0.91603) (4.20000,0.90909) (4.30000,0.90226) (4.40000,0.89552) (4.50000,0.88889) (4.60000,0.88235) (4.70000,0.87591) (4.80000,0.86957) (4.90000,0.86331) (5.00000,0.85714) (5.10000,0.85106) (5.20000,0.84507) (5.30000,0.83916) (5.40000,0.83333) (5.50000,0.82759) (5.60000,0.82192) (5.70000,0.81633) (5.80000,0.81081) (5.90000,0.80537) (6.00000,0.80000)};\node[below,align=center] at (3.5,-.65) {Radio $r/r_m$};\node[rotate=90] at (-1,2) {Velocidad / $(\Omega r_m)$};\draw[orange,thick] plot coordinates {(0.00000,1.50000) (0.10000,1.51667) (0.20000,1.53333) (0.30000,1.55000) (0.40000,1.56667) (0.50000,1.58333) (0.60000,1.60000) (0.70000,1.61667) (0.80000,1.63333) (0.90000,1.65000) (1.00000,1.66667) (1.10000,1.68333) (1.20000,1.70000) (1.30000,1.71667) (1.40000,1.73333) (1.50000,1.75000) (1.60000,1.76667) (1.70000,1.78333) (1.80000,1.80000) (1.90000,1.81667) (2.00000,1.83333) (2.10000,1.85000) (2.20000,1.86667) (2.30000,1.88333) (2.40000,1.90000) (2.50000,1.91667) (2.60000,1.93333) (2.70000,1.95000) (2.80000,1.96667) (2.90000,1.98333) (3.00000,2.00000) (3.10000,2.01667) (3.20000,2.03333) (3.30000,2.05000) (3.40000,2.06667) (3.50000,2.08333) (3.60000,2.10000) (3.70000,2.11667) (3.80000,2.13333) (3.90000,2.15000) (4.00000,2.16667) (4.10000,2.18333) (4.20000,2.20000) (4.30000,2.21667) (4.40000,2.23333) (4.50000,2.25000) (4.60000,2.26667) (4.70000,2.28333) (4.80000,2.30000) (4.90000,2.31667) (5.00000,2.33333) (5.10000,2.35000) (5.20000,2.36667) (5.30000,2.38333) (5.40000,2.40000) (5.50000,2.41667) (5.60000,2.43333) (5.70000,2.45000) (5.80000,2.46667) (5.90000,2.48333) (6.00000,2.50000)};\node[orange] at (5,2.9) {$U/(\Omega r_m)$};\node[blue] at (4.5,1.3) {$C_\theta/(\Omega r_m)$};
\end{tikzpicture}%
\fi
\fi
\endgroup
\par\nopagebreak
\studiumcaption{fig:radial}{Velocidad periférica y remolino libre calculados con escala común y normalización por Omega rm.}
\end{minipage}
\par\vspace{\baselineskip}
% studium:end-record eyJpZCI6ICJGSUctMDAxOSIsICJzZWN0aW9uIjogImMxMSJ9

\section{Vórtice forzado y distribución axial}

Para \(C_\theta=kr\), el momento angular específico crece como \(r^2\). Si \(h_0\) y \(s\) se mantienen uniformes,

\[
\fitmath{C_x\frac{dC_x}{dr}=-2k^2r,\qquad
C_x^2(r)=C_x^2(r_m)-2k^2(r^2-r_m^2).}
\]
La componente axial disminuye hacia punta en este caso. La expresión debe mantenerse positiva en todo el intervalo. Imponer simultáneamente vórtice forzado, \(C_x\) uniforme y \(h_0,s\) uniformes contradice la ecuación radial, salvo el caso trivial sin remolino.
Leyes mixtas, como \(C_\theta=ar+b/r\), permiten otros repartos. Su selección exige integrar la ecuación radial y comprobar la energía, las condiciones de contorno y el caudal. Una ley analítica no es físicamente admisible solo por ser suave.
% studium:end-record eyJpZCI6ICJQQVItMDA3NyIsICJzZWN0aW9uIjogImMxMSJ9

\section{Reacción y raíz limitante}

En un escalón axial repetitivo con \(C_x\) constante y vórtice libre a ambos lados del rotor,

\[
\fitmath{\mathcal R(r)=1-\frac{K_{v,e}+K_{v,s}}{2\Omega r^2}.}
\]
Para remolinos de suma positiva, la reacción disminuye hacia la raíz y puede llegar a cero o hacerse negativa. Elegir medio de reacción al radio medio no fija la reacción de toda la pala. La raíz puede ser limitante por difusión, mientras la punta puede serlo por Mach relativo y holguras.
La ley radial debe evaluarse con los radios reales de cada estación. Si el anillo cambia significativamente, no basta con identificar tubos de corriente por el mismo valor de radio. Un cálculo de curvatura de líneas de corriente o de equilibrio radial completo debe seguir su desplazamiento.
% studium:end-record eyJpZCI6ICJQQVItMDA3OCIsICJzZWN0aW9uIjogImMxMSJ9

\section{Método y comprobaciones}

Se elige una estación de referencia y un estado termodinámico, se propone una ley de remolino, se integra la ecuación radial con una condición de contorno y se ajusta el nivel de \(C_x\) para satisfacer el caudal integral. Después se construyen triángulos en varios radios y se comprueban reacción, difusión y Mach. Los problemas de este capítulo separan la integración de presión del cierre del trabajo para que ninguna hipótesis quede oculta.
Errores habituales: usar \(U^2/r\) en lugar de \(C_\theta^2/r\) en el equilibrio de la corriente absoluta, conservar el mismo \(K_v\) a través de un rotor que realiza trabajo y asumir \(C_x\) uniforme bajo cualquier ley torsional. Autocomprobación: explique por qué trabajo uniforme y reacción uniforme son condiciones distintas.
Formulario: \(\mathrm{d}p/dr=\rho C_\theta^2/r\), ecuación radial energética anterior y \(rC_\theta=K_v\) para vórtice libre. Referencias: ULE Diseño de compresores, pp. \(33-35\); MIT 16.50, lección 26, pp. \(1-3\), como contraste del efecto radial. Las integraciones y las curvas del manual son originales.
% studium:end-record eyJpZCI6ICJQQVItMDA3OSIsICJzZWN0aW9uIjogImMxMSJ9

\begin{tcolorbox}[title={Desarrollos del modelo}, studiumdef, breakable]

\noindent\textbf{Modelo desarrollado: Presión de vórtice libre}

\begin{equation}
\Delta p=\frac{K_v^{\,2}\,\rho}{2}\left(\frac{1}{r_m^2}-\frac{1}{r_t^2}\right)
\label{eq:EC-C11}
\end{equation}

\noindent\textit{Significado.} Diferencia de presión entre dos radios en un vórtice libre. $K_v=rC_\theta$ es el identificador del vórtice libre, en $\mathrm{m^2\,s^{-1}}$; no es el coeficiente de vórtice forzado, cuya unidad es $\mathrm{s^{-1}}$.

\noindent\textbf{Hipótesis}
\begin{itemize}
\item Equilibrio radial simple, axisimetría, densidad constante.
\end{itemize}

\noindent\textbf{Principios}
\begin{itemize}
\item Balance y definiciones físicas desarrollados en el capítulo; propiedades y dominio declarados en el problema.
\end{itemize}

\noindent\textbf{Aplicabilidad}
\begin{itemize}
\item Modelo educativo del caso C11-01-R
\end{itemize}

\noindent\textbf{Limitaciones}
\begin{itemize}
\item Modelo educativo; la aplicabilidad a una máquina real requiere revisión experta.
\end{itemize}

\noindent\textbf{Pasos}
\begin{enumerate}
\item
El momento angular constante fija $C_\theta$ en la punta.

\[
\fitmath{C_\theta=\frac{K_v}{r_t}}
\]
\item
Se integra el gradiente positivo de presión.

\[
\fitmath{\Delta p=\frac{K_v^{2} \rho \left(- \frac{1}{r_t^{2}} + \frac{1}{r_m^{2}}\right)}{2}}
\]
\end{enumerate}

{\noindent\footnotesize verificada dimensionalmente.\par}
\end{tcolorbox}

\begin{tcolorbox}[title={Problema resuelto 1}, studiumworked, breakable]

\noindent\textbf{Enunciado}
\(C11-01-R \cdot\) Presión de vórtice libre. Integre el equilibrio radial de un fluido de densidad constante y remolino K/r entre radio medio y punta. Datos: \(\rho=2.3\, \mathrm{kg}\,\mathrm{m}^{-3}\); \(K_v=36\, \mathrm{m}^{2}\,\mathrm{s}^{-1}\); \(r_m=0.29\, \mathrm{m}\); \(r_t=0.36\, \mathrm{m}\).

\noindent\textbf{Datos}
\begin{itemize}
\item $\rho = 2.3\,\mathrm{kg}\,\mathrm{m}^{-3}$ — densidad uniforme
\item $K_v = 36\,\mathrm{m}^{2}\,\mathrm{s}^{-1}$ — constante de remolino
\item $r_m = 0.29\,\mathrm{m}$ — radio referencia
\item $r_t = 0.36\,\mathrm{m}$ — radio exterior
\end{itemize}

\noindent\textbf{Incógnita}
$\Delta p$, en $\mathrm{Pa}$; también las magnitudes intermedias desarrolladas.

\noindent\textbf{Modelo físico}

\noindent\textit{Presión de vórtice libre}

\[
\fitmath{\Delta p=\frac{K_v^{2} \rho \left(- \frac{1}{r_t^{2}} + \frac{1}{r_m^{2}}\right)}{2}}
\]

\noindent\textbf{Hipótesis}
\begin{itemize}
\item Equilibrio radial simple, axisimetría, densidad constante.
\end{itemize}

\noindent\textbf{Desarrollo paso a paso}

\begin{enumerate}

\item El momento angular constante fija Ctheta en la punta.

\[\fitmath{\mathrm{ct}=\frac{\mathrm{kv}}{\mathrm{rt}}}\]

\[\fitmath{\mathrm{ct}=\frac{(36)}{(0.36)}\approx 100\,\mathrm{m/s}}\]

\item Se integra el gradiente positivo de presión.

\[\fitmath{\mathrm{dp}=\frac{\rho\cdot {\mathrm{kv}}^{2}\cdot \left(\frac{1}{{\mathrm{rm}}^{2}}-\frac{1}{{\mathrm{rt}}^{2}}\right)}{2}}\]

\[\fitmath{\mathrm{dp}=\frac{(2.3)\cdot {(36)}^{2}\cdot \left(\frac{1}{{(0.29)}^{2}}-\frac{1}{{(0.36)}^{2}}\right)}{2}\approx 6221.76\,\mathrm{Pa}}\]

\end{enumerate}


\noindent\textbf{Sustitución numérica}

\[
\fitmath{\frac{\left(\left(36\,\mathrm{m}^{2}\,\mathrm{s}^{-1}\right)\right)^{2} \left(2.3\,\mathrm{kg}\,\mathrm{m}^{-3}\right) \left(- \frac{1}{\left(0.36\,\mathrm{m}\right)^{2}} + \frac{1}{\left(0.29\,\mathrm{m}\right)^{2}}\right)}{2}}
\]

\noindent\textbf{Resultado}

\[
\fitmath{\Delta p=6.221759810\times10^{3}\,\mathrm{Pa}}
\]

\noindent\textbf{Interpretación}
La presión exterior aumenta. La aproximación de densidad constante debe justificarse antes de aplicarla a un gas real.

\noindent\textbf{Comprobación de los pasos}
\begin{itemize}
\item Paso 2: $100$ — reproducido
\item Paso 4: $6221.76$ — reproducido
\end{itemize}

\noindent\textbf{Comprobación de la respuesta} la aritmética del último paso reproduce la respuesta enunciada a la precisión con que está escrita. Es una reproducción aritmética; no valida el modelo físico ni sus hipótesis.

\noindent\textbf{Limitaciones}
\begin{itemize}
\item Modelo educativo; la aplicabilidad a una máquina real requiere revisión experta.
\end{itemize}

\noindent\textbf{Errores frecuentes}
\begin{itemize}
\item Invertir los radios y obtener caída de presión hacia fuera.
\end{itemize}
% studium:end-record eyJpZCI6ICJQUkItMDA4MSIsICJzZWN0aW9uIjogImMxMSJ9
\end{tcolorbox}

\begin{tcolorbox}[title={Problema resuelto 2}, studiumworked, breakable]

\noindent\textbf{Enunciado}
\(C11-02-R \cdot\) Trabajo uniforme con vórtice libre. Una fila móvil tiene constantes de vórtice libre distintas en entrada y salida. Calcule el trabajo y el incremento de temperatura total. Datos: \(\Omega=1080\, \mathrm{s}^{-1}\); \(K_e=11\, \mathrm{m}^{2}\,\mathrm{s}^{-1}\); \(K_s=39\, \mathrm{m}^{2}\,\mathrm{s}^{-1}\); \(c_p=1004.5\, \mathrm{J}\,\mathrm{kg}^{-1}\,\mathrm{K}^{-1}\).

\noindent\textbf{Datos}
\begin{itemize}
\item $\Omega = 1080\,\mathrm{s}^{-1}$ — velocidad angular
\item $K_e = 11\,\mathrm{m}^{2}\,\mathrm{s}^{-1}$ — constante entrada
\item $K_s = 39\,\mathrm{m}^{2}\,\mathrm{s}^{-1}$ — constante salida
\item $c_p = 1004.5\,\mathrm{J}\,\mathrm{kg}^{-1}\,\mathrm{K}^{-1}$ — calor específico
\end{itemize}

\noindent\textbf{Incógnita}
$\Delta T_0$, en $\mathrm{K}$; también las magnitudes intermedias desarrolladas.

\noindent\textbf{Modelo físico}

\noindent\textit{Trabajo uniforme con vórtice libre}

\[
\fitmath{\Delta T_0=\frac{\Omega \left(- K_e + K_s\right)}{c_p}}
\]

\noindent\textbf{Hipótesis}
\begin{itemize}
\item Misma velocidad angular y líneas de corriente del modelo ideal.
\end{itemize}

\noindent\textbf{Desarrollo paso a paso}

\begin{enumerate}

\item Euler cancela la dependencia radial r.

\[\fitmath{\mathrm{wc}=\Omega\cdot \left(\mathrm{ks}-\mathrm{ke}\right)}\]

\[\fitmath{\mathrm{wc}=(1080)\cdot \left((39)-(11)\right)\approx 30240\,\mathrm{J/kg}}\]

\item Se convierte el trabajo en incremento total.

\[\fitmath{\mathrm{dT}=\frac{\Omega\cdot \left(\mathrm{ks}-\mathrm{ke}\right)}{\mathrm{cp}}}\]

\[\fitmath{\mathrm{dT}=\frac{(1080)\cdot \left((39)-(11)\right)}{(1004.5)}\approx 30.10453\,\mathrm{K}}\]

\end{enumerate}


\noindent\textbf{Sustitución numérica}

\[
\fitmath{\frac{\left(1080\,\mathrm{s}^{-1}\right) \left(- \left(11\,\mathrm{m}^{2}\,\mathrm{s}^{-1}\right) + \left(39\,\mathrm{m}^{2}\,\mathrm{s}^{-1}\right)\right)}{\left(1004.5\,\mathrm{J}\,\mathrm{kg}^{-1}\,\mathrm{K}^{-1}\right)}}
\]

\noindent\textbf{Resultado}

\[
\fitmath{\Delta T_0=3.010452962\times10^{1}\,\mathrm{K}}
\]

\noindent\textbf{Interpretación}
El trabajo es uniforme con el radio, aunque la reacción no lo sea.

\noindent\textbf{Comprobación de los pasos}
\begin{itemize}
\item Paso 2: $30240$ — reproducido
\item Paso 4: $30.1045$ — reproducido
\end{itemize}

\noindent\textbf{Comprobación de la respuesta} la aritmética del último paso reproduce la respuesta enunciada a la precisión con que está escrita. Es una reproducción aritmética; no valida el modelo físico ni sus hipótesis.

\noindent\textbf{Limitaciones}
\begin{itemize}
\item Modelo educativo; la aplicabilidad a una máquina real requiere revisión experta.
\end{itemize}

\noindent\textbf{Errores frecuentes}
\begin{itemize}
\item Usar igual constante de vórtice a ambos lados de un rotor con trabajo.
\end{itemize}
% studium:end-record eyJpZCI6ICJQUkItMDA4MyIsICJzZWN0aW9uIjogImMxMSJ9
\end{tcolorbox}

\begin{tcolorbox}[title={Problema resuelto 3}, studiumworked, breakable]

\noindent\textbf{Enunciado}
\(C11-03-R \cdot\) Reacción hacia la raíz. Calcule la reacción reducida en la raíz para un escalón repetitivo de vórtice libre. Datos: \(K_e=16\, \mathrm{m}^{2}\,\mathrm{s}^{-1}\); \(K_s=49\, \mathrm{m}^{2}\,\mathrm{s}^{-1}\); \(\Omega=1120\, \mathrm{s}^{-1}\); \(r_h=0.205\, \mathrm{m}\).

\noindent\textbf{Datos}
\begin{itemize}
\item $K_e = 16\,\mathrm{m}^{2}\,\mathrm{s}^{-1}$ — constante tangencial entrada
\item $K_s = 49\,\mathrm{m}^{2}\,\mathrm{s}^{-1}$ — constante tangencial salida
\item $\Omega = 1120\,\mathrm{s}^{-1}$ — velocidad angular
\item $r_h = 0.205\,\mathrm{m}$ — radio raíz
\end{itemize}

\noindent\textbf{Incógnita}
$\mathcal{R}$ (adimensional); también las magnitudes intermedias desarrolladas.

\noindent\textbf{Modelo físico}

\noindent\textit{Reacción hacia la raíz}

\[
\fitmath{\mathcal R=1 - \frac{K_e + K_s}{2 \Omega r_h^{2}}}
\]

\noindent\textbf{Hipótesis}
\begin{itemize}
\item Criterio reducido válido para radio y Cx constantes en cada escalón repetitivo.
\end{itemize}

\noindent\textbf{Desarrollo paso a paso}

\begin{enumerate}

\item Se calcula la suma de componentes tangenciales en la raíz.

\[\fitmath{\mathrm{sum}=\frac{\mathrm{ke}+\mathrm{ks}}{\mathrm{rh}}}\]

\[\fitmath{\mathrm{sum}=\frac{(16)+(49)}{(0.205)}\approx 317.0732\,\mathrm{m/s}}\]

\item Se normaliza con dos veces U local.

\[\fitmath{\mathrm{rea}=1-\frac{\mathrm{ke}+\mathrm{ks}}{2\cdot \Omega\cdot {\mathrm{rh}}^{2}}}\]

\[\fitmath{\mathrm{rea}=1-\frac{(16)+(49)}{2\cdot (1120)\cdot {(0.205)}^{2}}\approx 0.3095096\,}\]

\end{enumerate}


\noindent\textbf{Sustitución numérica}

\[
\fitmath{1 - \frac{\left(16\,\mathrm{m}^{2}\,\mathrm{s}^{-1}\right) + \left(49\,\mathrm{m}^{2}\,\mathrm{s}^{-1}\right)}{2 \left(1120\,\mathrm{s}^{-1}\right) \left(0.205\,\mathrm{m}\right)^{2}}}
\]

\noindent\textbf{Resultado}

\[
\fitmath{\mathcal R=3.095096456\times10^{-1}}
\]

\noindent\textbf{Interpretación}
La raíz puede tener reacción mucho menor que el radio medio. Si resulta negativa se debe revisar el diseño, no cambiar el signo.

\noindent\textbf{Comprobación de los pasos}
\begin{itemize}
\item Paso 2: $317.073$ — reproducido
\item Paso 4: $0.30951$ — reproducido
\end{itemize}

\noindent\textbf{Comprobación de la respuesta} la aritmética del último paso reproduce la respuesta enunciada a la precisión con que está escrita. Es una reproducción aritmética; no valida el modelo físico ni sus hipótesis.

\noindent\textbf{Limitaciones}
\begin{itemize}
\item Modelo educativo; la aplicabilidad a una máquina real requiere revisión experta.
\end{itemize}

\noindent\textbf{Errores frecuentes}
\begin{itemize}
\item Asignar a toda la pala la reacción del radio medio.
\end{itemize}
% studium:end-record eyJpZCI6ICJQUkItMDA4NSIsICJzZWN0aW9uIjogImMxMSJ9
\end{tcolorbox}

\begin{tcolorbox}[title={Problema resuelto 4}, studiumworked, breakable]

\noindent\textbf{Enunciado}
\(C11-04-R \cdot\) Vórtice forzado admisible. Con entalpía total y entropía uniformes, calcule la velocidad axial exterior exigida por un remolino forzado. Datos: \(C_{x,m}=188\, \mathrm{m}\,\mathrm{s}^{-1}\); \(k_f=320\, \mathrm{s}^{-1}\); \(r_m=0.25\, \mathrm{m}\); \(r_t=0.33\, \mathrm{m}\).

\noindent\textbf{Datos}
\begin{itemize}
\item $C_{x,m} = 188\,\mathrm{m}\,\mathrm{s}^{-1}$ — velocidad axial al radio medio
\item $k_f = 320\,\mathrm{s}^{-1}$ — coeficiente $C_\theta=k_f r$
\item $r_m = 0.25\,\mathrm{m}$ — radio medio
\item $r_t = 0.33\,\mathrm{m}$ — radio exterior
\end{itemize}

\noindent\textbf{Incógnita}
$C_x$, en $\mathrm{m\,s^{-1}}$; también las magnitudes intermedias desarrolladas.

\noindent\textbf{Modelo físico}

\noindent\textit{Vórtice forzado admisible}

\[
\fitmath{C_x=\sqrt{C_{x,m}^{2} - 2 k_f^{2} \left(- r_m^{2} + r_t^{2}\right)}}
\]

\noindent\textbf{Hipótesis}
\begin{itemize}
\item Equilibrio radial simple, h0 y s uniformes.
\end{itemize}

\noindent\textbf{Desarrollo paso a paso}

\begin{enumerate}

\item El equilibrio energético fija el cuadrado axial exterior.

\[\fitmath{\mathrm{cx}_{2}={\mathrm{cxm}}^{2}-2\cdot {\mathrm{kv}}^{2}\cdot \left({\mathrm{rt}}^{2}-{\mathrm{rm}}^{2}\right)}\]

\[\fitmath{\mathrm{cx}_{2}={(188)}^{2}-2\cdot {(320)}^{2}\cdot \left({(0.33)}^{2}-{(0.25)}^{2}\right)\approx 25841.28\,\mathrm{m^{2}/s^{2}}}\]

\item Se selecciona la rama axial positiva.

\[\fitmath{\mathrm{cx}=\sqrt{{\mathrm{cxm}}^{2}-2\cdot {\mathrm{kv}}^{2}\cdot \left({\mathrm{rt}}^{2}-{\mathrm{rm}}^{2}\right)}}\]

\[\fitmath{\mathrm{cx}=\sqrt{{(188)}^{2}-2\cdot {(320)}^{2}\cdot \left({(0.33)}^{2}-{(0.25)}^{2}\right)}\approx 160.7522\,\mathrm{m/s}}\]

\end{enumerate}


\noindent\textbf{Sustitución numérica}

\[
\fitmath{\sqrt{\left(\left(188\,\mathrm{m}\,\mathrm{s}^{-1}\right)\right)^{2} - 2 \left(\left(320\,\mathrm{s}^{-1}\right)\right)^{2} \left(- \left(0.25\,\mathrm{m}\right)^{2} + \left(0.33\,\mathrm{m}\right)^{2}\right)}}
\]

\noindent\textbf{Resultado}

\[
\fitmath{C_x=1.607522317\times10^{2}\,\mathrm{m}\,\mathrm{s}^{-1}}
\]

\noindent\textbf{Interpretación}
El radicando debe ser positivo. Si no lo es, la ley propuesta no admite esta solución axial en todo el anillo.

\noindent\textbf{Comprobación de los pasos}
\begin{itemize}
\item Paso 2: $25841.3$ — reproducido
\item Paso 4: $160.752$ — reproducido
\end{itemize}

\noindent\textbf{Comprobación de la respuesta} la aritmética del último paso reproduce la respuesta enunciada a la precisión con que está escrita. Es una reproducción aritmética; no valida el modelo físico ni sus hipótesis.

\noindent\textbf{Limitaciones}
\begin{itemize}
\item Modelo educativo; la aplicabilidad a una máquina real requiere revisión experta.
\end{itemize}

\noindent\textbf{Errores frecuentes}
\begin{itemize}
\item Imponer Cx constante también para vórtice forzado.
\end{itemize}
% studium:end-record eyJpZCI6ICJQUkItMDA4NyIsICJzZWN0aW9uIjogImMxMSJ9
\end{tcolorbox}

\begin{tcolorbox}[title={Ejercicios propuestos}, studiumworked, breakable]

\noindent\textbf{1 (Básico, Desarrollo)}
\(C11-01-P \cdot\) Presión de vórtice libre. Ejercicio de transferencia. Integre el equilibrio radial de un fluido de densidad constante y remolino K/r entre radio medio y punta. Antes de consultar la solución, prediga el efecto de cambiar el primer dato respecto del caso resuelto de esta familia. Justifique la tendencia y explique qué restricción podría convertirse en limitante. Datos: \(\rho=2.4495\, \mathrm{kg}\,\mathrm{m}^{-3}\); \(K_v=36\, \mathrm{m}^{2}\,\mathrm{s}^{-1}\); \(r_m=0.29\, \mathrm{m}\); \(r_t=0.36\, \mathrm{m}\).

\noindent\textbf{2 (Intermedio, Desarrollo)}
\(C11-02-P \cdot\) Trabajo uniforme con vórtice libre. Ejercicio de transferencia. Una fila móvil tiene constantes de vórtice libre distintas en entrada y salida. Calcule el trabajo y el incremento de temperatura total. Antes de consultar la solución, prediga el efecto de cambiar el primer dato respecto del caso resuelto de esta familia. Justifique la tendencia y explique qué restricción podría convertirse en limitante. Datos: \(\Omega=1150.2\, \mathrm{s}^{-1}\); \(K_e=11\, \mathrm{m}^{2}\,\mathrm{s}^{-1}\); \(K_s=39\, \mathrm{m}^{2}\,\mathrm{s}^{-1}\); \(c_p=1004.5\, \mathrm{J}\,\mathrm{kg}^{-1}\,\mathrm{K}^{-1}\).

\noindent\textbf{3 (Avanzado, Desarrollo)}
\(C11-03-P \cdot\) Reacción hacia la raíz. Ejercicio de transferencia. Calcule la reacción reducida en la raíz para un escalón repetitivo de vórtice libre. Antes de consultar la solución, prediga el efecto de cambiar el primer dato respecto del caso resuelto de esta familia. Justifique la tendencia y explique qué restricción podría convertirse en limitante. Datos: \(K_e=17.04\, \mathrm{m}^{2}\,\mathrm{s}^{-1}\); \(K_s=49\, \mathrm{m}^{2}\,\mathrm{s}^{-1}\); \(\Omega=1120\, \mathrm{s}^{-1}\); \(r_h=0.205\, \mathrm{m}\).

\noindent\textbf{4 (Nivel de examen, Desarrollo)}
\(C11-04-P \cdot\) Vórtice forzado admisible. Ejercicio de transferencia. Con entalpía total y entropía uniformes, calcule la velocidad axial exterior exigida por un remolino forzado. Antes de consultar la solución, prediga el efecto de cambiar el primer dato respecto del caso resuelto de esta familia. Justifique la tendencia y explique qué restricción podría convertirse en limitante. Datos: \(C_{x,m}=200.22\, \mathrm{m}\,\mathrm{s}^{-1}\); \(k_f=320\, \mathrm{s}^{-1}\); \(r_m=0.25\, \mathrm{m}\); \(r_t=0.33\, \mathrm{m}\).
\end{tcolorbox}


\studiumchaptercheck

\chapter{Torsión, secciones y construcción geométrica de álabes}
\label{chap:c12}

\studiumchapterguide{Objetivo: desarrollar torsión, secciones y construcción geométrica de álabes, declarar hipótesis y comprobar condiciones físicas.}
% studium:end-record eyJpZCI6ICJQQVItMDA4MCIsICJzZWN0aW9uIjogImMxMiJ9

\section{Introducción y objetivos}

Una pala conecta secciones con direcciones de flujo distintas y debe transmitir cargas al disco. El objetivo es transformar leyes radiales de velocidad en ángulos de metal, generar una línea media y añadir espesor sin confundir esa construcción con un diseño aerodinámico validado. Pertenece al bloque IV y prepara la práctica de hoja de cálculo.
Se usará una familia polinómica educativa de perfiles, definida en este manual. No se reproducen tablas de coordenadas de las diapositivas ni se presenta la geometría obtenida como un perfil NACA. Su propósito es comprobar tangentes, continuidad y representación a varios radios.
% studium:end-record eyJpZCI6ICJQQVItMDA4MSIsICJzZWN0aW9uIjogImMxMiJ9

\section{De las velocidades a la torsión}

En cada radio se calcula \(U(r)=\Omega r\). Para una fila móvil, los ángulos del gas se obtienen de \(W_\theta=C_\theta-U\); para una fila fija, de \(C_\theta\). Con incidencia y desviación algebraicas prescritas:

\[
\fitmath{\kappa_e(r)=a_e(r)-i(r),\qquad \kappa_s(r)=a_s(r)-\delta(r).}
\]
Una variación de \(\kappa\) con el radio es torsión geométrica. La torsión de un rotor compresor no tiene por qué coincidir con la de un estator o con la de un rotor de turbina, aun si se usa la misma ley de vórtice libre. Cambian las estaciones, las referencias y el signo del trabajo.
Los ángulos del metal son una condición geométrica. En la práctica, el flujo no sigue exactamente esas tangentes y la desviación depende de carga y perfil. La elección inicial de \(i=\delta=0\) simplifica el ejemplo educativo y obliga a marcar la evaluación aerodinámica como pendiente.
% studium:end-record eyJpZCI6ICJQQVItMDA4MiIsICJzZWN0aW9uIjogImMxMiJ9

\section{Línea media con tangentes prescritas}

Se define \(\xi=x/c_x\), con \(0\leq\xi\leq1\), y \(y=c_xf(\xi)\). Para \(f(0)=0\), \(f(1)=d\), \(f'(0)=m_e=\tan\kappa_e\) y \(f'(1)=m_s=\tan\kappa_s\), un polinomio cúbico de Hermite es

\[
\fitmath{f(\xi)=(\xi^3-2\xi^2+\xi)m_e+(-2\xi^3+3\xi^2)d+(\xi^3-\xi^2)m_s.}
\]
Las cuatro condiciones determinan los cuatro coeficientes. El desplazamiento \(d\) fija la pendiente de la cuerda axial proyectada. Si se elige \(d=(m_e+m_s)/2\), el término cúbico se anula y queda una línea media cuadrática; no se convierte por ello en arco circular.
El calado de la cuerda es \(\chi=\arctan d\) y la cuerda real es \(c=c_x\sqrt{1+d^2}\). La figura~\ref{fig:geometria} muestra una construcción original y sus tangentes. Para ángulos próximos a noventa grados, una descripción de valor único \(y(x)\) puede ser inadecuada; se necesita una parametrización geométrica general.
% studium:end-record eyJpZCI6ICJQQVItMDA4MyIsICJzZWN0aW9uIjogImMxMiJ9

\par\vspace{\baselineskip}
\noindent\begin{minipage}{\linewidth}
\centering
\begingroup
\setbox0=\hbox{\begin{tikzpicture}
\draw[->,gray] (0,0)--(6.7,0) node[right] {$x/c_x$};\draw[->,gray] (0,-.3)--(0,2.7) node[above] {$y/c_x$};\draw[blue,very thick] (0,0) .. controls (1.8,1.5) and (3.8,2.1) .. (6,1.8);\draw[blue,very thick] (0,0) .. controls (1.8,.7) and (3.8,1.7) .. (6,1.8);\draw[dashed] (0,0) .. controls (1.8,1.1) and (3.8,1.9) .. (6,1.8);\draw[thin] (0,0)--(2,1.2);\draw[thin] (4.4,2.05)--(6.6,1.7);\node at (1.4,2.4) {$\kappa_e$};\node at (6,2.5) {$\kappa_s$};\draw[<->,red] (3,1.3)--(2.85,1.77) node[midway,left] {$2t_n$};
\end{tikzpicture}}%
\count0=\ht0
\advance\count0 by \dp0
\count2=\wd0
\count4=\dimexpr\textheight/100*38\relax
\count6=\linewidth
\ifnum\count0<1
\count0=1
\fi
\ifnum\count2<1
\count2=1
\fi
\divide\count0 by 4096
\divide\count2 by 4096
\divide\count4 by 4096
\divide\count6 by 4096
\ifnum\count0<1
\count0=1
\fi
\ifnum\count2<1
\count2=1
\fi
\ifnum\count4<1
\count4=1
\fi
\ifnum\count6<1
\count6=1
\fi
\count8=\count2
\ifnum\count0>\count4
\count8=\numexpr\count2*\count4/\count0\relax
\fi
\ifnum\count8>\count6
\resizebox{\linewidth}{!}{\begin{tikzpicture}
\draw[->,gray] (0,0)--(6.7,0) node[right] {$x/c_x$};\draw[->,gray] (0,-.3)--(0,2.7) node[above] {$y/c_x$};\draw[blue,very thick] (0,0) .. controls (1.8,1.5) and (3.8,2.1) .. (6,1.8);\draw[blue,very thick] (0,0) .. controls (1.8,.7) and (3.8,1.7) .. (6,1.8);\draw[dashed] (0,0) .. controls (1.8,1.1) and (3.8,1.9) .. (6,1.8);\draw[thin] (0,0)--(2,1.2);\draw[thin] (4.4,2.05)--(6.6,1.7);\node at (1.4,2.4) {$\kappa_e$};\node at (6,2.5) {$\kappa_s$};\draw[<->,red] (3,1.3)--(2.85,1.77) node[midway,left] {$2t_n$};
\end{tikzpicture}}%
\else
\ifnum\count0>\count4
\resizebox{!}{\dimexpr\textheight/100*38\relax}{\begin{tikzpicture}
\draw[->,gray] (0,0)--(6.7,0) node[right] {$x/c_x$};\draw[->,gray] (0,-.3)--(0,2.7) node[above] {$y/c_x$};\draw[blue,very thick] (0,0) .. controls (1.8,1.5) and (3.8,2.1) .. (6,1.8);\draw[blue,very thick] (0,0) .. controls (1.8,.7) and (3.8,1.7) .. (6,1.8);\draw[dashed] (0,0) .. controls (1.8,1.1) and (3.8,1.9) .. (6,1.8);\draw[thin] (0,0)--(2,1.2);\draw[thin] (4.4,2.05)--(6.6,1.7);\node at (1.4,2.4) {$\kappa_e$};\node at (6,2.5) {$\kappa_s$};\draw[<->,red] (3,1.3)--(2.85,1.77) node[midway,left] {$2t_n$};
\end{tikzpicture}}%
\else
\begin{tikzpicture}
\draw[->,gray] (0,0)--(6.7,0) node[right] {$x/c_x$};\draw[->,gray] (0,-.3)--(0,2.7) node[above] {$y/c_x$};\draw[blue,very thick] (0,0) .. controls (1.8,1.5) and (3.8,2.1) .. (6,1.8);\draw[blue,very thick] (0,0) .. controls (1.8,.7) and (3.8,1.7) .. (6,1.8);\draw[dashed] (0,0) .. controls (1.8,1.1) and (3.8,1.9) .. (6,1.8);\draw[thin] (0,0)--(2,1.2);\draw[thin] (4.4,2.05)--(6.6,1.7);\node at (1.4,2.4) {$\kappa_e$};\node at (6,2.5) {$\kappa_s$};\draw[<->,red] (3,1.3)--(2.85,1.77) node[midway,left] {$2t_n$};
\end{tikzpicture}%
\fi
\fi
\endgroup
\par\nopagebreak
\studiumcaption{fig:geometria}{Perfil educativo original; tangentes y espesor normal. No es un perfil comercial ni una predicción de rendimiento.}
\end{minipage}
\par\vspace{\baselineskip}
% studium:end-record eyJpZCI6ICJGSUctMDAxMiIsICJzZWN0aW9uIjogImMxMiJ9

\section{Espesor normal y coordenadas}


Se prescribe un semiespesor educativo



\[
\fitmath{t_n(\xi)=\frac{\varepsilon c_x}{2}\sin(\pi\xi).}
\]
La normal unitaria a la línea media es \(\mathbf n=(-f',1)/\sqrt{1+f'^2}\). Las superficies se obtienen de

\[
\fitmath{(x_\pm,y_\pm)=(x,y)\pm t_n\mathbf n.}
\]
El espesor máximo normal de esta ley es \(\varepsilon c_x\). Los extremos tienen espesor nulo y pueden formar cúspides; no representan radios de borde de ataque o salida adecuados para fabricación. Precisamente por ello se identifica como geometría de aprendizaje.
El desplazamiento normal modifica también la coordenada axial. Sumar espesor solo en la dirección \(y\) produce un espesor vertical, distinto del normal cuando la pendiente es grande. Deben comprobarse autocruces y distancia entre perfiles vecinos antes de estimar una garganta.
% studium:end-record eyJpZCI6ICJQQVItMDA4NCIsICJzZWN0aW9uIjogImMxMiJ9

\section{Apilamiento y representación tridimensional}

Se elige una línea de apilamiento, por ejemplo un punto común de referencia en cada sección. Unir bordes de ataque, centroides o puntos de cuerda produce geometrías distintas. El barrido axial y el desplazamiento tangencial de secciones pueden influir en carga y tensiones; no se añaden únicamente para mejorar el aspecto del dibujo.
Para una fila de \(Z_b\) álabes, el paso cambia como \(r\). Al representar varias secciones en un plano desarrollado se deben declarar la escala y el radio. Una vista superpuesta facilita comparar torsión, pero no permite medir directamente una holgura tridimensional. Una representación de sección no sustituye una superficie cerrada ni un análisis de manufacturabilidad.
% studium:end-record eyJpZCI6ICJQQVItMDA4NSIsICJzZWN0aW9uIjogImMxMiJ9

\section{Validación geométrica y límites}

Se comprueba que las coordenadas iniciales y finales coincidan con las condiciones, que las derivadas reproduzcan los ángulos prescritos y que las superficies no se crucen. Se calcula la solidez real después de redondear el número de álabes y se revisa la sensibilidad a la malla de puntos. La convergencia de una curva geométrica no valida su rendimiento aerodinámico.
Errores habituales: confundir calado con ángulo de salida, usar grados en funciones que esperan radianes, aplicar la ley radial solo a la velocidad periférica y conservar una componente tangencial incompatible con el trabajo. Autocomprobación: cambie el punto de apilamiento manteniendo cada sección; explique qué propiedades locales se conservan y qué geometría global cambia.
Formulario: \(\kappa=a-i\) en entrada, \(\kappa=a-\delta\) en salida, cúbica de Hermite y desplazamiento normal. Referencias: ULE Diseño de compresores, pp. \(28-35\); Esfuerzos y refrigeración, pp. \(2-7\); MIT 16.50, lección 26, contexto de torsión. La familia analítica de este capítulo es original y no tiene validación de perfil comercial.
% studium:end-record eyJpZCI6ICJQQVItMDA4NiIsICJzZWN0aW9uIjogImMxMiJ9

\begin{tcolorbox}[title={Desarrollos del modelo}, studiumdef, breakable]

\noindent\textbf{Modelo desarrollado: Torsión de entrada de rotor}

\begin{equation}
\Delta\beta=\arctan\frac{\Omega\,r_h}{C_x}-\arctan\frac{\Omega\,r_t}{C_x}
\label{eq:EC-C12}
\end{equation}

\noindent\textit{Significado.} Diferencia de ángulo relativo entre raíz y punta con prerrotación nula. El valor es negativo cuando la punta está más inclinada; no equivale al valor absoluto de la torsión geométrica.

\noindent\textbf{Hipótesis}
\begin{itemize}
\item Sin prerrotación, Cx uniforme, incidencia cero.
\end{itemize}

\noindent\textbf{Principios}
\begin{itemize}
\item Balance y definiciones físicas desarrollados en el capítulo; propiedades y dominio declarados en el problema.
\end{itemize}

\noindent\textbf{Aplicabilidad}
\begin{itemize}
\item Modelo educativo del caso C12-01-R
\end{itemize}

\noindent\textbf{Limitaciones}
\begin{itemize}
\item Modelo educativo; la aplicabilidad a una máquina real requiere revisión experta.
\end{itemize}

\noindent\textbf{Pasos}
\begin{enumerate}
\item
Se obtiene beta de raíz con componente tangencial -U.

\[
\fitmath{\beta_h=- \operatorname{atan}{\left(\frac{\Omega r_h}{C_x} \right)}}
\]
\item
Se obtiene beta de punta.

\[
\fitmath{\beta_t=- \operatorname{atan}{\left(\frac{\Omega r_t}{C_x} \right)}}
\]
\item
La torsión es una diferencia algebraica de orientaciones.

\[
\fitmath{\Delta\beta=\operatorname{atan}{\left(\frac{\Omega r_h}{C_x} \right)} - \operatorname{atan}{\left(\frac{\Omega r_t}{C_x} \right)}}
\]
\end{enumerate}

{\noindent\footnotesize verificada dimensionalmente.\par}
\end{tcolorbox}

\begin{tcolorbox}[title={Problema resuelto 1}, studiumworked, breakable]

\noindent\textbf{Enunciado}
\(C12-01-R \cdot\) Torsión de entrada de rotor. Un rotor sin prerrotación recibe velocidad axial uniforme. Calcule beta a raíz y punta y la diferencia de torsión de entrada. Datos: \(\Omega=1060\, \mathrm{s}^{-1}\); \(C_x=157\, \mathrm{m}\,\mathrm{s}^{-1}\); \(r_h=0.198\, \mathrm{m}\); \(r_t=0.314\, \mathrm{m}\).

\noindent\textbf{Datos}
\begin{itemize}
\item $\Omega = 1060\,\mathrm{s}^{-1}$ — velocidad angular
\item $C_x = 157\,\mathrm{m}\,\mathrm{s}^{-1}$ — velocidad axial
\item $r_h = 0.198\,\mathrm{m}$ — radio raíz
\item $r_t = 0.314\,\mathrm{m}$ — radio punta
\end{itemize}

\noindent\textbf{Incógnita}
$\Delta\beta$, en $\mathrm{rad}$; también las magnitudes intermedias desarrolladas.

\noindent\textbf{Modelo físico}

\noindent\textit{Torsión de entrada de rotor}

\[
\fitmath{\Delta\beta=\operatorname{atan}{\left(\frac{\Omega r_h}{C_x} \right)} - \operatorname{atan}{\left(\frac{\Omega r_t}{C_x} \right)}}
\]

\noindent\textbf{Hipótesis}
\begin{itemize}
\item Sin prerrotación, Cx uniforme, incidencia cero.
\end{itemize}

\noindent\textbf{Desarrollo paso a paso}

\begin{enumerate}

\item Se obtiene beta de raíz con componente tangencial -U.

\[\fitmath{\mathrm{bh}=\arctan\left(\frac{-\Omega\cdot \mathrm{rh}}{\mathrm{cx}}\right)}\]

\[\fitmath{\mathrm{bh}=\arctan\left(\frac{-(1060)\cdot (0.198)}{(157)}\right)\approx -0.9285466\,\mathrm{rad}}\]

\item Se obtiene beta de punta.

\[\fitmath{\mathrm{bt}=\arctan\left(\frac{-\Omega\cdot \mathrm{rt}}{\mathrm{cx}}\right)}\]

\[\fitmath{\mathrm{bt}=\arctan\left(\frac{-(1060)\cdot (0.314)}{(157)}\right)\approx -1.130045\,\mathrm{rad}}\]

\item La torsión es una diferencia algebraica de orientaciones.

\[\fitmath{\mathrm{twist}=\arctan\left(\frac{-\Omega\cdot \mathrm{rt}}{\mathrm{cx}}\right)-\arctan\left(\frac{-\Omega\cdot \mathrm{rh}}{\mathrm{cx}}\right)}\]

\[\fitmath{\mathrm{twist}=\arctan\left(\frac{-(1060)\cdot (0.314)}{(157)}\right)-\arctan\left(\frac{-(1060)\cdot (0.198)}{(157)}\right)\approx -0.2014989\,\mathrm{rad}}\]

\end{enumerate}


\noindent\textbf{Sustitución numérica}

\[
\fitmath{\operatorname{atan}{\left(\frac{\left(1060\,\mathrm{s}^{-1}\right) \left(0.198\,\mathrm{m}\right)}{\left(157\,\mathrm{m}\,\mathrm{s}^{-1}\right)} \right)} - \operatorname{atan}{\left(\frac{\left(1060\,\mathrm{s}^{-1}\right) \left(0.314\,\mathrm{m}\right)}{\left(157\,\mathrm{m}\,\mathrm{s}^{-1}\right)} \right)}}
\]

\noindent\textbf{Resultado}

\[
\fitmath{\Delta\beta=-2.014988500\times10^{-1}\,\mathrm{rad}}
\]

\noindent\textbf{Interpretación}
La orientación de punta es más negativa. La tangente del metal coincide solo bajo la hipótesis de incidencia prescrita igual a cero.

\noindent\textbf{Comprobación de los pasos}
\begin{itemize}
\item Paso 2: $-0.928547$ — reproducido
\item Paso 4: $-1.13005$ — reproducido
\item Paso 6: $-0.201499$ — reproducido
\end{itemize}

\noindent\textbf{Comprobación de la respuesta} la aritmética del último paso reproduce la respuesta enunciada a la precisión con que está escrita. Es una reproducción aritmética; no valida el modelo físico ni sus hipótesis.

\noindent\textbf{Limitaciones}
\begin{itemize}
\item Modelo educativo; la aplicabilidad a una máquina real requiere revisión experta.
\end{itemize}

\noindent\textbf{Errores frecuentes}
\begin{itemize}
\item Atribuir el mismo ángulo a todos los radios.
\end{itemize}
% studium:end-record eyJpZCI6ICJQUkItMDA4OSIsICJzZWN0aW9uIjogImMxMiJ9
\end{tcolorbox}

\begin{tcolorbox}[title={Problema resuelto 2}, studiumworked, breakable]

\noindent\textbf{Enunciado}
\(C12-02-R \cdot\) Cuerda y proyección. Una sección tiene cuerda axial y calado dados. Calcule cuerda real y solidez a paso fijado. Datos: \(c_x=0.044\, \mathrm{m}\); \(\chi=0.6457718\, \mathrm{rad}\); \(s=0.039\, \mathrm{m}\).

\noindent\textbf{Datos}
\begin{itemize}
\item $c_x = 0.044\,\mathrm{m}$ — cuerda axial
\item $\chi = 0.6457718232379\,\mathrm{rad}$ — calado
\item $s = 0.039\,\mathrm{m}$ — paso
\end{itemize}

\noindent\textbf{Incógnita}
$\sigma$ (adimensional); también las magnitudes intermedias desarrolladas.

\noindent\textbf{Modelo físico}

\noindent\textit{Cuerda y proyección}

\[
\fitmath{\sigma=\frac{c_x}{s \cos{\left(\chi \right)}}}
\]

\noindent\textbf{Hipótesis}
\begin{itemize}
\item Calado menor de noventa grados en módulo.
\end{itemize}

\noindent\textbf{Desarrollo paso a paso}

\begin{enumerate}

\item La cuerda axial es la proyección de la cuerda real.

\[\fitmath{\mathrm{chord}=\frac{\mathrm{cxg}}{\cos\left(\chi\right)}}\]

\[\fitmath{\mathrm{chord}=\frac{(0.044)}{\cos\left((0.6457718232)\right)}\approx 0.05509397\,\mathrm{m}}\]

\item La solidez usa cuerda real.

\[\fitmath{\mathrm{sol}=\frac{\mathrm{cxg}}{\cos\left(\chi\right)\cdot \mathrm{ss}}}\]

\[\fitmath{\mathrm{sol}=\frac{(0.044)}{\cos\left((0.6457718232)\right)\cdot (0.039)}\approx 1.412666\,}\]

\end{enumerate}


\noindent\textbf{Sustitución numérica}

\[
\fitmath{\frac{\left(0.044\,\mathrm{m}\right)}{\left(0.039\,\mathrm{m}\right) \cos{\left(\left(0.6457718232379\,\mathrm{rad}\right) \right)}}}
\]

\noindent\textbf{Resultado}

\[
\fitmath{\sigma=1.412665871\times10^{0}}
\]

\noindent\textbf{Interpretación}
Cuerda axial y cuerda real no pueden intercambiarse en la solidez sin cambiar su definición.

\noindent\textbf{Comprobación de los pasos}
\begin{itemize}
\item Paso 2: $0.055094$ — reproducido
\item Paso 4: $1.41267$ — reproducido
\end{itemize}

\noindent\textbf{Comprobación de la respuesta} la aritmética del último paso reproduce la respuesta enunciada a la precisión con que está escrita. Es una reproducción aritmética; no valida el modelo físico ni sus hipótesis.

\noindent\textbf{Limitaciones}
\begin{itemize}
\item Modelo educativo; la aplicabilidad a una máquina real requiere revisión experta.
\end{itemize}

\noindent\textbf{Errores frecuentes}
\begin{itemize}
\item Multiplicar por coseno al despejar la cuerda real.
\end{itemize}
% studium:end-record eyJpZCI6ICJQUkItMDA5MSIsICJzZWN0aW9uIjogImMxMiJ9
\end{tcolorbox}

\begin{tcolorbox}[title={Problema resuelto 3}, studiumworked, breakable]

\noindent\textbf{Enunciado}
\(C12-03-R \cdot\) Línea de Hermite. Para pendientes de entrada y salida dadas y desplazamiento d, calcule la ordenada de línea media al parámetro xi. Datos: \(c_x=0.047\, \mathrm{m}\); \(m_e=0.86\); \(m_s=-0.24\); \(\delta_y=0.31\); \(\xi=0.4\).

\noindent\textbf{Datos}
\begin{itemize}
\item $c_x = 0.047\,\mathrm{m}$ — cuerda axial
\item $m_e = 0.86$ — pendiente entrada
\item $m_s = -0.24$ — pendiente salida
\item $\delta_y = 0.31$ — desplazamiento final normalizado
\item $\xi = 0.4$ — parámetro xi
\end{itemize}

\noindent\textbf{Incógnita}
$y$, en $\mathrm{m}$; también las magnitudes intermedias desarrolladas.

\noindent\textbf{Modelo físico}

\noindent\textit{Línea de Hermite}

\[
\fitmath{y=c_x \left(\delta_y \left(- 2 \xi^{3} + 3 \xi^{2}\right) + m_e \left(\xi^{3} - 2 \xi^{2} + \xi\right) + m_s \left(\xi^{3} - \xi^{2}\right)\right)}
\]

\noindent\textbf{Hipótesis}
\begin{itemize}
\item Perfil educativo definido como función de x.
\end{itemize}

\noindent\textbf{Desarrollo paso a paso}

\begin{enumerate}

\item Se evalúa la línea media normalizada.

\[\fitmath{\mathrm{fy}=\left({\mathrm{xx}}^{3}-2\cdot {\mathrm{xx}}^{2}+\mathrm{xx}\right)\cdot \mathrm{me}+\left(-2\cdot {\mathrm{xx}}^{3}+3\cdot {\mathrm{xx}}^{2}\right)\cdot \mathrm{dd}+\left({\mathrm{xx}}^{3}-{\mathrm{xx}}^{2}\right)\cdot \mathrm{ms}}\]

\[\fitmath{\mathrm{fy}=\left({(0.4)}^{3}-2\cdot {(0.4)}^{2}+(0.4)\right)\cdot (0.86)+\left(-2\cdot {(0.4)}^{3}+3\cdot {(0.4)}^{2}\right)\cdot (0.31)+\left({(0.4)}^{3}-{(0.4)}^{2}\right)\cdot (-0.24)\approx 0.256\,}\]

\item Se devuelve la longitud física transversal.

\[\fitmath{\mathrm{yy}=\mathrm{cxg}\cdot \left(\left({\mathrm{xx}}^{3}-2\cdot {\mathrm{xx}}^{2}+\mathrm{xx}\right)\cdot \mathrm{me}+\left(-2\cdot {\mathrm{xx}}^{3}+3\cdot {\mathrm{xx}}^{2}\right)\cdot \mathrm{dd}+\left({\mathrm{xx}}^{3}-{\mathrm{xx}}^{2}\right)\cdot \mathrm{ms}\right)}\]

\[\fitmath{\mathrm{yy}=(0.047)\cdot \left(\left({(0.4)}^{3}-2\cdot {(0.4)}^{2}+(0.4)\right)\cdot (0.86)+\left(-2\cdot {(0.4)}^{3}+3\cdot {(0.4)}^{2}\right)\cdot (0.31)+\left({(0.4)}^{3}-{(0.4)}^{2}\right)\cdot (-0.24)\right)\approx 0.012032\,\mathrm{m}}\]

\end{enumerate}


\noindent\textbf{Sustitución numérica}

\[
\fitmath{\left(0.047\,\mathrm{m}\right) \left(\left(0.31\right) \left(- 2 \left(0.4\right)^{3} + 3 \left(0.4\right)^{2}\right) + \left(0.86\right) \left(\left(0.4\right)^{3} - 2 \left(0.4\right)^{2} + \left(0.4\right)\right) + \left(-0.24\right) \left(\left(0.4\right)^{3} - \left(0.4\right)^{2}\right)\right)}
\]

\noindent\textbf{Resultado}

\[
\fitmath{y=1.203200000\times10^{-2}\,\mathrm{m}}
\]

\noindent\textbf{Interpretación}
El resultado es una coordenada geométrica, no una predicción aerodinámica. Compruebe los cuatro extremos y pendientes del polinomio.

\noindent\textbf{Comprobación de los pasos}
\begin{itemize}
\item Paso 2: $0.256$ — reproducido
\item Paso 4: $0.012032$ — reproducido
\end{itemize}

\noindent\textbf{Comprobación de la respuesta} la aritmética del último paso reproduce la respuesta enunciada a la precisión con que está escrita. Es una reproducción aritmética; no valida el modelo físico ni sus hipótesis.

\noindent\textbf{Limitaciones}
\begin{itemize}
\item Modelo educativo; la aplicabilidad a una máquina real requiere revisión experta.
\end{itemize}

\noindent\textbf{Errores frecuentes}
\begin{itemize}
\item Confundir pendiente con ángulo en grados.
\end{itemize}
% studium:end-record eyJpZCI6ICJQUkItMDA5MyIsICJzZWN0aW9uIjogImMxMiJ9
\end{tcolorbox}

\begin{tcolorbox}[title={Problema resuelto 4}, studiumworked, breakable]

\noindent\textbf{Enunciado}
\(C12-04-R \cdot\) Desplazamiento normal por espesor. En un punto de una línea media se conocen pendiente y semiespesor normal. Halle las dos componentes del desplazamiento normal y su componente axial como resultado final. Datos: \(t_n=0.0024\, \mathrm{m}\); \(m_y=0.73\).

\noindent\textbf{Datos}
\begin{itemize}
\item $t_n = 0.0024\,\mathrm{m}$ — semiespesor normal local
\item $m_y = 0.73$ — pendiente local
\end{itemize}

\noindent\textbf{Incógnita}
$\Delta x$, en $\mathrm{m}$; también las magnitudes intermedias desarrolladas.

\noindent\textbf{Modelo físico}

\noindent\textit{Desplazamiento normal por espesor}

\[
\fitmath{\Delta x=- \frac{m_y t_n}{\sqrt{m_y^{2} + 1}}}
\]

\noindent\textbf{Hipótesis}
\begin{itemize}
\item Superficie obtenida por desplazamiento normal local.
\end{itemize}

\noindent\textbf{Desarrollo paso a paso}

\begin{enumerate}

\item La normal unitaria tiene componente transversal positiva.

\[\fitmath{\mathrm{dy}=\frac{\mathrm{tt}}{\sqrt{1+{\mathrm{slope}}^{2}}}}\]

\[\fitmath{\mathrm{dy}=\frac{(0.0024)}{\sqrt{1+{(0.73)}^{2}}}\approx 0.001938449\,\mathrm{m}}\]

\item La componente axial compensa la inclinación local.

\[\fitmath{\mathrm{dx}=\frac{-\mathrm{tt}\cdot \mathrm{slope}}{\sqrt{1+{\mathrm{slope}}^{2}}}}\]

\[\fitmath{\mathrm{dx}=\frac{-(0.0024)\cdot (0.73)}{\sqrt{1+{(0.73)}^{2}}}\approx -0.001415068\,\mathrm{m}}\]

\end{enumerate}


\noindent\textbf{Sustitución numérica}

\[
\fitmath{- \frac{\left(0.73\right) \left(0.0024\,\mathrm{m}\right)}{\sqrt{\left(0.73\right)^{2} + 1}}}
\]

\noindent\textbf{Resultado}

\[
\fitmath{\Delta x=-1.415067608\times10^{-3}\,\mathrm{m}}
\]

\noindent\textbf{Interpretación}
La longitud del vector desplazamiento es tt. Añadir solo dy produce un espesor vertical distinto.

\noindent\textbf{Comprobación de los pasos}
\begin{itemize}
\item Paso 2: $0.00193845$ — reproducido
\item Paso 4: $-0.00141507$ — reproducido
\end{itemize}

\noindent\textbf{Comprobación de la respuesta} la aritmética del último paso reproduce la respuesta enunciada a la precisión con que está escrita. Es una reproducción aritmética; no valida el modelo físico ni sus hipótesis.

\noindent\textbf{Limitaciones}
\begin{itemize}
\item Modelo educativo; la aplicabilidad a una máquina real requiere revisión experta.
\end{itemize}

\noindent\textbf{Errores frecuentes}
\begin{itemize}
\item Olvidar el desplazamiento axial al añadir espesor.
\end{itemize}
% studium:end-record eyJpZCI6ICJQUkItMDA5NSIsICJzZWN0aW9uIjogImMxMiJ9
\end{tcolorbox}

\begin{tcolorbox}[title={Ejercicios propuestos}, studiumworked, breakable]

\noindent\textbf{1 (Básico, Desarrollo)}
\(C12-01-P \cdot\) Torsión de entrada de rotor. Ejercicio de transferencia. Un rotor sin prerrotación recibe velocidad axial uniforme. Calcule beta a raíz y punta y la diferencia de torsión de entrada. Antes de consultar la solución, prediga el efecto de cambiar el primer dato respecto del caso resuelto de esta familia. Justifique la tendencia y explique qué restricción podría convertirse en limitante. Datos: \(\Omega=1128.9\, \mathrm{s}^{-1}\); \(C_x=157\, \mathrm{m}\,\mathrm{s}^{-1}\); \(r_h=0.198\, \mathrm{m}\); \(r_t=0.314\, \mathrm{m}\).

\noindent\textbf{2 (Intermedio, Desarrollo)}
\(C12-02-P \cdot\) Cuerda y proyección. Ejercicio de transferencia. Una sección tiene cuerda axial y calado dados. Calcule cuerda real y solidez a paso fijado. Antes de consultar la solución, prediga el efecto de cambiar el primer dato respecto del caso resuelto de esta familia. Justifique la tendencia y explique qué restricción podría convertirse en limitante. Datos: \(c_x=0.04686\, \mathrm{m}\); \(\chi=0.6457718\, \mathrm{rad}\); \(s=0.039\, \mathrm{m}\).

\noindent\textbf{3 (Avanzado, Desarrollo)}
\(C12-03-P \cdot\) Línea de Hermite. Ejercicio de transferencia. Para pendientes de entrada y salida dadas y desplazamiento d, calcule la ordenada de línea media al parámetro xi. Antes de consultar la solución, prediga el efecto de cambiar el primer dato respecto del caso resuelto de esta familia. Justifique la tendencia y explique qué restricción podría convertirse en limitante. Datos: \(c_x=0.050055\, \mathrm{m}\); \(m_e=0.86\); \(m_s=-0.24\); \(\delta_y=0.31\); \(\xi=0.4\).

\noindent\textbf{4 (Nivel de examen, Desarrollo)}
\(C12-04-P \cdot\) Desplazamiento normal por espesor. Ejercicio de transferencia. En un punto de una línea media se conocen pendiente y semiespesor normal. Halle las dos componentes del desplazamiento normal y su componente axial como resultado final. Antes de consultar la solución, prediga el efecto de cambiar el primer dato respecto del caso resuelto de esta familia. Justifique la tendencia y explique qué restricción podría convertirse en limitante. Datos: \(t_n=0.002556\, \mathrm{m}\); \(m_y=0.73\).
\end{tcolorbox}


\studiumchaptercheck

\chapter{Esfuerzos centrífugos, flexión y vibración}
\label{chap:c13}

\studiumchapterguide{Objetivo: desarrollar esfuerzos centrífugos, flexión y vibración, declarar hipótesis y comprobar condiciones físicas.}
% studium:end-record eyJpZCI6ICJQQVItMDA4NyIsICJzZWN0aW9uIjogImMxMyJ9

\section{Introducción y objetivos}

Las velocidades que mejoran el trabajo específico aumentan también las solicitaciones mecánicas. Una pala debe soportar su masa en rotación, fuerzas aerodinámicas y ciclos térmicos durante una vida determinada. El objetivo es deducir la tracción centrífuga, estudiar el efecto del estrechamiento y distinguir una comprobación elástica elemental de una evaluación de vida. Pertenece al bloque IV; el cálculo modal elemental se identifica como ampliación de apoyo.
No se deducen tensiones admisibles de gráficas históricas sin conocer material, temperatura y duración del ensayo. Los valores admisibles que aparecen en ejercicios son datos ficticios de diseño educativo y no recomendaciones para piezas reales.
% studium:end-record eyJpZCI6ICJQQVItMDA4OCIsICJzZWN0aW9uIjogImMxMyJ9

\section{Tracción centrífuga de una pala}

Considérese una sección a radio \(r\), de área resistente \(A_b(r)\), densidad material \(\rho_b\), que sostiene la parte exterior hasta \(r_t\). Un elemento de longitud radial \(d\xi\) exige fuerza centrípeta \(\rho_bA_b(\xi)\Omega^2\xi\,d\xi\). La tracción interior es

\[
\fitmath{F(r)=\rho_b\Omega^2\int_r^{r_t}A_b(\xi)\xi\,d\xi,\qquad
\sigma_c(r)=\frac{F(r)}{A_b(r)}.}
\]
Para área constante,

\[
\fitmath{\sigma_c(r)=\frac{\rho_b\Omega^2}{2}(r_t^2-r^2),\qquad
\sigma_{c,h}=\frac{\rho_b}{2}(U_t^2-U_h^2).}
\]
La raíz soporta toda la masa exterior y presenta la mayor tensión de este modelo uniforme. El área desaparece de la expresión de tensión porque masa y resistencia crecen en la misma proporción. Esto no significa que el área sea irrelevante para flexión, frecuencias, anclaje o refrigeración.
% studium:end-record eyJpZCI6ICJQQVItMDA4OSIsICJzZWN0aW9uIjogImMxMyJ9

\section{Estrechamiento y parámetro de anillo}

Si \(A_b(r)=A_h[1-k(r-r_h)/b]\), con \(0\leq k<1\), disminuye la masa exterior. La tensión de raíz se obtiene integrando esa ley, no multiplicando por una reducción arbitraria. La figura~\ref{fig:esfuerzos} compara soluciones calculadas para pala uniforme y estrechada. El máximo de una pala de sección variable debe buscarse en todo el radio; no se presupone su posición sin evaluar \(F/A_b\).
Para pala uniforme y área anular \(A_{an}=\pi(r_t^2-r_h^2)\),

\[
\fitmath{\sigma_{c,h}=\frac{\rho_b\Omega^2A_{an}}{2\pi}.}
\]
Sustituir \(\Omega=2\pi N/60\) permite relacionar tensión con \(A_{an}N^2\) si \(N\) se expresa en revoluciones por minuto. Cambiar a revoluciones por segundo modifica el factor numérico; una tabla de límites sin unidades es inutilizable.
% studium:end-record eyJpZCI6ICJQQVItMDA5MCIsICJzZWN0aW9uIjogImMxMyJ9

\par\vspace{\baselineskip}
\noindent\begin{minipage}{\linewidth}
\centering
\begingroup
\setbox0=\hbox{\begin{tikzpicture}
\draw[->,gray] (0,0)--(7,0);\draw[->,gray] (0,0)--(0,3.8);\draw (0.0,0)--(0.0,-.08) node[below,font=\scriptsize] {0.7};\draw (0,0.0)--(-.08,0.0) node[left,font=\scriptsize] {0};\draw (1.5,0)--(1.5,-.08) node[below,font=\scriptsize] {0.775};\draw (0,0.75)--(-.08,0.75) node[left,font=\scriptsize] {0.0625};\draw (3.0,0)--(3.0,-.08) node[below,font=\scriptsize] {0.85};\draw (0,1.5)--(-.08,1.5) node[left,font=\scriptsize] {0.125};\draw (4.5,0)--(4.5,-.08) node[below,font=\scriptsize] {0.925};\draw (0,2.25)--(-.08,2.25) node[left,font=\scriptsize] {0.188};\draw (6.0,0)--(6.0,-.08) node[below,font=\scriptsize] {1};\draw (0,3.0)--(-.08,3.0) node[left,font=\scriptsize] {0.25};\draw[blue,very thick] plot coordinates {(0.00000,3.06000) (0.10000,3.01785) (0.20000,2.97540) (0.30000,2.93265) (0.40000,2.88960) (0.50000,2.84625) (0.60000,2.80260) (0.70000,2.75865) (0.80000,2.71440) (0.90000,2.66985) (1.00000,2.62500) (1.10000,2.57985) (1.20000,2.53440) (1.30000,2.48865) (1.40000,2.44260) (1.50000,2.39625) (1.60000,2.34960) (1.70000,2.30265) (1.80000,2.25540) (1.90000,2.20785) (2.00000,2.16000) (2.10000,2.11185) (2.20000,2.06340) (2.30000,2.01465) (2.40000,1.96560) (2.50000,1.91625) (2.60000,1.86660) (2.70000,1.81665) (2.80000,1.76640) (2.90000,1.71585) (3.00000,1.66500) (3.10000,1.61385) (3.20000,1.56240) (3.30000,1.51065) (3.40000,1.45860) (3.50000,1.40625) (3.60000,1.35360) (3.70000,1.30065) (3.80000,1.24740) (3.90000,1.19385) (4.00000,1.14000) (4.10000,1.08585) (4.20000,1.03140) (4.30000,0.97665) (4.40000,0.92160) (4.50000,0.86625) (4.60000,0.81060) (4.70000,0.75465) (4.80000,0.69840) (4.90000,0.64185) (5.00000,0.58500) (5.10000,0.52785) (5.20000,0.47040) (5.30000,0.41265) (5.40000,0.35460) (5.50000,0.29625) (5.60000,0.23760) (5.70000,0.17865) (5.80000,0.11940) (5.90000,0.05985) (6.00000,0.00000)};\node[below,align=center] at (3.5,-.65) {Radio $r/r_t$};\node[rotate=90] at (-1,2) {$\sigma/(\rho_b\Omega^2r_t^2)$};\draw[teal,thick] plot coordinates {(0.00000,2.37960) (0.10000,2.35408) (0.20000,2.32819) (0.30000,2.30193) (0.40000,2.27531) (0.50000,2.24830) (0.60000,2.22092) (0.70000,2.19316) (0.80000,2.16501) (0.90000,2.13646) (1.00000,2.10753) (1.10000,2.07819) (1.20000,2.04845) (1.30000,2.01830) (1.40000,1.98773) (1.50000,1.95675) (1.60000,1.92534) (1.70000,1.89350) (1.80000,1.86122) (1.90000,1.82851) (2.00000,1.79535) (2.10000,1.76173) (2.20000,1.72766) (2.30000,1.69312) (2.40000,1.65811) (2.50000,1.62262) (2.60000,1.58664) (2.70000,1.55017) (2.80000,1.51320) (2.90000,1.47572) (3.00000,1.43772) (3.10000,1.39920) (3.20000,1.36014) (3.30000,1.32053) (3.40000,1.28037) (3.50000,1.23965) (3.60000,1.19836) (3.70000,1.15648) (3.80000,1.11400) (3.90000,1.07092) (4.00000,1.02722) (4.10000,0.98289) (4.20000,0.93792) (4.30000,0.89229) (4.40000,0.84598) (4.50000,0.79900) (4.60000,0.75131) (4.70000,0.70290) (4.80000,0.65377) (4.90000,0.60388) (5.00000,0.55323) (5.10000,0.50179) (5.20000,0.44955) (5.30000,0.39648) (5.40000,0.34257) (5.50000,0.28778) (5.60000,0.23211) (5.70000,0.17552) (5.80000,0.11799) (5.90000,0.05949) (6.00000,0.00000)};\node[blue] at (2.8,3.55) {Uniforme};\node[teal] at (2.8,2.85) {Estrechada};
\end{tikzpicture}}%
\count0=\ht0
\advance\count0 by \dp0
\count2=\wd0
\count4=\dimexpr\textheight/100*38\relax
\count6=\linewidth
\ifnum\count0<1
\count0=1
\fi
\ifnum\count2<1
\count2=1
\fi
\divide\count0 by 4096
\divide\count2 by 4096
\divide\count4 by 4096
\divide\count6 by 4096
\ifnum\count0<1
\count0=1
\fi
\ifnum\count2<1
\count2=1
\fi
\ifnum\count4<1
\count4=1
\fi
\ifnum\count6<1
\count6=1
\fi
\count8=\count2
\ifnum\count0>\count4
\count8=\numexpr\count2*\count4/\count0\relax
\fi
\ifnum\count8>\count6
\resizebox{\linewidth}{!}{\begin{tikzpicture}
\draw[->,gray] (0,0)--(7,0);\draw[->,gray] (0,0)--(0,3.8);\draw (0.0,0)--(0.0,-.08) node[below,font=\scriptsize] {0.7};\draw (0,0.0)--(-.08,0.0) node[left,font=\scriptsize] {0};\draw (1.5,0)--(1.5,-.08) node[below,font=\scriptsize] {0.775};\draw (0,0.75)--(-.08,0.75) node[left,font=\scriptsize] {0.0625};\draw (3.0,0)--(3.0,-.08) node[below,font=\scriptsize] {0.85};\draw (0,1.5)--(-.08,1.5) node[left,font=\scriptsize] {0.125};\draw (4.5,0)--(4.5,-.08) node[below,font=\scriptsize] {0.925};\draw (0,2.25)--(-.08,2.25) node[left,font=\scriptsize] {0.188};\draw (6.0,0)--(6.0,-.08) node[below,font=\scriptsize] {1};\draw (0,3.0)--(-.08,3.0) node[left,font=\scriptsize] {0.25};\draw[blue,very thick] plot coordinates {(0.00000,3.06000) (0.10000,3.01785) (0.20000,2.97540) (0.30000,2.93265) (0.40000,2.88960) (0.50000,2.84625) (0.60000,2.80260) (0.70000,2.75865) (0.80000,2.71440) (0.90000,2.66985) (1.00000,2.62500) (1.10000,2.57985) (1.20000,2.53440) (1.30000,2.48865) (1.40000,2.44260) (1.50000,2.39625) (1.60000,2.34960) (1.70000,2.30265) (1.80000,2.25540) (1.90000,2.20785) (2.00000,2.16000) (2.10000,2.11185) (2.20000,2.06340) (2.30000,2.01465) (2.40000,1.96560) (2.50000,1.91625) (2.60000,1.86660) (2.70000,1.81665) (2.80000,1.76640) (2.90000,1.71585) (3.00000,1.66500) (3.10000,1.61385) (3.20000,1.56240) (3.30000,1.51065) (3.40000,1.45860) (3.50000,1.40625) (3.60000,1.35360) (3.70000,1.30065) (3.80000,1.24740) (3.90000,1.19385) (4.00000,1.14000) (4.10000,1.08585) (4.20000,1.03140) (4.30000,0.97665) (4.40000,0.92160) (4.50000,0.86625) (4.60000,0.81060) (4.70000,0.75465) (4.80000,0.69840) (4.90000,0.64185) (5.00000,0.58500) (5.10000,0.52785) (5.20000,0.47040) (5.30000,0.41265) (5.40000,0.35460) (5.50000,0.29625) (5.60000,0.23760) (5.70000,0.17865) (5.80000,0.11940) (5.90000,0.05985) (6.00000,0.00000)};\node[below,align=center] at (3.5,-.65) {Radio $r/r_t$};\node[rotate=90] at (-1,2) {$\sigma/(\rho_b\Omega^2r_t^2)$};\draw[teal,thick] plot coordinates {(0.00000,2.37960) (0.10000,2.35408) (0.20000,2.32819) (0.30000,2.30193) (0.40000,2.27531) (0.50000,2.24830) (0.60000,2.22092) (0.70000,2.19316) (0.80000,2.16501) (0.90000,2.13646) (1.00000,2.10753) (1.10000,2.07819) (1.20000,2.04845) (1.30000,2.01830) (1.40000,1.98773) (1.50000,1.95675) (1.60000,1.92534) (1.70000,1.89350) (1.80000,1.86122) (1.90000,1.82851) (2.00000,1.79535) (2.10000,1.76173) (2.20000,1.72766) (2.30000,1.69312) (2.40000,1.65811) (2.50000,1.62262) (2.60000,1.58664) (2.70000,1.55017) (2.80000,1.51320) (2.90000,1.47572) (3.00000,1.43772) (3.10000,1.39920) (3.20000,1.36014) (3.30000,1.32053) (3.40000,1.28037) (3.50000,1.23965) (3.60000,1.19836) (3.70000,1.15648) (3.80000,1.11400) (3.90000,1.07092) (4.00000,1.02722) (4.10000,0.98289) (4.20000,0.93792) (4.30000,0.89229) (4.40000,0.84598) (4.50000,0.79900) (4.60000,0.75131) (4.70000,0.70290) (4.80000,0.65377) (4.90000,0.60388) (5.00000,0.55323) (5.10000,0.50179) (5.20000,0.44955) (5.30000,0.39648) (5.40000,0.34257) (5.50000,0.28778) (5.60000,0.23211) (5.70000,0.17552) (5.80000,0.11799) (5.90000,0.05949) (6.00000,0.00000)};\node[blue] at (2.8,3.55) {Uniforme};\node[teal] at (2.8,2.85) {Estrechada};
\end{tikzpicture}}%
\else
\ifnum\count0>\count4
\resizebox{!}{\dimexpr\textheight/100*38\relax}{\begin{tikzpicture}
\draw[->,gray] (0,0)--(7,0);\draw[->,gray] (0,0)--(0,3.8);\draw (0.0,0)--(0.0,-.08) node[below,font=\scriptsize] {0.7};\draw (0,0.0)--(-.08,0.0) node[left,font=\scriptsize] {0};\draw (1.5,0)--(1.5,-.08) node[below,font=\scriptsize] {0.775};\draw (0,0.75)--(-.08,0.75) node[left,font=\scriptsize] {0.0625};\draw (3.0,0)--(3.0,-.08) node[below,font=\scriptsize] {0.85};\draw (0,1.5)--(-.08,1.5) node[left,font=\scriptsize] {0.125};\draw (4.5,0)--(4.5,-.08) node[below,font=\scriptsize] {0.925};\draw (0,2.25)--(-.08,2.25) node[left,font=\scriptsize] {0.188};\draw (6.0,0)--(6.0,-.08) node[below,font=\scriptsize] {1};\draw (0,3.0)--(-.08,3.0) node[left,font=\scriptsize] {0.25};\draw[blue,very thick] plot coordinates {(0.00000,3.06000) (0.10000,3.01785) (0.20000,2.97540) (0.30000,2.93265) (0.40000,2.88960) (0.50000,2.84625) (0.60000,2.80260) (0.70000,2.75865) (0.80000,2.71440) (0.90000,2.66985) (1.00000,2.62500) (1.10000,2.57985) (1.20000,2.53440) (1.30000,2.48865) (1.40000,2.44260) (1.50000,2.39625) (1.60000,2.34960) (1.70000,2.30265) (1.80000,2.25540) (1.90000,2.20785) (2.00000,2.16000) (2.10000,2.11185) (2.20000,2.06340) (2.30000,2.01465) (2.40000,1.96560) (2.50000,1.91625) (2.60000,1.86660) (2.70000,1.81665) (2.80000,1.76640) (2.90000,1.71585) (3.00000,1.66500) (3.10000,1.61385) (3.20000,1.56240) (3.30000,1.51065) (3.40000,1.45860) (3.50000,1.40625) (3.60000,1.35360) (3.70000,1.30065) (3.80000,1.24740) (3.90000,1.19385) (4.00000,1.14000) (4.10000,1.08585) (4.20000,1.03140) (4.30000,0.97665) (4.40000,0.92160) (4.50000,0.86625) (4.60000,0.81060) (4.70000,0.75465) (4.80000,0.69840) (4.90000,0.64185) (5.00000,0.58500) (5.10000,0.52785) (5.20000,0.47040) (5.30000,0.41265) (5.40000,0.35460) (5.50000,0.29625) (5.60000,0.23760) (5.70000,0.17865) (5.80000,0.11940) (5.90000,0.05985) (6.00000,0.00000)};\node[below,align=center] at (3.5,-.65) {Radio $r/r_t$};\node[rotate=90] at (-1,2) {$\sigma/(\rho_b\Omega^2r_t^2)$};\draw[teal,thick] plot coordinates {(0.00000,2.37960) (0.10000,2.35408) (0.20000,2.32819) (0.30000,2.30193) (0.40000,2.27531) (0.50000,2.24830) (0.60000,2.22092) (0.70000,2.19316) (0.80000,2.16501) (0.90000,2.13646) (1.00000,2.10753) (1.10000,2.07819) (1.20000,2.04845) (1.30000,2.01830) (1.40000,1.98773) (1.50000,1.95675) (1.60000,1.92534) (1.70000,1.89350) (1.80000,1.86122) (1.90000,1.82851) (2.00000,1.79535) (2.10000,1.76173) (2.20000,1.72766) (2.30000,1.69312) (2.40000,1.65811) (2.50000,1.62262) (2.60000,1.58664) (2.70000,1.55017) (2.80000,1.51320) (2.90000,1.47572) (3.00000,1.43772) (3.10000,1.39920) (3.20000,1.36014) (3.30000,1.32053) (3.40000,1.28037) (3.50000,1.23965) (3.60000,1.19836) (3.70000,1.15648) (3.80000,1.11400) (3.90000,1.07092) (4.00000,1.02722) (4.10000,0.98289) (4.20000,0.93792) (4.30000,0.89229) (4.40000,0.84598) (4.50000,0.79900) (4.60000,0.75131) (4.70000,0.70290) (4.80000,0.65377) (4.90000,0.60388) (5.00000,0.55323) (5.10000,0.50179) (5.20000,0.44955) (5.30000,0.39648) (5.40000,0.34257) (5.50000,0.28778) (5.60000,0.23211) (5.70000,0.17552) (5.80000,0.11799) (5.90000,0.05949) (6.00000,0.00000)};\node[blue] at (2.8,3.55) {Uniforme};\node[teal] at (2.8,2.85) {Estrechada};
\end{tikzpicture}}%
\else
\begin{tikzpicture}
\draw[->,gray] (0,0)--(7,0);\draw[->,gray] (0,0)--(0,3.8);\draw (0.0,0)--(0.0,-.08) node[below,font=\scriptsize] {0.7};\draw (0,0.0)--(-.08,0.0) node[left,font=\scriptsize] {0};\draw (1.5,0)--(1.5,-.08) node[below,font=\scriptsize] {0.775};\draw (0,0.75)--(-.08,0.75) node[left,font=\scriptsize] {0.0625};\draw (3.0,0)--(3.0,-.08) node[below,font=\scriptsize] {0.85};\draw (0,1.5)--(-.08,1.5) node[left,font=\scriptsize] {0.125};\draw (4.5,0)--(4.5,-.08) node[below,font=\scriptsize] {0.925};\draw (0,2.25)--(-.08,2.25) node[left,font=\scriptsize] {0.188};\draw (6.0,0)--(6.0,-.08) node[below,font=\scriptsize] {1};\draw (0,3.0)--(-.08,3.0) node[left,font=\scriptsize] {0.25};\draw[blue,very thick] plot coordinates {(0.00000,3.06000) (0.10000,3.01785) (0.20000,2.97540) (0.30000,2.93265) (0.40000,2.88960) (0.50000,2.84625) (0.60000,2.80260) (0.70000,2.75865) (0.80000,2.71440) (0.90000,2.66985) (1.00000,2.62500) (1.10000,2.57985) (1.20000,2.53440) (1.30000,2.48865) (1.40000,2.44260) (1.50000,2.39625) (1.60000,2.34960) (1.70000,2.30265) (1.80000,2.25540) (1.90000,2.20785) (2.00000,2.16000) (2.10000,2.11185) (2.20000,2.06340) (2.30000,2.01465) (2.40000,1.96560) (2.50000,1.91625) (2.60000,1.86660) (2.70000,1.81665) (2.80000,1.76640) (2.90000,1.71585) (3.00000,1.66500) (3.10000,1.61385) (3.20000,1.56240) (3.30000,1.51065) (3.40000,1.45860) (3.50000,1.40625) (3.60000,1.35360) (3.70000,1.30065) (3.80000,1.24740) (3.90000,1.19385) (4.00000,1.14000) (4.10000,1.08585) (4.20000,1.03140) (4.30000,0.97665) (4.40000,0.92160) (4.50000,0.86625) (4.60000,0.81060) (4.70000,0.75465) (4.80000,0.69840) (4.90000,0.64185) (5.00000,0.58500) (5.10000,0.52785) (5.20000,0.47040) (5.30000,0.41265) (5.40000,0.35460) (5.50000,0.29625) (5.60000,0.23760) (5.70000,0.17865) (5.80000,0.11940) (5.90000,0.05985) (6.00000,0.00000)};\node[below,align=center] at (3.5,-.65) {Radio $r/r_t$};\node[rotate=90] at (-1,2) {$\sigma/(\rho_b\Omega^2r_t^2)$};\draw[teal,thick] plot coordinates {(0.00000,2.37960) (0.10000,2.35408) (0.20000,2.32819) (0.30000,2.30193) (0.40000,2.27531) (0.50000,2.24830) (0.60000,2.22092) (0.70000,2.19316) (0.80000,2.16501) (0.90000,2.13646) (1.00000,2.10753) (1.10000,2.07819) (1.20000,2.04845) (1.30000,2.01830) (1.40000,1.98773) (1.50000,1.95675) (1.60000,1.92534) (1.70000,1.89350) (1.80000,1.86122) (1.90000,1.82851) (2.00000,1.79535) (2.10000,1.76173) (2.20000,1.72766) (2.30000,1.69312) (2.40000,1.65811) (2.50000,1.62262) (2.60000,1.58664) (2.70000,1.55017) (2.80000,1.51320) (2.90000,1.47572) (3.00000,1.43772) (3.10000,1.39920) (3.20000,1.36014) (3.30000,1.32053) (3.40000,1.28037) (3.50000,1.23965) (3.60000,1.19836) (3.70000,1.15648) (3.80000,1.11400) (3.90000,1.07092) (4.00000,1.02722) (4.10000,0.98289) (4.20000,0.93792) (4.30000,0.89229) (4.40000,0.84598) (4.50000,0.79900) (4.60000,0.75131) (4.70000,0.70290) (4.80000,0.65377) (4.90000,0.60388) (5.00000,0.55323) (5.10000,0.50179) (5.20000,0.44955) (5.30000,0.39648) (5.40000,0.34257) (5.50000,0.28778) (5.60000,0.23211) (5.70000,0.17552) (5.80000,0.11799) (5.90000,0.05949) (6.00000,0.00000)};\node[blue] at (2.8,3.55) {Uniforme};\node[teal] at (2.8,2.85) {Estrechada};
\end{tikzpicture}%
\fi
\fi
\endgroup
\par\nopagebreak
\studiumcaption{fig:esfuerzos}{Tensión centrífuga calculada: pala uniforme y pala con sección linealmente estrechada (42 por ciento de reducción a punta). Normalización por rho \(Omega^{2} r_t^{2}\); sin flexión ni concentraciones.}
\end{minipage}
\par\vspace{\baselineskip}
% studium:end-record eyJpZCI6ICJGSUctMDAyMCIsICJzZWN0aW9uIjogImMxMyJ9

\section{Fuerzas aerodinámicas y flexión}

El cambio de cantidad de movimiento tangencial proporciona la fuerza total de fila. Para una distribución uniforme entre \(Z_b\) álabes, una estimación de fuerza por álabe es \(F_\theta\approx\dot m|\Delta C_\theta|/Z_b\). El momento de raíz depende de dónde actúe esa fuerza y de su distribución radial.
Con un modelo de viga esbelta lineal y un momento \(M_b\),

\[
\fitmath{\sigma_b=\frac{M_b y}{I_b}.}
\]
La tensión normal combina tracción y flexión con sus signos en cada fibra. Un desplazamiento del centro de masa puede generar un momento centrífugo que compense parte de la flexión nominal, pero también modifica otras condiciones. No constituye una compensación perfecta de cargas alternas ni de todos los puntos de operación.
Las raíces, cambios de sección y orificios producen concentraciones de tensión. Los modelos de barra y viga no resuelven esos detalles. Tampoco capturan anisotropía de monocristales, fluencia ni plasticidad local. Su función es orientar órdenes de magnitud y detectar incompatibilidades tempranas.
% studium:end-record eyJpZCI6ICJQQVItMDA5MSIsICJzZWN0aW9uIjogImMxMyJ9

\section{Disco y condiciones de contorno}

En un disco axisimétrico de espesor constante, la ecuación radial de equilibrio elástico incluye

\[
\fitmath{\frac{d\sigma_r}{dr}+\frac{\sigma_r-\sigma_\theta}{r}+\rho_b\Omega^2r=0.}
\]
No puede resolverse con esa ecuación solamente: se necesitan compatibilidad, ley constitutiva y condiciones de contorno. Un disco macizo, un disco con agujero y un disco cargado por palas tienen soluciones distintas. No se aplica la expresión de una pala uniforme al disco.
La tensión admisible depende del modo de fallo, temperatura, vida, defecto inicial y margen. Un cociente admisible sobre tensión nominal superior a uno no demuestra seguridad si se han omitido concentración, fatiga o fluencia. El informe científico conserva expresamente esa limitación.
% studium:end-record eyJpZCI6ICJQQVItMDA5MiIsICJzZWN0aW9uIjogImMxMyJ9

\section{Ampliación: frecuencia y excitación}

Como referencia educativa, una viga uniforme en voladizo sin rotación tiene primera frecuencia aproximada

\[
\fitmath{f_1=\frac{1.875^2}{2\pi b^2}\sqrt{\frac{EI_b}{\rho_bA_b}}.}
\]
El valor 1.875 procede de la primera raíz del problema modal del voladizo. La rotación introduce rigidización centrífuga y las condiciones de anclaje reales modifican la frecuencia. Una excitación de orden \(k\) tiene frecuencia \(kN/60\); comparar ambas sugiere posibles cruces, pero la respuesta depende de amortiguamiento y forzado.
Errores habituales: usar densidad del gas en una tensión centrífuga, introducir rpm donde se necesitan radianes por segundo, interpretar tensión elástica como vida y omitir la masa externa al corte. Autocomprobación: duplique \(\Omega\) y explique por qué se cuadruplica la tensión centrífuga ideal.
Formulario: integral de tracción, \(\sigma_b=M_by/I_b\), \(P=\Omega\tau\) y \(f_{exc}=kN/60\). Referencias: ULE Esfuerzos y refrigeración, pp. \(8-15\); MIT 16.50, lección 32, pp. \(1-2\). Deducciones y datos de práctica originales; la evaluación de vida y de piezas reales queda pendiente de revisión especializada.
% studium:end-record eyJpZCI6ICJQQVItMDA5MyIsICJzZWN0aW9uIjogImMxMyJ9

\begin{tcolorbox}[title={Desarrollos del modelo}, studiumdef, breakable]

\noindent\textbf{Modelo desarrollado: Tracción centrífuga uniforme}

\begin{equation}
\sigma_c=\frac{\rho\,\Omega^2}{2}\left(r_t^2-r_h^2\right)
\label{eq:EC-C13}
\end{equation}

\noindent\textit{Significado.} Tensión centrífuga en la raíz de una pala de sección constante. No incluye estrechamiento, flexión ni temperatura.

\noindent\textbf{Hipótesis}
\begin{itemize}
\item Barra radial uniforme, material homogéneo, régimen estacionario.
\end{itemize}

\noindent\textbf{Principios}
\begin{itemize}
\item Balance y definiciones físicas desarrollados en el capítulo; propiedades y dominio declarados en el problema.
\end{itemize}

\noindent\textbf{Aplicabilidad}
\begin{itemize}
\item Modelo educativo del caso C13-01-R
\end{itemize}

\noindent\textbf{Limitaciones}
\begin{itemize}
\item Modelo educativo; la aplicabilidad a una máquina real requiere revisión experta.
\end{itemize}

\noindent\textbf{Pasos}
\begin{enumerate}
\item
Se integra r entre raíz y punta.

\[
\fitmath{\frac{r_t^2 - r_h^2}{2}=- \frac{r_h^{2}}{2} + \frac{r_t^{2}}{2}}
\]
\item
Se multiplica por densidad y velocidad angular al cuadrado.

\[
\fitmath{\sigma=\frac{\Omega^{2} \rho \left(- r_h^{2} + r_t^{2}\right)}{2}}
\]
\end{enumerate}

{\noindent\footnotesize verificada dimensionalmente.\par}
\end{tcolorbox}

\begin{tcolorbox}[title={Problema resuelto 1}, studiumworked, breakable]

\noindent\textbf{Enunciado}
\(C13-01-R \cdot\) Tracción centrífuga uniforme. Calcule la tensión centrífuga nominal de raíz de una pala uniforme. Datos: \(\rho=8120\, \mathrm{kg}\,\mathrm{m}^{-3}\); \(\Omega=1170\, \mathrm{s}^{-1}\); \(r_h=0.245\, \mathrm{m}\); \(r_t=0.302\, \mathrm{m}\).

\noindent\textbf{Datos}
\begin{itemize}
\item $\rho = 8120\,\mathrm{kg}\,\mathrm{m}^{-3}$ — densidad material
\item $\Omega = 1170\,\mathrm{s}^{-1}$ — velocidad angular
\item $r_h = 0.245\,\mathrm{m}$ — radio raíz
\item $r_t = 0.302\,\mathrm{m}$ — radio punta
\end{itemize}

\noindent\textbf{Incógnita}
$\sigma$, en $\mathrm{Pa}$; también las magnitudes intermedias desarrolladas.

\noindent\textbf{Modelo físico}

\noindent\textit{Tracción centrífuga uniforme}

\[
\fitmath{\sigma=\frac{\Omega^{2} \rho \left(- r_h^{2} + r_t^{2}\right)}{2}}
\]

\noindent\textbf{Hipótesis}
\begin{itemize}
\item Barra radial uniforme, material homogéneo, régimen estacionario.
\end{itemize}

\noindent\textbf{Desarrollo paso a paso}

\begin{enumerate}

\item Se integra r entre raíz y punta.

\[\fitmath{\mathrm{integ}=\frac{{\mathrm{rt}}^{2}-{\mathrm{rh}}^{2}}{2}}\]

\[\fitmath{\mathrm{integ}=\frac{{(0.302)}^{2}-{(0.245)}^{2}}{2}\approx 0.0155895\,\mathrm{m^{2}}}\]

\item Se multiplica por densidad y velocidad angular al cuadrado.

\[\fitmath{\mathrm{stress}=\frac{\rho\cdot {\Omega}^{2}\cdot \left({\mathrm{rt}}^{2}-{\mathrm{rh}}^{2}\right)}{2}}\]

\[\fitmath{\mathrm{stress}=\frac{(8120)\cdot {(1170)}^{2}\cdot \left({(0.302)}^{2}-{(0.245)}^{2}\right)}{2}\approx 1.732846\times10^{8}\,\mathrm{Pa}}\]

\end{enumerate}


\noindent\textbf{Sustitución numérica}

\[
\fitmath{\frac{\left(\left(1170\,\mathrm{s}^{-1}\right)\right)^{2} \left(8120\,\mathrm{kg}\,\mathrm{m}^{-3}\right) \left(- \left(0.245\,\mathrm{m}\right)^{2} + \left(0.302\,\mathrm{m}\right)^{2}\right)}{2}}
\]

\noindent\textbf{Resultado}

\[
\fitmath{\sigma=1.732845884\times10^{8}\,\mathrm{Pa}}
\]

\noindent\textbf{Interpretación}
La tensión escala con el cuadrado de giro. No incluye concentraciones de raíz ni flexión.

\noindent\textbf{Comprobación de los pasos}
\begin{itemize}
\item Paso 2: $0.0155895$ — reproducido
\item Paso 4: $1.73285\times 10^{8}$ — reproducido
\end{itemize}

\noindent\textbf{Comprobación de la respuesta} la aritmética del último paso reproduce la respuesta enunciada a la precisión con que está escrita. Es una reproducción aritmética; no valida el modelo físico ni sus hipótesis.

\noindent\textbf{Limitaciones}
\begin{itemize}
\item Modelo educativo; la aplicabilidad a una máquina real requiere revisión experta.
\end{itemize}

\noindent\textbf{Errores frecuentes}
\begin{itemize}
\item Usar densidad del gas o rpm directamente.
\end{itemize}
% studium:end-record eyJpZCI6ICJQUkItMDA5NyIsICJzZWN0aW9uIjogImMxMyJ9
\end{tcolorbox}

\begin{tcolorbox}[title={Problema resuelto 2}, studiumworked, breakable]

\noindent\textbf{Enunciado}
\(C13-02-R \cdot\) Pala estrechada. La sección disminuye linealmente desde Ah en la raíz hasta Ah(1-k) en punta. Calcule la tensión centrífuga nominal de raíz. Datos: \(\rho=8050\, \mathrm{kg}\,\mathrm{m}^{-3}\); \(\Omega=1210\, \mathrm{s}^{-1}\); \(r_h=0.238\, \mathrm{m}\); \(r_t=0.301\, \mathrm{m}\); \(\alpha_t=0.42\).

\noindent\textbf{Datos}
\begin{itemize}
\item $\rho = 8050\,\mathrm{kg}\,\mathrm{m}^{-3}$ — densidad material
\item $\Omega = 1210\,\mathrm{s}^{-1}$ — velocidad angular
\item $r_h = 0.238\,\mathrm{m}$ — radio raíz
\item $r_t = 0.301\,\mathrm{m}$ — radio punta
\item $\alpha_t = 0.42$ — fracción de estrechamiento
\end{itemize}

\noindent\textbf{Incógnita}
$\sigma$, en $\mathrm{Pa}$; también las magnitudes intermedias desarrolladas.

\noindent\textbf{Modelo físico}

\noindent\textit{Pala estrechada}

\[
\fitmath{\sigma=\Omega^{2} \rho \left(- \frac{\alpha_t \left(- \frac{r_h^{3}}{3} - \frac{r_h \left(- r_h^{2} + r_t^{2}\right)}{2} + \frac{r_t^{3}}{3}\right)}{- r_h + r_t} - \frac{r_h^{2}}{2} + \frac{r_t^{2}}{2}\right)}
\]

\noindent\textbf{Hipótesis}
\begin{itemize}
\item Área lineal positiva en todo el intervalo; sin carga aerodinámica.
\end{itemize}

\noindent\textbf{Desarrollo paso a paso}

\begin{enumerate}

\item El ancho radial define la ley de área.

\[\fitmath{b=\mathrm{rt}-\mathrm{rh}}\]

\[\fitmath{b=(0.301)-(0.238)\approx 0.063\,\mathrm{m}}\]

\item La integral se divide por el área raíz; Ah se cancela.

\[\fitmath{\mathrm{stress}=\rho\cdot {\Omega}^{2}\cdot \left(\frac{{\mathrm{rt}}^{2}-{\mathrm{rh}}^{2}}{2}-\frac{\mathrm{kk}}{\mathrm{rt}-\mathrm{rh}}\cdot \left(\frac{{\mathrm{rt}}^{3}-{\mathrm{rh}}^{3}}{3}-\frac{\mathrm{rh}\cdot \left({\mathrm{rt}}^{2}-{\mathrm{rh}}^{2}\right)}{2}\right)\right)}\]

\[\fitmath{\mathrm{stress}=(8050)\cdot {(1210)}^{2}\cdot \left(\frac{{(0.301)}^{2}-{(0.238)}^{2}}{2}-\frac{(0.42)}{(0.301)-(0.238)}\cdot \left(\frac{{(0.301)}^{3}-{(0.238)}^{3}}{3}-\frac{(0.238)\cdot \left({(0.301)}^{2}-{(0.238)}^{2}\right)}{2}\right)\right)\approx 1.564486\times10^{8}\,\mathrm{Pa}}\]

\end{enumerate}


\noindent\textbf{Sustitución numérica}

\[
\fitmath{\left(\left(1210\,\mathrm{s}^{-1}\right)\right)^{2} \left(8050\,\mathrm{kg}\,\mathrm{m}^{-3}\right) \left(- \frac{\left(0.42\right) \left(- \frac{\left(0.238\,\mathrm{m}\right)^{3}}{3} - \frac{\left(0.238\,\mathrm{m}\right) \left(- \left(0.238\,\mathrm{m}\right)^{2} + \left(0.301\,\mathrm{m}\right)^{2}\right)}{2} + \frac{\left(0.301\,\mathrm{m}\right)^{3}}{3}\right)}{- \left(0.238\,\mathrm{m}\right) + \left(0.301\,\mathrm{m}\right)} - \frac{\left(0.238\,\mathrm{m}\right)^{2}}{2} + \frac{\left(0.301\,\mathrm{m}\right)^{2}}{2}\right)}
\]

\noindent\textbf{Resultado}

\[
\fitmath{\sigma=1.564486090\times10^{8}\,\mathrm{Pa}}
\]

\noindent\textbf{Interpretación}
La reducción de masa exterior reduce la tracción raíz respecto de una pala uniforme del mismo material y radios.

\noindent\textbf{Comprobación de los pasos}
\begin{itemize}
\item Paso 2: $0.063$ — reproducido
\item Paso 4: $1.56449\times 10^{8}$ — reproducido
\end{itemize}

\noindent\textbf{Comprobación de la respuesta} la aritmética del último paso reproduce la respuesta enunciada a la precisión con que está escrita. Es una reproducción aritmética; no valida el modelo físico ni sus hipótesis.

\noindent\textbf{Limitaciones}
\begin{itemize}
\item Modelo educativo; la aplicabilidad a una máquina real requiere revisión experta.
\end{itemize}

\noindent\textbf{Errores frecuentes}
\begin{itemize}
\item Reducir tensión multiplicando por el área de punta sin integrar.
\end{itemize}
% studium:end-record eyJpZCI6ICJQUkItMDA5OSIsICJzZWN0aW9uIjogImMxMyJ9
\end{tcolorbox}

\begin{tcolorbox}[title={Problema resuelto 3}, studiumworked, breakable]

\noindent\textbf{Enunciado}
\(C13-03-R \cdot\) Flexión nominal y tracción. Una pala tiene fuerza aerodinámica equivalente aplicada al brazo dado y tracción nominal conocida. Halle la máxima tensión normal en la fibra indicada. Datos: \(F_t=186\, \mathrm{N}\); \(L_b=0.038\, \mathrm{m}\); \(y=0.0042\, \mathrm{m}\); \(I_b=8.4\times 10^{-10}\, \mathrm{m}^{4}\); \(\sigma_N=1.78\times 10^{8}\, \mathrm{Pa}\).

\noindent\textbf{Datos}
\begin{itemize}
\item $F_t = 186\,\mathrm{N}$ — fuerza transversal equivalente
\item $L_b = 0.038\,\mathrm{m}$ — brazo respecto raíz
\item $y = 0.0042\,\mathrm{m}$ — distancia de fibra al eje neutro
\item $I_b = 8.4\times 10^{-10}\,\mathrm{m}^{4}$ — segundo momento de área
\item $\sigma_N = 178000000\,\mathrm{Pa}$ — tracción centrífuga nominal
\end{itemize}

\noindent\textbf{Incógnita}
$\sigma$, en $\mathrm{Pa}$; también las magnitudes intermedias desarrolladas.

\noindent\textbf{Modelo físico}

\noindent\textit{Flexión nominal y tracción}

\[
\fitmath{\sigma=\sigma_N + \frac{L_b F_t y}{I_b}}
\]

\noindent\textbf{Hipótesis}
\begin{itemize}
\item Viga elástica lineal, cargas y sección equivalentes.
\end{itemize}

\noindent\textbf{Desarrollo paso a paso}

\begin{enumerate}

\item La fuerza equivalente genera el momento de raíz.

\[\fitmath{\mathrm{moment}=\mathrm{force}\cdot \mathrm{arm}}\]

\[\fitmath{\mathrm{moment}=(186)\cdot (0.038)\approx 7.068\,\mathrm{N\,m}}\]

\item Se añade la flexión en la fibra de máxima tracción.

\[\fitmath{\mathrm{stress}=\mathrm{sig}+\frac{\mathrm{force}\cdot \mathrm{arm}\cdot \mathrm{yy}}{\mathrm{Ib}}}\]

\[\fitmath{\mathrm{stress}=(178000000)+\frac{(186)\cdot (0.038)\cdot (0.0042)}{(8.4\times10^{-10})}\approx 2.1334\times10^{8}\,\mathrm{Pa}}\]

\end{enumerate}


\noindent\textbf{Sustitución numérica}

\[
\fitmath{\left(178000000\,\mathrm{Pa}\right) + \frac{\left(0.038\,\mathrm{m}\right) \left(186\,\mathrm{N}\right) \left(0.0042\,\mathrm{m}\right)}{\left(8.4e-10\,\mathrm{m}^{4}\right)}}
\]

\noindent\textbf{Resultado}

\[
\fitmath{\sigma=2.133400000\times10^{8}\,\mathrm{Pa}}
\]

\noindent\textbf{Interpretación}
La fibra opuesta resta la contribución de flexión. La suma nominal no evalúa fatiga ni concentraciones.

\noindent\textbf{Comprobación de los pasos}
\begin{itemize}
\item Paso 2: $7.068$ — reproducido
\item Paso 4: $213340000$ — reproducido
\end{itemize}

\noindent\textbf{Comprobación de la respuesta} la aritmética del último paso reproduce la respuesta enunciada a la precisión con que está escrita. Es una reproducción aritmética; no valida el modelo físico ni sus hipótesis.

\noindent\textbf{Limitaciones}
\begin{itemize}
\item Modelo educativo; la aplicabilidad a una máquina real requiere revisión experta.
\end{itemize}

\noindent\textbf{Errores frecuentes}
\begin{itemize}
\item Usar momento polar en una fórmula de flexión sin justificarlo.
\end{itemize}
% studium:end-record eyJpZCI6ICJQUkItMDEwMSIsICJzZWN0aW9uIjogImMxMyJ9
\end{tcolorbox}

\begin{tcolorbox}[title={Problema resuelto 4}, studiumworked, breakable]

\noindent\textbf{Enunciado}
\(C13-04-R \cdot\) Frecuencia del voladizo educativo. Estime la primera frecuencia de una pala como viga uniforme no giratoria. Datos: \(E=1.85\times 10^{11}\, \mathrm{Pa}\); \(I_b=7.8\times 10^{-10}\, \mathrm{m}^{4}\); \(\rho=8030\, \mathrm{kg}\,\mathrm{m}^{-3}\); \(A_b=0.000165\, \mathrm{m}^{2}\); \(L=0.082\, \mathrm{m}\).

\noindent\textbf{Datos}
\begin{itemize}
\item $E = 185000000000\,\mathrm{Pa}$ — módulo elástico
\item $I_b = 7.8\times 10^{-10}\,\mathrm{m}^{4}$ — segundo momento
\item $\rho = 8030\,\mathrm{kg}\,\mathrm{m}^{-3}$ — densidad material
\item $A_b = 0.000165\,\mathrm{m}^{2}$ — área de sección
\item $L = 0.082\,\mathrm{m}$ — longitud
\end{itemize}

\noindent\textbf{Incógnita}
$f_1$, en $\mathrm{s^{-1}}$; también las magnitudes intermedias desarrolladas.

\noindent\textbf{Modelo físico}

\noindent\textit{Frecuencia del voladizo educativo}

\[
\fitmath{f_1=\frac{1.7578125 \sqrt{\frac{E I_b}{A_b \rho}}}{\pi L^{2}}}
\]

\noindent\textbf{Hipótesis}
\begin{itemize}
\item Voladizo uniforme Euler-Bernoulli, sin rotación ni amortiguamiento.
\end{itemize}

\noindent\textbf{Desarrollo paso a paso}

\begin{enumerate}

\item Se calcula la rigidez específica de flexión.

\[\fitmath{\mathrm{stiff}=\sqrt{\frac{\mathrm{Ee}\cdot \mathrm{Ib}}{\rho\cdot \mathrm{Ab}}}}\]

\[\fitmath{\mathrm{stiff}=\sqrt{\frac{(1.85\times10^{11})\cdot (7.8\times10^{-10})}{(8030)\cdot (0.000165)}}\approx 10.43598\,\mathrm{m^{2}/s}}\]

\item La primera raíz modal fija la frecuencia.

\[\fitmath{\mathrm{freq}=\frac{{1.875}^{2}}{2\cdot \pi\cdot {\mathrm{length}}^{2}}\cdot \sqrt{\frac{\mathrm{Ee}\cdot \mathrm{Ib}}{\rho\cdot \mathrm{Ab}}}}\]

\[\fitmath{\mathrm{freq}=\frac{{1.875}^{2}}{2\cdot \pi\cdot {(0.082)}^{2}}\cdot \sqrt{\frac{(1.85\times10^{11})\cdot (7.8\times10^{-10})}{(8030)\cdot (0.000165)}}\approx 868.4172\,\mathrm{s^{-1}}}\]

\end{enumerate}


\noindent\textbf{Sustitución numérica}

\[
\fitmath{\frac{1.7578125 \sqrt{\frac{\left(185000000000\,\mathrm{Pa}\right) \left(7.8e-10\,\mathrm{m}^{4}\right)}{\left(0.000165\,\mathrm{m}^{2}\right) \left(8030\,\mathrm{kg}\,\mathrm{m}^{-3}\right)}}}{\pi \left(0.082\,\mathrm{m}\right)^{2}}}
\]

\noindent\textbf{Resultado}

\[
\fitmath{f_1=8.684171505\times10^{2}\,\mathrm{s}^{-1}}
\]

\noindent\textbf{Interpretación}
Es una estimación no giratoria; no incluye rigidización centrífuga ni anclaje real. Compare después con órdenes de excitación definidos.

\noindent\textbf{Comprobación de los pasos}
\begin{itemize}
\item Paso 2: $10.436$ — reproducido
\item Paso 4: $868.417$ — reproducido
\end{itemize}

\noindent\textbf{Comprobación de la respuesta} la aritmética del último paso reproduce la respuesta enunciada a la precisión con que está escrita. Es una reproducción aritmética; no valida el modelo físico ni sus hipótesis.

\noindent\textbf{Limitaciones}
\begin{itemize}
\item Modelo educativo; la aplicabilidad a una máquina real requiere revisión experta.
\end{itemize}

\noindent\textbf{Errores frecuentes}
\begin{itemize}
\item Declarar ausencia de resonancia a partir de un único modelo estático.
\end{itemize}
% studium:end-record eyJpZCI6ICJQUkItMDEwMyIsICJzZWN0aW9uIjogImMxMyJ9
\end{tcolorbox}

\begin{tcolorbox}[title={Ejercicios propuestos}, studiumworked, breakable]

\noindent\textbf{1 (Básico, Desarrollo)}
\(C13-01-P \cdot\) Tracción centrífuga uniforme. Ejercicio de transferencia. Calcule la tensión centrífuga nominal de raíz de una pala uniforme. Antes de consultar la solución, prediga el efecto de cambiar el primer dato respecto del caso resuelto de esta familia. Justifique la tendencia y explique qué restricción podría convertirse en limitante. Datos: \(\rho=8647.8\, \mathrm{kg}\,\mathrm{m}^{-3}\); \(\Omega=1170\, \mathrm{s}^{-1}\); \(r_h=0.245\, \mathrm{m}\); \(r_t=0.302\, \mathrm{m}\).

\noindent\textbf{2 (Intermedio, Desarrollo)}
\(C13-02-P \cdot\) Pala estrechada. Ejercicio de transferencia. La sección disminuye linealmente desde Ah en la raíz hasta Ah(1-k) en punta. Calcule la tensión centrífuga nominal de raíz. Antes de consultar la solución, prediga el efecto de cambiar el primer dato respecto del caso resuelto de esta familia. Justifique la tendencia y explique qué restricción podría convertirse en limitante. Datos: \(\rho=8573.25\, \mathrm{kg}\,\mathrm{m}^{-3}\); \(\Omega=1210\, \mathrm{s}^{-1}\); \(r_h=0.238\, \mathrm{m}\); \(r_t=0.301\, \mathrm{m}\); \(\alpha_t=0.42\).

\noindent\textbf{3 (Avanzado, Desarrollo)}
\(C13-03-P \cdot\) Flexión nominal y tracción. Ejercicio de transferencia. Una pala tiene fuerza aerodinámica equivalente aplicada al brazo dado y tracción nominal conocida. Halle la máxima tensión normal en la fibra indicada. Antes de consultar la solución, prediga el efecto de cambiar el primer dato respecto del caso resuelto de esta familia. Justifique la tendencia y explique qué restricción podría convertirse en limitante. Datos: \(F_t=198.09\, \mathrm{N}\); \(L_b=0.038\, \mathrm{m}\); \(y=0.0042\, \mathrm{m}\); \(I_b=8.4\times 10^{-10}\, \mathrm{m}^{4}\); \(\sigma_N=1.78\times 10^{8}\, \mathrm{Pa}\).

\noindent\textbf{4 (Nivel de examen, Desarrollo)}
\(C13-04-P \cdot\) Frecuencia del voladizo educativo. Ejercicio de transferencia. Estime la primera frecuencia de una pala como viga uniforme no giratoria. Antes de consultar la solución, prediga el efecto de cambiar el primer dato respecto del caso resuelto de esta familia. Justifique la tendencia y explique qué restricción podría convertirse en limitante. Datos: \(E=1.97025\times 10^{11}\, \mathrm{Pa}\); \(I_b=7.8\times 10^{-10}\, \mathrm{m}^{4}\); \(\rho=8030\, \mathrm{kg}\,\mathrm{m}^{-3}\); \(A_b=0.000165\, \mathrm{m}^{2}\); \(L=0.082\, \mathrm{m}\).
\end{tcolorbox}


\studiumchaptercheck

\chapter{Transferencia térmica y refrigeración de álabes}
\label{chap:c14}

\studiumchapterguide{Objetivo: desarrollar transferencia térmica y refrigeración de álabes, declarar hipótesis y comprobar condiciones físicas.}
% studium:end-record eyJpZCI6ICJQQVItMDA5NCIsICJzZWN0aW9uIjogImMxNCJ9

\section{Introducción y objetivos}

La temperatura del gas puede superar ampliamente la temperatura admisible del metal. La refrigeración reduce la temperatura de la pala, pero consume aire que ya ha requerido trabajo de compresión e introduce mezcla y pérdidas. El objetivo es formular balances térmicos sencillos, diferenciar temperaturas relevantes y explicar el compromiso entre protección y rendimiento. Pertenece al bloque IV.
Los modelos del capítulo son educativos. Una red térmica de una dimensión no determina la vida de una pala ni el caudal de refrigeración seguro de un motor real. La distribución de coeficientes, propiedades y temperatura requiere ensayos o simulación validada.
% studium:end-record eyJpZCI6ICJQQVItMDA5NSIsICJzZWN0aW9uIjogImMxNCJ9

\section{Temperatura que ve la pared}

La temperatura estática no es, en general, la temperatura que alcanzaría una pared adiabática. Se introduce una temperatura de recuperación

\[
\fitmath{T_{aw}=T+r_{rec}\frac{V^2}{2c_p},}
\]
donde \(V\) es la velocidad en la referencia de la pared y \(r_{rec}\) el factor de recuperación. En un rotor debe usarse la referencia relativa. No se identifica automáticamente \(T_{aw}\) con \(T_0\); coincide si la recuperación es completa dentro del modelo.
El flujo térmico convectivo se escribe \(q''=h_g(T_{aw}-T_{w,g})\). El coeficiente \(h_g\) depende de Reynolds, Mach, turbulencia, curvatura, transición y rugosidad. El borde de ataque y las zonas de transición pueden tener demandas muy distintas de una media de superficie.
% studium:end-record eyJpZCI6ICJQQVItMDA5NiIsICJzZWN0aW9uIjogImMxNCJ9

\section{Red de resistencias térmicas}

Para una pared plana de espesor \(t\), conductividad \(k_m\), convección exterior \(h_g\) e interior \(h_c\), en régimen estacionario sin radiación:

\[
\fitmath{R''_{t}=\frac1{h_g}+\frac{t}{k_m}+\frac1{h_c},\qquad
q''=\frac{T_{aw}-T_c}{R''_{t}}.}
\]
Las temperaturas de las caras son

\[
\fitmath{T_{w,g}=T_{aw}-\frac{q''}{h_g},\qquad
T_{w,c}=T_c+\frac{q''}{h_c}.}
\]
Debe cumplirse \(T_c<T_{w,c}<T_{w,g}<T_{aw}\) cuando el calor fluye del gas al refrigerante. Si no se cumple, se revisan signos, temperaturas y resistencias. La figura~\ref{fig:refrigeracion} representa una sección original esquemática con convección interna, impacto y película; no reproduce una pieza de fabricante.
Un recubrimiento añade una resistencia \(t_{TBC}/k_{TBC}\), pero también introduce interfaces, tensiones y degradación. Reducir el flujo térmico no basta para asegurar su integridad. En una pared curva o con conductos la geometría de áreas debe tratarse de manera coherente.
% studium:end-record eyJpZCI6ICJQQVItMDA5NyIsICJzZWN0aW9uIjogImMxNCJ9

\par\vspace{\baselineskip}
\noindent\begin{minipage}{\linewidth}
\centering
\begingroup
\setbox0=\hbox{\begin{tikzpicture}
\draw[thick,fill=gray!15] (0,0) .. controls (0,2.5) and (2,3.3) .. (7,1) .. controls (3,1.5) and (1,.2) .. (0,0);\draw[blue,thick,->] (.6,.5)--(1,1.5)--(2.6,2)--(4.8,1.5);\node[blue,below] at (2.5,.4) {Convección interna};\draw[blue,thick,->] (1,1.5)--(.4,1.7);\node[blue,above] at (.4,2.3) {Impacto};\foreach \x/\y in {2.6/2.6,3.8/2.35,5/1.95}{\draw[blue,->] (\x-.1,\y-.5)--(\x,\y)--(\x+.7,\y-.1);}\node[blue] at (5.8,3.2) {Película};\draw[red,very thick,->] (-.4,3.4)--(2,3.4);\node[red,above] at (1.3,3.4) {Gas caliente};
\end{tikzpicture}}%
\count0=\ht0
\advance\count0 by \dp0
\count2=\wd0
\count4=\dimexpr\textheight/100*38\relax
\count6=\linewidth
\ifnum\count0<1
\count0=1
\fi
\ifnum\count2<1
\count2=1
\fi
\divide\count0 by 4096
\divide\count2 by 4096
\divide\count4 by 4096
\divide\count6 by 4096
\ifnum\count0<1
\count0=1
\fi
\ifnum\count2<1
\count2=1
\fi
\ifnum\count4<1
\count4=1
\fi
\ifnum\count6<1
\count6=1
\fi
\count8=\count2
\ifnum\count0>\count4
\count8=\numexpr\count2*\count4/\count0\relax
\fi
\ifnum\count8>\count6
\resizebox{\linewidth}{!}{\begin{tikzpicture}
\draw[thick,fill=gray!15] (0,0) .. controls (0,2.5) and (2,3.3) .. (7,1) .. controls (3,1.5) and (1,.2) .. (0,0);\draw[blue,thick,->] (.6,.5)--(1,1.5)--(2.6,2)--(4.8,1.5);\node[blue,below] at (2.5,.4) {Convección interna};\draw[blue,thick,->] (1,1.5)--(.4,1.7);\node[blue,above] at (.4,2.3) {Impacto};\foreach \x/\y in {2.6/2.6,3.8/2.35,5/1.95}{\draw[blue,->] (\x-.1,\y-.5)--(\x,\y)--(\x+.7,\y-.1);}\node[blue] at (5.8,3.2) {Película};\draw[red,very thick,->] (-.4,3.4)--(2,3.4);\node[red,above] at (1.3,3.4) {Gas caliente};
\end{tikzpicture}}%
\else
\ifnum\count0>\count4
\resizebox{!}{\dimexpr\textheight/100*38\relax}{\begin{tikzpicture}
\draw[thick,fill=gray!15] (0,0) .. controls (0,2.5) and (2,3.3) .. (7,1) .. controls (3,1.5) and (1,.2) .. (0,0);\draw[blue,thick,->] (.6,.5)--(1,1.5)--(2.6,2)--(4.8,1.5);\node[blue,below] at (2.5,.4) {Convección interna};\draw[blue,thick,->] (1,1.5)--(.4,1.7);\node[blue,above] at (.4,2.3) {Impacto};\foreach \x/\y in {2.6/2.6,3.8/2.35,5/1.95}{\draw[blue,->] (\x-.1,\y-.5)--(\x,\y)--(\x+.7,\y-.1);}\node[blue] at (5.8,3.2) {Película};\draw[red,very thick,->] (-.4,3.4)--(2,3.4);\node[red,above] at (1.3,3.4) {Gas caliente};
\end{tikzpicture}}%
\else
\begin{tikzpicture}
\draw[thick,fill=gray!15] (0,0) .. controls (0,2.5) and (2,3.3) .. (7,1) .. controls (3,1.5) and (1,.2) .. (0,0);\draw[blue,thick,->] (.6,.5)--(1,1.5)--(2.6,2)--(4.8,1.5);\node[blue,below] at (2.5,.4) {Convección interna};\draw[blue,thick,->] (1,1.5)--(.4,1.7);\node[blue,above] at (.4,2.3) {Impacto};\foreach \x/\y in {2.6/2.6,3.8/2.35,5/1.95}{\draw[blue,->] (\x-.1,\y-.5)--(\x,\y)--(\x+.7,\y-.1);}\node[blue] at (5.8,3.2) {Película};\draw[red,very thick,->] (-.4,3.4)--(2,3.4);\node[red,above] at (1.3,3.4) {Gas caliente};
\end{tikzpicture}%
\fi
\fi
\endgroup
\par\nopagebreak
\studiumcaption{fig:refrigeracion}{Convección interna, impacto y película: sección esquemática original, sin escala de fabricación.}
\end{minipage}
\par\vspace{\baselineskip}
% studium:end-record eyJpZCI6ICJGSUctMDAxNCIsICJzZWN0aW9uIjogImMxNCJ9

\section{Balance del refrigerante}

El calor absorbido se relaciona con el incremento de entalpía del caudal refrigerante:

\[
\fitmath{\dot Q=\dot m_c c_{p,c}(T_{c,s}-T_{c,e}).}
\]
Si se conocen flujo térmico medio y área mojada, \(\dot Q=q''A_w\). El área de intercambio no es el área anular de paso. El balance establece un caudal mínimo energético para el incremento de temperatura permitido; no asegura que los coeficientes de convección sean suficientes para mantener el metal dentro del límite.
La presión de refrigerante debe permitir atravesar los conductos e inyectarse en el gas. Sangrar aire más adelante en el compresor puede aumentar la presión disponible, pero también la temperatura de entrada del refrigerante. Por ello la selección de sangrado exige simultáneamente balance térmico y pérdidas de presión.
% studium:end-record eyJpZCI6ICJQQVItMDA5OCIsICJzZWN0aW9uIjogImMxNCJ9

\section{Película e impacto}


La efectividad adiabática de película se define como



\[
\fitmath{\eta_f=\frac{T_{aw,0}-T_{aw,f}}{T_{aw,0}-T_{c,e}}.}
\]
No es una eficiencia termodinámica. Relaciona temperaturas de recuperación sin y con inyección para unas condiciones geométricas y de flujo. Se calcula \(T_{aw,f}\) y después el calor con un coeficiente exterior compatible con la inyección. Mantener \(h_g\) inalterado es una hipótesis adicional: el chorro puede modificar la capa límite.
Más caudal de película no garantiza siempre mejor protección local. El desprendimiento del chorro y la mezcla pueden modificar la cobertura. La refrigeración por impacto concentra transferencia donde se necesita, pero consume presión y depende de separación y diámetro de chorros. La transpiración distribuye inyección a través de material poroso y plantea problemas de fabricación y durabilidad. Estas técnicas no se ordenan mediante una clasificación universal de eficiencia.
% studium:end-record eyJpZCI6ICJQQVItMDA5OSIsICJzZWN0aW9uIjogImMxNCJ9

\section{Tensiones térmicas y compromiso}

Un aumento uniforme de temperatura de una barra libre causa expansión, no tensión térmica de restricción. Si una barra uniaxial impide por completo esa expansión, el modelo elástico da \(\sigma=-E\alpha_T\Delta T\). Una pared con gradiente térmico, restricción biaxial y curvatura tiene otra solución; no se usa la fórmula de barra como ley de toda pala.
La refrigeración puede reducir la temperatura media y aumentar gradientes locales. Menor temperatura media y menor tensión térmica no son objetivos idénticos. La evaluación combina campo térmico, estado tensional, fatiga térmica, fluencia y oxidación. El consumo de aire afecta además el trabajo y la entropía de mezcla de la turbina.
Errores habituales: utilizar velocidad absoluta para recuperación de una pared móvil, omitir una resistencia, confundir efectividad de película con eficiencia del escalón y calcular tensión en una expansión libre. Autocomprobación: aumente \(h_c\) indefinidamente y determine qué resistencias siguen limitando el flujo de calor.
Formulario: red térmica anterior, \(\dot Q=\dot m_cc_{p,c}\Delta T_c\), definición de \(\eta_f\). Referencias: ULE Esfuerzos y refrigeración, pp. \(16-26\); MIT 16.50, lecciones 29 y 30, para contexto de balances y mecanismos. No se importan correlaciones numéricas de transferencia sin su rango de validez. Las geometrías y los problemas son originales.
% studium:end-record eyJpZCI6ICJQQVItMDEwMCIsICJzZWN0aW9uIjogImMxNCJ9

\begin{tcolorbox}[title={Desarrollos del modelo}, studiumdef, breakable]

\noindent\textbf{Modelo desarrollado: Red térmica de pared}

\begin{equation}
T_w=T_{aw}-\frac{T_{aw}-T_c}{h_g\left(\frac1{h_g}+\frac1{h_c}+\frac{t}{k_m}\right)}
\label{eq:EC-C14}
\end{equation}

\noindent\textit{Significado.} Temperatura de la cara caliente en una red térmica unidimensional. La temperatura motriz es la de pared adiabática $T_{aw}$, no una temperatura estática de gas no declarada.

\noindent\textbf{Hipótesis}
\begin{itemize}
\item Estacionario, una dimensión, sin radiación, coeficientes constantes.
\end{itemize}

\noindent\textbf{Principios}
\begin{itemize}
\item Balance y definiciones físicas desarrollados en el capítulo; propiedades y dominio declarados en el problema.
\end{itemize}

\noindent\textbf{Aplicabilidad}
\begin{itemize}
\item Modelo educativo del caso C14-01-R
\end{itemize}

\noindent\textbf{Limitaciones}
\begin{itemize}
\item Modelo educativo; la aplicabilidad a una máquina real requiere revisión experta.
\end{itemize}

\noindent\textbf{Pasos}
\begin{enumerate}
\item
Las resistencias se suman por unidad de área.

\[
\fitmath{R_{\mathrm{th}}=\frac{1}{h_g} + \frac{1}{h_c} + \frac{t}{k_m}}
\]
\item
El salto térmico impulsa el flujo.

\[
\fitmath{q''=\frac{- T_c + T_{aw}}{\frac{1}{h_g} + \frac{1}{h_c} + \frac{t}{k_m}}}
\]
\item
La caída exterior determina la cara caliente.

\[
\fitmath{T_w=T_{aw} - \frac{- T_c + T_{aw}}{h_g \left(\frac{1}{h_g} + \frac{1}{h_c} + \frac{t}{k_m}\right)}}
\]
\end{enumerate}

{\noindent\footnotesize verificada dimensionalmente.\par}
\end{tcolorbox}

\begin{tcolorbox}[title={Problema resuelto 1}, studiumworked, breakable]

\noindent\textbf{Enunciado}
\(C14-01-R \cdot\) Red térmica de pared. Calcule flujo térmico y temperatura exterior del metal en una pared plana con convección a ambos lados. Datos: \(T_{aw}=1620\, \mathrm{K}\); \(T_c=745\, \mathrm{K}\); \(h_g=1680\, \mathrm{W}\,\mathrm{m}^{-2}\,\mathrm{K}^{-1}\); \(h_c=2450\, \mathrm{W}\,\mathrm{m}^{-2}\,\mathrm{K}^{-1}\); \(t=0.0011\, \mathrm{m}\); \(k_m=21\, \mathrm{W}\,\mathrm{m}^{-1}\,\mathrm{K}^{-1}\).

\noindent\textbf{Datos}
\begin{itemize}
\item $T_{aw} = 1620\,\mathrm{K}$ — temperatura recuperación exterior
\item $T_c = 745\,\mathrm{K}$ — temperatura refrigerante local
\item $h_g = 1680\,\mathrm{W}\,\mathrm{m}^{-2}\,\mathrm{K}^{-1}$ — coeficiente exterior
\item $h_c = 2450\,\mathrm{W}\,\mathrm{m}^{-2}\,\mathrm{K}^{-1}$ — coeficiente interior
\item $t = 0.0011\,\mathrm{m}$ — espesor
\item $k_m = 21\,\mathrm{W}\,\mathrm{m}^{-1}\,\mathrm{K}^{-1}$ — conductividad metal
\end{itemize}

\noindent\textbf{Incógnita}
$T_w$, en $\mathrm{K}$; también las magnitudes intermedias desarrolladas.

\noindent\textbf{Modelo físico}

\noindent\textit{Red térmica de pared}

\[
\fitmath{T_w=T_{aw} - \frac{- T_c + T_{aw}}{h_g \left(\frac{1}{h_g} + \frac{1}{h_c} + \frac{t}{k_m}\right)}}
\]

\noindent\textbf{Hipótesis}
\begin{itemize}
\item Estacionario, una dimensión, sin radiación, coeficientes constantes.
\end{itemize}

\noindent\textbf{Desarrollo paso a paso}

\begin{enumerate}

\item Las resistencias se suman por unidad de área.

\[\fitmath{\mathrm{Res}=\frac{1}{\mathrm{hg}}+\frac{\mathrm{th}}{\mathrm{cond}}+\frac{1}{\mathrm{hc}}}\]

\[\fitmath{\mathrm{Res}=\frac{1}{(1680)}+\frac{(0.0011)}{(21)}+\frac{1}{(2450)}\approx 0.001055782\,\mathrm{m^{2}\,K/W}}\]

\item El salto térmico impulsa el flujo.

\[\fitmath{\mathrm{flux}=\frac{\mathrm{Tg}-\mathrm{Tc}}{\frac{1}{\mathrm{hg}}+\frac{\mathrm{th}}{\mathrm{cond}}+\frac{1}{\mathrm{hc}}}}\]

\[\fitmath{\mathrm{flux}=\frac{(1620)-(745)}{\frac{1}{(1680)}+\frac{(0.0011)}{(21)}+\frac{1}{(2450)}}\approx 828769.3\,\mathrm{W/m^{2}}}\]

\item La caída exterior determina la cara caliente.

\[\fitmath{\mathrm{Tw}=\mathrm{Tg}-\frac{\mathrm{Tg}-\mathrm{Tc}}{\mathrm{hg}\cdot \left(\frac{1}{\mathrm{hg}}+\frac{\mathrm{th}}{\mathrm{cond}}+\frac{1}{\mathrm{hc}}\right)}}\]

\[\fitmath{\mathrm{Tw}=(1620)-\frac{(1620)-(745)}{(1680)\cdot \left(\frac{1}{(1680)}+\frac{(0.0011)}{(21)}+\frac{1}{(2450)}\right)}\approx 1126.685\,\mathrm{K}}\]

\end{enumerate}


\noindent\textbf{Sustitución numérica}

\[
\fitmath{\left(1620\,\mathrm{K}\right) - \frac{- \left(745\,\mathrm{K}\right) + \left(1620\,\mathrm{K}\right)}{\left(1680\,\mathrm{W}\,\mathrm{m}^{-2}\,\mathrm{K}^{-1}\right) \left(\frac{1}{\left(1680\,\mathrm{W}\,\mathrm{m}^{-2}\,\mathrm{K}^{-1}\right)} + \frac{1}{\left(2450\,\mathrm{W}\,\mathrm{m}^{-2}\,\mathrm{K}^{-1}\right)} + \frac{\left(0.0011\,\mathrm{m}\right)}{\left(21\,\mathrm{W}\,\mathrm{m}^{-1}\,\mathrm{K}^{-1}\right)}\right)}}
\]

\noindent\textbf{Resultado}

\[
\fitmath{T_w=1.126684923\times10^{3}\,\mathrm{K}}
\]

\noindent\textbf{Interpretación}
La temperatura del metal debe quedar entre refrigerante y recuperación exterior. La diferencia de caras corresponde al flujo por conducción.

\noindent\textbf{Comprobación de los pasos}
\begin{itemize}
\item Paso 2: $0.00105578$ — reproducido
\item Paso 4: $828769$ — reproducido
\item Paso 6: $1126.68$ — reproducido
\end{itemize}

\noindent\textbf{Comprobación de la respuesta} la aritmética del último paso reproduce la respuesta enunciada a la precisión con que está escrita. Es una reproducción aritmética; no valida el modelo físico ni sus hipótesis.

\noindent\textbf{Limitaciones}
\begin{itemize}
\item Modelo educativo; la aplicabilidad a una máquina real requiere revisión experta.
\end{itemize}

\noindent\textbf{Errores frecuentes}
\begin{itemize}
\item Confundir área anular con área de intercambio.
\end{itemize}
% studium:end-record eyJpZCI6ICJQUkItMDEwNSIsICJzZWN0aW9uIjogImMxNCJ9
\end{tcolorbox}

\begin{tcolorbox}[title={Problema resuelto 2}, studiumworked, breakable]

\noindent\textbf{Enunciado}
\(C14-02-R \cdot\) Caudal energético de refrigeración. Una superficie recibe flujo térmico medio conocido. Calcule el caudal refrigerante que absorbería ese calor con el incremento permitido. Datos: \(q''=620000\, \mathrm{W}\,\mathrm{m}^{-2}\); \(A=0.034\, \mathrm{m}^{2}\); \(c_p=1090\, \mathrm{J}\,\mathrm{kg}^{-1}\,\mathrm{K}^{-1}\); \(\Delta T_c=185\, \mathrm{K}\).

\noindent\textbf{Datos}
\begin{itemize}
\item $q'' = 620000\,\mathrm{W}\,\mathrm{m}^{-2}$ — flujo térmico medio
\item $A = 0.034\,\mathrm{m}^{2}$ — área de intercambio
\item $c_p = 1090\,\mathrm{J}\,\mathrm{kg}^{-1}\,\mathrm{K}^{-1}$ — calor específico refrigerante
\item $\Delta T_c = 185\,\mathrm{K}$ — incremento de temperatura permitido
\end{itemize}

\noindent\textbf{Incógnita}
$\dot{m}_c$, en $\mathrm{kg\,s^{-1}}$; también las magnitudes intermedias desarrolladas.

\noindent\textbf{Modelo físico}

\noindent\textit{Caudal energético de refrigeración}

\[
\fitmath{\dot{m}_c=\frac{A q''}{c_p \Delta T_c}}
\]

\noindent\textbf{Hipótesis}
\begin{itemize}
\item Sin pérdidas térmicas adicionales; cp constante.
\end{itemize}

\noindent\textbf{Desarrollo paso a paso}

\begin{enumerate}

\item Se integra el flujo térmico uniforme sobre la superficie.

\[\fitmath{\mathrm{heat}=\mathrm{flux}\cdot \mathrm{area}}\]

\[\fitmath{\mathrm{heat}=(620000)\cdot (0.034)\approx 21080\,\mathrm{W}}\]

\item El balance de entalpía fija el caudal.

\[\fitmath{\mathrm{mc}=\frac{\mathrm{flux}\cdot \mathrm{area}}{\mathrm{cp}\cdot \mathrm{dt}}}\]

\[\fitmath{\mathrm{mc}=\frac{(620000)\cdot (0.034)}{(1090)\cdot (185)}\approx 0.1045376\,\mathrm{kg/s}}\]

\end{enumerate}


\noindent\textbf{Sustitución numérica}

\[
\fitmath{\frac{\left(0.034\,\mathrm{m}^{2}\right) \left(620000\,\mathrm{W}\,\mathrm{m}^{-2}\right)}{\left(1090\,\mathrm{J}\,\mathrm{kg}^{-1}\,\mathrm{K}^{-1}\right) \left(185\,\mathrm{K}\right)}}
\]

\noindent\textbf{Resultado}

\[
\fitmath{\dot{m}_c=1.045375651\times10^{-1}\,\mathrm{kg}\,\mathrm{s}^{-1}}
\]

\noindent\textbf{Interpretación}
El resultado es un mínimo del balance energético. La convección y la presión disponible deben comprobarse aparte.

\noindent\textbf{Comprobación de los pasos}
\begin{itemize}
\item Paso 2: $21080$ — reproducido
\item Paso 4: $0.104538$ — reproducido
\end{itemize}

\noindent\textbf{Comprobación de la respuesta} la aritmética del último paso reproduce la respuesta enunciada a la precisión con que está escrita. Es una reproducción aritmética; no valida el modelo físico ni sus hipótesis.

\noindent\textbf{Limitaciones}
\begin{itemize}
\item Modelo educativo; la aplicabilidad a una máquina real requiere revisión experta.
\end{itemize}

\noindent\textbf{Errores frecuentes}
\begin{itemize}
\item Tomar este mínimo como caudal seguro de una pala real.
\end{itemize}
% studium:end-record eyJpZCI6ICJQUkItMDEwNyIsICJzZWN0aW9uIjogImMxNCJ9
\end{tcolorbox}

\begin{tcolorbox}[title={Problema resuelto 3}, studiumworked, breakable]

\noindent\textbf{Enunciado}
\(C14-03-R \cdot\) Película y recuperación. Con efectividad de película prescrita, calcule la recuperación efectiva y el flujo térmico a una pared de temperatura dada. Datos: \(T_{aw}=1690\, \mathrm{K}\); \(T_c=775\, \mathrm{K}\); \(\eta_f=0.41\); \(T_w=1130\, \mathrm{K}\); \(h_g=1770\, \mathrm{W}\,\mathrm{m}^{-2}\,\mathrm{K}^{-1}\).

\noindent\textbf{Datos}
\begin{itemize}
\item $T_{aw} = 1690\,\mathrm{K}$ — recuperación sin película
\item $T_c = 775\,\mathrm{K}$ — temperatura refrigerante
\item $\eta_f = 0.41$ — efectividad de película
\item $T_w = 1130\,\mathrm{K}$ — temperatura pared
\item $h_g = 1770\,\mathrm{W}\,\mathrm{m}^{-2}\,\mathrm{K}^{-1}$ — coeficiente exterior compatible con inyección
\end{itemize}

\noindent\textbf{Incógnita}
$q''$, en $\mathrm{W\,m^{-2}}$; también las magnitudes intermedias desarrolladas.

\noindent\textbf{Modelo físico}

\noindent\textit{Película y recuperación}

\[
\fitmath{q''=h_g \left(T_{aw} - T_w - \eta_f \left(- T_c + T_{aw}\right)\right)}
\]

\noindent\textbf{Hipótesis}
\begin{itemize}
\item Efectividad y coeficiente dados para la misma geometría y punto.
\end{itemize}

\noindent\textbf{Desarrollo paso a paso}

\begin{enumerate}

\item Se despeja la recuperación con película.

\[\fitmath{\mathrm{Taw}=\mathrm{Tg}-\mathrm{eff}\cdot \left(\mathrm{Tg}-\mathrm{Tc}\right)}\]

\[\fitmath{\mathrm{Taw}=(1690)-(0.41)\cdot \left((1690)-(775)\right)\approx 1314.85\,\mathrm{K}}\]

\item Se aplica convección con el coeficiente especificado.

\[\fitmath{\mathrm{flux}=\mathrm{hg}\cdot \left(\mathrm{Tg}-\mathrm{eff}\cdot \left(\mathrm{Tg}-\mathrm{Tc}\right)-\mathrm{Tw}\right)}\]

\[\fitmath{\mathrm{flux}=(1770)\cdot \left((1690)-(0.41)\cdot \left((1690)-(775)\right)-(1130)\right)\approx 327184.5\,\mathrm{W/m^{2}}}\]

\end{enumerate}


\noindent\textbf{Sustitución numérica}

\[
\fitmath{\left(1770\,\mathrm{W}\,\mathrm{m}^{-2}\,\mathrm{K}^{-1}\right) \left(\left(1690\,\mathrm{K}\right) - \left(1130\,\mathrm{K}\right) - \left(0.41\right) \left(- \left(775\,\mathrm{K}\right) + \left(1690\,\mathrm{K}\right)\right)\right)}
\]

\noindent\textbf{Resultado}

\[
\fitmath{q''=3.271845000\times10^{5}\,\mathrm{W}\,\mathrm{m}^{-2}}
\]

\noindent\textbf{Interpretación}
La efectividad no sustituye al coeficiente de convección. Verifique que el flujo conserva el sentido esperado.

\noindent\textbf{Comprobación de los pasos}
\begin{itemize}
\item Paso 2: $1314.85$ — reproducido
\item Paso 4: $327185$ — reproducido
\end{itemize}

\noindent\textbf{Comprobación de la respuesta} la aritmética del último paso reproduce la respuesta enunciada a la precisión con que está escrita. Es una reproducción aritmética; no valida el modelo físico ni sus hipótesis.

\noindent\textbf{Limitaciones}
\begin{itemize}
\item Modelo educativo; la aplicabilidad a una máquina real requiere revisión experta.
\end{itemize}

\noindent\textbf{Errores frecuentes}
\begin{itemize}
\item Identificar efectividad de película y eficiencia de turbina.
\end{itemize}
% studium:end-record eyJpZCI6ICJQUkItMDEwOSIsICJzZWN0aW9uIjogImMxNCJ9
\end{tcolorbox}

\begin{tcolorbox}[title={Problema resuelto 4}, studiumworked, breakable]

\noindent\textbf{Enunciado}
\(C14-04-R \cdot\) Restricción térmica uniaxial. Una barra homogénea tiene impedida completamente su dilatación longitudinal. Determine la deformación térmica libre y la tensión elástica de compresión. Datos: \(E=1.74\times 10^{11}\, \mathrm{Pa}\); \(\alpha_{\mathrm{th}}=1.41\times 10^{-5}\, \mathrm{K}^{-1}\); \(\Delta T=128\, \mathrm{K}\).

\noindent\textbf{Datos}
\begin{itemize}
\item $E = 174000000000\,\mathrm{Pa}$ — módulo elástico
\item $\alpha_{\mathrm{th}} = 1.41\times 10^{-5}\,\mathrm{K}^{-1}$ — coeficiente dilatación
\item $\Delta T = 128\,\mathrm{K}$ — aumento uniforme de temperatura
\end{itemize}

\noindent\textbf{Incógnita}
$\sigma$, en $\mathrm{Pa}$; también las magnitudes intermedias desarrolladas.

\noindent\textbf{Modelo físico}

\noindent\textit{Restricción térmica uniaxial}

\[
\fitmath{\sigma=- E \alpha_{\mathrm{th}} \Delta T}
\]

\noindent\textbf{Hipótesis}
\begin{itemize}
\item Restricción total uniaxial, elasticidad lineal, sin relajación.
\end{itemize}

\noindent\textbf{Desarrollo paso a paso}

\begin{enumerate}

\item La expansión libre sería positiva.

\[\fitmath{\mathrm{strain}=\mathrm{alp}\cdot \mathrm{dt}}\]

\[\fitmath{\mathrm{strain}=(1.41\times10^{-5})\cdot (128)\approx 0.0018048\,}\]

\item La restricción exige tensión elástica de signo opuesto.

\[\fitmath{\mathrm{stress}=-\mathrm{Ee}\cdot \mathrm{alp}\cdot \mathrm{dt}}\]

\[\fitmath{\mathrm{stress}=-(1.74\times10^{11})\cdot (1.41\times10^{-5})\cdot (128)\approx -3.140352\times10^{8}\,\mathrm{Pa}}\]

\end{enumerate}


\noindent\textbf{Sustitución numérica}

\[
\fitmath{- \left(174000000000\,\mathrm{Pa}\right) \left(1.41e-05\,\mathrm{K}^{-1}\right) \left(128\,\mathrm{K}\right)}
\]

\noindent\textbf{Resultado}

\[
\fitmath{\sigma=-3.140352000\times10^{8}\,\mathrm{Pa}}
\]

\noindent\textbf{Interpretación}
El signo negativo indica compresión. Una barra libre tendría tensión térmica nula con la misma temperatura uniforme.

\noindent\textbf{Comprobación de los pasos}
\begin{itemize}
\item Paso 2: $0.0018048$ — reproducido
\item Paso 4: $-314035200$ — reproducido
\end{itemize}

\noindent\textbf{Comprobación de la respuesta} la aritmética del último paso reproduce la respuesta enunciada a la precisión con que está escrita. Es una reproducción aritmética; no valida el modelo físico ni sus hipótesis.

\noindent\textbf{Limitaciones}
\begin{itemize}
\item Modelo educativo; la aplicabilidad a una máquina real requiere revisión experta.
\end{itemize}

\noindent\textbf{Errores frecuentes}
\begin{itemize}
\item Aplicar esta tensión a toda pared con gradiente y curvatura.
\end{itemize}
% studium:end-record eyJpZCI6ICJQUkItMDExMSIsICJzZWN0aW9uIjogImMxNCJ9
\end{tcolorbox}

\begin{tcolorbox}[title={Ejercicios propuestos}, studiumworked, breakable]

\noindent\textbf{1 (Básico, Desarrollo)}
\(C14-01-P \cdot\) Red térmica de pared. Ejercicio de transferencia. Calcule flujo térmico y temperatura exterior del metal en una pared plana con convección a ambos lados. Antes de consultar la solución, prediga el efecto de cambiar el primer dato respecto del caso resuelto de esta familia. Justifique la tendencia y explique qué restricción podría convertirse en limitante. Datos: \(T_{aw}=1725.3\, \mathrm{K}\); \(T_c=745\, \mathrm{K}\); \(h_g=1680\, \mathrm{W}\,\mathrm{m}^{-2}\,\mathrm{K}^{-1}\); \(h_c=2450\, \mathrm{W}\,\mathrm{m}^{-2}\,\mathrm{K}^{-1}\); \(t=0.0011\, \mathrm{m}\); \(k_m=21\, \mathrm{W}\,\mathrm{m}^{-1}\,\mathrm{K}^{-1}\).

\noindent\textbf{2 (Intermedio, Desarrollo)}
\(C14-02-P \cdot\) Caudal energético de refrigeración. Ejercicio de transferencia. Una superficie recibe flujo térmico medio conocido. Calcule el caudal refrigerante que absorbería ese calor con el incremento permitido. Antes de consultar la solución, prediga el efecto de cambiar el primer dato respecto del caso resuelto de esta familia. Justifique la tendencia y explique qué restricción podría convertirse en limitante. Datos: \(q''=660300\, \mathrm{W}\,\mathrm{m}^{-2}\); \(A=0.034\, \mathrm{m}^{2}\); \(c_p=1090\, \mathrm{J}\,\mathrm{kg}^{-1}\,\mathrm{K}^{-1}\); \(\Delta T_c=185\, \mathrm{K}\).

\noindent\textbf{3 (Avanzado, Desarrollo)}
\(C14-03-P \cdot\) Película y recuperación. Ejercicio de transferencia. Con efectividad de película prescrita, calcule la recuperación efectiva y el flujo térmico a una pared de temperatura dada. Antes de consultar la solución, prediga el efecto de cambiar el primer dato respecto del caso resuelto de esta familia. Justifique la tendencia y explique qué restricción podría convertirse en limitante. Datos: \(T_{aw}=1799.85\, \mathrm{K}\); \(T_c=775\, \mathrm{K}\); \(\eta_f=0.41\); \(T_w=1130\, \mathrm{K}\); \(h_g=1770\, \mathrm{W}\,\mathrm{m}^{-2}\,\mathrm{K}^{-1}\).

\noindent\textbf{4 (Nivel de examen, Desarrollo)}
\(C14-04-P \cdot\) Restricción térmica uniaxial. Ejercicio de transferencia. Una barra homogénea tiene impedida completamente su dilatación longitudinal. Determine la deformación térmica libre y la tensión elástica de compresión. Antes de consultar la solución, prediga el efecto de cambiar el primer dato respecto del caso resuelto de esta familia. Justifique la tendencia y explique qué restricción podría convertirse en limitante. Datos: \(E=1.8531\times 10^{11}\, \mathrm{Pa}\); \(\alpha_{\mathrm{th}}=1.41\times 10^{-5}\, \mathrm{K}^{-1}\); \(\Delta T=128\, \mathrm{K}\).
\end{tcolorbox}


\studiumchaptercheck

\part{Compresores radiales}

\chapter{Compresores radiales: rodete, inductor y deslizamiento}
\label{chap:c15}

\studiumchapterguide{Objetivo: desarrollar compresores radiales: rodete, inductor y deslizamiento, declarar hipótesis y comprobar condiciones físicas.}
% studium:end-record eyJpZCI6ICJQQVItMDEwMSIsICJzZWN0aW9uIjogImMxNSJ9

\section{Introducción y objetivos}

Un compresor radial cambia de forma importante el radio de la corriente en el rotor. Esa variación permite un trabajo específico elevado, pero obliga a mantener la forma general de Euler y de la rotalpía. El objetivo es construir los triángulos del inductor y de la salida, distinguir geometría y deslizamiento y calcular trabajo, temperatura y presión. Pertenece al bloque V.
Las comparaciones de tamaño, caudal y rendimiento entre máquinas radiales y axiales dependen de aplicación y escala. Los valores orientativos de las diapositivas no se presentan como límites universales. La selección exige comparar requisitos de gasto, volumen disponible, relación de presión y régimen de giro.
% studium:end-record eyJpZCI6ICJQQVItMDEwMiIsICJzZWN0aW9uIjogImMxNSJ9

\section{Euler con radios distintos}

La entrada del rodete se denomina 1 y su salida 2. La dirección meridiana local se denota \(m\); en la salida puramente radial coincide con \(r\). El trabajo es

\[
\fitmath{w_c=U_2C_{\theta2}-U_1C_{\theta1},\qquad U_i=\Omega r_i.}
\]
Sin prerrotación, \(C_{\theta1}=0\). Si la salida ideal tiene álabes radiales y el gas los sigue, \(W_{\theta2}=0\), por lo que \(C_{\theta2,\infty}=U_2\) y \(w_{c,\infty}=U_2^2\). La igualdad depende de esas condiciones y del límite de guiado ideal.
Una prerrotación positiva reduce el trabajo para salida fijada, pero también modifica la velocidad relativa y el Mach del inductor. Una prerrotación contraria puede aumentar el trabajo y el Mach relativo. No es correcto asignarle un beneficio universal sin revisar ambos extremos.
% studium:end-record eyJpZCI6ICJQQVItMDEwMyIsICJzZWN0aW9uIjogImMxNSJ9

\section{Geometría de salida y convenio angular}

Se define un ángulo metálico \(\lambda_2\) medido desde la dirección meridiana hacia la tangente relativa positiva. En un álabe curvado hacia atrás respecto del giro, \(\lambda_2<0\), y el triángulo ideal satisface

\[
\fitmath{C_{\theta2,\infty}=U_2+C_{m2}\tan\lambda_2.}
\]
Otros textos miden un ángulo positivo desde la tangente y usan una cotangente. La conversión debe reconstruirse desde componentes; no se intercambia una cotangente por una tangente conservando el mismo número de grados.
La figura~\ref{fig:rodete} muestra un rodete radial esquemático original y las direcciones locales. La salida radial de la trayectoria no exige álabes rectos radiales: la curvatura del metal y la dirección meridiana son conceptos distintos.
% studium:end-record eyJpZCI6ICJQQVItMDEwNCIsICJzZWN0aW9uIjogImMxNSJ9

\par\vspace{\baselineskip}
\noindent\begin{minipage}{\linewidth}
\centering
\begingroup
\setbox0=\hbox{\begin{tikzpicture}
\draw[thick] (0,0) circle (2.1);\draw[thick] (0,0) circle (.65);\foreach \ang in {0,40,...,320}{\draw[blue,thick,rotate=\ang] (.65,0) .. controls (1.1,-.2) and (1.8,-.3) .. (2.1,-.3);}\draw[->,orange,thick] (1.1,2.1) arc (62:122:2.4);\node[above] at (0,2.6) {$\Omega>0$};\draw[->,thick] (2.1,0)--(3.8,0) node[right] {$C_{r2}$};\draw[->,thick] (2.1,0)--(2.1,1.4) node[above] {$C_{\theta2}$};
\end{tikzpicture}}%
\count0=\ht0
\advance\count0 by \dp0
\count2=\wd0
\count4=\dimexpr\textheight/100*38\relax
\count6=\linewidth
\ifnum\count0<1
\count0=1
\fi
\ifnum\count2<1
\count2=1
\fi
\divide\count0 by 4096
\divide\count2 by 4096
\divide\count4 by 4096
\divide\count6 by 4096
\ifnum\count0<1
\count0=1
\fi
\ifnum\count2<1
\count2=1
\fi
\ifnum\count4<1
\count4=1
\fi
\ifnum\count6<1
\count6=1
\fi
\count8=\count2
\ifnum\count0>\count4
\count8=\numexpr\count2*\count4/\count0\relax
\fi
\ifnum\count8>\count6
\resizebox{\linewidth}{!}{\begin{tikzpicture}
\draw[thick] (0,0) circle (2.1);\draw[thick] (0,0) circle (.65);\foreach \ang in {0,40,...,320}{\draw[blue,thick,rotate=\ang] (.65,0) .. controls (1.1,-.2) and (1.8,-.3) .. (2.1,-.3);}\draw[->,orange,thick] (1.1,2.1) arc (62:122:2.4);\node[above] at (0,2.6) {$\Omega>0$};\draw[->,thick] (2.1,0)--(3.8,0) node[right] {$C_{r2}$};\draw[->,thick] (2.1,0)--(2.1,1.4) node[above] {$C_{\theta2}$};
\end{tikzpicture}}%
\else
\ifnum\count0>\count4
\resizebox{!}{\dimexpr\textheight/100*38\relax}{\begin{tikzpicture}
\draw[thick] (0,0) circle (2.1);\draw[thick] (0,0) circle (.65);\foreach \ang in {0,40,...,320}{\draw[blue,thick,rotate=\ang] (.65,0) .. controls (1.1,-.2) and (1.8,-.3) .. (2.1,-.3);}\draw[->,orange,thick] (1.1,2.1) arc (62:122:2.4);\node[above] at (0,2.6) {$\Omega>0$};\draw[->,thick] (2.1,0)--(3.8,0) node[right] {$C_{r2}$};\draw[->,thick] (2.1,0)--(2.1,1.4) node[above] {$C_{\theta2}$};
\end{tikzpicture}}%
\else
\begin{tikzpicture}
\draw[thick] (0,0) circle (2.1);\draw[thick] (0,0) circle (.65);\foreach \ang in {0,40,...,320}{\draw[blue,thick,rotate=\ang] (.65,0) .. controls (1.1,-.2) and (1.8,-.3) .. (2.1,-.3);}\draw[->,orange,thick] (1.1,2.1) arc (62:122:2.4);\node[above] at (0,2.6) {$\Omega>0$};\draw[->,thick] (2.1,0)--(3.8,0) node[right] {$C_{r2}$};\draw[->,thick] (2.1,0)--(2.1,1.4) node[above] {$C_{\theta2}$};
\end{tikzpicture}%
\fi
\fi
\endgroup
\par\nopagebreak
\studiumcaption{fig:rodete}{Rodete radial esquemático original: giro, dirección radial y componente tangencial. Sin geometría de fabricante.}
\end{minipage}
\par\vspace{\baselineskip}
% studium:end-record eyJpZCI6ICJGSUctMDAxNSIsICJzZWN0aW9uIjogImMxNSJ9

\section{Deslizamiento y pérdidas}

Un número finito de álabes produce una salida que no coincide con el guiado ideal de infinitos álabes. Se define aquí

\[
\fitmath{\sigma_s=\frac{C_{\theta2}}{C_{\theta2,\infty}},\qquad
C_{\theta2}=\sigma_s C_{\theta2,\infty}.}
\]
Esta definición se utiliza solo cuando el denominador es positivo y no nulo. Algunas correlaciones definen en cambio \(1-\Delta C_\theta/U_2\); ambas coinciden para salida ideal radial, pero no para toda geometría curvada. El informe de cálculo debe declarar qué factor recibe.
El deslizamiento altera el remolino y el trabajo de Euler; no es una eficiencia isentrópica. No se multiplica por él una segunda vez si el \(C_{\theta2}\) dado ya es real. Las pérdidas viscosas determinan además qué presión se obtiene a partir del trabajo suministrado.
% studium:end-record eyJpZCI6ICJQQVItMDEwNSIsICJzZWN0aW9uIjogImMxNSJ9

\section{Inductor, Mach y rotalpía}

En la entrada, el radio varía entre cubo y punta. Sin prerrotación y con velocidad axial media uniforme,

\[
\fitmath{W_1(r)=\sqrt{C_{x1}^2+\Omega^2r^2}.}
\]
La punta del inductor puede ser limitante por Mach relativo aunque la velocidad absoluta sea moderada. La geometría del ojo y la incidencia deben evaluarse en varios radios. El radio de entrada no se sustituye por el radio de salida para calcular \(U_1\).
En el rotor adiabático,

\[
\fitmath{h_2-h_1=\frac{U_2^2-U_1^2}{2}-\frac{W_2^2-W_1^2}{2}.}
\]
Esta expresión muestra el término centrífugo. Puede aumentar la entalpía estática aunque el comportamiento de velocidad relativa no sea el de un rotor axial a radio constante. Suprimir el término \(U^2\) falsea el balance.
% studium:end-record eyJpZCI6ICJQQVItMDEwNiIsICJzZWN0aW9uIjogImMxNSJ9

\section{Cierre y límites}

Una vez calculado \(w_c\), \(T_{02}=T_{01}+w_c/c_p\). Una eficiencia total-total suministrada determina la relación de presión ideal equivalente mediante la misma definición de compresor. Se calculan luego temperatura estática, densidad y altura de salida. El difusor sigue siendo necesario para transformar parte de la gran velocidad absoluta de descarga en presión estática.
Errores habituales: omitir prerrotación no nula, confundir deslizamiento y eficiencia, mezclar las dos definiciones de factor y conservar entalpía total relativa cuando cambia el radio. Autocomprobación: compare un rodete radial ideal y uno con \(\sigma_s<1\), manteniendo velocidad periférica; explique qué disminuye y qué no queda determinado todavía.
Formulario: Euler general, \(C_{\theta2,\infty}=U_2+C_{m2}\tan\lambda_2\), \(w_c=\Delta(UC_\theta)\), rotalpía. Referencias: ULE Compresores radiales, pp. \(2-10\) y \(13-16\); MIT 16.50, lección 24, p. 4, para contraste de Euler radial. Los factores de los ejercicios son datos educativos prescritos, no correlaciones atribuidas a una fuente no consultada.
% studium:end-record eyJpZCI6ICJQQVItMDEwNyIsICJzZWN0aW9uIjogImMxNSJ9

\begin{tcolorbox}[title={Desarrollos del modelo}, studiumdef, breakable]

\noindent\textbf{Modelo desarrollado: Rodete radial con deslizamiento}

\begin{equation}
\Delta T_0=\frac{\sigma_s\,U^2}{c_p}
\label{eq:EC-C15}
\end{equation}

\noindent\textit{Significado.} Aumento de temperatura total con prerrotación nula. $\sigma_s$ es el factor de deslizamiento, no la solidez.

\noindent\textbf{Hipótesis}
\begin{itemize}
\item Sin prerrotación; salida ideal radial antes de aplicar slip.
\end{itemize}

\noindent\textbf{Principios}
\begin{itemize}
\item Balance y definiciones físicas desarrollados en el capítulo; propiedades y dominio declarados en el problema.
\end{itemize}

\noindent\textbf{Aplicabilidad}
\begin{itemize}
\item Modelo educativo del caso C15-01-R
\end{itemize}

\noindent\textbf{Limitaciones}
\begin{itemize}
\item Modelo educativo; la aplicabilidad a una máquina real requiere revisión experta.
\end{itemize}

\noindent\textbf{Pasos}
\begin{enumerate}
\item
Se reduce el remolino ideal con la definición indicada.

\[
\fitmath{C_\theta=\sigma_s U}
\]
\item
Euler usa el remolino real.

\[
\fitmath{w_c=\sigma_s U^{2}}
\]
\item
El balance energético fija el calentamiento total.

\[
\fitmath{\Delta T_0=\frac{\sigma_s U^{2}}{c_p}}
\]
\end{enumerate}

{\noindent\footnotesize verificada dimensionalmente.\par}
\end{tcolorbox}

\begin{tcolorbox}[title={Problema resuelto 1}, studiumworked, breakable]

\noindent\textbf{Enunciado}
\(C15-01-R \cdot\) Rodete radial con deslizamiento. Sin prerrotación y con álabes radiales ideales, calcule remolino real, trabajo y aumento total de temperatura. Datos: \(U=426\, \mathrm{m}\,\mathrm{s}^{-1}\); \(\sigma_s=0.87\); \(c_p=1004.5\, \mathrm{J}\,\mathrm{kg}^{-1}\,\mathrm{K}^{-1}\).

\noindent\textbf{Datos}
\begin{itemize}
\item $U = 426\,\mathrm{m}\,\mathrm{s}^{-1}$ — velocidad periférica salida
\item $\sigma_s = 0.87$ — factor definido como $C_\theta/C_{\theta,\mathrm{ideal}}$
\item $c_p = 1004.5\,\mathrm{J}\,\mathrm{kg}^{-1}\,\mathrm{K}^{-1}$ — calor específico
\end{itemize}

\noindent\textbf{Incógnita}
$\Delta T_0$, en $\mathrm{K}$; también las magnitudes intermedias desarrolladas.

\noindent\textbf{Modelo físico}

\noindent\textit{Rodete radial con deslizamiento}

\[
\fitmath{\Delta T_0=\frac{\sigma_s U^{2}}{c_p}}
\]

\noindent\textbf{Hipótesis}
\begin{itemize}
\item Sin prerrotación; salida ideal radial antes de aplicar slip.
\end{itemize}

\noindent\textbf{Desarrollo paso a paso}

\begin{enumerate}

\item Se reduce el remolino ideal con la definición indicada.

\[\fitmath{\mathrm{ct}=\mathrm{slip}\cdot \mathrm{uu}}\]

\[\fitmath{\mathrm{ct}=(0.87)\cdot (426)\approx 370.62\,\mathrm{m/s}}\]

\item Euler usa el remolino real.

\[\fitmath{\mathrm{wc}=\mathrm{slip}\cdot {\mathrm{uu}}^{2}}\]

\[\fitmath{\mathrm{wc}=(0.87)\cdot {(426)}^{2}\approx 157884.1\,\mathrm{J/kg}}\]

\item El balance energético fija el calentamiento total.

\[\fitmath{\mathrm{dT}=\frac{\mathrm{slip}\cdot {\mathrm{uu}}^{2}}{\mathrm{cp}}}\]

\[\fitmath{\mathrm{dT}=\frac{(0.87)\cdot {(426)}^{2}}{(1004.5)}\approx 157.1768\,\mathrm{K}}\]

\end{enumerate}


\noindent\textbf{Sustitución numérica}

\[
\fitmath{\frac{\left(0.87\right) \left(\left(426\,\mathrm{m}\,\mathrm{s}^{-1}\right)\right)^{2}}{\left(1004.5\,\mathrm{J}\,\mathrm{kg}^{-1}\,\mathrm{K}^{-1}\right)}}
\]

\noindent\textbf{Resultado}

\[
\fitmath{\Delta T_0=1.571768243\times10^{2}\,\mathrm{K}}
\]

\noindent\textbf{Interpretación}
El deslizamiento modifica trabajo. La presión necesita además una eficiencia o modelo de pérdida.

\noindent\textbf{Comprobación de los pasos}
\begin{itemize}
\item Paso 2: $370.62$ — reproducido
\item Paso 4: $157884$ — reproducido
\item Paso 6: $157.177$ — reproducido
\end{itemize}

\noindent\textbf{Comprobación de la respuesta} la aritmética del último paso reproduce la respuesta enunciada a la precisión con que está escrita. Es una reproducción aritmética; no valida el modelo físico ni sus hipótesis.

\noindent\textbf{Limitaciones}
\begin{itemize}
\item Modelo educativo; la aplicabilidad a una máquina real requiere revisión experta.
\end{itemize}

\noindent\textbf{Errores frecuentes}
\begin{itemize}
\item Multiplicar de nuevo por slip un trabajo que ya utiliza $C_\theta$ real.
\end{itemize}
% studium:end-record eyJpZCI6ICJQUkItMDExMyIsICJzZWN0aW9uIjogImMxNSJ9
\end{tcolorbox}

\begin{tcolorbox}[title={Problema resuelto 2}, studiumworked, breakable]

\noindent\textbf{Enunciado}
\(C15-02-R \cdot\) Curvatura hacia atrás. Un rodete tiene ángulo relativo desde la dirección meridiana negativo. Halle el trabajo con deslizamiento definido como cociente de remolinos. Datos: \(U=448\, \mathrm{m}\,\mathrm{s}^{-1}\); \(C_{m,s}=139\, \mathrm{m}\,\mathrm{s}^{-1}\); \(\lambda_s=-0.5061455\, \mathrm{rad}\); \(\sigma_s=0.89\).

\noindent\textbf{Datos}
\begin{itemize}
\item $U = 448\,\mathrm{m}\,\mathrm{s}^{-1}$ — velocidad periférica salida
\item $C_{m,s} = 139\,\mathrm{m}\,\mathrm{s}^{-1}$ — velocidad meridiana salida
\item $\lambda_s = -0.50614548307836\,\mathrm{rad}$ — ángulo relativo ideal desde meridiana
\item $\sigma_s = 0.89$ — factor de deslizamiento
\end{itemize}

\noindent\textbf{Incógnita}
$w_c$, en $\mathrm{J\,kg^{-1}}$; también las magnitudes intermedias desarrolladas.

\noindent\textbf{Modelo físico}

\noindent\textit{Curvatura hacia atrás}

\[
\fitmath{w_c=\sigma_s U \left(C_{m,s} \tan{\left(\lambda_s \right)} + U\right)}
\]

\noindent\textbf{Hipótesis}
\begin{itemize}
\item Factor prescrito positivo; salida ideal con remolino positivo.
\end{itemize}

\noindent\textbf{Desarrollo paso a paso}

\begin{enumerate}

\item El metal ideal determina la componente tangencial absoluta.

\[\fitmath{\mathrm{ctid}=\mathrm{uu}+\mathrm{cm}\cdot \tan\left(\mathrm{lam}\right)}\]

\[\fitmath{\mathrm{ctid}=(448)+(139)\cdot \tan\left((-0.5061454831)\right)\approx 370.951\,\mathrm{m/s}}\]

\item Se aplica el cociente de remolinos definido.

\[\fitmath{\mathrm{ct}=\mathrm{slip}\cdot \left(\mathrm{uu}+\mathrm{cm}\cdot \tan\left(\mathrm{lam}\right)\right)}\]

\[\fitmath{\mathrm{ct}=(0.89)\cdot \left((448)+(139)\cdot \tan\left((-0.5061454831)\right)\right)\approx 330.1464\,\mathrm{m/s}}\]

\item Sin prerrotación, Euler queda U2 Ctheta2.

\[\fitmath{\mathrm{wc}=\mathrm{uu}\cdot \mathrm{slip}\cdot \left(\mathrm{uu}+\mathrm{cm}\cdot \tan\left(\mathrm{lam}\right)\right)}\]

\[\fitmath{\mathrm{wc}=(448)\cdot (0.89)\cdot \left((448)+(139)\cdot \tan\left((-0.5061454831)\right)\right)\approx 147905.6\,\mathrm{J/kg}}\]

\end{enumerate}


\noindent\textbf{Sustitución numérica}

\[
\fitmath{\left(0.89\right) \left(448\,\mathrm{m}\,\mathrm{s}^{-1}\right) \left(\left(139\,\mathrm{m}\,\mathrm{s}^{-1}\right) \tan{\left(\left(-0.50614548307836\,\mathrm{rad}\right) \right)} + \left(448\,\mathrm{m}\,\mathrm{s}^{-1}\right)\right)}
\]

\noindent\textbf{Resultado}

\[
\fitmath{w_c=1.479055994\times10^{5}\,\mathrm{J}\,\mathrm{kg}^{-1}}
\]

\noindent\textbf{Interpretación}
La curvatura hacia atrás reduce el remolino ideal para estos datos. No intercambie esta definición con $1-dC_\theta/U$.

\noindent\textbf{Comprobación de los pasos}
\begin{itemize}
\item Paso 2: $370.951$ — reproducido
\item Paso 4: $330.146$ — reproducido
\item Paso 6: $147906$ — reproducido
\end{itemize}

\noindent\textbf{Comprobación de la respuesta} la aritmética del último paso reproduce la respuesta enunciada a la precisión con que está escrita. Es una reproducción aritmética; no valida el modelo físico ni sus hipótesis.

\noindent\textbf{Limitaciones}
\begin{itemize}
\item Modelo educativo; la aplicabilidad a una máquina real requiere revisión experta.
\end{itemize}

\noindent\textbf{Errores frecuentes}
\begin{itemize}
\item Usar una cotangente con el mismo ángulo desde meridiana.
\end{itemize}
% studium:end-record eyJpZCI6ICJQUkItMDExNSIsICJzZWN0aW9uIjogImMxNSJ9
\end{tcolorbox}

\begin{tcolorbox}[title={Problema resuelto 3}, studiumworked, breakable]

\noindent\textbf{Enunciado}
\(C15-03-R \cdot\) Inductor y Mach de punta. Un inductor sin prerrotación recibe gas con temperatura estática dada. Halle la velocidad relativa de punta y su Mach. Datos: \(\Omega=1920\, \mathrm{s}^{-1}\); \(r_t=0.102\, \mathrm{m}\); \(C_x=148\, \mathrm{m}\,\mathrm{s}^{-1}\); \(T_e=283\, \mathrm{K}\); \(\gamma=1.4\); \(R_g=287\, \mathrm{J}\,\mathrm{kg}^{-1}\,\mathrm{K}^{-1}\).

\noindent\textbf{Datos}
\begin{itemize}
\item $\Omega = 1920\,\mathrm{s}^{-1}$ — velocidad angular
\item $r_t = 0.102\,\mathrm{m}$ — radio punta inductor
\item $C_x = 148\,\mathrm{m}\,\mathrm{s}^{-1}$ — velocidad axial
\item $T_e = 283\,\mathrm{K}$ — temperatura estática
\item $\gamma = 1.4$ — razón de calores
\item $R_g = 287\,\mathrm{J}\,\mathrm{kg}^{-1}\,\mathrm{K}^{-1}$ — constante gas
\end{itemize}

\noindent\textbf{Incógnita}
$\mathrm{Ma}_r$ (adimensional); también las magnitudes intermedias desarrolladas.

\noindent\textbf{Modelo físico}

\noindent\textit{Inductor y Mach de punta}

\[
\fitmath{\mathrm{Ma}_r=\frac{\sqrt{\Omega^{2} r_t^{2} + C_x^{2}}}{\sqrt{R_g T_e \gamma}}}
\]

\noindent\textbf{Hipótesis}
\begin{itemize}
\item Sin prerrotación y velocidad axial uniforme.
\end{itemize}

\noindent\textbf{Desarrollo paso a paso}

\begin{enumerate}

\item La punta tiene mayor velocidad periférica que la raíz.

\[\fitmath{\mathrm{Ut}=\Omega\cdot \mathrm{rt}}\]

\[\fitmath{\mathrm{Ut}=(1920)\cdot (0.102)\approx 195.84\,\mathrm{m/s}}\]

\item Se forma el módulo relativo de entrada.

\[\fitmath{W=\sqrt{{\mathrm{cx}}^{2}+{\left(\Omega\cdot \mathrm{rt}\right)}^{2}}}\]

\[\fitmath{W=\sqrt{{(148)}^{2}+{\left((1920)\cdot (0.102)\right)}^{2}}\approx 245.4736\,\mathrm{m/s}}\]

\item El sonido corresponde al estado estático.

\[\fitmath{\mathrm{Mr}=\frac{\sqrt{{\mathrm{cx}}^{2}+{\left(\Omega\cdot \mathrm{rt}\right)}^{2}}}{\sqrt{\gamma\cdot R_g\cdot \mathrm{Te}}}}\]

\[\fitmath{\mathrm{Mr}=\frac{\sqrt{{(148)}^{2}+{\left((1920)\cdot (0.102)\right)}^{2}}}{\sqrt{(1.4)\cdot (287)\cdot (283)}}\approx 0.7279586\,}\]

\end{enumerate}


\noindent\textbf{Sustitución numérica}

\[
\fitmath{\frac{\sqrt{\left(\left(1920\,\mathrm{s}^{-1}\right)\right)^{2} \left(0.102\,\mathrm{m}\right)^{2} + \left(\left(148\,\mathrm{m}\,\mathrm{s}^{-1}\right)\right)^{2}}}{\sqrt{\left(287\,\mathrm{J}\,\mathrm{kg}^{-1}\,\mathrm{K}^{-1}\right) \left(283\,\mathrm{K}\right) \left(1.4\right)}}}
\]

\noindent\textbf{Resultado}

\[
\fitmath{\mathrm{Ma}_r=7.279585817\times10^{-1}}
\]

\noindent\textbf{Interpretación}
La temperatura estática se comparte entre referencias; la velocidad relativa aumenta por la rotación.

\noindent\textbf{Comprobación de los pasos}
\begin{itemize}
\item Paso 2: $195.84$ — reproducido
\item Paso 4: $245.474$ — reproducido
\item Paso 6: $0.727959$ — reproducido
\end{itemize}

\noindent\textbf{Comprobación de la respuesta} la aritmética del último paso reproduce la respuesta enunciada a la precisión con que está escrita. Es una reproducción aritmética; no valida el modelo físico ni sus hipótesis.

\noindent\textbf{Limitaciones}
\begin{itemize}
\item Modelo educativo; la aplicabilidad a una máquina real requiere revisión experta.
\end{itemize}

\noindent\textbf{Errores frecuentes}
\begin{itemize}
\item Usar radio exterior del rodete en el ojo del inductor.
\end{itemize}
% studium:end-record eyJpZCI6ICJQUkItMDExNyIsICJzZWN0aW9uIjogImMxNSJ9
\end{tcolorbox}

\begin{tcolorbox}[title={Problema resuelto 4}, studiumworked, breakable]

\noindent\textbf{Enunciado}
\(C15-04-R \cdot\) Prerrotación y trabajo. Manteniendo la descarga real, determine cuánto trabajo aporta el rodete con prerrotación positiva. Datos: \(\Omega=2080\, \mathrm{s}^{-1}\); \(r_e=0.079\, \mathrm{m}\); \(r_s=0.204\, \mathrm{m}\); \(C_{\theta e}=47\, \mathrm{m}\,\mathrm{s}^{-1}\); \(C_{\theta s}=351\, \mathrm{m}\,\mathrm{s}^{-1}\).

\noindent\textbf{Datos}
\begin{itemize}
\item $\Omega = 2080\,\mathrm{s}^{-1}$ — velocidad angular
\item $r_e = 0.079\,\mathrm{m}$ — radio entrada
\item $r_s = 0.204\,\mathrm{m}$ — radio salida
\item $C_{\theta e} = 47\,\mathrm{m}\,\mathrm{s}^{-1}$ — remolino entrada
\item $C_{\theta s} = 351\,\mathrm{m}\,\mathrm{s}^{-1}$ — remolino real salida
\end{itemize}

\noindent\textbf{Incógnita}
$w_c$, en $\mathrm{J\,kg^{-1}}$; también las magnitudes intermedias desarrolladas.

\noindent\textbf{Modelo físico}

\noindent\textit{Prerrotación y trabajo}

\[
\fitmath{w_c=\Omega \left(- C_{\theta e} r_e + C_{\theta s} r_s\right)}
\]

\noindent\textbf{Hipótesis}
\begin{itemize}
\item Euler general con valores medios compatibles.
\end{itemize}

\noindent\textbf{Desarrollo paso a paso}

\begin{enumerate}

\item Se calcula la contribución de descarga.

\[\fitmath{\mathrm{exit}=\Omega\cdot \mathrm{rs}\cdot \mathrm{cs}}\]

\[\fitmath{\mathrm{exit}=(2080)\cdot (0.204)\cdot (351)\approx 148936.3\,\mathrm{J/kg}}\]

\item Se resta la contribución de entrada.

\[\fitmath{\mathrm{wc}=\Omega\cdot \left(\mathrm{rs}\cdot \mathrm{cs}-\mathrm{re}\cdot \mathrm{ce}\right)}\]

\[\fitmath{\mathrm{wc}=(2080)\cdot \left((0.204)\cdot (351)-(0.079)\cdot (47)\right)\approx 141213.3\,\mathrm{J/kg}}\]

\end{enumerate}


\noindent\textbf{Sustitución numérica}

\[
\fitmath{\left(2080\,\mathrm{s}^{-1}\right) \left(- \left(47\,\mathrm{m}\,\mathrm{s}^{-1}\right) \left(0.079\,\mathrm{m}\right) + \left(351\,\mathrm{m}\,\mathrm{s}^{-1}\right) \left(0.204\,\mathrm{m}\right)\right)}
\]

\noindent\textbf{Resultado}

\[
\fitmath{w_c=1.412132800\times10^{5}\,\mathrm{J}\,\mathrm{kg}^{-1}}
\]

\noindent\textbf{Interpretación}
La prerrotación positiva reduce trabajo para descarga fijada. Su efecto en incidencia debe evaluarse aparte.

\noindent\textbf{Comprobación de los pasos}
\begin{itemize}
\item Paso 2: $148936$ — reproducido
\item Paso 4: $141213$ — reproducido
\end{itemize}

\noindent\textbf{Comprobación de la respuesta} la aritmética del último paso reproduce la respuesta enunciada a la precisión con que está escrita. Es una reproducción aritmética; no valida el modelo físico ni sus hipótesis.

\noindent\textbf{Limitaciones}
\begin{itemize}
\item Modelo educativo; la aplicabilidad a una máquina real requiere revisión experta.
\end{itemize}

\noindent\textbf{Errores frecuentes}
\begin{itemize}
\item Suponer prerrotación nula aunque el enunciado la proporcione.
\end{itemize}
% studium:end-record eyJpZCI6ICJQUkItMDExOSIsICJzZWN0aW9uIjogImMxNSJ9
\end{tcolorbox}

\begin{tcolorbox}[title={Ejercicios propuestos}, studiumworked, breakable]

\noindent\textbf{1 (Básico, Desarrollo)}
\(C15-01-P \cdot\) Rodete radial con deslizamiento. Ejercicio de transferencia. Sin prerrotación y con álabes radiales ideales, calcule remolino real, trabajo y aumento total de temperatura. Antes de consultar la solución, prediga el efecto de cambiar el primer dato respecto del caso resuelto de esta familia. Justifique la tendencia y explique qué restricción podría convertirse en limitante. Datos: \(U=453.69\, \mathrm{m}\,\mathrm{s}^{-1}\); \(\sigma_s=0.87\); \(c_p=1004.5\, \mathrm{J}\,\mathrm{kg}^{-1}\,\mathrm{K}^{-1}\).

\noindent\textbf{2 (Intermedio, Desarrollo)}
\(C15-02-P \cdot\) Curvatura hacia atrás. Ejercicio de transferencia. Un rodete tiene ángulo relativo desde la dirección meridiana negativo. Halle el trabajo con deslizamiento definido como cociente de remolinos. Antes de consultar la solución, prediga el efecto de cambiar el primer dato respecto del caso resuelto de esta familia. Justifique la tendencia y explique qué restricción podría convertirse en limitante. Datos: \(U=477.12\, \mathrm{m}\,\mathrm{s}^{-1}\); \(C_{m,s}=139\, \mathrm{m}\,\mathrm{s}^{-1}\); \(\lambda_s=-0.5061455\, \mathrm{rad}\); \(\sigma_s=0.89\).

\noindent\textbf{3 (Avanzado, Desarrollo)}
\(C15-03-P \cdot\) Inductor y Mach de punta. Ejercicio de transferencia. Un inductor sin prerrotación recibe gas con temperatura estática dada. Halle la velocidad relativa de punta y su Mach. Antes de consultar la solución, prediga el efecto de cambiar el primer dato respecto del caso resuelto de esta familia. Justifique la tendencia y explique qué restricción podría convertirse en limitante. Datos: \(\Omega=2044.8\, \mathrm{s}^{-1}\); \(r_t=0.102\, \mathrm{m}\); \(C_x=148\, \mathrm{m}\,\mathrm{s}^{-1}\); \(T_e=283\, \mathrm{K}\); \(\gamma=1.4\); \(R_g=287\, \mathrm{J}\,\mathrm{kg}^{-1}\,\mathrm{K}^{-1}\).

\noindent\textbf{4 (Nivel de examen, Desarrollo)}
\(C15-04-P \cdot\) Prerrotación y trabajo. Ejercicio de transferencia. Manteniendo la descarga real, determine cuánto trabajo aporta el rodete con prerrotación positiva. Antes de consultar la solución, prediga el efecto de cambiar el primer dato respecto del caso resuelto de esta familia. Justifique la tendencia y explique qué restricción podría convertirse en limitante. Datos: \(\Omega=2215.2\, \mathrm{s}^{-1}\); \(r_e=0.079\, \mathrm{m}\); \(r_s=0.204\, \mathrm{m}\); \(C_{\theta e}=47\, \mathrm{m}\,\mathrm{s}^{-1}\); \(C_{\theta s}=351\, \mathrm{m}\,\mathrm{s}^{-1}\).
\end{tcolorbox}


\studiumchaptercheck

\chapter{Difusores radiales, recuperación y geometría de paso}
\label{chap:c16}

\studiumchapterguide{Objetivo: desarrollar difusores radiales, recuperación y geometría de paso, declarar hipótesis y comprobar condiciones físicas.}
% studium:end-record eyJpZCI6ICJQQVItMDEwOCIsICJzZWN0aW9uIjogImMxNiJ9

\section{Introducción y objetivos}

La descarga del rodete contiene energía cinética que debe convertirse en presión estática antes de alimentar el componente siguiente. El difusor realiza esa transformación sin trabajo de eje. El objetivo es distinguir recuperación y pérdida, calcular el efecto de radio y ancho en una región sin álabes y dimensionar áreas sin confundir módulos y componentes. Pertenece al bloque V.
Los balances son medios y no resuelven el flujo no uniforme entre chorro y estela que abandona un rodete. Las predicciones de separación y de pérdida necesitan un modelo adicional, no contenido en la simple conservación de momento angular.
% studium:end-record eyJpZCI6ICJQQVItMDEwOSIsICJzZWN0aW9uIjogImMxNiJ9

\section{Difusor sin álabes}

Se denomina 2 a la entrada del difusor y 3 a su salida. Si no hay par externo apreciable,

\[
\fitmath{r_2C_{\theta2}=r_3C_{\theta3}.}
\]
La continuidad para ancho \(b(r)\), densidad media y factor de área disponible \(B\) es

\[
\fitmath{\dot m=2\pi r b(r)B(r)\rho(r)C_r(r).}
\]
Con densidad, ancho y bloqueo constantes, tanto \(C_r\) como \(C_\theta\) disminuyen como \(1/r\). El ángulo absoluto permanece constante y la trayectoria se aproxima a una espiral logarítmica. En gas compresible la densidad cambia y esa conclusión debe revisarse.
El trabajo es cero y, si el proceso es adiabático, \(h_{03}=h_{02}\). La reducción de \(C^2/2\) aumenta la entalpía estática. En un modelo ideal isentrópico se obtiene la presión con la relación de estados. En uno real hay además caída de presión total.
% studium:end-record eyJpZCI6ICJQQVItMDExMCIsICJzZWN0aW9uIjogImMxNiJ9

\section{Recuperación estática y pérdida total}


Para un modelo incompresible se define



\[
\fitmath{C_p=\frac{p_3-p_2}{\rho C_2^2/2},\qquad
\zeta=\frac{p_{02}-p_{03}}{\rho C_2^2/2}.}
\]
Con continuidad y energía mecánica compatible,

\[
\fitmath{C_p=1-\frac{C_3^2}{C_2^2}-\zeta.}
\]
La recuperación ideal corresponde a \(\zeta=0\). Una recuperación inferior a uno no es, por sí sola, la eficiencia del compresor. Puede quedar energía cinética a la salida y además existir pérdida. Si se calcula una \(\zeta<0\) en un difusor pasivo bajo estas hipótesis, se revisan datos y definiciones.
La figura~\ref{fig:difusor} muestra una curva calculada para un difusor sin álabes de ancho y densidad constantes. Se rotula como modelo ideal; no es una característica ensayada ni incluye separación. Cambiar el radio modifica la recuperación cinemática, pero aumentar indiscriminadamente longitud puede añadir fricción.
% studium:end-record eyJpZCI6ICJQQVItMDExMSIsICJzZWN0aW9uIjogImMxNiJ9

\par\vspace{\baselineskip}
\noindent\begin{minipage}{\linewidth}
\centering
\begingroup
\setbox0=\hbox{\begin{tikzpicture}
\draw[->,gray] (0,0)--(7,0);\draw[->,gray] (0,0)--(0,3.8);\draw (0.0,0)--(0.0,-.08) node[below,font=\scriptsize] {1};\draw (0,0.0)--(-.08,0.0) node[left,font=\scriptsize] {0};\draw (1.5,0)--(1.5,-.08) node[below,font=\scriptsize] {1.25};\draw (0,0.75)--(-.08,0.75) node[left,font=\scriptsize] {0.188};\draw (3.0,0)--(3.0,-.08) node[below,font=\scriptsize] {1.5};\draw (0,1.5)--(-.08,1.5) node[left,font=\scriptsize] {0.375};\draw (4.5,0)--(4.5,-.08) node[below,font=\scriptsize] {1.75};\draw (0,2.25)--(-.08,2.25) node[left,font=\scriptsize] {0.562};\draw (6.0,0)--(6.0,-.08) node[below,font=\scriptsize] {2};\draw (0,3.0)--(-.08,3.0) node[left,font=\scriptsize] {0.75};\draw[blue,very thick] plot coordinates {(0.00000,0.00000) (0.10000,0.13007) (0.20000,0.25390) (0.30000,0.37188) (0.40000,0.48438) (0.50000,0.59172) (0.60000,0.69421) (0.70000,0.79216) (0.80000,0.88581) (0.90000,0.97543) (1.00000,1.06122) (1.10000,1.14342) (1.20000,1.22222) (1.30000,1.29780) (1.40000,1.37034) (1.50000,1.44000) (1.60000,1.50693) (1.70000,1.57126) (1.80000,1.63314) (1.90000,1.69268) (2.00000,1.75000) (2.10000,1.80521) (2.20000,1.85842) (2.30000,1.90971) (2.40000,1.95918) (2.50000,2.00692) (2.60000,2.05300) (2.70000,2.09750) (2.80000,2.14050) (2.90000,2.18205) (3.00000,2.22222) (3.10000,2.26108) (3.20000,2.29868) (3.30000,2.33507) (3.40000,2.37030) (3.50000,2.40443) (3.60000,2.43750) (3.70000,2.46955) (3.80000,2.50062) (3.90000,2.53076) (4.00000,2.56000) (4.10000,2.58837) (4.20000,2.61592) (4.30000,2.64266) (4.40000,2.66864) (4.50000,2.69388) (4.60000,2.71841) (4.70000,2.74225) (4.80000,2.76543) (4.90000,2.78798) (5.00000,2.80992) (5.10000,2.83126) (5.20000,2.85204) (5.30000,2.87227) (5.40000,2.89197) (5.50000,2.91115) (5.60000,2.92985) (5.70000,2.94806) (5.80000,2.96581) (5.90000,2.98312) (6.00000,3.00000)};\node[below,align=center] at (3.5,-.65) {Razón $r_3/r_2$};\node[rotate=90] at (-1,2) {$C_p$};
\end{tikzpicture}}%
\count0=\ht0
\advance\count0 by \dp0
\count2=\wd0
\count4=\dimexpr\textheight/100*38\relax
\count6=\linewidth
\ifnum\count0<1
\count0=1
\fi
\ifnum\count2<1
\count2=1
\fi
\divide\count0 by 4096
\divide\count2 by 4096
\divide\count4 by 4096
\divide\count6 by 4096
\ifnum\count0<1
\count0=1
\fi
\ifnum\count2<1
\count2=1
\fi
\ifnum\count4<1
\count4=1
\fi
\ifnum\count6<1
\count6=1
\fi
\count8=\count2
\ifnum\count0>\count4
\count8=\numexpr\count2*\count4/\count0\relax
\fi
\ifnum\count8>\count6
\resizebox{\linewidth}{!}{\begin{tikzpicture}
\draw[->,gray] (0,0)--(7,0);\draw[->,gray] (0,0)--(0,3.8);\draw (0.0,0)--(0.0,-.08) node[below,font=\scriptsize] {1};\draw (0,0.0)--(-.08,0.0) node[left,font=\scriptsize] {0};\draw (1.5,0)--(1.5,-.08) node[below,font=\scriptsize] {1.25};\draw (0,0.75)--(-.08,0.75) node[left,font=\scriptsize] {0.188};\draw (3.0,0)--(3.0,-.08) node[below,font=\scriptsize] {1.5};\draw (0,1.5)--(-.08,1.5) node[left,font=\scriptsize] {0.375};\draw (4.5,0)--(4.5,-.08) node[below,font=\scriptsize] {1.75};\draw (0,2.25)--(-.08,2.25) node[left,font=\scriptsize] {0.562};\draw (6.0,0)--(6.0,-.08) node[below,font=\scriptsize] {2};\draw (0,3.0)--(-.08,3.0) node[left,font=\scriptsize] {0.75};\draw[blue,very thick] plot coordinates {(0.00000,0.00000) (0.10000,0.13007) (0.20000,0.25390) (0.30000,0.37188) (0.40000,0.48438) (0.50000,0.59172) (0.60000,0.69421) (0.70000,0.79216) (0.80000,0.88581) (0.90000,0.97543) (1.00000,1.06122) (1.10000,1.14342) (1.20000,1.22222) (1.30000,1.29780) (1.40000,1.37034) (1.50000,1.44000) (1.60000,1.50693) (1.70000,1.57126) (1.80000,1.63314) (1.90000,1.69268) (2.00000,1.75000) (2.10000,1.80521) (2.20000,1.85842) (2.30000,1.90971) (2.40000,1.95918) (2.50000,2.00692) (2.60000,2.05300) (2.70000,2.09750) (2.80000,2.14050) (2.90000,2.18205) (3.00000,2.22222) (3.10000,2.26108) (3.20000,2.29868) (3.30000,2.33507) (3.40000,2.37030) (3.50000,2.40443) (3.60000,2.43750) (3.70000,2.46955) (3.80000,2.50062) (3.90000,2.53076) (4.00000,2.56000) (4.10000,2.58837) (4.20000,2.61592) (4.30000,2.64266) (4.40000,2.66864) (4.50000,2.69388) (4.60000,2.71841) (4.70000,2.74225) (4.80000,2.76543) (4.90000,2.78798) (5.00000,2.80992) (5.10000,2.83126) (5.20000,2.85204) (5.30000,2.87227) (5.40000,2.89197) (5.50000,2.91115) (5.60000,2.92985) (5.70000,2.94806) (5.80000,2.96581) (5.90000,2.98312) (6.00000,3.00000)};\node[below,align=center] at (3.5,-.65) {Razón $r_3/r_2$};\node[rotate=90] at (-1,2) {$C_p$};
\end{tikzpicture}}%
\else
\ifnum\count0>\count4
\resizebox{!}{\dimexpr\textheight/100*38\relax}{\begin{tikzpicture}
\draw[->,gray] (0,0)--(7,0);\draw[->,gray] (0,0)--(0,3.8);\draw (0.0,0)--(0.0,-.08) node[below,font=\scriptsize] {1};\draw (0,0.0)--(-.08,0.0) node[left,font=\scriptsize] {0};\draw (1.5,0)--(1.5,-.08) node[below,font=\scriptsize] {1.25};\draw (0,0.75)--(-.08,0.75) node[left,font=\scriptsize] {0.188};\draw (3.0,0)--(3.0,-.08) node[below,font=\scriptsize] {1.5};\draw (0,1.5)--(-.08,1.5) node[left,font=\scriptsize] {0.375};\draw (4.5,0)--(4.5,-.08) node[below,font=\scriptsize] {1.75};\draw (0,2.25)--(-.08,2.25) node[left,font=\scriptsize] {0.562};\draw (6.0,0)--(6.0,-.08) node[below,font=\scriptsize] {2};\draw (0,3.0)--(-.08,3.0) node[left,font=\scriptsize] {0.75};\draw[blue,very thick] plot coordinates {(0.00000,0.00000) (0.10000,0.13007) (0.20000,0.25390) (0.30000,0.37188) (0.40000,0.48438) (0.50000,0.59172) (0.60000,0.69421) (0.70000,0.79216) (0.80000,0.88581) (0.90000,0.97543) (1.00000,1.06122) (1.10000,1.14342) (1.20000,1.22222) (1.30000,1.29780) (1.40000,1.37034) (1.50000,1.44000) (1.60000,1.50693) (1.70000,1.57126) (1.80000,1.63314) (1.90000,1.69268) (2.00000,1.75000) (2.10000,1.80521) (2.20000,1.85842) (2.30000,1.90971) (2.40000,1.95918) (2.50000,2.00692) (2.60000,2.05300) (2.70000,2.09750) (2.80000,2.14050) (2.90000,2.18205) (3.00000,2.22222) (3.10000,2.26108) (3.20000,2.29868) (3.30000,2.33507) (3.40000,2.37030) (3.50000,2.40443) (3.60000,2.43750) (3.70000,2.46955) (3.80000,2.50062) (3.90000,2.53076) (4.00000,2.56000) (4.10000,2.58837) (4.20000,2.61592) (4.30000,2.64266) (4.40000,2.66864) (4.50000,2.69388) (4.60000,2.71841) (4.70000,2.74225) (4.80000,2.76543) (4.90000,2.78798) (5.00000,2.80992) (5.10000,2.83126) (5.20000,2.85204) (5.30000,2.87227) (5.40000,2.89197) (5.50000,2.91115) (5.60000,2.92985) (5.70000,2.94806) (5.80000,2.96581) (5.90000,2.98312) (6.00000,3.00000)};\node[below,align=center] at (3.5,-.65) {Razón $r_3/r_2$};\node[rotate=90] at (-1,2) {$C_p$};
\end{tikzpicture}}%
\else
\begin{tikzpicture}
\draw[->,gray] (0,0)--(7,0);\draw[->,gray] (0,0)--(0,3.8);\draw (0.0,0)--(0.0,-.08) node[below,font=\scriptsize] {1};\draw (0,0.0)--(-.08,0.0) node[left,font=\scriptsize] {0};\draw (1.5,0)--(1.5,-.08) node[below,font=\scriptsize] {1.25};\draw (0,0.75)--(-.08,0.75) node[left,font=\scriptsize] {0.188};\draw (3.0,0)--(3.0,-.08) node[below,font=\scriptsize] {1.5};\draw (0,1.5)--(-.08,1.5) node[left,font=\scriptsize] {0.375};\draw (4.5,0)--(4.5,-.08) node[below,font=\scriptsize] {1.75};\draw (0,2.25)--(-.08,2.25) node[left,font=\scriptsize] {0.562};\draw (6.0,0)--(6.0,-.08) node[below,font=\scriptsize] {2};\draw (0,3.0)--(-.08,3.0) node[left,font=\scriptsize] {0.75};\draw[blue,very thick] plot coordinates {(0.00000,0.00000) (0.10000,0.13007) (0.20000,0.25390) (0.30000,0.37188) (0.40000,0.48438) (0.50000,0.59172) (0.60000,0.69421) (0.70000,0.79216) (0.80000,0.88581) (0.90000,0.97543) (1.00000,1.06122) (1.10000,1.14342) (1.20000,1.22222) (1.30000,1.29780) (1.40000,1.37034) (1.50000,1.44000) (1.60000,1.50693) (1.70000,1.57126) (1.80000,1.63314) (1.90000,1.69268) (2.00000,1.75000) (2.10000,1.80521) (2.20000,1.85842) (2.30000,1.90971) (2.40000,1.95918) (2.50000,2.00692) (2.60000,2.05300) (2.70000,2.09750) (2.80000,2.14050) (2.90000,2.18205) (3.00000,2.22222) (3.10000,2.26108) (3.20000,2.29868) (3.30000,2.33507) (3.40000,2.37030) (3.50000,2.40443) (3.60000,2.43750) (3.70000,2.46955) (3.80000,2.50062) (3.90000,2.53076) (4.00000,2.56000) (4.10000,2.58837) (4.20000,2.61592) (4.30000,2.64266) (4.40000,2.66864) (4.50000,2.69388) (4.60000,2.71841) (4.70000,2.74225) (4.80000,2.76543) (4.90000,2.78798) (5.00000,2.80992) (5.10000,2.83126) (5.20000,2.85204) (5.30000,2.87227) (5.40000,2.89197) (5.50000,2.91115) (5.60000,2.92985) (5.70000,2.94806) (5.80000,2.96581) (5.90000,2.98312) (6.00000,3.00000)};\node[below,align=center] at (3.5,-.65) {Razón $r_3/r_2$};\node[rotate=90] at (-1,2) {$C_p$};
\end{tikzpicture}%
\fi
\fi
\endgroup
\par\nopagebreak
\studiumcaption{fig:difusor}{Recuperación calculada \(C_p=1-(r2/r3)^{2}\) a ancho y densidad constantes. Sin fricción ni separación.}
\end{minipage}
\par\vspace{\baselineskip}
% studium:end-record eyJpZCI6ICJGSUctMDAxNiIsICJzZWN0aW9uIjogImMxNiJ9

\section{Difusor con álabes y garganta}

Una cascada difusora guía y desacelera la corriente, permitiendo mayor recuperación en menor diámetro bajo condiciones adecuadas. Introduce sensibilidad a incidencia, bloqueo y posibles ondas de choque. Su garganta debe dimensionarse con velocidad normal y área efectiva.
El número de álabes del difusor no tiene por qué coincidir con el del rodete. La interacción no estacionaria y la elección de espaciamientos requieren un análisis adicional. Una vista bidimensional no permite decidir por sí sola ruido, excitación o fatiga.
El giro que necesita el flujo puede estimarse con un balance tangencial y una geometría preliminar. Convertirlo en ángulo metálico exige incidencia y desviación; el paso del capítulo de cascadas sigue siendo necesario. La solución de un triángulo no garantiza que el canal admita el caudal sin separación.
% studium:end-record eyJpZCI6ICJQQVItMDExMiIsICJzZWN0aW9uIjogImMxNiJ9

\section{Ancho de salida y acoplamiento al rodete}


Dado el caudal, el radio, la densidad y la velocidad radial deseada,



\[
\fitmath{b_2=\frac{\dot m}{2\pi r_2\rho_2C_{r2}B_2}.}
\]
La temperatura estática usada para \(\rho_2\) debe derivarse de la total con el módulo completo \(C_2\), mientras la continuidad utiliza \(C_{r2}\). Este es un caso frecuente en el que dos velocidades distintas intervienen correctamente en dos ecuaciones del mismo corte.
Si la altura calculada es muy pequeña, la fricción de paredes, tolerancias y bloqueo pueden dominar. Si es muy grande, cambian geometría y distribución del flujo. La estimación media debe revisarse con criterios de esbeltez, difusión y compatibilidad con la carcasa.
% studium:end-record eyJpZCI6ICJQQVItMDExMyIsICJzZWN0aW9uIjogImMxNiJ9

\section{Errores y formulario}

Errores habituales: conservar \(C_\theta\) en lugar de \(rC_\theta\), utilizar \(C\) en lugar de \(C_r\) para el área cilíndrica y confundir presión total constante ideal con presión estática constante. Autocomprobación: en el modelo ideal de ancho constante, duplique el radio y deduzca la reducción de energía cinética.
Formulario: \(rC_\theta=\mathrm{constante}\), \(\dot m=2\pi rbB\rho C_r\), \(h_0=\mathrm{constante}\) y relación entre \(C_p\) y \(\zeta\). Referencias: ULE Compresores radiales, pp. 4, 7 y \(11-12\); el principio de ausencia de trabajo en fila fija se contrasta con MIT 16.50, lección 24. La integración de los modelos y todos los gráficos son originales; el modelo de pérdida se prescribe en cada ejercicio.
% studium:end-record eyJpZCI6ICJQQVItMDExNCIsICJzZWN0aW9uIjogImMxNiJ9

\begin{tcolorbox}[title={Desarrollos del modelo}, studiumdef, breakable]

\noindent\textbf{Modelo desarrollado: Difusor sin álabes ideal}

\begin{equation}
\Delta p=\frac{\rho\,(C_r^2+C_\theta^2)}{2}\left(1-\frac{r_e^2}{r_s^2}\right)
\label{eq:EC-C16}
\end{equation}

\noindent\textit{Significado.} Diferencia de presión entre entrada y salida en un flujo radial incompresible, de ancho constante y sin par.

\noindent\textbf{Hipótesis}
\begin{itemize}
\item Adiabático ideal, ancho y densidad constantes, sin par.
\end{itemize}

\noindent\textbf{Principios}
\begin{itemize}
\item Balance y definiciones físicas desarrollados en el capítulo; propiedades y dominio declarados en el problema.
\end{itemize}

\noindent\textbf{Aplicabilidad}
\begin{itemize}
\item Modelo educativo del caso C16-01-R
\end{itemize}

\noindent\textbf{Limitaciones}
\begin{itemize}
\item Modelo educativo; la aplicabilidad a una máquina real requiere revisión experta.
\end{itemize}

\noindent\textbf{Pasos}
\begin{enumerate}
\item
Continuidad y ausencia de par hacen escalar ambas componentes con r inverso.

\[
\fitmath{C_{\theta s}=\frac{C_\theta r_e}{r_s}}
\]
\item
Se evalúa el incremento de presión por deceleración ideal.

\[
\fitmath{\Delta p=\frac{\rho \left(C_r^{2} + C_\theta^{2}\right) \left(- \frac{r_e^{2}}{r_s^{2}} + 1\right)}{2}}
\]
\end{enumerate}

{\noindent\footnotesize verificada dimensionalmente.\par}
\end{tcolorbox}

\begin{tcolorbox}[title={Problema resuelto 1}, studiumworked, breakable]

\noindent\textbf{Enunciado}
\(C16-01-R \cdot\) Difusor sin álabes ideal. En un difusor de ancho y densidad constantes, calcule las componentes de salida y el incremento ideal de presión. Datos: \(r_e=0.21\, \mathrm{m}\); \(r_s=0.32\, \mathrm{m}\); \(C_r=128\, \mathrm{m}\,\mathrm{s}^{-1}\); \(C_\theta=337\, \mathrm{m}\,\mathrm{s}^{-1}\); \(\rho=2.9\, \mathrm{kg}\,\mathrm{m}^{-3}\).

\noindent\textbf{Datos}
\begin{itemize}
\item $r_e = 0.21\,\mathrm{m}$ — radio entrada
\item $r_s = 0.32\,\mathrm{m}$ — radio salida
\item $C_r = 128\,\mathrm{m}\,\mathrm{s}^{-1}$ — componente radial entrada
\item $C_\theta = 337\,\mathrm{m}\,\mathrm{s}^{-1}$ — componente tangencial entrada
\item $\rho = 2.9\,\mathrm{kg}\,\mathrm{m}^{-3}$ — densidad uniforme
\end{itemize}

\noindent\textbf{Incógnita}
$\Delta p$, en $\mathrm{Pa}$; también las magnitudes intermedias desarrolladas.

\noindent\textbf{Modelo físico}

\noindent\textit{Difusor sin álabes ideal}

\[
\fitmath{\Delta p=\frac{\rho \left(C_r^{2} + C_\theta^{2}\right) \left(- \frac{r_e^{2}}{r_s^{2}} + 1\right)}{2}}
\]

\noindent\textbf{Hipótesis}
\begin{itemize}
\item Adiabático ideal, ancho y densidad constantes, sin par.
\end{itemize}

\noindent\textbf{Desarrollo paso a paso}

\begin{enumerate}

\item Componente radial de salida

\[\fitmath{\mathrm{crs}=\frac{\mathrm{cr}\cdot \mathrm{re}}{\mathrm{rs}}}\]

\[\fitmath{\mathrm{crs}=\frac{(128)\cdot (0.21)}{(0.32)}\approx 84\,\mathrm{m/s}}\]

\item Continuidad y ausencia de par hacen escalar ambas componentes con r inverso.

\[\fitmath{\mathrm{cts}=\frac{\mathrm{ct}\cdot \mathrm{re}}{\mathrm{rs}}}\]

\[\fitmath{\mathrm{cts}=\frac{(337)\cdot (0.21)}{(0.32)}\approx 221.1562\,\mathrm{m/s}}\]

\item Se evalúa el incremento de presión por deceleración ideal.

\[\fitmath{\mathrm{dp}=\frac{\rho\cdot \left({\mathrm{cr}}^{2}+{\mathrm{ct}}^{2}\right)\cdot \left(1-{\left(\frac{\mathrm{re}}{\mathrm{rs}}\right)}^{2}\right)}{2}}\]

\[\fitmath{\mathrm{dp}=\frac{(2.9)\cdot \left({(128)}^{2}+{(337)}^{2}\right)\cdot \left(1-{\left(\frac{(0.21)}{(0.32)}\right)}^{2}\right)}{2}\approx 107281\,\mathrm{Pa}}\]

\end{enumerate}


\noindent\textbf{Sustitución numérica}

\[
\fitmath{\frac{\left(2.9\,\mathrm{kg}\,\mathrm{m}^{-3}\right) \left(\left(\left(128\,\mathrm{m}\,\mathrm{s}^{-1}\right)\right)^{2} + \left(\left(337\,\mathrm{m}\,\mathrm{s}^{-1}\right)\right)^{2}\right) \left(- \frac{\left(0.21\,\mathrm{m}\right)^{2}}{\left(0.32\,\mathrm{m}\right)^{2}} + 1\right)}{2}}
\]

\noindent\textbf{Resultado}

\[
\fitmath{\Delta p=1.072810240\times10^{5}\,\mathrm{Pa}}
\]

\noindent\textbf{Componente radial de salida}
\[
\fitmath{C_{rs}=C_{re}r_e/r_s=84\,\mathrm{m/s}}
\]

\noindent\textbf{Comprobación de continuidad.} $r_sC_{rs}=r_eC_{re}$ con ancho y densidad constantes; junto con $r_sC_{\theta s}=r_eC_{\theta e}$ reproduce la recuperación ideal de presión.

\noindent\textbf{Interpretación}
La presión aumenta y el momento angular se conserva. Es una integración incompresible ideal, no una curva ensayada.

\noindent\textbf{Comprobación de los pasos}
\begin{itemize}
\item Paso 2: $221.156$ — reproducido
\item Paso 4: $107281$ — reproducido
\end{itemize}

\noindent\textbf{Comprobación de la respuesta} la aritmética del último paso reproduce la respuesta enunciada a la precisión con que está escrita. Es una reproducción aritmética; no valida el modelo físico ni sus hipótesis.

\noindent\textbf{Limitaciones}
\begin{itemize}
\item Modelo educativo; la aplicabilidad a una máquina real requiere revisión experta.
\end{itemize}

\noindent\textbf{Errores frecuentes}
\begin{itemize}
\item Conservar $C_\theta$ en lugar de $rC_\theta$.
\end{itemize}
% studium:end-record eyJpZCI6ICJQUkItMDEyMSIsICJzZWN0aW9uIjogImMxNiJ9
\end{tcolorbox}

\begin{tcolorbox}[title={Problema resuelto 2}, studiumworked, breakable]

\noindent\textbf{Enunciado}
\(C16-02-R \cdot\) Recuperación con pérdida prescrita. Calcule el coeficiente de recuperación estática y el incremento de presión con una pérdida normalizada dada. Datos: \(C_e=364\, \mathrm{m}\,\mathrm{s}^{-1}\); \(C_s=216\, \mathrm{m}\,\mathrm{s}^{-1}\); \(\zeta_D=0.13\); \(\rho=3.1\, \mathrm{kg}\,\mathrm{m}^{-3}\).

\noindent\textbf{Datos}
\begin{itemize}
\item $C_e = 364\,\mathrm{m}\,\mathrm{s}^{-1}$ — velocidad entrada
\item $C_s = 216\,\mathrm{m}\,\mathrm{s}^{-1}$ — velocidad salida
\item $\zeta_D = 0.13$ — pérdida normalizada con dinámica de entrada
\item $\rho = 3.1\,\mathrm{kg}\,\mathrm{m}^{-3}$ — densidad
\end{itemize}

\noindent\textbf{Incógnita}
$\Delta p$, en $\mathrm{Pa}$; también las magnitudes intermedias desarrolladas.

\noindent\textbf{Modelo físico}

\noindent\textit{Recuperación con pérdida prescrita}

\[
\fitmath{\Delta p=\frac{\rho C_e^{2} \left(- \zeta_D + 1 - \frac{C_s^{2}}{C_e^{2}}\right)}{2}}
\]

\noindent\textbf{Hipótesis}
\begin{itemize}
\item Modelo incompresible, pérdida definida con velocidad entrada.
\end{itemize}

\noindent\textbf{Desarrollo paso a paso}

\begin{enumerate}

\item La recuperación resta energía residual y pérdida.

\[\fitmath{\mathrm{Cp}=1-{\left(\frac{\mathrm{vs}}{\mathrm{ve}}\right)}^{2}-\mathrm{loss}}\]

\[\fitmath{\mathrm{Cp}=1-{\left(\frac{(216)}{(364)}\right)}^{2}-(0.13)\approx 0.5178686\,}\]

\item Se dimensionaliza con la presión dinámica de entrada.

\[\fitmath{\mathrm{dp}=\frac{\rho\cdot {\mathrm{ve}}^{2}\cdot \left(1-{\left(\frac{\mathrm{vs}}{\mathrm{ve}}\right)}^{2}-\mathrm{loss}\right)}{2}}\]

\[\fitmath{\mathrm{dp}=\frac{(3.1)\cdot {(364)}^{2}\cdot \left(1-{\left(\frac{(216)}{(364)}\right)}^{2}-(0.13)\right)}{2}\approx 106354.1\,\mathrm{Pa}}\]

\end{enumerate}


\noindent\textbf{Sustitución numérica}

\[
\fitmath{\frac{\left(3.1\,\mathrm{kg}\,\mathrm{m}^{-3}\right) \left(\left(364\,\mathrm{m}\,\mathrm{s}^{-1}\right)\right)^{2} \left(- \left(0.13\right) + 1 - \frac{\left(\left(216\,\mathrm{m}\,\mathrm{s}^{-1}\right)\right)^{2}}{\left(\left(364\,\mathrm{m}\,\mathrm{s}^{-1}\right)\right)^{2}}\right)}{2}}
\]

\noindent\textbf{Resultado}

\[
\fitmath{\Delta p=1.063540560\times10^{5}\,\mathrm{Pa}}
\]

\noindent\textbf{Interpretación}
La pérdida total y la velocidad residual reducen recuperación por mecanismos distintos.

\noindent\textbf{Comprobación de los pasos}
\begin{itemize}
\item Paso 2: $0.517869$ — reproducido
\item Paso 4: $106354$ — reproducido
\end{itemize}

\noindent\textbf{Comprobación de la respuesta} la aritmética del último paso reproduce la respuesta enunciada a la precisión con que está escrita. Es una reproducción aritmética; no valida el modelo físico ni sus hipótesis.

\noindent\textbf{Limitaciones}
\begin{itemize}
\item Modelo educativo; la aplicabilidad a una máquina real requiere revisión experta.
\end{itemize}

\noindent\textbf{Errores frecuentes}
\begin{itemize}
\item Restar una pérdida normalizada con otro denominador.
\end{itemize}
% studium:end-record eyJpZCI6ICJQUkItMDEyMyIsICJzZWN0aW9uIjogImMxNiJ9
\end{tcolorbox}

\begin{tcolorbox}[title={Problema resuelto 3}, studiumworked, breakable]

\noindent\textbf{Enunciado}
\(C16-03-R \cdot\) Ancho del paso cilíndrico. Dimensione el ancho geométrico de salida radial con el caudal y área disponible dados. Datos: \(\dot m=8.7\, \mathrm{kg}\,\mathrm{s}^{-1}\); \(r=0.198\, \mathrm{m}\); \(\rho=3.4\, \mathrm{kg}\,\mathrm{m}^{-3}\); \(C_r=126\, \mathrm{m}\,\mathrm{s}^{-1}\); \(B=0.94\).

\noindent\textbf{Datos}
\begin{itemize}
\item $\dot m = 8.7\,\mathrm{kg}\,\mathrm{s}^{-1}$ — caudal
\item $r = 0.198\,\mathrm{m}$ — radio salida
\item $\rho = 3.4\,\mathrm{kg}\,\mathrm{m}^{-3}$ — densidad estática
\item $C_r = 126\,\mathrm{m}\,\mathrm{s}^{-1}$ — velocidad radial
\item $B = 0.94$ — fracción área disponible
\end{itemize}

\noindent\textbf{Incógnita}
$b$, en $\mathrm{m}$; también las magnitudes intermedias desarrolladas.

\noindent\textbf{Modelo físico}

\noindent\textit{Ancho del paso cilíndrico}

\[
\fitmath{b=\frac{\dot m}{2 \pi B C_r r \rho}}
\]

\noindent\textbf{Hipótesis}
\begin{itemize}
\item Paso cilíndrico uniforme y bloqueo explícito.
\end{itemize}

\noindent\textbf{Desarrollo paso a paso}

\begin{enumerate}

\item La continuidad fija el área cilíndrica efectiva.

\[\fitmath{\mathrm{ae}=\frac{\mathrm{md}}{\rho\cdot \mathrm{cr}}}\]

\[\fitmath{\mathrm{ae}=\frac{(8.7)}{(3.4)\cdot (126)}\approx 0.02030812\,\mathrm{m^{2}}}\]

\item La circunferencia y el bloqueo fijan el ancho.

\[\fitmath{\mathrm{width}=\frac{\mathrm{md}}{2\cdot \pi\cdot \mathrm{rad}\cdot \rho\cdot \mathrm{cr}\cdot \mathrm{bb}}}\]

\[\fitmath{\mathrm{width}=\frac{(8.7)}{2\cdot \pi\cdot (0.198)\cdot (3.4)\cdot (126)\cdot (0.94)}\approx 0.01736588\,\mathrm{m}}\]

\end{enumerate}


\noindent\textbf{Sustitución numérica}

\[
\fitmath{\frac{\left(8.7\,\mathrm{kg}\,\mathrm{s}^{-1}\right)}{2 \pi \left(0.94\right) \left(126\,\mathrm{m}\,\mathrm{s}^{-1}\right) \left(0.198\,\mathrm{m}\right) \left(3.4\,\mathrm{kg}\,\mathrm{m}^{-3}\right)}}
\]

\noindent\textbf{Resultado}

\[
\fitmath{b=1.736588330\times10^{-2}\,\mathrm{m}}
\]

\noindent\textbf{Interpretación}
Se usa componente radial, aunque el módulo absoluto sea mucho mayor por remolino.

\noindent\textbf{Comprobación de los pasos}
\begin{itemize}
\item Paso 2: $0.0203081$ — reproducido
\item Paso 4: $0.0173659$ — reproducido
\end{itemize}

\noindent\textbf{Comprobación de la respuesta} la aritmética del último paso reproduce la respuesta enunciada a la precisión con que está escrita. Es una reproducción aritmética; no valida el modelo físico ni sus hipótesis.

\noindent\textbf{Limitaciones}
\begin{itemize}
\item Modelo educativo; la aplicabilidad a una máquina real requiere revisión experta.
\end{itemize}

\noindent\textbf{Errores frecuentes}
\begin{itemize}
\item Usar el módulo absoluto en la continuidad radial.
\end{itemize}
% studium:end-record eyJpZCI6ICJQUkItMDEyNSIsICJzZWN0aW9uIjogImMxNiJ9
\end{tcolorbox}

\begin{tcolorbox}[title={Problema resuelto 4}, studiumworked, breakable]

\noindent\textbf{Enunciado}
\(C16-04-R \cdot\) Estado estático de descarga. La descarga del rodete tiene estado total y componentes radial y tangencial dadas. Halle temperatura estática y densidad para dimensionar el difusor. Datos: \(T_0=472\, \mathrm{K}\); \(p_0=342000\, \mathrm{Pa}\); \(C_r=141\, \mathrm{m}\,\mathrm{s}^{-1}\); \(C_\theta=328\, \mathrm{m}\,\mathrm{s}^{-1}\); \(c_p=1004.5\, \mathrm{J}\,\mathrm{kg}^{-1}\,\mathrm{K}^{-1}\); \(\gamma=1.4\); \(R_g=287\, \mathrm{J}\,\mathrm{kg}^{-1}\,\mathrm{K}^{-1}\).

\noindent\textbf{Datos}
\begin{itemize}
\item $T_0 = 472\,\mathrm{K}$ — temperatura total
\item $p_0 = 342000\,\mathrm{Pa}$ — presión total
\item $C_r = 141\,\mathrm{m}\,\mathrm{s}^{-1}$ — componente radial
\item $C_\theta = 328\,\mathrm{m}\,\mathrm{s}^{-1}$ — componente tangencial
\item $c_p = 1004.5\,\mathrm{J}\,\mathrm{kg}^{-1}\,\mathrm{K}^{-1}$ — calor específico
\item $\gamma = 1.4$ — razón de calores
\item $R_g = 287\,\mathrm{J}\,\mathrm{kg}^{-1}\,\mathrm{K}^{-1}$ — constante gas
\end{itemize}

\noindent\textbf{Incógnita}
$\rho$, en $\mathrm{kg\,m^{-3}}$; también las magnitudes intermedias desarrolladas.

\noindent\textbf{Modelo físico}

\noindent\textit{Estado estático de descarga}

\[
\fitmath{\rho=\frac{p_0 \left(\frac{T_0 - \frac{C_r^{2} + C_\theta^{2}}{2 c_p}}{T_0}\right)^{\frac{\gamma}{\gamma - 1}}}{R_g \left(T_0 - \frac{C_r^{2} + C_\theta^{2}}{2 c_p}\right)}}
\]

\noindent\textbf{Hipótesis}
\begin{itemize}
\item Gas perfecto constante; componentes ortogonales.
\end{itemize}

\noindent\textbf{Desarrollo paso a paso}

\begin{enumerate}

\item Ambas componentes contribuyen a energía cinética.

\[\fitmath{T=T_0-\frac{{\mathrm{cr}}^{2}+{\mathrm{ct}}^{2}}{2\cdot \mathrm{cp}}}\]

\[\fitmath{T=(472)-\frac{{(141)}^{2}+{(328)}^{2}}{2\cdot (1004.5)}\approx 408.553\,\mathrm{K}}\]

\item La presión estática deriva del remanso local.

\[\fitmath{p=p_0\cdot {\left(\frac{T_0-\frac{{\mathrm{cr}}^{2}+{\mathrm{ct}}^{2}}{2\cdot \mathrm{cp}}}{T_0}\right)}^{\frac{\gamma}{\gamma-1}}}\]

\[\fitmath{p=(342000)\cdot {\left(\frac{(472)-\frac{{(141)}^{2}+{(328)}^{2}}{2\cdot (1004.5)}}{(472)}\right)}^{\frac{(1.4)}{(1.4)-1}}\approx 206347.2\,\mathrm{Pa}}\]

\item Se calcula densidad a partir del mismo estado estático.

\[\fitmath{\rho=\frac{p_0\cdot {\left(\frac{T_0-\frac{{\mathrm{cr}}^{2}+{\mathrm{ct}}^{2}}{2\cdot \mathrm{cp}}}{T_0}\right)}^{\frac{\gamma}{\gamma-1}}}{R_g\cdot \left(T_0-\frac{{\mathrm{cr}}^{2}+{\mathrm{ct}}^{2}}{2\cdot \mathrm{cp}}\right)}}\]

\[\fitmath{\rho=\frac{(342000)\cdot {\left(\frac{(472)-\frac{{(141)}^{2}+{(328)}^{2}}{2\cdot (1004.5)}}{(472)}\right)}^{\frac{(1.4)}{(1.4)-1}}}{(287)\cdot \left((472)-\frac{{(141)}^{2}+{(328)}^{2}}{2\cdot (1004.5)}\right)}\approx 1.75982\,\mathrm{kg/m^{3}}}\]

\end{enumerate}


\noindent\textbf{Sustitución numérica}

\[
\fitmath{\frac{\left(342000\,\mathrm{Pa}\right) \left(\frac{\left(472\,\mathrm{K}\right) - \frac{\left(\left(141\,\mathrm{m}\,\mathrm{s}^{-1}\right)\right)^{2} + \left(\left(328\,\mathrm{m}\,\mathrm{s}^{-1}\right)\right)^{2}}{2 \left(1004.5\,\mathrm{J}\,\mathrm{kg}^{-1}\,\mathrm{K}^{-1}\right)}}{\left(472\,\mathrm{K}\right)}\right)^{\frac{\left(1.4\right)}{\left(1.4\right) - 1}}}{\left(287\,\mathrm{J}\,\mathrm{kg}^{-1}\,\mathrm{K}^{-1}\right) \left(\left(472\,\mathrm{K}\right) - \frac{\left(\left(141\,\mathrm{m}\,\mathrm{s}^{-1}\right)\right)^{2} + \left(\left(328\,\mathrm{m}\,\mathrm{s}^{-1}\right)\right)^{2}}{2 \left(1004.5\,\mathrm{J}\,\mathrm{kg}^{-1}\,\mathrm{K}^{-1}\right)}\right)}}
\]

\noindent\textbf{Resultado}

\[
\fitmath{\rho=1.759820021\times10^{0}\,\mathrm{kg}\,\mathrm{m}^{-3}}
\]

\noindent\textbf{Interpretación}
Usar solo Cr para temperatura sobreestima la densidad disponible. La continuidad sí utiliza Cr.

\noindent\textbf{Comprobación de los pasos}
\begin{itemize}
\item Paso 2: $408.553$ — reproducido
\item Paso 4: $206347$ — reproducido
\item Paso 6: $1.75982$ — reproducido
\end{itemize}

\noindent\textbf{Comprobación de la respuesta} la aritmética del último paso reproduce la respuesta enunciada a la precisión con que está escrita. Es una reproducción aritmética; no valida el modelo físico ni sus hipótesis.

\noindent\textbf{Limitaciones}
\begin{itemize}
\item Modelo educativo; la aplicabilidad a una máquina real requiere revisión experta.
\end{itemize}

\noindent\textbf{Errores frecuentes}
\begin{itemize}
\item Usar distintas definiciones de velocidad en energía sin justificarlo.
\end{itemize}
% studium:end-record eyJpZCI6ICJQUkItMDEyNyIsICJzZWN0aW9uIjogImMxNiJ9
\end{tcolorbox}

\begin{tcolorbox}[title={Ejercicios propuestos}, studiumworked, breakable]

\noindent\textbf{1 (Básico, Desarrollo)}
\(C16-01-P \cdot\) Difusor sin álabes ideal. Ejercicio de transferencia. En un difusor de ancho y densidad constantes, calcule las componentes de salida y el incremento ideal de presión. Antes de consultar la solución, prediga el efecto de cambiar el primer dato respecto del caso resuelto de esta familia. Justifique la tendencia y explique qué restricción podría convertirse en limitante. Datos: \(r_e=0.22365\, \mathrm{m}\); \(r_s=0.32\, \mathrm{m}\); \(C_r=128\, \mathrm{m}\,\mathrm{s}^{-1}\); \(C_\theta=337\, \mathrm{m}\,\mathrm{s}^{-1}\); \(\rho=2.9\, \mathrm{kg}\,\mathrm{m}^{-3}\).

\noindent\textbf{2 (Intermedio, Desarrollo)}
\(C16-02-P \cdot\) Recuperación con pérdida prescrita. Ejercicio de transferencia. Calcule el coeficiente de recuperación estática y el incremento de presión con una pérdida normalizada dada. Antes de consultar la solución, prediga el efecto de cambiar el primer dato respecto del caso resuelto de esta familia. Justifique la tendencia y explique qué restricción podría convertirse en limitante. Datos: \(C_e=387.66\, \mathrm{m}\,\mathrm{s}^{-1}\); \(C_s=216\, \mathrm{m}\,\mathrm{s}^{-1}\); \(\zeta_D=0.13\); \(\rho=3.1\, \mathrm{kg}\,\mathrm{m}^{-3}\).

\noindent\textbf{3 (Avanzado, Desarrollo)}
\(C16-03-P \cdot\) Ancho del paso cilíndrico. Ejercicio de transferencia. Dimensione el ancho geométrico de salida radial con el caudal y área disponible dados. Antes de consultar la solución, prediga el efecto de cambiar el primer dato respecto del caso resuelto de esta familia. Justifique la tendencia y explique qué restricción podría convertirse en limitante. Datos: \(\dot m=9.2655\, \mathrm{kg}\,\mathrm{s}^{-1}\); \(r=0.198\, \mathrm{m}\); \(\rho=3.4\, \mathrm{kg}\,\mathrm{m}^{-3}\); \(C_r=126\, \mathrm{m}\,\mathrm{s}^{-1}\); \(B=0.94\).

\noindent\textbf{4 (Nivel de examen, Desarrollo)}
\(C16-04-P \cdot\) Estado estático de descarga. Ejercicio de transferencia. La descarga del rodete tiene estado total y componentes radial y tangencial dadas. Halle temperatura estática y densidad para dimensionar el difusor. Antes de consultar la solución, prediga el efecto de cambiar el primer dato respecto del caso resuelto de esta familia. Justifique la tendencia y explique qué restricción podría convertirse en limitante. Datos: \(T_0=502.68\, \mathrm{K}\); \(p_0=342000\, \mathrm{Pa}\); \(C_r=141\, \mathrm{m}\,\mathrm{s}^{-1}\); \(C_\theta=328\, \mathrm{m}\,\mathrm{s}^{-1}\); \(c_p=1004.5\, \mathrm{J}\,\mathrm{kg}^{-1}\,\mathrm{K}^{-1}\); \(\gamma=1.4\); \(R_g=287\, \mathrm{J}\,\mathrm{kg}^{-1}\,\mathrm{K}^{-1}\).
\end{tcolorbox}


\studiumchaptercheck

\chapter{Prestaciones y límites de un compresor radial}
\label{chap:c17}

\studiumchapterguide{Objetivo: desarrollar prestaciones y límites de un compresor radial, declarar hipótesis y comprobar condiciones físicas.}
% studium:end-record eyJpZCI6ICJQQVItMDExNSIsICJzZWN0aW9uIjogImMxNyJ9

\section{Introducción y objetivos}

El compresor radial debe satisfacer al mismo tiempo trabajo, capacidad de paso, eficiencia, Mach y resistencia. El objetivo es integrar esos criterios y estudiar sensibilidad a diámetro, giro, deslizamiento y contrapresión. Pertenece al bloque V; las comparaciones de diseño se realizan con modelos educativos explícitos.
Un mapa real requiere medidas o una simulación contrastada en un conjunto de puntos. Este capítulo utiliza curvas calculadas del modelo de Euler y esquemas de restricciones. Ninguna curva sintética se presenta como línea de bombeo, contorno de eficiencia o prestación de un fabricante.
% studium:end-record eyJpZCI6ICJQQVItMDExNiIsICJzZWN0aW9uIjogImMxNyJ9

\section{Modelo de altura de Euler}


Para salida curvada y deslizamiento prescrito,



\[
\fitmath{w_c=U_2\sigma_s(U_2+C_{m2}\tan\lambda_2)-U_1C_{\theta1}.}
\]
Si \(U_2\), \(\sigma_s\), geometría y densidad de salida se mantienen constantes, \(C_{m2}\) crece con el caudal. Con curvatura hacia atrás, \(\tan\lambda_2<0\), y el trabajo del modelo disminuye al aumentar caudal. Con álabes radiales ideales, la contribución tangencial es independiente de caudal en ese modelo reducido.
Estas tendencias no son un mapa completo: la densidad, el deslizamiento y las pérdidas pueden cambiar con el punto operativo. La figura~\ref{fig:limites} dibuja únicamente esa relación cinemática y señala las hipótesis que permanecen fijas. No contiene una línea de estabilidad calculada.
% studium:end-record eyJpZCI6ICJQQVItMDExNyIsICJzZWN0aW9uIjogImMxNyJ9

\par\vspace{\baselineskip}
\noindent\begin{minipage}{\linewidth}
\centering
\begingroup
\setbox0=\hbox{\begin{tikzpicture}
\draw[->,gray] (0,0)--(7,0);\draw[->,gray] (0,0)--(0,3.8);\draw (0.0,0)--(0.0,-.08) node[below,font=\scriptsize] {0};\draw (0,0.0)--(-.08,0.0) node[left,font=\scriptsize] {0};\draw (1.5,0)--(1.5,-.08) node[below,font=\scriptsize] {0.125};\draw (0,0.75)--(-.08,0.75) node[left,font=\scriptsize] {0.214};\draw (3.0,0)--(3.0,-.08) node[below,font=\scriptsize] {0.25};\draw (0,1.5)--(-.08,1.5) node[left,font=\scriptsize] {0.429};\draw (4.5,0)--(4.5,-.08) node[below,font=\scriptsize] {0.375};\draw (0,2.25)--(-.08,2.25) node[left,font=\scriptsize] {0.643};\draw (6.0,0)--(6.0,-.08) node[below,font=\scriptsize] {0.5};\draw (0,3.0)--(-.08,3.0) node[left,font=\scriptsize] {0.857};\draw[blue,very thick] plot coordinates {(0.00000,3.08000) (0.10000,3.06518) (0.20000,3.05036) (0.30000,3.03554) (0.40000,3.02073) (0.50000,3.00591) (0.60000,2.99109) (0.70000,2.97627) (0.80000,2.96145) (0.90000,2.94663) (1.00000,2.93181) (1.10000,2.91699) (1.20000,2.90218) (1.30000,2.88736) (1.40000,2.87254) (1.50000,2.85772) (1.60000,2.84290) (1.70000,2.82808) (1.80000,2.81326) (1.90000,2.79845) (2.00000,2.78363) (2.10000,2.76881) (2.20000,2.75399) (2.30000,2.73917) (2.40000,2.72435) (2.50000,2.70953) (2.60000,2.69471) (2.70000,2.67990) (2.80000,2.66508) (2.90000,2.65026) (3.00000,2.63544) (3.10000,2.62062) (3.20000,2.60580) (3.30000,2.59098) (3.40000,2.57617) (3.50000,2.56135) (3.60000,2.54653) (3.70000,2.53171) (3.80000,2.51689) (3.90000,2.50207) (4.00000,2.48725) (4.10000,2.47244) (4.20000,2.45762) (4.30000,2.44280) (4.40000,2.42798) (4.50000,2.41316) (4.60000,2.39834) (4.70000,2.38352) (4.80000,2.36870) (4.90000,2.35389) (5.00000,2.33907) (5.10000,2.32425) (5.20000,2.30943) (5.30000,2.29461) (5.40000,2.27979) (5.50000,2.26497) (5.60000,2.25016) (5.70000,2.23534) (5.80000,2.22052) (5.90000,2.20570) (6.00000,2.19088)};\node[below,align=center] at (3.5,-.65) {$C_{m2}/U_2$};\node[rotate=90] at (-1,2) {$w_c/U_2^2$};
\end{tikzpicture}}%
\count0=\ht0
\advance\count0 by \dp0
\count2=\wd0
\count4=\dimexpr\textheight/100*38\relax
\count6=\linewidth
\ifnum\count0<1
\count0=1
\fi
\ifnum\count2<1
\count2=1
\fi
\divide\count0 by 4096
\divide\count2 by 4096
\divide\count4 by 4096
\divide\count6 by 4096
\ifnum\count0<1
\count0=1
\fi
\ifnum\count2<1
\count2=1
\fi
\ifnum\count4<1
\count4=1
\fi
\ifnum\count6<1
\count6=1
\fi
\count8=\count2
\ifnum\count0>\count4
\count8=\numexpr\count2*\count4/\count0\relax
\fi
\ifnum\count8>\count6
\resizebox{\linewidth}{!}{\begin{tikzpicture}
\draw[->,gray] (0,0)--(7,0);\draw[->,gray] (0,0)--(0,3.8);\draw (0.0,0)--(0.0,-.08) node[below,font=\scriptsize] {0};\draw (0,0.0)--(-.08,0.0) node[left,font=\scriptsize] {0};\draw (1.5,0)--(1.5,-.08) node[below,font=\scriptsize] {0.125};\draw (0,0.75)--(-.08,0.75) node[left,font=\scriptsize] {0.214};\draw (3.0,0)--(3.0,-.08) node[below,font=\scriptsize] {0.25};\draw (0,1.5)--(-.08,1.5) node[left,font=\scriptsize] {0.429};\draw (4.5,0)--(4.5,-.08) node[below,font=\scriptsize] {0.375};\draw (0,2.25)--(-.08,2.25) node[left,font=\scriptsize] {0.643};\draw (6.0,0)--(6.0,-.08) node[below,font=\scriptsize] {0.5};\draw (0,3.0)--(-.08,3.0) node[left,font=\scriptsize] {0.857};\draw[blue,very thick] plot coordinates {(0.00000,3.08000) (0.10000,3.06518) (0.20000,3.05036) (0.30000,3.03554) (0.40000,3.02073) (0.50000,3.00591) (0.60000,2.99109) (0.70000,2.97627) (0.80000,2.96145) (0.90000,2.94663) (1.00000,2.93181) (1.10000,2.91699) (1.20000,2.90218) (1.30000,2.88736) (1.40000,2.87254) (1.50000,2.85772) (1.60000,2.84290) (1.70000,2.82808) (1.80000,2.81326) (1.90000,2.79845) (2.00000,2.78363) (2.10000,2.76881) (2.20000,2.75399) (2.30000,2.73917) (2.40000,2.72435) (2.50000,2.70953) (2.60000,2.69471) (2.70000,2.67990) (2.80000,2.66508) (2.90000,2.65026) (3.00000,2.63544) (3.10000,2.62062) (3.20000,2.60580) (3.30000,2.59098) (3.40000,2.57617) (3.50000,2.56135) (3.60000,2.54653) (3.70000,2.53171) (3.80000,2.51689) (3.90000,2.50207) (4.00000,2.48725) (4.10000,2.47244) (4.20000,2.45762) (4.30000,2.44280) (4.40000,2.42798) (4.50000,2.41316) (4.60000,2.39834) (4.70000,2.38352) (4.80000,2.36870) (4.90000,2.35389) (5.00000,2.33907) (5.10000,2.32425) (5.20000,2.30943) (5.30000,2.29461) (5.40000,2.27979) (5.50000,2.26497) (5.60000,2.25016) (5.70000,2.23534) (5.80000,2.22052) (5.90000,2.20570) (6.00000,2.19088)};\node[below,align=center] at (3.5,-.65) {$C_{m2}/U_2$};\node[rotate=90] at (-1,2) {$w_c/U_2^2$};
\end{tikzpicture}}%
\else
\ifnum\count0>\count4
\resizebox{!}{\dimexpr\textheight/100*38\relax}{\begin{tikzpicture}
\draw[->,gray] (0,0)--(7,0);\draw[->,gray] (0,0)--(0,3.8);\draw (0.0,0)--(0.0,-.08) node[below,font=\scriptsize] {0};\draw (0,0.0)--(-.08,0.0) node[left,font=\scriptsize] {0};\draw (1.5,0)--(1.5,-.08) node[below,font=\scriptsize] {0.125};\draw (0,0.75)--(-.08,0.75) node[left,font=\scriptsize] {0.214};\draw (3.0,0)--(3.0,-.08) node[below,font=\scriptsize] {0.25};\draw (0,1.5)--(-.08,1.5) node[left,font=\scriptsize] {0.429};\draw (4.5,0)--(4.5,-.08) node[below,font=\scriptsize] {0.375};\draw (0,2.25)--(-.08,2.25) node[left,font=\scriptsize] {0.643};\draw (6.0,0)--(6.0,-.08) node[below,font=\scriptsize] {0.5};\draw (0,3.0)--(-.08,3.0) node[left,font=\scriptsize] {0.857};\draw[blue,very thick] plot coordinates {(0.00000,3.08000) (0.10000,3.06518) (0.20000,3.05036) (0.30000,3.03554) (0.40000,3.02073) (0.50000,3.00591) (0.60000,2.99109) (0.70000,2.97627) (0.80000,2.96145) (0.90000,2.94663) (1.00000,2.93181) (1.10000,2.91699) (1.20000,2.90218) (1.30000,2.88736) (1.40000,2.87254) (1.50000,2.85772) (1.60000,2.84290) (1.70000,2.82808) (1.80000,2.81326) (1.90000,2.79845) (2.00000,2.78363) (2.10000,2.76881) (2.20000,2.75399) (2.30000,2.73917) (2.40000,2.72435) (2.50000,2.70953) (2.60000,2.69471) (2.70000,2.67990) (2.80000,2.66508) (2.90000,2.65026) (3.00000,2.63544) (3.10000,2.62062) (3.20000,2.60580) (3.30000,2.59098) (3.40000,2.57617) (3.50000,2.56135) (3.60000,2.54653) (3.70000,2.53171) (3.80000,2.51689) (3.90000,2.50207) (4.00000,2.48725) (4.10000,2.47244) (4.20000,2.45762) (4.30000,2.44280) (4.40000,2.42798) (4.50000,2.41316) (4.60000,2.39834) (4.70000,2.38352) (4.80000,2.36870) (4.90000,2.35389) (5.00000,2.33907) (5.10000,2.32425) (5.20000,2.30943) (5.30000,2.29461) (5.40000,2.27979) (5.50000,2.26497) (5.60000,2.25016) (5.70000,2.23534) (5.80000,2.22052) (5.90000,2.20570) (6.00000,2.19088)};\node[below,align=center] at (3.5,-.65) {$C_{m2}/U_2$};\node[rotate=90] at (-1,2) {$w_c/U_2^2$};
\end{tikzpicture}}%
\else
\begin{tikzpicture}
\draw[->,gray] (0,0)--(7,0);\draw[->,gray] (0,0)--(0,3.8);\draw (0.0,0)--(0.0,-.08) node[below,font=\scriptsize] {0};\draw (0,0.0)--(-.08,0.0) node[left,font=\scriptsize] {0};\draw (1.5,0)--(1.5,-.08) node[below,font=\scriptsize] {0.125};\draw (0,0.75)--(-.08,0.75) node[left,font=\scriptsize] {0.214};\draw (3.0,0)--(3.0,-.08) node[below,font=\scriptsize] {0.25};\draw (0,1.5)--(-.08,1.5) node[left,font=\scriptsize] {0.429};\draw (4.5,0)--(4.5,-.08) node[below,font=\scriptsize] {0.375};\draw (0,2.25)--(-.08,2.25) node[left,font=\scriptsize] {0.643};\draw (6.0,0)--(6.0,-.08) node[below,font=\scriptsize] {0.5};\draw (0,3.0)--(-.08,3.0) node[left,font=\scriptsize] {0.857};\draw[blue,very thick] plot coordinates {(0.00000,3.08000) (0.10000,3.06518) (0.20000,3.05036) (0.30000,3.03554) (0.40000,3.02073) (0.50000,3.00591) (0.60000,2.99109) (0.70000,2.97627) (0.80000,2.96145) (0.90000,2.94663) (1.00000,2.93181) (1.10000,2.91699) (1.20000,2.90218) (1.30000,2.88736) (1.40000,2.87254) (1.50000,2.85772) (1.60000,2.84290) (1.70000,2.82808) (1.80000,2.81326) (1.90000,2.79845) (2.00000,2.78363) (2.10000,2.76881) (2.20000,2.75399) (2.30000,2.73917) (2.40000,2.72435) (2.50000,2.70953) (2.60000,2.69471) (2.70000,2.67990) (2.80000,2.66508) (2.90000,2.65026) (3.00000,2.63544) (3.10000,2.62062) (3.20000,2.60580) (3.30000,2.59098) (3.40000,2.57617) (3.50000,2.56135) (3.60000,2.54653) (3.70000,2.53171) (3.80000,2.51689) (3.90000,2.50207) (4.00000,2.48725) (4.10000,2.47244) (4.20000,2.45762) (4.30000,2.44280) (4.40000,2.42798) (4.50000,2.41316) (4.60000,2.39834) (4.70000,2.38352) (4.80000,2.36870) (4.90000,2.35389) (5.00000,2.33907) (5.10000,2.32425) (5.20000,2.30943) (5.30000,2.29461) (5.40000,2.27979) (5.50000,2.26497) (5.60000,2.25016) (5.70000,2.23534) (5.80000,2.22052) (5.90000,2.20570) (6.00000,2.19088)};\node[below,align=center] at (3.5,-.65) {$C_{m2}/U_2$};\node[rotate=90] at (-1,2) {$w_c/U_2^2$};
\end{tikzpicture}%
\fi
\fi
\endgroup
\par\nopagebreak
\studiumcaption{fig:limites}{Modelo cinemático de Euler: slip=0,88 y ángulo desde meridiana -30 grados, sin prerrotación. No predice estabilidad.}
\end{minipage}
\par\vspace{\baselineskip}
% studium:end-record eyJpZCI6ICJGSUctMDAxNyIsICJzZWN0aW9uIjogImMxNyJ9

\section{Velocidad periférica y sensibilidad}

Con salida radial, sin prerrotación y deslizamiento fijo, \(w_c=\sigma_sU_2^2\). Una variación pequeña satisface

\[
\fitmath{\frac{dw_c}{w_c}=\frac{d\sigma_s}{\sigma_s}+2\frac{dU_2}{U_2}.}
\]
Como \(U_2=\pi D_2N/60\), un aumento de diámetro o giro eleva el trabajo de forma cuadrática, pero también aumenta las solicitaciones y el Mach. No es válido evaluar el beneficio energético y mantener sin cambios todas las limitaciones.
Para un trabajo dado, \(U_2=\sqrt{w_c/\sigma_s}\). Una reducción del factor de deslizamiento exige mayor velocidad para conservar trabajo. La eficiencia termodinámica determina después qué relación de presión corresponde a ese trabajo. Dos factores con valores parecidos, como \(\sigma_s\) y \(\eta_c\), actúan en ecuaciones distintas.
% studium:end-record eyJpZCI6ICJQQVItMDExOCIsICJzZWN0aW9uIjogImMxNyJ9

\section{Bloqueo y estabilidad}

El inductor y la garganta del difusor pueden limitar el caudal por bloqueo. La condición relevante se expresa en la referencia de cada garganta. Hacia caudales bajos pueden aparecer separación y fenómenos de inestabilidad del sistema. El compresor radial tampoco es inmune al bombeo.
La línea operativa depende del sistema aguas abajo y del acoplamiento de ejes. Las pérdidas del difusor pueden modificar la característica global, incluso si el rodete conserva un trabajo de Euler alto. Una caída de presión total en un elemento pasivo se contabiliza en el balance del conjunto.
Las variables corregidas permiten comparar condiciones de entrada, pero no eliminan toda dependencia de Reynolds ni de composición. Una corrección de temperatura no convierte un ensayo de un gas en otro idéntico si cambian \(\gamma\), \(R_g\) o viscosidad.
% studium:end-record eyJpZCI6ICJQQVItMDExOSIsICJzZWN0aW9uIjogImMxNyJ9

\section{Evaluación de un punto de diseño}

Un expediente mínimo incluye caudal y estados de entrada, velocidad de giro, radios del inductor y salida, triángulos, factor de deslizamiento definido, modelo de pérdidas, temperaturas estáticas, Mach, anchos de paso y balance de potencia. Deben incluirse también los criterios que no pueden evaluarse por falta de datos.
La comparación de dos alternativas debe conservar la misma misión: igual caudal, relación de presión y condiciones de entrada, o declarar qué requisito cambia. Comparar solo eficiencias puede ocultar una gran diferencia de tamaño o de potencia. Comparar solo diámetro puede ocultar que una alternativa requiere un giro imposible para su material.
% studium:end-record eyJpZCI6ICJQQVItMDEyMCIsICJzZWN0aW9uIjogImMxNyJ9

\section{Incertidumbre y comprobación inversa}

Si un resultado depende de eficiencia y deslizamiento inciertos, se explora un intervalo y se recalcula la cadena. No se escriben más cifras significativas que las justificadas por los datos. La reproducción aritmética se guarda con precisión interna; la presentación se redondea según sensibilidad y uso.
Una comprobación inversa calcula el trabajo a partir de \(c_p\Delta T_0\), el remolino a partir de Euler y la presión isentrópica equivalente a partir de la eficiencia. Los residuos detectan incoherencias. No son evidencia experimental de que el factor de deslizamiento o el rendimiento elegidos representen una máquina real.
Errores habituales: declarar un mapa a partir de una sola ecuación, confundir bloqueo e inestabilidad, usar la eficiencia como factor sobre \(U^2\) y comparar puntos con condiciones de entrada distintas sin corrección. Autocomprobación: si \(\sigma_s\) disminuye y se conserva giro, explique la variación de trabajo, temperatura total y presión para eficiencia fija.
Formulario: Euler radial con deslizamiento, \(U=\pi DN/60\), \(\Delta T_0=w_c/c_p\), variables corregidas del capítulo 10. Referencias: ULE Compresores radiales, pp. 2 y \(5-16\); Generalidades, p. 22; contraste de Euler radial en MIT 16.50, lección 24, p. 4. Las curvas de sensibilidad son calculadas, con validez limitada al modelo indicado.
% studium:end-record eyJpZCI6ICJQQVItMDEyMSIsICJzZWN0aW9uIjogImMxNyJ9

\begin{tcolorbox}[title={Desarrollos del modelo}, studiumdef, breakable]

\noindent\textbf{Modelo desarrollado: Velocidad exigida por trabajo}

\begin{equation}
r=\frac{1}{\Omega}\sqrt{\frac{w_c}{\sigma_s}}
\label{eq:EC-C17}
\end{equation}

\noindent\textit{Significado.} Radio de salida exigido por un trabajo específico sin prerrotación.

\noindent\textbf{Hipótesis}
\begin{itemize}
\item Salida radial, sin prerrotación, slip constante.
\end{itemize}

\noindent\textbf{Principios}
\begin{itemize}
\item Balance y definiciones físicas desarrollados en el capítulo; propiedades y dominio declarados en el problema.
\end{itemize}

\noindent\textbf{Aplicabilidad}
\begin{itemize}
\item Modelo educativo del caso C17-01-R
\end{itemize}

\noindent\textbf{Limitaciones}
\begin{itemize}
\item Modelo educativo; la aplicabilidad a una máquina real requiere revisión experta.
\end{itemize}

\noindent\textbf{Pasos}
\begin{enumerate}
\item
La salida radial sin prerrotación da wc=slip U2 al cuadrado.

\[
\fitmath{U=\sqrt{\frac{w_c}{\sigma_s}}}
\]
\item
La geometría requerida sigue de U=Omega r.

\[
\fitmath{r=\frac{\sqrt{\frac{w_c}{\sigma_s}}}{\Omega}}
\]
\end{enumerate}

{\noindent\footnotesize verificada dimensionalmente.\par}
\end{tcolorbox}

\begin{tcolorbox}[title={Problema resuelto 1}, studiumworked, breakable]

\noindent\textbf{Enunciado}
\(C17-01-R \cdot\) Velocidad exigida por trabajo. Un rodete de salida radial debe suministrar trabajo dado con factor de deslizamiento prescrito. Calcule la velocidad periférica y el radio para el giro indicado. Datos: \(w_c=169000\, \mathrm{J}\,\mathrm{kg}^{-1}\); \(\sigma_s=0.86\); \(\Omega=2140\, \mathrm{s}^{-1}\).

\noindent\textbf{Datos}
\begin{itemize}
\item $w_c = 169000\,\mathrm{J}\,\mathrm{kg}^{-1}$ — trabajo requerido
\item $\sigma_s = 0.86$ — factor de deslizamiento
\item $\Omega = 2140\,\mathrm{s}^{-1}$ — velocidad angular
\end{itemize}

\noindent\textbf{Incógnita}
$r$, en $\mathrm{m}$; también las magnitudes intermedias desarrolladas.

\noindent\textbf{Modelo físico}

\noindent\textit{Velocidad exigida por trabajo}

\[
\fitmath{r=\frac{\sqrt{\frac{w_c}{\sigma_s}}}{\Omega}}
\]

\noindent\textbf{Hipótesis}
\begin{itemize}
\item Salida radial, sin prerrotación, slip constante.
\end{itemize}

\noindent\textbf{Desarrollo paso a paso}

\begin{enumerate}

\item La salida radial sin prerrotación da wc=slip U2 al cuadrado.

\[\fitmath{\mathrm{uu}=\sqrt{\frac{\mathrm{wc}}{\mathrm{slip}}}}\]

\[\fitmath{\mathrm{uu}=\sqrt{\frac{(169000)}{(0.86)}}\approx 443.2963\,\mathrm{m/s}}\]

\item La geometría requerida sigue de U=Omega r.

\[\fitmath{\mathrm{rad}=\frac{\sqrt{\frac{\mathrm{wc}}{\mathrm{slip}}}}{\Omega}}\]

\[\fitmath{\mathrm{rad}=\frac{\sqrt{\frac{(169000)}{(0.86)}}}{(2140)}\approx 0.2071478\,\mathrm{m}}\]

\end{enumerate}


\noindent\textbf{Sustitución numérica}

\[
\fitmath{\frac{\sqrt{\frac{\left(169000\,\mathrm{J}\,\mathrm{kg}^{-1}\right)}{\left(0.86\right)}}}{\left(2140\,\mathrm{s}^{-1}\right)}}
\]

\noindent\textbf{Resultado}

\[
\fitmath{r=2.071478134\times10^{-1}\,\mathrm{m}}
\]

\noindent\textbf{Interpretación}
El tamaño y el giro están ligados al trabajo. Las tensiones y el Mach deben comprobarse después.

\noindent\textbf{Comprobación de los pasos}
\begin{itemize}
\item Paso 2: $443.296$ — reproducido
\item Paso 4: $0.207148$ — reproducido
\end{itemize}

\noindent\textbf{Comprobación de la respuesta} la aritmética del último paso reproduce la respuesta enunciada a la precisión con que está escrita. Es una reproducción aritmética; no valida el modelo físico ni sus hipótesis.

\noindent\textbf{Limitaciones}
\begin{itemize}
\item Modelo educativo; la aplicabilidad a una máquina real requiere revisión experta.
\end{itemize}

\noindent\textbf{Errores frecuentes}
\begin{itemize}
\item Sustituir eficiencia isentrópica por factor de deslizamiento.
\end{itemize}
% studium:end-record eyJpZCI6ICJQUkItMDEyOSIsICJzZWN0aW9uIjogImMxNyJ9
\end{tcolorbox}

\begin{tcolorbox}[title={Problema resuelto 2}, studiumworked, breakable]

\noindent\textbf{Enunciado}
\(C17-02-R \cdot\) Presión y sensibilidad de eficiencia. Para trabajo de Euler fijado, calcule la relación de presión que produce el rendimiento suministrado. Datos: \(w_c=178000\, \mathrm{J}\,\mathrm{kg}^{-1}\); \(T_{0e}=297\, \mathrm{K}\); \(\eta=0.81\); \(c_p=1004.5\, \mathrm{J}\,\mathrm{kg}^{-1}\,\mathrm{K}^{-1}\); \(\gamma=1.4\).

\noindent\textbf{Datos}
\begin{itemize}
\item $w_c = 178000\,\mathrm{J}\,\mathrm{kg}^{-1}$ — trabajo de Euler comunicado
\item $T_{0e} = 297\,\mathrm{K}$ — temperatura total entrada
\item $\eta = 0.81$ — eficiencia total-total
\item $c_p = 1004.5\,\mathrm{J}\,\mathrm{kg}^{-1}\,\mathrm{K}^{-1}$ — calor específico
\item $\gamma = 1.4$ — razón de calores
\end{itemize}

\noindent\textbf{Incógnita}
$\pi_c$ (adimensional); también las magnitudes intermedias desarrolladas.

\noindent\textbf{Modelo físico}

\noindent\textit{Presión y sensibilidad de eficiencia}

\[
\fitmath{\pi_c=\left(1 + \frac{\eta w_c}{T_{0e} c_p}\right)^{\frac{\gamma}{\gamma - 1}}}
\]

\noindent\textbf{Hipótesis}
\begin{itemize}
\item Gas perfecto y trabajo fijado, sin refrigeración.
\end{itemize}

\noindent\textbf{Desarrollo paso a paso}

\begin{enumerate}

\item El trabajo real fija incremento total independientemente de eta.

\[\fitmath{\mathrm{dt}=\frac{\mathrm{wc}}{\mathrm{cp}}}\]

\[\fitmath{\mathrm{dt}=\frac{(178000)}{(1004.5)}\approx 177.2026\,\mathrm{K}}\]

\item La eficiencia fija el incremento isentrópico equivalente.

\[\fitmath{\mathrm{pr}={\left(1+\frac{\eta\cdot \mathrm{wc}}{\mathrm{cp}\cdot \mathrm{Tin}}\right)}^{\frac{\gamma}{\gamma-1}}}\]

\[\fitmath{\mathrm{pr}={\left(1+\frac{(0.81)\cdot (178000)}{(1004.5)\cdot (297)}\right)}^{\frac{(1.4)}{(1.4)-1}}\approx 3.974484\,}\]

\end{enumerate}


\noindent\textbf{Sustitución numérica}

\[
\fitmath{\left(1 + \frac{\left(0.81\right) \left(178000\,\mathrm{J}\,\mathrm{kg}^{-1}\right)}{\left(297\,\mathrm{K}\right) \left(1004.5\,\mathrm{J}\,\mathrm{kg}^{-1}\,\mathrm{K}^{-1}\right)}\right)^{\frac{\left(1.4\right)}{\left(1.4\right) - 1}}}
\]

\noindent\textbf{Resultado}

\[
\fitmath{\pi_c=3.974484195\times10^{0}}
\]

\noindent\textbf{Interpretación}
Variar eta a trabajo fijo cambia la presión, mientras el aumento real de temperatura permanece fijo.

\noindent\textbf{Comprobación de los pasos}
\begin{itemize}
\item Paso 2: $177.203$ — reproducido
\item Paso 4: $3.97448$ — reproducido
\end{itemize}

\noindent\textbf{Comprobación de la respuesta} la aritmética del último paso reproduce la respuesta enunciada a la precisión con que está escrita. Es una reproducción aritmética; no valida el modelo físico ni sus hipótesis.

\noindent\textbf{Limitaciones}
\begin{itemize}
\item Modelo educativo; la aplicabilidad a una máquina real requiere revisión experta.
\end{itemize}

\noindent\textbf{Errores frecuentes}
\begin{itemize}
\item Reducir el calentamiento real multiplicando por eficiencia.
\end{itemize}
% studium:end-record eyJpZCI6ICJQUkItMDEzMSIsICJzZWN0aW9uIjogImMxNyJ9
\end{tcolorbox}

\begin{tcolorbox}[title={Problema resuelto 3}, studiumworked, breakable]

\noindent\textbf{Enunciado}
\(C17-03-R \cdot\) Pendiente cinemática con caudal. Un modelo de rodete mantiene densidad, ancho, slip y giro constantes. Calcule la derivada del trabajo respecto al caudal. Datos: \(U=434\, \mathrm{m}\,\mathrm{s}^{-1}\); \(\sigma_s=0.88\); \(\lambda_s=-0.5410521\, \mathrm{rad}\); \(r=0.203\, \mathrm{m}\); \(b=0.017\, \mathrm{m}\); \(\rho=3.2\, \mathrm{kg}\,\mathrm{m}^{-3}\).

\noindent\textbf{Datos}
\begin{itemize}
\item $U = 434\,\mathrm{m}\,\mathrm{s}^{-1}$ — velocidad periférica
\item $\sigma_s = 0.88$ — factor de deslizamiento
\item $\lambda_s = -0.54105206811824\,\mathrm{rad}$ — ángulo ideal desde meridiana
\item $r = 0.203\,\mathrm{m}$ — radio salida
\item $b = 0.017\,\mathrm{m}$ — ancho salida
\item $\rho = 3.2\,\mathrm{kg}\,\mathrm{m}^{-3}$ — densidad salida
\end{itemize}

\noindent\textbf{Incógnita}
$\mathrm{d}w_c/\mathrm{d}\dot{m}$, en $\mathrm{J\,s\,kg^{-2}}$; también las magnitudes intermedias desarrolladas.

\noindent\textbf{Modelo físico}

\noindent\textit{Pendiente cinemática con caudal}

\[
\fitmath{\frac{\mathrm{d}w_c}{\mathrm{d}\dot{m}}=\frac{\sigma_s U \tan{\left(\lambda_s \right)}}{2 \pi r \rho b}}
\]

\noindent\textbf{Hipótesis}
\begin{itemize}
\item Densidad, ancho, slip, radio y giro fijos; sin prerrotación.
\end{itemize}

\noindent\textbf{Desarrollo paso a paso}

\begin{enumerate}

\item La continuidad da la derivada de velocidad meridiana por caudal.

\[\fitmath{\mathrm{dcm}=\frac{1}{2\cdot \pi\cdot \mathrm{rad}\cdot \mathrm{width}\cdot \rho}}\]

\[\fitmath{\mathrm{dcm}=\frac{1}{2\cdot \pi\cdot (0.203)\cdot (0.017)\cdot (3.2)}\approx 14.41203\,\mathrm{m/kg}}\]

\item La derivada del modelo de Euler incorpora el ángulo.

\[\fitmath{\mathrm{slope}=\frac{\mathrm{uu}\cdot \mathrm{slip}\cdot \tan\left(\mathrm{lam}\right)}{2\cdot \pi\cdot \mathrm{rad}\cdot \mathrm{width}\cdot \rho}}\]

\[\fitmath{\mathrm{slope}=\frac{(434)\cdot (0.88)\cdot \tan\left((-0.5410520681)\right)}{2\cdot \pi\cdot (0.203)\cdot (0.017)\cdot (3.2)}\approx -3307.283\,\mathrm{J\,s/kg^{2}}}\]

\end{enumerate}


\noindent\textbf{Sustitución numérica}

\[
\fitmath{\frac{\left(0.88\right) \left(434\,\mathrm{m}\,\mathrm{s}^{-1}\right) \tan{\left(\left(-0.54105206811824\,\mathrm{rad}\right) \right)}}{2 \pi \left(0.203\,\mathrm{m}\right) \left(3.2\,\mathrm{kg}\,\mathrm{m}^{-3}\right) \left(0.017\,\mathrm{m}\right)}}
\]

\noindent\textbf{Resultado}

\[
\fitmath{\frac{\mathrm{d}w_c}{\mathrm{d}\dot{m}}=-3.307282833\times10^{3}\,\mathrm{J}\,\mathrm{s}\,\mathrm{kg}^{-2}}
\]

\noindent\textbf{Interpretación}
La pendiente negativa proviene de curvatura hacia atrás. No es una predicción de la línea de bombeo.

\noindent\textbf{Comprobación de los pasos}
\begin{itemize}
\item Paso 2: $14.412$ — reproducido
\item Paso 4: $-3307.28$ — reproducido
\end{itemize}

\noindent\textbf{Comprobación de la respuesta} la aritmética del último paso reproduce la respuesta enunciada a la precisión con que está escrita. Es una reproducción aritmética; no valida el modelo físico ni sus hipótesis.

\noindent\textbf{Limitaciones}
\begin{itemize}
\item Modelo educativo; la aplicabilidad a una máquina real requiere revisión experta.
\end{itemize}

\noindent\textbf{Errores frecuentes}
\begin{itemize}
\item Presentar la pendiente de Euler como estabilidad global.
\end{itemize}
% studium:end-record eyJpZCI6ICJQUkItMDEzMyIsICJzZWN0aW9uIjogImMxNyJ9
\end{tcolorbox}

\begin{tcolorbox}[title={Problema resuelto 4}, studiumworked, breakable]

\noindent\textbf{Enunciado}
\(C17-04-R \cdot\) Balance inverso de rendimiento. De un ensayo educativo se conocen aumento total de temperatura y relación de presión. Calcule la eficiencia total-total y juzgue admisibilidad. Datos: \(T_{0e}=303\, \mathrm{K}\); \(T_{0s}=469\, \mathrm{K}\); \(\pi_c=3.85\); \(\gamma=1.4\).

\noindent\textbf{Datos}
\begin{itemize}
\item $T_{0e} = 303\,\mathrm{K}$ — temperatura total entrada
\item $T_{0s} = 469\,\mathrm{K}$ — temperatura total salida
\item $\pi_c = 3.85$ — relación de presión total
\item $\gamma = 1.4$ — razón de calores
\end{itemize}

\noindent\textbf{Incógnita}
$\eta$ (adimensional); también las magnitudes intermedias desarrolladas.

\noindent\textbf{Modelo físico}

\noindent\textit{Balance inverso de rendimiento}

\[
\fitmath{\eta=\frac{T_{0e} \left(\pi_c^{\frac{\gamma - 1}{\gamma}} - 1\right)}{- T_{0e} + T_{0s}}}
\]

\noindent\textbf{Hipótesis}
\begin{itemize}
\item Gas perfecto constante, sin calor externo; mediciones medias representativas.
\end{itemize}

\noindent\textbf{Desarrollo paso a paso}

\begin{enumerate}

\item Se calcula la temperatura isentrópica equivalente.

\[\fitmath{\mathrm{Tis}=\mathrm{Tin}\cdot {\mathrm{pr}}^{\frac{\gamma-1}{\gamma}}}\]

\[\fitmath{\mathrm{Tis}=(303)\cdot {(3.85)}^{\frac{(1.4)-1}{(1.4)}}\approx 445.3661\,\mathrm{K}}\]

\item Se divide aumento ideal por aumento real.

\[\fitmath{\mathrm{eta}=\frac{\mathrm{Tin}\cdot \left({\mathrm{pr}}^{\frac{\gamma-1}{\gamma}}-1\right)}{\mathrm{Tout}-\mathrm{Tin}}}\]

\[\fitmath{\mathrm{eta}=\frac{(303)\cdot \left({(3.85)}^{\frac{(1.4)-1}{(1.4)}}-1\right)}{(469)-(303)}\approx 0.8576269\,}\]

\end{enumerate}


\noindent\textbf{Sustitución numérica}

\[
\fitmath{\frac{\left(303\,\mathrm{K}\right) \left(\left(3.85\right)^{\frac{\left(1.4\right) - 1}{\left(1.4\right)}} - 1\right)}{- \left(303\,\mathrm{K}\right) + \left(469\,\mathrm{K}\right)}}
\]

\noindent\textbf{Resultado}

\[
\fitmath{\eta=8.576268981\times10^{-1}}
\]

\noindent\textbf{Interpretación}
El valor debe estar entre cero y uno bajo las hipótesis adiabáticas. Un valor admisible no prueba precisión experimental.

\noindent\textbf{Comprobación de los pasos}
\begin{itemize}
\item Paso 2: $445.366$ — reproducido
\item Paso 4: $0.857627$ — reproducido
\end{itemize}

\noindent\textbf{Comprobación de la respuesta} la aritmética del último paso reproduce la respuesta enunciada a la precisión con que está escrita. Es una reproducción aritmética; no valida el modelo físico ni sus hipótesis.

\noindent\textbf{Limitaciones}
\begin{itemize}
\item Modelo educativo; la aplicabilidad a una máquina real requiere revisión experta.
\end{itemize}

\noindent\textbf{Errores frecuentes}
\begin{itemize}
\item Confundir presión estática medida con total sin corrección.
\end{itemize}
% studium:end-record eyJpZCI6ICJQUkItMDEzNSIsICJzZWN0aW9uIjogImMxNyJ9
\end{tcolorbox}

\begin{tcolorbox}[title={Ejercicios propuestos}, studiumworked, breakable]

\noindent\textbf{1 (Básico, Desarrollo)}
\(C17-01-P \cdot\) Velocidad exigida por trabajo. Ejercicio de transferencia. Un rodete de salida radial debe suministrar trabajo dado con factor de deslizamiento prescrito. Calcule la velocidad periférica y el radio para el giro indicado. Antes de consultar la solución, prediga el efecto de cambiar el primer dato respecto del caso resuelto de esta familia. Justifique la tendencia y explique qué restricción podría convertirse en limitante. Datos: \(w_c=179985\, \mathrm{J}\,\mathrm{kg}^{-1}\); \(\sigma_s=0.86\); \(\Omega=2140\, \mathrm{s}^{-1}\).

\noindent\textbf{2 (Intermedio, Desarrollo)}
\(C17-02-P \cdot\) Presión y sensibilidad de eficiencia. Ejercicio de transferencia. Para trabajo de Euler fijado, calcule la relación de presión que produce el rendimiento suministrado. Antes de consultar la solución, prediga el efecto de cambiar el primer dato respecto del caso resuelto de esta familia. Justifique la tendencia y explique qué restricción podría convertirse en limitante. Datos: \(w_c=189570\, \mathrm{J}\,\mathrm{kg}^{-1}\); \(T_{0e}=297\, \mathrm{K}\); \(\eta=0.81\); \(c_p=1004.5\, \mathrm{J}\,\mathrm{kg}^{-1}\,\mathrm{K}^{-1}\); \(\gamma=1.4\).

\noindent\textbf{3 (Avanzado, Desarrollo)}
\(C17-03-P \cdot\) Pendiente cinemática con caudal. Ejercicio de transferencia. Un modelo de rodete mantiene densidad, ancho, slip y giro constantes. Calcule la derivada del trabajo respecto al caudal. Antes de consultar la solución, prediga el efecto de cambiar el primer dato respecto del caso resuelto de esta familia. Justifique la tendencia y explique qué restricción podría convertirse en limitante. Datos: \(U=462.21\, \mathrm{m}\,\mathrm{s}^{-1}\); \(\sigma_s=0.88\); \(\lambda_s=-0.5410521\, \mathrm{rad}\); \(r=0.203\, \mathrm{m}\); \(b=0.017\, \mathrm{m}\); \(\rho=3.2\, \mathrm{kg}\,\mathrm{m}^{-3}\).

\noindent\textbf{4 (Nivel de examen, ERROR\_IDENTIFICATION)}
\(C17-04-P \cdot\) Balance inverso de rendimiento. Ejercicio de transferencia. De un ensayo educativo se conocen aumento total de temperatura y relación de presión. Calcule la eficiencia total-total y juzgue admisibilidad. Antes de consultar la solución, prediga el efecto de cambiar el primer dato respecto del caso resuelto de esta familia. Justifique la tendencia y explique qué restricción podría convertirse en limitante. Datos: \(T_{0e}=322.695\, \mathrm{K}\); \(T_{0s}=469\, \mathrm{K}\); \(\pi_c=3.85\); \(\gamma=1.4\).
\end{tcolorbox}


\studiumchaptercheck

\part{Prácticas y ampliaciones identificadas}

\chapter{Práctica reproducible: secciones a distintos radios}
\label{chap:c18}

\studiumchapterguide{Objetivo: desarrollar práctica reproducible: secciones a distintos radios, declarar hipótesis y comprobar condiciones físicas.}
% studium:end-record eyJpZCI6ICJQQVItMDEyMiIsICJzZWN0aW9uIjogImMxOCJ9

\section{Introducción y objetivos}

El documento docente del trabajo de hoja de cálculo solicita representación de secciones a varios radios y contempla cuatro configuraciones: rotor compresor, estator compresor, estator turbina y rotor turbina. Este capítulo enseña el procedimiento con una familia geométrica original, sin sustituir el trabajo personal ni proporcionar una entrega particular de un grupo.
La práctica desarrolla entradas, ecuaciones, coordenadas, gráficos, comprobaciones e interpretación. El programa reproducible incluido genera tablas CSV y figuras a partir de un archivo de parámetros. Esas tablas pueden importarse en una hoja de cálculo y permiten comprobar cada fórmula sin depender de una aplicación concreta.
% studium:end-record eyJpZCI6ICJQQVItMDEyMyIsICJzZWN0aW9uIjogImMxOCJ9

\section{Entradas y separación de funciones}

Se registran \(\Omega\), radios de raíz, medio y punta, velocidades axiales, constantes de remolino por estación, cuerda axial, espesor y número de álabes. Se añaden incidencia y desviación prescritas. Los parámetros termodinámicos solo son necesarios si se pretende calcular densidad o Mach; no se inventan cuando el objetivo es puramente geométrico.
La configuración decide qué velocidades se usan. En rotor compresor, \(a_e=\beta_1\), \(a_s=\beta_2\). En estator compresor, \(a_e=\alpha_2\), \(a_s=\alpha_3\). En estator turbina, \(a_e=\alpha_1\), \(a_s=\alpha_2\). En rotor turbina, \(a_e=\beta_2\), \(a_s=\beta_3\). No se obtiene una de las cuatro cambiando solo el título de una misma tabla.
% studium:end-record eyJpZCI6ICJQQVItMDEyNCIsICJzZWN0aW9uIjogImMxOCJ9

\section{Cálculo por radio}

Con vórtice libre en cada estación, \(C_{\theta j}=K_{v,j}/r\). Se calculan \(U=\Omega r\), \(W_{\theta j}=C_{\theta j}-U\), \(\alpha_j=\arctan(C_{\theta j}/C_{xj})\) y \(\beta_j=\arctan(W_{\theta j}/C_{xj})\). Cada fila registra radio, componentes, módulos y ángulos; la comprobación vectorial debe producir residuo cero dentro de tolerancia.
Después se deducen ángulos metálicos con los valores de incidencia y desviación elegidos y se construye la línea media del capítulo 12. La figura~\ref{fig:practica} superpone secciones calculadas con escala común y radios identificados. No representa un diseño de fabricación ni una pala de un motor concreto.
% studium:end-record eyJpZCI6ICJQQVItMDEyNSIsICJzZWN0aW9uIjogImMxOCJ9

\par\vspace{\baselineskip}
\noindent\begin{minipage}{\linewidth}
\centering
\begingroup
\setbox0=\hbox{\begin{tikzpicture}
\draw[->,gray] (-.3,0)--(4.2,0) node[right] {$x$ (escala 0,75 cm por cm)};\draw[->,gray] (0,.4)--(0,-5.8) node[below] {$y$ (misma escala)};\draw[blue,thick] plot coordinates {(0.00000,0.00000) (0.12871,-0.11105) (0.25701,-0.21883) (0.38477,-0.32344) (0.51187,-0.42502) (0.63819,-0.52370) (0.76364,-0.61959) (0.88813,-0.71284) (1.01159,-0.80357) (1.13397,-0.89191) (1.25522,-0.97797) (1.37533,-1.06188) (1.49429,-1.14373) (1.61212,-1.22363) (1.72885,-1.30167) (1.84453,-1.37792) (1.95923,-1.45243) (2.07303,-1.52526) (2.18602,-1.59643) (2.29832,-1.66596) (2.41004,-1.73384) (2.52132,-1.80004) (2.63228,-1.86452) (2.74307,-1.92721) (2.85383,-1.98803) (2.96470,-2.04688) (3.07580,-2.10364) (3.18727,-2.15816) (3.29923,-2.21030) (3.41178,-2.25989) (3.52500,-2.30675)};\draw[blue,thick] plot coordinates {(0.00000,-0.00000) (0.10629,-0.13293) (0.21299,-0.26291) (0.32023,-0.38983) (0.42813,-0.51357) (0.53681,-0.63400) (0.64636,-0.75098) (0.75687,-0.86439) (0.86841,-0.97409) (0.98103,-1.07997) (1.09478,-1.18191) (1.20967,-1.27977) (1.32571,-1.37347) (1.44288,-1.46291) (1.56115,-1.54798) (1.68047,-1.62863) (1.80077,-1.70479) (1.92197,-1.77641) (2.04398,-1.84347) (2.16668,-1.90595) (2.28996,-1.96387) (2.41368,-2.01724) (2.53772,-2.06611) (2.66193,-2.11054) (2.78617,-2.15063) (2.91030,-2.18647) (3.03420,-2.21818) (3.15773,-2.24591) (3.28077,-2.26980) (3.40322,-2.29002) (3.52500,-2.30675)};\node[blue,right] at (3.7,-2.3067) {$r=0.215$ m};\draw[teal,thick] plot coordinates {(0.00000,0.00000) (0.13026,-0.15738) (0.26025,-0.31215) (0.38983,-0.46441) (0.51885,-0.61426) (0.64721,-0.76181) (0.77477,-0.90715) (0.90143,-1.05040) (1.02709,-1.19163) (1.15167,-1.33095) (1.27510,-1.46846) (1.39732,-1.60423) (1.51830,-1.73834) (1.63800,-1.87086) (1.75643,-2.00185) (1.87359,-2.13137) (1.98950,-2.25945) (2.10420,-2.38612) (2.21774,-2.51141) (2.33020,-2.63531) (2.44164,-2.75783) (2.55217,-2.87894) (2.66189,-2.99860) (2.77090,-3.11678) (2.87934,-3.23341) (2.98731,-3.34843) (3.09497,-3.46175) (3.20243,-3.57327) (3.30984,-3.68290) (3.41732,-3.79051) (3.52500,-3.89599)};\draw[teal,thick] plot coordinates {(0.00000,-0.00000) (0.10474,-0.17553) (0.20975,-0.34863) (0.31517,-0.51918) (0.42115,-0.68710) (0.52779,-0.85228) (0.63523,-1.01461) (0.74357,-1.17400) (0.85291,-1.33035) (0.96333,-1.48356) (1.07490,-1.63355) (1.18768,-1.78022) (1.30170,-1.92351) (1.41700,-2.06334) (1.53357,-2.19965) (1.65141,-2.33239) (1.77050,-2.46152) (1.89080,-2.58701) (2.01226,-2.70884) (2.13480,-2.82700) (2.25836,-2.94151) (2.38283,-3.05238) (2.50811,-3.15964) (2.63410,-3.26334) (2.76066,-3.36354) (2.88769,-3.46031) (3.01503,-3.55374) (3.14257,-3.64391) (3.27016,-3.73093) (3.39768,-3.81492) (3.52500,-3.89599)};\node[teal,right] at (3.7,-3.8960) {$r=0.265$ m};\draw[orange,thick] plot coordinates {(0.00000,0.00000) (0.13113,-0.20092) (0.26204,-0.39968) (0.39258,-0.59637) (0.52262,-0.79106) (0.65201,-0.98386) (0.78062,-1.17484) (0.90834,-1.36408) (1.03506,-1.55166) (1.16067,-1.73767) (1.28510,-1.92216) (1.40827,-2.10520) (1.53013,-2.28686) (1.65063,-2.46719) (1.76976,-2.64623) (1.88750,-2.82401) (2.00386,-3.00056) (2.11886,-3.17591) (2.23254,-3.35006) (2.34495,-3.52300) (2.45616,-3.69472) (2.56625,-3.86521) (2.67531,-4.03443) (2.78345,-4.20234) (2.89078,-4.36889) (2.99742,-4.53401) (3.10350,-4.69763) (3.20916,-4.85969) (3.31454,-5.02009) (3.41977,-5.17874) (3.52500,-5.33554)};\draw[orange,thick] plot coordinates {(0.00000,-0.00000) (0.10387,-0.21635) (0.20796,-0.43061) (0.31242,-0.64269) (0.41738,-0.85252) (0.52299,-1.06001) (0.62938,-1.26507) (0.73666,-1.46762) (0.84494,-1.66758) (0.95433,-1.86488) (1.06490,-2.05944) (1.17673,-2.25120) (1.28987,-2.44010) (1.40437,-2.62609) (1.52024,-2.80913) (1.63750,-2.98917) (1.75614,-3.16619) (1.87614,-3.34018) (1.99746,-3.51113) (2.12005,-3.67903) (2.24384,-3.84390) (2.36875,-4.00576) (2.49469,-4.16464) (2.62155,-4.32060) (2.74922,-4.47366) (2.87758,-4.62391) (3.00650,-4.77141) (3.13584,-4.91623) (3.26546,-5.05846) (3.39523,-5.19820) (3.52500,-5.33554)};\node[orange,right] at (3.7,-5.3355) {$r=0.315$ m};
\end{tikzpicture}}%
\count0=\ht0
\advance\count0 by \dp0
\count2=\wd0
\count4=\dimexpr\textheight/100*38\relax
\count6=\linewidth
\ifnum\count0<1
\count0=1
\fi
\ifnum\count2<1
\count2=1
\fi
\divide\count0 by 4096
\divide\count2 by 4096
\divide\count4 by 4096
\divide\count6 by 4096
\ifnum\count0<1
\count0=1
\fi
\ifnum\count2<1
\count2=1
\fi
\ifnum\count4<1
\count4=1
\fi
\ifnum\count6<1
\count6=1
\fi
\count8=\count2
\ifnum\count0>\count4
\count8=\numexpr\count2*\count4/\count0\relax
\fi
\ifnum\count8>\count6
\resizebox{\linewidth}{!}{\begin{tikzpicture}
\draw[->,gray] (-.3,0)--(4.2,0) node[right] {$x$ (escala 0,75 cm por cm)};\draw[->,gray] (0,.4)--(0,-5.8) node[below] {$y$ (misma escala)};\draw[blue,thick] plot coordinates {(0.00000,0.00000) (0.12871,-0.11105) (0.25701,-0.21883) (0.38477,-0.32344) (0.51187,-0.42502) (0.63819,-0.52370) (0.76364,-0.61959) (0.88813,-0.71284) (1.01159,-0.80357) (1.13397,-0.89191) (1.25522,-0.97797) (1.37533,-1.06188) (1.49429,-1.14373) (1.61212,-1.22363) (1.72885,-1.30167) (1.84453,-1.37792) (1.95923,-1.45243) (2.07303,-1.52526) (2.18602,-1.59643) (2.29832,-1.66596) (2.41004,-1.73384) (2.52132,-1.80004) (2.63228,-1.86452) (2.74307,-1.92721) (2.85383,-1.98803) (2.96470,-2.04688) (3.07580,-2.10364) (3.18727,-2.15816) (3.29923,-2.21030) (3.41178,-2.25989) (3.52500,-2.30675)};\draw[blue,thick] plot coordinates {(0.00000,-0.00000) (0.10629,-0.13293) (0.21299,-0.26291) (0.32023,-0.38983) (0.42813,-0.51357) (0.53681,-0.63400) (0.64636,-0.75098) (0.75687,-0.86439) (0.86841,-0.97409) (0.98103,-1.07997) (1.09478,-1.18191) (1.20967,-1.27977) (1.32571,-1.37347) (1.44288,-1.46291) (1.56115,-1.54798) (1.68047,-1.62863) (1.80077,-1.70479) (1.92197,-1.77641) (2.04398,-1.84347) (2.16668,-1.90595) (2.28996,-1.96387) (2.41368,-2.01724) (2.53772,-2.06611) (2.66193,-2.11054) (2.78617,-2.15063) (2.91030,-2.18647) (3.03420,-2.21818) (3.15773,-2.24591) (3.28077,-2.26980) (3.40322,-2.29002) (3.52500,-2.30675)};\node[blue,right] at (3.7,-2.3067) {$r=0.215$ m};\draw[teal,thick] plot coordinates {(0.00000,0.00000) (0.13026,-0.15738) (0.26025,-0.31215) (0.38983,-0.46441) (0.51885,-0.61426) (0.64721,-0.76181) (0.77477,-0.90715) (0.90143,-1.05040) (1.02709,-1.19163) (1.15167,-1.33095) (1.27510,-1.46846) (1.39732,-1.60423) (1.51830,-1.73834) (1.63800,-1.87086) (1.75643,-2.00185) (1.87359,-2.13137) (1.98950,-2.25945) (2.10420,-2.38612) (2.21774,-2.51141) (2.33020,-2.63531) (2.44164,-2.75783) (2.55217,-2.87894) (2.66189,-2.99860) (2.77090,-3.11678) (2.87934,-3.23341) (2.98731,-3.34843) (3.09497,-3.46175) (3.20243,-3.57327) (3.30984,-3.68290) (3.41732,-3.79051) (3.52500,-3.89599)};\draw[teal,thick] plot coordinates {(0.00000,-0.00000) (0.10474,-0.17553) (0.20975,-0.34863) (0.31517,-0.51918) (0.42115,-0.68710) (0.52779,-0.85228) (0.63523,-1.01461) (0.74357,-1.17400) (0.85291,-1.33035) (0.96333,-1.48356) (1.07490,-1.63355) (1.18768,-1.78022) (1.30170,-1.92351) (1.41700,-2.06334) (1.53357,-2.19965) (1.65141,-2.33239) (1.77050,-2.46152) (1.89080,-2.58701) (2.01226,-2.70884) (2.13480,-2.82700) (2.25836,-2.94151) (2.38283,-3.05238) (2.50811,-3.15964) (2.63410,-3.26334) (2.76066,-3.36354) (2.88769,-3.46031) (3.01503,-3.55374) (3.14257,-3.64391) (3.27016,-3.73093) (3.39768,-3.81492) (3.52500,-3.89599)};\node[teal,right] at (3.7,-3.8960) {$r=0.265$ m};\draw[orange,thick] plot coordinates {(0.00000,0.00000) (0.13113,-0.20092) (0.26204,-0.39968) (0.39258,-0.59637) (0.52262,-0.79106) (0.65201,-0.98386) (0.78062,-1.17484) (0.90834,-1.36408) (1.03506,-1.55166) (1.16067,-1.73767) (1.28510,-1.92216) (1.40827,-2.10520) (1.53013,-2.28686) (1.65063,-2.46719) (1.76976,-2.64623) (1.88750,-2.82401) (2.00386,-3.00056) (2.11886,-3.17591) (2.23254,-3.35006) (2.34495,-3.52300) (2.45616,-3.69472) (2.56625,-3.86521) (2.67531,-4.03443) (2.78345,-4.20234) (2.89078,-4.36889) (2.99742,-4.53401) (3.10350,-4.69763) (3.20916,-4.85969) (3.31454,-5.02009) (3.41977,-5.17874) (3.52500,-5.33554)};\draw[orange,thick] plot coordinates {(0.00000,-0.00000) (0.10387,-0.21635) (0.20796,-0.43061) (0.31242,-0.64269) (0.41738,-0.85252) (0.52299,-1.06001) (0.62938,-1.26507) (0.73666,-1.46762) (0.84494,-1.66758) (0.95433,-1.86488) (1.06490,-2.05944) (1.17673,-2.25120) (1.28987,-2.44010) (1.40437,-2.62609) (1.52024,-2.80913) (1.63750,-2.98917) (1.75614,-3.16619) (1.87614,-3.34018) (1.99746,-3.51113) (2.12005,-3.67903) (2.24384,-3.84390) (2.36875,-4.00576) (2.49469,-4.16464) (2.62155,-4.32060) (2.74922,-4.47366) (2.87758,-4.62391) (3.00650,-4.77141) (3.13584,-4.91623) (3.26546,-5.05846) (3.39523,-5.19820) (3.52500,-5.33554)};\node[orange,right] at (3.7,-5.3355) {$r=0.315$ m};
\end{tikzpicture}}%
\else
\ifnum\count0>\count4
\resizebox{!}{\dimexpr\textheight/100*38\relax}{\begin{tikzpicture}
\draw[->,gray] (-.3,0)--(4.2,0) node[right] {$x$ (escala 0,75 cm por cm)};\draw[->,gray] (0,.4)--(0,-5.8) node[below] {$y$ (misma escala)};\draw[blue,thick] plot coordinates {(0.00000,0.00000) (0.12871,-0.11105) (0.25701,-0.21883) (0.38477,-0.32344) (0.51187,-0.42502) (0.63819,-0.52370) (0.76364,-0.61959) (0.88813,-0.71284) (1.01159,-0.80357) (1.13397,-0.89191) (1.25522,-0.97797) (1.37533,-1.06188) (1.49429,-1.14373) (1.61212,-1.22363) (1.72885,-1.30167) (1.84453,-1.37792) (1.95923,-1.45243) (2.07303,-1.52526) (2.18602,-1.59643) (2.29832,-1.66596) (2.41004,-1.73384) (2.52132,-1.80004) (2.63228,-1.86452) (2.74307,-1.92721) (2.85383,-1.98803) (2.96470,-2.04688) (3.07580,-2.10364) (3.18727,-2.15816) (3.29923,-2.21030) (3.41178,-2.25989) (3.52500,-2.30675)};\draw[blue,thick] plot coordinates {(0.00000,-0.00000) (0.10629,-0.13293) (0.21299,-0.26291) (0.32023,-0.38983) (0.42813,-0.51357) (0.53681,-0.63400) (0.64636,-0.75098) (0.75687,-0.86439) (0.86841,-0.97409) (0.98103,-1.07997) (1.09478,-1.18191) (1.20967,-1.27977) (1.32571,-1.37347) (1.44288,-1.46291) (1.56115,-1.54798) (1.68047,-1.62863) (1.80077,-1.70479) (1.92197,-1.77641) (2.04398,-1.84347) (2.16668,-1.90595) (2.28996,-1.96387) (2.41368,-2.01724) (2.53772,-2.06611) (2.66193,-2.11054) (2.78617,-2.15063) (2.91030,-2.18647) (3.03420,-2.21818) (3.15773,-2.24591) (3.28077,-2.26980) (3.40322,-2.29002) (3.52500,-2.30675)};\node[blue,right] at (3.7,-2.3067) {$r=0.215$ m};\draw[teal,thick] plot coordinates {(0.00000,0.00000) (0.13026,-0.15738) (0.26025,-0.31215) (0.38983,-0.46441) (0.51885,-0.61426) (0.64721,-0.76181) (0.77477,-0.90715) (0.90143,-1.05040) (1.02709,-1.19163) (1.15167,-1.33095) (1.27510,-1.46846) (1.39732,-1.60423) (1.51830,-1.73834) (1.63800,-1.87086) (1.75643,-2.00185) (1.87359,-2.13137) (1.98950,-2.25945) (2.10420,-2.38612) (2.21774,-2.51141) (2.33020,-2.63531) (2.44164,-2.75783) (2.55217,-2.87894) (2.66189,-2.99860) (2.77090,-3.11678) (2.87934,-3.23341) (2.98731,-3.34843) (3.09497,-3.46175) (3.20243,-3.57327) (3.30984,-3.68290) (3.41732,-3.79051) (3.52500,-3.89599)};\draw[teal,thick] plot coordinates {(0.00000,-0.00000) (0.10474,-0.17553) (0.20975,-0.34863) (0.31517,-0.51918) (0.42115,-0.68710) (0.52779,-0.85228) (0.63523,-1.01461) (0.74357,-1.17400) (0.85291,-1.33035) (0.96333,-1.48356) (1.07490,-1.63355) (1.18768,-1.78022) (1.30170,-1.92351) (1.41700,-2.06334) (1.53357,-2.19965) (1.65141,-2.33239) (1.77050,-2.46152) (1.89080,-2.58701) (2.01226,-2.70884) (2.13480,-2.82700) (2.25836,-2.94151) (2.38283,-3.05238) (2.50811,-3.15964) (2.63410,-3.26334) (2.76066,-3.36354) (2.88769,-3.46031) (3.01503,-3.55374) (3.14257,-3.64391) (3.27016,-3.73093) (3.39768,-3.81492) (3.52500,-3.89599)};\node[teal,right] at (3.7,-3.8960) {$r=0.265$ m};\draw[orange,thick] plot coordinates {(0.00000,0.00000) (0.13113,-0.20092) (0.26204,-0.39968) (0.39258,-0.59637) (0.52262,-0.79106) (0.65201,-0.98386) (0.78062,-1.17484) (0.90834,-1.36408) (1.03506,-1.55166) (1.16067,-1.73767) (1.28510,-1.92216) (1.40827,-2.10520) (1.53013,-2.28686) (1.65063,-2.46719) (1.76976,-2.64623) (1.88750,-2.82401) (2.00386,-3.00056) (2.11886,-3.17591) (2.23254,-3.35006) (2.34495,-3.52300) (2.45616,-3.69472) (2.56625,-3.86521) (2.67531,-4.03443) (2.78345,-4.20234) (2.89078,-4.36889) (2.99742,-4.53401) (3.10350,-4.69763) (3.20916,-4.85969) (3.31454,-5.02009) (3.41977,-5.17874) (3.52500,-5.33554)};\draw[orange,thick] plot coordinates {(0.00000,-0.00000) (0.10387,-0.21635) (0.20796,-0.43061) (0.31242,-0.64269) (0.41738,-0.85252) (0.52299,-1.06001) (0.62938,-1.26507) (0.73666,-1.46762) (0.84494,-1.66758) (0.95433,-1.86488) (1.06490,-2.05944) (1.17673,-2.25120) (1.28987,-2.44010) (1.40437,-2.62609) (1.52024,-2.80913) (1.63750,-2.98917) (1.75614,-3.16619) (1.87614,-3.34018) (1.99746,-3.51113) (2.12005,-3.67903) (2.24384,-3.84390) (2.36875,-4.00576) (2.49469,-4.16464) (2.62155,-4.32060) (2.74922,-4.47366) (2.87758,-4.62391) (3.00650,-4.77141) (3.13584,-4.91623) (3.26546,-5.05846) (3.39523,-5.19820) (3.52500,-5.33554)};\node[orange,right] at (3.7,-5.3355) {$r=0.315$ m};
\end{tikzpicture}}%
\else
\begin{tikzpicture}
\draw[->,gray] (-.3,0)--(4.2,0) node[right] {$x$ (escala 0,75 cm por cm)};\draw[->,gray] (0,.4)--(0,-5.8) node[below] {$y$ (misma escala)};\draw[blue,thick] plot coordinates {(0.00000,0.00000) (0.12871,-0.11105) (0.25701,-0.21883) (0.38477,-0.32344) (0.51187,-0.42502) (0.63819,-0.52370) (0.76364,-0.61959) (0.88813,-0.71284) (1.01159,-0.80357) (1.13397,-0.89191) (1.25522,-0.97797) (1.37533,-1.06188) (1.49429,-1.14373) (1.61212,-1.22363) (1.72885,-1.30167) (1.84453,-1.37792) (1.95923,-1.45243) (2.07303,-1.52526) (2.18602,-1.59643) (2.29832,-1.66596) (2.41004,-1.73384) (2.52132,-1.80004) (2.63228,-1.86452) (2.74307,-1.92721) (2.85383,-1.98803) (2.96470,-2.04688) (3.07580,-2.10364) (3.18727,-2.15816) (3.29923,-2.21030) (3.41178,-2.25989) (3.52500,-2.30675)};\draw[blue,thick] plot coordinates {(0.00000,-0.00000) (0.10629,-0.13293) (0.21299,-0.26291) (0.32023,-0.38983) (0.42813,-0.51357) (0.53681,-0.63400) (0.64636,-0.75098) (0.75687,-0.86439) (0.86841,-0.97409) (0.98103,-1.07997) (1.09478,-1.18191) (1.20967,-1.27977) (1.32571,-1.37347) (1.44288,-1.46291) (1.56115,-1.54798) (1.68047,-1.62863) (1.80077,-1.70479) (1.92197,-1.77641) (2.04398,-1.84347) (2.16668,-1.90595) (2.28996,-1.96387) (2.41368,-2.01724) (2.53772,-2.06611) (2.66193,-2.11054) (2.78617,-2.15063) (2.91030,-2.18647) (3.03420,-2.21818) (3.15773,-2.24591) (3.28077,-2.26980) (3.40322,-2.29002) (3.52500,-2.30675)};\node[blue,right] at (3.7,-2.3067) {$r=0.215$ m};\draw[teal,thick] plot coordinates {(0.00000,0.00000) (0.13026,-0.15738) (0.26025,-0.31215) (0.38983,-0.46441) (0.51885,-0.61426) (0.64721,-0.76181) (0.77477,-0.90715) (0.90143,-1.05040) (1.02709,-1.19163) (1.15167,-1.33095) (1.27510,-1.46846) (1.39732,-1.60423) (1.51830,-1.73834) (1.63800,-1.87086) (1.75643,-2.00185) (1.87359,-2.13137) (1.98950,-2.25945) (2.10420,-2.38612) (2.21774,-2.51141) (2.33020,-2.63531) (2.44164,-2.75783) (2.55217,-2.87894) (2.66189,-2.99860) (2.77090,-3.11678) (2.87934,-3.23341) (2.98731,-3.34843) (3.09497,-3.46175) (3.20243,-3.57327) (3.30984,-3.68290) (3.41732,-3.79051) (3.52500,-3.89599)};\draw[teal,thick] plot coordinates {(0.00000,-0.00000) (0.10474,-0.17553) (0.20975,-0.34863) (0.31517,-0.51918) (0.42115,-0.68710) (0.52779,-0.85228) (0.63523,-1.01461) (0.74357,-1.17400) (0.85291,-1.33035) (0.96333,-1.48356) (1.07490,-1.63355) (1.18768,-1.78022) (1.30170,-1.92351) (1.41700,-2.06334) (1.53357,-2.19965) (1.65141,-2.33239) (1.77050,-2.46152) (1.89080,-2.58701) (2.01226,-2.70884) (2.13480,-2.82700) (2.25836,-2.94151) (2.38283,-3.05238) (2.50811,-3.15964) (2.63410,-3.26334) (2.76066,-3.36354) (2.88769,-3.46031) (3.01503,-3.55374) (3.14257,-3.64391) (3.27016,-3.73093) (3.39768,-3.81492) (3.52500,-3.89599)};\node[teal,right] at (3.7,-3.8960) {$r=0.265$ m};\draw[orange,thick] plot coordinates {(0.00000,0.00000) (0.13113,-0.20092) (0.26204,-0.39968) (0.39258,-0.59637) (0.52262,-0.79106) (0.65201,-0.98386) (0.78062,-1.17484) (0.90834,-1.36408) (1.03506,-1.55166) (1.16067,-1.73767) (1.28510,-1.92216) (1.40827,-2.10520) (1.53013,-2.28686) (1.65063,-2.46719) (1.76976,-2.64623) (1.88750,-2.82401) (2.00386,-3.00056) (2.11886,-3.17591) (2.23254,-3.35006) (2.34495,-3.52300) (2.45616,-3.69472) (2.56625,-3.86521) (2.67531,-4.03443) (2.78345,-4.20234) (2.89078,-4.36889) (2.99742,-4.53401) (3.10350,-4.69763) (3.20916,-4.85969) (3.31454,-5.02009) (3.41977,-5.17874) (3.52500,-5.33554)};\draw[orange,thick] plot coordinates {(0.00000,-0.00000) (0.10387,-0.21635) (0.20796,-0.43061) (0.31242,-0.64269) (0.41738,-0.85252) (0.52299,-1.06001) (0.62938,-1.26507) (0.73666,-1.46762) (0.84494,-1.66758) (0.95433,-1.86488) (1.06490,-2.05944) (1.17673,-2.25120) (1.28987,-2.44010) (1.40437,-2.62609) (1.52024,-2.80913) (1.63750,-2.98917) (1.75614,-3.16619) (1.87614,-3.34018) (1.99746,-3.51113) (2.12005,-3.67903) (2.24384,-3.84390) (2.36875,-4.00576) (2.49469,-4.16464) (2.62155,-4.32060) (2.74922,-4.47366) (2.87758,-4.62391) (3.00650,-4.77141) (3.13584,-4.91623) (3.26546,-5.05846) (3.39523,-5.19820) (3.52500,-5.33554)};\node[orange,right] at (3.7,-5.3355) {$r=0.315$ m};
\end{tikzpicture}%
\fi
\fi
\endgroup
\par\nopagebreak
\studiumcaption{fig:practica}{Secciones calculadas con la familia de Hermite del ejemplo reproducible, rotor compresor, radios 0,215; 0,265 y 0,315 m. Escala igual en ambos ejes.}
\end{minipage}
\par\vspace{\baselineskip}
% studium:end-record eyJpZCI6ICJGSUctMDAyMSIsICJzZWN0aW9uIjogImMxOCJ9

\section{Coordenadas y fórmulas de hoja}

Una estructura útil tiene hojas de entradas, estaciones, secciones, coordenadas y comprobaciones. Cada coordenada depende únicamente de entradas y fórmulas; los resultados no se introducen como valores manuales. Se convierte cada ángulo de grados a radianes antes de aplicar funciones trigonométricas. La función de dos argumentos debe verificarse en la aplicación usada porque el orden de sus argumentos varía entre productos.
La columna de \(\xi\) recorre de cero a uno. Se evalúan \(f\), \(f'\), \(t_n\) y las dos superficies. Se grafica una dispersión XY con líneas, no un gráfico de categorías. Los ejes deben tener igual escala cuando se quiera interpretar la geometría. El paso se representa desplazando una copia de la sección en la dirección tangencial desarrollada.
Un cambio en velocidad de giro modifica las betas de las filas móviles; las alfas no cambian si se mantienen las velocidades absolutas prescritas. Este es un control funcional importante. Si todos los gráficos cambian de la misma forma, es posible que se haya utilizado la referencia incorrecta.
% studium:end-record eyJpZCI6ICJQQVItMDEyNiIsICJzZWN0aW9uIjogImMxOCJ9

\section{Validación y convergencia}

Se comprueban extremos de la línea media, tangentes, espesor máximo, ausencia de puntos no finitos y consistencia de unidades. La interpolación se refina duplicando el número de puntos y comparando coordenadas a igual parámetro. Para la familia analítica la evaluación puntual es exacta salvo redondeo; la discretización afecta la visualización y cualquier estimación de área o curvatura por diferencias.
Se comparan los ángulos de raíz y punta con los triángulos del modelo radial. Se comprueba paso y solidez en cada radio, manteniendo un número entero de álabes. Los autocruces de superficies o la intersección con perfiles vecinos invalidan una geometría aunque todas las celdas contengan números.
La geometría educativa termina en bordes de espesor cero. Un trabajo de diseño real debería introducir radios, espesor de salida, restricciones de fabricación y verificación aerodinámica. El manual conserva ese punto como revisión pendiente; no atribuye aptitud a una superficie por haber sido dibujada correctamente.
% studium:end-record eyJpZCI6ICJQQVItMDEyNyIsICJzZWN0aW9uIjogImMxOCJ9

\section{Interpretación y documentación del trabajo}

Un informe útil explica por qué cambian las secciones con el radio y qué hipótesis originan la torsión. Debe incluir un manual de entradas con unidades, valores admisibles, mensajes de error y una prueba reproducible. La contribución de cada miembro del grupo se documenta con tareas realmente realizadas, sin firmas ni aportaciones inventadas.
Como ejercicio autónomo, el lector puede cambiar la ley radial o imponer una distribución de desviación y comparar el resultado. Debe justificar qué propiedades conserva y volver a comprobar equilibrio, caudal y trabajo si el nuevo objetivo pretende ser aerodinámico. Cambiar una fórmula geométrica no actualiza por sí solo todo el modelo físico.
Errores habituales: usar un gráfico con ejes de distinta escala para medir ángulos, mezclar filas móviles y fijas, confundir una tabla de coordenadas con una validación aerodinámica y ocultar parámetros ajustados manualmente. Autocomprobación: anule \(\Omega\) y compruebe que velocidades absoluta y relativa coinciden en cada estación.
Formulario: leyes radiales y familia de Hermite de los capítulos \(11-12\). Referencias: enunciado oficial de trabajo, p. 1; ULE Esfuerzos y refrigeración, pp. \(2-7\); Diseño de compresores, pp. \(28-35\). La familia de perfiles, los datos, los programas y todas las figuras de práctica son originales.
% studium:end-record eyJpZCI6ICJQQVItMDEyOCIsICJzZWN0aW9uIjogImMxOCJ9

\begin{tcolorbox}[title={Desarrollos del modelo}, studiumdef, breakable]

\noindent\textbf{Modelo desarrollado: Rotor compresor a radio elegido}

\begin{equation}
\kappa=-i+\arctan\frac{K_v/r-\Omega\,r}{C_x}
\label{eq:EC-C18}
\end{equation}

\noindent\textit{Significado.} Ángulo metálico en función de la incidencia, el vórtice libre y la velocidad de giro.

\noindent\textbf{Hipótesis}
\begin{itemize}
\item Vórtice libre por estación, incidencia prescrita.
\end{itemize}

\noindent\textbf{Principios}
\begin{itemize}
\item Balance y definiciones físicas desarrollados en el capítulo; propiedades y dominio declarados en el problema.
\end{itemize}

\noindent\textbf{Aplicabilidad}
\begin{itemize}
\item Modelo educativo del caso C18-01-R
\end{itemize}

\noindent\textbf{Limitaciones}
\begin{itemize}
\item Modelo educativo; la aplicabilidad a una máquina real requiere revisión experta.
\end{itemize}

\noindent\textbf{Pasos}
\begin{enumerate}
\item
La velocidad tangencial absoluta sigue vórtice libre.

\[
\fitmath{C_\theta=\frac{K_e}{r}}
\]
\item
Se forma el ángulo del gas en referencia rotor.

\[
\fitmath{\beta=\operatorname{atan}{\left(\frac{- \Omega r + \frac{K_e}{r}}{C_x} \right)}}
\]
\item
Se despeja metal a partir de incidencia definida.

\[
\fitmath{\kappa_e=- i + \operatorname{atan}{\left(\frac{- \Omega r + \frac{K_e}{r}}{C_x} \right)}}
\]
\end{enumerate}

{\noindent\footnotesize verificada dimensionalmente.\par}
\end{tcolorbox}

\begin{tcolorbox}[title={Problema resuelto 1}, studiumworked, breakable]

\noindent\textbf{Enunciado}
\(C18-01-R \cdot\) Rotor compresor a radio elegido. Aplique la plantilla radial educativa a la entrada de un rotor compresor. Calcule U, beta y ángulo metálico con incidencia prescrita. Datos: \(r=0.271\, \mathrm{m}\); \(\Omega=1010\, \mathrm{s}^{-1}\); \(K_e=14\, \mathrm{m}^{2}\,\mathrm{s}^{-1}\); \(C_x=159\, \mathrm{m}\,\mathrm{s}^{-1}\); \(i=-0.01745329\, \mathrm{rad}\).

\noindent\textbf{Datos}
\begin{itemize}
\item $r = 0.271\,\mathrm{m}$ — radio elegido
\item $\Omega = 1010\,\mathrm{s}^{-1}$ — velocidad angular
\item $K_e = 14\,\mathrm{m}^{2}\,\mathrm{s}^{-1}$ — constante remolino de entrada
\item $C_x = 159\,\mathrm{m}\,\mathrm{s}^{-1}$ — velocidad axial
\item $i = -0.017453292519943\,\mathrm{rad}$ — incidencia algebraica prescrita
\end{itemize}

\noindent\textbf{Incógnita}
$\kappa_e$, en $\mathrm{rad}$; también las magnitudes intermedias desarrolladas.

\noindent\textbf{Modelo físico}

\noindent\textit{Rotor compresor a radio elegido}

\[
\fitmath{\kappa_e=- i + \operatorname{atan}{\left(\frac{- \Omega r + \frac{K_e}{r}}{C_x} \right)}}
\]

\noindent\textbf{Hipótesis}
\begin{itemize}
\item Vórtice libre por estación, incidencia prescrita.
\end{itemize}

\noindent\textbf{Desarrollo paso a paso}

\begin{enumerate}

\item Velocidad periférica

\[\fitmath{U=\Omega\cdot \mathrm{rad}}\]

\[\fitmath{U=(1010)\cdot (0.271)\approx 273.71\,\mathrm{m/s}}\]

\item La velocidad tangencial absoluta sigue vórtice libre.

\[\fitmath{\mathrm{ct}=\frac{\mathrm{kv}}{\mathrm{rad}}}\]

\[\fitmath{\mathrm{ct}=\frac{(14)}{(0.271)}\approx 51.66052\,\mathrm{m/s}}\]

\item Se forma el ángulo del gas en referencia rotor.

\[\fitmath{\beta=\arctan\left(\frac{\frac{\mathrm{kv}}{\mathrm{rad}}-\Omega\cdot \mathrm{rad}}{\mathrm{cx}}\right)}\]

\[\fitmath{\beta=\arctan\left(\frac{\frac{(14)}{(0.271)}-(1010)\cdot (0.271)}{(159)}\right)\approx -0.9493752\,\mathrm{rad}}\]

\item Se despeja metal a partir de incidencia definida.

\[\fitmath{\mathrm{kap}=\arctan\left(\frac{\frac{\mathrm{kv}}{\mathrm{rad}}-\Omega\cdot \mathrm{rad}}{\mathrm{cx}}\right)-\mathrm{inc}}\]

\[\fitmath{\mathrm{kap}=\arctan\left(\frac{\frac{(14)}{(0.271)}-(1010)\cdot (0.271)}{(159)}\right)-(-0.01745329252)\approx -0.9319219\,\mathrm{rad}}\]

\end{enumerate}


\noindent\textbf{Sustitución numérica}

\[
\fitmath{- \left(-0.017453292519943\,\mathrm{rad}\right) + \operatorname{atan}{\left(\frac{- \left(1010\,\mathrm{s}^{-1}\right) \left(0.271\,\mathrm{m}\right) + \frac{\left(14\,\mathrm{m}^{2}\,\mathrm{s}^{-1}\right)}{\left(0.271\,\mathrm{m}\right)}}{\left(159\,\mathrm{m}\,\mathrm{s}^{-1}\right)} \right)}}
\]

\noindent\textbf{Resultado}

\[
\fitmath{\kappa_e=-9.319219110\times10^{-1}\,\mathrm{rad}}
\]

\noindent\textbf{Velocidad periférica}
\[
\fitmath{U=\Omega r=273.71\,\mathrm{m/s}}
\]

\noindent\textbf{Comprobación cinemática.} Insertar este $U$ en $\beta=\arctan[(C_\theta-U)/C_x]$ reproduce el ángulo desarrollado.

\noindent\textbf{Interpretación}
La misma tabla debe actualizarse al variar r u Omega. El metal es una propuesta geométrica educativa.

\noindent\textbf{Comprobación de los pasos}
\begin{itemize}
\item Paso 2: $51.6605$ — reproducido
\item Paso 4: $-0.949375$ — reproducido
\item Paso 6: $-0.931922$ — reproducido
\end{itemize}

\noindent\textbf{Comprobación de la respuesta} la aritmética del último paso reproduce la respuesta enunciada a la precisión con que está escrita. Es una reproducción aritmética; no valida el modelo físico ni sus hipótesis.

\noindent\textbf{Limitaciones}
\begin{itemize}
\item Modelo educativo; la aplicabilidad a una máquina real requiere revisión experta.
\end{itemize}

\noindent\textbf{Errores frecuentes}
\begin{itemize}
\item Usar alfa para construir un rotor.
\end{itemize}
% studium:end-record eyJpZCI6ICJQUkItMDEzNyIsICJzZWN0aW9uIjogImMxOCJ9
\end{tcolorbox}

\begin{tcolorbox}[title={Problema resuelto 2}, studiumworked, breakable]

\noindent\textbf{Enunciado}
\(C18-02-R \cdot\) Estator compresor y orientación de salida. Con remolino de salida prescrito, halle el ángulo absoluto y el ángulo metálico para un estator compresor. Datos: \(r=0.282\, \mathrm{m}\); \(K_s=9\, \mathrm{m}^{2}\,\mathrm{s}^{-1}\); \(C_x=162\, \mathrm{m}\,\mathrm{s}^{-1}\); \(\delta=0.03839724\, \mathrm{rad}\).

\noindent\textbf{Datos}
\begin{itemize}
\item $r = 0.282\,\mathrm{m}$ — radio
\item $K_s = 9\,\mathrm{m}^{2}\,\mathrm{s}^{-1}$ — constante remolino salida
\item $C_x = 162\,\mathrm{m}\,\mathrm{s}^{-1}$ — velocidad axial
\item $\delta = 0.038397243543875\,\mathrm{rad}$ — desviación algebraica prescrita
\end{itemize}

\noindent\textbf{Incógnita}
$\kappa_e$, en $\mathrm{rad}$; también las magnitudes intermedias desarrolladas.

\noindent\textbf{Modelo físico}

\noindent\textit{Estator compresor y orientación de salida}

\[
\fitmath{\kappa_e=- \delta + \operatorname{atan}{\left(\frac{K_s}{C_x r} \right)}}
\]

\noindent\textbf{Hipótesis}
\begin{itemize}
\item Velocidades absolutas prescritas, desviación como dato.
\end{itemize}

\noindent\textbf{Desarrollo paso a paso}

\begin{enumerate}

\item La fila fija utiliza velocidad absoluta.

\[\fitmath{\alpha=\arctan\left(\frac{\mathrm{kv}}{\mathrm{rad}\cdot \mathrm{cx}}\right)}\]

\[\fitmath{\alpha=\arctan\left(\frac{(9)}{(0.282)\cdot (162)}\right)\approx 0.1945146\,\mathrm{rad}}\]

\item La definición delta=gas-metal permite despejar metal.

\[\fitmath{\mathrm{kap}=\arctan\left(\frac{\mathrm{kv}}{\mathrm{rad}\cdot \mathrm{cx}}\right)-\mathrm{dev}}\]

\[\fitmath{\mathrm{kap}=\arctan\left(\frac{(9)}{(0.282)\cdot (162)}\right)-(0.03839724354)\approx 0.1561174\,\mathrm{rad}}\]

\end{enumerate}


\noindent\textbf{Sustitución numérica}

\[
\fitmath{- \left(0.038397243543875\,\mathrm{rad}\right) + \operatorname{atan}{\left(\frac{\left(9\,\mathrm{m}^{2}\,\mathrm{s}^{-1}\right)}{\left(162\,\mathrm{m}\,\mathrm{s}^{-1}\right) \left(0.282\,\mathrm{m}\right)} \right)}}
\]

\noindent\textbf{Resultado}

\[
\fitmath{\kappa_e=1.561173538\times10^{-1}\,\mathrm{rad}}
\]

\noindent\textbf{Interpretación}
No aparece Omega porque la fila está fija; solo cambiaría el resultado si el cambio de operación modifica los datos absolutos.

\noindent\textbf{Comprobación de los pasos}
\begin{itemize}
\item Paso 2: $0.194515$ — reproducido
\item Paso 4: $0.156117$ — reproducido
\end{itemize}

\noindent\textbf{Comprobación de la respuesta} la aritmética del último paso reproduce la respuesta enunciada a la precisión con que está escrita. Es una reproducción aritmética; no valida el modelo físico ni sus hipótesis.

\noindent\textbf{Limitaciones}
\begin{itemize}
\item Modelo educativo; la aplicabilidad a una máquina real requiere revisión experta.
\end{itemize}

\noindent\textbf{Errores frecuentes}
\begin{itemize}
\item Restar U a una fila fija.
\end{itemize}
% studium:end-record eyJpZCI6ICJQUkItMDEzOSIsICJzZWN0aW9uIjogImMxOCJ9
\end{tcolorbox}

\begin{tcolorbox}[title={Problema resuelto 3}, studiumworked, breakable]

\noindent\textbf{Enunciado}
\(C18-03-R \cdot\) Estator turbina y calado educativo. Una sección educativa del estator de turbina tiene pendientes metálicas dadas. Con d igual al promedio de pendientes, calcule calado y cuerda real. Datos: \(m_e=0.12\); \(m_s=1.38\); \(c_x=0.046\, \mathrm{m}\).

\noindent\textbf{Datos}
\begin{itemize}
\item $m_e = 0.12$ — pendiente metal entrada
\item $m_s = 1.38$ — pendiente metal salida
\item $c_x = 0.046\,\mathrm{m}$ — cuerda axial
\end{itemize}

\noindent\textbf{Incógnita}
$c$, en $\mathrm{m}$; también las magnitudes intermedias desarrolladas.

\noindent\textbf{Modelo físico}

\noindent\textit{Estator turbina y calado educativo}

\[
\fitmath{c=c_x \sqrt{\left(\frac{m_e}{2} + \frac{m_s}{2}\right)^{2} + 1}}
\]

\noindent\textbf{Hipótesis}
\begin{itemize}
\item Línea media Hermite educativa; x monótona.
\end{itemize}

\noindent\textbf{Desarrollo paso a paso}

\begin{enumerate}

\item El desplazamiento normalizado se prescribe como promedio de pendientes.

\[\fitmath{\mathrm{dd}=\frac{\mathrm{me}+\mathrm{ms}}{2}}\]

\[\fitmath{\mathrm{dd}=\frac{(0.12)+(1.38)}{2}\approx 0.75\,}\]

\item El calado de cuerda es atan d.

\[\fitmath{\chi=\arctan\left(\frac{\mathrm{me}+\mathrm{ms}}{2}\right)}\]

\[\fitmath{\chi=\arctan\left(\frac{(0.12)+(1.38)}{2}\right)\approx 0.6435011\,\mathrm{rad}}\]

\item La cuerda real se deduce del desplazamiento final.

\[\fitmath{\mathrm{chord}=\mathrm{cxg}\cdot \sqrt{1+{\left(\frac{\mathrm{me}+\mathrm{ms}}{2}\right)}^{2}}}\]

\[\fitmath{\mathrm{chord}=(0.046)\cdot \sqrt{1+{\left(\frac{(0.12)+(1.38)}{2}\right)}^{2}}\approx 0.0575\,\mathrm{m}}\]

\end{enumerate}


\noindent\textbf{Sustitución numérica}

\[
\fitmath{\left(0.046\,\mathrm{m}\right) \sqrt{\left(\frac{\left(0.12\right)}{2} + \frac{\left(1.38\right)}{2}\right)^{2} + 1}}
\]

\noindent\textbf{Resultado}

\[
\fitmath{c=5.750000000\times10^{-2}\,\mathrm{m}}
\]

\noindent\textbf{Interpretación}
Promediar pendientes no equivale a promediar ángulos. La ley es geométrica y no una correlación de mínima pérdida.

\noindent\textbf{Comprobación de los pasos}
\begin{itemize}
\item Paso 2: $0.75$ — reproducido
\item Paso 4: $0.643501$ — reproducido
\item Paso 6: $0.0575$ — reproducido
\end{itemize}

\noindent\textbf{Comprobación de la respuesta} la aritmética del último paso reproduce la respuesta enunciada a la precisión con que está escrita. Es una reproducción aritmética; no valida el modelo físico ni sus hipótesis.

\noindent\textbf{Limitaciones}
\begin{itemize}
\item Modelo educativo; la aplicabilidad a una máquina real requiere revisión experta.
\end{itemize}

\noindent\textbf{Errores frecuentes}
\begin{itemize}
\item Presentar esta sección como un perfil de fabricante.
\end{itemize}
% studium:end-record eyJpZCI6ICJQUkItMDE0MSIsICJzZWN0aW9uIjogImMxOCJ9
\end{tcolorbox}

\begin{tcolorbox}[title={Problema resuelto 4}, studiumworked, breakable]

\noindent\textbf{Enunciado}
\(C18-04-R \cdot\) Rotor turbina y espesor por radio. En la salida del rotor de turbina, calcule beta y la coordenada transversal de la superficie superior en un punto de pendiente conocida. Datos: \(r=0.294\, \mathrm{m}\); \(\Omega=1120\, \mathrm{s}^{-1}\); \(K_s=-7\, \mathrm{m}^{2}\,\mathrm{s}^{-1}\); \(C_x=201\, \mathrm{m}\,\mathrm{s}^{-1}\); \(y=0.014\, \mathrm{m}\); \(t_n=0.002\, \mathrm{m}\); \(m_y=-0.62\).

\noindent\textbf{Datos}
\begin{itemize}
\item $r = 0.294\,\mathrm{m}$ — radio
\item $\Omega = 1120\,\mathrm{s}^{-1}$ — velocidad angular
\item $K_s = -7\,\mathrm{m}^{2}\,\mathrm{s}^{-1}$ — constante remolino salida
\item $C_x = 201\,\mathrm{m}\,\mathrm{s}^{-1}$ — velocidad axial
\item $y = 0.014\,\mathrm{m}$ — ordenada línea media en el punto
\item $t_n = 0.002\,\mathrm{m}$ — semiespesor normal
\item $m_y = -0.62$ — pendiente de línea media en el punto
\end{itemize}

\noindent\textbf{Incógnita}
$y_s$, en $\mathrm{m}$; también las magnitudes intermedias desarrolladas.

\noindent\textbf{Modelo físico}

\noindent\textit{Rotor turbina y espesor por radio}

\[
\fitmath{y_s=\frac{t_n}{\sqrt{m_y^{2} + 1}} + y}
\]

\noindent\textbf{Hipótesis}
\begin{itemize}
\item Vórtice libre y geometría normal local educativa.
\end{itemize}

\noindent\textbf{Desarrollo paso a paso}

\begin{enumerate}

\item El ángulo del gas se forma en referencia móvil.

\[\fitmath{\beta=\arctan\left(\frac{\frac{\mathrm{kv}}{\mathrm{rad}}-\Omega\cdot \mathrm{rad}}{\mathrm{cx}}\right)}\]

\[\fitmath{\beta=\arctan\left(\frac{\frac{(-7)}{(0.294)}-(1120)\cdot (0.294)}{(201)}\right)\approx -1.053286\,\mathrm{rad}}\]

\item La superficie superior desplaza la ordenada según la normal.

\[\fitmath{\mathrm{ys}=\mathrm{yy}+\frac{\mathrm{tt}}{\sqrt{1+{\mathrm{slope}}^{2}}}}\]

\[\fitmath{\mathrm{ys}=(0.014)+\frac{(0.002)}{\sqrt{1+{(-0.62)}^{2}}}\approx 0.01569981\,\mathrm{m}}\]

\end{enumerate}


\noindent\textbf{Sustitución numérica}

\[
\fitmath{\frac{\left(0.002\,\mathrm{m}\right)}{\sqrt{\left(-0.62\right)^{2} + 1}} + \left(0.014\,\mathrm{m}\right)}
\]

\noindent\textbf{Resultado}

\[
\fitmath{y_s=1.569980538\times10^{-2}\,\mathrm{m}}
\]

\noindent\textbf{Interpretación}
El dibujo completo debe desplazar también x. Esta salida tiene remolino negativo definido, sin cambiar el sentido de giro.

\noindent\textbf{Comprobación de los pasos}
\begin{itemize}
\item Paso 2: $-1.05329$ — reproducido
\item Paso 4: $0.0156998$ — reproducido
\end{itemize}

\noindent\textbf{Comprobación de la respuesta} la aritmética del último paso reproduce la respuesta enunciada a la precisión con que está escrita. Es una reproducción aritmética; no valida el modelo físico ni sus hipótesis.

\noindent\textbf{Limitaciones}
\begin{itemize}
\item Modelo educativo; la aplicabilidad a una máquina real requiere revisión experta.
\end{itemize}

\noindent\textbf{Errores frecuentes}
\begin{itemize}
\item Cambiar de signo U para representar una turbina.
\end{itemize}
% studium:end-record eyJpZCI6ICJQUkItMDE0MyIsICJzZWN0aW9uIjogImMxOCJ9
\end{tcolorbox}

\begin{tcolorbox}[title={Ejercicios propuestos}, studiumworked, breakable]

\noindent\textbf{1 (Básico, Desarrollo)}
\(C18-01-P \cdot\) Rotor compresor a radio elegido. Ejercicio de transferencia. Aplique la plantilla radial educativa a la entrada de un rotor compresor. Calcule U, beta y ángulo metálico con incidencia prescrita. Antes de consultar la solución, prediga el efecto de cambiar el primer dato respecto del caso resuelto de esta familia. Justifique la tendencia y explique qué restricción podría convertirse en limitante. Datos: \(r=0.288615\, \mathrm{m}\); \(\Omega=1010\, \mathrm{s}^{-1}\); \(K_e=14\, \mathrm{m}^{2}\,\mathrm{s}^{-1}\); \(C_x=159\, \mathrm{m}\,\mathrm{s}^{-1}\); \(i=-0.01745329\, \mathrm{rad}\).

\noindent\textbf{2 (Intermedio, Desarrollo)}
\(C18-02-P \cdot\) Estator compresor y orientación de salida. Ejercicio de transferencia. Con remolino de salida prescrito, halle el ángulo absoluto y el ángulo metálico para un estator compresor. Antes de consultar la solución, prediga el efecto de cambiar el primer dato respecto del caso resuelto de esta familia. Justifique la tendencia y explique qué restricción podría convertirse en limitante. Datos: \(r=0.30033\, \mathrm{m}\); \(K_s=9\, \mathrm{m}^{2}\,\mathrm{s}^{-1}\); \(C_x=162\, \mathrm{m}\,\mathrm{s}^{-1}\); \(\delta=0.03839724\, \mathrm{rad}\).

\noindent\textbf{3 (Avanzado, Desarrollo)}
\(C18-03-P \cdot\) Estator turbina y calado educativo. Ejercicio de transferencia. Una sección educativa del estator de turbina tiene pendientes metálicas dadas. Con d igual al promedio de pendientes, calcule calado y cuerda real. Antes de consultar la solución, prediga el efecto de cambiar el primer dato respecto del caso resuelto de esta familia. Justifique la tendencia y explique qué restricción podría convertirse en limitante. Datos: \(m_e=0.1278\); \(m_s=1.38\); \(c_x=0.046\, \mathrm{m}\).

\noindent\textbf{4 (Nivel de examen, Desarrollo)}
\(C18-04-P \cdot\) Rotor turbina y espesor por radio. Ejercicio de transferencia. En la salida del rotor de turbina, calcule beta y la coordenada transversal de la superficie superior en un punto de pendiente conocida. Antes de consultar la solución, prediga el efecto de cambiar el primer dato respecto del caso resuelto de esta familia. Justifique la tendencia y explique qué restricción podría convertirse en limitante. Datos: \(r=0.31311\, \mathrm{m}\); \(\Omega=1120\, \mathrm{s}^{-1}\); \(K_s=-7\, \mathrm{m}^{2}\,\mathrm{s}^{-1}\); \(C_x=201\, \mathrm{m}\,\mathrm{s}^{-1}\); \(y=0.014\, \mathrm{m}\); \(t_n=0.002\, \mathrm{m}\); \(m_y=-0.62\).
\end{tcolorbox}


\studiumchaptercheck

\chapter{Simulacros originales y estrategia de resolución}
\label{chap:c19}

\studiumchapterguide{Objetivo: desarrollar simulacros originales y estrategia de resolución, declarar hipótesis y comprobar condiciones físicas.}
% studium:end-record eyJpZCI6ICJQQVItMDEyOSIsICJzZWN0aW9uIjogImMxOSJ9

\section{Alcance y uso}

Estos dos simulacros son originales y no representan exámenes oficiales de la Universidad de León. Se construyen a partir de las competencias del temario y de las familias generales de cálculo identificadas en los materiales privados: escalones axiales, estados totales y estáticos, geometría radial y prestaciones de un compresor radial. Los datos, la secuencia de preguntas y las soluciones se han generado de forma independiente.
La guía 2026/2027, pp. \(1-2\), sitúa la asignatura en cuarto curso y establece resultados de aprendizaje de diseño y análisis de prestaciones de turbomáquinas axiales y radiales. Su evaluación combina pruebas mixtas y trabajos, según pp. \(3-4\). Los tiempos y puntuaciones que siguen son una propuesta de entrenamiento de este manual; no se atribuyen al profesorado.
% studium:end-record eyJpZCI6ICJQQVItMDEzMCIsICJzZWN0aW9uIjogImMxOSJ9

\section{Simulacro A: escalones y consistencia física}

Tiempo orientativo: 120 minutos. Tres problemas: A1 compresor axial, 4 puntos; A2 turbina axial, 4 puntos; A3 distribución radial y tensión, 2 puntos. Use las constantes suministradas y justifique cada hipótesis. Los ejercicios correspondientes son C19-01-P, C19-02-P y C19-03-P, en ese orden. Las soluciones se encuentran en el cuaderno separado para permitir un intento sin respuestas.
En A1 se debe cerrar trabajo, temperatura, presión y área. Una presión final numéricamente correcta con velocidad equivocada en continuidad no demuestra que el procedimiento sea correcto. En A2 se conserva el signo del remolino de salida y se distingue trabajo positivo extraído de trabajo comunicado al fluido. En A3 se conectan un modelo aerodinámico radial y un modelo mecánico nominal; no se deduce vida del material.
Criterio orientativo: 25 por ciento de cada problema por planteamiento y convenio, 35 por ciento por desarrollo y despejes, 25 por ciento por sustitución y unidades y 15 por ciento por interpretación y comprobaciones. Una inconsistencia física no se compensa por mantener muchas cifras decimales.
% studium:end-record eyJpZCI6ICJQQVItMDEzMSIsICJzZWN0aW9uIjogImMxOSJ9

\section{Simulacro B: diseño y límites}

Tiempo orientativo: 120 minutos. Tres problemas: B1 compresor radial, 4 puntos; B2 difusión y geometría de cascada, 3 puntos; B3 refrigeración, 3 puntos. Corresponden a C19-04-P, C19-05-P y C19-06-P. Deben resolverse desde sus datos, sin trasladar resultados del simulacro A.
B1 exige mantener radios distintos y una definición explícita de deslizamiento. B2 distingue un indicador de difusión de una validación aerodinámica y revisa el efecto del número entero de álabes. B3 separa balance de calor, temperatura de pared y mínimo energético de refrigerante. En todos los casos se pide indicar qué verificación requeriría datos adicionales o juicio experto.
% studium:end-record eyJpZCI6ICJQQVItMDEzMiIsICJzZWN0aW9uIjogImMxOSJ9

\section{Estrategia, errores y autoevaluación}

Antes de calcular, haga una tabla de datos y marque con claridad qué presiones y temperaturas son totales. Dibuje estaciones, sentido de giro y componentes tangenciales. Enumere incógnitas y ecuaciones para evitar sobredefinir el problema. Mantenga unidades coherentes y posponga el redondeo hasta presentar resultados.
Reserve una revisión final para signos, dimensiones y límites: temperatura estática positiva, eficiencias admisibles, áreas y radios positivos, continuidad y energía. Si una raíz matemática produce flujo invertido o una temperatura no admisible, no se selecciona por cercanía al resultado esperado. Se vuelve al dominio físico del modelo.
La autoevaluación debe distinguir un error de aritmética de un error de modelo. Anote qué cambió al corregirlo y vuelva a resolver una familia de transferencia. Consulte la solución desarrollada solo después de intentar la predicción de tendencias. Los problemas propuestos del resto del manual mantienen esa misma separación entre enunciado y solución.
Referencias: guía docente ULE 2026/2027, pp. \(1-4\); presentación y documentos técnicos aportados para la correspondencia temática. Los materiales de Wuolah se utilizaron únicamente para clasificar competencias de estudio. Ningún enunciado, dibujo ni solución de esos documentos se reproduce o transforma aquí.
% studium:end-record eyJpZCI6ICJQQVItMDEzMyIsICJzZWN0aW9uIjogImMxOSJ9

\begin{tcolorbox}[title={Ejercicios propuestos}, studiumworked, breakable]

\noindent\textbf{1 (Nivel de examen, Desarrollo)}
\(C19-01-P \cdot A1 \cdot\) Cierre de compresor axial. Un rotor compresor al radio medio recibe Cx y remolinos dados. Obtenga trabajo, temperatura total final, presión total final y área geométrica de entrada. Resultado final: área de entrada. Datos: \(T_{0e}=302\, \mathrm{K}\); \(p_{0e}=118000\, \mathrm{Pa}\); \(C_x=156\, \mathrm{m}\,\mathrm{s}^{-1}\); \(C_{\theta e}=38\, \mathrm{m}\,\mathrm{s}^{-1}\); \(C_{\theta s}=147\, \mathrm{m}\,\mathrm{s}^{-1}\); \(U=291\, \mathrm{m}\,\mathrm{s}^{-1}\); \(\dot m=21\, \mathrm{kg}\,\mathrm{s}^{-1}\); \(\eta=0.87\); \(B=0.96\); \(c_p=1004.5\, \mathrm{J}\,\mathrm{kg}^{-1}\,\mathrm{K}^{-1}\); \(\gamma=1.4\); \(R_g=287\, \mathrm{J}\,\mathrm{kg}^{-1}\,\mathrm{K}^{-1}\).

\noindent\textbf{2 (Nivel de examen, Desarrollo)}
\(C19-02-P \cdot A2 \cdot\) Turbina con descarga axial. Diseñe un escalón repetitivo de turbina que descarga axialmente. Calcule carga, trabajo, temperatura total de salida y presión total de salida. Datos: \(T_{0e}=1490\, \mathrm{K}\); \(p_{0e}=940000\, \mathrm{Pa}\); \(\mathcal R=0.48\); \(U=314\, \mathrm{m}\,\mathrm{s}^{-1}\); \(\eta=0.895\); \(c_p=1156.697\, \mathrm{J}\,\mathrm{kg}^{-1}\,\mathrm{K}^{-1}\); \(\gamma=1.33\).

\noindent\textbf{3 (Nivel de examen, Desarrollo)}
\(C19-03-P \cdot A3 \cdot\) Distribución radial y tracción. Un escalón compresor tiene leyes libres en ambas estaciones. Calcule trabajo radialmente uniforme, reacción en raíz y tensión centrífuga nominal de pala uniforme. Resultado final: tensión. Datos: \(\Omega=1090\, \mathrm{s}^{-1}\); \(K_e=14\, \mathrm{m}^{2}\,\mathrm{s}^{-1}\); \(K_s=44\, \mathrm{m}^{2}\,\mathrm{s}^{-1}\); \(r_h=0.218\, \mathrm{m}\); \(r_t=0.302\, \mathrm{m}\); \(\rho=8050\, \mathrm{kg}\,\mathrm{m}^{-3}\).

\noindent\textbf{4 (Nivel de examen, Desarrollo)}
\(C19-04-P \cdot B1 \cdot\) Rodete radial y prestación. Un rodete de salida radial sin prerrotación tiene radio y giro dados. Con slip y eficiencia prescritos, halle trabajo, temperatura total y relación de presión. Datos: \(\Omega=2070\, \mathrm{s}^{-1}\); \(r=0.206\, \mathrm{m}\); \(\sigma_s=0.875\); \(\eta=0.825\); \(T_{0e}=296\, \mathrm{K}\); \(c_p=1004.5\, \mathrm{J}\,\mathrm{kg}^{-1}\,\mathrm{K}^{-1}\); \(\gamma=1.4\).

\noindent\textbf{5 (Nivel de examen, Desarrollo)}
\(C19-05-P \cdot B2 \cdot\) Cascada y elección entera. Una fila compresora tiene velocidades relativas y cambio tangencial dados. Se ha elegido un número entero de álabes; calcule paso, solidez real y difusión. Datos: \(r=0.279\, \mathrm{m}\); \(c=0.054\, \mathrm{m}\); \(Z=43\); \(W_e=296\, \mathrm{m}\,\mathrm{s}^{-1}\); \(W_s=224\, \mathrm{m}\,\mathrm{s}^{-1}\); \(|\Delta W_\theta|=103\, \mathrm{m}\,\mathrm{s}^{-1}\).

\noindent\textbf{6 (Nivel de examen, Desarrollo)}
\(C19-06-P \cdot B3 \cdot\) Balance térmico integrado. Una pared se refrigera por convección interior. Calcule flujo por área, temperatura caliente y caudal mínimo energético para el incremento permitido. Resultado final: caudal. Datos: \(T_{aw}=1580\, \mathrm{K}\); \(T_c=720\, \mathrm{K}\); \(h_g=1540\, \mathrm{W}\,\mathrm{m}^{-2}\,\mathrm{K}^{-1}\); \(h_c=2320\, \mathrm{W}\,\mathrm{m}^{-2}\,\mathrm{K}^{-1}\); \(t=0.0012\, \mathrm{m}\); \(k_m=23\, \mathrm{W}\,\mathrm{m}^{-1}\,\mathrm{K}^{-1}\); \(A=0.028\, \mathrm{m}^{2}\); \(c_p=1080\, \mathrm{J}\,\mathrm{kg}^{-1}\,\mathrm{K}^{-1}\); \(\Delta T_c=165\, \mathrm{K}\).
\end{tcolorbox}


\studiumchaptercheck

\appendix
\chapter{Notación y unidades}

\begingroup
\setlength{\LTpre}{0pt}
\setlength{\LTpost}{0pt}
\begin{longtable}{@{}>{\centering\arraybackslash}p{0.17\linewidth}>{\raggedright\arraybackslash}p{0.55\linewidth}>{\raggedright\arraybackslash}p{0.2\linewidth}@{}}
\textbf{Símbolo} & \textbf{Significado} & \textbf{Unidades} \\
\hline
\endfirsthead
\textbf{Símbolo} & \textbf{Significado} & \textbf{Unidades} \\
\hline
\endhead
$B$ & Fracción de área disponible; el bloqueo es 1-B & adimensional \\
$C$ & Velocidad absoluta & $\mathrm{m\,s^{-1}}$ \\
$C_{\theta}$ & Componente tangencial absoluta algebraica & $\mathrm{m\,s^{-1}}$ \\
$C_{x}$ & Componente axial absoluta; coincide con la relativa local & $\mathrm{m\,s^{-1}}$ \\
$K_v$ & Constante de vórtice libre $rC_\theta$ (vórtice libre) & $\mathrm{m^2\,s^{-1}}$ \\
$k_f$ & Coeficiente de vórtice forzado $C_\theta=k_f r$ & $\mathrm{s^{-1}}$ \\
$\Omega$ & Velocidad angular en radianes por segundo & $\mathrm{s^{-1}}$ \\
$R_{g}$ & Constante del gas; distinta de la reacción & $\mathrm{J\,kg^{-1}\,K^{-1}}$ \\
$T_{0}$ & Temperatura total absoluta & $\mathrm{K}$ \\
$U$ & Velocidad periférica, positiva en el giro declarado & $\mathrm{m\,s^{-1}}$ \\
$W$ & Velocidad relativa al rotor & $\mathrm{m\,s^{-1}}$ \\
$\alpha$ & Ángulo absoluto desde el eje hacia el giro positivo & $\mathrm{rad}$ \\
$\beta$ & Ángulo relativo algebraico desde el eje hacia el giro positivo & $\mathrm{rad}$ \\
$\delta$ & Desviación: ángulo de gas menos ángulo metálico de salida & $\mathrm{rad}$ \\
$h$ & Entalpía estática & $\mathrm{J\,kg^{-1}}$ \\
$h_{0}$ & Entalpía total absoluta & $\mathrm{J\,kg^{-1}}$ \\
$i$ & Incidencia: ángulo de gas menos ángulo metálico de entrada & $\mathrm{rad}$ \\
$\mathcal{R}$ & Reacción definida como cociente de cambios estáticos & adimensional \\
$p_{0}$ & Presión total absoluta & $\mathrm{Pa}$ \\
$\phi$ & Coeficiente de flujo $C_x/U$ & adimensional \\
$\psi_c$ & Carga de compresor $w_c/U^2$ & adimensional \\
$\psi_t$ & Carga de turbina $w_t/U^2$ & adimensional \\
$r$ & Radio desde el eje & $\mathrm{m}$ \\
$\sigma$ & Solidez cuerda sobre paso; no tensión en este contexto & adimensional \\
$\sigma_{c}$ & Tensión centrífuga nominal & $\mathrm{Pa}$ \\
$\sigma_{s}$ & Deslizamiento definido como cociente de remolinos real e ideal & adimensional \\
$w_{f}$ & Trabajo específico comunicado al fluido; positivo en compresor & $\mathrm{J\,kg^{-1}}$ \\
$w_{t}$ & Trabajo específico extraído positivo en turbina & $\mathrm{J\,kg^{-1}}$ \\
\end{longtable}
\endgroup
\par\medskip\noindent\textbf{Advertencia de notación.} El identificador \texttt{kv} designa dos magnitudes distintas en este manual: $K_v=rC_\theta$ para el vórtice libre, en $\mathrm{m^2\,s^{-1}}$, y $k_f$ para el vórtice forzado, en $\mathrm{s^{-1}}$. No se intercambian. La eficiencia politrópica $\eta_p$ de un compresor no es el rendimiento propulsivo de una hélice; las fórmulas equivalentes de otras asignaturas usan definiciones distintas. Compruebe siempre la definición escrita antes de copiar un valor.\par\medskip
\chapter{Formulario y condiciones de uso}

\textbf{Convención.} Las fórmulas siguientes son las mismas relaciones ya deducidas en los capítulos. Se han reescrito con la notación del apéndice para evitar identificadores internos. Antes de usar una de ellas, compruebe sus hipótesis, signos y unidades; una sustitución numérica correcta no demuestra que el modelo sea válido.

\noindent\textbf{Modelo C01 · Estado estático desde remanso}

\[
\fitmath{\rho=\frac{p_0}{R_g T_0}\left(1+\frac{\gamma-1}{2}\mathrm{Ma}^2\right)^{-\frac{1}{\gamma-1}}}
\]

\noindent\textbf{Modelo C02 · Par de un compresor}

\[
\fitmath{\dot W=\Omega\,\dot m\,r\,(C_{\theta s}-C_{\theta e})}
\]

\noindent\textbf{Modelo C03 · Compresión total-total}

\[
\fitmath{w_c=\frac{T_{0e}\,c_p}{\eta_{tt}}\left(\pi_c^{\frac{\gamma-1}{\gamma}}-1\right)}
\]

\noindent\textbf{Modelo C04 · Acción ideal}

\[
\fitmath{\Delta T_0=\frac{2\,U\,(C_\theta-U)}{c_p}}
\]

\noindent\textbf{Modelo C05 · Triángulos desde carga y reacción}

\[
\fitmath{W_{\theta 3}=U\left(-\frac{\psi}{2}-\mathcal R\right)}
\]

\noindent\textbf{Modelo C06 · Área crítica}

\[
\fitmath{A^*=\frac{\dot m\,\sqrt{T_{0e}}}{p_{0e}}\left(\frac{2}{\gamma+1}\right)^{-\frac{\gamma+1}{2(\gamma-1)}}\sqrt{\frac{R_g}{\gamma}}}
\]

\noindent\textbf{Modelo C07 · Trabajo y presión del escalón}

\[
\fitmath{\pi_c=\left(1+\frac{\eta_{tt}\,\psi\,U^2}{T_{0e}\,c_p}\right)^{\frac{\gamma}{\gamma-1}}}
\]

\noindent\textbf{Modelo C08 · Indicador de difusión}

\[
\fitmath{D=\frac{\Delta V_\theta}{2\,\sigma\,V_e}+1-\frac{V_s}{V_e}}
\]

\noindent\textbf{Modelo C09 · Compresión de varios escalones}

\[
\fitmath{\dot W=\frac{T_{0e}\,c_p\,\dot m}{\eta_m}\left(\pi_c^{\frac{\gamma-1}{\eta_p\,\gamma}}-1\right)}
\]

\noindent\textbf{Modelo C10 · Incidencia durante variación de caudal}

\[
\fitmath{i=-\kappa-\arctan\frac{U}{C_x}}
\]

\noindent\textbf{Modelo C11 · Presión de vórtice libre}

\[
\fitmath{\Delta p=\frac{K_v^{\,2}\,\rho}{2}\left(\frac{1}{r_m^2}-\frac{1}{r_t^2}\right)}
\]

\noindent\textbf{Modelo C12 · Torsión de entrada de rotor}

\[
\fitmath{\Delta\beta=\arctan\frac{\Omega\,r_h}{C_x}-\arctan\frac{\Omega\,r_t}{C_x}}
\]

\noindent\textbf{Modelo C13 · Tracción centrífuga uniforme}

\[
\fitmath{\sigma_c=\frac{\rho\,\Omega^2}{2}\left(r_t^2-r_h^2\right)}
\]

\noindent\textbf{Modelo C14 · Red térmica de pared}

\[
\fitmath{T_w=T_{aw}-\frac{T_{aw}-T_c}{h_g\left(\frac1{h_g}+\frac1{h_c}+\frac{t}{k_m}\right)}}
\]

\noindent\textbf{Modelo C15 · Rodete radial con deslizamiento}

\[
\fitmath{\Delta T_0=\frac{\sigma_s\,U^2}{c_p}}
\]

\noindent\textbf{Modelo C16 · Difusor sin álabes ideal}

\[
\fitmath{\Delta p=\frac{\rho\,(C_r^2+C_\theta^2)}{2}\left(1-\frac{r_e^2}{r_s^2}\right)}
\]

\noindent\textbf{Modelo C17 · Velocidad exigida por trabajo}

\[
\fitmath{r=\frac{1}{\Omega}\sqrt{\frac{w_c}{\sigma_s}}}
\]

\noindent\textbf{Modelo C18 · Rotor compresor a radio elegido}

\[
\fitmath{\kappa=-i+\arctan\frac{K_v/r-\Omega\,r}{C_x}}
\]

\chapter{Soluciones de los ejercicios}
\section{Capítulo 1: Fundamentos, clasificación y estados del gas}
\subsection*{C01-01-P}

\noindent\textbf{Enunciado}
\(C01-01-P \cdot\) Estado estático desde remanso. Ejercicio de transferencia. Una estación de ensayo tiene los datos totales y Mach indicados. Obtenga temperatura, presión y densidad estáticas; la incógnita final es la densidad. Antes de consultar la solución, prediga el efecto de cambiar el primer dato respecto del caso resuelto de esta familia. Justifique la tendencia y explique qué restricción podría convertirse en limitante. Datos: \(T_0=338.67\, \mathrm{K}\); \(p_0=154000\, \mathrm{Pa}\); \(\mathrm{Ma}=0.62\); \(\gamma=1.4\); \(R_g=287\, \mathrm{J}\,\mathrm{kg}^{-1}\,\mathrm{K}^{-1}\).

\noindent\textbf{Datos}
\begin{itemize}
\item $T_0 = 338.67\,\mathrm{K}$ — temperatura total
\item $p_0 = 154000\,\mathrm{Pa}$ — presión total
\item $\mathrm{Ma} = 0.62$ — Mach absoluto
\item $\gamma = 1.4$ — razón de calores
\item $R_g = 287\,\mathrm{J}\,\mathrm{kg}^{-1}\,\mathrm{K}^{-1}$ — constante del gas
\end{itemize}

\noindent\textbf{Incógnita}
$\rho$, en $\mathrm{kg\,m^{-3}}$; también las magnitudes intermedias desarrolladas.

\noindent\textbf{Modelo físico}

\noindent\textit{Estado estático desde remanso}

\[
\fitmath{\rho=\frac{p_0 \left(\frac{\mathrm{Ma}^{2} \left(\gamma - 1\right)}{2} + 1\right)^{- \frac{1}{\gamma - 1}}}{R_g T_0}}
\]

\noindent\textbf{Hipótesis}
\begin{itemize}
\item Gas perfecto y propiedades constantes.
\item Remanso definido localmente; no se presume isentropía entre estaciones.
\end{itemize}

\noindent\textbf{Desarrollo paso a paso}

\begin{enumerate}

\item El remanso local fija la temperatura estática.

\[\fitmath{T=\frac{T_0}{1+\frac{\left(\gamma-1\right)\cdot {\mathrm{Ma}}^{2}}{2}}}\]

\[\fitmath{T=\frac{(338.67)}{1+\frac{\left((1.4)-1\right)\cdot {(0.62)}^{2}}{2}}\approx 314.4919\,\mathrm{K}}\]

\item Se aplica el frenado isentrópico local para recuperar la presión.

\[\fitmath{p=\frac{p_0}{{\left(1+\frac{\left(\gamma-1\right)\cdot {\mathrm{Ma}}^{2}}{2}\right)}^{\frac{\gamma}{\gamma-1}}}}\]

\[\fitmath{p=\frac{(154000)}{{\left(1+\frac{\left((1.4)-1\right)\cdot {(0.62)}^{2}}{2}\right)}^{\frac{(1.4)}{(1.4)-1}}}\approx 118832.5\,\mathrm{Pa}}\]

\item La ecuación de estado usa las dos magnitudes estáticas.

\[\fitmath{\rho=\frac{p_0}{R_g\cdot T_0\cdot {\left(1+\frac{\left(\gamma-1\right)\cdot {\mathrm{Ma}}^{2}}{2}\right)}^{\frac{1}{\gamma-1}}}}\]

\[\fitmath{\rho=\frac{(154000)}{(287)\cdot (338.67)\cdot {\left(1+\frac{\left((1.4)-1\right)\cdot {(0.62)}^{2}}{2}\right)}^{\frac{1}{(1.4)-1}}}\approx 1.31657\,\mathrm{kg/m^{3}}}\]

\end{enumerate}


\noindent\textbf{Sustitución numérica}

\[
\fitmath{\frac{\left(154000\,\mathrm{Pa}\right) \left(\frac{\left(0.62\right)^{2} \left(\left(1.4\right) - 1\right)}{2} + 1\right)^{- \frac{1}{\left(1.4\right) - 1}}}{\left(287\,\mathrm{J}\,\mathrm{kg}^{-1}\,\mathrm{K}^{-1}\right) \left(338.67\,\mathrm{K}\right)}}
\]

\noindent\textbf{Resultado}

\[
\fitmath{\rho=1.316569586\times10^{0}\,\mathrm{kg}\,\mathrm{m}^{-3}}
\]

\noindent\textbf{Interpretación}
La densidad total no sustituye a la estática en el área de paso. Verifique que T y p son menores que sus valores totales.

\noindent\textbf{Comprobación de los pasos}
\begin{itemize}
\item Paso 2: $314.492$ — reproducido
\item Paso 4: $118832$ — reproducido
\item Paso 6: $1.31657$ — reproducido
\end{itemize}

\noindent\textbf{Comprobación de la respuesta} la aritmética del último paso reproduce la respuesta enunciada a la precisión con que está escrita. Es una reproducción aritmética; no valida el modelo físico ni sus hipótesis.

\noindent\textbf{Limitaciones}
\begin{itemize}
\item Modelo educativo; la aplicabilidad a una máquina real requiere revisión experta.
\end{itemize}

\noindent\textbf{Errores frecuentes}
\begin{itemize}
\item Introducir Celsius en un cociente.
\item Usar presión total y temperatura estática en la misma ecuación de estado.
\end{itemize}
% studium:end-record eyJpZCI6ICJQUkItMDAwMiIsICJzZWN0aW9uIjogImMwMSJ9
\subsection*{C01-02-P}

\noindent\textbf{Enunciado}
\(C01-02-P \cdot\) Difusor adiabático. Ejercicio de transferencia. Un difusor reduce el módulo de velocidad. Halle el aumento de temperatura estática y la temperatura de salida. Antes de consultar la solución, prediga el efecto de cambiar el primer dato respecto del caso resuelto de esta familia. Justifique la tendencia y explique qué restricción podría convertirse en limitante. Datos: \(T_e=309.915\, \mathrm{K}\); \(V_e=218\, \mathrm{m}\,\mathrm{s}^{-1}\); \(V_s=94\, \mathrm{m}\,\mathrm{s}^{-1}\); \(c_p=1004.5\, \mathrm{J}\,\mathrm{kg}^{-1}\,\mathrm{K}^{-1}\).

\noindent\textbf{Datos}
\begin{itemize}
\item $T_e = 309.915\,\mathrm{K}$ — temperatura estática de entrada
\item $V_e = 218\,\mathrm{m}\,\mathrm{s}^{-1}$ — velocidad de entrada
\item $V_s = 94\,\mathrm{m}\,\mathrm{s}^{-1}$ — velocidad de salida
\item $c_p = 1004.5\,\mathrm{J}\,\mathrm{kg}^{-1}\,\mathrm{K}^{-1}$ — calor específico
\end{itemize}

\noindent\textbf{Incógnita}
$T_s$, en $\mathrm{K}$; también las magnitudes intermedias desarrolladas.

\noindent\textbf{Modelo físico}

\noindent\textit{Difusor adiabático}

\[
\fitmath{T_s=T_e + \frac{V_e^{2} - V_s^{2}}{2 c_p}}
\]

\noindent\textbf{Hipótesis}
\begin{itemize}
\item Estacionario, adiabático, sin trabajo ni cambio de cota.
\end{itemize}

\noindent\textbf{Desarrollo paso a paso}

\begin{enumerate}

\item Sin trabajo ni calor, la entalpía total permanece constante.

\[\fitmath{\mathrm{dT}=\frac{{\mathrm{Ve}}^{2}-{\mathrm{Vs}}^{2}}{2\cdot \mathrm{cp}}}\]

\[\fitmath{\mathrm{dT}=\frac{{(218)}^{2}-{(94)}^{2}}{2\cdot (1004.5)}\approx 19.25734\,\mathrm{K}}\]

\item Se suma el aumento estático al estado de entrada.

\[\fitmath{\mathrm{Ts}=\mathrm{Te}+\frac{{\mathrm{Ve}}^{2}-{\mathrm{Vs}}^{2}}{2\cdot \mathrm{cp}}}\]

\[\fitmath{\mathrm{Ts}=(309.915)+\frac{{(218)}^{2}-{(94)}^{2}}{2\cdot (1004.5)}\approx 329.1723\,\mathrm{K}}\]

\end{enumerate}


\noindent\textbf{Sustitución numérica}

\[
\fitmath{\left(309.915\,\mathrm{K}\right) + \frac{\left(\left(218\,\mathrm{m}\,\mathrm{s}^{-1}\right)\right)^{2} - \left(\left(94\,\mathrm{m}\,\mathrm{s}^{-1}\right)\right)^{2}}{2 \left(1004.5\,\mathrm{J}\,\mathrm{kg}^{-1}\,\mathrm{K}^{-1}\right)}}
\]

\noindent\textbf{Resultado}

\[
\fitmath{T_s=3.291723420\times10^{2}\,\mathrm{K}}
\]

\noindent\textbf{Interpretación}
La temperatura estática aumenta aunque no se suministra trabajo al gas. La presión final no puede calcularse sin una condición de pérdida.

\noindent\textbf{Comprobación de los pasos}
\begin{itemize}
\item Paso 2: $19.2573$ — reproducido
\item Paso 4: $329.172$ — reproducido
\end{itemize}

\noindent\textbf{Comprobación de la respuesta} la aritmética del último paso reproduce la respuesta enunciada a la precisión con que está escrita. Es una reproducción aritmética; no valida el modelo físico ni sus hipótesis.

\noindent\textbf{Limitaciones}
\begin{itemize}
\item Modelo educativo; la aplicabilidad a una máquina real requiere revisión experta.
\end{itemize}

\noindent\textbf{Errores frecuentes}
\begin{itemize}
\item Interpretar el aumento estático como trabajo de eje.
\end{itemize}
% studium:end-record eyJpZCI6ICJQUkItMDAwNCIsICJzZWN0aW9uIjogImMwMSJ9
\subsection*{C01-03-P}

\noindent\textbf{Enunciado}
\(C01-03-P \cdot\) Entropía en fila fija. Ejercicio de transferencia. Una fila fija conserva temperatura total y reduce la presión total con la razón indicada. Halle la generación específica de entropía. Antes de consultar la solución, prediga el efecto de cambiar el primer dato respecto del caso resuelto de esta familia. Justifique la tendencia y explique qué restricción podría convertirse en limitante. Datos: \(\frac{p_{0s}}{p_{0e}}=0.947\); \(R_g=287\, \mathrm{J}\,\mathrm{kg}^{-1}\,\mathrm{K}^{-1}\).

\noindent\textbf{Datos}
\begin{itemize}
\item $\frac{p_{0s}}{p_{0e}} = 0.947$ — presión total salida/entrada
\item $R_g = 287\,\mathrm{J}\,\mathrm{kg}^{-1}\,\mathrm{K}^{-1}$ — constante del gas
\end{itemize}

\noindent\textbf{Incógnita}
$\Delta s$, en $\mathrm{J\,kg^{-1}\,K^{-1}}$; también las magnitudes intermedias desarrolladas.

\noindent\textbf{Modelo físico}

\noindent\textit{Entropía en fila fija}

\[
\fitmath{\Delta s=- R_g \log{\left(\frac{p_{0s}}{p_{0e}} \right)}}
\]

\noindent\textbf{Hipótesis}
\begin{itemize}
\item Gas perfecto, adiabático, sin trabajo ni mezcla de especies.
\end{itemize}

\noindent\textbf{Desarrollo paso a paso}

\begin{enumerate}

\item La igualdad de temperaturas totales elimina el término térmico de la entropía.

\[\fitmath{L=-\ln\left(\mathrm{ratio}\right)}\]

\[\fitmath{L=-\ln\left((0.947)\right)\approx 0.05445619\,}\]

\item El descenso de presión produce generación positiva.

\[\fitmath{\mathrm{ds}=-R_g\cdot \ln\left(\mathrm{ratio}\right)}\]

\[\fitmath{\mathrm{ds}=-(287)\cdot \ln\left((0.947)\right)\approx 15.62893\,\mathrm{J/(kg\,K)}}\]

\end{enumerate}


\noindent\textbf{Sustitución numérica}

\[
\fitmath{- \left(287\,\mathrm{J}\,\mathrm{kg}^{-1}\,\mathrm{K}^{-1}\right) \log{\left(\left(0.947\right) \right)}}
\]

\noindent\textbf{Resultado}

\[
\fitmath{\Delta s=1.562892532\times10^{1}\,\mathrm{J}\,\mathrm{kg}^{-1}\,\mathrm{K}^{-1}}
\]

\noindent\textbf{Interpretación}
El signo debe ser positivo para la fila adiabática pasiva. No se ha perdido energía total hacia el ambiente.

\noindent\textbf{Comprobación de los pasos}
\begin{itemize}
\item Paso 2: $0.0544562$ — reproducido
\item Paso 4: $15.6289$ — reproducido
\end{itemize}

\noindent\textbf{Comprobación de la respuesta} la aritmética del último paso reproduce la respuesta enunciada a la precisión con que está escrita. Es una reproducción aritmética; no valida el modelo físico ni sus hipótesis.

\noindent\textbf{Limitaciones}
\begin{itemize}
\item Modelo educativo; la aplicabilidad a una máquina real requiere revisión experta.
\end{itemize}

\noindent\textbf{Errores frecuentes}
\begin{itemize}
\item Invertir la razón y obtener entropía negativa.
\end{itemize}
% studium:end-record eyJpZCI6ICJQUkItMDAwNiIsICJzZWN0aW9uIjogImMwMSJ9
\subsection*{C01-04-P}

\noindent\textbf{Enunciado}
\(C01-04-P \cdot\) Anillo con remolino. Ejercicio de transferencia. Dimensione el área frontal anular a partir de un módulo absoluto y un ángulo con el eje. Halle también la altura a radio medio fijado; esa altura es la incógnita final. Antes de consultar la solución, prediga el efecto de cambiar el primer dato respecto del caso resuelto de esta familia. Justifique la tendencia y explique qué restricción podría convertirse en limitante. Datos: \(\dot m=20.874\, \mathrm{kg}\,\mathrm{s}^{-1}\); \(\rho=1.38\, \mathrm{kg}\,\mathrm{m}^{-3}\); \(C_e=186\, \mathrm{m}\,\mathrm{s}^{-1}\); \(\alpha_e=0.5585054\, \mathrm{rad}\); \(r_m=0.29\, \mathrm{m}\).

\noindent\textbf{Datos}
\begin{itemize}
\item $\dot m = 20.874\,\mathrm{kg}\,\mathrm{s}^{-1}$ — caudal
\item $\rho = 1.38\,\mathrm{kg}\,\mathrm{m}^{-3}$ — densidad estática
\item $C_e = 186\,\mathrm{m}\,\mathrm{s}^{-1}$ — módulo absoluto
\item $\alpha_e = 0.55850536063819\,\mathrm{rad}$ — ángulo absoluto con el eje
\item $r_m = 0.29\,\mathrm{m}$ — radio medio aritmético
\end{itemize}

\noindent\textbf{Incógnita}
$b$, en $\mathrm{m}$; también las magnitudes intermedias desarrolladas.

\noindent\textbf{Modelo físico}

\noindent\textit{Anillo con remolino}

\[
\fitmath{b=\frac{\dot m}{2 \pi C_e \rho r_m \cos{\left(\alpha_e \right)}}}
\]

\noindent\textbf{Hipótesis}
\begin{itemize}
\item Flujo axial medio uniforme, bloqueo despreciable.
\end{itemize}

\noindent\textbf{Desarrollo paso a paso}

\begin{enumerate}

\item La proyección axial atraviesa el área frontal.

\[\fitmath{\mathrm{Cx}=\mathrm{Vel}\cdot \cos\left(\mathrm{ang}\right)}\]

\[\fitmath{\mathrm{Cx}=(186)\cdot \cos\left((0.5585053606)\right)\approx 157.7369\,\mathrm{m/s}}\]

\item La continuidad fija el área anular.

\[\fitmath{\mathrm{Aa}=\frac{\mathrm{md}}{\rho\cdot \mathrm{Vel}\cdot \cos\left(\mathrm{ang}\right)}}\]

\[\fitmath{\mathrm{Aa}=\frac{(20.874)}{(1.38)\cdot (186)\cdot \cos\left((0.5585053606)\right)}\approx 0.09589438\,\mathrm{m^{2}}}\]

\item La identidad geométrica A=2 pi rm b fija la altura.

\[\fitmath{b=\frac{\mathrm{md}}{\rho\cdot \mathrm{Vel}\cdot \cos\left(\mathrm{ang}\right)\cdot 2\cdot \pi\cdot \mathrm{rm}}}\]

\[\fitmath{b=\frac{(20.874)}{(1.38)\cdot (186)\cdot \cos\left((0.5585053606)\right)\cdot 2\cdot \pi\cdot (0.29)}\approx 0.05262781\,\mathrm{m}}\]

\end{enumerate}


\noindent\textbf{Sustitución numérica}

\[
\fitmath{\frac{\left(20.874\,\mathrm{kg}\,\mathrm{s}^{-1}\right)}{2 \pi \left(186\,\mathrm{m}\,\mathrm{s}^{-1}\right) \left(1.38\,\mathrm{kg}\,\mathrm{m}^{-3}\right) \left(0.29\,\mathrm{m}\right) \cos{\left(\left(0.55850536063819\,\mathrm{rad}\right) \right)}}}
\]

\noindent\textbf{Resultado}

\[
\fitmath{b=5.262781008\times10^{-2}\,\mathrm{m}}
\]

\noindent\textbf{Interpretación}
Usar el módulo de velocidad reduciría indebidamente el área. Reconstruya el caudal con el área obtenida.

\noindent\textbf{Comprobación de los pasos}
\begin{itemize}
\item Paso 2: $157.737$ — reproducido
\item Paso 4: $0.0958944$ — reproducido
\item Paso 6: $0.0526278$ — reproducido
\end{itemize}

\noindent\textbf{Comprobación de la respuesta} la aritmética del último paso reproduce la respuesta enunciada a la precisión con que está escrita. Es una reproducción aritmética; no valida el modelo físico ni sus hipótesis.

\noindent\textbf{Limitaciones}
\begin{itemize}
\item Modelo educativo; la aplicabilidad a una máquina real requiere revisión experta.
\end{itemize}

\noindent\textbf{Errores frecuentes}
\begin{itemize}
\item Usar el seno del ángulo axial.
\item Confundir altura y diámetro.
\end{itemize}
% studium:end-record eyJpZCI6ICJQUkItMDAwOCIsICJzZWN0aW9uIjogImMwMSJ9
\section{Capítulo 2: Momento angular, Euler y triángulos de velocidades}
\subsection*{C02-01-P}

\noindent\textbf{Enunciado}
\(C02-01-P \cdot\) Par de un compresor. Ejercicio de transferencia. Calcule el par comunicado al fluido y la potencia ideal de un rotor a radio constante. La incógnita final es la potencia. Antes de consultar la solución, prediga el efecto de cambiar el primer dato respecto del caso resuelto de esta familia. Justifique la tendencia y explique qué restricción podría convertirse en limitante. Datos: \(\dot m=23.43\, \mathrm{kg}\,\mathrm{s}^{-1}\); \(r=0.31\, \mathrm{m}\); \(\Omega=970\, \mathrm{s}^{-1}\); \(C_{\theta e}=24\, \mathrm{m}\,\mathrm{s}^{-1}\); \(C_{\theta s}=109\, \mathrm{m}\,\mathrm{s}^{-1}\).

\noindent\textbf{Datos}
\begin{itemize}
\item $\dot m = 23.43\,\mathrm{kg}\,\mathrm{s}^{-1}$ — caudal
\item $r = 0.31\,\mathrm{m}$ — radio
\item $\Omega = 970\,\mathrm{s}^{-1}$ — velocidad angular en rad/s
\item $C_{\theta e} = 24\,\mathrm{m}\,\mathrm{s}^{-1}$ — remolino de entrada
\item $C_{\theta s} = 109\,\mathrm{m}\,\mathrm{s}^{-1}$ — remolino de salida
\end{itemize}

\noindent\textbf{Incógnita}
$\dot{W}$, en $\mathrm{W}$; también las magnitudes intermedias desarrolladas.

\noindent\textbf{Modelo físico}

\noindent\textit{Par de un compresor}

\[
\fitmath{\dot{W}=\Omega \dot m r \left(- C_{\theta e} + C_{\theta s}\right)}
\]

\noindent\textbf{Hipótesis}
\begin{itemize}
\item Estacionario y sin pares adicionales.
\item Radio constante y flujo uniforme en los cortes.
\end{itemize}

\noindent\textbf{Desarrollo paso a paso}

\begin{enumerate}

\item Se calcula el flujo neto de momento angular.

\[\fitmath{\tau=\mathrm{md}\cdot \mathrm{rad}\cdot \left(\mathrm{cs}-\mathrm{ce}\right)}\]

\[\fitmath{\tau=(23.43)\cdot (0.31)\cdot \left((109)-(24)\right)\approx 617.3805\,\mathrm{N\,m}}\]

\item Se multiplica el par por la velocidad angular.

\[\fitmath{\mathrm{Pow}=\mathrm{md}\cdot \mathrm{rad}\cdot \left(\mathrm{cs}-\mathrm{ce}\right)\cdot \Omega}\]

\[\fitmath{\mathrm{Pow}=(23.43)\cdot (0.31)\cdot \left((109)-(24)\right)\cdot (970)\approx 598859.1\,\mathrm{W}}\]

\end{enumerate}


\noindent\textbf{Sustitución numérica}

\[
\fitmath{\left(970\,\mathrm{s}^{-1}\right) \left(23.43\,\mathrm{kg}\,\mathrm{s}^{-1}\right) \left(0.31\,\mathrm{m}\right) \left(- \left(24\,\mathrm{m}\,\mathrm{s}^{-1}\right) + \left(109\,\mathrm{m}\,\mathrm{s}^{-1}\right)\right)}
\]

\noindent\textbf{Resultado}

\[
\fitmath{\dot{W}=5.988590850\times10^{5}\,\mathrm{W}}
\]

\noindent\textbf{Interpretación}
La potencia sobre el fluido es positiva; la reacción sobre el rotor es negativa.

\noindent\textbf{Comprobación de los pasos}
\begin{itemize}
\item Paso 2: $617.38$ — reproducido
\item Paso 4: $598859$ — reproducido
\end{itemize}

\noindent\textbf{Comprobación de la respuesta} la aritmética del último paso reproduce la respuesta enunciada a la precisión con que está escrita. Es una reproducción aritmética; no valida el modelo físico ni sus hipótesis.

\noindent\textbf{Limitaciones}
\begin{itemize}
\item Modelo educativo; la aplicabilidad a una máquina real requiere revisión experta.
\end{itemize}

\noindent\textbf{Errores frecuentes}
\begin{itemize}
\item Usar velocidades relativas en el volumen fijo.
\end{itemize}
% studium:end-record eyJpZCI6ICJQUkItMDAxMCIsICJzZWN0aW9uIjogImMwMiJ9
\subsection*{C02-02-P}

\noindent\textbf{Enunciado}
\(C02-02-P \cdot\) Triángulo y referencia relativa. Ejercicio de transferencia. Una estación tiene las componentes absolutas y velocidad periférica indicadas. Obtenga la componente relativa, el módulo relativo y la temperatura total relativa. Antes de consultar la solución, prediga el efecto de cambiar el primer dato respecto del caso resuelto de esta familia. Justifique la tendencia y explique qué restricción podría convertirse en limitante. Datos: \(C_x=168.27\, \mathrm{m}\,\mathrm{s}^{-1}\); \(C_\theta=76\, \mathrm{m}\,\mathrm{s}^{-1}\); \(U=284\, \mathrm{m}\,\mathrm{s}^{-1}\); \(T_e=306\, \mathrm{K}\); \(c_p=1004.5\, \mathrm{J}\,\mathrm{kg}^{-1}\,\mathrm{K}^{-1}\).

\noindent\textbf{Datos}
\begin{itemize}
\item $C_x = 168.27\,\mathrm{m}\,\mathrm{s}^{-1}$ — componente axial
\item $C_\theta = 76\,\mathrm{m}\,\mathrm{s}^{-1}$ — componente tangencial absoluta
\item $U = 284\,\mathrm{m}\,\mathrm{s}^{-1}$ — velocidad periférica
\item $T_e = 306\,\mathrm{K}$ — temperatura estática
\item $c_p = 1004.5\,\mathrm{J}\,\mathrm{kg}^{-1}\,\mathrm{K}^{-1}$ — calor específico
\end{itemize}

\noindent\textbf{Incógnita}
$T_{0r}$, en $\mathrm{K}$; también las magnitudes intermedias desarrolladas.

\noindent\textbf{Modelo físico}

\noindent\textit{Triángulo y referencia relativa}

\[
\fitmath{T_{0r}=T_e + \frac{C_x^{2} + \left(C_\theta - U\right)^{2}}{2 c_p}}
\]

\noindent\textbf{Hipótesis}
\begin{itemize}
\item Mismo punto y referencia de rotor con U tangencial.
\end{itemize}

\noindent\textbf{Desarrollo paso a paso}

\begin{enumerate}

\item La transformación vectorial resta U a la componente tangencial.

\[\fitmath{\mathrm{wt}=\mathrm{ct}-\mathrm{uu}}\]

\[\fitmath{\mathrm{wt}=(76)-(284)\approx -208\,\mathrm{m/s}}\]

\item El módulo usa componentes ortogonales.

\[\fitmath{\mathrm{Wr}=\sqrt{{\mathrm{cx}}^{2}+{\left(\mathrm{ct}-\mathrm{uu}\right)}^{2}}}\]

\[\fitmath{\mathrm{Wr}=\sqrt{{(168.27)}^{2}+{\left((76)-(284)\right)}^{2}}\approx 267.5421\,\mathrm{m/s}}\]

\item El remanso relativo se construye con W.

\[\fitmath{\mathrm{Tzr}=\mathrm{Te}+\frac{{\mathrm{cx}}^{2}+{\left(\mathrm{ct}-\mathrm{uu}\right)}^{2}}{2\cdot \mathrm{cp}}}\]

\[\fitmath{\mathrm{Tzr}=(306)+\frac{{(168.27)}^{2}+{\left((76)-(284)\right)}^{2}}{2\cdot (1004.5)}\approx 341.6291\,\mathrm{K}}\]

\end{enumerate}


\noindent\textbf{Sustitución numérica}

\[
\fitmath{\left(306\,\mathrm{K}\right) + \frac{\left(\left(168.27\,\mathrm{m}\,\mathrm{s}^{-1}\right)\right)^{2} + \left(\left(76\,\mathrm{m}\,\mathrm{s}^{-1}\right) - \left(284\,\mathrm{m}\,\mathrm{s}^{-1}\right)\right)^{2}}{2 \left(1004.5\,\mathrm{J}\,\mathrm{kg}^{-1}\,\mathrm{K}^{-1}\right)}}
\]

\noindent\textbf{Resultado}

\[
\fitmath{T_{0r}=3.416290657\times10^{2}\,\mathrm{K}}
\]

\noindent\textbf{Interpretación}
El beta relativo es negativo. La temperatura total relativa no tiene por qué coincidir con la absoluta.

\noindent\textbf{Comprobación de los pasos}
\begin{itemize}
\item Paso 2: $-208$ — reproducido
\item Paso 4: $267.542$ — reproducido
\item Paso 6: $341.629$ — reproducido
\end{itemize}

\noindent\textbf{Comprobación de la respuesta} la aritmética del último paso reproduce la respuesta enunciada a la precisión con que está escrita. Es una reproducción aritmética; no valida el modelo físico ni sus hipótesis.

\noindent\textbf{Limitaciones}
\begin{itemize}
\item Modelo educativo; la aplicabilidad a una máquina real requiere revisión experta.
\end{itemize}

\noindent\textbf{Errores frecuentes}
\begin{itemize}
\item Restar módulos de velocidades.
\end{itemize}
% studium:end-record eyJpZCI6ICJQUkItMDAxMiIsICJzZWN0aW9uIjogImMwMiJ9
\subsection*{C02-03-P}

\noindent\textbf{Enunciado}
\(C02-03-P \cdot\) Euler con radios distintos. Ejercicio de transferencia. Un rotor radial trabaja con prerrotación no nula. Calcule el trabajo comunicado al fluido y el incremento de temperatura total. Antes de consultar la solución, prediga el efecto de cambiar el primer dato respecto del caso resuelto de esta familia. Justifique la tendencia y explique qué restricción podría convertirse en limitante. Datos: \(\Omega=1757.25\, \mathrm{s}^{-1}\); \(r_e=0.075\, \mathrm{m}\); \(r_s=0.19\, \mathrm{m}\); \(C_{\theta e}=31\, \mathrm{m}\,\mathrm{s}^{-1}\); \(C_{\theta s}=247\, \mathrm{m}\,\mathrm{s}^{-1}\); \(c_p=1004.5\, \mathrm{J}\,\mathrm{kg}^{-1}\,\mathrm{K}^{-1}\).

\noindent\textbf{Datos}
\begin{itemize}
\item $\Omega = 1757.25\,\mathrm{s}^{-1}$ — velocidad angular
\item $r_e = 0.075\,\mathrm{m}$ — radio entrada
\item $r_s = 0.19\,\mathrm{m}$ — radio salida
\item $C_{\theta e} = 31\,\mathrm{m}\,\mathrm{s}^{-1}$ — remolino entrada
\item $C_{\theta s} = 247\,\mathrm{m}\,\mathrm{s}^{-1}$ — remolino salida
\item $c_p = 1004.5\,\mathrm{J}\,\mathrm{kg}^{-1}\,\mathrm{K}^{-1}$ — calor específico
\end{itemize}

\noindent\textbf{Incógnita}
$\Delta T_0$, en $\mathrm{K}$; también las magnitudes intermedias desarrolladas.

\noindent\textbf{Modelo físico}

\noindent\textit{Euler con radios distintos}

\[
\fitmath{\Delta T_0=\frac{\Omega \left(- C_{\theta e} r_e + C_{\theta s} r_s\right)}{c_p}}
\]

\noindent\textbf{Hipótesis}
\begin{itemize}
\item Adiabático, propiedades constantes.
\item Valores medios compatibles con el flujo de momento angular.
\end{itemize}

\noindent\textbf{Desarrollo paso a paso}

\begin{enumerate}

\item La velocidad periférica de salida depende de su radio.

\[\fitmath{\mathrm{Us}=\Omega\cdot \mathrm{rs}}\]

\[\fitmath{\mathrm{Us}=(1757.25)\cdot (0.19)\approx 333.8775\,\mathrm{m/s}}\]

\item Se conserva el término de prerrotación en Euler.

\[\fitmath{\mathrm{wc}=\Omega\cdot \left(\mathrm{rs}\cdot \mathrm{cs}-\mathrm{re}\cdot \mathrm{ce}\right)}\]

\[\fitmath{\mathrm{wc}=(1757.25)\cdot \left((0.19)\cdot (247)-(0.075)\cdot (31)\right)\approx 78382.14\,\mathrm{J/kg}}\]

\item Se transforma trabajo en aumento de remanso.

\[\fitmath{\mathrm{dTz}=\frac{\Omega\cdot \left(\mathrm{rs}\cdot \mathrm{cs}-\mathrm{re}\cdot \mathrm{ce}\right)}{\mathrm{cp}}}\]

\[\fitmath{\mathrm{dTz}=\frac{(1757.25)\cdot \left((0.19)\cdot (247)-(0.075)\cdot (31)\right)}{(1004.5)}\approx 78.031\,\mathrm{K}}\]

\end{enumerate}


\noindent\textbf{Sustitución numérica}

\[
\fitmath{\frac{\left(1757.25\,\mathrm{s}^{-1}\right) \left(- \left(31\,\mathrm{m}\,\mathrm{s}^{-1}\right) \left(0.075\,\mathrm{m}\right) + \left(247\,\mathrm{m}\,\mathrm{s}^{-1}\right) \left(0.19\,\mathrm{m}\right)\right)}{\left(1004.5\,\mathrm{J}\,\mathrm{kg}^{-1}\,\mathrm{K}^{-1}\right)}}
\]

\noindent\textbf{Resultado}

\[
\fitmath{\Delta T_0=7.803099676\times10^{1}\,\mathrm{K}}
\]

\noindent\textbf{Interpretación}
Omitir el remolino de entrada sobreestima el trabajo en este caso. Compare ambas evaluaciones.

\noindent\textbf{Comprobación de los pasos}
\begin{itemize}
\item Paso 2: $333.877$ — reproducido
\item Paso 4: $78382.1$ — reproducido
\item Paso 6: $78.031$ — reproducido
\end{itemize}

\noindent\textbf{Comprobación de la respuesta} la aritmética del último paso reproduce la respuesta enunciada a la precisión con que está escrita. Es una reproducción aritmética; no valida el modelo físico ni sus hipótesis.

\noindent\textbf{Limitaciones}
\begin{itemize}
\item Modelo educativo; la aplicabilidad a una máquina real requiere revisión experta.
\end{itemize}

\noindent\textbf{Errores frecuentes}
\begin{itemize}
\item Usar el radio de salida para ambos extremos.
\end{itemize}
% studium:end-record eyJpZCI6ICJQUkItMDAxNCIsICJzZWN0aW9uIjogImMwMiJ9
\subsection*{C02-04-P}

\noindent\textbf{Enunciado}
\(C02-04-P \cdot\) Comprobación de rotalpía. Ejercicio de transferencia. Para un rotor se conocen estados cinemáticos de entrada y salida. Calcule el cambio de entalpía estática exigido por rotalpía. Antes de consultar la solución, prediga el efecto de cambiar el primer dato respecto del caso resuelto de esta familia. Justifique la tendencia y explique qué restricción podría convertirse en limitante. Datos: \(U_e=204.48\, \mathrm{m}\,\mathrm{s}^{-1}\); \(U_s=331\, \mathrm{m}\,\mathrm{s}^{-1}\); \(W_e=246\, \mathrm{m}\,\mathrm{s}^{-1}\); \(W_s=214\, \mathrm{m}\,\mathrm{s}^{-1}\).

\noindent\textbf{Datos}
\begin{itemize}
\item $U_e = 204.48\,\mathrm{m}\,\mathrm{s}^{-1}$ — velocidad periférica entrada
\item $U_s = 331\,\mathrm{m}\,\mathrm{s}^{-1}$ — velocidad periférica salida
\item $W_e = 246\,\mathrm{m}\,\mathrm{s}^{-1}$ — velocidad relativa entrada
\item $W_s = 214\,\mathrm{m}\,\mathrm{s}^{-1}$ — velocidad relativa salida
\end{itemize}

\noindent\textbf{Incógnita}
$\Delta h$, en $\mathrm{J\,kg^{-1}}$; también las magnitudes intermedias desarrolladas.

\noindent\textbf{Modelo físico}

\noindent\textit{Comprobación de rotalpía}

\[
\fitmath{\Delta h=- \frac{U_e^{2}}{2} + \frac{U_s^{2}}{2} + \frac{W_e^{2}}{2} - \frac{W_s^{2}}{2}}
\]

\noindent\textbf{Hipótesis}
\begin{itemize}
\item Rotor adiabático, estacionario en referencia móvil, sin cambio de cota.
\end{itemize}

\noindent\textbf{Desarrollo paso a paso}

\begin{enumerate}

\item Se calcula la contribución de cambio de radio.

\[\fitmath{\mathrm{du}=\frac{{\mathrm{us}}^{2}-{\mathrm{ue}}^{2}}{2}}\]

\[\fitmath{\mathrm{du}=\frac{{(331)}^{2}-{(204.48)}^{2}}{2}\approx 33874.46\,\mathrm{J/kg}}\]

\item Se evalúa por separado la variación de energía cinética relativa.

\[\fitmath{\mathrm{ref}=\frac{{\mathrm{ws}}^{2}-{\mathrm{we}}^{2}}{2}}\]

\[\fitmath{\mathrm{ref}=\frac{{(214)}^{2}-{(246)}^{2}}{2}\approx -7360\,\mathrm{J/kg}}\]

\item De I constante se despeja la variación estática.

\[\fitmath{\mathrm{dh}=\frac{{\mathrm{us}}^{2}-{\mathrm{ue}}^{2}-{\mathrm{ws}}^{2}+{\mathrm{we}}^{2}}{2}}\]

\[\fitmath{\mathrm{dh}=\frac{{(331)}^{2}-{(204.48)}^{2}-{(214)}^{2}+{(246)}^{2}}{2}\approx 41234.46\,\mathrm{J/kg}}\]

\end{enumerate}


\noindent\textbf{Sustitución numérica}

\[
\fitmath{- \frac{\left(\left(204.48\,\mathrm{m}\,\mathrm{s}^{-1}\right)\right)^{2}}{2} + \frac{\left(\left(331\,\mathrm{m}\,\mathrm{s}^{-1}\right)\right)^{2}}{2} + \frac{\left(\left(246\,\mathrm{m}\,\mathrm{s}^{-1}\right)\right)^{2}}{2} - \frac{\left(\left(214\,\mathrm{m}\,\mathrm{s}^{-1}\right)\right)^{2}}{2}}
\]

\noindent\textbf{Resultado}

\[
\fitmath{\Delta h=4.123446480\times10^{4}\,\mathrm{J}\,\mathrm{kg}^{-1}}
\]

\noindent\textbf{Interpretación}
El término centrífugo no puede omitirse cuando U cambia. Reintroduzca dh en la identidad de rotalpía y compruebe residuo cero.

\noindent\textbf{Comprobación de los pasos}
\begin{itemize}
\item Paso 2: $33874.5$ — reproducido
\item Paso 4: $-7360$ — reproducido
\item Paso 6: $41234.5$ — reproducido
\end{itemize}

\noindent\textbf{Comprobación de la respuesta} la aritmética del último paso reproduce la respuesta enunciada a la precisión con que está escrita. Es una reproducción aritmética; no valida el modelo físico ni sus hipótesis.

\noindent\textbf{Limitaciones}
\begin{itemize}
\item Modelo educativo; la aplicabilidad a una máquina real requiere revisión experta.
\end{itemize}

\noindent\textbf{Errores frecuentes}
\begin{itemize}
\item Conservar h+W al cuadrado/2 a radio variable.
\end{itemize}
% studium:end-record eyJpZCI6ICJQUkItMDAxNiIsICJzZWN0aW9uIjogImMwMiJ9
\section{Capítulo 3: Rendimientos, pérdidas y semejanza}
\subsection*{C03-01-P}

\noindent\textbf{Enunciado}
\(C03-01-P \cdot\) Compresión total-total. Ejercicio de transferencia. Calcule temperatura total de salida y trabajo de una compresión con relación de presión y rendimiento dados. Antes de consultar la solución, prediga el efecto de cambiar el primer dato respecto del caso resuelto de esta familia. Justifique la tendencia y explique qué restricción podría convertirse en limitante. Datos: \(T_{0e}=313.11\, \mathrm{K}\); \(\pi_c=4.6\); \(\eta=0.84\); \(\gamma=1.4\); \(c_p=1004.5\, \mathrm{J}\,\mathrm{kg}^{-1}\,\mathrm{K}^{-1}\).

\noindent\textbf{Datos}
\begin{itemize}
\item $T_{0e} = 313.11\,\mathrm{K}$ — temperatura total entrada
\item $\pi_c = 4.6$ — relación presión total
\item $\eta = 0.84$ — rendimiento total-total
\item $\gamma = 1.4$ — razón de calores
\item $c_p = 1004.5\,\mathrm{J}\,\mathrm{kg}^{-1}\,\mathrm{K}^{-1}$ — calor específico
\end{itemize}

\noindent\textbf{Incógnita}
$w_c$, en $\mathrm{J\,kg^{-1}}$; también las magnitudes intermedias desarrolladas.

\noindent\textbf{Modelo físico}

\noindent\textit{Compresión total-total}

\[
\fitmath{w_c=\frac{T_{0e} c_p \left(\pi_c^{\frac{\gamma - 1}{\gamma}} - 1\right)}{\eta}}
\]

\noindent\textbf{Hipótesis}
\begin{itemize}
\item Gas caloríficamente perfecto y proceso adiabático.
\end{itemize}

\noindent\textbf{Desarrollo paso a paso}

\begin{enumerate}

\item Temperatura total de salida

\[\fitmath{\mathrm{Tout}=\mathrm{Tin}\cdot \left(1+\frac{{\mathrm{pr}}^{\frac{\gamma-1}{\gamma}}-1}{\mathrm{eta}}\right)}\]

\[\fitmath{\mathrm{Tout}=(313.11)\cdot \left(1+\frac{{(4.6)}^{\frac{(1.4)-1}{(1.4)}}-1}{(0.84)}\right)\approx 516.8304\,\mathrm{K}}\]

\item La referencia isentrópica utiliza la misma presión de salida.

\[\fitmath{\mathrm{dTs}=\mathrm{Tin}\cdot \left({\mathrm{pr}}^{\frac{\gamma-1}{\gamma}}-1\right)}\]

\[\fitmath{\mathrm{dTs}=(313.11)\cdot \left({(4.6)}^{\frac{(1.4)-1}{(1.4)}}-1\right)\approx 171.1252\,\mathrm{K}}\]

\item El trabajo real es mayor que el ideal.

\[\fitmath{\mathrm{wc}=\frac{\mathrm{cp}\cdot \mathrm{Tin}\cdot \left({\mathrm{pr}}^{\frac{\gamma-1}{\gamma}}-1\right)}{\mathrm{eta}}}\]

\[\fitmath{\mathrm{wc}=\frac{(1004.5)\cdot (313.11)\cdot \left({(4.6)}^{\frac{(1.4)-1}{(1.4)}}-1\right)}{(0.84)}\approx 204637.2\,\mathrm{J/kg}}\]

\end{enumerate}


\noindent\textbf{Sustitución numérica}

\[
\fitmath{\frac{\left(313.11\,\mathrm{K}\right) \left(1004.5\,\mathrm{J}\,\mathrm{kg}^{-1}\,\mathrm{K}^{-1}\right) \left(\left(4.6\right)^{\frac{\left(1.4\right) - 1}{\left(1.4\right)}} - 1\right)}{\left(0.84\right)}}
\]

\noindent\textbf{Resultado}

\[
\fitmath{w_c=2.046371686\times10^{5}\,\mathrm{J}\,\mathrm{kg}^{-1}}
\]

\noindent\textbf{Temperatura total de salida}
\[
\fitmath{T_{0s}=T_{0e}+w_c/c_p=516.8304267\,\mathrm{K}}
\]

\noindent\textbf{Comprobación energética.} $c_p(T_{0s}-T_{0e})=w_c$; no confundir el incremento isentrópico con la temperatura real de salida.

\noindent\textbf{Interpretación}
La eficiencia divide el trabajo ideal en compresión. El resultado debe superar el coste isentrópico.

\noindent\textbf{Comprobación de los pasos}
\begin{itemize}
\item Paso 2: $171.125$ — reproducido
\item Paso 4: $204637$ — reproducido
\end{itemize}

\noindent\textbf{Comprobación de la respuesta} la aritmética del último paso reproduce la respuesta enunciada a la precisión con que está escrita. Es una reproducción aritmética; no valida el modelo físico ni sus hipótesis.

\noindent\textbf{Limitaciones}
\begin{itemize}
\item Modelo educativo; la aplicabilidad a una máquina real requiere revisión experta.
\end{itemize}

\noindent\textbf{Errores frecuentes}
\begin{itemize}
\item Multiplicar por eta para calcular el trabajo real.
\end{itemize}
% studium:end-record eyJpZCI6ICJQUkItMDAxOCIsICJzZWN0aW9uIjogImMwMyJ9
\subsection*{C03-02-P}

\noindent\textbf{Enunciado}
\(C03-02-P \cdot\) Expansión total-total. Ejercicio de transferencia. Determine el trabajo positivo extraído por una turbina con razón de presión salida/entrada dada. Antes de consultar la solución, prediga el efecto de cambiar el primer dato respecto del caso resuelto de esta familia. Justifique la tendencia y explique qué restricción podría convertirse en limitante. Datos: \(T_{0e}=1528.275\, \mathrm{K}\); \(\pi_t=0.43\); \(\eta=0.89\); \(\gamma=1.33\); \(c_p=1156.697\, \mathrm{J}\,\mathrm{kg}^{-1}\,\mathrm{K}^{-1}\).

\noindent\textbf{Datos}
\begin{itemize}
\item $T_{0e} = 1528.275\,\mathrm{K}$ — temperatura total entrada
\item $\pi_t = 0.43$ — presión total salida/entrada
\item $\eta = 0.89$ — rendimiento total-total
\item $\gamma = 1.33$ — razón de calores
\item $c_p = 1156.696969697\,\mathrm{J}\,\mathrm{kg}^{-1}\,\mathrm{K}^{-1}$ — calor específico compatible con R=287
\end{itemize}

\noindent\textbf{Incógnita}
$w_t$, en $\mathrm{J\,kg^{-1}}$; también las magnitudes intermedias desarrolladas.

\noindent\textbf{Modelo físico}

\noindent\textit{Expansión total-total}

\[
\fitmath{w_t=T_{0e} c_p \eta \left(1 - \pi_t^{\frac{\gamma - 1}{\gamma}}\right)}
\]

\noindent\textbf{Hipótesis}
\begin{itemize}
\item Sin sangrados ni refrigeración, gas perfecto constante.
\end{itemize}

\noindent\textbf{Desarrollo paso a paso}

\begin{enumerate}

\item La caída ideal se refiere a igual presión final.

\[\fitmath{\mathrm{dTs}=\mathrm{Tin}\cdot \left(1-{\mathrm{pr}}^{\frac{\gamma-1}{\gamma}}\right)}\]

\[\fitmath{\mathrm{dTs}=(1528.275)\cdot \left(1-{(0.43)}^{\frac{(1.33)-1}{(1.33)}}\right)\approx 288.7434\,\mathrm{K}}\]

\item La caída real es una fracción de la ideal.

\[\fitmath{\mathrm{wt}=\mathrm{eta}\cdot \mathrm{cp}\cdot \mathrm{Tin}\cdot \left(1-{\mathrm{pr}}^{\frac{\gamma-1}{\gamma}}\right)}\]

\[\fitmath{\mathrm{wt}=(0.89)\cdot (1156.69697)\cdot (1528.275)\cdot \left(1-{(0.43)}^{\frac{(1.33)-1}{(1.33)}}\right)\approx 297249.9\,\mathrm{J/kg}}\]

\end{enumerate}


\noindent\textbf{Sustitución numérica}

\[
\fitmath{\left(1528.275\,\mathrm{K}\right) \left(1156.696969697\,\mathrm{J}\,\mathrm{kg}^{-1}\,\mathrm{K}^{-1}\right) \left(0.89\right) \left(1 - \left(0.43\right)^{\frac{\left(1.33\right) - 1}{\left(1.33\right)}}\right)}
\]

\noindent\textbf{Resultado}

\[
\fitmath{w_t=2.972499007\times10^{5}\,\mathrm{J}\,\mathrm{kg}^{-1}}
\]

\noindent\textbf{Interpretación}
El trabajo útil idealizado es positivo y menor que la caída isentrópica disponible.

\noindent\textbf{Comprobación de los pasos}
\begin{itemize}
\item Paso 2: $288.743$ — reproducido
\item Paso 4: $297250$ — reproducido
\end{itemize}

\noindent\textbf{Comprobación de la respuesta} la aritmética del último paso reproduce la respuesta enunciada a la precisión con que está escrita. Es una reproducción aritmética; no valida el modelo físico ni sus hipótesis.

\noindent\textbf{Limitaciones}
\begin{itemize}
\item Modelo educativo; la aplicabilidad a una máquina real requiere revisión experta.
\end{itemize}

\noindent\textbf{Errores frecuentes}
\begin{itemize}
\item Usar presión entrada/salida sin transformar la fórmula.
\end{itemize}
% studium:end-record eyJpZCI6ICJQUkItMDAyMCIsICJzZWN0aW9uIjogImMwMyJ9
\subsection*{C03-03-P}

\noindent\textbf{Enunciado}
\(C03-03-P \cdot\) Conversión politrópica. Ejercicio de transferencia. Un compresor tiene eficiencia politrópica uniforme. Calcule la razón de temperaturas y su eficiencia isentrópica global. Antes de consultar la solución, prediga el efecto de cambiar el primer dato respecto del caso resuelto de esta familia. Justifique la tendencia y explique qué restricción podría convertirse en limitante. Datos: \(\pi_c=13.206\); \(\eta_p=0.905\); \(\gamma=1.4\).

\noindent\textbf{Datos}
\begin{itemize}
\item $\pi_c = 13.206$ — relación presión total
\item $\eta_p = 0.905$ — eficiencia politrópica
\item $\gamma = 1.4$ — razón de calores
\end{itemize}

\noindent\textbf{Incógnita}
$\eta$ (adimensional); también las magnitudes intermedias desarrolladas.

\noindent\textbf{Modelo físico}

\noindent\textit{Conversión politrópica}

\[
\fitmath{\eta=\frac{\pi_c^{\frac{\gamma - 1}{\gamma}} - 1}{\pi_c^{\frac{\gamma - 1}{\eta_p \gamma}} - 1}}
\]

\noindent\textbf{Hipótesis}
\begin{itemize}
\item Eficiencia politrópica y propiedades constantes.
\end{itemize}

\noindent\textbf{Desarrollo paso a paso}

\begin{enumerate}

\item Se integra la relación diferencial politrópica.

\[\fitmath{\mathrm{Tr}={\mathrm{pr}}^{\frac{\gamma-1}{\gamma\cdot \mathrm{ep}}}}\]

\[\fitmath{\mathrm{Tr}={(13.206)}^{\frac{(1.4)-1}{(1.4)\cdot (0.905)}}\approx 2.258576\,}\]

\item Se compara el trabajo isentrópico global con el real.

\[\fitmath{\eta=\frac{{\mathrm{pr}}^{\frac{\gamma-1}{\gamma}}-1}{{\mathrm{pr}}^{\frac{\gamma-1}{\gamma\cdot \mathrm{ep}}}-1}}\]

\[\fitmath{\eta=\frac{{(13.206)}^{\frac{(1.4)-1}{(1.4)}}-1}{{(13.206)}^{\frac{(1.4)-1}{(1.4)\cdot (0.905)}}-1}\approx 0.8663416\,}\]

\end{enumerate}


\noindent\textbf{Sustitución numérica}

\[
\fitmath{\frac{\left(13.206\right)^{\frac{\left(1.4\right) - 1}{\left(1.4\right)}} - 1}{\left(13.206\right)^{\frac{\left(1.4\right) - 1}{\left(0.905\right) \left(1.4\right)}} - 1}}
\]

\noindent\textbf{Resultado}

\[
\fitmath{\eta=8.663416108\times10^{-1}}
\]

\noindent\textbf{Interpretación}
La eficiencia global es menor que la politrópica en este modelo de compresión.

\noindent\textbf{Comprobación de los pasos}
\begin{itemize}
\item Paso 2: $2.25858$ — reproducido
\item Paso 4: $0.866342$ — reproducido
\end{itemize}

\noindent\textbf{Comprobación de la respuesta} la aritmética del último paso reproduce la respuesta enunciada a la precisión con que está escrita. Es una reproducción aritmética; no valida el modelo físico ni sus hipótesis.

\noindent\textbf{Limitaciones}
\begin{itemize}
\item Modelo educativo; la aplicabilidad a una máquina real requiere revisión experta.
\end{itemize}

\noindent\textbf{Errores frecuentes}
\begin{itemize}
\item Elevar la eficiencia de escalón al número de escalones.
\end{itemize}
% studium:end-record eyJpZCI6ICJQUkItMDAyMiIsICJzZWN0aW9uIjogImMwMyJ9
\subsection*{C03-04-P}

\noindent\textbf{Enunciado}
\(C03-04-P \cdot\) Corrección de caudal. Ejercicio de transferencia. Convierta el caudal medido a las referencias de presión y temperatura especificadas. Antes de consultar la solución, prediga el efecto de cambiar el primer dato respecto del caso resuelto de esta familia. Justifique la tendencia y explique qué restricción podría convertirse en limitante. Datos: \(\dot m=33.7605\, \mathrm{kg}\,\mathrm{s}^{-1}\); \(T_{0e}=271\, \mathrm{K}\); \(p_{0e}=82500\, \mathrm{Pa}\); \(T_{\mathrm{ref}}=288.15\, \mathrm{K}\); \(p_{\mathrm{ref}}=101325\, \mathrm{Pa}\).

\noindent\textbf{Datos}
\begin{itemize}
\item $\dot m = 33.7605\,\mathrm{kg}\,\mathrm{s}^{-1}$ — caudal real
\item $T_{0e} = 271\,\mathrm{K}$ — temperatura total entrada
\item $p_{0e} = 82500\,\mathrm{Pa}$ — presión total entrada
\item $T_{\mathrm{ref}} = 288.15\,\mathrm{K}$ — temperatura referencia
\item $p_{\mathrm{ref}} = 101325\,\mathrm{Pa}$ — presión referencia
\end{itemize}

\noindent\textbf{Incógnita}
$\dot{m}_c$, en $\mathrm{kg\,s^{-1}}$; también las magnitudes intermedias desarrolladas.

\noindent\textbf{Modelo físico}

\noindent\textit{Corrección de caudal}

\[
\fitmath{\dot{m}_c=\frac{\dot m p_{\mathrm{ref}} \sqrt{\frac{T_{0e}}{T_{\mathrm{ref}}}}}{p_{0e}}}
\]

\noindent\textbf{Hipótesis}
\begin{itemize}
\item Mismo gas y referencias declaradas.
\end{itemize}

\noindent\textbf{Desarrollo paso a paso}

\begin{enumerate}

\item La corrección térmica es la raíz de la razón de temperaturas.

\[\fitmath{\theta=\sqrt{\frac{\mathrm{Tin}}{\mathrm{Tr}}}}\]

\[\fitmath{\theta=\sqrt{\frac{(271)}{(288.15)}}\approx 0.9697847\,}\]

\item La presión corrige la capacidad de paso del gas.

\[\fitmath{\mathrm{mc}=\frac{\mathrm{md}\cdot \sqrt{\frac{\mathrm{Tin}}{\mathrm{Tr}}}}{\frac{\mathrm{pin}}{\mathrm{pr}}}}\]

\[\fitmath{\mathrm{mc}=\frac{(33.7605)\cdot \sqrt{\frac{(271)}{(288.15)}}}{\frac{(82500)}{(101325)}}\approx 40.21118\,\mathrm{kg/s}}\]

\end{enumerate}


\noindent\textbf{Sustitución numérica}

\[
\fitmath{\frac{\left(33.7605\,\mathrm{kg}\,\mathrm{s}^{-1}\right) \left(101325\,\mathrm{Pa}\right) \sqrt{\frac{\left(271\,\mathrm{K}\right)}{\left(288.15\,\mathrm{K}\right)}}}{\left(82500\,\mathrm{Pa}\right)}}
\]

\noindent\textbf{Resultado}

\[
\fitmath{\dot{m}_c=4.021118460\times10^{1}\,\mathrm{kg}\,\mathrm{s}^{-1}}
\]

\noindent\textbf{Interpretación}
El caudal corregido es una variable de comparación. No es una nueva medición de masa real.

\noindent\textbf{Comprobación de los pasos}
\begin{itemize}
\item Paso 2: $0.969785$ — reproducido
\item Paso 4: $40.2112$ — reproducido
\end{itemize}

\noindent\textbf{Comprobación de la respuesta} la aritmética del último paso reproduce la respuesta enunciada a la precisión con que está escrita. Es una reproducción aritmética; no valida el modelo físico ni sus hipótesis.

\noindent\textbf{Limitaciones}
\begin{itemize}
\item Modelo educativo; la aplicabilidad a una máquina real requiere revisión experta.
\end{itemize}

\noindent\textbf{Errores frecuentes}
\begin{itemize}
\item Invertir la corrección de presión.
\end{itemize}
% studium:end-record eyJpZCI6ICJQUkItMDAyNCIsICJzZWN0aW9uIjogImMwMyJ9
\section{Capítulo 4: Turbinas axiales de acción}
\subsection*{C04-01-P}

\noindent\textbf{Enunciado}
\(C04-01-P \cdot\) Acción ideal. Ejercicio de transferencia. Una turbina de acción ideal recibe las componentes y velocidad periférica dadas. Halle remolino de salida, trabajo y descenso de temperatura total. Antes de consultar la solución, prediga el efecto de cambiar el primer dato respecto del caso resuelto de esta familia. Justifique la tendencia y explique qué restricción podría convertirse en limitante. Datos: \(C_\theta=492.03\, \mathrm{m}\,\mathrm{s}^{-1}\); \(U=211\, \mathrm{m}\,\mathrm{s}^{-1}\); \(c_p=1156.697\, \mathrm{J}\,\mathrm{kg}^{-1}\,\mathrm{K}^{-1}\).

\noindent\textbf{Datos}
\begin{itemize}
\item $C_\theta = 492.03\,\mathrm{m}\,\mathrm{s}^{-1}$ — remolino entrada rotor
\item $U = 211\,\mathrm{m}\,\mathrm{s}^{-1}$ — velocidad periférica
\item $c_p = 1156.696969697\,\mathrm{J}\,\mathrm{kg}^{-1}\,\mathrm{K}^{-1}$ — calor específico
\end{itemize}

\noindent\textbf{Incógnita}
$\Delta T_0$, en $\mathrm{K}$; también las magnitudes intermedias desarrolladas.

\noindent\textbf{Modelo físico}

\noindent\textit{Acción ideal}

\[
\fitmath{\Delta T_0=\frac{2 U \left(C_\theta - U\right)}{c_p}}
\]

\noindent\textbf{Hipótesis}
\begin{itemize}
\item Acción ideal, radio y componente axial constantes.
\end{itemize}

\noindent\textbf{Desarrollo paso a paso}

\begin{enumerate}

\item La simetría relativa da la componente tangencial de salida.

\[\fitmath{\mathrm{cts}=2\cdot \mathrm{uu}-\mathrm{ct}}\]

\[\fitmath{\mathrm{cts}=2\cdot (211)-(492.03)\approx -70.03\,\mathrm{m/s}}\]

\item Euler proporciona trabajo positivo de turbina.

\[\fitmath{\mathrm{wt}=2\cdot \mathrm{uu}\cdot \left(\mathrm{ct}-\mathrm{uu}\right)}\]

\[\fitmath{\mathrm{wt}=2\cdot (211)\cdot \left((492.03)-(211)\right)\approx 118594.7\,\mathrm{J/kg}}\]

\item Se divide por el calor específico.

\[\fitmath{\mathrm{dTz}=\frac{2\cdot \mathrm{uu}\cdot \left(\mathrm{ct}-\mathrm{uu}\right)}{\mathrm{cp}}}\]

\[\fitmath{\mathrm{dTz}=\frac{2\cdot (211)\cdot \left((492.03)-(211)\right)}{(1156.69697)}\approx 102.5287\,\mathrm{K}}\]

\end{enumerate}


\noindent\textbf{Sustitución numérica}

\[
\fitmath{\frac{2 \left(211\,\mathrm{m}\,\mathrm{s}^{-1}\right) \left(\left(492.03\,\mathrm{m}\,\mathrm{s}^{-1}\right) - \left(211\,\mathrm{m}\,\mathrm{s}^{-1}\right)\right)}{\left(1156.696969697\,\mathrm{J}\,\mathrm{kg}^{-1}\,\mathrm{K}^{-1}\right)}}
\]

\noindent\textbf{Resultado}

\[
\fitmath{\Delta T_0=1.025287202\times10^{2}\,\mathrm{K}}
\]

\noindent\textbf{Interpretación}
El remolino de salida es negativo con estos datos. Mantener su signo es necesario para cerrar Euler.

\noindent\textbf{Comprobación de los pasos}
\begin{itemize}
\item Paso 2: $-70.03$ — reproducido
\item Paso 4: $118595$ — reproducido
\item Paso 6: $102.529$ — reproducido
\end{itemize}

\noindent\textbf{Comprobación de la respuesta} la aritmética del último paso reproduce la respuesta enunciada a la precisión con que está escrita. Es una reproducción aritmética; no valida el modelo físico ni sus hipótesis.

\noindent\textbf{Limitaciones}
\begin{itemize}
\item Modelo educativo; la aplicabilidad a una máquina real requiere revisión experta.
\end{itemize}

\noindent\textbf{Errores frecuentes}
\begin{itemize}
\item Tomar valor absoluto del remolino antes de aplicar Euler.
\end{itemize}
% studium:end-record eyJpZCI6ICJQUkItMDAyNiIsICJzZWN0aW9uIjogImMwNCJ9
\subsection*{C04-02-P}

\noindent\textbf{Enunciado}
\(C04-02-P \cdot\) Óptimo de utilización. Ejercicio de transferencia. Para una velocidad de entrada al rotor y ángulo fijados, obtenga la velocidad periférica de máximo trabajo ideal y ese trabajo. Antes de consultar la solución, prediga el efecto de cambiar el primer dato respecto del caso resuelto de esta familia. Justifique la tendencia y explique qué restricción podría convertirse en limitante. Datos: \(C_e=628.35\, \mathrm{m}\,\mathrm{s}^{-1}\); \(\alpha_e=1.117011\, \mathrm{rad}\).

\noindent\textbf{Datos}
\begin{itemize}
\item $C_e = 628.35\,\mathrm{m}\,\mathrm{s}^{-1}$ — módulo entrada rotor
\item $\alpha_e = 1.1170107212764\,\mathrm{rad}$ — ángulo absoluto axial
\end{itemize}

\noindent\textbf{Incógnita}
$w_{\max}$, en $\mathrm{J\,kg^{-1}}$; también las magnitudes intermedias desarrolladas.

\noindent\textbf{Modelo físico}

\noindent\textit{Óptimo de utilización}

\[
\fitmath{w_{\max}=\frac{C_e^{2} \sin^{2}{\left(\alpha_e \right)}}{2}}
\]

\noindent\textbf{Hipótesis}
\begin{itemize}
\item Velocidad y orientación incidentes fijadas; rotor de acción ideal.
\end{itemize}

\noindent\textbf{Desarrollo paso a paso}

\begin{enumerate}

\item Se proyecta la velocidad tangencial.

\[\fitmath{\mathrm{ct}=\mathrm{Vel}\cdot \sin\left(\mathrm{ang}\right)}\]

\[\fitmath{\mathrm{ct}=(628.35)\cdot \sin\left((1.117010721)\right)\approx 564.7572\,\mathrm{m/s}}\]

\item El óptimo de la parábola está en U=Ctheta/2.

\[\fitmath{\mathrm{uopt}=\frac{\mathrm{Vel}\cdot \sin\left(\mathrm{ang}\right)}{2}}\]

\[\fitmath{\mathrm{uopt}=\frac{(628.35)\cdot \sin\left((1.117010721)\right)}{2}\approx 282.3786\,\mathrm{m/s}}\]

\item Se evalúa el máximo ideal.

\[\fitmath{\mathrm{wmax}=\frac{{\mathrm{Vel}}^{2}\cdot {\sin\left(\mathrm{ang}\right)}^{2}}{2}}\]

\[\fitmath{\mathrm{wmax}=\frac{{(628.35)}^{2}\cdot {\sin\left((1.117010721)\right)}^{2}}{2}\approx 159475.4\,\mathrm{J/kg}}\]

\end{enumerate}


\noindent\textbf{Sustitución numérica}

\[
\fitmath{\frac{\left(\left(628.35\,\mathrm{m}\,\mathrm{s}^{-1}\right)\right)^{2} \sin^{2}{\left(\left(1.1170107212764\,\mathrm{rad}\right) \right)}}{2}}
\]

\noindent\textbf{Resultado}

\[
\fitmath{w_{\max}=1.594753695\times10^{5}\,\mathrm{J}\,\mathrm{kg}^{-1}}
\]

\noindent\textbf{Interpretación}
Queda velocidad axial a la salida incluso en el óptimo. La utilización cinética no es una eficiencia total-total.

\noindent\textbf{Comprobación de los pasos}
\begin{itemize}
\item Paso 2: $564.757$ — reproducido
\item Paso 4: $282.379$ — reproducido
\item Paso 6: $159475$ — reproducido
\end{itemize}

\noindent\textbf{Comprobación de la respuesta} la aritmética del último paso reproduce la respuesta enunciada a la precisión con que está escrita. Es una reproducción aritmética; no valida el modelo físico ni sus hipótesis.

\noindent\textbf{Limitaciones}
\begin{itemize}
\item Modelo educativo; la aplicabilidad a una máquina real requiere revisión experta.
\end{itemize}

\noindent\textbf{Errores frecuentes}
\begin{itemize}
\item Suponer que el óptimo elimina toda velocidad de salida.
\end{itemize}
% studium:end-record eyJpZCI6ICJQUkItMDAyOCIsICJzZWN0aW9uIjogImMwNCJ9
\subsection*{C04-03-P}

\noindent\textbf{Enunciado}
\(C04-03-P \cdot\) Rotor con reducción de velocidad relativa. Ejercicio de transferencia. Un rotor recibe W tangencial de entrada positiva y descarga con dirección relativa simétrica y módulo reducido por k. Halle la potencia extraída. Antes de consultar la solución, prediga el efecto de cambiar el primer dato respecto del caso resuelto de esta familia. Justifique la tendencia y explique qué restricción podría convertirse en limitante. Datos: \(C_x=187.44\, \mathrm{m}\,\mathrm{s}^{-1}\); \(C_\theta=493\, \mathrm{m}\,\mathrm{s}^{-1}\); \(U=238\, \mathrm{m}\,\mathrm{s}^{-1}\); \(k_W=0.92\); \(\dot m=17.8\, \mathrm{kg}\,\mathrm{s}^{-1}\).

\noindent\textbf{Datos}
\begin{itemize}
\item $C_x = 187.44\,\mathrm{m}\,\mathrm{s}^{-1}$ — componente axial entrada
\item $C_\theta = 493\,\mathrm{m}\,\mathrm{s}^{-1}$ — remolino entrada
\item $U = 238\,\mathrm{m}\,\mathrm{s}^{-1}$ — velocidad periférica
\item $k_W = 0.92$ — cociente Wsal/Went
\item $\dot m = 17.8\,\mathrm{kg}\,\mathrm{s}^{-1}$ — caudal
\end{itemize}

\noindent\textbf{Incógnita}
$\dot{W}$, en $\mathrm{W}$; también las magnitudes intermedias desarrolladas.

\noindent\textbf{Modelo físico}

\noindent\textit{Rotor con reducción de velocidad relativa}

\[
\fitmath{\dot{W}=\dot m U \left(C_\theta + k_W \left(C_\theta - U\right) - U\right)}
\]

\noindent\textbf{Hipótesis}
\begin{itemize}
\item Radio constante; k y dirección relativa son datos del modelo.
\end{itemize}

\noindent\textbf{Desarrollo paso a paso}

\begin{enumerate}

\item La salida tangencial absoluta combina U y la componente relativa reducida.

\[\fitmath{\mathrm{cts}=\mathrm{uu}-\mathrm{kk}\cdot \left(\mathrm{ct}-\mathrm{uu}\right)}\]

\[\fitmath{\mathrm{cts}=(238)-(0.92)\cdot \left((493)-(238)\right)\approx 3.4\,\mathrm{m/s}}\]

\item Se calcula el trabajo de Euler.

\[\fitmath{\mathrm{wt}=\mathrm{uu}\cdot \left(\mathrm{ct}-\left(\mathrm{uu}-\mathrm{kk}\cdot \left(\mathrm{ct}-\mathrm{uu}\right)\right)\right)}\]

\[\fitmath{\mathrm{wt}=(238)\cdot \left((493)-\left((238)-(0.92)\cdot \left((493)-(238)\right)\right)\right)\approx 116524.8\,\mathrm{J/kg}}\]

\item El caudal convierte trabajo en potencia.

\[\fitmath{\mathrm{Pow}=\mathrm{md}\cdot \mathrm{uu}\cdot \left(\mathrm{ct}-\left(\mathrm{uu}-\mathrm{kk}\cdot \left(\mathrm{ct}-\mathrm{uu}\right)\right)\right)}\]

\[\fitmath{\mathrm{Pow}=(17.8)\cdot (238)\cdot \left((493)-\left((238)-(0.92)\cdot \left((493)-(238)\right)\right)\right)\approx 2074141\,\mathrm{W}}\]

\end{enumerate}


\noindent\textbf{Sustitución numérica}

\[
\fitmath{\left(17.8\,\mathrm{kg}\,\mathrm{s}^{-1}\right) \left(238\,\mathrm{m}\,\mathrm{s}^{-1}\right) \left(\left(493\,\mathrm{m}\,\mathrm{s}^{-1}\right) + \left(0.92\right) \left(\left(493\,\mathrm{m}\,\mathrm{s}^{-1}\right) - \left(238\,\mathrm{m}\,\mathrm{s}^{-1}\right)\right) - \left(238\,\mathrm{m}\,\mathrm{s}^{-1}\right)\right)}
\]

\noindent\textbf{Resultado}

\[
\fitmath{\dot{W}=2.074141440\times10^{6}\,\mathrm{W}}
\]

\noindent\textbf{Interpretación}
La componente axial de salida es k veces la de entrada al mantener el ángulo relativo simétrico. No se impone además acción estática perfecta.

\noindent\textbf{Comprobación de los pasos}
\begin{itemize}
\item Paso 2: $3.4$ — reproducido
\item Paso 4: $116525$ — reproducido
\item Paso 6: $2.07414\times 10^{6}$ — reproducido
\end{itemize}

\noindent\textbf{Comprobación de la respuesta} la aritmética del último paso reproduce la respuesta enunciada a la precisión con que está escrita. Es una reproducción aritmética; no valida el modelo físico ni sus hipótesis.

\noindent\textbf{Limitaciones}
\begin{itemize}
\item Modelo educativo; la aplicabilidad a una máquina real requiere revisión experta.
\end{itemize}

\noindent\textbf{Errores frecuentes}
\begin{itemize}
\item Mantener simultáneamente Cx constante y escala k de todo el vector.
\end{itemize}
% studium:end-record eyJpZCI6ICJQUkItMDAzMCIsICJzZWN0aW9uIjogImMwNCJ9
\subsection*{C04-04-P}

\noindent\textbf{Enunciado}
\(C04-04-P \cdot\) Carga de cascada de turbina. Ejercicio de transferencia. Estime el parámetro de carga Z de una fila con ángulos algebraicos dados, en el modelo de baja compresibilidad definido en el capítulo. Antes de consultar la solución, prediga el efecto de cambiar el primer dato respecto del caso resuelto de esta familia. Justifique la tendencia y explique qué restricción podría convertirse en limitante. Datos: \(\alpha_e=0.8736246\, \mathrm{rad}\); \(\alpha_s=-1.012291\, \mathrm{rad}\); \(s=0.024\, \mathrm{m}\); \(c_x=0.049\, \mathrm{m}\).

\noindent\textbf{Datos}
\begin{itemize}
\item $\alpha_e = 0.87362455708576\,\mathrm{rad}$ — ángulo entrada en referencia de fila
\item $\alpha_s = -1.0122909661567\,\mathrm{rad}$ — ángulo salida
\item $s = 0.024\,\mathrm{m}$ — paso
\item $c_x = 0.049\,\mathrm{m}$ — cuerda axial
\end{itemize}

\noindent\textbf{Incógnita}
$Z$ (adimensional); también las magnitudes intermedias desarrolladas.

\noindent\textbf{Modelo físico}

\noindent\textit{Carga de cascada de turbina}

\[
\fitmath{Z=\frac{2 s \left(\tan{\left(\alpha_e \right)} - \tan{\left(\alpha_s \right)}\right) \cos^{2}{\left(\alpha_s \right)}}{c_x}}
\]

\noindent\textbf{Hipótesis}
\begin{itemize}
\item Densidad y velocidad axial aproximadamente constantes; ángulos algebraicos.
\end{itemize}

\noindent\textbf{Desarrollo paso a paso}

\begin{enumerate}

\item Se calcula la razón de paso sobre cuerda axial.

\[\fitmath{\mathrm{sc}=\frac{\mathrm{ss}}{\mathrm{cc}}}\]

\[\fitmath{\mathrm{sc}=\frac{(0.024)}{(0.049)}\approx 0.4897959\,}\]

\item Se evalúa la carga a partir del cambio tangencial normalizado.

\[\fitmath{Z=2\cdot \frac{\mathrm{ss}}{\mathrm{cc}}\cdot {\cos\left(\mathrm{ps}\right)}^{2}\cdot \left(\tan\left(\mathrm{pe}\right)-\tan\left(\mathrm{ps}\right)\right)}\]

\[\fitmath{Z=2\cdot \frac{(0.024)}{(0.049)}\cdot {\cos\left((-1.012290966)\right)}^{2}\cdot \left(\tan\left((0.8736245571)\right)-\tan\left((-1.012290966)\right)\right)\approx 0.7686973\,}\]

\end{enumerate}


\noindent\textbf{Sustitución numérica}

\[
\fitmath{\frac{2 \left(0.024\,\mathrm{m}\right) \left(\tan{\left(\left(0.87362455708576\,\mathrm{rad}\right) \right)} - \tan{\left(\left(-1.0122909661567\,\mathrm{rad}\right) \right)}\right) \cos^{2}{\left(\left(-1.0122909661567\,\mathrm{rad}\right) \right)}}{\left(0.049\,\mathrm{m}\right)}}
\]

\noindent\textbf{Resultado}

\[
\fitmath{Z=7.686972619\times10^{-1}}
\]

\noindent\textbf{Interpretación}
El resultado es un indicador de carga del modelo. No permite asignar una pérdida ni una eficiencia sin correlación contrastada.

\noindent\textbf{Comprobación de los pasos}
\begin{itemize}
\item Paso 2: $0.489796$ — reproducido
\item Paso 4: $0.768697$ — reproducido
\end{itemize}

\noindent\textbf{Comprobación de la respuesta} la aritmética del último paso reproduce la respuesta enunciada a la precisión con que está escrita. Es una reproducción aritmética; no valida el modelo físico ni sus hipótesis.

\noindent\textbf{Limitaciones}
\begin{itemize}
\item Modelo educativo; la aplicabilidad a una máquina real requiere revisión experta.
\end{itemize}

\noindent\textbf{Errores frecuentes}
\begin{itemize}
\item Usar cuerda real donde el parámetro pide cuerda axial.
\end{itemize}
% studium:end-record eyJpZCI6ICJQUkItMDAzMiIsICJzZWN0aW9uIjogImMwNCJ9
\section{Capítulo 5: Turbinas de reacción y rendimiento del escalón}
\subsection*{C05-01-P}

\noindent\textbf{Enunciado}
\(C05-01-P \cdot\) Triángulos desde carga y reacción. Ejercicio de transferencia. Diseñe las componentes del escalón repetitivo con carga, reacción y velocidad periférica dadas. Obtenga la componente relativa tangencial a la salida del rotor. Antes de consultar la solución, prediga el efecto de cambiar el primer dato respecto del caso resuelto de esta familia. Justifique la tendencia y explique qué restricción podría convertirse en limitante. Datos: \(U=324.825\, \mathrm{m}\,\mathrm{s}^{-1}\); \(\psi=1.42\); \(\mathcal R=0.46\).

\noindent\textbf{Datos}
\begin{itemize}
\item $U = 324.825\,\mathrm{m}\,\mathrm{s}^{-1}$ — velocidad periférica
\item $\psi = 1.42$ — carga turbina
\item $\mathcal R = 0.46$ — reacción
\end{itemize}

\noindent\textbf{Incógnita}
$W_{\theta 3}$, en $\mathrm{m\,s^{-1}}$; también las magnitudes intermedias desarrolladas.

\noindent\textbf{Modelo físico}

\noindent\textit{Triángulos desde carga y reacción}

\[
\fitmath{W_{\theta 3}=U \left(- \frac{\psi}{2} - \mathcal R\right)}
\]

\noindent\textbf{Hipótesis}
\begin{itemize}
\item Radio y componente axial constantes; igual energía cinética en extremos del escalón.
\end{itemize}

\noindent\textbf{Desarrollo paso a paso}

\begin{enumerate}

\item La suma y diferencia de remolinos dan la entrada.

\[\fitmath{\mathrm{ct}_{2}=\mathrm{uu}\cdot \left(1-\mathrm{rea}+\frac{\psi}{2}\right)}\]

\[\fitmath{\mathrm{ct}_{2}=(324.825)\cdot \left(1-(0.46)+\frac{(1.42)}{2}\right)\approx 406.0312\,\mathrm{m/s}}\]

\item Se obtiene la descarga absoluta.

\[\fitmath{\mathrm{ct}_{3}=\mathrm{uu}\cdot \left(1-\mathrm{rea}-\frac{\psi}{2}\right)}\]

\[\fitmath{\mathrm{ct}_{3}=(324.825)\cdot \left(1-(0.46)-\frac{(1.42)}{2}\right)\approx -55.22025\,\mathrm{m/s}}\]

\item La referencia relativa resta U a la descarga.

\[\fitmath{\mathrm{wt}_{3}=\mathrm{uu}\cdot \left(-\mathrm{rea}-\frac{\psi}{2}\right)}\]

\[\fitmath{\mathrm{wt}_{3}=(324.825)\cdot \left(-(0.46)-\frac{(1.42)}{2}\right)\approx -380.0453\,\mathrm{m/s}}\]

\end{enumerate}


\noindent\textbf{Sustitución numérica}

\[
\fitmath{\left(324.825\,\mathrm{m}\,\mathrm{s}^{-1}\right) \left(- \frac{\left(1.42\right)}{2} - \left(0.46\right)\right)}
\]

\noindent\textbf{Resultado}

\[
\fitmath{W_{\theta 3}=-3.800452500\times10^{2}\,\mathrm{m}\,\mathrm{s}^{-1}}
\]

\noindent\textbf{Interpretación}
La componente relativa de salida es negativa; su módulo no debe sustituir su signo en la suma vectorial.

\noindent\textbf{Comprobación de los pasos}
\begin{itemize}
\item Paso 2: $406.031$ — reproducido
\item Paso 4: $-55.2202$ — reproducido
\item Paso 6: $-380.045$ — reproducido
\end{itemize}

\noindent\textbf{Comprobación de la respuesta} la aritmética del último paso reproduce la respuesta enunciada a la precisión con que está escrita. Es una reproducción aritmética; no valida el modelo físico ni sus hipótesis.

\noindent\textbf{Limitaciones}
\begin{itemize}
\item Modelo educativo; la aplicabilidad a una máquina real requiere revisión experta.
\end{itemize}

\noindent\textbf{Errores frecuentes}
\begin{itemize}
\item Elegir carga, reacción y descarga axial como tres datos independientes.
\end{itemize}
% studium:end-record eyJpZCI6ICJQUkItMDAzNCIsICJzZWN0aW9uIjogImMwNSJ9
\subsection*{C05-02-P}

\noindent\textbf{Enunciado}
\(C05-02-P \cdot\) Reacción general no repetitiva. Ejercicio de transferencia. Calcule la reacción estática de una turbina cuando las velocidades absolutas de entrada y salida del escalón difieren. Antes de consultar la solución, prediga el efecto de cambiar el primer dato respecto del caso resuelto de esta familia. Justifique la tendencia y explique qué restricción podría convertirse en limitante. Datos: \(U=305.655\, \mathrm{m}\,\mathrm{s}^{-1}\); \(C_{\theta2}=421\, \mathrm{m}\,\mathrm{s}^{-1}\); \(C_{\theta3}=-38\, \mathrm{m}\,\mathrm{s}^{-1}\); \(C_1=127\, \mathrm{m}\,\mathrm{s}^{-1}\); \(C_2=467\, \mathrm{m}\,\mathrm{s}^{-1}\); \(C_3=203\, \mathrm{m}\,\mathrm{s}^{-1}\).

\noindent\textbf{Datos}
\begin{itemize}
\item $U = 305.655\,\mathrm{m}\,\mathrm{s}^{-1}$ — velocidad periférica
\item $C_{\theta 2} = 421\,\mathrm{m}\,\mathrm{s}^{-1}$ — remolino entrada rotor
\item $C_{\theta 3} = -38\,\mathrm{m}\,\mathrm{s}^{-1}$ — remolino salida rotor
\item $c_{1} = 127\,\mathrm{m}\,\mathrm{s}^{-1}$ — módulo entrada escalón
\item $c_{2} = 467\,\mathrm{m}\,\mathrm{s}^{-1}$ — módulo entrada rotor
\item $c_{3} = 203\,\mathrm{m}\,\mathrm{s}^{-1}$ — módulo salida escalón
\end{itemize}

\noindent\textbf{Incógnita}
$\mathcal{R}$ (adimensional); también las magnitudes intermedias desarrolladas.

\noindent\textbf{Modelo físico}

\noindent\textit{Reacción general no repetitiva}

\[
\fitmath{\mathcal R=\frac{- \frac{c_{2}^{2}}{2} + \frac{c_{3}^{2}}{2} + U \left(C_{\theta 2} - C_{\theta 3}\right)}{- \frac{c_{1}^{2}}{2} + \frac{c_{3}^{2}}{2} + U \left(C_{\theta 2} - C_{\theta 3}\right)}}
\]

\noindent\textbf{Hipótesis}
\begin{itemize}
\item Rotor y estator adiabáticos; radios constantes para Euler.
\end{itemize}

\noindent\textbf{Desarrollo paso a paso}

\begin{enumerate}

\item Euler da el trabajo extraído.

\[\fitmath{\mathrm{wt}=\mathrm{uu}\cdot \left(\mathrm{ct}_{2}-\mathrm{ct}_{3}\right)}\]

\[\fitmath{\mathrm{wt}=(305.655)\cdot \left((421)-(-38)\right)\approx 140295.6\,\mathrm{J/kg}}\]

\item La caída estática del rotor incluye el cambio cinético.

\[\fitmath{\mathrm{dhr}=\mathrm{uu}\cdot \left(\mathrm{ct}_{2}-\mathrm{ct}_{3}\right)+\frac{{c_{3}}^{2}-{c_{2}}^{2}}{2}}\]

\[\fitmath{\mathrm{dhr}=(305.655)\cdot \left((421)-(-38)\right)+\frac{{(203)}^{2}-{(467)}^{2}}{2}\approx 51855.64\,\mathrm{J/kg}}\]

\item La caída estática del escalón tiene extremos 1 y 3.

\[\fitmath{\mathrm{rea}=\frac{\mathrm{uu}\cdot \left(\mathrm{ct}_{2}-\mathrm{ct}_{3}\right)+\frac{{c_{3}}^{2}-{c_{2}}^{2}}{2}}{\mathrm{uu}\cdot \left(\mathrm{ct}_{2}-\mathrm{ct}_{3}\right)+\frac{{c_{3}}^{2}-{c_{1}}^{2}}{2}}}\]

\[\fitmath{\mathrm{rea}=\frac{(305.655)\cdot \left((421)-(-38)\right)+\frac{{(203)}^{2}-{(467)}^{2}}{2}}{(305.655)\cdot \left((421)-(-38)\right)+\frac{{(203)}^{2}-{(127)}^{2}}{2}}\approx 0.3392903\,}\]

\end{enumerate}


\noindent\textbf{Sustitución numérica}

\[
\fitmath{\frac{- \frac{\left(\left(467\,\mathrm{m}\,\mathrm{s}^{-1}\right)\right)^{2}}{2} + \frac{\left(\left(203\,\mathrm{m}\,\mathrm{s}^{-1}\right)\right)^{2}}{2} + \left(305.655\,\mathrm{m}\,\mathrm{s}^{-1}\right) \left(\left(421\,\mathrm{m}\,\mathrm{s}^{-1}\right) - \left(-38\,\mathrm{m}\,\mathrm{s}^{-1}\right)\right)}{- \frac{\left(\left(127\,\mathrm{m}\,\mathrm{s}^{-1}\right)\right)^{2}}{2} + \frac{\left(\left(203\,\mathrm{m}\,\mathrm{s}^{-1}\right)\right)^{2}}{2} + \left(305.655\,\mathrm{m}\,\mathrm{s}^{-1}\right) \left(\left(421\,\mathrm{m}\,\mathrm{s}^{-1}\right) - \left(-38\,\mathrm{m}\,\mathrm{s}^{-1}\right)\right)}}
\]

\noindent\textbf{Resultado}

\[
\fitmath{\mathcal R=3.392902552\times10^{-1}}
\]

\noindent\textbf{Interpretación}
La reacción general difiere de la fórmula reducida si cambia la energía cinética entre extremos.

\noindent\textbf{Comprobación de los pasos}
\begin{itemize}
\item Paso 2: $140296$ — reproducido
\item Paso 4: $51855.6$ — reproducido
\item Paso 6: $0.33929$ — reproducido
\end{itemize}

\noindent\textbf{Comprobación de la respuesta} la aritmética del último paso reproduce la respuesta enunciada a la precisión con que está escrita. Es una reproducción aritmética; no valida el modelo físico ni sus hipótesis.

\noindent\textbf{Limitaciones}
\begin{itemize}
\item Modelo educativo; la aplicabilidad a una máquina real requiere revisión experta.
\end{itemize}

\noindent\textbf{Errores frecuentes}
\begin{itemize}
\item Usar caída total en el denominador de una reacción definida como estática.
\end{itemize}
% studium:end-record eyJpZCI6ICJQUkItMDAzNiIsICJzZWN0aW9uIjogImMwNSJ9
\subsection*{C05-03-P}

\noindent\textbf{Enunciado}
\(C05-03-P \cdot\) Presión de descarga de turbina. Ejercicio de transferencia. Dados trabajo y rendimiento total-total de un escalón, halle la presión total de salida. Antes de consultar la solución, prediga el efecto de cambiar el primer dato respecto del caso resuelto de esta familia. Justifique la tendencia y explique qué restricción podría convertirse en limitante. Datos: \(T_{0e}=1608.15\, \mathrm{K}\); \(p_{0e}=1240000\, \mathrm{Pa}\); \(w_t=156000\, \mathrm{J}\,\mathrm{kg}^{-1}\); \(\eta=0.885\); \(c_p=1156.697\, \mathrm{J}\,\mathrm{kg}^{-1}\,\mathrm{K}^{-1}\); \(\gamma=1.33\).

\noindent\textbf{Datos}
\begin{itemize}
\item $T_{0e} = 1608.15\,\mathrm{K}$ — temperatura total entrada
\item $p_{0e} = 1240000\,\mathrm{Pa}$ — presión total entrada
\item $w_t = 156000\,\mathrm{J}\,\mathrm{kg}^{-1}$ — trabajo extraído
\item $\eta = 0.885$ — rendimiento total-total
\item $c_p = 1156.696969697\,\mathrm{J}\,\mathrm{kg}^{-1}\,\mathrm{K}^{-1}$ — calor específico
\item $\gamma = 1.33$ — razón de calores
\end{itemize}

\noindent\textbf{Incógnita}
$p_{0s}$, en $\mathrm{Pa}$; también las magnitudes intermedias desarrolladas.

\noindent\textbf{Modelo físico}

\noindent\textit{Presión de descarga de turbina}

\[
\fitmath{p_{0s}=p_{0e} \left(1 - \frac{w_t}{T_{0e} c_p \eta}\right)^{\frac{\gamma}{\gamma - 1}}}
\]

\noindent\textbf{Hipótesis}
\begin{itemize}
\item Sin refrigeración ni sangrado; gas perfecto constante.
\end{itemize}

\noindent\textbf{Desarrollo paso a paso}

\begin{enumerate}

\item Se determina la caída de temperatura ideal equivalente.

\[\fitmath{\mathrm{dTs}=\frac{\mathrm{wt}}{\eta\cdot \mathrm{cp}}}\]

\[\fitmath{\mathrm{dTs}=\frac{(156000)}{(0.885)\cdot (1156.69697)}\approx 152.3918\,\mathrm{K}}\]

\item Se calcula la razón de temperaturas isentrópicas.

\[\fitmath{\mathrm{rr}=1-\frac{\mathrm{wt}}{\eta\cdot \mathrm{cp}\cdot \mathrm{Tin}}}\]

\[\fitmath{\mathrm{rr}=1-\frac{(156000)}{(0.885)\cdot (1156.69697)\cdot (1608.15)}\approx 0.9052378\,}\]

\item La razón de presiones corresponde al estado ideal equivalente.

\[\fitmath{\mathrm{pout}=\mathrm{pin}\cdot {\left(1-\frac{\mathrm{wt}}{\eta\cdot \mathrm{cp}\cdot \mathrm{Tin}}\right)}^{\frac{\gamma}{\gamma-1}}}\]

\[\fitmath{\mathrm{pout}=(1240000)\cdot {\left(1-\frac{(156000)}{(0.885)\cdot (1156.69697)\cdot (1608.15)}\right)}^{\frac{(1.33)}{(1.33)-1}}\approx 830160.7\,\mathrm{Pa}}\]

\end{enumerate}


\noindent\textbf{Sustitución numérica}

\[
\fitmath{\left(1240000\,\mathrm{Pa}\right) \left(1 - \frac{\left(156000\,\mathrm{J}\,\mathrm{kg}^{-1}\right)}{\left(1608.15\,\mathrm{K}\right) \left(1156.696969697\,\mathrm{J}\,\mathrm{kg}^{-1}\,\mathrm{K}^{-1}\right) \left(0.885\right)}\right)^{\frac{\left(1.33\right)}{\left(1.33\right) - 1}}}
\]

\noindent\textbf{Resultado}

\[
\fitmath{p_{0s}=8.301606971\times10^{5}\,\mathrm{Pa}}
\]

\noindent\textbf{Interpretación}
La base debe ser positiva. La presión resultante es total absoluta, no estática.

\noindent\textbf{Comprobación de los pasos}
\begin{itemize}
\item Paso 2: $152.392$ — reproducido
\item Paso 4: $0.905238$ — reproducido
\item Paso 6: $830161$ — reproducido
\end{itemize}

\noindent\textbf{Comprobación de la respuesta} la aritmética del último paso reproduce la respuesta enunciada a la precisión con que está escrita. Es una reproducción aritmética; no valida el modelo físico ni sus hipótesis.

\noindent\textbf{Limitaciones}
\begin{itemize}
\item Modelo educativo; la aplicabilidad a una máquina real requiere revisión experta.
\end{itemize}

\noindent\textbf{Errores frecuentes}
\begin{itemize}
\item Introducir eficiencia total-estática en una expresión total-total.
\end{itemize}
% studium:end-record eyJpZCI6ICJQUkItMDAzOCIsICJzZWN0aW9uIjogImMwNSJ9
\subsection*{C05-04-P}

\noindent\textbf{Enunciado}
\(C05-04-P \cdot\) Descarga axial y potencia. Ejercicio de transferencia. Una turbina repetitiva de reacción dada descarga sin remolino. Calcule su carga y potencia útil. Antes de consultar la solución, prediga el efecto de cambiar el primer dato respecto del caso resuelto de esta familia. Justifique la tendencia y explique qué restricción podría convertirse en limitante. Datos: \(\mathcal R=0.565\); \(U=322\, \mathrm{m}\,\mathrm{s}^{-1}\); \(\dot m=26\, \mathrm{kg}\,\mathrm{s}^{-1}\); \(\eta_m=0.975\).

\noindent\textbf{Datos}
\begin{itemize}
\item $\mathcal R = 0.565$ — reacción
\item $U = 322\,\mathrm{m}\,\mathrm{s}^{-1}$ — velocidad periférica
\item $\dot m = 26\,\mathrm{kg}\,\mathrm{s}^{-1}$ — caudal
\item $\eta_m = 0.975$ — rendimiento mecánico
\end{itemize}

\noindent\textbf{Incógnita}
$\dot{W}$, en $\mathrm{W}$; también las magnitudes intermedias desarrolladas.

\noindent\textbf{Modelo físico}

\noindent\textit{Descarga axial y potencia}

\[
\fitmath{\dot{W}=2 \eta_m \dot m U^{2} \left(1 - \mathcal R\right)}
\]

\noindent\textbf{Hipótesis}
\begin{itemize}
\item Escalón repetitivo, radio y Cx constantes, eficiencia mecánica separada.
\end{itemize}

\noindent\textbf{Desarrollo paso a paso}

\begin{enumerate}

\item La descarga axial liga la carga a la reacción.

\[\fitmath{\mathrm{psi}=2\cdot \left(1-\mathrm{rea}\right)}\]

\[\fitmath{\mathrm{psi}=2\cdot \left(1-(0.565)\right)\approx 0.87\,}\]

\item El trabajo se obtiene de la carga.

\[\fitmath{\mathrm{wt}=2\cdot \left(1-\mathrm{rea}\right)\cdot {\mathrm{uu}}^{2}}\]

\[\fitmath{\mathrm{wt}=2\cdot \left(1-(0.565)\right)\cdot {(322)}^{2}\approx 90205.08\,\mathrm{J/kg}}\]

\item Se aplica una sola vez el rendimiento mecánico.

\[\fitmath{\mathrm{Pow}=\mathrm{em}\cdot \mathrm{md}\cdot 2\cdot \left(1-\mathrm{rea}\right)\cdot {\mathrm{uu}}^{2}}\]

\[\fitmath{\mathrm{Pow}=(0.975)\cdot (26)\cdot 2\cdot \left(1-(0.565)\right)\cdot {(322)}^{2}\approx 2286699\,\mathrm{W}}\]

\end{enumerate}


\noindent\textbf{Sustitución numérica}

\[
\fitmath{2 \left(0.975\right) \left(26\,\mathrm{kg}\,\mathrm{s}^{-1}\right) \left(\left(322\,\mathrm{m}\,\mathrm{s}^{-1}\right)\right)^{2} \left(1 - \left(0.565\right)\right)}
\]

\noindent\textbf{Resultado}

\[
\fitmath{\dot{W}=2.286698778\times10^{6}\,\mathrm{W}}
\]

\noindent\textbf{Interpretación}
La condición de descarga impide variar carga y reacción independientemente.

\noindent\textbf{Comprobación de los pasos}
\begin{itemize}
\item Paso 2: $0.87$ — reproducido
\item Paso 4: $90205.1$ — reproducido
\item Paso 6: $2.2867\times 10^{6}$ — reproducido
\end{itemize}

\noindent\textbf{Comprobación de la respuesta} la aritmética del último paso reproduce la respuesta enunciada a la precisión con que está escrita. Es una reproducción aritmética; no valida el modelo físico ni sus hipótesis.

\noindent\textbf{Limitaciones}
\begin{itemize}
\item Modelo educativo; la aplicabilidad a una máquina real requiere revisión experta.
\end{itemize}

\noindent\textbf{Errores frecuentes}
\begin{itemize}
\item Multiplicar por eficiencia mecánica dos veces.
\end{itemize}
% studium:end-record eyJpZCI6ICJQUkItMDA0MCIsICJzZWN0aW9uIjogImMwNSJ9
\section{Capítulo 6: Caudal, bloqueo sónico y diseño de turbinas}
\subsection*{C06-01-P}

\noindent\textbf{Enunciado}
\(C06-01-P \cdot\) Área crítica. Ejercicio de transferencia. Determine el gasto máximo por área y el área crítica necesaria para el caudal dado. Antes de consultar la solución, prediga el efecto de cambiar el primer dato respecto del caso resuelto de esta familia. Justifique la tendencia y explique qué restricción podría convertirse en limitante. Datos: \(p_{0e}=724200\, \mathrm{Pa}\); \(T_{0e}=1320\, \mathrm{K}\); \(\gamma=1.33\); \(R_g=287\, \mathrm{J}\,\mathrm{kg}^{-1}\,\mathrm{K}^{-1}\); \(\dot m=24\, \mathrm{kg}\,\mathrm{s}^{-1}\).

\noindent\textbf{Datos}
\begin{itemize}
\item $p_{0e} = 724200\,\mathrm{Pa}$ — presión total garganta de referencia
\item $T_{0e} = 1320\,\mathrm{K}$ — temperatura total
\item $\gamma = 1.33$ — razón de calores
\item $R_g = 287\,\mathrm{J}\,\mathrm{kg}^{-1}\,\mathrm{K}^{-1}$ — constante del gas
\item $\dot m = 24\,\mathrm{kg}\,\mathrm{s}^{-1}$ — caudal
\end{itemize}

\noindent\textbf{Incógnita}
$A^{*}$, en $\mathrm{m^{2}}$; también las magnitudes intermedias desarrolladas.

\noindent\textbf{Modelo físico}

\noindent\textit{Área crítica}

\[
\fitmath{A^{*}=\frac{\sqrt{T_{0e}} \dot m \left(\frac{2}{\gamma + 1}\right)^{- \frac{\gamma + 1}{2 \gamma - 2}}}{p_{0e} \sqrt{\frac{\gamma}{R_g}}}}
\]

\noindent\textbf{Hipótesis}
\begin{itemize}
\item Garganta ideal unidimensional isentrópica en la referencia indicada.
\end{itemize}

\noindent\textbf{Desarrollo paso a paso}

\begin{enumerate}

\item Se evalúa el cuadrado del factor crítico a M=1.

\[\fitmath{F=\frac{\gamma}{R_g}\cdot {\left(\frac{2}{\gamma+1}\right)}^{\frac{\gamma+1}{\gamma-1}}}\]

\[\fitmath{F=\frac{(1.33)}{(287)}\cdot {\left(\frac{2}{(1.33)+1}\right)}^{\frac{(1.33)+1}{(1.33)-1}}\approx 0.001576408\,\mathrm{s^{2}\,K/m^{2}}}\]

\item Se calcula el flujo másico máximo por área.

\[\fitmath{\mathrm{flux}=\frac{\mathrm{pin}}{\sqrt{\mathrm{Tin}}}\cdot \sqrt{\frac{\gamma}{R_g}}\cdot {\left(\frac{2}{\gamma+1}\right)}^{\frac{\gamma+1}{2\cdot \left(\gamma-1\right)}}}\]

\[\fitmath{\mathrm{flux}=\frac{(724200)}{\sqrt{(1320)}}\cdot \sqrt{\frac{(1.33)}{(287)}}\cdot {\left(\frac{2}{(1.33)+1}\right)}^{\frac{(1.33)+1}{2\cdot \left((1.33)-1\right)}}\approx 791.4178\,\mathrm{kg/(m^{2}\,s)}}\]

\item Se dimensiona el área normal efectiva.

\[\fitmath{\mathrm{Ast}=\frac{\mathrm{md}}{\frac{\mathrm{pin}}{\sqrt{\mathrm{Tin}}}\cdot \sqrt{\frac{\gamma}{R_g}}\cdot {\left(\frac{2}{\gamma+1}\right)}^{\frac{\gamma+1}{2\cdot \left(\gamma-1\right)}}}}\]

\[\fitmath{\mathrm{Ast}=\frac{(24)}{\frac{(724200)}{\sqrt{(1320)}}\cdot \sqrt{\frac{(1.33)}{(287)}}\cdot {\left(\frac{2}{(1.33)+1}\right)}^{\frac{(1.33)+1}{2\cdot \left((1.33)-1\right)}}}\approx 0.03032532\,\mathrm{m^{2}}}\]

\end{enumerate}


\noindent\textbf{Sustitución numérica}

\[
\fitmath{\frac{\sqrt{\left(1320\,\mathrm{K}\right)} \left(24\,\mathrm{kg}\,\mathrm{s}^{-1}\right) \left(\frac{2}{\left(1.33\right) + 1}\right)^{- \frac{\left(1.33\right) + 1}{2 \left(1.33\right) - 2}}}{\left(724200\,\mathrm{Pa}\right) \sqrt{\frac{\left(1.33\right)}{\left(287\,\mathrm{J}\,\mathrm{kg}^{-1}\,\mathrm{K}^{-1}\right)}}}}
\]

\noindent\textbf{Resultado}

\[
\fitmath{A^{*}=3.032532263\times10^{-2}\,\mathrm{m}^{2}}
\]

\noindent\textbf{Interpretación}
El área es normal al flujo y efectiva. Un coeficiente de descarga inferior a uno exigiría mayor área geométrica.

\noindent\textbf{Comprobación de los pasos}
\begin{itemize}
\item Paso 2: $0.00157641$ — reproducido
\item Paso 4: $791.418$ — reproducido
\item Paso 6: $0.0303253$ — reproducido
\end{itemize}

\noindent\textbf{Comprobación de la respuesta} la aritmética del último paso reproduce la respuesta enunciada a la precisión con que está escrita. Es una reproducción aritmética; no valida el modelo físico ni sus hipótesis.

\noindent\textbf{Limitaciones}
\begin{itemize}
\item Modelo educativo; la aplicabilidad a una máquina real requiere revisión experta.
\end{itemize}

\noindent\textbf{Errores frecuentes}
\begin{itemize}
\item Usar área frontal anular como garganta sin proyección.
\end{itemize}
% studium:end-record eyJpZCI6ICJQUkItMDA0MiIsICJzZWN0aW9uIjogImMwNiJ9
\subsection*{C06-02-P}

\noindent\textbf{Enunciado}
\(C06-02-P \cdot\) Cociente crítico. Ejercicio de transferencia. Calcule la presión crítica estática de una tobera convergente ideal. Antes de consultar la solución, prediga el efecto de cambiar el primer dato respecto del caso resuelto de esta familia. Justifique la tendencia y explique qué restricción podría convertirse en limitante. Datos: \(p_{0e}=974475\, \mathrm{Pa}\); \(\gamma=1.32\).

\noindent\textbf{Datos}
\begin{itemize}
\item $p_{0e} = 974475\,\mathrm{Pa}$ — presión total depósito
\item $\gamma = 1.32$ — razón de calores
\end{itemize}

\noindent\textbf{Incógnita}
$p^{*}$, en $\mathrm{Pa}$; también las magnitudes intermedias desarrolladas.

\noindent\textbf{Modelo físico}

\noindent\textit{Cociente crítico}

\[
\fitmath{p^{*}=p_{0e} \left(\frac{2}{\gamma + 1}\right)^{\frac{\gamma}{\gamma - 1}}}
\]

\noindent\textbf{Hipótesis}
\begin{itemize}
\item Depósito y tobera convergente ideal; gas perfecto.
\end{itemize}

\noindent\textbf{Desarrollo paso a paso}

\begin{enumerate}

\item Se aplica la razón isentrópica a Mach uno.

\[\fitmath{\mathrm{ratio}={\left(\frac{2}{\gamma+1}\right)}^{\frac{\gamma}{\gamma-1}}}\]

\[\fitmath{\mathrm{ratio}={\left(\frac{2}{(1.32)+1}\right)}^{\frac{(1.32)}{(1.32)-1}}\approx 0.5421392\,}\]

\item Se obtiene la presión de garganta crítica.

\[\fitmath{\mathrm{pstar}=\mathrm{pin}\cdot {\left(\frac{2}{\gamma+1}\right)}^{\frac{\gamma}{\gamma-1}}}\]

\[\fitmath{\mathrm{pstar}=(974475)\cdot {\left(\frac{2}{(1.32)+1}\right)}^{\frac{(1.32)}{(1.32)-1}}\approx 528301.1\,\mathrm{Pa}}\]

\end{enumerate}


\noindent\textbf{Sustitución numérica}

\[
\fitmath{\left(974475\,\mathrm{Pa}\right) \left(\frac{2}{\left(1.32\right) + 1}\right)^{\frac{\left(1.32\right)}{\left(1.32\right) - 1}}}
\]

\noindent\textbf{Resultado}

\[
\fitmath{p^{*}=5.283010773\times10^{5}\,\mathrm{Pa}}
\]

\noindent\textbf{Interpretación}
Una contrapresión inferior permite bloqueo en la tobera ideal indicada; no define toda la expansión de una cascada real.

\noindent\textbf{Comprobación de los pasos}
\begin{itemize}
\item Paso 2: $0.542139$ — reproducido
\item Paso 4: $528301$ — reproducido
\end{itemize}

\noindent\textbf{Comprobación de la respuesta} la aritmética del último paso reproduce la respuesta enunciada a la precisión con que está escrita. Es una reproducción aritmética; no valida el modelo físico ni sus hipótesis.

\noindent\textbf{Limitaciones}
\begin{itemize}
\item Modelo educativo; la aplicabilidad a una máquina real requiere revisión experta.
\end{itemize}

\noindent\textbf{Errores frecuentes}
\begin{itemize}
\item Interpretar la razón crítica como pérdida de presión total.
\end{itemize}
% studium:end-record eyJpZCI6ICJQUkItMDA0NCIsICJzZWN0aW9uIjogImMwNiJ9
\subsection*{C06-03-P}

\noindent\textbf{Enunciado}
\(C06-03-P \cdot\) Reparto entero de escalones. Ejercicio de transferencia. Se ha elegido un número entero de escalones de igual trabajo. Calcule la carga real que cierra exactamente el trabajo global. Antes de consultar la solución, prediga el efecto de cambiar el primer dato respecto del caso resuelto de esta familia. Justifique la tendencia y explique qué restricción podría convertirse en limitante. Datos: \(w_t=408960\, \mathrm{J}\,\mathrm{kg}^{-1}\); \(N_{\mathrm{esc}}=3\); \(U=338\, \mathrm{m}\,\mathrm{s}^{-1}\).

\noindent\textbf{Datos}
\begin{itemize}
\item $w_t = 408960\,\mathrm{J}\,\mathrm{kg}^{-1}$ — trabajo total
\item $N_{\mathrm{esc}} = 3$ — número de escalones
\item $U = 338\,\mathrm{m}\,\mathrm{s}^{-1}$ — velocidad periférica común de referencia
\end{itemize}

\noindent\textbf{Incógnita}
$\psi$ (adimensional); también las magnitudes intermedias desarrolladas.

\noindent\textbf{Modelo físico}

\noindent\textit{Reparto entero de escalones}

\[
\fitmath{\psi=\frac{w_t}{N_{\mathrm{esc}} U^{2}}}
\]

\noindent\textbf{Hipótesis}
\begin{itemize}
\item Trabajo igual por escalón y misma U de referencia.
\end{itemize}

\noindent\textbf{Desarrollo paso a paso}

\begin{enumerate}

\item El reparto conserva la suma del trabajo.

\[\fitmath{\mathrm{we}=\frac{\mathrm{wt}}{\mathrm{nn}}}\]

\[\fitmath{\mathrm{we}=\frac{(408960)}{(3)}\approx 136320\,\mathrm{J/kg}}\]

\item Se calcula la carga efectiva tras elegir el entero.

\[\fitmath{\psi=\frac{\mathrm{wt}}{\mathrm{nn}\cdot {\mathrm{uu}}^{2}}}\]

\[\fitmath{\psi=\frac{(408960)}{(3)\cdot {(338)}^{2}}\approx 1.193236\,}\]

\end{enumerate}


\noindent\textbf{Sustitución numérica}

\[
\fitmath{\frac{\left(408960\,\mathrm{J}\,\mathrm{kg}^{-1}\right)}{\left(3\right) \left(\left(338\,\mathrm{m}\,\mathrm{s}^{-1}\right)\right)^{2}}}
\]

\noindent\textbf{Resultado}

\[
\fitmath{\psi=1.193235531\times10^{0}}
\]

\noindent\textbf{Interpretación}
El entero elegido obliga a recalcular psi. La factibilidad aerodinámica se revisa después.

\noindent\textbf{Comprobación de los pasos}
\begin{itemize}
\item Paso 2: $136320$ — reproducido
\item Paso 4: $1.19324$ — reproducido
\end{itemize}

\noindent\textbf{Comprobación de la respuesta} la aritmética del último paso reproduce la respuesta enunciada a la precisión con que está escrita. Es una reproducción aritmética; no valida el modelo físico ni sus hipótesis.

\noindent\textbf{Limitaciones}
\begin{itemize}
\item Modelo educativo; la aplicabilidad a una máquina real requiere revisión experta.
\end{itemize}

\noindent\textbf{Errores frecuentes}
\begin{itemize}
\item Redondear n y conservar una carga previa incompatible.
\end{itemize}
% studium:end-record eyJpZCI6ICJQUkItMDA0NiIsICJzZWN0aW9uIjogImMwNiJ9
\subsection*{C06-04-P}

\noindent\textbf{Enunciado}
\(C06-04-P \cdot\) Caudal de garganta oblicua. Ejercicio de transferencia. Una sección axial opera con módulo de velocidad sónico y ángulo dado. Halle el caudal usando su proyección axial. Antes de consultar la solución, prediga el efecto de cambiar el primer dato respecto del caso resuelto de esta familia. Justifique la tendencia y explique qué restricción podría convertirse en limitante. Datos: \(\rho=2.0448\, \mathrm{kg}\,\mathrm{m}^{-3}\); \(C=626\, \mathrm{m}\,\mathrm{s}^{-1}\); \(\alpha_e=1.064651\, \mathrm{rad}\); \(A=0.045\, \mathrm{m}^{2}\).

\noindent\textbf{Datos}
\begin{itemize}
\item $\rho = 2.0448\,\mathrm{kg}\,\mathrm{m}^{-3}$ — densidad estática
\item $C = 626\,\mathrm{m}\,\mathrm{s}^{-1}$ — módulo sónico local
\item $\alpha_e = 1.0646508437165\,\mathrm{rad}$ — ángulo con el eje
\item $A = 0.045\,\mathrm{m}^{2}$ — área anular frontal
\end{itemize}

\noindent\textbf{Incógnita}
$\dot{m}$, en $\mathrm{kg\,s^{-1}}$; también las magnitudes intermedias desarrolladas.

\noindent\textbf{Modelo físico}

\noindent\textit{Caudal de garganta oblicua}

\[
\fitmath{\dot m=A \rho C \cos{\left(\alpha_e \right)}}
\]

\noindent\textbf{Hipótesis}
\begin{itemize}
\item Corte axial con densidad y velocidad uniformes.
\end{itemize}

\noindent\textbf{Desarrollo paso a paso}

\begin{enumerate}

\item Se proyecta el módulo sobre la normal axial.

\[\fitmath{\mathrm{cx}=\mathrm{vel}\cdot \cos\left(\mathrm{ang}\right)}\]

\[\fitmath{\mathrm{cx}=(626)\cdot \cos\left((1.064650844)\right)\approx 303.4908\,\mathrm{m/s}}\]

\item La continuidad usa esa componente.

\[\fitmath{\mathrm{md}=\rho\cdot \mathrm{vel}\cdot \cos\left(\mathrm{ang}\right)\cdot \mathrm{area}}\]

\[\fitmath{\mathrm{md}=(2.0448)\cdot (626)\cdot \cos\left((1.064650844)\right)\cdot (0.045)\approx 27.92601\,\mathrm{kg/s}}\]

\end{enumerate}


\noindent\textbf{Sustitución numérica}

\[
\fitmath{\left(0.045\,\mathrm{m}^{2}\right) \left(2.0448\,\mathrm{kg}\,\mathrm{m}^{-3}\right) \left(626\,\mathrm{m}\,\mathrm{s}^{-1}\right) \cos{\left(\left(1.0646508437165\,\mathrm{rad}\right) \right)}}
\]

\noindent\textbf{Resultado}

\[
\fitmath{\dot m=2.792601150\times10^{1}\,\mathrm{kg}\,\mathrm{s}^{-1}}
\]

\noindent\textbf{Interpretación}
Mach axial y Mach del módulo son distintos. La garganta real puede tener otra orientación y área.

\noindent\textbf{Comprobación de los pasos}
\begin{itemize}
\item Paso 2: $303.491$ — reproducido
\item Paso 4: $27.926$ — reproducido
\end{itemize}

\noindent\textbf{Comprobación de la respuesta} la aritmética del último paso reproduce la respuesta enunciada a la precisión con que está escrita. Es una reproducción aritmética; no valida el modelo físico ni sus hipótesis.

\noindent\textbf{Limitaciones}
\begin{itemize}
\item Modelo educativo; la aplicabilidad a una máquina real requiere revisión experta.
\end{itemize}

\noindent\textbf{Errores frecuentes}
\begin{itemize}
\item Omitir el coseno por estar el módulo a Mach uno.
\end{itemize}
% studium:end-record eyJpZCI6ICJQUkItMDA0OCIsICJzZWN0aW9uIjogImMwNiJ9
\section{Capítulo 7: Escalón de compresor axial al radio medio}
\subsection*{C07-01-P}

\noindent\textbf{Enunciado}
\(C07-01-P \cdot\) Trabajo y presión del escalón. Ejercicio de transferencia. Dados los parámetros del escalón compresor, calcule trabajo, incremento total de temperatura y relación de presión. Antes de consultar la solución, prediga el efecto de cambiar el primer dato respecto del caso resuelto de esta familia. Justifique la tendencia y explique qué restricción podría convertirse en limitante. Datos: \(U=290.745\, \mathrm{m}\,\mathrm{s}^{-1}\); \(\psi=0.39\); \(T_{0e}=301\, \mathrm{K}\); \(\eta=0.88\); \(c_p=1004.5\, \mathrm{J}\,\mathrm{kg}^{-1}\,\mathrm{K}^{-1}\); \(\gamma=1.4\).

\noindent\textbf{Datos}
\begin{itemize}
\item $U = 290.745\,\mathrm{m}\,\mathrm{s}^{-1}$ — velocidad periférica
\item $\psi = 0.39$ — carga
\item $T_{0e} = 301\,\mathrm{K}$ — temperatura total entrada
\item $\eta = 0.88$ — eficiencia total-total
\item $c_p = 1004.5\,\mathrm{J}\,\mathrm{kg}^{-1}\,\mathrm{K}^{-1}$ — calor específico
\item $\gamma = 1.4$ — razón de calores
\end{itemize}

\noindent\textbf{Incógnita}
$\pi_c$ (adimensional); también las magnitudes intermedias desarrolladas.

\noindent\textbf{Modelo físico}

\noindent\textit{Trabajo y presión del escalón}

\[
\fitmath{\pi_c=\left(1 + \frac{\eta \psi U^{2}}{T_{0e} c_p}\right)^{\frac{\gamma}{\gamma - 1}}}
\]

\noindent\textbf{Hipótesis}
\begin{itemize}
\item Escalón adiabático; eficiencia global y gas perfecto constantes.
\end{itemize}

\noindent\textbf{Desarrollo paso a paso}

\begin{enumerate}

\item La carga define el trabajo comunicado.

\[\fitmath{\mathrm{wc}=\mathrm{psi}\cdot {\mathrm{uu}}^{2}}\]

\[\fitmath{\mathrm{wc}=(0.39)\cdot {(290.745)}^{2}\approx 32967.74\,\mathrm{J/kg}}\]

\item El balance energético fija el incremento térmico.

\[\fitmath{\mathrm{dT}=\frac{\mathrm{psi}\cdot {\mathrm{uu}}^{2}}{\mathrm{cp}}}\]

\[\fitmath{\mathrm{dT}=\frac{(0.39)\cdot {(290.745)}^{2}}{(1004.5)}\approx 32.82005\,\mathrm{K}}\]

\item El trabajo ideal equivalente fija la relación de presión.

\[\fitmath{\mathrm{pr}={\left(1+\frac{\mathrm{eta}\cdot \mathrm{psi}\cdot {\mathrm{uu}}^{2}}{\mathrm{cp}\cdot \mathrm{Tin}}\right)}^{\frac{\gamma}{\gamma-1}}}\]

\[\fitmath{\mathrm{pr}={\left(1+\frac{(0.88)\cdot (0.39)\cdot {(290.745)}^{2}}{(1004.5)\cdot (301)}\right)}^{\frac{(1.4)}{(1.4)-1}}\approx 1.378068\,}\]

\end{enumerate}


\noindent\textbf{Sustitución numérica}

\[
\fitmath{\left(1 + \frac{\left(0.88\right) \left(0.39\right) \left(\left(290.745\,\mathrm{m}\,\mathrm{s}^{-1}\right)\right)^{2}}{\left(301\,\mathrm{K}\right) \left(1004.5\,\mathrm{J}\,\mathrm{kg}^{-1}\,\mathrm{K}^{-1}\right)}\right)^{\frac{\left(1.4\right)}{\left(1.4\right) - 1}}}
\]

\noindent\textbf{Resultado}

\[
\fitmath{\pi_c=1.378068393\times10^{0}}
\]

\noindent\textbf{Interpretación}
La relación debe ser mayor que uno. Compare el trabajo reconstruido desde la temperatura con Euler.

\noindent\textbf{Comprobación de los pasos}
\begin{itemize}
\item Paso 2: $32967.7$ — reproducido
\item Paso 4: $32.82$ — reproducido
\item Paso 6: $1.37807$ — reproducido
\end{itemize}

\noindent\textbf{Comprobación de la respuesta} la aritmética del último paso reproduce la respuesta enunciada a la precisión con que está escrita. Es una reproducción aritmética; no valida el modelo físico ni sus hipótesis.

\noindent\textbf{Limitaciones}
\begin{itemize}
\item Modelo educativo; la aplicabilidad a una máquina real requiere revisión experta.
\end{itemize}

\noindent\textbf{Errores frecuentes}
\begin{itemize}
\item Aplicar la razón de temperaturas reales como si fuera isentrópica.
\end{itemize}
% studium:end-record eyJpZCI6ICJQUkItMDA1MCIsICJzZWN0aW9uIjogImMwNyJ9
\subsection*{C07-02-P}

\noindent\textbf{Enunciado}
\(C07-02-P \cdot\) Reacción y remolinos. Ejercicio de transferencia. Construya un rotor compresor repetitivo a partir de reacción y carga. La incógnita final es el ángulo relativo de salida, algebraico en radianes. Antes de consultar la solución, prediga el efecto de cambiar el primer dato respecto del caso resuelto de esta familia. Justifique la tendencia y explique qué restricción podría convertirse en limitante. Datos: \(U=315.24\, \mathrm{m}\,\mathrm{s}^{-1}\); \(\psi=0.43\); \(\mathcal R=0.55\); \(C_x=164\, \mathrm{m}\,\mathrm{s}^{-1}\).

\noindent\textbf{Datos}
\begin{itemize}
\item $U = 315.24\,\mathrm{m}\,\mathrm{s}^{-1}$ — velocidad periférica
\item $\psi = 0.43$ — carga
\item $\mathcal R = 0.55$ — reacción
\item $C_x = 164\,\mathrm{m}\,\mathrm{s}^{-1}$ — velocidad axial
\end{itemize}

\noindent\textbf{Incógnita}
$\beta_2$, en $\mathrm{rad}$; también las magnitudes intermedias desarrolladas.

\noindent\textbf{Modelo físico}

\noindent\textit{Reacción y remolinos}

\[
\fitmath{\beta_2=\operatorname{atan}{\left(\frac{U \left(\frac{\psi}{2} - \mathcal R\right)}{C_x} \right)}}
\]

\noindent\textbf{Hipótesis}
\begin{itemize}
\item Radio y velocidad axial constantes; C3 al cuadrado igual a C1 al cuadrado.
\end{itemize}

\noindent\textbf{Desarrollo paso a paso}

\begin{enumerate}

\item La suma y diferencia fijan el remolino de entrada.

\[\fitmath{ct1=\mathrm{uu}\cdot \left(1-\mathrm{rea}-\frac{\mathrm{psi}}{2}\right)}\]

\[\fitmath{ct1=(315.24)\cdot \left(1-(0.55)-\frac{(0.43)}{2}\right)\approx 74.0814\,\mathrm{m/s}}\]

\item Se calcula la componente absoluta de salida.

\[\fitmath{ct2=\mathrm{uu}\cdot \left(1-\mathrm{rea}+\frac{\mathrm{psi}}{2}\right)}\]

\[\fitmath{ct2=(315.24)\cdot \left(1-(0.55)+\frac{(0.43)}{2}\right)\approx 209.6346\,\mathrm{m/s}}\]

\item Se resta U y se obtiene el ángulo desde el eje.

\[\fitmath{beta2=\arctan\left(\frac{\mathrm{uu}\cdot \left(-\mathrm{rea}+\frac{\mathrm{psi}}{2}\right)}{\mathrm{cx}}\right)}\]

\[\fitmath{beta2=\arctan\left(\frac{(315.24)\cdot \left(-(0.55)+\frac{(0.43)}{2}\right)}{(164)}\right)\approx -0.5721\,\mathrm{rad}}\]

\end{enumerate}


\noindent\textbf{Sustitución numérica}

\[
\fitmath{\operatorname{atan}{\left(\frac{\left(315.24\,\mathrm{m}\,\mathrm{s}^{-1}\right) \left(\frac{\left(0.43\right)}{2} - \left(0.55\right)\right)}{\left(164\,\mathrm{m}\,\mathrm{s}^{-1}\right)} \right)}}
\]

\noindent\textbf{Resultado}

\[
\fitmath{\beta_2=-5.721000360\times10^{-1}\,\mathrm{rad}}
\]

\noindent\textbf{Interpretación}
La beta negativa se conserva en la construcción. La reacción dada solo coincide con la definición bajo la hipótesis repetitiva.

\noindent\textbf{Comprobación de los pasos}
\begin{itemize}
\item Paso 2: $74.0814$ — reproducido
\item Paso 4: $209.635$ — reproducido
\item Paso 6: $-0.5721$ — reproducido
\end{itemize}

\noindent\textbf{Comprobación de la respuesta} la aritmética del último paso reproduce la respuesta enunciada a la precisión con que está escrita. Es una reproducción aritmética; no valida el modelo físico ni sus hipótesis.

\noindent\textbf{Limitaciones}
\begin{itemize}
\item Modelo educativo; la aplicabilidad a una máquina real requiere revisión experta.
\end{itemize}

\noindent\textbf{Errores frecuentes}
\begin{itemize}
\item Usar el módulo positivo de beta de otro convenio.
\end{itemize}
% studium:end-record eyJpZCI6ICJQUkItMDA1MiIsICJzZWN0aW9uIjogImMwNyJ9
\subsection*{C07-03-P}

\noindent\textbf{Enunciado}
\(C07-03-P \cdot\) Mach relativo. Ejercicio de transferencia. Calcule temperatura estática de entrada y Mach relativo del rotor a partir de remanso y componentes absolutas. Antes de consultar la solución, prediga el efecto de cambiar el primer dato respecto del caso resuelto de esta familia. Justifique la tendencia y explique qué restricción podría convertirse en limitante. Datos: \(T_0=329.085\, \mathrm{K}\); \(C_x=172\, \mathrm{m}\,\mathrm{s}^{-1}\); \(C_\theta=43\, \mathrm{m}\,\mathrm{s}^{-1}\); \(U=319\, \mathrm{m}\,\mathrm{s}^{-1}\); \(c_p=1004.5\, \mathrm{J}\,\mathrm{kg}^{-1}\,\mathrm{K}^{-1}\); \(\gamma=1.4\); \(R_g=287\, \mathrm{J}\,\mathrm{kg}^{-1}\,\mathrm{K}^{-1}\).

\noindent\textbf{Datos}
\begin{itemize}
\item $T_0 = 329.085\,\mathrm{K}$ — temperatura total absoluta
\item $C_x = 172\,\mathrm{m}\,\mathrm{s}^{-1}$ — componente axial
\item $C_\theta = 43\,\mathrm{m}\,\mathrm{s}^{-1}$ — remolino absoluto
\item $U = 319\,\mathrm{m}\,\mathrm{s}^{-1}$ — velocidad periférica
\item $c_p = 1004.5\,\mathrm{J}\,\mathrm{kg}^{-1}\,\mathrm{K}^{-1}$ — calor específico
\item $\gamma = 1.4$ — razón de calores
\item $R_g = 287\,\mathrm{J}\,\mathrm{kg}^{-1}\,\mathrm{K}^{-1}$ — constante del gas
\end{itemize}

\noindent\textbf{Incógnita}
$\mathrm{Ma}_r$ (adimensional); también las magnitudes intermedias desarrolladas.

\noindent\textbf{Modelo físico}

\noindent\textit{Mach relativo}

\[
\fitmath{\mathrm{Ma}_r=\frac{\sqrt{C_x^{2} + \left(C_\theta - U\right)^{2}}}{\sqrt{R_g \gamma \left(T_0 - \frac{C_\theta^{2} + C_x^{2}}{2 c_p}\right)}}}
\]

\noindent\textbf{Hipótesis}
\begin{itemize}
\item Gas perfecto y referencia local de rotor.
\end{itemize}

\noindent\textbf{Desarrollo paso a paso}

\begin{enumerate}

\item La energía cinética absoluta determina la temperatura estática.

\[\fitmath{T=T_0-\frac{{\mathrm{cx}}^{2}+{\mathrm{ct}}^{2}}{2\cdot \mathrm{cp}}}\]

\[\fitmath{T=(329.085)-\frac{{(172)}^{2}+{(43)}^{2}}{2\cdot (1004.5)}\approx 313.4389\,\mathrm{K}}\]

\item El módulo relativo se construye vectorialmente.

\[\fitmath{W=\sqrt{{\mathrm{cx}}^{2}+{\left(\mathrm{ct}-\mathrm{uu}\right)}^{2}}}\]

\[\fitmath{W=\sqrt{{(172)}^{2}+{\left((43)-(319)\right)}^{2}}\approx 325.2076\,\mathrm{m/s}}\]

\item El sonido depende de temperatura estática.

\[\fitmath{\mathrm{Mr}=\frac{\sqrt{{\mathrm{cx}}^{2}+{\left(\mathrm{ct}-\mathrm{uu}\right)}^{2}}}{\sqrt{\gamma\cdot R_g\cdot \left(T_0-\frac{{\mathrm{cx}}^{2}+{\mathrm{ct}}^{2}}{2\cdot \mathrm{cp}}\right)}}}\]

\[\fitmath{\mathrm{Mr}=\frac{\sqrt{{(172)}^{2}+{\left((43)-(319)\right)}^{2}}}{\sqrt{(1.4)\cdot (287)\cdot \left((329.085)-\frac{{(172)}^{2}+{(43)}^{2}}{2\cdot (1004.5)}\right)}}\approx 0.9163878\,}\]

\end{enumerate}


\noindent\textbf{Sustitución numérica}

\[
\fitmath{\frac{\sqrt{\left(\left(172\,\mathrm{m}\,\mathrm{s}^{-1}\right)\right)^{2} + \left(\left(43\,\mathrm{m}\,\mathrm{s}^{-1}\right) - \left(319\,\mathrm{m}\,\mathrm{s}^{-1}\right)\right)^{2}}}{\sqrt{\left(287\,\mathrm{J}\,\mathrm{kg}^{-1}\,\mathrm{K}^{-1}\right) \left(1.4\right) \left(\left(329.085\,\mathrm{K}\right) - \frac{\left(\left(43\,\mathrm{m}\,\mathrm{s}^{-1}\right)\right)^{2} + \left(\left(172\,\mathrm{m}\,\mathrm{s}^{-1}\right)\right)^{2}}{2 \left(1004.5\,\mathrm{J}\,\mathrm{kg}^{-1}\,\mathrm{K}^{-1}\right)}\right)}}}
\]

\noindent\textbf{Resultado}

\[
\fitmath{\mathrm{Ma}_r=9.163878079\times10^{-1}}
\]

\noindent\textbf{Interpretación}
El Mach relativo puede superar al absoluto. La propiedad acústica utiliza la misma temperatura estática en ambas referencias.

\noindent\textbf{Comprobación de los pasos}
\begin{itemize}
\item Paso 2: $313.439$ — reproducido
\item Paso 4: $325.208$ — reproducido
\item Paso 6: $0.916388$ — reproducido
\end{itemize}

\noindent\textbf{Comprobación de la respuesta} la aritmética del último paso reproduce la respuesta enunciada a la precisión con que está escrita. Es una reproducción aritmética; no valida el modelo físico ni sus hipótesis.

\noindent\textbf{Limitaciones}
\begin{itemize}
\item Modelo educativo; la aplicabilidad a una máquina real requiere revisión experta.
\end{itemize}

\noindent\textbf{Errores frecuentes}
\begin{itemize}
\item Usar temperatura total en la velocidad del sonido.
\end{itemize}
% studium:end-record eyJpZCI6ICJQUkItMDA1NCIsICJzZWN0aW9uIjogImMwNyJ9
\subsection*{C07-04-P}

\noindent\textbf{Enunciado}
\(C07-04-P \cdot\) Cierre de reacción por entalpías. Ejercicio de transferencia. Un rotor tiene los remolinos indicados y velocidad axial constante; el escalón es repetitivo. Calcule su reacción desde incrementos estáticos. Antes de consultar la solución, prediga el efecto de cambiar el primer dato respecto del caso resuelto de esta familia. Justifique la tendencia y explique qué restricción podría convertirse en limitante. Datos: \(U=302.46\, \mathrm{m}\,\mathrm{s}^{-1}\); \(C_{\theta e}=61\, \mathrm{m}\,\mathrm{s}^{-1}\); \(C_{\theta s}=172\, \mathrm{m}\,\mathrm{s}^{-1}\).

\noindent\textbf{Datos}
\begin{itemize}
\item $U = 302.46\,\mathrm{m}\,\mathrm{s}^{-1}$ — velocidad periférica
\item $C_{\theta e} = 61\,\mathrm{m}\,\mathrm{s}^{-1}$ — remolino entrada rotor
\item $C_{\theta s} = 172\,\mathrm{m}\,\mathrm{s}^{-1}$ — remolino salida rotor
\end{itemize}

\noindent\textbf{Incógnita}
$\mathcal{R}$ (adimensional); también las magnitudes intermedias desarrolladas.

\noindent\textbf{Modelo físico}

\noindent\textit{Cierre de reacción por entalpías}

\[
\fitmath{\mathcal R=\frac{\frac{C_{\theta e}^{2}}{2} - \frac{C_{\theta s}^{2}}{2} + U \left(- C_{\theta e} + C_{\theta s}\right)}{U \left(- C_{\theta e} + C_{\theta s}\right)}}
\]

\noindent\textbf{Hipótesis}
\begin{itemize}
\item Radio y Cx constantes, igual velocidad absoluta entre extremos del escalón.
\end{itemize}

\noindent\textbf{Desarrollo paso a paso}

\begin{enumerate}

\item Euler da el incremento total de entalpía.

\[\fitmath{\mathrm{dhz}=\mathrm{uu}\cdot \left(\mathrm{cs}-\mathrm{ce}\right)}\]

\[\fitmath{\mathrm{dhz}=(302.46)\cdot \left((172)-(61)\right)\approx 33573.06\,\mathrm{J/kg}}\]

\item Se resta el incremento de energía cinética absoluta.

\[\fitmath{\mathrm{dh}=\mathrm{uu}\cdot \left(\mathrm{cs}-\mathrm{ce}\right)-\frac{{\mathrm{cs}}^{2}-{\mathrm{ce}}^{2}}{2}}\]

\[\fitmath{\mathrm{dh}=(302.46)\cdot \left((172)-(61)\right)-\frac{{(172)}^{2}-{(61)}^{2}}{2}\approx 20641.56\,\mathrm{J/kg}}\]

\item Se forma la razón de incremento estático del rotor y del escalón.

\[\fitmath{\mathrm{rea}=\frac{\mathrm{uu}\cdot \left(\mathrm{cs}-\mathrm{ce}\right)-\frac{{\mathrm{cs}}^{2}-{\mathrm{ce}}^{2}}{2}}{\mathrm{uu}\cdot \left(\mathrm{cs}-\mathrm{ce}\right)}}\]

\[\fitmath{\mathrm{rea}=\frac{(302.46)\cdot \left((172)-(61)\right)-\frac{{(172)}^{2}-{(61)}^{2}}{2}}{(302.46)\cdot \left((172)-(61)\right)}\approx 0.6148251\,}\]

\end{enumerate}


\noindent\textbf{Sustitución numérica}

\[
\fitmath{\frac{\frac{\left(\left(61\,\mathrm{m}\,\mathrm{s}^{-1}\right)\right)^{2}}{2} - \frac{\left(\left(172\,\mathrm{m}\,\mathrm{s}^{-1}\right)\right)^{2}}{2} + \left(302.46\,\mathrm{m}\,\mathrm{s}^{-1}\right) \left(- \left(61\,\mathrm{m}\,\mathrm{s}^{-1}\right) + \left(172\,\mathrm{m}\,\mathrm{s}^{-1}\right)\right)}{\left(302.46\,\mathrm{m}\,\mathrm{s}^{-1}\right) \left(- \left(61\,\mathrm{m}\,\mathrm{s}^{-1}\right) + \left(172\,\mathrm{m}\,\mathrm{s}^{-1}\right)\right)}}
\]

\noindent\textbf{Resultado}

\[
\fitmath{\mathcal R=6.148251008\times10^{-1}}
\]

\noindent\textbf{Interpretación}
El resultado debe coincidir con 1-(ce+cs)/(2U). Esa identidad es una comprobación algebraica, no experimental.

\noindent\textbf{Comprobación de los pasos}
\begin{itemize}
\item Paso 2: $33573.1$ — reproducido
\item Paso 4: $20641.6$ — reproducido
\item Paso 6: $0.614825$ — reproducido
\end{itemize}

\noindent\textbf{Comprobación de la respuesta} la aritmética del último paso reproduce la respuesta enunciada a la precisión con que está escrita. Es una reproducción aritmética; no valida el modelo físico ni sus hipótesis.

\noindent\textbf{Limitaciones}
\begin{itemize}
\item Modelo educativo; la aplicabilidad a una máquina real requiere revisión experta.
\end{itemize}

\noindent\textbf{Errores frecuentes}
\begin{itemize}
\item Confundir el incremento estático de rotor y el total.
\end{itemize}
% studium:end-record eyJpZCI6ICJQUkItMDA1NiIsICJzZWN0aW9uIjogImMwNyJ9
\section{Capítulo 8: Cascadas: difusión, incidencia, desviación y pérdidas}
\subsection*{C08-01-P}

\noindent\textbf{Enunciado}
\(C08-01-P \cdot\) Indicador de difusión. Ejercicio de transferencia. Estime D para una fila móvil con módulos relativos y cambio tangencial dados. Antes de consultar la solución, prediga el efecto de cambiar el primer dato respecto del caso resuelto de esta familia. Justifique la tendencia y explique qué restricción podría convertirse en limitante. Datos: \(W_e=300.33\, \mathrm{m}\,\mathrm{s}^{-1}\); \(W_s=216\, \mathrm{m}\,\mathrm{s}^{-1}\); \(|\Delta W_\theta|=97\, \mathrm{m}\,\mathrm{s}^{-1}\); \(\sigma=1.24\).

\noindent\textbf{Datos}
\begin{itemize}
\item $W_e = 300.33\,\mathrm{m}\,\mathrm{s}^{-1}$ — módulo relativo entrada
\item $W_s = 216\,\mathrm{m}\,\mathrm{s}^{-1}$ — módulo relativo salida
\item $|\Delta W_\theta| = 97\,\mathrm{m}\,\mathrm{s}^{-1}$ — magnitud de cambio tangencial relativo
\item $\sigma = 1.24$ — solidez
\end{itemize}

\noindent\textbf{Incógnita}
$D$ (adimensional); también las magnitudes intermedias desarrolladas.

\noindent\textbf{Modelo físico}

\noindent\textit{Indicador de difusión}

\[
\fitmath{D=\frac{|\Delta W_\theta|}{2 \sigma W_e} + 1 - \frac{W_s}{W_e}}
\]

\noindent\textbf{Hipótesis}
\begin{itemize}
\item Factor de difusión definido en el capítulo, referencia relativa.
\end{itemize}

\noindent\textbf{Desarrollo paso a paso}

\begin{enumerate}

\item Se calcula la contribución de deceleración.

\[\fitmath{\mathrm{dec}=1-\frac{\mathrm{vs}}{\mathrm{ve}}}\]

\[\fitmath{\mathrm{dec}=1-\frac{(216)}{(300.33)}\approx 0.2807911\,}\]

\item La carga tangencial añade la segunda contribución.

\[\fitmath{D=1-\frac{\mathrm{vs}}{\mathrm{ve}}+\frac{\mathrm{dv}}{2\cdot \mathrm{sol}\cdot \mathrm{ve}}}\]

\[\fitmath{D=1-\frac{(216)}{(300.33)}+\frac{(97)}{2\cdot (1.24)\cdot (300.33)}\approx 0.4110242\,}\]

\end{enumerate}


\noindent\textbf{Sustitución numérica}

\[
\fitmath{\frac{\left(97\,\mathrm{m}\,\mathrm{s}^{-1}\right)}{2 \left(1.24\right) \left(300.33\,\mathrm{m}\,\mathrm{s}^{-1}\right)} + 1 - \frac{\left(216\,\mathrm{m}\,\mathrm{s}^{-1}\right)}{\left(300.33\,\mathrm{m}\,\mathrm{s}^{-1}\right)}}
\]

\noindent\textbf{Resultado}

\[
\fitmath{D=4.110242174\times10^{-1}}
\]

\noindent\textbf{Interpretación}
D es un indicador preliminar. La aceptabilidad exige criterio de perfil y régimen declarado.

\noindent\textbf{Comprobación de los pasos}
\begin{itemize}
\item Paso 2: $0.280791$ — reproducido
\item Paso 4: $0.411024$ — reproducido
\end{itemize}

\noindent\textbf{Comprobación de la respuesta} la aritmética del último paso reproduce la respuesta enunciada a la precisión con que está escrita. Es una reproducción aritmética; no valida el modelo físico ni sus hipótesis.

\noindent\textbf{Limitaciones}
\begin{itemize}
\item Modelo educativo; la aplicabilidad a una máquina real requiere revisión experta.
\end{itemize}

\noindent\textbf{Errores frecuentes}
\begin{itemize}
\item Usar velocidades absolutas para un rotor.
\end{itemize}
% studium:end-record eyJpZCI6ICJQUkItMDA1OCIsICJzZWN0aW9uIjogImMwOCJ9
\subsection*{C08-02-P}

\noindent\textbf{Enunciado}
\(C08-02-P \cdot\) Solidez mínima del modelo. Ejercicio de transferencia. Un diseño impone D no superior al valor indicado. Calcule la solidez que satisface exactamente el límite y compruebe la positividad del denominador. Antes de consultar la solución, prediga el efecto de cambiar el primer dato respecto del caso resuelto de esta familia. Justifique la tendencia y explique qué restricción podría convertirse en limitante. Datos: \(W_e=320.565\, \mathrm{m}\,\mathrm{s}^{-1}\); \(W_s=232\, \mathrm{m}\,\mathrm{s}^{-1}\); \(|\Delta W_\theta|=112\, \mathrm{m}\,\mathrm{s}^{-1}\); \(D_{\mathrm{lim}}=0.47\).

\noindent\textbf{Datos}
\begin{itemize}
\item $W_e = 320.565\,\mathrm{m}\,\mathrm{s}^{-1}$ — módulo relativo entrada
\item $W_s = 232\,\mathrm{m}\,\mathrm{s}^{-1}$ — módulo relativo salida
\item $|\Delta W_\theta| = 112\,\mathrm{m}\,\mathrm{s}^{-1}$ — cambio tangencial de magnitud positiva
\item $D_{\mathrm{lim}} = 0.47$ — límite educativo de difusión
\end{itemize}

\noindent\textbf{Incógnita}
$\sigma$ (adimensional); también las magnitudes intermedias desarrolladas.

\noindent\textbf{Modelo físico}

\noindent\textit{Solidez mínima del modelo}

\[
\fitmath{\sigma=\frac{|\Delta W_\theta|}{2 W_e \left(D_{\mathrm{lim}} - 1 + \frac{W_s}{W_e}\right)}}
\]

\noindent\textbf{Hipótesis}
\begin{itemize}
\item Criterio de difusión prescrito como dato; sin garantía de pérdida real.
\end{itemize}

\noindent\textbf{Desarrollo paso a paso}

\begin{enumerate}

\item Se determina el margen restante después de la deceleración.

\[\fitmath{\mathrm{den}=\mathrm{lim}-1+\frac{\mathrm{vs}}{\mathrm{ve}}}\]

\[\fitmath{\mathrm{den}=(0.47)-1+\frac{(232)}{(320.565)}\approx 0.1937222\,}\]

\item Se despeja la solidez.

\[\fitmath{\mathrm{sol}=\frac{\mathrm{dv}}{2\cdot \mathrm{ve}\cdot \left(\mathrm{lim}-1+\frac{\mathrm{vs}}{\mathrm{ve}}\right)}}\]

\[\fitmath{\mathrm{sol}=\frac{(112)}{2\cdot (320.565)\cdot \left((0.47)-1+\frac{(232)}{(320.565)}\right)}\approx 0.9017633\,}\]

\end{enumerate}


\noindent\textbf{Sustitución numérica}

\[
\fitmath{\frac{\left(112\,\mathrm{m}\,\mathrm{s}^{-1}\right)}{2 \left(320.565\,\mathrm{m}\,\mathrm{s}^{-1}\right) \left(\left(0.47\right) - 1 + \frac{\left(232\,\mathrm{m}\,\mathrm{s}^{-1}\right)}{\left(320.565\,\mathrm{m}\,\mathrm{s}^{-1}\right)}\right)}}
\]

\noindent\textbf{Resultado}

\[
\fitmath{\sigma=9.017633499\times10^{-1}}
\]

\noindent\textbf{Interpretación}
Si el denominador fuera no positivo, este giro y cociente de módulos no podrían cumplir el límite por aumentar solidez.

\noindent\textbf{Comprobación de los pasos}
\begin{itemize}
\item Paso 2: $0.193722$ — reproducido
\item Paso 4: $0.901763$ — reproducido
\end{itemize}

\noindent\textbf{Comprobación de la respuesta} la aritmética del último paso reproduce la respuesta enunciada a la precisión con que está escrita. Es una reproducción aritmética; no valida el modelo físico ni sus hipótesis.

\noindent\textbf{Limitaciones}
\begin{itemize}
\item Modelo educativo; la aplicabilidad a una máquina real requiere revisión experta.
\end{itemize}

\noindent\textbf{Errores frecuentes}
\begin{itemize}
\item Aceptar una solidez negativa.
\end{itemize}
% studium:end-record eyJpZCI6ICJQUkItMDA2MCIsICJzZWN0aW9uIjogImMwOCJ9
\subsection*{C08-03-P}

\noindent\textbf{Enunciado}
\(C08-03-P \cdot\) Número entero y solidez real. Ejercicio de transferencia. Se eligen Z álabes enteros para un radio y una cuerda dados. Calcule paso y solidez real tras la elección. Antes de consultar la solución, prediga el efecto de cambiar el primer dato respecto del caso resuelto de esta familia. Justifique la tendencia y explique qué restricción podría convertirse en limitante. Datos: \(r=0.280095\, \mathrm{m}\); \(c=0.052\, \mathrm{m}\); \(Z=41\).

\noindent\textbf{Datos}
\begin{itemize}
\item $r = 0.280095\,\mathrm{m}$ — radio
\item $c = 0.052\,\mathrm{m}$ — cuerda
\item $Z = 41$ — número entero de álabes
\end{itemize}

\noindent\textbf{Incógnita}
$\sigma$ (adimensional); también las magnitudes intermedias desarrolladas.

\noindent\textbf{Modelo físico}

\noindent\textit{Número entero y solidez real}

\[
\fitmath{\sigma=\frac{c Z}{2 \pi r}}
\]

\noindent\textbf{Hipótesis}
\begin{itemize}
\item Fila uniforme al radio de referencia.
\end{itemize}

\noindent\textbf{Desarrollo paso a paso}

\begin{enumerate}

\item La circunferencia se reparte en pasos iguales.

\[\fitmath{\mathrm{ss}=\frac{2\cdot \pi\cdot \mathrm{rad}}{\mathrm{zz}}}\]

\[\fitmath{\mathrm{ss}=\frac{2\cdot \pi\cdot (0.280095)}{(41)}\approx 0.04292412\,\mathrm{m}}\]

\item La solidez usa el paso realmente obtenido.

\[\fitmath{\mathrm{sol}=\frac{\mathrm{cc}\cdot \mathrm{zz}}{2\cdot \pi\cdot \mathrm{rad}}}\]

\[\fitmath{\mathrm{sol}=\frac{(0.052)\cdot (41)}{2\cdot \pi\cdot (0.280095)}\approx 1.21144\,}\]

\end{enumerate}


\noindent\textbf{Sustitución numérica}

\[
\fitmath{\frac{\left(0.052\,\mathrm{m}\right) \left(41\right)}{2 \pi \left(0.280095\,\mathrm{m}\right)}}
\]

\noindent\textbf{Resultado}

\[
\fitmath{\sigma=1.211440185\times10^{0}}
\]

\noindent\textbf{Interpretación}
El redondeo del número de álabes modifica la solidez; debe evaluarse después el parámetro de carga.

\noindent\textbf{Comprobación de los pasos}
\begin{itemize}
\item Paso 2: $0.0429241$ — reproducido
\item Paso 4: $1.21144$ — reproducido
\end{itemize}

\noindent\textbf{Comprobación de la respuesta} la aritmética del último paso reproduce la respuesta enunciada a la precisión con que está escrita. Es una reproducción aritmética; no valida el modelo físico ni sus hipótesis.

\noindent\textbf{Limitaciones}
\begin{itemize}
\item Modelo educativo; la aplicabilidad a una máquina real requiere revisión experta.
\end{itemize}

\noindent\textbf{Errores frecuentes}
\begin{itemize}
\item Conservar exactamente la solidez objetivo después del redondeo.
\end{itemize}
% studium:end-record eyJpZCI6ICJQUkItMDA2MiIsICJzZWN0aW9uIjogImMwOCJ9
\subsection*{C08-04-P}

\noindent\textbf{Enunciado}
\(C08-04-P \cdot\) Incidencia y desviación algebraicas. Ejercicio de transferencia. En una fila fija se miden ángulos de gas y se conocen los metálicos. Determine incidencia y desviación; la incógnita final es delta. Antes de consultar la solución, prediga el efecto de cambiar el primer dato respecto del caso resuelto de esta familia. Justifique la tendencia y explique qué restricción podría convertirse en limitante. Datos: \(\alpha_e=0.9108001\, \mathrm{rad}\); \(\alpha_s=0.3839724\, \mathrm{rad}\); \(\kappa_e=0.8028515\, \mathrm{rad}\); \(\kappa_s=0.3141593\, \mathrm{rad}\).

\noindent\textbf{Datos}
\begin{itemize}
\item $\alpha_e = 0.91080007015324\,\mathrm{rad}$ — ángulo gas entrada
\item $\alpha_s = 0.38397243543875\,\mathrm{rad}$ — ángulo gas salida
\item $\kappa_e = 0.80285145591739\,\mathrm{rad}$ — ángulo metal entrada
\item $\kappa_s = 0.31415926535898\,\mathrm{rad}$ — ángulo metal salida
\end{itemize}

\noindent\textbf{Incógnita}
$\delta$, en $\mathrm{rad}$; también las magnitudes intermedias desarrolladas.

\noindent\textbf{Modelo físico}

\noindent\textit{Incidencia y desviación algebraicas}

\[
\fitmath{\delta=\alpha_s - \kappa_s}
\]

\noindent\textbf{Hipótesis}
\begin{itemize}
\item Ángulos desde el eje y mismo sentido positivo.
\end{itemize}

\noindent\textbf{Desarrollo paso a paso}

\begin{enumerate}

\item La incidencia compara gas y metal en entrada.

\[\fitmath{\mathrm{inc}=\mathrm{ae}-\mathrm{ke}}\]

\[\fitmath{\mathrm{inc}=(0.9108000702)-(0.8028514559)\approx 0.1079486\,\mathrm{rad}}\]

\item La desviación compara gas y metal en salida según nuestro convenio.

\[\fitmath{\mathrm{delta}=\mathrm{asv}-\mathrm{ks}}\]

\[\fitmath{\mathrm{delta}=(0.3839724354)-(0.3141592654)\approx 0.06981317\,\mathrm{rad}}\]

\end{enumerate}


\noindent\textbf{Sustitución numérica}

\[
\fitmath{\left(0.38397243543875\,\mathrm{rad}\right) - \left(0.31415926535898\,\mathrm{rad}\right)}
\]

\noindent\textbf{Resultado}

\[
\fitmath{\delta=6.981317008\times10^{-2}\,\mathrm{rad}}
\]

\noindent\textbf{Interpretación}
Los signos dependen de una definición explícita. Para expresar en grados multiplique los radianes por 180/pi.

\noindent\textbf{Comprobación de los pasos}
\begin{itemize}
\item Paso 2: $0.107949$ — reproducido
\item Paso 4: $0.0698132$ — reproducido
\end{itemize}

\noindent\textbf{Comprobación de la respuesta} la aritmética del último paso reproduce la respuesta enunciada a la precisión con que está escrita. Es una reproducción aritmética; no valida el modelo físico ni sus hipótesis.

\noindent\textbf{Limitaciones}
\begin{itemize}
\item Modelo educativo; la aplicabilidad a una máquina real requiere revisión experta.
\end{itemize}

\noindent\textbf{Errores frecuentes}
\begin{itemize}
\item Comparar ángulo absoluto con beta relativa.
\end{itemize}
% studium:end-record eyJpZCI6ICJQUkItMDA2NCIsICJzZWN0aW9uIjogImMwOCJ9
\section{Capítulo 9: Compresión multiescalón, anillo y acoplamiento}
\subsection*{C09-01-P}

\noindent\textbf{Enunciado}
\(C09-01-P \cdot\) Compresión de varios escalones. Ejercicio de transferencia. Una compresión politrópica global alimenta un eje. Halle temperatura final y potencia requerida en el eje. Antes de consultar la solución, prediga el efecto de cambiar el primer dato respecto del caso resuelto de esta familia. Justifique la tendencia y explique qué restricción podría convertirse en limitante. Datos: \(T_{0e}=307.785\, \mathrm{K}\); \(\pi_c=9.3\); \(\eta_p=0.9\); \(\gamma=1.4\); \(c_p=1004.5\, \mathrm{J}\,\mathrm{kg}^{-1}\,\mathrm{K}^{-1}\); \(\dot m=23\, \mathrm{kg}\,\mathrm{s}^{-1}\); \(\eta_m=0.98\).

\noindent\textbf{Datos}
\begin{itemize}
\item $T_{0e} = 307.785\,\mathrm{K}$ — temperatura total entrada
\item $\pi_c = 9.3$ — relación global de presión
\item $\eta_p = 0.9$ — eficiencia politrópica
\item $\gamma = 1.4$ — razón de calores
\item $c_p = 1004.5\,\mathrm{J}\,\mathrm{kg}^{-1}\,\mathrm{K}^{-1}$ — calor específico
\item $\dot m = 23\,\mathrm{kg}\,\mathrm{s}^{-1}$ — caudal
\item $\eta_m = 0.98$ — rendimiento mecánico compresor
\end{itemize}

\noindent\textbf{Incógnita}
$\dot{W}$, en $\mathrm{W}$; también las magnitudes intermedias desarrolladas.

\noindent\textbf{Modelo físico}

\noindent\textit{Compresión de varios escalones}

\[
\fitmath{\dot{W}=\frac{T_{0e} c_p \dot m \left(\pi_c^{\frac{\gamma - 1}{\eta_p \gamma}} - 1\right)}{\eta_m}}
\]

\noindent\textbf{Hipótesis}
\begin{itemize}
\item Propiedades y eficiencia politrópica constantes, sin sangrado.
\end{itemize}

\noindent\textbf{Desarrollo paso a paso}

\begin{enumerate}

\item Se integra el modelo politrópico global.

\[\fitmath{\mathrm{Tout}=\mathrm{Tin}\cdot {\mathrm{pr}}^{\frac{\gamma-1}{\gamma\cdot \mathrm{ep}}}}\]

\[\fitmath{\mathrm{Tout}=(307.785)\cdot {(9.3)}^{\frac{(1.4)-1}{(1.4)\cdot (0.9)}}\approx 624.7444\,\mathrm{K}}\]

\item La potencia de eje supera la potencia comunicada al fluido.

\[\fitmath{\mathrm{Pow}=\frac{\mathrm{md}\cdot \mathrm{cp}\cdot \mathrm{Tin}\cdot \left({\mathrm{pr}}^{\frac{\gamma-1}{\gamma\cdot \mathrm{ep}}}-1\right)}{\mathrm{em}}}\]

\[\fitmath{\mathrm{Pow}=\frac{(23)\cdot (1004.5)\cdot (307.785)\cdot \left({(9.3)}^{\frac{(1.4)-1}{(1.4)\cdot (0.9)}}-1\right)}{(0.98)}\approx 7472317\,\mathrm{W}}\]

\end{enumerate}


\noindent\textbf{Sustitución numérica}

\[
\fitmath{\frac{\left(307.785\,\mathrm{K}\right) \left(1004.5\,\mathrm{J}\,\mathrm{kg}^{-1}\,\mathrm{K}^{-1}\right) \left(23\,\mathrm{kg}\,\mathrm{s}^{-1}\right) \left(\left(9.3\right)^{\frac{\left(1.4\right) - 1}{\left(0.9\right) \left(1.4\right)}} - 1\right)}{\left(0.98\right)}}
\]

\noindent\textbf{Resultado}

\[
\fitmath{\dot{W}=7.472317275\times10^{6}\,\mathrm{W}}
\]

\noindent\textbf{Interpretación}
No se multiplica la eficiencia de cada escalón. El caudal y el incremento total de entalpía fijan la potencia.

\noindent\textbf{Comprobación de los pasos}
\begin{itemize}
\item Paso 2: $624.744$ — reproducido
\item Paso 4: $7.47232\times 10^{6}$ — reproducido
\end{itemize}

\noindent\textbf{Comprobación de la respuesta} la aritmética del último paso reproduce la respuesta enunciada a la precisión con que está escrita. Es una reproducción aritmética; no valida el modelo físico ni sus hipótesis.

\noindent\textbf{Limitaciones}
\begin{itemize}
\item Modelo educativo; la aplicabilidad a una máquina real requiere revisión experta.
\end{itemize}

\noindent\textbf{Errores frecuentes}
\begin{itemize}
\item Usar el producto de eficiencias como eficiencia global.
\end{itemize}
% studium:end-record eyJpZCI6ICJQUkItMDA2NiIsICJzZWN0aW9uIjogImMwOSJ9
\subsection*{C09-02-P}

\noindent\textbf{Enunciado}
\(C09-02-P \cdot\) Altura de anillo con bloqueo. Ejercicio de transferencia. Dimensione el área geométrica y la altura del anillo con un factor B de área disponible. Antes de consultar la solución, prediga el efecto de cambiar el primer dato respecto del caso resuelto de esta familia. Justifique la tendencia y explique qué restricción podría convertirse en limitante. Datos: \(\dot m=28.755\, \mathrm{kg}\,\mathrm{s}^{-1}\); \(\rho=3.6\, \mathrm{kg}\,\mathrm{m}^{-3}\); \(C_x=151\, \mathrm{m}\,\mathrm{s}^{-1}\); \(B=0.955\); \(r_m=0.285\, \mathrm{m}\).

\noindent\textbf{Datos}
\begin{itemize}
\item $\dot m = 28.755\,\mathrm{kg}\,\mathrm{s}^{-1}$ — caudal
\item $\rho = 3.6\,\mathrm{kg}\,\mathrm{m}^{-3}$ — densidad estática
\item $C_x = 151\,\mathrm{m}\,\mathrm{s}^{-1}$ — velocidad axial efectiva
\item $B = 0.955$ — fracción de área disponible
\item $r_m = 0.285\,\mathrm{m}$ — radio medio aritmético
\end{itemize}

\noindent\textbf{Incógnita}
$b$, en $\mathrm{m}$; también las magnitudes intermedias desarrolladas.

\noindent\textbf{Modelo físico}

\noindent\textit{Altura de anillo con bloqueo}

\[
\fitmath{b=\frac{\dot m}{2 \pi B C_x \rho r_m}}
\]

\noindent\textbf{Hipótesis}
\begin{itemize}
\item Flujo uniforme en el área disponible.
\end{itemize}

\noindent\textbf{Desarrollo paso a paso}

\begin{enumerate}

\item La continuidad se aplica al área efectiva B A.

\[\fitmath{\mathrm{area}=\frac{\mathrm{md}}{\rho\cdot \mathrm{cx}\cdot \mathrm{bb}}}\]

\[\fitmath{\mathrm{area}=\frac{(28.755)}{(3.6)\cdot (151)\cdot (0.955)}\approx 0.0553899\,\mathrm{m^{2}}}\]

\item Se convierte el área geométrica en altura.

\[\fitmath{b=\frac{\mathrm{md}}{\rho\cdot \mathrm{cx}\cdot \mathrm{bb}\cdot 2\cdot \pi\cdot \mathrm{rm}}}\]

\[\fitmath{b=\frac{(28.755)}{(3.6)\cdot (151)\cdot (0.955)\cdot 2\cdot \pi\cdot (0.285)}\approx 0.03093184\,\mathrm{m}}\]

\end{enumerate}


\noindent\textbf{Sustitución numérica}

\[
\fitmath{\frac{\left(28.755\,\mathrm{kg}\,\mathrm{s}^{-1}\right)}{2 \pi \left(0.955\right) \left(151\,\mathrm{m}\,\mathrm{s}^{-1}\right) \left(3.6\,\mathrm{kg}\,\mathrm{m}^{-3}\right) \left(0.285\,\mathrm{m}\right)}}
\]

\noindent\textbf{Resultado}

\[
\fitmath{b=3.093184491\times10^{-2}\,\mathrm{m}}
\]

\noindent\textbf{Interpretación}
B inferior a uno exige mayor área geométrica. La raíz debe permanecer positiva al restar b/2 al radio medio.

\noindent\textbf{Comprobación de los pasos}
\begin{itemize}
\item Paso 2: $0.0553899$ — reproducido
\item Paso 4: $0.0309318$ — reproducido
\end{itemize}

\noindent\textbf{Comprobación de la respuesta} la aritmética del último paso reproduce la respuesta enunciada a la precisión con que está escrita. Es una reproducción aritmética; no valida el modelo físico ni sus hipótesis.

\noindent\textbf{Limitaciones}
\begin{itemize}
\item Modelo educativo; la aplicabilidad a una máquina real requiere revisión experta.
\end{itemize}

\noindent\textbf{Errores frecuentes}
\begin{itemize}
\item Confundir B con la fracción bloqueada 1-B.
\end{itemize}
% studium:end-record eyJpZCI6ICJQUkItMDA2OCIsICJzZWN0aW9uIjogImMwOSJ9
\subsection*{C09-03-P}

\noindent\textbf{Enunciado}
\(C09-03-P \cdot\) Trabajo por escalón y razón de presión. Ejercicio de transferencia. Un escalón tardío recibe gas más caliente. Calcule su relación de presión para un trabajo y eficiencia fijos. Antes de consultar la solución, prediga el efecto de cambiar el primer dato respecto del caso resuelto de esta familia. Justifique la tendencia y explique qué restricción podría convertirse en limitante. Datos: \(T_{0e}=518.655\, \mathrm{K}\); \(w_c=29800\, \mathrm{J}\,\mathrm{kg}^{-1}\); \(\eta=0.87\); \(\gamma=1.4\); \(c_p=1004.5\, \mathrm{J}\,\mathrm{kg}^{-1}\,\mathrm{K}^{-1}\).

\noindent\textbf{Datos}
\begin{itemize}
\item $T_{0e} = 518.655\,\mathrm{K}$ — temperatura total entrada
\item $w_c = 29800\,\mathrm{J}\,\mathrm{kg}^{-1}$ — trabajo específico
\item $\eta = 0.87$ — eficiencia total-total
\item $\gamma = 1.4$ — razón de calores
\item $c_p = 1004.5\,\mathrm{J}\,\mathrm{kg}^{-1}\,\mathrm{K}^{-1}$ — calor específico
\end{itemize}

\noindent\textbf{Incógnita}
$\pi_c$ (adimensional); también las magnitudes intermedias desarrolladas.

\noindent\textbf{Modelo físico}

\noindent\textit{Trabajo por escalón y razón de presión}

\[
\fitmath{\pi_c=\left(1 + \frac{\eta w_c}{T_{0e} c_p}\right)^{\frac{\gamma}{\gamma - 1}}}
\]

\noindent\textbf{Hipótesis}
\begin{itemize}
\item Gas perfecto constante y escalón adiabático.
\end{itemize}

\noindent\textbf{Desarrollo paso a paso}

\begin{enumerate}

\item Se calcula el aumento relativo de temperatura ideal.

\[\fitmath{\mathrm{rr}=\frac{\mathrm{eta}\cdot \mathrm{wc}}{\mathrm{cp}\cdot \mathrm{Tin}}}\]

\[\fitmath{\mathrm{rr}=\frac{(0.87)\cdot (29800)}{(1004.5)\cdot (518.655)}\approx 0.04976305\,}\]

\item Se recupera la relación de presión.

\[\fitmath{\mathrm{pr}={\left(1+\frac{\mathrm{eta}\cdot \mathrm{wc}}{\mathrm{cp}\cdot \mathrm{Tin}}\right)}^{\frac{\gamma}{\gamma-1}}}\]

\[\fitmath{\mathrm{pr}={\left(1+\frac{(0.87)\cdot (29800)}{(1004.5)\cdot (518.655)}\right)}^{\frac{(1.4)}{(1.4)-1}}\approx 1.185276\,}\]

\end{enumerate}


\noindent\textbf{Sustitución numérica}

\[
\fitmath{\left(1 + \frac{\left(0.87\right) \left(29800\,\mathrm{J}\,\mathrm{kg}^{-1}\right)}{\left(518.655\,\mathrm{K}\right) \left(1004.5\,\mathrm{J}\,\mathrm{kg}^{-1}\,\mathrm{K}^{-1}\right)}\right)^{\frac{\left(1.4\right)}{\left(1.4\right) - 1}}}
\]

\noindent\textbf{Resultado}

\[
\fitmath{\pi_c=1.185275999\times10^{0}}
\]

\noindent\textbf{Interpretación}
A mayor temperatura de entrada y mismo trabajo, la razón de presión es menor.

\noindent\textbf{Comprobación de los pasos}
\begin{itemize}
\item Paso 2: $0.0497631$ — reproducido
\item Paso 4: $1.18528$ — reproducido
\end{itemize}

\noindent\textbf{Comprobación de la respuesta} la aritmética del último paso reproduce la respuesta enunciada a la precisión con que está escrita. Es una reproducción aritmética; no valida el modelo físico ni sus hipótesis.

\noindent\textbf{Limitaciones}
\begin{itemize}
\item Modelo educativo; la aplicabilidad a una máquina real requiere revisión experta.
\end{itemize}

\noindent\textbf{Errores frecuentes}
\begin{itemize}
\item Asignar igual razón de presión a todos los escalones de igual trabajo.
\end{itemize}
% studium:end-record eyJpZCI6ICJQUkItMDA3MCIsICJzZWN0aW9uIjogImMwOSJ9
\subsection*{C09-04-P}

\noindent\textbf{Enunciado}
\(C09-04-P \cdot\) Acoplamiento de eje con caudales distintos. Ejercicio de transferencia. El compresor y la turbina comparten eje con rendimiento mecánico dado. Calcule el trabajo que debe extraer cada kilogramo de gas en la turbina. Antes de consultar la solución, prediga el efecto de cambiar el primer dato respecto del caso resuelto de esta familia. Justifique la tendencia y explique qué restricción podría convertirse en limitante. Datos: \(\dot{m}_c=26.625\, \mathrm{kg}\,\mathrm{s}^{-1}\); \(\dot{m}_t=25.6\, \mathrm{kg}\,\mathrm{s}^{-1}\); \(w_c=214000\, \mathrm{J}\,\mathrm{kg}^{-1}\); \(\dot{W}_{\mathrm{acc}}=85000\, \mathrm{W}\); \(\eta_m=0.975\).

\noindent\textbf{Datos}
\begin{itemize}
\item $\dot{m}_c = 26.625\,\mathrm{kg}\,\mathrm{s}^{-1}$ — caudal compresor
\item $\dot{m}_t = 25.6\,\mathrm{kg}\,\mathrm{s}^{-1}$ — caudal turbina
\item $w_c = 214000\,\mathrm{J}\,\mathrm{kg}^{-1}$ — trabajo comunicado al aire
\item $\dot{W}_{\mathrm{acc}} = 85000\,\mathrm{W}$ — potencia de accesorios
\item $\eta_m = 0.975$ — eficiencia de transmisión desde turbina
\end{itemize}

\noindent\textbf{Incógnita}
$w_t$, en $\mathrm{J\,kg^{-1}}$; también las magnitudes intermedias desarrolladas.

\noindent\textbf{Modelo físico}

\noindent\textit{Acoplamiento de eje con caudales distintos}

\[
\fitmath{w_t=\frac{\dot{W}_{\mathrm{acc}} + \dot{m}_c w_c}{\eta_m \dot{m}_t}}
\]

\noindent\textbf{Hipótesis}
\begin{itemize}
\item wc ya excluye pérdidas de transmisión; em se aplica una vez al lado turbina.
\end{itemize}

\noindent\textbf{Desarrollo paso a paso}

\begin{enumerate}

\item Se suma la demanda de potencia de aire y accesorios.

\[\fitmath{\mathrm{need}=\mathrm{mc}\cdot \mathrm{wc}+\mathrm{acc}}\]

\[\fitmath{\mathrm{need}=(26.625)\cdot (214000)+(85000)\approx 5782750\,\mathrm{W}}\]

\item El caudal de turbina y la transmisión fijan su trabajo.

\[\fitmath{\mathrm{wt}=\frac{\mathrm{mc}\cdot \mathrm{wc}+\mathrm{acc}}{\mathrm{em}\cdot \mathrm{mt}}}\]

\[\fitmath{\mathrm{wt}=\frac{(26.625)\cdot (214000)+(85000)}{(0.975)\cdot (25.6)}\approx 231680.7\,\mathrm{J/kg}}\]

\end{enumerate}


\noindent\textbf{Sustitución numérica}

\[
\fitmath{\frac{\left(85000\,\mathrm{W}\right) + \left(26.625\,\mathrm{kg}\,\mathrm{s}^{-1}\right) \left(214000\,\mathrm{J}\,\mathrm{kg}^{-1}\right)}{\left(0.975\right) \left(25.6\,\mathrm{kg}\,\mathrm{s}^{-1}\right)}}
\]

\noindent\textbf{Resultado}

\[
\fitmath{w_t=2.316806891\times10^{5}\,\mathrm{J}\,\mathrm{kg}^{-1}}
\]

\noindent\textbf{Interpretación}
Usar el caudal del compresor en la turbina omite la diferencia de masas. El reparto de accesorios está explicitado en este modelo.

\noindent\textbf{Comprobación de los pasos}
\begin{itemize}
\item Paso 2: $5782750$ — reproducido
\item Paso 4: $231681$ — reproducido
\end{itemize}

\noindent\textbf{Comprobación de la respuesta} la aritmética del último paso reproduce la respuesta enunciada a la precisión con que está escrita. Es una reproducción aritmética; no valida el modelo físico ni sus hipótesis.

\noindent\textbf{Limitaciones}
\begin{itemize}
\item Modelo educativo; la aplicabilidad a una máquina real requiere revisión experta.
\end{itemize}

\noindent\textbf{Errores frecuentes}
\begin{itemize}
\item Suponer igualdad de caudales sin revisar combustible y sangrados.
\end{itemize}
% studium:end-record eyJpZCI6ICJQUkItMDA3MiIsICJzZWN0aW9uIjogImMwOSJ9
\section{Capítulo 10: Arranque, pérdida rotativa, bombeo y mapas}
\subsection*{C10-01-P}

\noindent\textbf{Enunciado}
\(C10-01-P \cdot\) Incidencia durante variación de caudal. Ejercicio de transferencia. Un rotor sin prerrotación conserva velocidad periférica y ángulo metálico. Calcule su beta de entrada y la incidencia a la nueva velocidad axial. Antes de consultar la solución, prediga el efecto de cambiar el primer dato respecto del caso resuelto de esta familia. Justifique la tendencia y explique qué restricción podría convertirse en limitante. Datos: \(U=263.055\, \mathrm{m}\,\mathrm{s}^{-1}\); \(C_x=119\, \mathrm{m}\,\mathrm{s}^{-1}\); \(\kappa_e=-1.064651\, \mathrm{rad}\).

\noindent\textbf{Datos}
\begin{itemize}
\item $U = 263.055\,\mathrm{m}\,\mathrm{s}^{-1}$ — velocidad periférica
\item $C_x = 119\,\mathrm{m}\,\mathrm{s}^{-1}$ — velocidad axial nueva
\item $\kappa_e = -1.0646508437165\,\mathrm{rad}$ — ángulo metálico de entrada
\end{itemize}

\noindent\textbf{Incógnita}
$i$, en $\mathrm{rad}$; también las magnitudes intermedias desarrolladas.

\noindent\textbf{Modelo físico}

\noindent\textit{Incidencia durante variación de caudal}

\[
\fitmath{i=- \kappa_e - \operatorname{atan}{\left(\frac{U}{C_x} \right)}}
\]

\noindent\textbf{Hipótesis}
\begin{itemize}
\item Velocidad axial positiva, sin prerrotación, metal fijo.
\end{itemize}

\noindent\textbf{Desarrollo paso a paso}

\begin{enumerate}

\item El triángulo relativo utiliza tangencial -U.

\[\fitmath{\beta=\arctan\left(\frac{-\mathrm{uu}}{\mathrm{cx}}\right)}\]

\[\fitmath{\beta=\arctan\left(\frac{-(263.055)}{(119)}\right)\approx -1.145968\,\mathrm{rad}}\]

\item Se compara con la dirección del metal.

\[\fitmath{\mathrm{inc}=\arctan\left(\frac{-\mathrm{uu}}{\mathrm{cx}}\right)-\mathrm{kap}}\]

\[\fitmath{\mathrm{inc}=\arctan\left(\frac{-(263.055)}{(119)}\right)-(-1.064650844)\approx -0.0813167\,\mathrm{rad}}\]

\end{enumerate}


\noindent\textbf{Sustitución numérica}

\[
\fitmath{- \left(-1.0646508437165\,\mathrm{rad}\right) - \operatorname{atan}{\left(\frac{\left(263.055\,\mathrm{m}\,\mathrm{s}^{-1}\right)}{\left(119\,\mathrm{m}\,\mathrm{s}^{-1}\right)} \right)}}
\]

\noindent\textbf{Resultado}

\[
\fitmath{i=-8.131670125\times10^{-2}\,\mathrm{rad}}
\]

\noindent\textbf{Interpretación}
El signo corresponde exclusivamente al convenio algebraico de este manual. Debe compararse con una incidencia nominal calculada igual.

\noindent\textbf{Comprobación de los pasos}
\begin{itemize}
\item Paso 2: $-1.14597$ — reproducido
\item Paso 4: $-0.0813167$ — reproducido
\end{itemize}

\noindent\textbf{Comprobación de la respuesta} la aritmética del último paso reproduce la respuesta enunciada a la precisión con que está escrita. Es una reproducción aritmética; no valida el modelo físico ni sus hipótesis.

\noindent\textbf{Limitaciones}
\begin{itemize}
\item Modelo educativo; la aplicabilidad a una máquina real requiere revisión experta.
\end{itemize}

\noindent\textbf{Errores frecuentes}
\begin{itemize}
\item Conservar beta al cambiar Cx/U.
\end{itemize}
% studium:end-record eyJpZCI6ICJQUkItMDA3NCIsICJzZWN0aW9uIjogImMxMCJ9
\subsection*{C10-02-P}

\noindent\textbf{Enunciado}
\(C10-02-P \cdot\) Margen de caudal definido. Ejercicio de transferencia. A igual velocidad corregida, un punto operativo y un límite prescrito tienen los caudales indicados. Calcule el margen según la definición del capítulo. Antes de consultar la solución, prediga el efecto de cambiar el primer dato respecto del caso resuelto de esta familia. Justifique la tendencia y explique qué restricción podría convertirse en limitante. Datos: \(\dot{m}_{c,\mathrm{op}}=30.246\, \mathrm{kg}\,\mathrm{s}^{-1}\); \(\dot{m}_{c,\mathrm{lim}}=23.9\, \mathrm{kg}\,\mathrm{s}^{-1}\).

\noindent\textbf{Datos}
\begin{itemize}
\item $\dot{m}_{c,\mathrm{op}} = 30.246\,\mathrm{kg}\,\mathrm{s}^{-1}$ — caudal corregido operativo
\item $\dot{m}_{c,\mathrm{lim}} = 23.9\,\mathrm{kg}\,\mathrm{s}^{-1}$ — caudal corregido límite
\end{itemize}

\noindent\textbf{Incógnita}
$\mathrm{SM}_m$ (adimensional); también las magnitudes intermedias desarrolladas.

\noindent\textbf{Modelo físico}

\noindent\textit{Margen de caudal definido}

\[
\fitmath{\mathrm{SM}_m=\frac{- 100 \dot{m}_{c,\mathrm{lim}} + 100 \dot{m}_{c,\mathrm{op}}}{\dot{m}_{c,\mathrm{op}}}}
\]

\noindent\textbf{Hipótesis}
\begin{itemize}
\item Caudales corregidos a las mismas referencias.
\end{itemize}

\noindent\textbf{Desarrollo paso a paso}

\begin{enumerate}

\item Se calcula la separación de caudal.

\[\fitmath{\mathrm{dm}=\mathrm{op}-\mathrm{lim}}\]

\[\fitmath{\mathrm{dm}=(30.246)-(23.9)\approx 6.346\,\mathrm{kg/s}}\]

\item Se normaliza por el caudal operativo y se expresa en porcentaje.

\[\fitmath{\mathrm{SM}=\frac{100\cdot \left(\mathrm{op}-\mathrm{lim}\right)}{\mathrm{op}}}\]

\[\fitmath{\mathrm{SM}=\frac{100\cdot \left((30.246)-(23.9)\right)}{(30.246)}\approx 20.98129\,}\]

\end{enumerate}


\noindent\textbf{Sustitución numérica}

\[
\fitmath{\frac{- 100 \left(23.9\,\mathrm{kg}\,\mathrm{s}^{-1}\right) + 100 \left(30.246\,\mathrm{kg}\,\mathrm{s}^{-1}\right)}{\left(30.246\,\mathrm{kg}\,\mathrm{s}^{-1}\right)}}
\]

\noindent\textbf{Resultado}

\[
\fitmath{\mathrm{SM}_m=2.098128678\times10^{1}}
\]

\noindent\textbf{Interpretación}
El porcentaje solo es comparable con la misma definición y velocidad corregida. El límite es un dato, no una predicción del modelo.

\noindent\textbf{Comprobación de los pasos}
\begin{itemize}
\item Paso 2: $6.346$ — reproducido
\item Paso 4: $20.9813$ — reproducido
\end{itemize}

\noindent\textbf{Comprobación de la respuesta} la aritmética del último paso reproduce la respuesta enunciada a la precisión con que está escrita. Es una reproducción aritmética; no valida el modelo físico ni sus hipótesis.

\noindent\textbf{Limitaciones}
\begin{itemize}
\item Modelo educativo; la aplicabilidad a una máquina real requiere revisión experta.
\end{itemize}

\noindent\textbf{Errores frecuentes}
\begin{itemize}
\item Normalizar con el límite y llamar igual al margen resultante.
\end{itemize}
% studium:end-record eyJpZCI6ICJQUkItMDA3NiIsICJzZWN0aW9uIjogImMxMCJ9
\subsection*{C10-03-P}

\noindent\textbf{Enunciado}
\(C10-03-P \cdot\) Corrección de velocidad. Ejercicio de transferencia. Halle la velocidad corregida y la velocidad periférica física al radio dado. La incógnita final es U real. Antes de consultar la solución, prediga el efecto de cambiar el primer dato respecto del caso resuelto de esta familia. Justifique la tendencia y explique qué restricción podría convertirse en limitante. Datos: \(N=209.45\, \mathrm{s}^{-1}\); \(T_{0e}=259\, \mathrm{K}\); \(T_{\mathrm{ref}}=288.15\, \mathrm{K}\); \(r=0.276\, \mathrm{m}\).

\noindent\textbf{Datos}
\begin{itemize}
\item $N = 209.45\,\mathrm{s}^{-1}$ — giro físico en revoluciones por segundo; equivale a 12567 rpm
\item $T_{0e} = 259\,\mathrm{K}$ — temperatura total entrada
\item $T_{\mathrm{ref}} = 288.15\,\mathrm{K}$ — temperatura referencia
\item $r = 0.276\,\mathrm{m}$ — radio
\end{itemize}

\noindent\textbf{Incógnita}
$U$, en $\mathrm{m\,s^{-1}}$; también las magnitudes intermedias desarrolladas.

\noindent\textbf{Modelo físico}

\noindent\textit{Corrección de velocidad}

\[
\fitmath{U=2 \pi N r}
\]

\noindent\textbf{Hipótesis}
\begin{itemize}
\item Giro en revoluciones por segundo y radio fijo.
\end{itemize}

\noindent\textbf{Desarrollo paso a paso}

\begin{enumerate}

\item La corrección compara velocidades a la referencia térmica.

\[\fitmath{\mathrm{Nc}=\frac{\mathrm{nn}}{\sqrt{\frac{\mathrm{Tin}}{\mathrm{Tr}}}}}\]

\[\fitmath{\mathrm{Nc}=\frac{(209.45)}{\sqrt{\frac{(259)}{(288.15)}}}\approx 220.9224\,\mathrm{s^{-1}}}\]

\item El giro en revoluciones por segundo se convierte a velocidad angular para U.

\[\fitmath{U=2\cdot \pi\cdot \mathrm{nn}\cdot \mathrm{rad}}\]

\[\fitmath{U=2\cdot \pi\cdot (209.45)\cdot (0.276)\approx 363.2196\,\mathrm{m/s}}\]

\end{enumerate}


\noindent\textbf{Sustitución numérica}

\[
\fitmath{2 \pi \left(209.45\,\mathrm{s}^{-1}\right) \left(0.276\,\mathrm{m}\right)}
\]

\noindent\textbf{Resultado}

\[
\fitmath{U=3.632196329\times10^{2}\,\mathrm{m}\,\mathrm{s}^{-1}}
\]

\noindent\textbf{Interpretación}
El dato se expresa en revoluciones por segundo; U se calcula con N físico, no con N corregido.

\noindent\textbf{Comprobación de los pasos}
\begin{itemize}
\item Paso 2: $220.922$ — reproducido
\item Paso 4: $363.22$ — reproducido
\end{itemize}

\noindent\textbf{Comprobación de la respuesta} la aritmética del último paso reproduce la respuesta enunciada a la precisión con que está escrita. Es una reproducción aritmética; no valida el modelo físico ni sus hipótesis.

\noindent\textbf{Limitaciones}
\begin{itemize}
\item Modelo educativo; la aplicabilidad a una máquina real requiere revisión experta.
\end{itemize}

\noindent\textbf{Errores frecuentes}
\begin{itemize}
\item Usar N corregido en una velocidad mecánica real.
\end{itemize}
% studium:end-record eyJpZCI6ICJQUkItMDIwOCIsICJzZWN0aW9uIjogImMxMCJ9
\subsection*{C10-04-P}

\noindent\textbf{Enunciado}
\(C10-04-P \cdot\) Estabilidad lineal del modelo ampliado. Ejercicio de transferencia. Para el modelo de dos estados del capítulo, calcule traza y determinante de la matriz lineal. El resultado final es el determinante; interprete estabilidad. Antes de consultar la solución, prediga el efecto de cambiar el primer dato respecto del caso resuelto de esta familia. Justifique la tendencia y explique qué restricción podría convertirse en limitante. Datos: \(a_c=-1917\, \mathrm{Pa}\,\mathrm{s}\,\mathrm{kg}^{-1}\); \(b_v=0.00016\, \mathrm{kg}\,\mathrm{s}^{-1}\,\mathrm{Pa}^{-1}\); \(L_c=260\, \mathrm{Pa}\,\mathrm{s}^{2}\,\mathrm{kg}^{-1}\); \(K_p=0.0018\, \mathrm{kg}\,\mathrm{Pa}^{-1}\).

\noindent\textbf{Datos}
\begin{itemize}
\item $a_c = -1917\,\mathrm{Pa}\,\mathrm{s}\,\mathrm{kg}^{-1}$ — pendiente de compresor
\item $b_v = 0.00016\,\mathrm{kg}\,\mathrm{s}^{-1}\,\mathrm{Pa}^{-1}$ — pendiente de descarga
\item $L_c = 260\,\mathrm{Pa}\,\mathrm{s}^{2}\,\mathrm{kg}^{-1}$ — inercia del modelo
\item $K_p = 0.0018\,\mathrm{kg}\,\mathrm{Pa}^{-1}$ — capacidad del modelo
\end{itemize}

\noindent\textbf{Incógnita}
$\det(J)$, en $\mathrm{s^{-2}}$; también las magnitudes intermedias desarrolladas.

\noindent\textbf{Modelo físico}

\noindent\textit{Estabilidad lineal del modelo ampliado}

\[
\fitmath{\det(J)=\frac{- a_c b_v + 1}{K_p L_c}}
\]

\noindent\textbf{Hipótesis}
\begin{itemize}
\item Linealización local; parámetros positivos L,K; modelo pedagógico sin retardo.
\end{itemize}

\noindent\textbf{Desarrollo paso a paso}

\begin{enumerate}

\item La traza suma los términos diagonales.

\[\fitmath{\mathrm{tr}=\frac{\mathrm{aa}}{\mathrm{ll}}-\frac{\mathrm{bb}}{\mathrm{kk}}}\]

\[\fitmath{\mathrm{tr}=\frac{(-1917)}{(260)}-\frac{(0.00016)}{(0.0018)}\approx -7.461966\,\mathrm{s^{-1}}}\]

\item El determinante incluye el acoplamiento cruzado.

\[\fitmath{\mathrm{det}=\frac{1-\mathrm{aa}\cdot \mathrm{bb}}{\mathrm{ll}\cdot \mathrm{kk}}}\]

\[\fitmath{\mathrm{det}=\frac{1-(-1917)\cdot (0.00016)}{(260)\cdot (0.0018)}\approx 2.792137\,\mathrm{s^{-2}}}\]

\end{enumerate}


\noindent\textbf{Sustitución numérica}

\[
\fitmath{\frac{- \left(-1917\,\mathrm{Pa}\,\mathrm{s}\,\mathrm{kg}^{-1}\right) \left(0.00016\,\mathrm{kg}\,\mathrm{s}^{-1}\,\mathrm{Pa}^{-1}\right) + 1}{\left(0.0018\,\mathrm{kg}\,\mathrm{Pa}^{-1}\right) \left(260\,\mathrm{Pa}\,\mathrm{s}^{2}\,\mathrm{kg}^{-1}\right)}}
\]

\noindent\textbf{Resultado}

\[
\fitmath{\det(J)=2.792136752\times10^{0}\,\mathrm{s}^{-2}}
\]

\noindent\textbf{Interpretación}
Traza negativa y determinante positivo dan estabilidad lineal del modelo de dos estados. Esto no predice pérdida rotativa ni estabilidad de un motor real.

\noindent\textbf{Comprobación de los pasos}
\begin{itemize}
\item Paso 2: $-7.46197$ — reproducido
\item Paso 4: $2.79214$ — reproducido
\end{itemize}

\noindent\textbf{Comprobación de la respuesta} la aritmética del último paso reproduce la respuesta enunciada a la precisión con que está escrita. Es una reproducción aritmética; no valida el modelo físico ni sus hipótesis.

\noindent\textbf{Limitaciones}
\begin{itemize}
\item Modelo educativo; la aplicabilidad a una máquina real requiere revisión experta.
\end{itemize}

\noindent\textbf{Errores frecuentes}
\begin{itemize}
\item Tomar L como longitud sin unidades equivalentes.
\end{itemize}
% studium:end-record eyJpZCI6ICJQUkItMDA4MCIsICJzZWN0aW9uIjogImMxMCJ9
\section{Capítulo 11: Equilibrio radial y leyes de remolino}
\subsection*{C11-01-P}

\noindent\textbf{Enunciado}
\(C11-01-P \cdot\) Presión de vórtice libre. Ejercicio de transferencia. Integre el equilibrio radial de un fluido de densidad constante y remolino K/r entre radio medio y punta. Antes de consultar la solución, prediga el efecto de cambiar el primer dato respecto del caso resuelto de esta familia. Justifique la tendencia y explique qué restricción podría convertirse en limitante. Datos: \(\rho=2.4495\, \mathrm{kg}\,\mathrm{m}^{-3}\); \(K_v=36\, \mathrm{m}^{2}\,\mathrm{s}^{-1}\); \(r_m=0.29\, \mathrm{m}\); \(r_t=0.36\, \mathrm{m}\).

\noindent\textbf{Datos}
\begin{itemize}
\item $\rho = 2.4495\,\mathrm{kg}\,\mathrm{m}^{-3}$ — densidad uniforme
\item $K_v = 36\,\mathrm{m}^{2}\,\mathrm{s}^{-1}$ — constante de remolino
\item $r_m = 0.29\,\mathrm{m}$ — radio referencia
\item $r_t = 0.36\,\mathrm{m}$ — radio exterior
\end{itemize}

\noindent\textbf{Incógnita}
$\Delta p$, en $\mathrm{Pa}$; también las magnitudes intermedias desarrolladas.

\noindent\textbf{Modelo físico}

\noindent\textit{Presión de vórtice libre}

\[
\fitmath{\Delta p=\frac{K_v^{2} \rho \left(- \frac{1}{r_t^{2}} + \frac{1}{r_m^{2}}\right)}{2}}
\]

\noindent\textbf{Hipótesis}
\begin{itemize}
\item Equilibrio radial simple, axisimetría, densidad constante.
\end{itemize}

\noindent\textbf{Desarrollo paso a paso}

\begin{enumerate}

\item El momento angular constante fija Ctheta en la punta.

\[\fitmath{\mathrm{ct}=\frac{\mathrm{kv}}{\mathrm{rt}}}\]

\[\fitmath{\mathrm{ct}=\frac{(36)}{(0.36)}\approx 100\,\mathrm{m/s}}\]

\item Se integra el gradiente positivo de presión.

\[\fitmath{\mathrm{dp}=\frac{\rho\cdot {\mathrm{kv}}^{2}\cdot \left(\frac{1}{{\mathrm{rm}}^{2}}-\frac{1}{{\mathrm{rt}}^{2}}\right)}{2}}\]

\[\fitmath{\mathrm{dp}=\frac{(2.4495)\cdot {(36)}^{2}\cdot \left(\frac{1}{{(0.29)}^{2}}-\frac{1}{{(0.36)}^{2}}\right)}{2}\approx 6626.174\,\mathrm{Pa}}\]

\end{enumerate}


\noindent\textbf{Sustitución numérica}

\[
\fitmath{\frac{\left(\left(36\,\mathrm{m}^{2}\,\mathrm{s}^{-1}\right)\right)^{2} \left(2.4495\,\mathrm{kg}\,\mathrm{m}^{-3}\right) \left(- \frac{1}{\left(0.36\,\mathrm{m}\right)^{2}} + \frac{1}{\left(0.29\,\mathrm{m}\right)^{2}}\right)}{2}}
\]

\noindent\textbf{Resultado}

\[
\fitmath{\Delta p=6.626174197\times10^{3}\,\mathrm{Pa}}
\]

\noindent\textbf{Interpretación}
La presión exterior aumenta. La aproximación de densidad constante debe justificarse antes de aplicarla a un gas real.

\noindent\textbf{Comprobación de los pasos}
\begin{itemize}
\item Paso 2: $100$ — reproducido
\item Paso 4: $6626.17$ — reproducido
\end{itemize}

\noindent\textbf{Comprobación de la respuesta} la aritmética del último paso reproduce la respuesta enunciada a la precisión con que está escrita. Es una reproducción aritmética; no valida el modelo físico ni sus hipótesis.

\noindent\textbf{Limitaciones}
\begin{itemize}
\item Modelo educativo; la aplicabilidad a una máquina real requiere revisión experta.
\end{itemize}

\noindent\textbf{Errores frecuentes}
\begin{itemize}
\item Invertir los radios y obtener caída de presión hacia fuera.
\end{itemize}
% studium:end-record eyJpZCI6ICJQUkItMDA4MiIsICJzZWN0aW9uIjogImMxMSJ9
\subsection*{C11-02-P}

\noindent\textbf{Enunciado}
\(C11-02-P \cdot\) Trabajo uniforme con vórtice libre. Ejercicio de transferencia. Una fila móvil tiene constantes de vórtice libre distintas en entrada y salida. Calcule el trabajo y el incremento de temperatura total. Antes de consultar la solución, prediga el efecto de cambiar el primer dato respecto del caso resuelto de esta familia. Justifique la tendencia y explique qué restricción podría convertirse en limitante. Datos: \(\Omega=1150.2\, \mathrm{s}^{-1}\); \(K_e=11\, \mathrm{m}^{2}\,\mathrm{s}^{-1}\); \(K_s=39\, \mathrm{m}^{2}\,\mathrm{s}^{-1}\); \(c_p=1004.5\, \mathrm{J}\,\mathrm{kg}^{-1}\,\mathrm{K}^{-1}\).

\noindent\textbf{Datos}
\begin{itemize}
\item $\Omega = 1150.2\,\mathrm{s}^{-1}$ — velocidad angular
\item $K_e = 11\,\mathrm{m}^{2}\,\mathrm{s}^{-1}$ — constante entrada
\item $K_s = 39\,\mathrm{m}^{2}\,\mathrm{s}^{-1}$ — constante salida
\item $c_p = 1004.5\,\mathrm{J}\,\mathrm{kg}^{-1}\,\mathrm{K}^{-1}$ — calor específico
\end{itemize}

\noindent\textbf{Incógnita}
$\Delta T_0$, en $\mathrm{K}$; también las magnitudes intermedias desarrolladas.

\noindent\textbf{Modelo físico}

\noindent\textit{Trabajo uniforme con vórtice libre}

\[
\fitmath{\Delta T_0=\frac{\Omega \left(- K_e + K_s\right)}{c_p}}
\]

\noindent\textbf{Hipótesis}
\begin{itemize}
\item Misma velocidad angular y líneas de corriente del modelo ideal.
\end{itemize}

\noindent\textbf{Desarrollo paso a paso}

\begin{enumerate}

\item Euler cancela la dependencia radial r.

\[\fitmath{\mathrm{wc}=\Omega\cdot \left(\mathrm{ks}-\mathrm{ke}\right)}\]

\[\fitmath{\mathrm{wc}=(1150.2)\cdot \left((39)-(11)\right)\approx 32205.6\,\mathrm{J/kg}}\]

\item Se convierte el trabajo en incremento total.

\[\fitmath{\mathrm{dT}=\frac{\Omega\cdot \left(\mathrm{ks}-\mathrm{ke}\right)}{\mathrm{cp}}}\]

\[\fitmath{\mathrm{dT}=\frac{(1150.2)\cdot \left((39)-(11)\right)}{(1004.5)}\approx 32.06132\,\mathrm{K}}\]

\end{enumerate}


\noindent\textbf{Sustitución numérica}

\[
\fitmath{\frac{\left(1150.2\,\mathrm{s}^{-1}\right) \left(- \left(11\,\mathrm{m}^{2}\,\mathrm{s}^{-1}\right) + \left(39\,\mathrm{m}^{2}\,\mathrm{s}^{-1}\right)\right)}{\left(1004.5\,\mathrm{J}\,\mathrm{kg}^{-1}\,\mathrm{K}^{-1}\right)}}
\]

\noindent\textbf{Resultado}

\[
\fitmath{\Delta T_0=3.206132404\times10^{1}\,\mathrm{K}}
\]

\noindent\textbf{Interpretación}
El trabajo es uniforme con el radio, aunque la reacción no lo sea.

\noindent\textbf{Comprobación de los pasos}
\begin{itemize}
\item Paso 2: $32205.6$ — reproducido
\item Paso 4: $32.0613$ — reproducido
\end{itemize}

\noindent\textbf{Comprobación de la respuesta} la aritmética del último paso reproduce la respuesta enunciada a la precisión con que está escrita. Es una reproducción aritmética; no valida el modelo físico ni sus hipótesis.

\noindent\textbf{Limitaciones}
\begin{itemize}
\item Modelo educativo; la aplicabilidad a una máquina real requiere revisión experta.
\end{itemize}

\noindent\textbf{Errores frecuentes}
\begin{itemize}
\item Usar igual constante de vórtice a ambos lados de un rotor con trabajo.
\end{itemize}
% studium:end-record eyJpZCI6ICJQUkItMDA4NCIsICJzZWN0aW9uIjogImMxMSJ9
\subsection*{C11-03-P}

\noindent\textbf{Enunciado}
\(C11-03-P \cdot\) Reacción hacia la raíz. Ejercicio de transferencia. Calcule la reacción reducida en la raíz para un escalón repetitivo de vórtice libre. Antes de consultar la solución, prediga el efecto de cambiar el primer dato respecto del caso resuelto de esta familia. Justifique la tendencia y explique qué restricción podría convertirse en limitante. Datos: \(K_e=17.04\, \mathrm{m}^{2}\,\mathrm{s}^{-1}\); \(K_s=49\, \mathrm{m}^{2}\,\mathrm{s}^{-1}\); \(\Omega=1120\, \mathrm{s}^{-1}\); \(r_h=0.205\, \mathrm{m}\).

\noindent\textbf{Datos}
\begin{itemize}
\item $K_e = 17.04\,\mathrm{m}^{2}\,\mathrm{s}^{-1}$ — constante tangencial entrada
\item $K_s = 49\,\mathrm{m}^{2}\,\mathrm{s}^{-1}$ — constante tangencial salida
\item $\Omega = 1120\,\mathrm{s}^{-1}$ — velocidad angular
\item $r_h = 0.205\,\mathrm{m}$ — radio raíz
\end{itemize}

\noindent\textbf{Incógnita}
$\mathcal{R}$ (adimensional); también las magnitudes intermedias desarrolladas.

\noindent\textbf{Modelo físico}

\noindent\textit{Reacción hacia la raíz}

\[
\fitmath{\mathcal R=1 - \frac{K_e + K_s}{2 \Omega r_h^{2}}}
\]

\noindent\textbf{Hipótesis}
\begin{itemize}
\item Criterio reducido válido para radio y Cx constantes en cada escalón repetitivo.
\end{itemize}

\noindent\textbf{Desarrollo paso a paso}

\begin{enumerate}

\item Se calcula la suma de componentes tangenciales en la raíz.

\[\fitmath{\mathrm{sum}=\frac{\mathrm{ke}+\mathrm{ks}}{\mathrm{rh}}}\]

\[\fitmath{\mathrm{sum}=\frac{(17.04)+(49)}{(0.205)}\approx 322.1463\,\mathrm{m/s}}\]

\item Se normaliza con dos veces U local.

\[\fitmath{\mathrm{rea}=1-\frac{\mathrm{ke}+\mathrm{ks}}{2\cdot \Omega\cdot {\mathrm{rh}}^{2}}}\]

\[\fitmath{\mathrm{rea}=1-\frac{(17.04)+(49)}{2\cdot (1120)\cdot {(0.205)}^{2}}\approx 0.2984618\,}\]

\end{enumerate}


\noindent\textbf{Sustitución numérica}

\[
\fitmath{1 - \frac{\left(17.04\,\mathrm{m}^{2}\,\mathrm{s}^{-1}\right) + \left(49\,\mathrm{m}^{2}\,\mathrm{s}^{-1}\right)}{2 \left(1120\,\mathrm{s}^{-1}\right) \left(0.205\,\mathrm{m}\right)^{2}}}
\]

\noindent\textbf{Resultado}

\[
\fitmath{\mathcal R=2.984617999\times10^{-1}}
\]

\noindent\textbf{Interpretación}
La raíz puede tener reacción mucho menor que el radio medio. Si resulta negativa se debe revisar el diseño, no cambiar el signo.

\noindent\textbf{Comprobación de los pasos}
\begin{itemize}
\item Paso 2: $322.146$ — reproducido
\item Paso 4: $0.298462$ — reproducido
\end{itemize}

\noindent\textbf{Comprobación de la respuesta} la aritmética del último paso reproduce la respuesta enunciada a la precisión con que está escrita. Es una reproducción aritmética; no valida el modelo físico ni sus hipótesis.

\noindent\textbf{Limitaciones}
\begin{itemize}
\item Modelo educativo; la aplicabilidad a una máquina real requiere revisión experta.
\end{itemize}

\noindent\textbf{Errores frecuentes}
\begin{itemize}
\item Asignar a toda la pala la reacción del radio medio.
\end{itemize}
% studium:end-record eyJpZCI6ICJQUkItMDA4NiIsICJzZWN0aW9uIjogImMxMSJ9
\subsection*{C11-04-P}

\noindent\textbf{Enunciado}
\(C11-04-P \cdot\) Vórtice forzado admisible. Ejercicio de transferencia. Con entalpía total y entropía uniformes, calcule la velocidad axial exterior exigida por un remolino forzado. Antes de consultar la solución, prediga el efecto de cambiar el primer dato respecto del caso resuelto de esta familia. Justifique la tendencia y explique qué restricción podría convertirse en limitante. Datos: \(C_{x,m}=200.22\, \mathrm{m}\,\mathrm{s}^{-1}\); \(k_f=320\, \mathrm{s}^{-1}\); \(r_m=0.25\, \mathrm{m}\); \(r_t=0.33\, \mathrm{m}\).

\noindent\textbf{Datos}
\begin{itemize}
\item $C_{x,m} = 200.22\,\mathrm{m}\,\mathrm{s}^{-1}$ — velocidad axial al radio medio
\item $k_f = 320\,\mathrm{s}^{-1}$ — coeficiente $C_\theta=k_f r$
\item $r_m = 0.25\,\mathrm{m}$ — radio medio
\item $r_t = 0.33\,\mathrm{m}$ — radio exterior
\end{itemize}

\noindent\textbf{Incógnita}
$C_x$, en $\mathrm{m\,s^{-1}}$; también las magnitudes intermedias desarrolladas.

\noindent\textbf{Modelo físico}

\noindent\textit{Vórtice forzado admisible}

\[
\fitmath{C_x=\sqrt{C_{x,m}^{2} - 2 k_f^{2} \left(- r_m^{2} + r_t^{2}\right)}}
\]

\noindent\textbf{Hipótesis}
\begin{itemize}
\item Equilibrio radial simple, h0 y s uniformes.
\end{itemize}

\noindent\textbf{Desarrollo paso a paso}

\begin{enumerate}

\item El equilibrio energético fija el cuadrado axial exterior.

\[\fitmath{\mathrm{cx}_{2}={\mathrm{cxm}}^{2}-2\cdot {\mathrm{kv}}^{2}\cdot \left({\mathrm{rt}}^{2}-{\mathrm{rm}}^{2}\right)}\]

\[\fitmath{\mathrm{cx}_{2}={(200.22)}^{2}-2\cdot {(320)}^{2}\cdot \left({(0.33)}^{2}-{(0.25)}^{2}\right)\approx 30585.33\,\mathrm{m^{2}/s^{2}}}\]

\item Se selecciona la rama axial positiva.

\[\fitmath{\mathrm{cx}=\sqrt{{\mathrm{cxm}}^{2}-2\cdot {\mathrm{kv}}^{2}\cdot \left({\mathrm{rt}}^{2}-{\mathrm{rm}}^{2}\right)}}\]

\[\fitmath{\mathrm{cx}=\sqrt{{(200.22)}^{2}-2\cdot {(320)}^{2}\cdot \left({(0.33)}^{2}-{(0.25)}^{2}\right)}\approx 174.8866\,\mathrm{m/s}}\]

\end{enumerate}


\noindent\textbf{Sustitución numérica}

\[
\fitmath{\sqrt{\left(\left(200.22\,\mathrm{m}\,\mathrm{s}^{-1}\right)\right)^{2} - 2 \left(\left(320\,\mathrm{s}^{-1}\right)\right)^{2} \left(- \left(0.25\,\mathrm{m}\right)^{2} + \left(0.33\,\mathrm{m}\right)^{2}\right)}}
\]

\noindent\textbf{Resultado}

\[
\fitmath{C_x=1.748866158\times10^{2}\,\mathrm{m}\,\mathrm{s}^{-1}}
\]

\noindent\textbf{Interpretación}
El radicando debe ser positivo. Si no lo es, la ley propuesta no admite esta solución axial en todo el anillo.

\noindent\textbf{Comprobación de los pasos}
\begin{itemize}
\item Paso 2: $30585.3$ — reproducido
\item Paso 4: $174.887$ — reproducido
\end{itemize}

\noindent\textbf{Comprobación de la respuesta} la aritmética del último paso reproduce la respuesta enunciada a la precisión con que está escrita. Es una reproducción aritmética; no valida el modelo físico ni sus hipótesis.

\noindent\textbf{Limitaciones}
\begin{itemize}
\item Modelo educativo; la aplicabilidad a una máquina real requiere revisión experta.
\end{itemize}

\noindent\textbf{Errores frecuentes}
\begin{itemize}
\item Imponer Cx constante también para vórtice forzado.
\end{itemize}
% studium:end-record eyJpZCI6ICJQUkItMDA4OCIsICJzZWN0aW9uIjogImMxMSJ9
\section{Capítulo 12: Torsión, secciones y construcción geométrica de álabes}
\subsection*{C12-01-P}

\noindent\textbf{Enunciado}
\(C12-01-P \cdot\) Torsión de entrada de rotor. Ejercicio de transferencia. Un rotor sin prerrotación recibe velocidad axial uniforme. Calcule beta a raíz y punta y la diferencia de torsión de entrada. Antes de consultar la solución, prediga el efecto de cambiar el primer dato respecto del caso resuelto de esta familia. Justifique la tendencia y explique qué restricción podría convertirse en limitante. Datos: \(\Omega=1128.9\, \mathrm{s}^{-1}\); \(C_x=157\, \mathrm{m}\,\mathrm{s}^{-1}\); \(r_h=0.198\, \mathrm{m}\); \(r_t=0.314\, \mathrm{m}\).

\noindent\textbf{Datos}
\begin{itemize}
\item $\Omega = 1128.9\,\mathrm{s}^{-1}$ — velocidad angular
\item $C_x = 157\,\mathrm{m}\,\mathrm{s}^{-1}$ — velocidad axial
\item $r_h = 0.198\,\mathrm{m}$ — radio raíz
\item $r_t = 0.314\,\mathrm{m}$ — radio punta
\end{itemize}

\noindent\textbf{Incógnita}
$\Delta\beta$, en $\mathrm{rad}$; también las magnitudes intermedias desarrolladas.

\noindent\textbf{Modelo físico}

\noindent\textit{Torsión de entrada de rotor}

\[
\fitmath{\Delta\beta=\operatorname{atan}{\left(\frac{\Omega r_h}{C_x} \right)} - \operatorname{atan}{\left(\frac{\Omega r_t}{C_x} \right)}}
\]

\noindent\textbf{Hipótesis}
\begin{itemize}
\item Sin prerrotación, Cx uniforme, incidencia cero.
\end{itemize}

\noindent\textbf{Desarrollo paso a paso}

\begin{enumerate}

\item Se obtiene beta de raíz con componente tangencial -U.

\[\fitmath{\mathrm{bh}=\arctan\left(\frac{-\Omega\cdot \mathrm{rh}}{\mathrm{cx}}\right)}\]

\[\fitmath{\mathrm{bh}=\arctan\left(\frac{-(1128.9)\cdot (0.198)}{(157)}\right)\approx -0.9584674\,\mathrm{rad}}\]

\item Se obtiene beta de punta.

\[\fitmath{\mathrm{bt}=\arctan\left(\frac{-\Omega\cdot \mathrm{rt}}{\mathrm{cx}}\right)}\]

\[\fitmath{\mathrm{bt}=\arctan\left(\frac{-(1128.9)\cdot (0.314)}{(157)}\right)\approx -1.153855\,\mathrm{rad}}\]

\item La torsión es una diferencia algebraica de orientaciones.

\[\fitmath{\mathrm{twist}=\arctan\left(\frac{-\Omega\cdot \mathrm{rt}}{\mathrm{cx}}\right)-\arctan\left(\frac{-\Omega\cdot \mathrm{rh}}{\mathrm{cx}}\right)}\]

\[\fitmath{\mathrm{twist}=\arctan\left(\frac{-(1128.9)\cdot (0.314)}{(157)}\right)-\arctan\left(\frac{-(1128.9)\cdot (0.198)}{(157)}\right)\approx -0.1953875\,\mathrm{rad}}\]

\end{enumerate}


\noindent\textbf{Sustitución numérica}

\[
\fitmath{\operatorname{atan}{\left(\frac{\left(1128.9\,\mathrm{s}^{-1}\right) \left(0.198\,\mathrm{m}\right)}{\left(157\,\mathrm{m}\,\mathrm{s}^{-1}\right)} \right)} - \operatorname{atan}{\left(\frac{\left(1128.9\,\mathrm{s}^{-1}\right) \left(0.314\,\mathrm{m}\right)}{\left(157\,\mathrm{m}\,\mathrm{s}^{-1}\right)} \right)}}
\]

\noindent\textbf{Resultado}

\[
\fitmath{\Delta\beta=-1.953874696\times10^{-1}\,\mathrm{rad}}
\]

\noindent\textbf{Interpretación}
La orientación de punta es más negativa. La tangente del metal coincide solo bajo la hipótesis de incidencia prescrita igual a cero.

\noindent\textbf{Comprobación de los pasos}
\begin{itemize}
\item Paso 2: $-0.958467$ — reproducido
\item Paso 4: $-1.15385$ — reproducido
\item Paso 6: $-0.195387$ — reproducido
\end{itemize}

\noindent\textbf{Comprobación de la respuesta} la aritmética del último paso reproduce la respuesta enunciada a la precisión con que está escrita. Es una reproducción aritmética; no valida el modelo físico ni sus hipótesis.

\noindent\textbf{Limitaciones}
\begin{itemize}
\item Modelo educativo; la aplicabilidad a una máquina real requiere revisión experta.
\end{itemize}

\noindent\textbf{Errores frecuentes}
\begin{itemize}
\item Atribuir el mismo ángulo a todos los radios.
\end{itemize}
% studium:end-record eyJpZCI6ICJQUkItMDA5MCIsICJzZWN0aW9uIjogImMxMiJ9
\subsection*{C12-02-P}

\noindent\textbf{Enunciado}
\(C12-02-P \cdot\) Cuerda y proyección. Ejercicio de transferencia. Una sección tiene cuerda axial y calado dados. Calcule cuerda real y solidez a paso fijado. Antes de consultar la solución, prediga el efecto de cambiar el primer dato respecto del caso resuelto de esta familia. Justifique la tendencia y explique qué restricción podría convertirse en limitante. Datos: \(c_x=0.04686\, \mathrm{m}\); \(\chi=0.6457718\, \mathrm{rad}\); \(s=0.039\, \mathrm{m}\).

\noindent\textbf{Datos}
\begin{itemize}
\item $c_x = 0.04686\,\mathrm{m}$ — cuerda axial
\item $\chi = 0.6457718232379\,\mathrm{rad}$ — calado
\item $s = 0.039\,\mathrm{m}$ — paso
\end{itemize}

\noindent\textbf{Incógnita}
$\sigma$ (adimensional); también las magnitudes intermedias desarrolladas.

\noindent\textbf{Modelo físico}

\noindent\textit{Cuerda y proyección}

\[
\fitmath{\sigma=\frac{c_x}{s \cos{\left(\chi \right)}}}
\]

\noindent\textbf{Hipótesis}
\begin{itemize}
\item Calado menor de noventa grados en módulo.
\end{itemize}

\noindent\textbf{Desarrollo paso a paso}

\begin{enumerate}

\item La cuerda axial es la proyección de la cuerda real.

\[\fitmath{\mathrm{chord}=\frac{\mathrm{cxg}}{\cos\left(\chi\right)}}\]

\[\fitmath{\mathrm{chord}=\frac{(0.04686)}{\cos\left((0.6457718232)\right)}\approx 0.05867508\,\mathrm{m}}\]

\item La solidez usa cuerda real.

\[\fitmath{\mathrm{sol}=\frac{\mathrm{cxg}}{\cos\left(\chi\right)\cdot \mathrm{ss}}}\]

\[\fitmath{\mathrm{sol}=\frac{(0.04686)}{\cos\left((0.6457718232)\right)\cdot (0.039)}\approx 1.504489\,}\]

\end{enumerate}


\noindent\textbf{Sustitución numérica}

\[
\fitmath{\frac{\left(0.04686\,\mathrm{m}\right)}{\left(0.039\,\mathrm{m}\right) \cos{\left(\left(0.6457718232379\,\mathrm{rad}\right) \right)}}}
\]

\noindent\textbf{Resultado}

\[
\fitmath{\sigma=1.504489152\times10^{0}}
\]

\noindent\textbf{Interpretación}
Cuerda axial y cuerda real no pueden intercambiarse en la solidez sin cambiar su definición.

\noindent\textbf{Comprobación de los pasos}
\begin{itemize}
\item Paso 2: $0.0586751$ — reproducido
\item Paso 4: $1.50449$ — reproducido
\end{itemize}

\noindent\textbf{Comprobación de la respuesta} la aritmética del último paso reproduce la respuesta enunciada a la precisión con que está escrita. Es una reproducción aritmética; no valida el modelo físico ni sus hipótesis.

\noindent\textbf{Limitaciones}
\begin{itemize}
\item Modelo educativo; la aplicabilidad a una máquina real requiere revisión experta.
\end{itemize}

\noindent\textbf{Errores frecuentes}
\begin{itemize}
\item Multiplicar por coseno al despejar la cuerda real.
\end{itemize}
% studium:end-record eyJpZCI6ICJQUkItMDA5MiIsICJzZWN0aW9uIjogImMxMiJ9
\subsection*{C12-03-P}

\noindent\textbf{Enunciado}
\(C12-03-P \cdot\) Línea de Hermite. Ejercicio de transferencia. Para pendientes de entrada y salida dadas y desplazamiento d, calcule la ordenada de línea media al parámetro xi. Antes de consultar la solución, prediga el efecto de cambiar el primer dato respecto del caso resuelto de esta familia. Justifique la tendencia y explique qué restricción podría convertirse en limitante. Datos: \(c_x=0.050055\, \mathrm{m}\); \(m_e=0.86\); \(m_s=-0.24\); \(\delta_y=0.31\); \(\xi=0.4\).

\noindent\textbf{Datos}
\begin{itemize}
\item $c_x = 0.050055\,\mathrm{m}$ — cuerda axial
\item $m_e = 0.86$ — pendiente entrada
\item $m_s = -0.24$ — pendiente salida
\item $\delta_y = 0.31$ — desplazamiento final normalizado
\item $\xi = 0.4$ — parámetro xi
\end{itemize}

\noindent\textbf{Incógnita}
$y$, en $\mathrm{m}$; también las magnitudes intermedias desarrolladas.

\noindent\textbf{Modelo físico}

\noindent\textit{Línea de Hermite}

\[
\fitmath{y=c_x \left(\delta_y \left(- 2 \xi^{3} + 3 \xi^{2}\right) + m_e \left(\xi^{3} - 2 \xi^{2} + \xi\right) + m_s \left(\xi^{3} - \xi^{2}\right)\right)}
\]

\noindent\textbf{Hipótesis}
\begin{itemize}
\item Perfil educativo definido como función de x.
\end{itemize}

\noindent\textbf{Desarrollo paso a paso}

\begin{enumerate}

\item Se evalúa la línea media normalizada.

\[\fitmath{\mathrm{fy}=\left({\mathrm{xx}}^{3}-2\cdot {\mathrm{xx}}^{2}+\mathrm{xx}\right)\cdot \mathrm{me}+\left(-2\cdot {\mathrm{xx}}^{3}+3\cdot {\mathrm{xx}}^{2}\right)\cdot \mathrm{dd}+\left({\mathrm{xx}}^{3}-{\mathrm{xx}}^{2}\right)\cdot \mathrm{ms}}\]

\[\fitmath{\mathrm{fy}=\left({(0.4)}^{3}-2\cdot {(0.4)}^{2}+(0.4)\right)\cdot (0.86)+\left(-2\cdot {(0.4)}^{3}+3\cdot {(0.4)}^{2}\right)\cdot (0.31)+\left({(0.4)}^{3}-{(0.4)}^{2}\right)\cdot (-0.24)\approx 0.256\,}\]

\item Se devuelve la longitud física transversal.

\[\fitmath{\mathrm{yy}=\mathrm{cxg}\cdot \left(\left({\mathrm{xx}}^{3}-2\cdot {\mathrm{xx}}^{2}+\mathrm{xx}\right)\cdot \mathrm{me}+\left(-2\cdot {\mathrm{xx}}^{3}+3\cdot {\mathrm{xx}}^{2}\right)\cdot \mathrm{dd}+\left({\mathrm{xx}}^{3}-{\mathrm{xx}}^{2}\right)\cdot \mathrm{ms}\right)}\]

\[\fitmath{\mathrm{yy}=(0.050055)\cdot \left(\left({(0.4)}^{3}-2\cdot {(0.4)}^{2}+(0.4)\right)\cdot (0.86)+\left(-2\cdot {(0.4)}^{3}+3\cdot {(0.4)}^{2}\right)\cdot (0.31)+\left({(0.4)}^{3}-{(0.4)}^{2}\right)\cdot (-0.24)\right)\approx 0.01281408\,\mathrm{m}}\]

\end{enumerate}


\noindent\textbf{Sustitución numérica}

\[
\fitmath{\left(0.050055\,\mathrm{m}\right) \left(\left(0.31\right) \left(- 2 \left(0.4\right)^{3} + 3 \left(0.4\right)^{2}\right) + \left(0.86\right) \left(\left(0.4\right)^{3} - 2 \left(0.4\right)^{2} + \left(0.4\right)\right) + \left(-0.24\right) \left(\left(0.4\right)^{3} - \left(0.4\right)^{2}\right)\right)}
\]

\noindent\textbf{Resultado}

\[
\fitmath{y=1.281408000\times10^{-2}\,\mathrm{m}}
\]

\noindent\textbf{Interpretación}
El resultado es una coordenada geométrica, no una predicción aerodinámica. Compruebe los cuatro extremos y pendientes del polinomio.

\noindent\textbf{Comprobación de los pasos}
\begin{itemize}
\item Paso 2: $0.256$ — reproducido
\item Paso 4: $0.0128141$ — reproducido
\end{itemize}

\noindent\textbf{Comprobación de la respuesta} la aritmética del último paso reproduce la respuesta enunciada a la precisión con que está escrita. Es una reproducción aritmética; no valida el modelo físico ni sus hipótesis.

\noindent\textbf{Limitaciones}
\begin{itemize}
\item Modelo educativo; la aplicabilidad a una máquina real requiere revisión experta.
\end{itemize}

\noindent\textbf{Errores frecuentes}
\begin{itemize}
\item Confundir pendiente con ángulo en grados.
\end{itemize}
% studium:end-record eyJpZCI6ICJQUkItMDA5NCIsICJzZWN0aW9uIjogImMxMiJ9
\subsection*{C12-04-P}

\noindent\textbf{Enunciado}
\(C12-04-P \cdot\) Desplazamiento normal por espesor. Ejercicio de transferencia. En un punto de una línea media se conocen pendiente y semiespesor normal. Halle las dos componentes del desplazamiento normal y su componente axial como resultado final. Antes de consultar la solución, prediga el efecto de cambiar el primer dato respecto del caso resuelto de esta familia. Justifique la tendencia y explique qué restricción podría convertirse en limitante. Datos: \(t_n=0.002556\, \mathrm{m}\); \(m_y=0.73\).

\noindent\textbf{Datos}
\begin{itemize}
\item $t_n = 0.002556\,\mathrm{m}$ — semiespesor normal local
\item $m_y = 0.73$ — pendiente local
\end{itemize}

\noindent\textbf{Incógnita}
$\Delta x$, en $\mathrm{m}$; también las magnitudes intermedias desarrolladas.

\noindent\textbf{Modelo físico}

\noindent\textit{Desplazamiento normal por espesor}

\[
\fitmath{\Delta x=- \frac{m_y t_n}{\sqrt{m_y^{2} + 1}}}
\]

\noindent\textbf{Hipótesis}
\begin{itemize}
\item Superficie obtenida por desplazamiento normal local.
\end{itemize}

\noindent\textbf{Desarrollo paso a paso}

\begin{enumerate}

\item La normal unitaria tiene componente transversal positiva.

\[\fitmath{\mathrm{dy}=\frac{\mathrm{tt}}{\sqrt{1+{\mathrm{slope}}^{2}}}}\]

\[\fitmath{\mathrm{dy}=\frac{(0.002556)}{\sqrt{1+{(0.73)}^{2}}}\approx 0.002064448\,\mathrm{m}}\]

\item La componente axial compensa la inclinación local.

\[\fitmath{\mathrm{dx}=\frac{-\mathrm{tt}\cdot \mathrm{slope}}{\sqrt{1+{\mathrm{slope}}^{2}}}}\]

\[\fitmath{\mathrm{dx}=\frac{-(0.002556)\cdot (0.73)}{\sqrt{1+{(0.73)}^{2}}}\approx -0.001507047\,\mathrm{m}}\]

\end{enumerate}


\noindent\textbf{Sustitución numérica}

\[
\fitmath{- \frac{\left(0.73\right) \left(0.002556\,\mathrm{m}\right)}{\sqrt{\left(0.73\right)^{2} + 1}}}
\]

\noindent\textbf{Resultado}

\[
\fitmath{\Delta x=-1.507047002\times10^{-3}\,\mathrm{m}}
\]

\noindent\textbf{Interpretación}
La longitud del vector desplazamiento es tt. Añadir solo dy produce un espesor vertical distinto.

\noindent\textbf{Comprobación de los pasos}
\begin{itemize}
\item Paso 2: $0.00206445$ — reproducido
\item Paso 4: $-0.00150705$ — reproducido
\end{itemize}

\noindent\textbf{Comprobación de la respuesta} la aritmética del último paso reproduce la respuesta enunciada a la precisión con que está escrita. Es una reproducción aritmética; no valida el modelo físico ni sus hipótesis.

\noindent\textbf{Limitaciones}
\begin{itemize}
\item Modelo educativo; la aplicabilidad a una máquina real requiere revisión experta.
\end{itemize}

\noindent\textbf{Errores frecuentes}
\begin{itemize}
\item Olvidar el desplazamiento axial al añadir espesor.
\end{itemize}
% studium:end-record eyJpZCI6ICJQUkItMDA5NiIsICJzZWN0aW9uIjogImMxMiJ9
\section{Capítulo 13: Esfuerzos centrífugos, flexión y vibración}
\subsection*{C13-01-P}

\noindent\textbf{Enunciado}
\(C13-01-P \cdot\) Tracción centrífuga uniforme. Ejercicio de transferencia. Calcule la tensión centrífuga nominal de raíz de una pala uniforme. Antes de consultar la solución, prediga el efecto de cambiar el primer dato respecto del caso resuelto de esta familia. Justifique la tendencia y explique qué restricción podría convertirse en limitante. Datos: \(\rho=8647.8\, \mathrm{kg}\,\mathrm{m}^{-3}\); \(\Omega=1170\, \mathrm{s}^{-1}\); \(r_h=0.245\, \mathrm{m}\); \(r_t=0.302\, \mathrm{m}\).

\noindent\textbf{Datos}
\begin{itemize}
\item $\rho = 8647.8\,\mathrm{kg}\,\mathrm{m}^{-3}$ — densidad material
\item $\Omega = 1170\,\mathrm{s}^{-1}$ — velocidad angular
\item $r_h = 0.245\,\mathrm{m}$ — radio raíz
\item $r_t = 0.302\,\mathrm{m}$ — radio punta
\end{itemize}

\noindent\textbf{Incógnita}
$\sigma$, en $\mathrm{Pa}$; también las magnitudes intermedias desarrolladas.

\noindent\textbf{Modelo físico}

\noindent\textit{Tracción centrífuga uniforme}

\[
\fitmath{\sigma=\frac{\Omega^{2} \rho \left(- r_h^{2} + r_t^{2}\right)}{2}}
\]

\noindent\textbf{Hipótesis}
\begin{itemize}
\item Barra radial uniforme, material homogéneo, régimen estacionario.
\end{itemize}

\noindent\textbf{Desarrollo paso a paso}

\begin{enumerate}

\item Se integra r entre raíz y punta.

\[\fitmath{\mathrm{integ}=\frac{{\mathrm{rt}}^{2}-{\mathrm{rh}}^{2}}{2}}\]

\[\fitmath{\mathrm{integ}=\frac{{(0.302)}^{2}-{(0.245)}^{2}}{2}\approx 0.0155895\,\mathrm{m^{2}}}\]

\item Se multiplica por densidad y velocidad angular al cuadrado.

\[\fitmath{\mathrm{stress}=\frac{\rho\cdot {\Omega}^{2}\cdot \left({\mathrm{rt}}^{2}-{\mathrm{rh}}^{2}\right)}{2}}\]

\[\fitmath{\mathrm{stress}=\frac{(8647.8)\cdot {(1170)}^{2}\cdot \left({(0.302)}^{2}-{(0.245)}^{2}\right)}{2}\approx 1.845481\times10^{8}\,\mathrm{Pa}}\]

\end{enumerate}


\noindent\textbf{Sustitución numérica}

\[
\fitmath{\frac{\left(\left(1170\,\mathrm{s}^{-1}\right)\right)^{2} \left(8647.8\,\mathrm{kg}\,\mathrm{m}^{-3}\right) \left(- \left(0.245\,\mathrm{m}\right)^{2} + \left(0.302\,\mathrm{m}\right)^{2}\right)}{2}}
\]

\noindent\textbf{Resultado}

\[
\fitmath{\sigma=1.845480866\times10^{8}\,\mathrm{Pa}}
\]

\noindent\textbf{Interpretación}
La tensión escala con el cuadrado de giro. No incluye concentraciones de raíz ni flexión.

\noindent\textbf{Comprobación de los pasos}
\begin{itemize}
\item Paso 2: $0.0155895$ — reproducido
\item Paso 4: $1.84548\times 10^{8}$ — reproducido
\end{itemize}

\noindent\textbf{Comprobación de la respuesta} la aritmética del último paso reproduce la respuesta enunciada a la precisión con que está escrita. Es una reproducción aritmética; no valida el modelo físico ni sus hipótesis.

\noindent\textbf{Limitaciones}
\begin{itemize}
\item Modelo educativo; la aplicabilidad a una máquina real requiere revisión experta.
\end{itemize}

\noindent\textbf{Errores frecuentes}
\begin{itemize}
\item Usar densidad del gas o rpm directamente.
\end{itemize}
% studium:end-record eyJpZCI6ICJQUkItMDA5OCIsICJzZWN0aW9uIjogImMxMyJ9
\subsection*{C13-02-P}

\noindent\textbf{Enunciado}
\(C13-02-P \cdot\) Pala estrechada. Ejercicio de transferencia. La sección disminuye linealmente desde Ah en la raíz hasta Ah(1-k) en punta. Calcule la tensión centrífuga nominal de raíz. Antes de consultar la solución, prediga el efecto de cambiar el primer dato respecto del caso resuelto de esta familia. Justifique la tendencia y explique qué restricción podría convertirse en limitante. Datos: \(\rho=8573.25\, \mathrm{kg}\,\mathrm{m}^{-3}\); \(\Omega=1210\, \mathrm{s}^{-1}\); \(r_h=0.238\, \mathrm{m}\); \(r_t=0.301\, \mathrm{m}\); \(\alpha_t=0.42\).

\noindent\textbf{Datos}
\begin{itemize}
\item $\rho = 8573.25\,\mathrm{kg}\,\mathrm{m}^{-3}$ — densidad material
\item $\Omega = 1210\,\mathrm{s}^{-1}$ — velocidad angular
\item $r_h = 0.238\,\mathrm{m}$ — radio raíz
\item $r_t = 0.301\,\mathrm{m}$ — radio punta
\item $\alpha_t = 0.42$ — fracción de estrechamiento
\end{itemize}

\noindent\textbf{Incógnita}
$\sigma$, en $\mathrm{Pa}$; también las magnitudes intermedias desarrolladas.

\noindent\textbf{Modelo físico}

\noindent\textit{Pala estrechada}

\[
\fitmath{\sigma=\Omega^{2} \rho \left(- \frac{\alpha_t \left(- \frac{r_h^{3}}{3} - \frac{r_h \left(- r_h^{2} + r_t^{2}\right)}{2} + \frac{r_t^{3}}{3}\right)}{- r_h + r_t} - \frac{r_h^{2}}{2} + \frac{r_t^{2}}{2}\right)}
\]

\noindent\textbf{Hipótesis}
\begin{itemize}
\item Área lineal positiva en todo el intervalo; sin carga aerodinámica.
\end{itemize}

\noindent\textbf{Desarrollo paso a paso}

\begin{enumerate}

\item El ancho radial define la ley de área.

\[\fitmath{b=\mathrm{rt}-\mathrm{rh}}\]

\[\fitmath{b=(0.301)-(0.238)\approx 0.063\,\mathrm{m}}\]

\item La integral se divide por el área raíz; Ah se cancela.

\[\fitmath{\mathrm{stress}=\rho\cdot {\Omega}^{2}\cdot \left(\frac{{\mathrm{rt}}^{2}-{\mathrm{rh}}^{2}}{2}-\frac{\mathrm{kk}}{\mathrm{rt}-\mathrm{rh}}\cdot \left(\frac{{\mathrm{rt}}^{3}-{\mathrm{rh}}^{3}}{3}-\frac{\mathrm{rh}\cdot \left({\mathrm{rt}}^{2}-{\mathrm{rh}}^{2}\right)}{2}\right)\right)}\]

\[\fitmath{\mathrm{stress}=(8573.25)\cdot {(1210)}^{2}\cdot \left(\frac{{(0.301)}^{2}-{(0.238)}^{2}}{2}-\frac{(0.42)}{(0.301)-(0.238)}\cdot \left(\frac{{(0.301)}^{3}-{(0.238)}^{3}}{3}-\frac{(0.238)\cdot \left({(0.301)}^{2}-{(0.238)}^{2}\right)}{2}\right)\right)\approx 1.666178\times10^{8}\,\mathrm{Pa}}\]

\end{enumerate}


\noindent\textbf{Sustitución numérica}

\[
\fitmath{\left(\left(1210\,\mathrm{s}^{-1}\right)\right)^{2} \left(8573.25\,\mathrm{kg}\,\mathrm{m}^{-3}\right) \left(- \frac{\left(0.42\right) \left(- \frac{\left(0.238\,\mathrm{m}\right)^{3}}{3} - \frac{\left(0.238\,\mathrm{m}\right) \left(- \left(0.238\,\mathrm{m}\right)^{2} + \left(0.301\,\mathrm{m}\right)^{2}\right)}{2} + \frac{\left(0.301\,\mathrm{m}\right)^{3}}{3}\right)}{- \left(0.238\,\mathrm{m}\right) + \left(0.301\,\mathrm{m}\right)} - \frac{\left(0.238\,\mathrm{m}\right)^{2}}{2} + \frac{\left(0.301\,\mathrm{m}\right)^{2}}{2}\right)}
\]

\noindent\textbf{Resultado}

\[
\fitmath{\sigma=1.666177686\times10^{8}\,\mathrm{Pa}}
\]

\noindent\textbf{Interpretación}
La reducción de masa exterior reduce la tracción raíz respecto de una pala uniforme del mismo material y radios.

\noindent\textbf{Comprobación de los pasos}
\begin{itemize}
\item Paso 2: $0.063$ — reproducido
\item Paso 4: $1.66618\times 10^{8}$ — reproducido
\end{itemize}

\noindent\textbf{Comprobación de la respuesta} la aritmética del último paso reproduce la respuesta enunciada a la precisión con que está escrita. Es una reproducción aritmética; no valida el modelo físico ni sus hipótesis.

\noindent\textbf{Limitaciones}
\begin{itemize}
\item Modelo educativo; la aplicabilidad a una máquina real requiere revisión experta.
\end{itemize}

\noindent\textbf{Errores frecuentes}
\begin{itemize}
\item Reducir tensión multiplicando por el área de punta sin integrar.
\end{itemize}
% studium:end-record eyJpZCI6ICJQUkItMDEwMCIsICJzZWN0aW9uIjogImMxMyJ9
\subsection*{C13-03-P}

\noindent\textbf{Enunciado}
\(C13-03-P \cdot\) Flexión nominal y tracción. Ejercicio de transferencia. Una pala tiene fuerza aerodinámica equivalente aplicada al brazo dado y tracción nominal conocida. Halle la máxima tensión normal en la fibra indicada. Antes de consultar la solución, prediga el efecto de cambiar el primer dato respecto del caso resuelto de esta familia. Justifique la tendencia y explique qué restricción podría convertirse en limitante. Datos: \(F_t=198.09\, \mathrm{N}\); \(L_b=0.038\, \mathrm{m}\); \(y=0.0042\, \mathrm{m}\); \(I_b=8.4\times 10^{-10}\, \mathrm{m}^{4}\); \(\sigma_N=1.78\times 10^{8}\, \mathrm{Pa}\).

\noindent\textbf{Datos}
\begin{itemize}
\item $F_t = 198.09\,\mathrm{N}$ — fuerza transversal equivalente
\item $L_b = 0.038\,\mathrm{m}$ — brazo respecto raíz
\item $y = 0.0042\,\mathrm{m}$ — distancia de fibra al eje neutro
\item $I_b = 8.4\times 10^{-10}\,\mathrm{m}^{4}$ — segundo momento de área
\item $\sigma_N = 178000000\,\mathrm{Pa}$ — tracción centrífuga nominal
\end{itemize}

\noindent\textbf{Incógnita}
$\sigma$, en $\mathrm{Pa}$; también las magnitudes intermedias desarrolladas.

\noindent\textbf{Modelo físico}

\noindent\textit{Flexión nominal y tracción}

\[
\fitmath{\sigma=\sigma_N + \frac{L_b F_t y}{I_b}}
\]

\noindent\textbf{Hipótesis}
\begin{itemize}
\item Viga elástica lineal, cargas y sección equivalentes.
\end{itemize}

\noindent\textbf{Desarrollo paso a paso}

\begin{enumerate}

\item La fuerza equivalente genera el momento de raíz.

\[\fitmath{\mathrm{moment}=\mathrm{force}\cdot \mathrm{arm}}\]

\[\fitmath{\mathrm{moment}=(198.09)\cdot (0.038)\approx 7.52742\,\mathrm{N\,m}}\]

\item Se añade la flexión en la fibra de máxima tracción.

\[\fitmath{\mathrm{stress}=\mathrm{sig}+\frac{\mathrm{force}\cdot \mathrm{arm}\cdot \mathrm{yy}}{\mathrm{Ib}}}\]

\[\fitmath{\mathrm{stress}=(178000000)+\frac{(198.09)\cdot (0.038)\cdot (0.0042)}{(8.4\times10^{-10})}\approx 2.156371\times10^{8}\,\mathrm{Pa}}\]

\end{enumerate}


\noindent\textbf{Sustitución numérica}

\[
\fitmath{\left(178000000\,\mathrm{Pa}\right) + \frac{\left(0.038\,\mathrm{m}\right) \left(198.09\,\mathrm{N}\right) \left(0.0042\,\mathrm{m}\right)}{\left(8.4e-10\,\mathrm{m}^{4}\right)}}
\]

\noindent\textbf{Resultado}

\[
\fitmath{\sigma=2.156371000\times10^{8}\,\mathrm{Pa}}
\]

\noindent\textbf{Interpretación}
La fibra opuesta resta la contribución de flexión. La suma nominal no evalúa fatiga ni concentraciones.

\noindent\textbf{Comprobación de los pasos}
\begin{itemize}
\item Paso 2: $7.52742$ — reproducido
\item Paso 4: $215637100$ — reproducido
\end{itemize}

\noindent\textbf{Comprobación de la respuesta} la aritmética del último paso reproduce la respuesta enunciada a la precisión con que está escrita. Es una reproducción aritmética; no valida el modelo físico ni sus hipótesis.

\noindent\textbf{Limitaciones}
\begin{itemize}
\item Modelo educativo; la aplicabilidad a una máquina real requiere revisión experta.
\end{itemize}

\noindent\textbf{Errores frecuentes}
\begin{itemize}
\item Usar momento polar en una fórmula de flexión sin justificarlo.
\end{itemize}
% studium:end-record eyJpZCI6ICJQUkItMDEwMiIsICJzZWN0aW9uIjogImMxMyJ9
\subsection*{C13-04-P}

\noindent\textbf{Enunciado}
\(C13-04-P \cdot\) Frecuencia del voladizo educativo. Ejercicio de transferencia. Estime la primera frecuencia de una pala como viga uniforme no giratoria. Antes de consultar la solución, prediga el efecto de cambiar el primer dato respecto del caso resuelto de esta familia. Justifique la tendencia y explique qué restricción podría convertirse en limitante. Datos: \(E=1.97025\times 10^{11}\, \mathrm{Pa}\); \(I_b=7.8\times 10^{-10}\, \mathrm{m}^{4}\); \(\rho=8030\, \mathrm{kg}\,\mathrm{m}^{-3}\); \(A_b=0.000165\, \mathrm{m}^{2}\); \(L=0.082\, \mathrm{m}\).

\noindent\textbf{Datos}
\begin{itemize}
\item $E = 197025000000\,\mathrm{Pa}$ — módulo elástico
\item $I_b = 7.8\times 10^{-10}\,\mathrm{m}^{4}$ — segundo momento
\item $\rho = 8030\,\mathrm{kg}\,\mathrm{m}^{-3}$ — densidad material
\item $A_b = 0.000165\,\mathrm{m}^{2}$ — área de sección
\item $L = 0.082\,\mathrm{m}$ — longitud
\end{itemize}

\noindent\textbf{Incógnita}
$f_1$, en $\mathrm{s^{-1}}$; también las magnitudes intermedias desarrolladas.

\noindent\textbf{Modelo físico}

\noindent\textit{Frecuencia del voladizo educativo}

\[
\fitmath{f_1=\frac{1.7578125 \sqrt{\frac{E I_b}{A_b \rho}}}{\pi L^{2}}}
\]

\noindent\textbf{Hipótesis}
\begin{itemize}
\item Voladizo uniforme Euler-Bernoulli, sin rotación ni amortiguamiento.
\end{itemize}

\noindent\textbf{Desarrollo paso a paso}

\begin{enumerate}

\item Se calcula la rigidez específica de flexión.

\[\fitmath{\mathrm{stiff}=\sqrt{\frac{\mathrm{Ee}\cdot \mathrm{Ib}}{\rho\cdot \mathrm{Ab}}}}\]

\[\fitmath{\mathrm{stiff}=\sqrt{\frac{(1.97025\times10^{11})\cdot (7.8\times10^{-10})}{(8030)\cdot (0.000165)}}\approx 10.76981\,\mathrm{m^{2}/s}}\]

\item La primera raíz modal fija la frecuencia.

\[\fitmath{\mathrm{freq}=\frac{{1.875}^{2}}{2\cdot \pi\cdot {\mathrm{length}}^{2}}\cdot \sqrt{\frac{\mathrm{Ee}\cdot \mathrm{Ib}}{\rho\cdot \mathrm{Ab}}}}\]

\[\fitmath{\mathrm{freq}=\frac{{1.875}^{2}}{2\cdot \pi\cdot {(0.082)}^{2}}\cdot \sqrt{\frac{(1.97025\times10^{11})\cdot (7.8\times10^{-10})}{(8030)\cdot (0.000165)}}\approx 896.1964\,\mathrm{s^{-1}}}\]

\end{enumerate}


\noindent\textbf{Sustitución numérica}

\[
\fitmath{\frac{1.7578125 \sqrt{\frac{\left(197025000000\,\mathrm{Pa}\right) \left(7.8e-10\,\mathrm{m}^{4}\right)}{\left(0.000165\,\mathrm{m}^{2}\right) \left(8030\,\mathrm{kg}\,\mathrm{m}^{-3}\right)}}}{\pi \left(0.082\,\mathrm{m}\right)^{2}}}
\]

\noindent\textbf{Resultado}

\[
\fitmath{f_1=8.961964014\times10^{2}\,\mathrm{s}^{-1}}
\]

\noindent\textbf{Interpretación}
Es una estimación no giratoria; no incluye rigidización centrífuga ni anclaje real. Compare después con órdenes de excitación definidos.

\noindent\textbf{Comprobación de los pasos}
\begin{itemize}
\item Paso 2: $10.7698$ — reproducido
\item Paso 4: $896.196$ — reproducido
\end{itemize}

\noindent\textbf{Comprobación de la respuesta} la aritmética del último paso reproduce la respuesta enunciada a la precisión con que está escrita. Es una reproducción aritmética; no valida el modelo físico ni sus hipótesis.

\noindent\textbf{Limitaciones}
\begin{itemize}
\item Modelo educativo; la aplicabilidad a una máquina real requiere revisión experta.
\end{itemize}

\noindent\textbf{Errores frecuentes}
\begin{itemize}
\item Declarar ausencia de resonancia a partir de un único modelo estático.
\end{itemize}
% studium:end-record eyJpZCI6ICJQUkItMDEwNCIsICJzZWN0aW9uIjogImMxMyJ9
\section{Capítulo 14: Transferencia térmica y refrigeración de álabes}
\subsection*{C14-01-P}

\noindent\textbf{Enunciado}
\(C14-01-P \cdot\) Red térmica de pared. Ejercicio de transferencia. Calcule flujo térmico y temperatura exterior del metal en una pared plana con convección a ambos lados. Antes de consultar la solución, prediga el efecto de cambiar el primer dato respecto del caso resuelto de esta familia. Justifique la tendencia y explique qué restricción podría convertirse en limitante. Datos: \(T_{aw}=1725.3\, \mathrm{K}\); \(T_c=745\, \mathrm{K}\); \(h_g=1680\, \mathrm{W}\,\mathrm{m}^{-2}\,\mathrm{K}^{-1}\); \(h_c=2450\, \mathrm{W}\,\mathrm{m}^{-2}\,\mathrm{K}^{-1}\); \(t=0.0011\, \mathrm{m}\); \(k_m=21\, \mathrm{W}\,\mathrm{m}^{-1}\,\mathrm{K}^{-1}\).

\noindent\textbf{Datos}
\begin{itemize}
\item $T_{aw} = 1725.3\,\mathrm{K}$ — temperatura recuperación exterior
\item $T_c = 745\,\mathrm{K}$ — temperatura refrigerante local
\item $h_g = 1680\,\mathrm{W}\,\mathrm{m}^{-2}\,\mathrm{K}^{-1}$ — coeficiente exterior
\item $h_c = 2450\,\mathrm{W}\,\mathrm{m}^{-2}\,\mathrm{K}^{-1}$ — coeficiente interior
\item $t = 0.0011\,\mathrm{m}$ — espesor
\item $k_m = 21\,\mathrm{W}\,\mathrm{m}^{-1}\,\mathrm{K}^{-1}$ — conductividad metal
\end{itemize}

\noindent\textbf{Incógnita}
$T_w$, en $\mathrm{K}$; también las magnitudes intermedias desarrolladas.

\noindent\textbf{Modelo físico}

\noindent\textit{Red térmica de pared}

\[
\fitmath{T_w=T_{aw} - \frac{- T_c + T_{aw}}{h_g \left(\frac{1}{h_g} + \frac{1}{h_c} + \frac{t}{k_m}\right)}}
\]

\noindent\textbf{Hipótesis}
\begin{itemize}
\item Estacionario, una dimensión, sin radiación, coeficientes constantes.
\end{itemize}

\noindent\textbf{Desarrollo paso a paso}

\begin{enumerate}

\item Las resistencias se suman por unidad de área.

\[\fitmath{\mathrm{Res}=\frac{1}{\mathrm{hg}}+\frac{\mathrm{th}}{\mathrm{cond}}+\frac{1}{\mathrm{hc}}}\]

\[\fitmath{\mathrm{Res}=\frac{1}{(1680)}+\frac{(0.0011)}{(21)}+\frac{1}{(2450)}\approx 0.001055782\,\mathrm{m^{2}\,K/W}}\]

\item El salto térmico impulsa el flujo.

\[\fitmath{\mathrm{flux}=\frac{\mathrm{Tg}-\mathrm{Tc}}{\frac{1}{\mathrm{hg}}+\frac{\mathrm{th}}{\mathrm{cond}}+\frac{1}{\mathrm{hc}}}}\]

\[\fitmath{\mathrm{flux}=\frac{(1725.3)-(745)}{\frac{1}{(1680)}+\frac{(0.0011)}{(21)}+\frac{1}{(2450)}}\approx 928505.8\,\mathrm{W/m^{2}}}\]

\item La caída exterior determina la cara caliente.

\[\fitmath{\mathrm{Tw}=\mathrm{Tg}-\frac{\mathrm{Tg}-\mathrm{Tc}}{\mathrm{hg}\cdot \left(\frac{1}{\mathrm{hg}}+\frac{\mathrm{th}}{\mathrm{cond}}+\frac{1}{\mathrm{hc}}\right)}}\]

\[\fitmath{\mathrm{Tw}=(1725.3)-\frac{(1725.3)-(745)}{(1680)\cdot \left(\frac{1}{(1680)}+\frac{(0.0011)}{(21)}+\frac{1}{(2450)}\right)}\approx 1172.618\,\mathrm{K}}\]

\end{enumerate}


\noindent\textbf{Sustitución numérica}

\[
\fitmath{\left(1725.3\,\mathrm{K}\right) - \frac{- \left(745\,\mathrm{K}\right) + \left(1725.3\,\mathrm{K}\right)}{\left(1680\,\mathrm{W}\,\mathrm{m}^{-2}\,\mathrm{K}^{-1}\right) \left(\frac{1}{\left(1680\,\mathrm{W}\,\mathrm{m}^{-2}\,\mathrm{K}^{-1}\right)} + \frac{1}{\left(2450\,\mathrm{W}\,\mathrm{m}^{-2}\,\mathrm{K}^{-1}\right)} + \frac{\left(0.0011\,\mathrm{m}\right)}{\left(21\,\mathrm{W}\,\mathrm{m}^{-1}\,\mathrm{K}^{-1}\right)}\right)}}
\]

\noindent\textbf{Resultado}

\[
\fitmath{T_w=1.172617977\times10^{3}\,\mathrm{K}}
\]

\noindent\textbf{Interpretación}
La temperatura del metal debe quedar entre refrigerante y recuperación exterior. La diferencia de caras corresponde al flujo por conducción.

\noindent\textbf{Comprobación de los pasos}
\begin{itemize}
\item Paso 2: $0.00105578$ — reproducido
\item Paso 4: $928506$ — reproducido
\item Paso 6: $1172.62$ — reproducido
\end{itemize}

\noindent\textbf{Comprobación de la respuesta} la aritmética del último paso reproduce la respuesta enunciada a la precisión con que está escrita. Es una reproducción aritmética; no valida el modelo físico ni sus hipótesis.

\noindent\textbf{Limitaciones}
\begin{itemize}
\item Modelo educativo; la aplicabilidad a una máquina real requiere revisión experta.
\end{itemize}

\noindent\textbf{Errores frecuentes}
\begin{itemize}
\item Confundir área anular con área de intercambio.
\end{itemize}
% studium:end-record eyJpZCI6ICJQUkItMDEwNiIsICJzZWN0aW9uIjogImMxNCJ9
\subsection*{C14-02-P}

\noindent\textbf{Enunciado}
\(C14-02-P \cdot\) Caudal energético de refrigeración. Ejercicio de transferencia. Una superficie recibe flujo térmico medio conocido. Calcule el caudal refrigerante que absorbería ese calor con el incremento permitido. Antes de consultar la solución, prediga el efecto de cambiar el primer dato respecto del caso resuelto de esta familia. Justifique la tendencia y explique qué restricción podría convertirse en limitante. Datos: \(q''=660300\, \mathrm{W}\,\mathrm{m}^{-2}\); \(A=0.034\, \mathrm{m}^{2}\); \(c_p=1090\, \mathrm{J}\,\mathrm{kg}^{-1}\,\mathrm{K}^{-1}\); \(\Delta T_c=185\, \mathrm{K}\).

\noindent\textbf{Datos}
\begin{itemize}
\item $q'' = 660300\,\mathrm{W}\,\mathrm{m}^{-2}$ — flujo térmico medio
\item $A = 0.034\,\mathrm{m}^{2}$ — área de intercambio
\item $c_p = 1090\,\mathrm{J}\,\mathrm{kg}^{-1}\,\mathrm{K}^{-1}$ — calor específico refrigerante
\item $\Delta T_c = 185\,\mathrm{K}$ — incremento de temperatura permitido
\end{itemize}

\noindent\textbf{Incógnita}
$\dot{m}_c$, en $\mathrm{kg\,s^{-1}}$; también las magnitudes intermedias desarrolladas.

\noindent\textbf{Modelo físico}

\noindent\textit{Caudal energético de refrigeración}

\[
\fitmath{\dot{m}_c=\frac{A q''}{c_p \Delta T_c}}
\]

\noindent\textbf{Hipótesis}
\begin{itemize}
\item Sin pérdidas térmicas adicionales; cp constante.
\end{itemize}

\noindent\textbf{Desarrollo paso a paso}

\begin{enumerate}

\item Se integra el flujo térmico uniforme sobre la superficie.

\[\fitmath{\mathrm{heat}=\mathrm{flux}\cdot \mathrm{area}}\]

\[\fitmath{\mathrm{heat}=(660300)\cdot (0.034)\approx 22450.2\,\mathrm{W}}\]

\item El balance de entalpía fija el caudal.

\[\fitmath{\mathrm{mc}=\frac{\mathrm{flux}\cdot \mathrm{area}}{\mathrm{cp}\cdot \mathrm{dt}}}\]

\[\fitmath{\mathrm{mc}=\frac{(660300)\cdot (0.034)}{(1090)\cdot (185)}\approx 0.1113325\,\mathrm{kg/s}}\]

\end{enumerate}


\noindent\textbf{Sustitución numérica}

\[
\fitmath{\frac{\left(0.034\,\mathrm{m}^{2}\right) \left(660300\,\mathrm{W}\,\mathrm{m}^{-2}\right)}{\left(1090\,\mathrm{J}\,\mathrm{kg}^{-1}\,\mathrm{K}^{-1}\right) \left(185\,\mathrm{K}\right)}}
\]

\noindent\textbf{Resultado}

\[
\fitmath{\dot{m}_c=1.113325068\times10^{-1}\,\mathrm{kg}\,\mathrm{s}^{-1}}
\]

\noindent\textbf{Interpretación}
El resultado es un mínimo del balance energético. La convección y la presión disponible deben comprobarse aparte.

\noindent\textbf{Comprobación de los pasos}
\begin{itemize}
\item Paso 2: $22450.2$ — reproducido
\item Paso 4: $0.111333$ — reproducido
\end{itemize}

\noindent\textbf{Comprobación de la respuesta} la aritmética del último paso reproduce la respuesta enunciada a la precisión con que está escrita. Es una reproducción aritmética; no valida el modelo físico ni sus hipótesis.

\noindent\textbf{Limitaciones}
\begin{itemize}
\item Modelo educativo; la aplicabilidad a una máquina real requiere revisión experta.
\end{itemize}

\noindent\textbf{Errores frecuentes}
\begin{itemize}
\item Tomar este mínimo como caudal seguro de una pala real.
\end{itemize}
% studium:end-record eyJpZCI6ICJQUkItMDEwOCIsICJzZWN0aW9uIjogImMxNCJ9
\subsection*{C14-03-P}

\noindent\textbf{Enunciado}
\(C14-03-P \cdot\) Película y recuperación. Ejercicio de transferencia. Con efectividad de película prescrita, calcule la recuperación efectiva y el flujo térmico a una pared de temperatura dada. Antes de consultar la solución, prediga el efecto de cambiar el primer dato respecto del caso resuelto de esta familia. Justifique la tendencia y explique qué restricción podría convertirse en limitante. Datos: \(T_{aw}=1799.85\, \mathrm{K}\); \(T_c=775\, \mathrm{K}\); \(\eta_f=0.41\); \(T_w=1130\, \mathrm{K}\); \(h_g=1770\, \mathrm{W}\,\mathrm{m}^{-2}\,\mathrm{K}^{-1}\).

\noindent\textbf{Datos}
\begin{itemize}
\item $T_{aw} = 1799.85\,\mathrm{K}$ — recuperación sin película
\item $T_c = 775\,\mathrm{K}$ — temperatura refrigerante
\item $\eta_f = 0.41$ — efectividad de película
\item $T_w = 1130\,\mathrm{K}$ — temperatura pared
\item $h_g = 1770\,\mathrm{W}\,\mathrm{m}^{-2}\,\mathrm{K}^{-1}$ — coeficiente exterior compatible con inyección
\end{itemize}

\noindent\textbf{Incógnita}
$q''$, en $\mathrm{W\,m^{-2}}$; también las magnitudes intermedias desarrolladas.

\noindent\textbf{Modelo físico}

\noindent\textit{Película y recuperación}

\[
\fitmath{q''=h_g \left(T_{aw} - T_w - \eta_f \left(- T_c + T_{aw}\right)\right)}
\]

\noindent\textbf{Hipótesis}
\begin{itemize}
\item Efectividad y coeficiente dados para la misma geometría y punto.
\end{itemize}

\noindent\textbf{Desarrollo paso a paso}

\begin{enumerate}

\item Se despeja la recuperación con película.

\[\fitmath{\mathrm{Taw}=\mathrm{Tg}-\mathrm{eff}\cdot \left(\mathrm{Tg}-\mathrm{Tc}\right)}\]

\[\fitmath{\mathrm{Taw}=(1799.85)-(0.41)\cdot \left((1799.85)-(775)\right)\approx 1379.661\,\mathrm{K}}\]

\item Se aplica convección con el coeficiente especificado.

\[\fitmath{\mathrm{flux}=\mathrm{hg}\cdot \left(\mathrm{Tg}-\mathrm{eff}\cdot \left(\mathrm{Tg}-\mathrm{Tc}\right)-\mathrm{Tw}\right)}\]

\[\fitmath{\mathrm{flux}=(1770)\cdot \left((1799.85)-(0.41)\cdot \left((1799.85)-(775)\right)-(1130)\right)\approx 441900.9\,\mathrm{W/m^{2}}}\]

\end{enumerate}


\noindent\textbf{Sustitución numérica}

\[
\fitmath{\left(1770\,\mathrm{W}\,\mathrm{m}^{-2}\,\mathrm{K}^{-1}\right) \left(\left(1799.85\,\mathrm{K}\right) - \left(1130\,\mathrm{K}\right) - \left(0.41\right) \left(- \left(775\,\mathrm{K}\right) + \left(1799.85\,\mathrm{K}\right)\right)\right)}
\]

\noindent\textbf{Resultado}

\[
\fitmath{q''=4.419008550\times10^{5}\,\mathrm{W}\,\mathrm{m}^{-2}}
\]

\noindent\textbf{Interpretación}
La efectividad no sustituye al coeficiente de convección. Verifique que el flujo conserva el sentido esperado.

\noindent\textbf{Comprobación de los pasos}
\begin{itemize}
\item Paso 2: $1379.66$ — reproducido
\item Paso 4: $441901$ — reproducido
\end{itemize}

\noindent\textbf{Comprobación de la respuesta} la aritmética del último paso reproduce la respuesta enunciada a la precisión con que está escrita. Es una reproducción aritmética; no valida el modelo físico ni sus hipótesis.

\noindent\textbf{Limitaciones}
\begin{itemize}
\item Modelo educativo; la aplicabilidad a una máquina real requiere revisión experta.
\end{itemize}

\noindent\textbf{Errores frecuentes}
\begin{itemize}
\item Identificar efectividad de película y eficiencia de turbina.
\end{itemize}
% studium:end-record eyJpZCI6ICJQUkItMDExMCIsICJzZWN0aW9uIjogImMxNCJ9
\subsection*{C14-04-P}

\noindent\textbf{Enunciado}
\(C14-04-P \cdot\) Restricción térmica uniaxial. Ejercicio de transferencia. Una barra homogénea tiene impedida completamente su dilatación longitudinal. Determine la deformación térmica libre y la tensión elástica de compresión. Antes de consultar la solución, prediga el efecto de cambiar el primer dato respecto del caso resuelto de esta familia. Justifique la tendencia y explique qué restricción podría convertirse en limitante. Datos: \(E=1.8531\times 10^{11}\, \mathrm{Pa}\); \(\alpha_{\mathrm{th}}=1.41\times 10^{-5}\, \mathrm{K}^{-1}\); \(\Delta T=128\, \mathrm{K}\).

\noindent\textbf{Datos}
\begin{itemize}
\item $E = 185310000000\,\mathrm{Pa}$ — módulo elástico
\item $\alpha_{\mathrm{th}} = 1.41\times 10^{-5}\,\mathrm{K}^{-1}$ — coeficiente dilatación
\item $\Delta T = 128\,\mathrm{K}$ — aumento uniforme de temperatura
\end{itemize}

\noindent\textbf{Incógnita}
$\sigma$, en $\mathrm{Pa}$; también las magnitudes intermedias desarrolladas.

\noindent\textbf{Modelo físico}

\noindent\textit{Restricción térmica uniaxial}

\[
\fitmath{\sigma=- E \alpha_{\mathrm{th}} \Delta T}
\]

\noindent\textbf{Hipótesis}
\begin{itemize}
\item Restricción total uniaxial, elasticidad lineal, sin relajación.
\end{itemize}

\noindent\textbf{Desarrollo paso a paso}

\begin{enumerate}

\item La expansión libre sería positiva.

\[\fitmath{\mathrm{strain}=\mathrm{alp}\cdot \mathrm{dt}}\]

\[\fitmath{\mathrm{strain}=(1.41\times10^{-5})\cdot (128)\approx 0.0018048\,}\]

\item La restricción exige tensión elástica de signo opuesto.

\[\fitmath{\mathrm{stress}=-\mathrm{Ee}\cdot \mathrm{alp}\cdot \mathrm{dt}}\]

\[\fitmath{\mathrm{stress}=-(1.8531\times10^{11})\cdot (1.41\times10^{-5})\cdot (128)\approx -3.344475\times10^{8}\,\mathrm{Pa}}\]

\end{enumerate}


\noindent\textbf{Sustitución numérica}

\[
\fitmath{- \left(185310000000\,\mathrm{Pa}\right) \left(1.41e-05\,\mathrm{K}^{-1}\right) \left(128\,\mathrm{K}\right)}
\]

\noindent\textbf{Resultado}

\[
\fitmath{\sigma=-3.344474880\times10^{8}\,\mathrm{Pa}}
\]

\noindent\textbf{Interpretación}
El signo negativo indica compresión. Una barra libre tendría tensión térmica nula con la misma temperatura uniforme.

\noindent\textbf{Comprobación de los pasos}
\begin{itemize}
\item Paso 2: $0.0018048$ — reproducido
\item Paso 4: $-334447488$ — reproducido
\end{itemize}

\noindent\textbf{Comprobación de la respuesta} la aritmética del último paso reproduce la respuesta enunciada a la precisión con que está escrita. Es una reproducción aritmética; no valida el modelo físico ni sus hipótesis.

\noindent\textbf{Limitaciones}
\begin{itemize}
\item Modelo educativo; la aplicabilidad a una máquina real requiere revisión experta.
\end{itemize}

\noindent\textbf{Errores frecuentes}
\begin{itemize}
\item Aplicar esta tensión a toda pared con gradiente y curvatura.
\end{itemize}
% studium:end-record eyJpZCI6ICJQUkItMDExMiIsICJzZWN0aW9uIjogImMxNCJ9
\section{Capítulo 15: Compresores radiales: rodete, inductor y deslizamiento}
\subsection*{C15-01-P}

\noindent\textbf{Enunciado}
\(C15-01-P \cdot\) Rodete radial con deslizamiento. Ejercicio de transferencia. Sin prerrotación y con álabes radiales ideales, calcule remolino real, trabajo y aumento total de temperatura. Antes de consultar la solución, prediga el efecto de cambiar el primer dato respecto del caso resuelto de esta familia. Justifique la tendencia y explique qué restricción podría convertirse en limitante. Datos: \(U=453.69\, \mathrm{m}\,\mathrm{s}^{-1}\); \(\sigma_s=0.87\); \(c_p=1004.5\, \mathrm{J}\,\mathrm{kg}^{-1}\,\mathrm{K}^{-1}\).

\noindent\textbf{Datos}
\begin{itemize}
\item $U = 453.69\,\mathrm{m}\,\mathrm{s}^{-1}$ — velocidad periférica salida
\item $\sigma_s = 0.87$ — factor definido como $C_\theta/C_{\theta,\mathrm{ideal}}$
\item $c_p = 1004.5\,\mathrm{J}\,\mathrm{kg}^{-1}\,\mathrm{K}^{-1}$ — calor específico
\end{itemize}

\noindent\textbf{Incógnita}
$\Delta T_0$, en $\mathrm{K}$; también las magnitudes intermedias desarrolladas.

\noindent\textbf{Modelo físico}

\noindent\textit{Rodete radial con deslizamiento}

\[
\fitmath{\Delta T_0=\frac{\sigma_s U^{2}}{c_p}}
\]

\noindent\textbf{Hipótesis}
\begin{itemize}
\item Sin prerrotación; salida ideal radial antes de aplicar slip.
\end{itemize}

\noindent\textbf{Desarrollo paso a paso}

\begin{enumerate}

\item Se reduce el remolino ideal con la definición indicada.

\[\fitmath{\mathrm{ct}=\mathrm{slip}\cdot \mathrm{uu}}\]

\[\fitmath{\mathrm{ct}=(0.87)\cdot (453.69)\approx 394.7103\,\mathrm{m/s}}\]

\item Euler usa el remolino real.

\[\fitmath{\mathrm{wc}=\mathrm{slip}\cdot {\mathrm{uu}}^{2}}\]

\[\fitmath{\mathrm{wc}=(0.87)\cdot {(453.69)}^{2}\approx 179076.1\,\mathrm{J/kg}}\]

\item El balance energético fija el calentamiento total.

\[\fitmath{\mathrm{dT}=\frac{\mathrm{slip}\cdot {\mathrm{uu}}^{2}}{\mathrm{cp}}}\]

\[\fitmath{\mathrm{dT}=\frac{(0.87)\cdot {(453.69)}^{2}}{(1004.5)}\approx 178.2739\,\mathrm{K}}\]

\end{enumerate}


\noindent\textbf{Sustitución numérica}

\[
\fitmath{\frac{\left(0.87\right) \left(\left(453.69\,\mathrm{m}\,\mathrm{s}^{-1}\right)\right)^{2}}{\left(1004.5\,\mathrm{J}\,\mathrm{kg}^{-1}\,\mathrm{K}^{-1}\right)}}
\]

\noindent\textbf{Resultado}

\[
\fitmath{\Delta T_0=1.782738835\times10^{2}\,\mathrm{K}}
\]

\noindent\textbf{Interpretación}
El deslizamiento modifica trabajo. La presión necesita además una eficiencia o modelo de pérdida.

\noindent\textbf{Comprobación de los pasos}
\begin{itemize}
\item Paso 2: $394.71$ — reproducido
\item Paso 4: $179076$ — reproducido
\item Paso 6: $178.274$ — reproducido
\end{itemize}

\noindent\textbf{Comprobación de la respuesta} la aritmética del último paso reproduce la respuesta enunciada a la precisión con que está escrita. Es una reproducción aritmética; no valida el modelo físico ni sus hipótesis.

\noindent\textbf{Limitaciones}
\begin{itemize}
\item Modelo educativo; la aplicabilidad a una máquina real requiere revisión experta.
\end{itemize}

\noindent\textbf{Errores frecuentes}
\begin{itemize}
\item Multiplicar de nuevo por slip un trabajo que ya utiliza $C_\theta$ real.
\end{itemize}
% studium:end-record eyJpZCI6ICJQUkItMDExNCIsICJzZWN0aW9uIjogImMxNSJ9
\subsection*{C15-02-P}

\noindent\textbf{Enunciado}
\(C15-02-P \cdot\) Curvatura hacia atrás. Ejercicio de transferencia. Un rodete tiene ángulo relativo desde la dirección meridiana negativo. Halle el trabajo con deslizamiento definido como cociente de remolinos. Antes de consultar la solución, prediga el efecto de cambiar el primer dato respecto del caso resuelto de esta familia. Justifique la tendencia y explique qué restricción podría convertirse en limitante. Datos: \(U=477.12\, \mathrm{m}\,\mathrm{s}^{-1}\); \(C_{m,s}=139\, \mathrm{m}\,\mathrm{s}^{-1}\); \(\lambda_s=-0.5061455\, \mathrm{rad}\); \(\sigma_s=0.89\).

\noindent\textbf{Datos}
\begin{itemize}
\item $U = 477.12\,\mathrm{m}\,\mathrm{s}^{-1}$ — velocidad periférica salida
\item $C_{m,s} = 139\,\mathrm{m}\,\mathrm{s}^{-1}$ — velocidad meridiana salida
\item $\lambda_s = -0.50614548307836\,\mathrm{rad}$ — ángulo relativo ideal desde meridiana
\item $\sigma_s = 0.89$ — factor de deslizamiento
\end{itemize}

\noindent\textbf{Incógnita}
$w_c$, en $\mathrm{J\,kg^{-1}}$; también las magnitudes intermedias desarrolladas.

\noindent\textbf{Modelo físico}

\noindent\textit{Curvatura hacia atrás}

\[
\fitmath{w_c=\sigma_s U \left(C_{m,s} \tan{\left(\lambda_s \right)} + U\right)}
\]

\noindent\textbf{Hipótesis}
\begin{itemize}
\item Factor prescrito positivo; salida ideal con remolino positivo.
\end{itemize}

\noindent\textbf{Desarrollo paso a paso}

\begin{enumerate}

\item El metal ideal determina la componente tangencial absoluta.

\[\fitmath{\mathrm{ctid}=\mathrm{uu}+\mathrm{cm}\cdot \tan\left(\mathrm{lam}\right)}\]

\[\fitmath{\mathrm{ctid}=(477.12)+(139)\cdot \tan\left((-0.5061454831)\right)\approx 400.071\,\mathrm{m/s}}\]

\item Se aplica el cociente de remolinos definido.

\[\fitmath{\mathrm{ct}=\mathrm{slip}\cdot \left(\mathrm{uu}+\mathrm{cm}\cdot \tan\left(\mathrm{lam}\right)\right)}\]

\[\fitmath{\mathrm{ct}=(0.89)\cdot \left((477.12)+(139)\cdot \tan\left((-0.5061454831)\right)\right)\approx 356.0632\,\mathrm{m/s}}\]

\item Sin prerrotación, Euler queda U2 Ctheta2.

\[\fitmath{\mathrm{wc}=\mathrm{uu}\cdot \mathrm{slip}\cdot \left(\mathrm{uu}+\mathrm{cm}\cdot \tan\left(\mathrm{lam}\right)\right)}\]

\[\fitmath{\mathrm{wc}=(477.12)\cdot (0.89)\cdot \left((477.12)+(139)\cdot \tan\left((-0.5061454831)\right)\right)\approx 169884.9\,\mathrm{J/kg}}\]

\end{enumerate}


\noindent\textbf{Sustitución numérica}

\[
\fitmath{\left(0.89\right) \left(477.12\,\mathrm{m}\,\mathrm{s}^{-1}\right) \left(\left(139\,\mathrm{m}\,\mathrm{s}^{-1}\right) \tan{\left(\left(-0.50614548307836\,\mathrm{rad}\right) \right)} + \left(477.12\,\mathrm{m}\,\mathrm{s}^{-1}\right)\right)}
\]

\noindent\textbf{Resultado}

\[
\fitmath{w_c=1.698848870\times10^{5}\,\mathrm{J}\,\mathrm{kg}^{-1}}
\]

\noindent\textbf{Interpretación}
La curvatura hacia atrás reduce el remolino ideal para estos datos. No intercambie esta definición con $1-dC_\theta/U$.

\noindent\textbf{Comprobación de los pasos}
\begin{itemize}
\item Paso 2: $400.071$ — reproducido
\item Paso 4: $356.063$ — reproducido
\item Paso 6: $169885$ — reproducido
\end{itemize}

\noindent\textbf{Comprobación de la respuesta} la aritmética del último paso reproduce la respuesta enunciada a la precisión con que está escrita. Es una reproducción aritmética; no valida el modelo físico ni sus hipótesis.

\noindent\textbf{Limitaciones}
\begin{itemize}
\item Modelo educativo; la aplicabilidad a una máquina real requiere revisión experta.
\end{itemize}

\noindent\textbf{Errores frecuentes}
\begin{itemize}
\item Usar una cotangente con el mismo ángulo desde meridiana.
\end{itemize}
% studium:end-record eyJpZCI6ICJQUkItMDExNiIsICJzZWN0aW9uIjogImMxNSJ9
\subsection*{C15-03-P}

\noindent\textbf{Enunciado}
\(C15-03-P \cdot\) Inductor y Mach de punta. Ejercicio de transferencia. Un inductor sin prerrotación recibe gas con temperatura estática dada. Halle la velocidad relativa de punta y su Mach. Antes de consultar la solución, prediga el efecto de cambiar el primer dato respecto del caso resuelto de esta familia. Justifique la tendencia y explique qué restricción podría convertirse en limitante. Datos: \(\Omega=2044.8\, \mathrm{s}^{-1}\); \(r_t=0.102\, \mathrm{m}\); \(C_x=148\, \mathrm{m}\,\mathrm{s}^{-1}\); \(T_e=283\, \mathrm{K}\); \(\gamma=1.4\); \(R_g=287\, \mathrm{J}\,\mathrm{kg}^{-1}\,\mathrm{K}^{-1}\).

\noindent\textbf{Datos}
\begin{itemize}
\item $\Omega = 2044.8\,\mathrm{s}^{-1}$ — velocidad angular
\item $r_t = 0.102\,\mathrm{m}$ — radio punta inductor
\item $C_x = 148\,\mathrm{m}\,\mathrm{s}^{-1}$ — velocidad axial
\item $T_e = 283\,\mathrm{K}$ — temperatura estática
\item $\gamma = 1.4$ — razón de calores
\item $R_g = 287\,\mathrm{J}\,\mathrm{kg}^{-1}\,\mathrm{K}^{-1}$ — constante gas
\end{itemize}

\noindent\textbf{Incógnita}
$\mathrm{Ma}_r$ (adimensional); también las magnitudes intermedias desarrolladas.

\noindent\textbf{Modelo físico}

\noindent\textit{Inductor y Mach de punta}

\[
\fitmath{\mathrm{Ma}_r=\frac{\sqrt{\Omega^{2} r_t^{2} + C_x^{2}}}{\sqrt{R_g T_e \gamma}}}
\]

\noindent\textbf{Hipótesis}
\begin{itemize}
\item Sin prerrotación y velocidad axial uniforme.
\end{itemize}

\noindent\textbf{Desarrollo paso a paso}

\begin{enumerate}

\item La punta tiene mayor velocidad periférica que la raíz.

\[\fitmath{\mathrm{Ut}=\Omega\cdot \mathrm{rt}}\]

\[\fitmath{\mathrm{Ut}=(2044.8)\cdot (0.102)\approx 208.5696\,\mathrm{m/s}}\]

\item Se forma el módulo relativo de entrada.

\[\fitmath{W=\sqrt{{\mathrm{cx}}^{2}+{\left(\Omega\cdot \mathrm{rt}\right)}^{2}}}\]

\[\fitmath{W=\sqrt{{(148)}^{2}+{\left((2044.8)\cdot (0.102)\right)}^{2}}\approx 255.7446\,\mathrm{m/s}}\]

\item El sonido corresponde al estado estático.

\[\fitmath{\mathrm{Mr}=\frac{\sqrt{{\mathrm{cx}}^{2}+{\left(\Omega\cdot \mathrm{rt}\right)}^{2}}}{\sqrt{\gamma\cdot R_g\cdot \mathrm{Te}}}}\]

\[\fitmath{\mathrm{Mr}=\frac{\sqrt{{(148)}^{2}+{\left((2044.8)\cdot (0.102)\right)}^{2}}}{\sqrt{(1.4)\cdot (287)\cdot (283)}}\approx 0.7584173\,}\]

\end{enumerate}


\noindent\textbf{Sustitución numérica}

\[
\fitmath{\frac{\sqrt{\left(\left(2044.8\,\mathrm{s}^{-1}\right)\right)^{2} \left(0.102\,\mathrm{m}\right)^{2} + \left(\left(148\,\mathrm{m}\,\mathrm{s}^{-1}\right)\right)^{2}}}{\sqrt{\left(287\,\mathrm{J}\,\mathrm{kg}^{-1}\,\mathrm{K}^{-1}\right) \left(283\,\mathrm{K}\right) \left(1.4\right)}}}
\]

\noindent\textbf{Resultado}

\[
\fitmath{\mathrm{Ma}_r=7.584172706\times10^{-1}}
\]

\noindent\textbf{Interpretación}
La temperatura estática se comparte entre referencias; la velocidad relativa aumenta por la rotación.

\noindent\textbf{Comprobación de los pasos}
\begin{itemize}
\item Paso 2: $208.57$ — reproducido
\item Paso 4: $255.745$ — reproducido
\item Paso 6: $0.758417$ — reproducido
\end{itemize}

\noindent\textbf{Comprobación de la respuesta} la aritmética del último paso reproduce la respuesta enunciada a la precisión con que está escrita. Es una reproducción aritmética; no valida el modelo físico ni sus hipótesis.

\noindent\textbf{Limitaciones}
\begin{itemize}
\item Modelo educativo; la aplicabilidad a una máquina real requiere revisión experta.
\end{itemize}

\noindent\textbf{Errores frecuentes}
\begin{itemize}
\item Usar radio exterior del rodete en el ojo del inductor.
\end{itemize}
% studium:end-record eyJpZCI6ICJQUkItMDExOCIsICJzZWN0aW9uIjogImMxNSJ9
\subsection*{C15-04-P}

\noindent\textbf{Enunciado}
\(C15-04-P \cdot\) Prerrotación y trabajo. Ejercicio de transferencia. Manteniendo la descarga real, determine cuánto trabajo aporta el rodete con prerrotación positiva. Antes de consultar la solución, prediga el efecto de cambiar el primer dato respecto del caso resuelto de esta familia. Justifique la tendencia y explique qué restricción podría convertirse en limitante. Datos: \(\Omega=2215.2\, \mathrm{s}^{-1}\); \(r_e=0.079\, \mathrm{m}\); \(r_s=0.204\, \mathrm{m}\); \(C_{\theta e}=47\, \mathrm{m}\,\mathrm{s}^{-1}\); \(C_{\theta s}=351\, \mathrm{m}\,\mathrm{s}^{-1}\).

\noindent\textbf{Datos}
\begin{itemize}
\item $\Omega = 2215.2\,\mathrm{s}^{-1}$ — velocidad angular
\item $r_e = 0.079\,\mathrm{m}$ — radio entrada
\item $r_s = 0.204\,\mathrm{m}$ — radio salida
\item $C_{\theta e} = 47\,\mathrm{m}\,\mathrm{s}^{-1}$ — remolino entrada
\item $C_{\theta s} = 351\,\mathrm{m}\,\mathrm{s}^{-1}$ — remolino real salida
\end{itemize}

\noindent\textbf{Incógnita}
$w_c$, en $\mathrm{J\,kg^{-1}}$; también las magnitudes intermedias desarrolladas.

\noindent\textbf{Modelo físico}

\noindent\textit{Prerrotación y trabajo}

\[
\fitmath{w_c=\Omega \left(- C_{\theta e} r_e + C_{\theta s} r_s\right)}
\]

\noindent\textbf{Hipótesis}
\begin{itemize}
\item Euler general con valores medios compatibles.
\end{itemize}

\noindent\textbf{Desarrollo paso a paso}

\begin{enumerate}

\item Se calcula la contribución de descarga.

\[\fitmath{\mathrm{exit}=\Omega\cdot \mathrm{rs}\cdot \mathrm{cs}}\]

\[\fitmath{\mathrm{exit}=(2215.2)\cdot (0.204)\cdot (351)\approx 158617.2\,\mathrm{J/kg}}\]

\item Se resta la contribución de entrada.

\[\fitmath{\mathrm{wc}=\Omega\cdot \left(\mathrm{rs}\cdot \mathrm{cs}-\mathrm{re}\cdot \mathrm{ce}\right)}\]

\[\fitmath{\mathrm{wc}=(2215.2)\cdot \left((0.204)\cdot (351)-(0.079)\cdot (47)\right)\approx 150392.1\,\mathrm{J/kg}}\]

\end{enumerate}


\noindent\textbf{Sustitución numérica}

\[
\fitmath{\left(2215.2\,\mathrm{s}^{-1}\right) \left(- \left(47\,\mathrm{m}\,\mathrm{s}^{-1}\right) \left(0.079\,\mathrm{m}\right) + \left(351\,\mathrm{m}\,\mathrm{s}^{-1}\right) \left(0.204\,\mathrm{m}\right)\right)}
\]

\noindent\textbf{Resultado}

\[
\fitmath{w_c=1.503921432\times10^{5}\,\mathrm{J}\,\mathrm{kg}^{-1}}
\]

\noindent\textbf{Interpretación}
La prerrotación positiva reduce trabajo para descarga fijada. Su efecto en incidencia debe evaluarse aparte.

\noindent\textbf{Comprobación de los pasos}
\begin{itemize}
\item Paso 2: $158617$ — reproducido
\item Paso 4: $150392$ — reproducido
\end{itemize}

\noindent\textbf{Comprobación de la respuesta} la aritmética del último paso reproduce la respuesta enunciada a la precisión con que está escrita. Es una reproducción aritmética; no valida el modelo físico ni sus hipótesis.

\noindent\textbf{Limitaciones}
\begin{itemize}
\item Modelo educativo; la aplicabilidad a una máquina real requiere revisión experta.
\end{itemize}

\noindent\textbf{Errores frecuentes}
\begin{itemize}
\item Suponer prerrotación nula aunque el enunciado la proporcione.
\end{itemize}
% studium:end-record eyJpZCI6ICJQUkItMDEyMCIsICJzZWN0aW9uIjogImMxNSJ9
\section{Capítulo 16: Difusores radiales, recuperación y geometría de paso}
\subsection*{C16-01-P}

\noindent\textbf{Enunciado}
\(C16-01-P \cdot\) Difusor sin álabes ideal. Ejercicio de transferencia. En un difusor de ancho y densidad constantes, calcule las componentes de salida y el incremento ideal de presión. Antes de consultar la solución, prediga el efecto de cambiar el primer dato respecto del caso resuelto de esta familia. Justifique la tendencia y explique qué restricción podría convertirse en limitante. Datos: \(r_e=0.22365\, \mathrm{m}\); \(r_s=0.32\, \mathrm{m}\); \(C_r=128\, \mathrm{m}\,\mathrm{s}^{-1}\); \(C_\theta=337\, \mathrm{m}\,\mathrm{s}^{-1}\); \(\rho=2.9\, \mathrm{kg}\,\mathrm{m}^{-3}\).

\noindent\textbf{Datos}
\begin{itemize}
\item $r_e = 0.22365\,\mathrm{m}$ — radio entrada
\item $r_s = 0.32\,\mathrm{m}$ — radio salida
\item $C_r = 128\,\mathrm{m}\,\mathrm{s}^{-1}$ — componente radial entrada
\item $C_\theta = 337\,\mathrm{m}\,\mathrm{s}^{-1}$ — componente tangencial entrada
\item $\rho = 2.9\,\mathrm{kg}\,\mathrm{m}^{-3}$ — densidad uniforme
\end{itemize}

\noindent\textbf{Incógnita}
$\Delta p$, en $\mathrm{Pa}$; también las magnitudes intermedias desarrolladas.

\noindent\textbf{Modelo físico}

\noindent\textit{Difusor sin álabes ideal}

\[
\fitmath{\Delta p=\frac{\rho \left(C_r^{2} + C_\theta^{2}\right) \left(- \frac{r_e^{2}}{r_s^{2}} + 1\right)}{2}}
\]

\noindent\textbf{Hipótesis}
\begin{itemize}
\item Adiabático ideal, ancho y densidad constantes, sin par.
\end{itemize}

\noindent\textbf{Desarrollo paso a paso}

\begin{enumerate}

\item Componente radial de salida

\[\fitmath{\mathrm{crs}=\frac{\mathrm{cr}\cdot \mathrm{re}}{\mathrm{rs}}}\]

\[\fitmath{\mathrm{crs}=\frac{(128)\cdot (0.22365)}{(0.32)}\approx 89.46\,\mathrm{m/s}}\]

\item Continuidad y ausencia de par hacen escalar ambas componentes con r inverso.

\[\fitmath{\mathrm{cts}=\frac{\mathrm{ct}\cdot \mathrm{re}}{\mathrm{rs}}}\]

\[\fitmath{\mathrm{cts}=\frac{(337)\cdot (0.22365)}{(0.32)}\approx 235.5314\,\mathrm{m/s}}\]

\item Se evalúa el incremento de presión por deceleración ideal.

\[\fitmath{\mathrm{dp}=\frac{\rho\cdot \left({\mathrm{cr}}^{2}+{\mathrm{ct}}^{2}\right)\cdot \left(1-{\left(\frac{\mathrm{re}}{\mathrm{rs}}\right)}^{2}\right)}{2}}\]

\[\fitmath{\mathrm{dp}=\frac{(2.9)\cdot \left({(128)}^{2}+{(337)}^{2}\right)\cdot \left(1-{\left(\frac{(0.22365)}{(0.32)}\right)}^{2}\right)}{2}\approx 96388.55\,\mathrm{Pa}}\]

\end{enumerate}


\noindent\textbf{Sustitución numérica}

\[
\fitmath{\frac{\left(2.9\,\mathrm{kg}\,\mathrm{m}^{-3}\right) \left(\left(\left(128\,\mathrm{m}\,\mathrm{s}^{-1}\right)\right)^{2} + \left(\left(337\,\mathrm{m}\,\mathrm{s}^{-1}\right)\right)^{2}\right) \left(- \frac{\left(0.22365\,\mathrm{m}\right)^{2}}{\left(0.32\,\mathrm{m}\right)^{2}} + 1\right)}{2}}
\]

\noindent\textbf{Resultado}

\[
\fitmath{\Delta p=9.638855435\times10^{4}\,\mathrm{Pa}}
\]

\noindent\textbf{Componente radial de salida}
\[
\fitmath{C_{rs}=C_{re}r_e/r_s=89.46\,\mathrm{m/s}}
\]

\noindent\textbf{Comprobación de continuidad.} $r_sC_{rs}=r_eC_{re}$ con ancho y densidad constantes; junto con $r_sC_{\theta s}=r_eC_{\theta e}$ reproduce la recuperación ideal de presión.

\noindent\textbf{Interpretación}
La presión aumenta y el momento angular se conserva. Es una integración incompresible ideal, no una curva ensayada.

\noindent\textbf{Comprobación de los pasos}
\begin{itemize}
\item Paso 2: $235.531$ — reproducido
\item Paso 4: $96388.6$ — reproducido
\end{itemize}

\noindent\textbf{Comprobación de la respuesta} la aritmética del último paso reproduce la respuesta enunciada a la precisión con que está escrita. Es una reproducción aritmética; no valida el modelo físico ni sus hipótesis.

\noindent\textbf{Limitaciones}
\begin{itemize}
\item Modelo educativo; la aplicabilidad a una máquina real requiere revisión experta.
\end{itemize}

\noindent\textbf{Errores frecuentes}
\begin{itemize}
\item Conservar $C_\theta$ en lugar de $rC_\theta$.
\end{itemize}
% studium:end-record eyJpZCI6ICJQUkItMDEyMiIsICJzZWN0aW9uIjogImMxNiJ9
\subsection*{C16-02-P}

\noindent\textbf{Enunciado}
\(C16-02-P \cdot\) Recuperación con pérdida prescrita. Ejercicio de transferencia. Calcule el coeficiente de recuperación estática y el incremento de presión con una pérdida normalizada dada. Antes de consultar la solución, prediga el efecto de cambiar el primer dato respecto del caso resuelto de esta familia. Justifique la tendencia y explique qué restricción podría convertirse en limitante. Datos: \(C_e=387.66\, \mathrm{m}\,\mathrm{s}^{-1}\); \(C_s=216\, \mathrm{m}\,\mathrm{s}^{-1}\); \(\zeta_D=0.13\); \(\rho=3.1\, \mathrm{kg}\,\mathrm{m}^{-3}\).

\noindent\textbf{Datos}
\begin{itemize}
\item $C_e = 387.66\,\mathrm{m}\,\mathrm{s}^{-1}$ — velocidad entrada
\item $C_s = 216\,\mathrm{m}\,\mathrm{s}^{-1}$ — velocidad salida
\item $\zeta_D = 0.13$ — pérdida normalizada con dinámica de entrada
\item $\rho = 3.1\,\mathrm{kg}\,\mathrm{m}^{-3}$ — densidad
\end{itemize}

\noindent\textbf{Incógnita}
$\Delta p$, en $\mathrm{Pa}$; también las magnitudes intermedias desarrolladas.

\noindent\textbf{Modelo físico}

\noindent\textit{Recuperación con pérdida prescrita}

\[
\fitmath{\Delta p=\frac{\rho C_e^{2} \left(- \zeta_D + 1 - \frac{C_s^{2}}{C_e^{2}}\right)}{2}}
\]

\noindent\textbf{Hipótesis}
\begin{itemize}
\item Modelo incompresible, pérdida definida con velocidad entrada.
\end{itemize}

\noindent\textbf{Desarrollo paso a paso}

\begin{enumerate}

\item La recuperación resta energía residual y pérdida.

\[\fitmath{\mathrm{Cp}=1-{\left(\frac{\mathrm{vs}}{\mathrm{ve}}\right)}^{2}-\mathrm{loss}}\]

\[\fitmath{\mathrm{Cp}=1-{\left(\frac{(216)}{(387.66)}\right)}^{2}-(0.13)\approx 0.5595401\,}\]

\item Se dimensionaliza con la presión dinámica de entrada.

\[\fitmath{\mathrm{dp}=\frac{\rho\cdot {\mathrm{ve}}^{2}\cdot \left(1-{\left(\frac{\mathrm{vs}}{\mathrm{ve}}\right)}^{2}-\mathrm{loss}\right)}{2}}\]

\[\fitmath{\mathrm{dp}=\frac{(3.1)\cdot {(387.66)}^{2}\cdot \left(1-{\left(\frac{(216)}{(387.66)}\right)}^{2}-(0.13)\right)}{2}\approx 130336.2\,\mathrm{Pa}}\]

\end{enumerate}


\noindent\textbf{Sustitución numérica}

\[
\fitmath{\frac{\left(3.1\,\mathrm{kg}\,\mathrm{m}^{-3}\right) \left(\left(387.66\,\mathrm{m}\,\mathrm{s}^{-1}\right)\right)^{2} \left(- \left(0.13\right) + 1 - \frac{\left(\left(216\,\mathrm{m}\,\mathrm{s}^{-1}\right)\right)^{2}}{\left(\left(387.66\,\mathrm{m}\,\mathrm{s}^{-1}\right)\right)^{2}}\right)}{2}}
\]

\noindent\textbf{Resultado}

\[
\fitmath{\Delta p=1.303361516\times10^{5}\,\mathrm{Pa}}
\]

\noindent\textbf{Interpretación}
La pérdida total y la velocidad residual reducen recuperación por mecanismos distintos.

\noindent\textbf{Comprobación de los pasos}
\begin{itemize}
\item Paso 2: $0.55954$ — reproducido
\item Paso 4: $130336$ — reproducido
\end{itemize}

\noindent\textbf{Comprobación de la respuesta} la aritmética del último paso reproduce la respuesta enunciada a la precisión con que está escrita. Es una reproducción aritmética; no valida el modelo físico ni sus hipótesis.

\noindent\textbf{Limitaciones}
\begin{itemize}
\item Modelo educativo; la aplicabilidad a una máquina real requiere revisión experta.
\end{itemize}

\noindent\textbf{Errores frecuentes}
\begin{itemize}
\item Restar una pérdida normalizada con otro denominador.
\end{itemize}
% studium:end-record eyJpZCI6ICJQUkItMDEyNCIsICJzZWN0aW9uIjogImMxNiJ9
\subsection*{C16-03-P}

\noindent\textbf{Enunciado}
\(C16-03-P \cdot\) Ancho del paso cilíndrico. Ejercicio de transferencia. Dimensione el ancho geométrico de salida radial con el caudal y área disponible dados. Antes de consultar la solución, prediga el efecto de cambiar el primer dato respecto del caso resuelto de esta familia. Justifique la tendencia y explique qué restricción podría convertirse en limitante. Datos: \(\dot m=9.2655\, \mathrm{kg}\,\mathrm{s}^{-1}\); \(r=0.198\, \mathrm{m}\); \(\rho=3.4\, \mathrm{kg}\,\mathrm{m}^{-3}\); \(C_r=126\, \mathrm{m}\,\mathrm{s}^{-1}\); \(B=0.94\).

\noindent\textbf{Datos}
\begin{itemize}
\item $\dot m = 9.2655\,\mathrm{kg}\,\mathrm{s}^{-1}$ — caudal
\item $r = 0.198\,\mathrm{m}$ — radio salida
\item $\rho = 3.4\,\mathrm{kg}\,\mathrm{m}^{-3}$ — densidad estática
\item $C_r = 126\,\mathrm{m}\,\mathrm{s}^{-1}$ — velocidad radial
\item $B = 0.94$ — fracción área disponible
\end{itemize}

\noindent\textbf{Incógnita}
$b$, en $\mathrm{m}$; también las magnitudes intermedias desarrolladas.

\noindent\textbf{Modelo físico}

\noindent\textit{Ancho del paso cilíndrico}

\[
\fitmath{b=\frac{\dot m}{2 \pi B C_r r \rho}}
\]

\noindent\textbf{Hipótesis}
\begin{itemize}
\item Paso cilíndrico uniforme y bloqueo explícito.
\end{itemize}

\noindent\textbf{Desarrollo paso a paso}

\begin{enumerate}

\item La continuidad fija el área cilíndrica efectiva.

\[\fitmath{\mathrm{ae}=\frac{\mathrm{md}}{\rho\cdot \mathrm{cr}}}\]

\[\fitmath{\mathrm{ae}=\frac{(9.2655)}{(3.4)\cdot (126)}\approx 0.02162815\,\mathrm{m^{2}}}\]

\item La circunferencia y el bloqueo fijan el ancho.

\[\fitmath{\mathrm{width}=\frac{\mathrm{md}}{2\cdot \pi\cdot \mathrm{rad}\cdot \rho\cdot \mathrm{cr}\cdot \mathrm{bb}}}\]

\[\fitmath{\mathrm{width}=\frac{(9.2655)}{2\cdot \pi\cdot (0.198)\cdot (3.4)\cdot (126)\cdot (0.94)}\approx 0.01849467\,\mathrm{m}}\]

\end{enumerate}


\noindent\textbf{Sustitución numérica}

\[
\fitmath{\frac{\left(9.2655\,\mathrm{kg}\,\mathrm{s}^{-1}\right)}{2 \pi \left(0.94\right) \left(126\,\mathrm{m}\,\mathrm{s}^{-1}\right) \left(0.198\,\mathrm{m}\right) \left(3.4\,\mathrm{kg}\,\mathrm{m}^{-3}\right)}}
\]

\noindent\textbf{Resultado}

\[
\fitmath{b=1.849466572\times10^{-2}\,\mathrm{m}}
\]

\noindent\textbf{Interpretación}
Se usa componente radial, aunque el módulo absoluto sea mucho mayor por remolino.

\noindent\textbf{Comprobación de los pasos}
\begin{itemize}
\item Paso 2: $0.0216282$ — reproducido
\item Paso 4: $0.0184947$ — reproducido
\end{itemize}

\noindent\textbf{Comprobación de la respuesta} la aritmética del último paso reproduce la respuesta enunciada a la precisión con que está escrita. Es una reproducción aritmética; no valida el modelo físico ni sus hipótesis.

\noindent\textbf{Limitaciones}
\begin{itemize}
\item Modelo educativo; la aplicabilidad a una máquina real requiere revisión experta.
\end{itemize}

\noindent\textbf{Errores frecuentes}
\begin{itemize}
\item Usar el módulo absoluto en la continuidad radial.
\end{itemize}
% studium:end-record eyJpZCI6ICJQUkItMDEyNiIsICJzZWN0aW9uIjogImMxNiJ9
\subsection*{C16-04-P}

\noindent\textbf{Enunciado}
\(C16-04-P \cdot\) Estado estático de descarga. Ejercicio de transferencia. La descarga del rodete tiene estado total y componentes radial y tangencial dadas. Halle temperatura estática y densidad para dimensionar el difusor. Antes de consultar la solución, prediga el efecto de cambiar el primer dato respecto del caso resuelto de esta familia. Justifique la tendencia y explique qué restricción podría convertirse en limitante. Datos: \(T_0=502.68\, \mathrm{K}\); \(p_0=342000\, \mathrm{Pa}\); \(C_r=141\, \mathrm{m}\,\mathrm{s}^{-1}\); \(C_\theta=328\, \mathrm{m}\,\mathrm{s}^{-1}\); \(c_p=1004.5\, \mathrm{J}\,\mathrm{kg}^{-1}\,\mathrm{K}^{-1}\); \(\gamma=1.4\); \(R_g=287\, \mathrm{J}\,\mathrm{kg}^{-1}\,\mathrm{K}^{-1}\).

\noindent\textbf{Datos}
\begin{itemize}
\item $T_0 = 502.68\,\mathrm{K}$ — temperatura total
\item $p_0 = 342000\,\mathrm{Pa}$ — presión total
\item $C_r = 141\,\mathrm{m}\,\mathrm{s}^{-1}$ — componente radial
\item $C_\theta = 328\,\mathrm{m}\,\mathrm{s}^{-1}$ — componente tangencial
\item $c_p = 1004.5\,\mathrm{J}\,\mathrm{kg}^{-1}\,\mathrm{K}^{-1}$ — calor específico
\item $\gamma = 1.4$ — razón de calores
\item $R_g = 287\,\mathrm{J}\,\mathrm{kg}^{-1}\,\mathrm{K}^{-1}$ — constante gas
\end{itemize}

\noindent\textbf{Incógnita}
$\rho$, en $\mathrm{kg\,m^{-3}}$; también las magnitudes intermedias desarrolladas.

\noindent\textbf{Modelo físico}

\noindent\textit{Estado estático de descarga}

\[
\fitmath{\rho=\frac{p_0 \left(\frac{T_0 - \frac{C_r^{2} + C_\theta^{2}}{2 c_p}}{T_0}\right)^{\frac{\gamma}{\gamma - 1}}}{R_g \left(T_0 - \frac{C_r^{2} + C_\theta^{2}}{2 c_p}\right)}}
\]

\noindent\textbf{Hipótesis}
\begin{itemize}
\item Gas perfecto constante; componentes ortogonales.
\end{itemize}

\noindent\textbf{Desarrollo paso a paso}

\begin{enumerate}

\item Ambas componentes contribuyen a energía cinética.

\[\fitmath{T=T_0-\frac{{\mathrm{cr}}^{2}+{\mathrm{ct}}^{2}}{2\cdot \mathrm{cp}}}\]

\[\fitmath{T=(502.68)-\frac{{(141)}^{2}+{(328)}^{2}}{2\cdot (1004.5)}\approx 439.233\,\mathrm{K}}\]

\item La presión estática deriva del remanso local.

\[\fitmath{p=p_0\cdot {\left(\frac{T_0-\frac{{\mathrm{cr}}^{2}+{\mathrm{ct}}^{2}}{2\cdot \mathrm{cp}}}{T_0}\right)}^{\frac{\gamma}{\gamma-1}}}\]

\[\fitmath{p=(342000)\cdot {\left(\frac{(502.68)-\frac{{(141)}^{2}+{(328)}^{2}}{2\cdot (1004.5)}}{(502.68)}\right)}^{\frac{(1.4)}{(1.4)-1}}\approx 213274\,\mathrm{Pa}}\]

\item Se calcula densidad a partir del mismo estado estático.

\[\fitmath{\rho=\frac{p_0\cdot {\left(\frac{T_0-\frac{{\mathrm{cr}}^{2}+{\mathrm{ct}}^{2}}{2\cdot \mathrm{cp}}}{T_0}\right)}^{\frac{\gamma}{\gamma-1}}}{R_g\cdot \left(T_0-\frac{{\mathrm{cr}}^{2}+{\mathrm{ct}}^{2}}{2\cdot \mathrm{cp}}\right)}}\]

\[\fitmath{\rho=\frac{(342000)\cdot {\left(\frac{(502.68)-\frac{{(141)}^{2}+{(328)}^{2}}{2\cdot (1004.5)}}{(502.68)}\right)}^{\frac{(1.4)}{(1.4)-1}}}{(287)\cdot \left((502.68)-\frac{{(141)}^{2}+{(328)}^{2}}{2\cdot (1004.5)}\right)}\approx 1.691847\,\mathrm{kg/m^{3}}}\]

\end{enumerate}


\noindent\textbf{Sustitución numérica}

\[
\fitmath{\frac{\left(342000\,\mathrm{Pa}\right) \left(\frac{\left(502.68\,\mathrm{K}\right) - \frac{\left(\left(141\,\mathrm{m}\,\mathrm{s}^{-1}\right)\right)^{2} + \left(\left(328\,\mathrm{m}\,\mathrm{s}^{-1}\right)\right)^{2}}{2 \left(1004.5\,\mathrm{J}\,\mathrm{kg}^{-1}\,\mathrm{K}^{-1}\right)}}{\left(502.68\,\mathrm{K}\right)}\right)^{\frac{\left(1.4\right)}{\left(1.4\right) - 1}}}{\left(287\,\mathrm{J}\,\mathrm{kg}^{-1}\,\mathrm{K}^{-1}\right) \left(\left(502.68\,\mathrm{K}\right) - \frac{\left(\left(141\,\mathrm{m}\,\mathrm{s}^{-1}\right)\right)^{2} + \left(\left(328\,\mathrm{m}\,\mathrm{s}^{-1}\right)\right)^{2}}{2 \left(1004.5\,\mathrm{J}\,\mathrm{kg}^{-1}\,\mathrm{K}^{-1}\right)}\right)}}
\]

\noindent\textbf{Resultado}

\[
\fitmath{\rho=1.691846742\times10^{0}\,\mathrm{kg}\,\mathrm{m}^{-3}}
\]

\noindent\textbf{Interpretación}
Usar solo Cr para temperatura sobreestima la densidad disponible. La continuidad sí utiliza Cr.

\noindent\textbf{Comprobación de los pasos}
\begin{itemize}
\item Paso 2: $439.233$ — reproducido
\item Paso 4: $213274$ — reproducido
\item Paso 6: $1.69185$ — reproducido
\end{itemize}

\noindent\textbf{Comprobación de la respuesta} la aritmética del último paso reproduce la respuesta enunciada a la precisión con que está escrita. Es una reproducción aritmética; no valida el modelo físico ni sus hipótesis.

\noindent\textbf{Limitaciones}
\begin{itemize}
\item Modelo educativo; la aplicabilidad a una máquina real requiere revisión experta.
\end{itemize}

\noindent\textbf{Errores frecuentes}
\begin{itemize}
\item Usar distintas definiciones de velocidad en energía sin justificarlo.
\end{itemize}
% studium:end-record eyJpZCI6ICJQUkItMDEyOCIsICJzZWN0aW9uIjogImMxNiJ9
\section{Capítulo 17: Prestaciones y límites de un compresor radial}
\subsection*{C17-01-P}

\noindent\textbf{Enunciado}
\(C17-01-P \cdot\) Velocidad exigida por trabajo. Ejercicio de transferencia. Un rodete de salida radial debe suministrar trabajo dado con factor de deslizamiento prescrito. Calcule la velocidad periférica y el radio para el giro indicado. Antes de consultar la solución, prediga el efecto de cambiar el primer dato respecto del caso resuelto de esta familia. Justifique la tendencia y explique qué restricción podría convertirse en limitante. Datos: \(w_c=179985\, \mathrm{J}\,\mathrm{kg}^{-1}\); \(\sigma_s=0.86\); \(\Omega=2140\, \mathrm{s}^{-1}\).

\noindent\textbf{Datos}
\begin{itemize}
\item $w_c = 179985\,\mathrm{J}\,\mathrm{kg}^{-1}$ — trabajo requerido
\item $\sigma_s = 0.86$ — factor de deslizamiento
\item $\Omega = 2140\,\mathrm{s}^{-1}$ — velocidad angular
\end{itemize}

\noindent\textbf{Incógnita}
$r$, en $\mathrm{m}$; también las magnitudes intermedias desarrolladas.

\noindent\textbf{Modelo físico}

\noindent\textit{Velocidad exigida por trabajo}

\[
\fitmath{r=\frac{\sqrt{\frac{w_c}{\sigma_s}}}{\Omega}}
\]

\noindent\textbf{Hipótesis}
\begin{itemize}
\item Salida radial, sin prerrotación, slip constante.
\end{itemize}

\noindent\textbf{Desarrollo paso a paso}

\begin{enumerate}

\item La salida radial sin prerrotación da wc=slip U2 al cuadrado.

\[\fitmath{\mathrm{uu}=\sqrt{\frac{\mathrm{wc}}{\mathrm{slip}}}}\]

\[\fitmath{\mathrm{uu}=\sqrt{\frac{(179985)}{(0.86)}}\approx 457.4766\,\mathrm{m/s}}\]

\item La geometría requerida sigue de U=Omega r.

\[\fitmath{\mathrm{rad}=\frac{\sqrt{\frac{\mathrm{wc}}{\mathrm{slip}}}}{\Omega}}\]

\[\fitmath{\mathrm{rad}=\frac{\sqrt{\frac{(179985)}{(0.86)}}}{(2140)}\approx 0.2137741\,\mathrm{m}}\]

\end{enumerate}


\noindent\textbf{Sustitución numérica}

\[
\fitmath{\frac{\sqrt{\frac{\left(179985\,\mathrm{J}\,\mathrm{kg}^{-1}\right)}{\left(0.86\right)}}}{\left(2140\,\mathrm{s}^{-1}\right)}}
\]

\noindent\textbf{Resultado}

\[
\fitmath{r=2.137741347\times10^{-1}\,\mathrm{m}}
\]

\noindent\textbf{Interpretación}
El tamaño y el giro están ligados al trabajo. Las tensiones y el Mach deben comprobarse después.

\noindent\textbf{Comprobación de los pasos}
\begin{itemize}
\item Paso 2: $457.477$ — reproducido
\item Paso 4: $0.213774$ — reproducido
\end{itemize}

\noindent\textbf{Comprobación de la respuesta} la aritmética del último paso reproduce la respuesta enunciada a la precisión con que está escrita. Es una reproducción aritmética; no valida el modelo físico ni sus hipótesis.

\noindent\textbf{Limitaciones}
\begin{itemize}
\item Modelo educativo; la aplicabilidad a una máquina real requiere revisión experta.
\end{itemize}

\noindent\textbf{Errores frecuentes}
\begin{itemize}
\item Sustituir eficiencia isentrópica por factor de deslizamiento.
\end{itemize}
% studium:end-record eyJpZCI6ICJQUkItMDEzMCIsICJzZWN0aW9uIjogImMxNyJ9
\subsection*{C17-02-P}

\noindent\textbf{Enunciado}
\(C17-02-P \cdot\) Presión y sensibilidad de eficiencia. Ejercicio de transferencia. Para trabajo de Euler fijado, calcule la relación de presión que produce el rendimiento suministrado. Antes de consultar la solución, prediga el efecto de cambiar el primer dato respecto del caso resuelto de esta familia. Justifique la tendencia y explique qué restricción podría convertirse en limitante. Datos: \(w_c=189570\, \mathrm{J}\,\mathrm{kg}^{-1}\); \(T_{0e}=297\, \mathrm{K}\); \(\eta=0.81\); \(c_p=1004.5\, \mathrm{J}\,\mathrm{kg}^{-1}\,\mathrm{K}^{-1}\); \(\gamma=1.4\).

\noindent\textbf{Datos}
\begin{itemize}
\item $w_c = 189570\,\mathrm{J}\,\mathrm{kg}^{-1}$ — trabajo de Euler comunicado
\item $T_{0e} = 297\,\mathrm{K}$ — temperatura total entrada
\item $\eta = 0.81$ — eficiencia total-total
\item $c_p = 1004.5\,\mathrm{J}\,\mathrm{kg}^{-1}\,\mathrm{K}^{-1}$ — calor específico
\item $\gamma = 1.4$ — razón de calores
\end{itemize}

\noindent\textbf{Incógnita}
$\pi_c$ (adimensional); también las magnitudes intermedias desarrolladas.

\noindent\textbf{Modelo físico}

\noindent\textit{Presión y sensibilidad de eficiencia}

\[
\fitmath{\pi_c=\left(1 + \frac{\eta w_c}{T_{0e} c_p}\right)^{\frac{\gamma}{\gamma - 1}}}
\]

\noindent\textbf{Hipótesis}
\begin{itemize}
\item Gas perfecto y trabajo fijado, sin refrigeración.
\end{itemize}

\noindent\textbf{Desarrollo paso a paso}

\begin{enumerate}

\item El trabajo real fija incremento total independientemente de eta.

\[\fitmath{\mathrm{dt}=\frac{\mathrm{wc}}{\mathrm{cp}}}\]

\[\fitmath{\mathrm{dt}=\frac{(189570)}{(1004.5)}\approx 188.7208\,\mathrm{K}}\]

\item La eficiencia fija el incremento isentrópico equivalente.

\[\fitmath{\mathrm{pr}={\left(1+\frac{\eta\cdot \mathrm{wc}}{\mathrm{cp}\cdot \mathrm{Tin}}\right)}^{\frac{\gamma}{\gamma-1}}}\]

\[\fitmath{\mathrm{pr}={\left(1+\frac{(0.81)\cdot (189570)}{(1004.5)\cdot (297)}\right)}^{\frac{(1.4)}{(1.4)-1}}\approx 4.276969\,}\]

\end{enumerate}


\noindent\textbf{Sustitución numérica}

\[
\fitmath{\left(1 + \frac{\left(0.81\right) \left(189570\,\mathrm{J}\,\mathrm{kg}^{-1}\right)}{\left(297\,\mathrm{K}\right) \left(1004.5\,\mathrm{J}\,\mathrm{kg}^{-1}\,\mathrm{K}^{-1}\right)}\right)^{\frac{\left(1.4\right)}{\left(1.4\right) - 1}}}
\]

\noindent\textbf{Resultado}

\[
\fitmath{\pi_c=4.276969338\times10^{0}}
\]

\noindent\textbf{Interpretación}
Variar eta a trabajo fijo cambia la presión, mientras el aumento real de temperatura permanece fijo.

\noindent\textbf{Comprobación de los pasos}
\begin{itemize}
\item Paso 2: $188.721$ — reproducido
\item Paso 4: $4.27697$ — reproducido
\end{itemize}

\noindent\textbf{Comprobación de la respuesta} la aritmética del último paso reproduce la respuesta enunciada a la precisión con que está escrita. Es una reproducción aritmética; no valida el modelo físico ni sus hipótesis.

\noindent\textbf{Limitaciones}
\begin{itemize}
\item Modelo educativo; la aplicabilidad a una máquina real requiere revisión experta.
\end{itemize}

\noindent\textbf{Errores frecuentes}
\begin{itemize}
\item Reducir el calentamiento real multiplicando por eficiencia.
\end{itemize}
% studium:end-record eyJpZCI6ICJQUkItMDEzMiIsICJzZWN0aW9uIjogImMxNyJ9
\subsection*{C17-03-P}

\noindent\textbf{Enunciado}
\(C17-03-P \cdot\) Pendiente cinemática con caudal. Ejercicio de transferencia. Un modelo de rodete mantiene densidad, ancho, slip y giro constantes. Calcule la derivada del trabajo respecto al caudal. Antes de consultar la solución, prediga el efecto de cambiar el primer dato respecto del caso resuelto de esta familia. Justifique la tendencia y explique qué restricción podría convertirse en limitante. Datos: \(U=462.21\, \mathrm{m}\,\mathrm{s}^{-1}\); \(\sigma_s=0.88\); \(\lambda_s=-0.5410521\, \mathrm{rad}\); \(r=0.203\, \mathrm{m}\); \(b=0.017\, \mathrm{m}\); \(\rho=3.2\, \mathrm{kg}\,\mathrm{m}^{-3}\).

\noindent\textbf{Datos}
\begin{itemize}
\item $U = 462.21\,\mathrm{m}\,\mathrm{s}^{-1}$ — velocidad periférica
\item $\sigma_s = 0.88$ — factor de deslizamiento
\item $\lambda_s = -0.54105206811824\,\mathrm{rad}$ — ángulo ideal desde meridiana
\item $r = 0.203\,\mathrm{m}$ — radio salida
\item $b = 0.017\,\mathrm{m}$ — ancho salida
\item $\rho = 3.2\,\mathrm{kg}\,\mathrm{m}^{-3}$ — densidad salida
\end{itemize}

\noindent\textbf{Incógnita}
$\mathrm{d}w_c/\mathrm{d}\dot{m}$, en $\mathrm{J\,s\,kg^{-2}}$; también las magnitudes intermedias desarrolladas.

\noindent\textbf{Modelo físico}

\noindent\textit{Pendiente cinemática con caudal}

\[
\fitmath{\frac{\mathrm{d}w_c}{\mathrm{d}\dot{m}}=\frac{\sigma_s U \tan{\left(\lambda_s \right)}}{2 \pi r \rho b}}
\]

\noindent\textbf{Hipótesis}
\begin{itemize}
\item Densidad, ancho, slip, radio y giro fijos; sin prerrotación.
\end{itemize}

\noindent\textbf{Desarrollo paso a paso}

\begin{enumerate}

\item La continuidad da la derivada de velocidad meridiana por caudal.

\[\fitmath{\mathrm{dcm}=\frac{1}{2\cdot \pi\cdot \mathrm{rad}\cdot \mathrm{width}\cdot \rho}}\]

\[\fitmath{\mathrm{dcm}=\frac{1}{2\cdot \pi\cdot (0.203)\cdot (0.017)\cdot (3.2)}\approx 14.41203\,\mathrm{m/kg}}\]

\item La derivada del modelo de Euler incorpora el ángulo.

\[\fitmath{\mathrm{slope}=\frac{\mathrm{uu}\cdot \mathrm{slip}\cdot \tan\left(\mathrm{lam}\right)}{2\cdot \pi\cdot \mathrm{rad}\cdot \mathrm{width}\cdot \rho}}\]

\[\fitmath{\mathrm{slope}=\frac{(462.21)\cdot (0.88)\cdot \tan\left((-0.5410520681)\right)}{2\cdot \pi\cdot (0.203)\cdot (0.017)\cdot (3.2)}\approx -3522.256\,\mathrm{J\,s/kg^{2}}}\]

\end{enumerate}


\noindent\textbf{Sustitución numérica}

\[
\fitmath{\frac{\left(0.88\right) \left(462.21\,\mathrm{m}\,\mathrm{s}^{-1}\right) \tan{\left(\left(-0.54105206811824\,\mathrm{rad}\right) \right)}}{2 \pi \left(0.203\,\mathrm{m}\right) \left(3.2\,\mathrm{kg}\,\mathrm{m}^{-3}\right) \left(0.017\,\mathrm{m}\right)}}
\]

\noindent\textbf{Resultado}

\[
\fitmath{\frac{\mathrm{d}w_c}{\mathrm{d}\dot{m}}=-3.522256217\times10^{3}\,\mathrm{J}\,\mathrm{s}\,\mathrm{kg}^{-2}}
\]

\noindent\textbf{Interpretación}
La pendiente negativa proviene de curvatura hacia atrás. No es una predicción de la línea de bombeo.

\noindent\textbf{Comprobación de los pasos}
\begin{itemize}
\item Paso 2: $14.412$ — reproducido
\item Paso 4: $-3522.26$ — reproducido
\end{itemize}

\noindent\textbf{Comprobación de la respuesta} la aritmética del último paso reproduce la respuesta enunciada a la precisión con que está escrita. Es una reproducción aritmética; no valida el modelo físico ni sus hipótesis.

\noindent\textbf{Limitaciones}
\begin{itemize}
\item Modelo educativo; la aplicabilidad a una máquina real requiere revisión experta.
\end{itemize}

\noindent\textbf{Errores frecuentes}
\begin{itemize}
\item Presentar la pendiente de Euler como estabilidad global.
\end{itemize}
% studium:end-record eyJpZCI6ICJQUkItMDEzNCIsICJzZWN0aW9uIjogImMxNyJ9
\subsection*{C17-04-P}

\noindent\textbf{Enunciado}
\(C17-04-P \cdot\) Balance inverso de rendimiento. Ejercicio de transferencia. De un ensayo educativo se conocen aumento total de temperatura y relación de presión. Calcule la eficiencia total-total y juzgue admisibilidad. Antes de consultar la solución, prediga el efecto de cambiar el primer dato respecto del caso resuelto de esta familia. Justifique la tendencia y explique qué restricción podría convertirse en limitante. Datos: \(T_{0e}=322.695\, \mathrm{K}\); \(T_{0s}=469\, \mathrm{K}\); \(\pi_c=3.85\); \(\gamma=1.4\).

\noindent\textbf{Datos}
\begin{itemize}
\item $T_{0e} = 322.695\,\mathrm{K}$ — temperatura total entrada
\item $T_{0s} = 469\,\mathrm{K}$ — temperatura total salida
\item $\pi_c = 3.85$ — relación de presión total
\item $\gamma = 1.4$ — razón de calores
\end{itemize}

\noindent\textbf{Incógnita}
$\eta$ (adimensional); también las magnitudes intermedias desarrolladas.

\noindent\textbf{Modelo físico}

\noindent\textit{Balance inverso de rendimiento}

\[
\fitmath{\eta=\frac{T_{0e} \left(\pi_c^{\frac{\gamma - 1}{\gamma}} - 1\right)}{- T_{0e} + T_{0s}}}
\]

\noindent\textbf{Hipótesis}
\begin{itemize}
\item Gas perfecto constante, sin calor externo; mediciones medias representativas.
\end{itemize}

\noindent\textbf{Desarrollo paso a paso}

\begin{enumerate}

\item Se calcula la temperatura isentrópica equivalente.

\[\fitmath{\mathrm{Tis}=\mathrm{Tin}\cdot {\mathrm{pr}}^{\frac{\gamma-1}{\gamma}}}\]

\[\fitmath{\mathrm{Tis}=(322.695)\cdot {(3.85)}^{\frac{(1.4)-1}{(1.4)}}\approx 474.3149\,\mathrm{K}}\]

\item Se divide aumento ideal por aumento real.

\[\fitmath{\eta=\frac{\mathrm{Tin}\cdot \left({\mathrm{pr}}^{\frac{\gamma-1}{\gamma}}-1\right)}{\mathrm{Tout}-\mathrm{Tin}}}\]

\[\fitmath{\eta=\frac{(322.695)\cdot \left({(3.85)}^{\frac{(1.4)-1}{(1.4)}}-1\right)}{(469)-(322.695)}\approx 1.036327\,}\]

\end{enumerate}


\noindent\textbf{Sustitución numérica}

\[
\fitmath{\frac{\left(322.695\,\mathrm{K}\right) \left(\left(3.85\right)^{\frac{\left(1.4\right) - 1}{\left(1.4\right)}} - 1\right)}{- \left(322.695\,\mathrm{K}\right) + \left(469\,\mathrm{K}\right)}}
\]

\noindent\textbf{Resultado}

\[
\fitmath{\eta=1.036327257\times10^{0}}
\]

\noindent\textbf{Interpretación}
La eficiencia calculada supera uno. Estos datos no son compatibles con el modelo adiabático: revise identificación estática-total, calibración y calor externo. El resultado es un diagnóstico de inconsistencia, no una prestación admisible.

\noindent\textbf{Comprobación de los pasos}
\begin{itemize}
\item Paso 2: $474.315$ — reproducido
\item Paso 4: $1.03633$ — reproducido
\end{itemize}

\noindent\textbf{Comprobación de la respuesta} la aritmética del último paso reproduce la respuesta enunciada a la precisión con que está escrita. Es una reproducción aritmética; no valida el modelo físico ni sus hipótesis.

\noindent\textbf{Limitaciones}
\begin{itemize}
\item Modelo educativo; la aplicabilidad a una máquina real requiere revisión experta.
\end{itemize}

\noindent\textbf{Errores frecuentes}
\begin{itemize}
\item Confundir presión estática medida con total sin corrección.
\end{itemize}
% studium:end-record eyJpZCI6ICJQUkItMDE1MSIsICJzZWN0aW9uIjogImMxNyJ9
\section{Capítulo 18: Práctica reproducible: secciones a distintos radios}
\subsection*{C18-01-P}

\noindent\textbf{Enunciado}
\(C18-01-P \cdot\) Rotor compresor a radio elegido. Ejercicio de transferencia. Aplique la plantilla radial educativa a la entrada de un rotor compresor. Calcule U, beta y ángulo metálico con incidencia prescrita. Antes de consultar la solución, prediga el efecto de cambiar el primer dato respecto del caso resuelto de esta familia. Justifique la tendencia y explique qué restricción podría convertirse en limitante. Datos: \(r=0.288615\, \mathrm{m}\); \(\Omega=1010\, \mathrm{s}^{-1}\); \(K_e=14\, \mathrm{m}^{2}\,\mathrm{s}^{-1}\); \(C_x=159\, \mathrm{m}\,\mathrm{s}^{-1}\); \(i=-0.01745329\, \mathrm{rad}\).

\noindent\textbf{Datos}
\begin{itemize}
\item $r = 0.288615\,\mathrm{m}$ — radio elegido
\item $\Omega = 1010\,\mathrm{s}^{-1}$ — velocidad angular
\item $K_e = 14\,\mathrm{m}^{2}\,\mathrm{s}^{-1}$ — constante remolino de entrada
\item $C_x = 159\,\mathrm{m}\,\mathrm{s}^{-1}$ — velocidad axial
\item $i = -0.017453292519943\,\mathrm{rad}$ — incidencia algebraica prescrita
\end{itemize}

\noindent\textbf{Incógnita}
$\kappa_e$, en $\mathrm{rad}$; también las magnitudes intermedias desarrolladas.

\noindent\textbf{Modelo físico}

\noindent\textit{Rotor compresor a radio elegido}

\[
\fitmath{\kappa_e=- i + \operatorname{atan}{\left(\frac{- \Omega r + \frac{K_e}{r}}{C_x} \right)}}
\]

\noindent\textbf{Hipótesis}
\begin{itemize}
\item Vórtice libre por estación, incidencia prescrita.
\end{itemize}

\noindent\textbf{Desarrollo paso a paso}

\begin{enumerate}

\item Velocidad periférica

\[\fitmath{U=\Omega\cdot \mathrm{rad}}\]

\[\fitmath{U=(1010)\cdot (0.288615)\approx 291.5011\,\mathrm{m/s}}\]

\item La velocidad tangencial absoluta sigue vórtice libre.

\[\fitmath{\mathrm{ct}=\frac{\mathrm{kv}}{\mathrm{rad}}}\]

\[\fitmath{\mathrm{ct}=\frac{(14)}{(0.288615)}\approx 48.50753\,\mathrm{m/s}}\]

\item Se forma el ángulo del gas en referencia rotor.

\[\fitmath{\beta=\arctan\left(\frac{\frac{\mathrm{kv}}{\mathrm{rad}}-\Omega\cdot \mathrm{rad}}{\mathrm{cx}}\right)}\]

\[\fitmath{\beta=\arctan\left(\frac{\frac{(14)}{(0.288615)}-(1010)\cdot (0.288615)}{(159)}\right)\approx -0.9913775\,\mathrm{rad}}\]

\item Se despeja metal a partir de incidencia definida.

\[\fitmath{\mathrm{kap}=\arctan\left(\frac{\frac{\mathrm{kv}}{\mathrm{rad}}-\Omega\cdot \mathrm{rad}}{\mathrm{cx}}\right)-\mathrm{inc}}\]

\[\fitmath{\mathrm{kap}=\arctan\left(\frac{\frac{(14)}{(0.288615)}-(1010)\cdot (0.288615)}{(159)}\right)-(-0.01745329252)\approx -0.9739242\,\mathrm{rad}}\]

\end{enumerate}


\noindent\textbf{Sustitución numérica}

\[
\fitmath{- \left(-0.017453292519943\,\mathrm{rad}\right) + \operatorname{atan}{\left(\frac{- \left(1010\,\mathrm{s}^{-1}\right) \left(0.288615\,\mathrm{m}\right) + \frac{\left(14\,\mathrm{m}^{2}\,\mathrm{s}^{-1}\right)}{\left(0.288615\,\mathrm{m}\right)}}{\left(159\,\mathrm{m}\,\mathrm{s}^{-1}\right)} \right)}}
\]

\noindent\textbf{Resultado}

\[
\fitmath{\kappa_e=-9.739241832\times10^{-1}\,\mathrm{rad}}
\]

\noindent\textbf{Velocidad periférica}
\[
\fitmath{U=\Omega r=291.50115\,\mathrm{m/s}}
\]

\noindent\textbf{Comprobación cinemática.} Insertar este $U$ en $\beta=\arctan[(C_\theta-U)/C_x]$ reproduce el ángulo desarrollado.

\noindent\textbf{Interpretación}
La misma tabla debe actualizarse al variar r u Omega. El metal es una propuesta geométrica educativa.

\noindent\textbf{Comprobación de los pasos}
\begin{itemize}
\item Paso 2: $48.5075$ — reproducido
\item Paso 4: $-0.991377$ — reproducido
\item Paso 6: $-0.973924$ — reproducido
\end{itemize}

\noindent\textbf{Comprobación de la respuesta} la aritmética del último paso reproduce la respuesta enunciada a la precisión con que está escrita. Es una reproducción aritmética; no valida el modelo físico ni sus hipótesis.

\noindent\textbf{Limitaciones}
\begin{itemize}
\item Modelo educativo; la aplicabilidad a una máquina real requiere revisión experta.
\end{itemize}

\noindent\textbf{Errores frecuentes}
\begin{itemize}
\item Usar alfa para construir un rotor.
\end{itemize}
% studium:end-record eyJpZCI6ICJQUkItMDEzOCIsICJzZWN0aW9uIjogImMxOCJ9
\subsection*{C18-02-P}

\noindent\textbf{Enunciado}
\(C18-02-P \cdot\) Estator compresor y orientación de salida. Ejercicio de transferencia. Con remolino de salida prescrito, halle el ángulo absoluto y el ángulo metálico para un estator compresor. Antes de consultar la solución, prediga el efecto de cambiar el primer dato respecto del caso resuelto de esta familia. Justifique la tendencia y explique qué restricción podría convertirse en limitante. Datos: \(r=0.30033\, \mathrm{m}\); \(K_s=9\, \mathrm{m}^{2}\,\mathrm{s}^{-1}\); \(C_x=162\, \mathrm{m}\,\mathrm{s}^{-1}\); \(\delta=0.03839724\, \mathrm{rad}\).

\noindent\textbf{Datos}
\begin{itemize}
\item $r = 0.30033\,\mathrm{m}$ — radio
\item $K_s = 9\,\mathrm{m}^{2}\,\mathrm{s}^{-1}$ — constante remolino salida
\item $C_x = 162\,\mathrm{m}\,\mathrm{s}^{-1}$ — velocidad axial
\item $\delta = 0.038397243543875\,\mathrm{rad}$ — desviación algebraica prescrita
\end{itemize}

\noindent\textbf{Incógnita}
$\kappa_e$, en $\mathrm{rad}$; también las magnitudes intermedias desarrolladas.

\noindent\textbf{Modelo físico}

\noindent\textit{Estator compresor y orientación de salida}

\[
\fitmath{\kappa_e=- \delta + \operatorname{atan}{\left(\frac{K_s}{C_x r} \right)}}
\]

\noindent\textbf{Hipótesis}
\begin{itemize}
\item Velocidades absolutas prescritas, desviación como dato.
\end{itemize}

\noindent\textbf{Desarrollo paso a paso}

\begin{enumerate}

\item La fila fija utiliza velocidad absoluta.

\[\fitmath{\alpha=\arctan\left(\frac{\mathrm{kv}}{\mathrm{rad}\cdot \mathrm{cx}}\right)}\]

\[\fitmath{\alpha=\arctan\left(\frac{(9)}{(0.30033)\cdot (162)}\right)\approx 0.1829141\,\mathrm{rad}}\]

\item La definición delta=gas-metal permite despejar metal.

\[\fitmath{\mathrm{kap}=\arctan\left(\frac{\mathrm{kv}}{\mathrm{rad}\cdot \mathrm{cx}}\right)-\mathrm{dev}}\]

\[\fitmath{\mathrm{kap}=\arctan\left(\frac{(9)}{(0.30033)\cdot (162)}\right)-(0.03839724354)\approx 0.1445168\,\mathrm{rad}}\]

\end{enumerate}


\noindent\textbf{Sustitución numérica}

\[
\fitmath{- \left(0.038397243543875\,\mathrm{rad}\right) + \operatorname{atan}{\left(\frac{\left(9\,\mathrm{m}^{2}\,\mathrm{s}^{-1}\right)}{\left(162\,\mathrm{m}\,\mathrm{s}^{-1}\right) \left(0.30033\,\mathrm{m}\right)} \right)}}
\]

\noindent\textbf{Resultado}

\[
\fitmath{\kappa_e=1.445168334\times10^{-1}\,\mathrm{rad}}
\]

\noindent\textbf{Interpretación}
No aparece Omega porque la fila está fija; solo cambiaría el resultado si el cambio de operación modifica los datos absolutos.

\noindent\textbf{Comprobación de los pasos}
\begin{itemize}
\item Paso 2: $0.182914$ — reproducido
\item Paso 4: $0.144517$ — reproducido
\end{itemize}

\noindent\textbf{Comprobación de la respuesta} la aritmética del último paso reproduce la respuesta enunciada a la precisión con que está escrita. Es una reproducción aritmética; no valida el modelo físico ni sus hipótesis.

\noindent\textbf{Limitaciones}
\begin{itemize}
\item Modelo educativo; la aplicabilidad a una máquina real requiere revisión experta.
\end{itemize}

\noindent\textbf{Errores frecuentes}
\begin{itemize}
\item Restar U a una fila fija.
\end{itemize}
% studium:end-record eyJpZCI6ICJQUkItMDE0MCIsICJzZWN0aW9uIjogImMxOCJ9
\subsection*{C18-03-P}

\noindent\textbf{Enunciado}
\(C18-03-P \cdot\) Estator turbina y calado educativo. Ejercicio de transferencia. Una sección educativa del estator de turbina tiene pendientes metálicas dadas. Con d igual al promedio de pendientes, calcule calado y cuerda real. Antes de consultar la solución, prediga el efecto de cambiar el primer dato respecto del caso resuelto de esta familia. Justifique la tendencia y explique qué restricción podría convertirse en limitante. Datos: \(m_e=0.1278\); \(m_s=1.38\); \(c_x=0.046\, \mathrm{m}\).

\noindent\textbf{Datos}
\begin{itemize}
\item $m_e = 0.1278$ — pendiente metal entrada
\item $m_s = 1.38$ — pendiente metal salida
\item $c_x = 0.046\,\mathrm{m}$ — cuerda axial
\end{itemize}

\noindent\textbf{Incógnita}
$c$, en $\mathrm{m}$; también las magnitudes intermedias desarrolladas.

\noindent\textbf{Modelo físico}

\noindent\textit{Estator turbina y calado educativo}

\[
\fitmath{c=c_x \sqrt{\left(\frac{m_e}{2} + \frac{m_s}{2}\right)^{2} + 1}}
\]

\noindent\textbf{Hipótesis}
\begin{itemize}
\item Línea media Hermite educativa; x monótona.
\end{itemize}

\noindent\textbf{Desarrollo paso a paso}

\begin{enumerate}

\item El desplazamiento normalizado se prescribe como promedio de pendientes.

\[\fitmath{\mathrm{dd}=\frac{\mathrm{me}+\mathrm{ms}}{2}}\]

\[\fitmath{\mathrm{dd}=\frac{(0.1278)+(1.38)}{2}\approx 0.7539\,}\]

\item El calado de cuerda es atan d.

\[\fitmath{\mathrm{chi}=\arctan\left(\frac{\mathrm{me}+\mathrm{ms}}{2}\right)}\]

\[\fitmath{\mathrm{chi}=\arctan\left(\frac{(0.1278)+(1.38)}{2}\right)\approx 0.6459924\,\mathrm{rad}}\]

\item La cuerda real se deduce del desplazamiento final.

\[\fitmath{\mathrm{chord}=\mathrm{cxg}\cdot \sqrt{1+{\left(\frac{\mathrm{me}+\mathrm{ms}}{2}\right)}^{2}}}\]

\[\fitmath{\mathrm{chord}=(0.046)\cdot \sqrt{1+{\left(\frac{(0.1278)+(1.38)}{2}\right)}^{2}}\approx 0.05760782\,\mathrm{m}}\]

\end{enumerate}


\noindent\textbf{Sustitución numérica}

\[
\fitmath{\left(0.046\,\mathrm{m}\right) \sqrt{\left(\frac{\left(0.1278\right)}{2} + \frac{\left(1.38\right)}{2}\right)^{2} + 1}}
\]

\noindent\textbf{Resultado}

\[
\fitmath{c=5.760781878\times10^{-2}\,\mathrm{m}}
\]

\noindent\textbf{Interpretación}
Promediar pendientes no equivale a promediar ángulos. La ley es geométrica y no una correlación de mínima pérdida.

\noindent\textbf{Comprobación de los pasos}
\begin{itemize}
\item Paso 2: $0.7539$ — reproducido
\item Paso 4: $0.645992$ — reproducido
\item Paso 6: $0.0576078$ — reproducido
\end{itemize}

\noindent\textbf{Comprobación de la respuesta} la aritmética del último paso reproduce la respuesta enunciada a la precisión con que está escrita. Es una reproducción aritmética; no valida el modelo físico ni sus hipótesis.

\noindent\textbf{Limitaciones}
\begin{itemize}
\item Modelo educativo; la aplicabilidad a una máquina real requiere revisión experta.
\end{itemize}

\noindent\textbf{Errores frecuentes}
\begin{itemize}
\item Presentar esta sección como un perfil de fabricante.
\end{itemize}
% studium:end-record eyJpZCI6ICJQUkItMDE0MiIsICJzZWN0aW9uIjogImMxOCJ9
\subsection*{C18-04-P}

\noindent\textbf{Enunciado}
\(C18-04-P \cdot\) Rotor turbina y espesor por radio. Ejercicio de transferencia. En la salida del rotor de turbina, calcule beta y la coordenada transversal de la superficie superior en un punto de pendiente conocida. Antes de consultar la solución, prediga el efecto de cambiar el primer dato respecto del caso resuelto de esta familia. Justifique la tendencia y explique qué restricción podría convertirse en limitante. Datos: \(r=0.31311\, \mathrm{m}\); \(\Omega=1120\, \mathrm{s}^{-1}\); \(K_s=-7\, \mathrm{m}^{2}\,\mathrm{s}^{-1}\); \(C_x=201\, \mathrm{m}\,\mathrm{s}^{-1}\); \(y=0.014\, \mathrm{m}\); \(t_n=0.002\, \mathrm{m}\); \(m_y=-0.62\).

\noindent\textbf{Datos}
\begin{itemize}
\item $r = 0.31311\,\mathrm{m}$ — radio
\item $\Omega = 1120\,\mathrm{s}^{-1}$ — velocidad angular
\item $K_s = -7\,\mathrm{m}^{2}\,\mathrm{s}^{-1}$ — constante remolino salida
\item $C_x = 201\,\mathrm{m}\,\mathrm{s}^{-1}$ — velocidad axial
\item $y = 0.014\,\mathrm{m}$ — ordenada línea media en el punto
\item $t_n = 0.002\,\mathrm{m}$ — semiespesor normal
\item $m_y = -0.62$ — pendiente de línea media en el punto
\end{itemize}

\noindent\textbf{Incógnita}
$y_s$, en $\mathrm{m}$; también las magnitudes intermedias desarrolladas.

\noindent\textbf{Modelo físico}

\noindent\textit{Rotor turbina y espesor por radio}

\[
\fitmath{y_s=\frac{t_n}{\sqrt{m_y^{2} + 1}} + y}
\]

\noindent\textbf{Hipótesis}
\begin{itemize}
\item Vórtice libre y geometría normal local educativa.
\end{itemize}

\noindent\textbf{Desarrollo paso a paso}

\begin{enumerate}

\item El ángulo del gas se forma en referencia móvil.

\[\fitmath{\beta=\arctan\left(\frac{\frac{\mathrm{kv}}{\mathrm{rad}}-\Omega\cdot \mathrm{rad}}{\mathrm{cx}}\right)}\]

\[\fitmath{\beta=\arctan\left(\frac{\frac{(-7)}{(0.31311)}-(1120)\cdot (0.31311)}{(201)}\right)\approx -1.07658\,\mathrm{rad}}\]

\item La superficie superior desplaza la ordenada según la normal.

\[\fitmath{\mathrm{ys}=\mathrm{yy}+\frac{\mathrm{tt}}{\sqrt{1+{\mathrm{slope}}^{2}}}}\]

\[\fitmath{\mathrm{ys}=(0.014)+\frac{(0.002)}{\sqrt{1+{(-0.62)}^{2}}}\approx 0.01569981\,\mathrm{m}}\]

\end{enumerate}


\noindent\textbf{Sustitución numérica}

\[
\fitmath{\frac{\left(0.002\,\mathrm{m}\right)}{\sqrt{\left(-0.62\right)^{2} + 1}} + \left(0.014\,\mathrm{m}\right)}
\]

\noindent\textbf{Resultado}

\[
\fitmath{y_s=1.569980538\times10^{-2}\,\mathrm{m}}
\]

\noindent\textbf{Interpretación}
El dibujo completo debe desplazar también x. Esta salida tiene remolino negativo definido, sin cambiar el sentido de giro.

\noindent\textbf{Comprobación de los pasos}
\begin{itemize}
\item Paso 2: $-1.07658$ — reproducido
\item Paso 4: $0.0156998$ — reproducido
\end{itemize}

\noindent\textbf{Comprobación de la respuesta} la aritmética del último paso reproduce la respuesta enunciada a la precisión con que está escrita. Es una reproducción aritmética; no valida el modelo físico ni sus hipótesis.

\noindent\textbf{Limitaciones}
\begin{itemize}
\item Modelo educativo; la aplicabilidad a una máquina real requiere revisión experta.
\end{itemize}

\noindent\textbf{Errores frecuentes}
\begin{itemize}
\item Cambiar de signo U para representar una turbina.
\end{itemize}
% studium:end-record eyJpZCI6ICJQUkItMDE0NCIsICJzZWN0aW9uIjogImMxOCJ9
\section{Capítulo 19: Simulacros originales y estrategia de resolución}
\subsection*{C19-01-P}

\noindent\textbf{Enunciado}
\(C19-01-P \cdot A1 \cdot\) Cierre de compresor axial. Un rotor compresor al radio medio recibe Cx y remolinos dados. Obtenga trabajo, temperatura total final, presión total final y área geométrica de entrada. Resultado final: área de entrada. Datos: \(T_{0e}=302\, \mathrm{K}\); \(p_{0e}=118000\, \mathrm{Pa}\); \(C_x=156\, \mathrm{m}\,\mathrm{s}^{-1}\); \(C_{\theta e}=38\, \mathrm{m}\,\mathrm{s}^{-1}\); \(C_{\theta s}=147\, \mathrm{m}\,\mathrm{s}^{-1}\); \(U=291\, \mathrm{m}\,\mathrm{s}^{-1}\); \(\dot m=21\, \mathrm{kg}\,\mathrm{s}^{-1}\); \(\eta=0.87\); \(B=0.96\); \(c_p=1004.5\, \mathrm{J}\,\mathrm{kg}^{-1}\,\mathrm{K}^{-1}\); \(\gamma=1.4\); \(R_g=287\, \mathrm{J}\,\mathrm{kg}^{-1}\,\mathrm{K}^{-1}\).

\noindent\textbf{Datos}
\begin{itemize}
\item $T_{0e} = 302\,\mathrm{K}$ — temperatura total entrada
\item $p_{0e} = 118000\,\mathrm{Pa}$ — presión total entrada
\item $C_x = 156\,\mathrm{m}\,\mathrm{s}^{-1}$ — componente axial
\item $C_{\theta e} = 38\,\mathrm{m}\,\mathrm{s}^{-1}$ — remolino entrada rotor
\item $C_{\theta s} = 147\,\mathrm{m}\,\mathrm{s}^{-1}$ — remolino salida rotor
\item $U = 291\,\mathrm{m}\,\mathrm{s}^{-1}$ — velocidad periférica
\item $\dot m = 21\,\mathrm{kg}\,\mathrm{s}^{-1}$ — caudal
\item $\eta = 0.87$ — eficiencia total-total
\item $B = 0.96$ — fracción disponible
\item $c_p = 1004.5\,\mathrm{J}\,\mathrm{kg}^{-1}\,\mathrm{K}^{-1}$ — calor específico
\item $\gamma = 1.4$ — razón de calores
\item $R_g = 287\,\mathrm{J}\,\mathrm{kg}^{-1}\,\mathrm{K}^{-1}$ — constante gas
\end{itemize}

\noindent\textbf{Incógnita}
$A$, en $\mathrm{m^{2}}$; también las magnitudes intermedias desarrolladas.

\noindent\textbf{Modelo físico}

\noindent\textit{A1 \(\cdot\) Cierre de compresor axial}

\[
\fitmath{A=\frac{R_g \dot m \left(\frac{T_{0e} - \frac{C_{\theta e}^{2} + C_x^{2}}{2 c_p}}{T_{0e}}\right)^{- \frac{\gamma}{\gamma - 1}} \left(T_{0e} - \frac{C_{\theta e}^{2} + C_x^{2}}{2 c_p}\right)}{B C_x p_{0e}}}
\]

\noindent\textbf{Hipótesis}
\begin{itemize}
\item Gas perfecto, radio medio y Cx constantes, escalón adiabático.
\end{itemize}

\noindent\textbf{Desarrollo paso a paso}

\begin{enumerate}

\item Euler fija el trabajo.

\[\fitmath{\mathrm{wc}=\mathrm{uu}\cdot \left(\mathrm{cs}-\mathrm{ce}\right)}\]

\[\fitmath{\mathrm{wc}=(291)\cdot \left((147)-(38)\right)\approx 31719\,\mathrm{J/kg}}\]

\item El balance de energía fija la temperatura total de salida.

\[\fitmath{\mathrm{Tout}=\mathrm{Tin}+\frac{\mathrm{uu}\cdot \left(\mathrm{cs}-\mathrm{ce}\right)}{\mathrm{cp}}}\]

\[\fitmath{\mathrm{Tout}=(302)+\frac{(291)\cdot \left((147)-(38)\right)}{(1004.5)}\approx 333.5769\,\mathrm{K}}\]

\item La eficiencia determina presión total final.

\[\fitmath{\mathrm{pout}=\mathrm{pin}\cdot {\left(1+\frac{\eta\cdot \mathrm{uu}\cdot \left(\mathrm{cs}-\mathrm{ce}\right)}{\mathrm{cp}\cdot \mathrm{Tin}}\right)}^{\frac{\gamma}{\gamma-1}}}\]

\[\fitmath{\mathrm{pout}=(118000)\cdot {\left(1+\frac{(0.87)\cdot (291)\cdot \left((147)-(38)\right)}{(1004.5)\cdot (302)}\right)}^{\frac{(1.4)}{(1.4)-1}}\approx 160037.6\,\mathrm{Pa}}\]

\item La temperatura estática de entrada usa el módulo absoluto.

\[\fitmath{\mathrm{Te}=\mathrm{Tin}-\frac{{\mathrm{cx}}^{2}+{\mathrm{ce}}^{2}}{2\cdot \mathrm{cp}}}\]

\[\fitmath{\mathrm{Te}=(302)-\frac{{(156)}^{2}+{(38)}^{2}}{2\cdot (1004.5)}\approx 289.1677\,\mathrm{K}}\]

\item La continuidad utiliza densidad estática y componente axial.

\[\fitmath{\mathrm{area}=\frac{\mathrm{md}\cdot R_g\cdot \left(\mathrm{Tin}-\frac{{\mathrm{cx}}^{2}+{\mathrm{ce}}^{2}}{2\cdot \mathrm{cp}}\right)}{\mathrm{cx}\cdot \mathrm{bb}\cdot \mathrm{pin}\cdot {\left(\frac{\mathrm{Tin}-\frac{{\mathrm{cx}}^{2}+{\mathrm{ce}}^{2}}{2\cdot \mathrm{cp}}}{\mathrm{Tin}}\right)}^{\frac{\gamma}{\gamma-1}}}}\]

\[\fitmath{\mathrm{area}=\frac{(21)\cdot (287)\cdot \left((302)-\frac{{(156)}^{2}+{(38)}^{2}}{2\cdot (1004.5)}\right)}{(156)\cdot (0.96)\cdot (118000)\cdot {\left(\frac{(302)-\frac{{(156)}^{2}+{(38)}^{2}}{2\cdot (1004.5)}}{(302)}\right)}^{\frac{(1.4)}{(1.4)-1}}}\approx 0.1148082\,\mathrm{m^{2}}}\]

\end{enumerate}


\noindent\textbf{Sustitución numérica}

\[
\fitmath{\frac{\left(287\,\mathrm{J}\,\mathrm{kg}^{-1}\,\mathrm{K}^{-1}\right) \left(21\,\mathrm{kg}\,\mathrm{s}^{-1}\right) \left(\frac{\left(302\,\mathrm{K}\right) - \frac{\left(\left(38\,\mathrm{m}\,\mathrm{s}^{-1}\right)\right)^{2} + \left(\left(156\,\mathrm{m}\,\mathrm{s}^{-1}\right)\right)^{2}}{2 \left(1004.5\,\mathrm{J}\,\mathrm{kg}^{-1}\,\mathrm{K}^{-1}\right)}}{\left(302\,\mathrm{K}\right)}\right)^{- \frac{\left(1.4\right)}{\left(1.4\right) - 1}} \left(\left(302\,\mathrm{K}\right) - \frac{\left(\left(38\,\mathrm{m}\,\mathrm{s}^{-1}\right)\right)^{2} + \left(\left(156\,\mathrm{m}\,\mathrm{s}^{-1}\right)\right)^{2}}{2 \left(1004.5\,\mathrm{J}\,\mathrm{kg}^{-1}\,\mathrm{K}^{-1}\right)}\right)}{\left(0.96\right) \left(156\,\mathrm{m}\,\mathrm{s}^{-1}\right) \left(118000\,\mathrm{Pa}\right)}}
\]

\noindent\textbf{Resultado}

\[
\fitmath{A=1.148082328\times10^{-1}\,\mathrm{m}^{2}}
\]

\noindent\textbf{Interpretación}
Reconstruya masa con rho estática y B A. La presión final no determina la pérdida de cada fila por separado.

\noindent\textbf{Comprobación de los pasos}
\begin{itemize}
\item Paso 2: $31719$ — reproducido
\item Paso 4: $333.577$ — reproducido
\item Paso 6: $160038$ — reproducido
\item Paso 8: $289.168$ — reproducido
\item Paso 10: $0.114808$ — reproducido
\end{itemize}

\noindent\textbf{Comprobación de la respuesta} la aritmética del último paso reproduce la respuesta enunciada a la precisión con que está escrita. Es una reproducción aritmética; no valida el modelo físico ni sus hipótesis.

\noindent\textbf{Limitaciones}
\begin{itemize}
\item Modelo educativo; la aplicabilidad a una máquina real requiere revisión experta.
\end{itemize}

\noindent\textbf{Errores frecuentes}
\begin{itemize}
\item Usar densidad total para el anillo.
\end{itemize}
% studium:end-record eyJpZCI6ICJQUkItMDE0NSIsICJzZWN0aW9uIjogImMxOSJ9
\begin{center}\begin{tikzpicture}\draw[thick] (0,0)--(7,0);\draw[thick] (0,2.5)--(7,2.5);\node[draw,fill=orange!20,minimum width=2cm,minimum height=1.5cm] at (3.5,1.25) {Rotor};\draw[->,blue,thick] (.2,1.2)--(2,1.2) node[midway,above] {$C_x,C_{\theta e}$};\draw[->,blue,thick] (5,1.2)--(6.8,1.2) node[midway,above] {$C_x,C_{\theta s}$};\draw[->,orange,thick] (3.5,2.8)--(3.5,3.6) node[above] {$U>0$};\node[below,align=center] at (3.5,-0.7) {$w_c=3.1719\times10^{4}\,\mathrm{J}\,\mathrm{kg}^{-1},\quad T_{0s}=3.3358\times10^{2}\,\mathrm{K}$};\node[below,align=center] at (3.5,-1.2999999999999998) {$p_{0s}=1.6004\times10^{5}\,\mathrm{Pa},\quad T_e=2.8917\times10^{2}\,\mathrm{K}$};\node[below,align=center] at (3.5,-1.9) {$A=1.1481\times10^{-1}\,\mathrm{m}^{2}$};\end{tikzpicture}\end{center}
\subsection*{C19-02-P}

\noindent\textbf{Enunciado}
\(C19-02-P \cdot A2 \cdot\) Turbina con descarga axial. Diseñe un escalón repetitivo de turbina que descarga axialmente. Calcule carga, trabajo, temperatura total de salida y presión total de salida. Datos: \(T_{0e}=1490\, \mathrm{K}\); \(p_{0e}=940000\, \mathrm{Pa}\); \(\mathcal R=0.48\); \(U=314\, \mathrm{m}\,\mathrm{s}^{-1}\); \(\eta=0.895\); \(c_p=1156.697\, \mathrm{J}\,\mathrm{kg}^{-1}\,\mathrm{K}^{-1}\); \(\gamma=1.33\).

\noindent\textbf{Datos}
\begin{itemize}
\item $T_{0e} = 1490\,\mathrm{K}$ — temperatura total entrada
\item $p_{0e} = 940000\,\mathrm{Pa}$ — presión total entrada
\item $\mathcal R = 0.48$ — reacción
\item $U = 314\,\mathrm{m}\,\mathrm{s}^{-1}$ — velocidad periférica
\item $\eta = 0.895$ — eficiencia total-total
\item $c_p = 1156.696969697\,\mathrm{J}\,\mathrm{kg}^{-1}\,\mathrm{K}^{-1}$ — calor específico
\item $\gamma = 1.33$ — razón de calores
\end{itemize}

\noindent\textbf{Incógnita}
$p_{0s}$, en $\mathrm{Pa}$; también las magnitudes intermedias desarrolladas.

\noindent\textbf{Modelo físico}

\noindent\textit{A2 \(\cdot\) Turbina con descarga axial}

\[
\fitmath{p_{0s}=p_{0e} \left(1 - \frac{U^{2} \left(2 - 2 \mathcal R\right)}{T_{0e} c_p \eta}\right)^{\frac{\gamma}{\gamma - 1}}}
\]

\noindent\textbf{Hipótesis}
\begin{itemize}
\item Repetitivo, mismo radio y Cx constantes, sin refrigeración.
\end{itemize}

\noindent\textbf{Desarrollo paso a paso}

\begin{enumerate}

\item La descarga axial liga carga y reacción.

\[\fitmath{\mathrm{psi}=2\cdot \left(1-\mathrm{rea}\right)}\]

\[\fitmath{\mathrm{psi}=2\cdot \left(1-(0.48)\right)\approx 1.04\,}\]

\item Euler da el trabajo extraído.

\[\fitmath{\mathrm{wt}=2\cdot \left(1-\mathrm{rea}\right)\cdot {\mathrm{uu}}^{2}}\]

\[\fitmath{\mathrm{wt}=2\cdot \left(1-(0.48)\right)\cdot {(314)}^{2}\approx 102539.8\,\mathrm{J/kg}}\]

\item El remanso disminuye por el trabajo.

\[\fitmath{\mathrm{Tout}=\mathrm{Tin}-\frac{2\cdot \left(1-\mathrm{rea}\right)\cdot {\mathrm{uu}}^{2}}{\mathrm{cp}}}\]

\[\fitmath{\mathrm{Tout}=(1490)-\frac{2\cdot \left(1-(0.48)\right)\cdot {(314)}^{2}}{(1156.69697)}\approx 1401.351\,\mathrm{K}}\]

\item El estado ideal equivalente permite hallar presión total.

\[\fitmath{\mathrm{pout}=\mathrm{pin}\cdot {\left(1-\frac{2\cdot \left(1-\mathrm{rea}\right)\cdot {\mathrm{uu}}^{2}}{\mathrm{eta}\cdot \mathrm{cp}\cdot \mathrm{Tin}}\right)}^{\frac{\gamma}{\gamma-1}}}\]

\[\fitmath{\mathrm{pout}=(940000)\cdot {\left(1-\frac{2\cdot \left(1-(0.48)\right)\cdot {(314)}^{2}}{(0.895)\cdot (1156.69697)\cdot (1490)}\right)}^{\frac{(1.33)}{(1.33)-1}}\approx 712401.5\,\mathrm{Pa}}\]

\end{enumerate}


\noindent\textbf{Sustitución numérica}

\[
\fitmath{\left(940000\,\mathrm{Pa}\right) \left(1 - \frac{\left(\left(314\,\mathrm{m}\,\mathrm{s}^{-1}\right)\right)^{2} \left(2 - 2 \left(0.48\right)\right)}{\left(1490\,\mathrm{K}\right) \left(1156.696969697\,\mathrm{J}\,\mathrm{kg}^{-1}\,\mathrm{K}^{-1}\right) \left(0.895\right)}\right)^{\frac{\left(1.33\right)}{\left(1.33\right) - 1}}}
\]

\noindent\textbf{Resultado}

\[
\fitmath{p_{0s}=7.124015323\times10^{5}\,\mathrm{Pa}}
\]

\noindent\textbf{Interpretación}
La razón de presiones es menor que uno y la base isentrópica debe permanecer positiva. No se ha calculado la presión estática sin una velocidad de salida.

\noindent\textbf{Comprobación de los pasos}
\begin{itemize}
\item Paso 2: $1.04$ — reproducido
\item Paso 4: $102540$ — reproducido
\item Paso 6: $1401.35$ — reproducido
\item Paso 8: $712402$ — reproducido
\end{itemize}

\noindent\textbf{Comprobación de la respuesta} la aritmética del último paso reproduce la respuesta enunciada a la precisión con que está escrita. Es una reproducción aritmética; no valida el modelo físico ni sus hipótesis.

\noindent\textbf{Limitaciones}
\begin{itemize}
\item Modelo educativo; la aplicabilidad a una máquina real requiere revisión experta.
\end{itemize}

\noindent\textbf{Errores frecuentes}
\begin{itemize}
\item Imponer una carga incompatible con reacción y descarga axial.
\end{itemize}
% studium:end-record eyJpZCI6ICJQUkItMDE0NiIsICJzZWN0aW9uIjogImMxOSJ9
\begin{center}\begin{tikzpicture}\node[draw,fill=blue!15,minimum height=1.5cm] (s) at (2,1) {Estator};\node[draw,fill=orange!20,minimum height=1.5cm] (r) at (5,1) {Rotor};\draw[->] (0,1)--(s);\draw[->] (s)--(r) node[midway,above] {2};\draw[->] (r)--(7,1) node[right] {$C_{\theta3}=0$};\node at (0,1.6) {1};\node at (7,.4) {3};\draw[->,orange,thick] (5,2)--(5,2.8) node[above] {$U$};\node[below,align=center] at (3.5,-0.7) {$\psi=1.0400\times10^{0}\,,\quad w_t=1.0254\times10^{5}\,\mathrm{J}\,\mathrm{kg}^{-1}$};\node[below,align=center] at (3.5,-1.2999999999999998) {$T_{0s}=1.4014\times10^{3}\,\mathrm{K},\quad p_{0s}=7.1240\times10^{5}\,\mathrm{Pa}$};\end{tikzpicture}\end{center}
\subsection*{C19-03-P}

\noindent\textbf{Enunciado}
\(C19-03-P \cdot A3 \cdot\) Distribución radial y tracción. Un escalón compresor tiene leyes libres en ambas estaciones. Calcule trabajo radialmente uniforme, reacción en raíz y tensión centrífuga nominal de pala uniforme. Resultado final: tensión. Datos: \(\Omega=1090\, \mathrm{s}^{-1}\); \(K_e=14\, \mathrm{m}^{2}\,\mathrm{s}^{-1}\); \(K_s=44\, \mathrm{m}^{2}\,\mathrm{s}^{-1}\); \(r_h=0.218\, \mathrm{m}\); \(r_t=0.302\, \mathrm{m}\); \(\rho=8050\, \mathrm{kg}\,\mathrm{m}^{-3}\).

\noindent\textbf{Datos}
\begin{itemize}
\item $\Omega = 1090\,\mathrm{s}^{-1}$ — velocidad angular
\item $K_e = 14\,\mathrm{m}^{2}\,\mathrm{s}^{-1}$ — constante entrada
\item $K_s = 44\,\mathrm{m}^{2}\,\mathrm{s}^{-1}$ — constante salida
\item $r_h = 0.218\,\mathrm{m}$ — radio raíz
\item $r_t = 0.302\,\mathrm{m}$ — radio punta
\item $\rho = 8050\,\mathrm{kg}\,\mathrm{m}^{-3}$ — densidad material
\end{itemize}

\noindent\textbf{Incógnita}
$\sigma$, en $\mathrm{Pa}$; también las magnitudes intermedias desarrolladas.

\noindent\textbf{Modelo físico}

\noindent\textit{A3 \(\cdot\) Distribución radial y tracción}

\[
\fitmath{\sigma=\frac{\Omega^{2} \rho \left(- r_h^{2} + r_t^{2}\right)}{2}}
\]

\noindent\textbf{Hipótesis}
\begin{itemize}
\item Reacción reducida de escalón repetitivo; pala mecánica uniforme.
\end{itemize}

\noindent\textbf{Desarrollo paso a paso}

\begin{enumerate}

\item Euler hace uniforme el trabajo para estas leyes libres.

\[\fitmath{\mathrm{wc}=\Omega\cdot \left(\mathrm{ks}-\mathrm{ke}\right)}\]

\[\fitmath{\mathrm{wc}=(1090)\cdot \left((44)-(14)\right)\approx 32700\,\mathrm{J/kg}}\]

\item La reacción se evalúa con U local en raíz.

\[\fitmath{\mathrm{rea}=1-\frac{\mathrm{ke}+\mathrm{ks}}{2\cdot \Omega\cdot {\mathrm{rh}}^{2}}}\]

\[\fitmath{\mathrm{rea}=1-\frac{(14)+(44)}{2\cdot (1090)\cdot {(0.218)}^{2}}\approx 0.440167\,}\]

\item La tracción nominal integra toda la masa exterior.

\[\fitmath{\mathrm{stress}=\frac{\rho\cdot {\Omega}^{2}\cdot \left({\mathrm{rt}}^{2}-{\mathrm{rh}}^{2}\right)}{2}}\]

\[\fitmath{\mathrm{stress}=\frac{(8050)\cdot {(1090)}^{2}\cdot \left({(0.302)}^{2}-{(0.218)}^{2}\right)}{2}\approx 2.088822\times10^{8}\,\mathrm{Pa}}\]

\end{enumerate}


\noindent\textbf{Sustitución numérica}

\[
\fitmath{\frac{\left(\left(1090\,\mathrm{s}^{-1}\right)\right)^{2} \left(8050\,\mathrm{kg}\,\mathrm{m}^{-3}\right) \left(- \left(0.218\,\mathrm{m}\right)^{2} + \left(0.302\,\mathrm{m}\right)^{2}\right)}{2}}
\]

\noindent\textbf{Resultado}

\[
\fitmath{\sigma=2.088822372\times10^{8}\,\mathrm{Pa}}
\]

\noindent\textbf{Interpretación}
La admisibilidad aerodinámica y mecánica son comprobaciones distintas. No se ha suministrado una curva de vida ni concentración de raíz.

\noindent\textbf{Comprobación de los pasos}
\begin{itemize}
\item Paso 2: $32700$ — reproducido
\item Paso 4: $0.440167$ — reproducido
\item Paso 6: $2.08882\times 10^{8}$ — reproducido
\end{itemize}

\noindent\textbf{Comprobación de la respuesta} la aritmética del último paso reproduce la respuesta enunciada a la precisión con que está escrita. Es una reproducción aritmética; no valida el modelo físico ni sus hipótesis.

\noindent\textbf{Limitaciones}
\begin{itemize}
\item Modelo educativo; la aplicabilidad a una máquina real requiere revisión experta.
\end{itemize}

\noindent\textbf{Errores frecuentes}
\begin{itemize}
\item Usar la reacción media en toda la pala.
\end{itemize}
% studium:end-record eyJpZCI6ICJQUkItMDE0NyIsICJzZWN0aW9uIjogImMxOSJ9
\begin{center}\begin{tikzpicture}\draw[fill=gray!15,thick] (1,0) rectangle (2,3.7);\draw[thick] (.5,0)--(2.5,0);\foreach \x in {.6,.9,...,2.4}{\draw (\x,0)--(\x-.25,-.25);}\draw[->,orange] (1.5,2)--(1.5,4.5) node[above] {Inercia centrífuga};\draw[<->] (3.2,0)--(3.2,3.7) node[midway,right] {$r_t-r_h$};\node[left] at (1,0) {$r_h$};\node[left] at (1,3.7) {$r_t$};\node at (5.5,2.8) {$rC_\theta=K_v$};\node at (5.5,1.7) {$U=\Omega r$};\node[below,align=center] at (3.5,-0.7) {$w_c=3.2700\times10^{4}\,\mathrm{J}\,\mathrm{kg}^{-1},\quad \mathcal R=4.4017\times10^{-1}\,$};\node[below,align=center] at (3.5,-1.2999999999999998) {$\sigma=2.0888\times10^{8}\,\mathrm{Pa}$};\end{tikzpicture}\end{center}
\subsection*{C19-04-P}

\noindent\textbf{Enunciado}
\(C19-04-P \cdot B1 \cdot\) Rodete radial y prestación. Un rodete de salida radial sin prerrotación tiene radio y giro dados. Con slip y eficiencia prescritos, halle trabajo, temperatura total y relación de presión. Datos: \(\Omega=2070\, \mathrm{s}^{-1}\); \(r=0.206\, \mathrm{m}\); \(\sigma_s=0.875\); \(\eta=0.825\); \(T_{0e}=296\, \mathrm{K}\); \(c_p=1004.5\, \mathrm{J}\,\mathrm{kg}^{-1}\,\mathrm{K}^{-1}\); \(\gamma=1.4\).

\noindent\textbf{Datos}
\begin{itemize}
\item $\Omega = 2070\,\mathrm{s}^{-1}$ — velocidad angular
\item $r = 0.206\,\mathrm{m}$ — radio salida
\item $\sigma_s = 0.875$ — factor $C_\theta/C_{\theta,\mathrm{ideal}}$
\item $\eta = 0.825$ — eficiencia total-total
\item $T_{0e} = 296\,\mathrm{K}$ — temperatura total entrada
\item $c_p = 1004.5\,\mathrm{J}\,\mathrm{kg}^{-1}\,\mathrm{K}^{-1}$ — calor específico
\item $\gamma = 1.4$ — razón de calores
\end{itemize}

\noindent\textbf{Incógnita}
$\pi_c$ (adimensional); también las magnitudes intermedias desarrolladas.

\noindent\textbf{Modelo físico}

\noindent\textit{B1 \(\cdot\) Rodete radial y prestación}

\[
\fitmath{\pi_c=\left(\frac{\Omega^{2} \eta r^{2} \sigma_s}{T_{0e} c_p} + 1\right)^{\frac{\gamma}{\gamma - 1}}}
\]

\noindent\textbf{Hipótesis}
\begin{itemize}
\item Gas perfecto, salida radial ideal antes de slip; sin prerrotación.
\end{itemize}

\noindent\textbf{Desarrollo paso a paso}

\begin{enumerate}

\item La velocidad periférica depende del radio de salida.

\[\fitmath{\mathrm{uu}=\Omega\cdot \mathrm{rad}}\]

\[\fitmath{\mathrm{uu}=(2070)\cdot (0.206)\approx 426.42\,\mathrm{m/s}}\]

\item El deslizamiento modifica el remolino y el trabajo.

\[\fitmath{\mathrm{wc}=\mathrm{slip}\cdot {\left(\Omega\cdot \mathrm{rad}\right)}^{2}}\]

\[\fitmath{\mathrm{wc}=(0.875)\cdot {\left((2070)\cdot (0.206)\right)}^{2}\approx 159104.8\,\mathrm{J/kg}}\]

\item La temperatura total aumenta según el trabajo real.

\[\fitmath{\mathrm{Tout}=\mathrm{Tin}+\frac{\mathrm{slip}\cdot {\left(\Omega\cdot \mathrm{rad}\right)}^{2}}{\mathrm{cp}}}\]

\[\fitmath{\mathrm{Tout}=(296)+\frac{(0.875)\cdot {\left((2070)\cdot (0.206)\right)}^{2}}{(1004.5)}\approx 454.392\,\mathrm{K}}\]

\item La eficiencia enlaza trabajo con razón de presión.

\[\fitmath{\mathrm{pr}={\left(1+\frac{\mathrm{eta}\cdot \mathrm{slip}\cdot {\left(\Omega\cdot \mathrm{rad}\right)}^{2}}{\mathrm{cp}\cdot \mathrm{Tin}}\right)}^{\frac{\gamma}{\gamma-1}}}\]

\[\fitmath{\mathrm{pr}={\left(1+\frac{(0.825)\cdot (0.875)\cdot {\left((2070)\cdot (0.206)\right)}^{2}}{(1004.5)\cdot (296)}\right)}^{\frac{(1.4)}{(1.4)-1}}\approx 3.595949\,}\]

\end{enumerate}


\noindent\textbf{Sustitución numérica}

\[
\fitmath{\left(\frac{\left(\left(2070\,\mathrm{s}^{-1}\right)\right)^{2} \left(0.825\right) \left(0.206\,\mathrm{m}\right)^{2} \left(0.875\right)}{\left(296\,\mathrm{K}\right) \left(1004.5\,\mathrm{J}\,\mathrm{kg}^{-1}\,\mathrm{K}^{-1}\right)} + 1\right)^{\frac{\left(1.4\right)}{\left(1.4\right) - 1}}}
\]

\noindent\textbf{Resultado}

\[
\fitmath{\pi_c=3.595948824\times10^{0}}
\]

\noindent\textbf{Interpretación}
El resultado no es un mapa. Deben revisarse Mach de inductor, difusión, giro y tensión para juzgar realizabilidad.

\noindent\textbf{Comprobación de los pasos}
\begin{itemize}
\item Paso 2: $426.42$ — reproducido
\item Paso 4: $159105$ — reproducido
\item Paso 6: $454.392$ — reproducido
\item Paso 8: $3.59595$ — reproducido
\end{itemize}

\noindent\textbf{Comprobación de la respuesta} la aritmética del último paso reproduce la respuesta enunciada a la precisión con que está escrita. Es una reproducción aritmética; no valida el modelo físico ni sus hipótesis.

\noindent\textbf{Limitaciones}
\begin{itemize}
\item Modelo educativo; la aplicabilidad a una máquina real requiere revisión experta.
\end{itemize}

\noindent\textbf{Errores frecuentes}
\begin{itemize}
\item Confundir slip y rendimiento isentrópico.
\end{itemize}
% studium:end-record eyJpZCI6ICJQUkItMDE0OCIsICJzZWN0aW9uIjogImMxOSJ9
\begin{center}\begin{tikzpicture}\draw[thick] (2.5,1.7) circle (1.5);\draw[thick] (2.5,1.7) circle (.45);\foreach \a in {0,45,...,315}{\draw[blue,rotate around={\a:(2.5,1.7)}] (2.95,1.7)--(4,1.7);}\draw[->,orange,thick] (3,3.6) arc (70:140:1.8);\node at (2.5,4) {$\Omega>0$};\draw[->] (4,1.7)--(6,1.7) node[right] {$C_{r2}$};\draw[->] (4,1.7)--(4,2.8) node[right] {$C_{\theta2}$};\node at (5.7,.5) {Prerrotación nula};\node[below,align=center] at (3.5,-0.7) {$U=4.2642\times10^{2}\,\mathrm{m}\,\mathrm{s}^{-1},\quad w_c=1.5910\times10^{5}\,\mathrm{J}\,\mathrm{kg}^{-1}$};\node[below,align=center] at (3.5,-1.2999999999999998) {$T_{0s}=4.5439\times10^{2}\,\mathrm{K},\quad \pi_c=3.5959\times10^{0}\,$};\end{tikzpicture}\end{center}
\subsection*{C19-05-P}

\noindent\textbf{Enunciado}
\(C19-05-P \cdot B2 \cdot\) Cascada y elección entera. Una fila compresora tiene velocidades relativas y cambio tangencial dados. Se ha elegido un número entero de álabes; calcule paso, solidez real y difusión. Datos: \(r=0.279\, \mathrm{m}\); \(c=0.054\, \mathrm{m}\); \(Z=43\); \(W_e=296\, \mathrm{m}\,\mathrm{s}^{-1}\); \(W_s=224\, \mathrm{m}\,\mathrm{s}^{-1}\); \(|\Delta W_\theta|=103\, \mathrm{m}\,\mathrm{s}^{-1}\).

\noindent\textbf{Datos}
\begin{itemize}
\item $r = 0.279\,\mathrm{m}$ — radio
\item $c = 0.054\,\mathrm{m}$ — cuerda real
\item $Z = 43$ — número de álabes
\item $W_e = 296\,\mathrm{m}\,\mathrm{s}^{-1}$ — módulo relativo entrada
\item $W_s = 224\,\mathrm{m}\,\mathrm{s}^{-1}$ — módulo relativo salida
\item $|\Delta W_\theta| = 103\,\mathrm{m}\,\mathrm{s}^{-1}$ — magnitud de cambio tangencial
\end{itemize}

\noindent\textbf{Incógnita}
$D$ (adimensional); también las magnitudes intermedias desarrolladas.

\noindent\textbf{Modelo físico}

\noindent\textit{B2 \(\cdot\) Cascada y elección entera}

\[
\fitmath{D=1 - \frac{W_s}{W_e} + \frac{\pi |\Delta W_\theta| r}{c W_e Z}}
\]

\noindent\textbf{Hipótesis}
\begin{itemize}
\item Indicador de difusión definido en referencia relativa.
\end{itemize}

\noindent\textbf{Desarrollo paso a paso}

\begin{enumerate}

\item El entero fija el paso real.

\[\fitmath{\mathrm{ss}=\frac{2\cdot \pi\cdot \mathrm{rad}}{\mathrm{zz}}}\]

\[\fitmath{\mathrm{ss}=\frac{2\cdot \pi\cdot (0.279)}{(43)}\approx 0.04076764\,\mathrm{m}}\]

\item Se calcula la solidez alcanzada.

\[\fitmath{\mathrm{sol}=\frac{\mathrm{chord}\cdot \mathrm{zz}}{2\cdot \pi\cdot \mathrm{rad}}}\]

\[\fitmath{\mathrm{sol}=\frac{(0.054)\cdot (43)}{2\cdot \pi\cdot (0.279)}\approx 1.32458\,}\]

\item La difusión utiliza esa solidez, no una anterior al redondeo.

\[\fitmath{D=1-\frac{\mathrm{vs}}{\mathrm{ve}}+\frac{\mathrm{dv}}{2\cdot \frac{\mathrm{chord}\cdot \mathrm{zz}}{2\cdot \pi\cdot \mathrm{rad}}\cdot \mathrm{ve}}}\]

\[\fitmath{D=1-\frac{(224)}{(296)}+\frac{(103)}{2\cdot \frac{(0.054)\cdot (43)}{2\cdot \pi\cdot (0.279)}\cdot (296)}\approx 0.3745955\,}\]

\end{enumerate}


\noindent\textbf{Sustitución numérica}

\[
\fitmath{1 - \frac{\left(224\,\mathrm{m}\,\mathrm{s}^{-1}\right)}{\left(296\,\mathrm{m}\,\mathrm{s}^{-1}\right)} + \frac{\pi \left(103\,\mathrm{m}\,\mathrm{s}^{-1}\right) \left(0.279\,\mathrm{m}\right)}{\left(0.054\,\mathrm{m}\right) \left(296\,\mathrm{m}\,\mathrm{s}^{-1}\right) \left(43\right)}}
\]

\noindent\textbf{Resultado}

\[
\fitmath{D=3.745954502\times10^{-1}}
\]

\noindent\textbf{Interpretación}
Compare con un límite solo si se ha suministrado su rango de validez. D no determina una eficiencia del perfil.

\noindent\textbf{Comprobación de los pasos}
\begin{itemize}
\item Paso 2: $0.0407676$ — reproducido
\item Paso 4: $1.32458$ — reproducido
\item Paso 6: $0.374595$ — reproducido
\end{itemize}

\noindent\textbf{Comprobación de la respuesta} la aritmética del último paso reproduce la respuesta enunciada a la precisión con que está escrita. Es una reproducción aritmética; no valida el modelo físico ni sus hipótesis.

\noindent\textbf{Limitaciones}
\begin{itemize}
\item Modelo educativo; la aplicabilidad a una máquina real requiere revisión experta.
\end{itemize}

\noindent\textbf{Errores frecuentes}
\begin{itemize}
\item Mantener solidez objetivo tras redondear Z.
\end{itemize}
% studium:end-record eyJpZCI6ICJQUkItMDE0OSIsICJzZWN0aW9uIjogImMxOSJ9
\begin{center}\begin{tikzpicture}\foreach \y in {0,1.8,3.6}{\draw[blue,thick] (0,\y) .. controls (1,\y+.8) and (2.8,\y+.8) .. (4,\y+.3);\draw[blue,thick] (0,\y) .. controls (1,\y+.2) and (2.8,\y+.3) .. (4,\y+.3);}\draw[<->] (-.5,0)--(-.5,1.8) node[midway,left] {$s$};\draw[<->] (0,4.6)--(4,4.9) node[midway,above] {$c$};\node[align=center] at (6,2) {$Z_b$ entero\\$V_e,V_s,|\Delta V_\theta|$};\node[below,align=center] at (3.5,-0.7) {$s=4.0768\times10^{-2}\,\mathrm{m},\quad \sigma=1.3246\times10^{0}\,$};\node[below,align=center] at (3.5,-1.2999999999999998) {$D=3.7460\times10^{-1}\,$};\end{tikzpicture}\end{center}
\subsection*{C19-06-P}

\noindent\textbf{Enunciado}
\(C19-06-P \cdot B3 \cdot\) Balance térmico integrado. Una pared se refrigera por convección interior. Calcule flujo por área, temperatura caliente y caudal mínimo energético para el incremento permitido. Resultado final: caudal. Datos: \(T_{aw}=1580\, \mathrm{K}\); \(T_c=720\, \mathrm{K}\); \(h_g=1540\, \mathrm{W}\,\mathrm{m}^{-2}\,\mathrm{K}^{-1}\); \(h_c=2320\, \mathrm{W}\,\mathrm{m}^{-2}\,\mathrm{K}^{-1}\); \(t=0.0012\, \mathrm{m}\); \(k_m=23\, \mathrm{W}\,\mathrm{m}^{-1}\,\mathrm{K}^{-1}\); \(A=0.028\, \mathrm{m}^{2}\); \(c_p=1080\, \mathrm{J}\,\mathrm{kg}^{-1}\,\mathrm{K}^{-1}\); \(\Delta T_c=165\, \mathrm{K}\).

\noindent\textbf{Datos}
\begin{itemize}
\item $T_{aw} = 1580\,\mathrm{K}$ — temperatura recuperación exterior
\item $T_c = 720\,\mathrm{K}$ — temperatura refrigerante local media del modelo
\item $h_g = 1540\,\mathrm{W}\,\mathrm{m}^{-2}\,\mathrm{K}^{-1}$ — coeficiente exterior
\item $h_c = 2320\,\mathrm{W}\,\mathrm{m}^{-2}\,\mathrm{K}^{-1}$ — coeficiente interior
\item $t = 0.0012\,\mathrm{m}$ — espesor metal
\item $k_m = 23\,\mathrm{W}\,\mathrm{m}^{-1}\,\mathrm{K}^{-1}$ — conductividad
\item $A = 0.028\,\mathrm{m}^{2}$ — superficie intercambio
\item $c_p = 1080\,\mathrm{J}\,\mathrm{kg}^{-1}\,\mathrm{K}^{-1}$ — calor específico refrigerante
\item $\Delta T_c = 165\,\mathrm{K}$ — incremento permitido
\end{itemize}

\noindent\textbf{Incógnita}
$\dot{m}_c$, en $\mathrm{kg\,s^{-1}}$; también las magnitudes intermedias desarrolladas.

\noindent\textbf{Modelo físico}

\noindent\textit{B3 \(\cdot\) Balance térmico integrado}

\[
\fitmath{\dot{m}_c=\frac{A \left(- T_c + T_{aw}\right)}{c_p \Delta T_c \left(\frac{1}{h_g} + \frac{1}{h_c} + \frac{t}{k_m}\right)}}
\]

\noindent\textbf{Hipótesis}
\begin{itemize}
\item Red estacionaria unidimensional, propiedades constantes, sin radiación.
\end{itemize}

\noindent\textbf{Desarrollo paso a paso}

\begin{enumerate}

\item Se suma la resistencia térmica por área.

\[\fitmath{\mathrm{Res}=\frac{1}{\mathrm{hg}}+\frac{\mathrm{th}}{\mathrm{cond}}+\frac{1}{\mathrm{hc}}}\]

\[\fitmath{\mathrm{Res}=\frac{1}{(1540)}+\frac{(0.0012)}{(23)}+\frac{1}{(2320)}\approx 0.001132559\,\mathrm{m^{2}\,K/W}}\]

\item Se calcula el flujo estacionario.

\[\fitmath{\mathrm{flux}=\frac{\mathrm{Tg}-\mathrm{Tc}}{\frac{1}{\mathrm{hg}}+\frac{\mathrm{th}}{\mathrm{cond}}+\frac{1}{\mathrm{hc}}}}\]

\[\fitmath{\mathrm{flux}=\frac{(1580)-(720)}{\frac{1}{(1540)}+\frac{(0.0012)}{(23)}+\frac{1}{(2320)}}\approx 759342.3\,\mathrm{W/m^{2}}}\]

\item La cara caliente sigue del salto convectivo exterior.

\[\fitmath{\mathrm{Tw}=\mathrm{Tg}-\frac{\mathrm{Tg}-\mathrm{Tc}}{\mathrm{hg}\cdot \left(\frac{1}{\mathrm{hg}}+\frac{\mathrm{th}}{\mathrm{cond}}+\frac{1}{\mathrm{hc}}\right)}}\]

\[\fitmath{\mathrm{Tw}=(1580)-\frac{(1580)-(720)}{(1540)\cdot \left(\frac{1}{(1540)}+\frac{(0.0012)}{(23)}+\frac{1}{(2320)}\right)}\approx 1086.921\,\mathrm{K}}\]

\item La energía absorbida determina el mínimo de caudal.

\[\fitmath{\mathrm{mc}=\frac{\mathrm{area}\cdot \left(\mathrm{Tg}-\mathrm{Tc}\right)}{\left(\frac{1}{\mathrm{hg}}+\frac{\mathrm{th}}{\mathrm{cond}}+\frac{1}{\mathrm{hc}}\right)\cdot \mathrm{cp}\cdot \mathrm{dt}}}\]

\[\fitmath{\mathrm{mc}=\frac{(0.028)\cdot \left((1580)-(720)\right)}{\left(\frac{1}{(1540)}+\frac{(0.0012)}{(23)}+\frac{1}{(2320)}\right)\cdot (1080)\cdot (165)}\approx 0.119313\,\mathrm{kg/s}}\]

\end{enumerate}


\noindent\textbf{Sustitución numérica}

\[
\fitmath{\frac{\left(0.028\,\mathrm{m}^{2}\right) \left(- \left(720\,\mathrm{K}\right) + \left(1580\,\mathrm{K}\right)\right)}{\left(1080\,\mathrm{J}\,\mathrm{kg}^{-1}\,\mathrm{K}^{-1}\right) \left(165\,\mathrm{K}\right) \left(\frac{1}{\left(1540\,\mathrm{W}\,\mathrm{m}^{-2}\,\mathrm{K}^{-1}\right)} + \frac{1}{\left(2320\,\mathrm{W}\,\mathrm{m}^{-2}\,\mathrm{K}^{-1}\right)} + \frac{\left(0.0012\,\mathrm{m}\right)}{\left(23\,\mathrm{W}\,\mathrm{m}^{-1}\,\mathrm{K}^{-1}\right)}\right)}}
\]

\noindent\textbf{Resultado}

\[
\fitmath{\dot{m}_c=1.193130451\times10^{-1}\,\mathrm{kg}\,\mathrm{s}^{-1}}
\]

\noindent\textbf{Interpretación}
La temperatura Tc es un valor medio local prescrito para esta red. El mínimo energético no garantiza coeficiente interior ni presión de suministro.

\noindent\textbf{Comprobación de los pasos}
\begin{itemize}
\item Paso 2: $0.00113256$ — reproducido
\item Paso 4: $759342$ — reproducido
\item Paso 6: $1086.92$ — reproducido
\item Paso 8: $0.119313$ — reproducido
\end{itemize}

\noindent\textbf{Comprobación de la respuesta} la aritmética del último paso reproduce la respuesta enunciada a la precisión con que está escrita. Es una reproducción aritmética; no valida el modelo físico ni sus hipótesis.

\noindent\textbf{Limitaciones}
\begin{itemize}
\item Modelo educativo; la aplicabilidad a una máquina real requiere revisión experta.
\end{itemize}

\noindent\textbf{Errores frecuentes}
\begin{itemize}
\item Tratar el balance energético como diseño de refrigeración seguro.
\end{itemize}
% studium:end-record eyJpZCI6ICJQUkItMDE1MCIsICJzZWN0aW9uIjogImMxOSJ9
\begin{center}\begin{tikzpicture}\node (a) at (0,1) {$T_{aw}$};\node[draw,minimum height=1cm] (b) at (1.8,1) {$1/h_g$};\node[draw,minimum height=1cm] (c) at (3.8,1) {$t/k_m$};\node[draw,minimum height=1cm] (d) at (5.8,1) {$1/h_c$};\node (e) at (7.3,1) {$T_c$};\draw[->] (a)--(b);\draw[->] (b)--(c);\draw[->] (c)--(d);\draw[->] (d)--(e);\node at (2.8,1.8) {$T_{w,g}$};\draw[->,red,thick] (1,2.5)--(6.5,2.5) node[midway,above] {$q^{\prime\prime}$};\node[below,align=center] at (3.5,-0.7) {$R_{\mathrm{th}}=1.1326\times10^{-3}\,\mathrm{m}^{2}\,\mathrm{K}\,\mathrm{W}^{-1},\quad q''=7.5934\times10^{5}\,\mathrm{W}\,\mathrm{m}^{-2}$};\node[below,align=center] at (3.5,-1.2999999999999998) {$T_w=1.0869\times10^{3}\,\mathrm{K},\quad \dot{m}_c=1.1931\times10^{-1}\,\mathrm{kg}\,\mathrm{s}^{-1}$};\end{tikzpicture}\end{center}
\chapter{Verificación y revisión pendiente}
Esta edición conserva un expediente reproducible separado con el inventario de fuentes, las llamadas reales a Studium, los controles de problemas y la auditoría integrada. El estado de la auditoría es de revisión requerida; el libro no constituye una certificación científica.

Se han reproducido los pasos de los 150 problemas. Las evaluaciones con unidades y dominio incluyen un ejercicio de diagnóstico con rendimiento calculado superior a uno: su solución identifica datos incompatibles con el modelo, no acepta esa prestación como admisible. Se han comprobado 15 identidades, derivadas e integrales y ocho residuos seleccionados de conservación. Las 18 deducciones registradas de modelos de ejemplo tienen comprobación dimensional, con especialización explícita de parámetros adimensionales donde lo exige el verificador.

Siguen pendientes la revisión experta de las hipótesis físicas y correlaciones, la evaluación aerodinámica de los perfiles educativos y la comprobación de vida de materiales. El modelo lineal de estabilidad no predice la línea de bombeo de un motor. Los esquemas de mapas no son datos experimentales. Los modelos térmicos no constituyen un diseño seguro de refrigeración.

La fundamentación formal automática de todas las ecuaciones del cuerpo aún es incompleta. El informe de ecuaciones conserva cada hallazgo por párrafo. La diversidad de fuentes citadas no se interpreta como validación semántica de cada cifra. Una comprobación de unidades no demuestra la validez de una ley física ni que sus hipótesis se cumplan.

El código del framework no se ha modificado. Se ha fijado la revisión 08ef1f42892d (hash completo en el expediente), identificada como versión alfa 1.0.0a1. Antes de redactar se ejecutaron 85 pruebas básicas, todas satisfactorias. Los cambios editoriales de esta publicación se aplican mediante un programa externo sobre el LaTeX generado y se documentan en el informe editorial.
\chapter{Plan de estudio}
\begin{longtable}{p{.15\linewidth}p{.32\linewidth}p{.40\linewidth}}
\textbf{Tramo} & \textbf{Lectura} & \textbf{Criterio para avanzar}\\\hline
1 & Capítulos 1--3 & Distinguir estados, referencias y rendimientos; cerrar Euler y unidades.\\
2 & Capítulos 4--6 & Diseñar triángulos de turbina y comprobar reacción y capacidad de gasto.\\
3 & Capítulos 7--10 & Cerrar escalón compresor, difusión, anillo y condiciones de arranque.\\
4 & Capítulos 11--14 & Integrar equilibrio radial, torsión, tensiones y balances térmicos.\\
5 & Capítulos 15--17 & Mantener radios distintos, definición de deslizamiento y límites del difusor.\\
6 & Capítulos 18--19 & Reproducir las cuatro configuraciones y resolver ambos simulacros sin consultar respuestas.\\
\end{longtable}
Al revisar un ejercicio, anote por separado errores de convenio, de modelo, de unidades y de cálculo. Vuelva a resolver el caso de transferencia de la misma familia y explique la tendencia antes de operar. Si un modelo requiere un dato que no se proporciona, identifique la falta de información en lugar de inventar un coeficiente.
\backmatter
\begin{thebibliography}{99}
\bibitem{guia} Universidad de León. \emph{Guía docente 2026/2027, 0710340: Turbomáquinas}. Documento de cinco páginas aportado por el usuario. Temario p. 2, metodología y evaluación pp. 3--4. Lectura del documento; no se afirma una validación criptográfica de su firma.
\bibitem{ule} Universidad de León. \emph{Material docente de Turbomáquinas}: presentación 2026/2027; enunciado del trabajo de hoja de cálculo; Generalidades; Diseño de escalón de compresor axial; Turbinas axiales; Esfuerzos y refrigeración; Compresores radiales. Siete PDF aportados, 144 páginas. Las referencias a diapositivas usan la página del PDF.
\bibitem{mit24} Martinez-Sanchez, M. \emph{16.50 Introduction to Propulsion Systems}, Spring 2012, lecciones 24, 26, 27, 28, 29, 30 y 32. MIT OpenCourseWare. Notas del curso basadas en el trabajo docente de Jack L. Kerrebrock. Consultadas para contraste de principios y modelos; imágenes de terceros excluidas. \url{https://ocw.mit.edu/courses/16-50-introduction-to-propulsion-systems-spring-2012/pages/lecture-notes/}. Licencia y condiciones: \url{https://ocw.mit.edu/terms/}.
\bibitem{hill} Hill, P. G.; Peterson, C. R. \emph{Mechanics and Thermodynamics of Propulsion}, 2.ª ed. La guía 2026/2027 identifica Pearson India Education, 2010, ISBN 978-81-317-2951-9. El catálogo editorial confirma título, autores, edición e ISBN; su lista de 2020 indica 2009 como año de esa referencia. Se conserva la fecha de la guía con esta discrepancia pendiente de comprobar en el ejemplar. La edición internacional de 1992, Addison-Wesley, ISBN 978-0-201-14659-2, es la citada en las diapositivas. Metadatos contrastados; libros no consultados íntegramente. \url{https://www.pearson.com/en-us/subject-catalog/p/mechanics-and-thermodynamics-of-propulsion/P200000003466/9780201146592}.
\bibitem{mattingly} Mattingly, J. D.; Heiser, W. H.; Pratt, D. T. \emph{Aircraft Engine Design}, 2.ª ed. AIAA, Reston, 2002. ISBN 978-1-56347-538-2. DOI: \url{https://doi.org/10.2514/4.861444}. Bibliografía recomendada por la guía; identidad contrastada con el registro editorial en Crossref. Texto comercial no consultado íntegramente.
\bibitem{oates} Oates, G. C. (ed.). \emph{Aerothermodynamics of Aircraft Engine Components}. AIAA, New York, 1985. ISBN 978-0-915928-97-2. DOI: \url{https://doi.org/10.2514/4.861338}. Bibliografía de la guía; metadatos contrastados con Crossref. Oates figura como editor. Texto comercial no consultado íntegramente.
\bibitem{lewis} Lewis, R. I. \emph{Turbomachinery Performance Analysis}. Arnold, London, 1996. ISBN 978-0-340-63191-1, según la guía. Autor, título, editorial y año contrastados en la página del autor: \url{https://www.staff.ncl.ac.uk/r.i.lewis/arnold.htm}. Texto no consultado íntegramente.
\bibitem{powers} \emph{Journal of Propulsion and Power}. AIAA. ISSN electrónico 1533-3876. La entrada sin título asociada a Joseph M. Powers en la guía se identifica por este ISSN; corresponde a una revista, no a un libro de Powers. \url{https://doi.org/10.2514/jpp}. No se atribuye una fórmula a un artículo no seleccionado y leído.
\bibitem{cumpsty} Cumpsty, N. A. \emph{Compressor Aerodynamics}. Reimpresión de Krieger, 2004, de la edición Longman de 1989. ISBN 978-1-57524-247-7. Referencia auxiliar de la presentación. Metadatos y disponibilidad abierta contrastados en Imperial College London: \url{https://hdl.handle.net/10044/1/128140}. Se consultó el registro bibliográfico, no el texto completo; no se usa como testigo de una fórmula concreta.
\bibitem{rr} Rolls-Royce. \emph{The Jet Engine}, 5.ª ed. Wiley, 2015. ISBN 978-1-119-06599-9. Referencia auxiliar de la presentación. Datos contrastados en el catálogo editorial: \url{https://www.wiley-vch.de/en/areas-interest/engineering/the-jet-engine-978-1-119-06599-9}. Texto no consultado íntegramente.
\bibitem{wuolah} \emph{Exámenes} y \emph{Problemas}, documentos estudiantiles anónimos distribuidos por Wuolah, 13 y 11 páginas. Referencias privadas no oficiales; únicamente análisis de competencias, sin reproducción de contenidos protegidos ni aceptación de resultados manuscritos.
\end{thebibliography}

\end{document}
